\documentclass[11pt,a4paper,oneside]{book}
\usepackage[T1]{fontenc}
\usepackage[utf8]{inputenc}
\usepackage[indonesian]{babel}
\usepackage{lmodern}
\usepackage{amsmath,amsthm,amssymb,amscd}
\usepackage{graphicx}
\usepackage[a4paper,margin=22mm,headheight=15pt]{geometry}
\usepackage{microtype}
\usepackage[hidelinks,pdfusetitle]{hyperref}
\usepackage{bookmark}
% language package supplied by unit wrapper

% UTF-8 input supplied by unit wrapper

%\usepackage[utf8]{inputenc}

\usepackage{ifthen}

%\usepackage{fancyhdr}

%Aus Grundvorspann

%Packages

\usepackage{amscd}

\usepackage{amssymb}

\usepackage{epsfig} %Graphik

\usepackage{verbatim,moreverb} %

\usepackage{enumitem}

\usepackage{eurosym}

\usepackage{epigraph}

\setcounter{MaxMatrixCols}{20}

%Mathematische Einzelbefehle

%[[Projekt:Semantische Vorlagen/Mathematischer Modus/Latex]]

\newcommand{\mathdisp}[2]{\[ #1 \, #2 \]}

\newcommand{\mathdisplayteile}[3]{\[ #1 \] \[ #2 \, #3 \]}

\newcommand{\mathind}[3]{$ #1$, $#2 $#3}

\newcommand{\mathbed}[8]{$ #1 ${\ifthenelse{\equal{#2}{}}{, }{ #2 }}$ #3 ${\ifthenelse{\equal{#5}{}}{}{{\ifthenelse{\equal{#4}{}}{, }{ #4 }}}$ #5 $}{\ifthenelse{\equal{#7}{}}{}{{\ifthenelse{\equal{#6}{}}{, }{ #6 }}}$ #7 $}#8}

\newcommand{\mathbeddisp}[8]{\[ #1 {\ifthenelse{\equal{#2}{}}{,\, }{ \text{ #2 } }} #3 {\ifthenelse{\equal{#5}{}}{}{{\ifthenelse{\equal{#4}{}}{,\, }{ \text{ #4 } }}} #5 }{\ifthenelse{\equal{#7}{}}{}{{\ifthenelse{\equal{#6}{}}{,\, }{ \text{ #6 } }}} #7 }#8 \]}

\newcommand{\mathbedtermdisp}[8]{\[ #1\] #2 $ #3 $#4 $ #5 $#6$ #7 $#8}

\newcommand{\mathkon}[4]{$ #1 $ #2 $ #3 $#4}

\newcommand{\mathkor}[5]{{\ifthenelse{\equal{#1}{}}{}{#1 }}$ #2 $ #3 $ #4 $#5}

\newcommand{\mathkork}[5]{{\ifthenelse{\equal{#1}{}}{}{#1 }}$ #2 $ #3 $ #4 $#5}

\newcommand{\mathlist}[5]{$ #1 $\ifthenelse{\equal{#2}{}}{,}{ #2} $ #3 $\ifthenelse{\equal{#4}{}}{,}{ #4} $ #5 $}

\newcommand{\mathlistdisp}[5]{\[  #1  \ifthenelse{\equal{#2}{}}{,}{ \text{ #2 }}  #3  \ifthenelse{\equal{#4}{}}{,}{\text{ #4 } }  #5 \] }

\newcommand{\alignier}[2]{\begin{align*} #1 \, #2 \end{align*}}

\newcommand{\mathlistvierdisp}[7]{\[ #1,\, #3, \, #5 \text{ #6 } #7  \]}

%[[Projekt:Semantische Vorlagen/Mathematische Vergleichsvorlagen/Latex]]

\newcommand{\mavergleichskette}[4]{$ #1 #2 #3#4$}

\newcommand{\mavergleichskettedisp}[4]{\[  #1 #2 #3 #4 \]}

\newcommand{\mavergleichskettedisplang}[4]{\[  #1 #2 #3 #4 \]}

\newcommand{\mavergleichskettealign}[4]{ \begin{eqnarray*}  #1 #2 #3 #4 \end{eqnarray*}  }

\newcommand{\vergleichskette}[9]{ #1 \, #2 \, #3 \ifthenelse
{\equal {#4}{ }} {} {\, #4 \, #5 }  \ifthenelse {\equal {#6}{ } }{}
{\, #6 \, #7 } \ifthenelse {\equal {#8}{ }}{} {\, #8 \, #9 }  }

\newcommand{\vergleichskettelang}[9]{ #1 \, #2 \, #3 \ifthenelse
{\equal {#4}{ }} {} {\, #4 \, #5 }  \ifthenelse {\equal {#6}{ } }{}
{\, #6 \, #7 } \ifthenelse {\equal {#8}{ }}{} {\, #8 \, #9 }  }

\newcommand{\vergleichskettefortsetzung}[8]{ \, #1 \, #2 \ifthenelse
{\equal {#3}{ }} {} {\, #3 \, #4 }  \ifthenelse {\equal {#5}{ } }{}
{\, #5 \, #6 } \ifthenelse {\equal {#7}{ }}{} {\, #7 \, #8 } }

\newcommand{\vergleichskettealign}[9]{ #1 & #2 & #3 \ifthenelse
{\equal {#4}{ }} {} {\cr & #4 & #5 }  \ifthenelse {\equal {#6}{ }
}{} {\cr & #6 & #7 } \ifthenelse {\equal {#8}{ }}{} {\cr & #8 & #9 }
}

\newcommand{\vergleichskettefortsetzungalign}[8]{ \cr & #1 & #2 \ifthenelse
{\equal {#3}{ }} {} {\cr & #3 & #4 }  \ifthenelse {\equal {#5}{ }
}{} {\cr & #5 & #6 } \ifthenelse {\equal {#7}{ }}{} {\cr & #7 & #8 }
}

\newcommand{\mavergleichskettealigndrucklinks}[4]{ \begin{eqnarray*}  #1 #2 #3 #4 \end{eqnarray*}  }

\newcommand{\vergleichskettealigndrucklinks}[9]{  &  & #1 \cr & #2 & #3 \ifthenelse
{\equal {#4}{ }} {} {\cr & #4 & #5 }  \ifthenelse {\equal {#6}{ }
}{} {\cr & #6 & #7 } \ifthenelse {\equal {#8}{ }}{} {\cr & #8 & #9 }
}

\newcommand{\mavergleichskettealignhandlinks}[4]{ \begin{eqnarray*}  #1 #2 #3 #4 \end{eqnarray*}  }

\newcommand{\vergleichskettealignhandlinks}[9]{ #1 & #2 & #3 \ifthenelse
{\equal {#4}{ }} {} {\cr & #4 & #5 }  \ifthenelse {\equal {#6}{ }
}{} {\cr & #6 & #7 } \ifthenelse {\equal {#8}{ }}{} {\cr & #8 & #9 }
}

\newcommand{\mavergleichskettedisphandlinks}[4]{ \[ #1 #2 #3 #4 \] }

\newcommand{\vergleichskettedisphandlinks}[9]{  #1 \, #2 \, #3 \ifthenelse
{\equal {#4}{ }} {} {\, #4 \, #5 }  \ifthenelse {\equal {#6}{ } }{}
{\, #6 \, #7 } \ifthenelse {\equal {#8}{ }}{} {\, #8 \, #9 }  }

\newcommand{\zeilemitteil}[2]{ \ifthenelse {\equal {#2} {} } {#1} { #1 \cr & & #2}   }

\newcommand{\mavergleichskettek}[4]{$ #1 #2 #3#4$}

\newcommand{\vergleichskettek}[9]{ #1 \, #2 \, #3 \ifthenelse
{\equal {#4}{ }} {} {\, #4 \, #5 }  \ifthenelse {\equal {#6}{ } }{}
{\, #6 \, #7 } \ifthenelse {\equal {#8}{ }}{} {\, #8 \, #9 }  }

\newcommand{\vergleichskettefortsetzungk}[8]{ \, #1 \, #2 \ifthenelse
{\equal {#3}{ }} {} {\, #3 \, #4 }  \ifthenelse {\equal {#5}{ } }{}
{\, #5 \, #6 } \ifthenelse {\equal {#7}{ }}{} {\, #7 \, #8 } }

%[[Projekt:Semantische Vorlagen/Mathematische Strukturvorlagen/Latex]]

\newcommand{\N}   {{\mathbb N}}
\newcommand{\Z}   {{\mathbb Z}}
\newcommand{\Q}   {{\mathbb Q}}
\newcommand{\R}   {{\mathbb R}}
\newcommand{\C}   {{\mathbb C}}
\newcommand{\Complex}  {{\mathbb C}}

\newcommand{\defeq}{:=}

\newcommand{\defeqr}{=:}

\newcommand{\maabb}[4]{${\ifthenelse {\equal {#1}{} } {} {#1 \colon }}\,  #2 \rightarrow #3 #4 $}

\newcommand{\maabbele}[6]{$ {\ifthenelse {\equal {#1}{} } {} {#1 \colon }}\,   #2 \rightarrow  #3 , \, #4 \mapsto #5 #6 $}

\newcommand{\maabbdisp}[4]{\[ {\ifthenelse {\equal {#1}{} } {} {#1 \colon }}\,  #2 \longrightarrow #3 #4 \]}

\newcommand{\maabbeledisp}[6]{\[ {\ifthenelse {\equal {#1}{} } {} {#1 \colon }}\,   #2 \longrightarrow  #3 , \, #4 \longmapsto #5 #6 \]}

%Variante bei ifequal-Problem

\newcommand{\maabbnamedisp}[4]{\[ #1 \colon \, #2 \longrightarrow #3 #4 \]}

\newcommand{\maabbeledispvar}[6]{\[ {\ifthenelse {\equal {#1}{} } {} {#1 \colon }}\,   #2 \longrightarrow  #3 , \, #4 \longmapsto #5 #6  \]}

\newcommand{\maabbelementzeiledisplay}[6]{\[ {\ifthenelse {\equal {#1}{} } {} {#1 \colon }}\, #2 \longrightarrow #3 , \]  	\[ #4 \longmapsto #5 #6 \]}

\newcommand{\maabbelementdoppelzeiledisplay}[6]{\[ {\ifthenelse {\equal {#1}{} } {} {#1 \colon }}\, #2 \longrightarrow #3 , \] \[ #4 \longmapsto \] \[  #5 #6 \]}

\newcommand{ \arccot  } {\operatorname{arccot } }

\newcommand{\betrag}[1]{\left| #1 \right|}

\newcommand{\strukturgarbe}{ {\mathcal O} }

\newcommand{\schnittring}[1]{\Gamma (#1 , {\mathcal O})}

\newcommand{\tensor}{\otimes}

%[[Projekt:Semantische Vorlagen/Textgestaltung/Latex]]

\newcommand{\einrueckung}[1]{\par \hspace{1cm}#1 \par}

\newcommand{\zusatz}[2]{(#1)#2}

\newcommand{\zusatzklammer}[3]{(#1#2)#3}

\newcommand{\zusatzgs}[3]{\-- #1 \-- #3}

\newcommand{\zusatzfussnote}[3]{#3\footnote{#1#2}}

%Feinabstimmung

\newcommand{\latexdruckeinskurz}{\!}

\newcommand{\latexdruckzweikurz}{\! \!}

\newcommand{\latexdruckdreikurz}{\! \! \!}

%[[Projekt:Semantische Vorlagen/Satzumgebungen/Latex]]

\newtheorem{fakt}{Fakta}[section] %so wird abschnittsweise %durchnummeriert. Bei [] wird % einfach nummeriert.

\newtheorem{Fakt}{fakt}[Fakt]

\newtheorem{Proposition}[fakt]{Proposisi}

\newtheorem{Korollar}[fakt]{Korolari}

\newtheorem{Lemma}[fakt]{Lema}

\newtheorem{Satz}[fakt]{Teorema}

\newtheorem{Theorem}[fakt]{Teorema}

\theoremstyle{definition}

\newtheorem{Axiom}[fakt]{Aksioma}

\newtheorem{aufgabe}[fakt]{Soal}

\newtheorem{Aufgabe}[fakt]{Soal}

\newtheorem{klausuraufgabe}{Soal ujian}

\newtheorem{Klausuraufgabe}{Soal ujian}

\newtheorem{beispiel}[fakt]{Contoh}

\newtheorem{Beispiel}[fakt]{Contoh}

\newtheorem{Konstruktion}[fakt]{Konstruksi}

\newtheorem{konstruktion}[fakt]{Konstruksi}

\newtheorem{Verfahren}[fakt]{Verfahren}

\newtheorem{verfahren}[fakt]{Verfahren}

\newtheorem{bemerkung}[fakt]{Catatan}

\newtheorem{Bemerkung}[fakt]{Catatan}

\newtheorem{definition}[fakt]{Definisi}

\newtheorem{Definition}[fakt]{Definisi}

\newtheorem{notation}[fakt]{Notation}

\newtheorem{Notation}[fakt]{Notation}

\newtheorem{Frage}[fakt]{Frage}

\newtheorem{frage}[fakt]{Frage}

\newtheorem{problem}[fakt]{Problem}

\newtheorem{Problem}[fakt]{Problem}

\newtheorem{situation}[fakt]{Situation}

\newtheorem{Situation}[fakt]{Situation}

%Einlesegestaltung

%[[Projekt:Semantische Vorlagen/Umgebungsdesign/Latex]]

%Aufgabengestaltung

\newcommand{\inputaufgabe}[4]{\aufgabevorskip \begin{aufgabe} \punkte {#1} \aufgabepunktskip #2 #3 #4\end{aufgabe} \aufgabenachskip}

\newcommand{\inputaufgabegibtloesung}[4]{\aufgabevorskip \begin{aufgabe}\hspace{-1.8 mm}* \punkte {#1} \aufgabepunktskip #2 #3 #4\end{aufgabe} \aufgabenachskip}

\newcommand{\inputaufgabepunktenummervonhand}[3]{\aufgabevorskip {\bf Aufgabe #1} (#2 Punkte) \aufgabepunktskip #3  \aufgabenachskip}

\newcommand{\inputaufgabeloesung}[2]{ \begin{aufgabe} #1 \end{aufgabe} \aufgabenachskip

\bigskip

Lösung

\bigskip #2

}

\newcommand{\inputaufgabeklausurloesung}[3]{\par \newpage \begin{aufgabe} {\ifthenelse {\equal {#1}{}}{} {\ifthenelse {\equal {#1}{1}} {(1 Punkt)} {(#1 Punkte)}}} \par \bigskip #2 \end{aufgabe} \underline{L\"osung:} \par \bigskip #3 }

\newcommand{\inputaufgabeloesungvar}[2]{\aufgabevorskip Aufgabe: #1 \aufgabenachskip
Lösung: #2}

\newcommand{\inputaufgabepunkteloesung}[3]{\aufgabevorskip \begin{aufgabe} {\ifthenelse {\equal {#1}{}}{} {\ifthenelse {\equal {#1}{1}} {(1 Punkt)} {(#1 Punkte)}}}  \aufgabepunktskip #2 \end{aufgabe} \aufgabenachskip \loesung #3}

\newcommand{\loesung}[1]{\underline{Lösung:} #1}

\newcommand{\punkte}[1]{\ifthenelse {\equal {#1}{}}{} {\ifthenelse {\equal {#1}{1}} {(1 poin)} {(#1 poin)}}}

\newcommand{\inputexercise}[4]{\aufgabevorskip \begin{exercise} \points {#1} \aufgabepunktskip #2 #3 #4\end{exercise} \aufgabenachskip}

\newcommand{\points}[1]{\ifthenelse {\equal {#1}{}}{} {\ifthenelse {\equal {#1}{1}} {(1 point)} {(#1 points)}}}

%Beispielgestaltung

\newcommand{\inputbeispiel}[2]
{\beispielvorskip
\begin{beispiel}
\beispielbenennung{#1}
\beispielbenennungnachskip #2
\end{beispiel}
\beispielnachskip}

\newcommand{\beispielbenennung}[1]{ \ifthenelse {\equal {#1}{} \OR \equal {#1}{   }  }{}{(#1)}     }

%Bemerkunggestaltung

\newcommand{\inputbemerkung}[2]
{\bemerkungvorskip
\begin{bemerkung}
\bemerkungbenennung{#1}
\bemerkungbenennungnachskip #2
\end{bemerkung}
\bemerkungnachskip}

\newcommand{\bemerkungbenennung}[1]{  \ifthenelse {\equal {#1}{} \OR \equal {#1}{   }  }{}{(#1)}}

\newcommand{\inputverfahren}[2]
{\bemerkungvorskip
\begin{verfahren}
\bemerkungbenennung{#1}
\bemerkungbenennungnachskip #2
\end{verfahren}
\bemerkungnachskip}

\newcommand{\inputkonstruktion}[2]
{\bemerkungvorskip
\begin{konstruktion}
\bemerkungbenennung{#1}
\bemerkungbenennungnachskip #2
\end{konstruktion}
\bemerkungnachskip}

\newcommand{\inputfrage}[2]
{\bemerkungvorskip
\begin{frage}
\problembenennung{#1}
\bemerkungbenennungnachskip #2
\end{frage}
\bemerkungnachskip}

\newcommand{\fragebenennung}[1]{  \ifthenelse {\equal {#1}{} \OR \equal {#1}{   }  }{}{(#1)}}

\newcommand{\inputproblem}[2]
{\bemerkungvorskip
\begin{problem}
\problembenennung{#1}
\bemerkungbenennungnachskip #2
\end{problem}
\bemerkungnachskip}

\newcommand{\problembenennung}[1]{  \ifthenelse {\equal {#1}{} \OR \equal {#1}{   }  }{}{(#1)}}

%Situationsgestaltung

\newcommand{\inputsituation}[2]
{\begin{situation}
\situationbenennung{#1} #2
\end{situation}
}

\newcommand{\situationbenennung}[1]{\ifthenelse {\equal {#1}{} \OR \equal {#1}{  }  }{}{(#1)}}

%Definitiongestaltung

\newcommand{\inputdefinition}[2]
{\definitionvorskip
\begin{definition}
\definitionbenennung{#1}
\definitionbenennungnachskip #2
\end{definition}
\definitionnachskip}

\newcommand{\inputaxiom}[2]
{\definitionvorskip
\begin{Axiom}
\definitionbenennung{#1}
\definitionbenennungnachskip #2
\end{Axiom}
\definitionnachskip}

\newcommand{\definitionbenennung}[1]{\ifthenelse {\equal {#1}{} \OR \equal {#1}{  }  }{}{(#1)}}

\newcommand{\inputnotation}[2]
{\definitionvorskip \begin{notation} \definitionbenennung{#1} \definitionbenennungnachskip #2 \end{notation} \definitionnachskip}

%Fakt- und Beweisgestaltung

\renewcommand{\proofname}{\hspace{0cm}{\it Bukti}}

\newcommand{\beweis}[1]{\begin{proof} #1 \end{proof}}

\newcommand{\inputfakt}[4]
{\faktvorskip
\begin{#2}

%\label{#1}   (Absetzen, sonst  kann das folgende dahinter rutschen)

\faktbenennung{#3}
\faktbenennungnachskip #4
\end{#2}
\faktnachskip
}

\newcommand{\inputfaktbeweis}[5]
{\faktvorskip
\begin{#2}

%\label{#1}   (Absetzen, sonst  kann das folgende dahinter rutschen)

\faktbenennung{#3}
\faktbenennungnachskip #4
\end{#2}
\faktnachskip
\beweis{#5}}

\newcommand{\inputfaktbeweisnichtvorgefuehrt}[5]{\inputfaktbeweis {#1} {#2} {#3} {#4} {Dieser Beweis wurde in der Vorlesung nicht vorgef\"uhrt.} }

\newcommand{\inputfaktbeweistrivial}[4]{\inputfaktbeweis {#1} {#2} {#3} {#4} {Das ist trivial.} }

\newcommand{\inputfaktuebergangbeweis}[6]
{\faktvorskip
\begin{#2}

%\label{#1} (Absetzen, sonst kann das
% folgende dahinter rutschen)

\faktbenennung{#3}
\faktbenennungnachskip #4
\end{#2}
\faktnachskip #5 \faktnachskip
\beweis{#6}}

\newcommand{\inputbeweis}[1]{\beweis{#1} }

\newcommand{\faktbenennung}[1]{\ifthenelse {\equal {#1}{} \OR \equal {#1}{  }  }{}{(#1)}}

%Gestaltung der Faktstruktur

\newcommand{\faktsituation}[1]{\faktsituationskip #1}

\newcommand{\faktvoraussetzung}[1]{\faktvoraussetzungskip #1}

\newcommand{\faktvoraussetzungleer}[1]{ \hspace{-0,15cm} }

\newcommand{\faktvoraussetzungpos}[1]{\faktvoraussetzungskip #1}

\newcommand{\faktuebergang}[1]{\faktuebergangskip #1}

\newcommand{\faktuebergangleer}[1]{\faktuebergangskip \hspace{-0,15cm} }

\newcommand{\faktfolgerung}[1]{\faktfolgerungskip #1}

\newcommand{\faktzusatz}[1]{\faktzusatzskip #1}

%[[Projekt:Semantische Vorlagen/Beweisgestaltung/Latex]]

\newcommand{\teilbeweis}[5]{#1#2#3#4#5}

\newcommand{\fallunterscheidung}[5]{#1 #2\ifthenelse {\equal {#3}{}} {} { #3}\ifthenelse {\equal {#4}{}} {} { #4}\ifthenelse {\equal {#5}{}} {} { #5}}

\newcommand{\fallunterscheidungzwei}[2]{#1 #2}

\newcommand{\fallunterscheidungdrei}[3]{#1 #2 #3 }

\newcommand{\fallunterscheidungvier}[4]{#1 #2 #3 #4 }

\newcommand{\fallunterscheidungfuenf}[5]{#1 #2 #3 #4 #5 }

%[[Projekt:Semantische Vorlagen/Stichwortumsetzung/Latex]]

\newcommand{\betonung}[2]{\emph{#1}#2}

\newcommand{\stichwort}[2]{\emph{#1}#2}

\newcommand{\stichwortpraemath}[3]{$#1$-\emph{#2}#3}

\newcommand{\definitionswort}[2]{\emph{#1}#2}

\newcommand{\definitionswortpraemath}[3]{$#1$-\emph{#2}#3}

\newcommand{\definitionswortteil}[2]{\emph{#1}#2}

\newcommand{\definitionswortenp}[2]{\emph{#1}#2}

\newcommand{\definitionsverweis}[3]{#1#3}

\newcommand{\definitionsverweisanfuehrung}[3]{\anfuehrung {#1} {#3} }

%Anführungszeichen

\newcommand{\anfuehrung}[2]{"`{#1}"'#2}

\newcommand{\anfuehrungenglisch}[2]{``{#1}''#2}

%[[Projekt:Semantische Auflistungsvorlagen/Latex]]

%Aufzählungen

\newcommand{\aufzaehlungeins}[1]{\begin{enumerate} \item  #1 \end{enumerate}}

\newcommand{\aufzaehlungzwei}[2]{\begin{enumerate} \item  #1 \item  #2 \end{enumerate}}

\newcommand{\aufzaehlungdrei}[3]{\begin{enumerate} \item  #1 \item  #2 \item  #3  \end{enumerate}}

\newcommand{\aufzaehlungvier}[4]{\begin{enumerate} \item  #1 \item  #2 \item  #3 \item  #4 \end{enumerate}}

\newcommand{\aufzaehlungfuenf}[5]{\begin{enumerate} \item  #1 \item  #2 \item  #3 \item  #4 \item  #5  \end{enumerate}}

\newcommand{\aufzaehlungsechs}[6]{\begin{enumerate} \item  #1 \item  #2 \item  #3 \item  #4 \item  #5 \item  #6 \end{enumerate}}

\newcommand{\aufzaehlungsieben}[7]{\begin{enumerate} \item  #1 \item  #2 \item  #3 \item  #4 \item  #5 \item #6 \item  #7 \end{enumerate}}

\newcommand{\aufzaehlungacht}[8]{\begin{enumerate} \item  #1 \item  #2 \item  #3 \item  #4 \item  #5 \item #6 \item  #7 \item #8 \end{enumerate}}

\newcommand{\aufzaehlungneun}[9]{\begin{enumerate} \item  #1 \item  #2 \item  #3 \item  #4 \item  #5 \item #6 \item  #7  \item  #8  \item  #9 \end{enumerate}}

\newcommand{\itemfuenf}[5]{\item  #1 \item  #2 \item  #3 \item  #4 \item  #5}

\newcommand{\itemsechs}[6]{\item  #1 \item  #2 \item  #3 \item  #4 \item  #5 \item  #6}

\newcommand{\itemsieben}[7]{\item  #1 \item  #2 \item  #3 \item  #4 \item  #5 \item  #6 \item #7}

\newcommand{\aufzaehlungeinsabc}[1]{\begin{enumerate}[label=(\alph*)]  #1 \end{enumerate}}

\newcommand{\aufzaehlungzweiabc}[2]{\begin{enumerate}[label=(\alph*)] \item  #1 \item  #2 \end{enumerate}}

\newcommand{\aufzaehlungdreiabc}[3]{\begin{enumerate}[label=(\alph*)] \item  #1 \item  #2 \item  #3  \end{enumerate}}

\newcommand{\aufzaehlungvierabc}[4]{\begin{enumerate}[label=(\alph*)] \item  #1 \item  #2 \item  #3 \item  #4 \end{enumerate}}

\newcommand{\aufzaehlungfuenfabc}[5]{\begin{enumerate}[label=(\alph*)] \item  #1 \item  #2 \item  #3 \item  #4 \item  #5  \end{enumerate}}

\newcommand{\aufzaehlungsechsabc}[6]{\begin{enumerate}[label=(\alph*)] \item  #1 \item  #2 \item  #3 \item  #4 \item  #5 \item  #6 \end{enumerate}}

\newcommand{\aufzaehlungsiebenabc}[7]{\begin{enumerate}[label=(\alph*)] \item  #1 \item  #2 \item  #3 \item  #4 \item  #5 \item #6 \item  #7 \end{enumerate}}

\newcommand{\aufzaehlungachtabc}[8]{\begin{enumerate}[label=(\alph*)] \item  #1 \item  #2 \item  #3 \item  #4 \item  #5 \item #6 \item  #7 \item #8 \end{enumerate}}

\newcommand{\aufzaehlungneunabc}[9]{\begin{enumerate}[label=(\alph*)] \item  #1 \item  #2 \item  #3 \item  #4 \item  #5 \item #6 \item  #7  \item  #8  \item  #9 \end{enumerate}}

\newcommand{\aufzaehlungzweireihe}[2]{\begin{enumerate} #1 #2 \end{enumerate}}

\newcommand{\aufzaehlungdreibeweis}[3]{(1) #1 \par (2) #2 \par (3) #3 }

\newcommand{\aufzaehlungvierbeweis}[4]{(1) #1 \par (2) #2 \par (3) #3  \par (4) #4}

\newcommand{\aufzaehlungfuenfbeweis}[5]{(1) #1 \par (2) #2 \par (3) #3  \par (4) #4  \par (5) #5  }

%Auflistungen

\newcommand{\listitem}{\par \listskip \listdot}

\newcommand{\auflistungzwei}[2]{\par \listskip \listdot  #1 \par \listskip \listdot  #2 \par \listskip}

\newcommand{\auflistungdrei}[3]{\par \listskip \listdot  #1 \par \listskip \listdot  #2\par \listskip \listdot  #3\par\listskip }

\newcommand{\auflistungvier}[4]{\par \listskip \listdot  #1 \par \listskip \listdot  #2\par \listskip \listdot  #3\par \listskip \listdot  #4  \par \listskip }

\newcommand{\auflistungfuenf}[5]{\par \listskip \listdot  #1 \par \listskip \listdot  #2\par \listskip \listdot  #3\par \listskip \listdot  #4  \par \listskip  \listdot  #5 \par \listskip   }

\newcommand{\auflistungsechs}[6]{\par \listskip \listdot  #1 \par \listskip \listdot  #2\par \listskip \listdot  #3\par \listskip \listdot  #4  \par \listskip  \listdot  #5 \par \listskip \listdot  #6 \par \listskip}

\newcommand{\auflistungsieben}[7]{\par \listskip \listdot  #1 \par \listskip \listdot  #2\par \listskip \listdot  #3\par \listskip \listdot  #4  \par \listskip  \listdot  #5 \par \listskip \listdot  #6 \par \listskip \listdot  #7 \par \listskip }

\newcommand{\auflistungacht}[8]{\par \listskip \listdot  #1 \par \listskip \listdot  #2\par \listskip \listdot  #3\par \listskip \listdot  #4  \par \listskip  \listdot  #5 \par \listskip \listdot  #6 \par \listskip \listdot  #7 \par \listskip \listdot  #8  \par \listskip}

\newcommand{\auflistungneun}[9]{\par \listskip \listdot  #1 \par \listskip \listdot  #2\par \listskip \listdot  #3\par \listskip \listdot  #4  \par \listskip  \listdot  #5 \par \listskip \listdot  #6 \par \listskip \listdot  #7 \par \listskip \listdot  #8  \par \listskip  \listdot  #9 \par \listskip  }

\newcommand{\listitemfuenf}[5]{\listitem  #1 \listitem  #2 \listitem  #3 \listitem  #4 \listitem  #5}

\newcommand{\auflistungzweireihe}[2]{ #1 #2 }

%Gestaltung der Auflistungen

\newcommand{\listdot}{$\bullet \,$}

%Bildgestaltung

\newcommand{\inputbild}[1]{\includegraphics[scale]{#1}}

%Tabellengestaltung

%[[Projekt:Semantische Vorlagen/Zeilenkonfiguration/Latex]]

\newcommand{\zeileundzwei}[2]{#1 & #2}

\newcommand{\zeileunddrei}[3]{#1 & #2 & #3}

\newcommand{\zeileundvier}[4]{#1 & #2 & #3 & #4}

\newcommand{\zeileundfuenf}[5]{#1 & #2 & #3 & #4 & #5}

\newcommand{\zeileundsechs}[6]{#1 & #2 & #3 & #4 & #5 & #6}

\newcommand{\zeileundsieben}[7]{#1 & #2 & #3 & #4 & #5 & #6 & #7}

\newcommand{\zeileundacht}[8]{#1 & #2 & #3 & #4 & #5 & #6 & #7 & #8}

\newcommand{\zeileundneun}[9]{#1 & #2 & #3 & #4 & #5 & #6 & #7 & #8 & #9}

\newcommand{\mazeileundzwei}[2]{$#1$ & $#2$}

\newcommand{\mazeileunddrei}[3]{$#1$ & $ #2$ & $#3$}

\newcommand{\mazeileundvier}[4]{$#1$ & $#2$ & $#3$ & $#4$}

\newcommand{\mazeileundfuenf}[5]{$#1$ & $#2$ & $#3$ & $#4$ & $#5$}

\newcommand{\mazeileundsechs}[6]{$#1$ & $#2$ & $#3$ & $#4$ & $#5$ & $#6$}

\newcommand{\mazeileundsieben}[7]{$#1$ & $#2$ & $#3$ & $#4$ & $#5$ & $#6$ & $#7$}

\newcommand{\mazeileundacht}[8]{$#1$ & $#2$ & $#3$ & $#4$ & $#5$ & $#6$ & $#7$ & $#8$}

\newcommand{\mazeileundneun}[9]{$#1$ & $#2$ & $#3$ & $#4$ & $#5$ & $#6$ & $#7$ & $#8$ & $#9$}

%besser als Listen?
%Liste mit und (&) bzw. Bruch.

\newcommand {\listesechsmaund}[6]{$#1$ & $#2$& $#3$ & $#4$ & $#5$ & $#6$}

\newcommand {\listesiebenmaund}[7]{$#1$ & $#2$& $#3$ & $#4$ & $#5$ & $#6$ & $#7$}

\newcommand{\listezweiund}[2]{#1}

\newcommand{\listedreiund}[3]{#1}

\newcommand{\listesiebenund}[7]{#1}

\newcommand{\listesechsbruch}[6]{#1 \\ \hline #2 \\ \hline #3 \\ \hline #4 \\ \hline #5 \\ \hline #6}

\newcommand{\leitzeileunddrei}[3]{#1}

\newcommand{\madoppelzeileunddrei}[6]{$#1$&$#2$&$#3$\\ \hline $#4$&$#5$&$#6$}

%[[Projekt:Semantische Vorlagen/Tabellen/Latex]]

\newcommand{\tabelleviervier}[4]{\par \bigskip \begin{tabular}{|c|c|c|c|} \hline #1 \\ \hline #2 \\ \hline #3 \\ \hline #4 \\ \hline \end{tabular}}

\newcommand{\tabellefuenfvier}[5]{\par \bigskip \begin{tabular}{|c|c|c|c|c|} \hline #1 \\ \hline #2 \\ \hline #3 \\ \hline #4 \\ \hline #5  \\ \hline \end{tabular}}

%Mit math. Einträgen - ohne Leitzeile

\newcommand{\matabellezweivier}[5]{\begin{center} \begin{tabular}{|c|c|} \hline #2 \\ \hline #3 \\ \hline #4 \\ \hline #5  \\ \hline \end{tabular} \end{center} }

\newcommand{\matabellezweifuenf}[6]{\begin{center} \begin{tabular}{|c|c|} \hline     #2 \\ \hline #3 \\ \hline #4 \\ \hline #5 \\ \hline #6  \\ \hline \end{tabular} \end{center} }

\newcommand{\matabellezweisechs}[7]{\begin{center} \begin{tabular}{|c|c|} \hline     #2 \\ \hline #3 \\ \hline #4 \\ \hline #5 \\ \hline #6 \\ \hline #7   \\ \hline \end{tabular} \end{center} }

\newcommand{\matabellezweisieben}[8]{\begin{center} \begin{tabular}{|c|c|} \hline     #2 \\ \hline #3 \\ \hline #4 \\ \hline #5 \\ \hline #6 \\ \hline #7 \\ \hline #8  \\ \hline \end{tabular} \end{center} }

%Mit math. Einträgen - Mit Leitzeile

\newcommand{\matabellesiebenzwoelf}[3]{\begin{center} \begin{tabular}{|c||c|c|c|c|c|c|c|}  \hline  #1   \\ \hline \hline #2 \\ \hline #3 \\ \hline \end{tabular} \end{center} }

\newcommand{\matabellezwoelfzwoelf}[3]{\begin{center} \begin{tabular}{|c||c|c|c|c|c|c|c|c|c|c|c|c|}  \hline  #1   \\ \hline \hline #2 \\ \hline #3 \\ \hline \end{tabular} \end{center} }

\newcommand{\matabelledreineun}[5]{\begin{center} \begin{tabular}{|c|c|c|}\hline  #1 \\ \hline  #2 \\ \hline #3 \\ \hline   #4 \\ \hline  #5 \\ \hline \end{tabular} \end{center} }

\newcommand{\matabelledreizehn}[6]{\begin{center} \begin{tabular}{|c|c|c|}\hline  #1 \\ \hline  #2 \\ \hline #3 \\ \hline   #4 \\ \hline  #5 \\ \hline  #6 \\ \hline \end{tabular} \end{center} }

\newcommand{\matabelledreielf}[6]{\begin{center} \begin{tabular}{|c|c|c|}\hline  #1 \\ \hline  #2 \\ \hline #3 \\ \hline   #4 \\ \hline  #5 \\ \hline  #6 \\ \hline \end{tabular} \end{center} }

%[[Projekt:Semantische Vorlagen/Klausur/Punktetabellen/Latex]]

\newcommand{\punktetabellezehn}
{\vspace{3ex} \begin{center} \renewcommand{\arraystretch}{1.5} \tabcolsep=0.78ex \begin{tabular}{@{}|c|c|c|c|c|c|c|c|c|c|c|c|c|}
\hline Aufgabe & 1 & 2 & 3 & 4 & 5 & 6 & 7 & 8 & 9 & 10 & $\sum$
\\ \hline m\"ogliche Pkt. & \aeins & \azwei & \adrei & \avier & \afuenf & \asechs   & \asieben & \aacht  & \aneun & \azehn & \aelf
\\ \hline erhaltene Pkt.  & \phantom{11} & \phantom{11} & \phantom{11} & \phantom{11} & \phantom{11}& \phantom{11} & \phantom{11} & \phantom{11} & \phantom{11} & & \\ \hline \end{tabular} \end{center} \vspace{2ex} }

\newcommand{\punktetabelleelf}
{\vspace{3ex} \begin{center} \renewcommand{\arraystretch}{1.5} \tabcolsep=0.78ex \begin{tabular}{@{}|c|c|c|c|c|c|c|c|c|c|c|c|c|c|}
\hline Aufgabe & 1 & 2 & 3 & 4 & 5 & 6 & 7 & 8 & 9 & 10 & 11 & $\sum$
\\ \hline m\"ogliche Pkt. & \aeins & \azwei & \adrei & \avier & \afuenf & \asechs   & \asieben & \aacht  & \aneun & \azehn & \aelf & \azwoelf
\\ \hline erhaltene Pkt.  & \phantom{11} & \phantom{11} & \phantom{11} & \phantom{11} & \phantom{11}& \phantom{11} & \phantom{11} & \phantom{11} & \phantom{11} & \phantom{11} & & \\ \hline \end{tabular} \end{center} \vspace{2ex} }

\newcommand{\punktetabellezwoelf}
{\vspace{3ex} \begin{center} \renewcommand{\arraystretch}{1.5} \tabcolsep=0.78ex \begin{tabular}{@{}|c|c|c|c|c|c|c|c|c|c|c|c|c|c|c|}
\hline Aufgabe & 1 & 2 & 3 & 4 & 5 & 6 & 7 & 8 & 9 & 10 & 11 & 12    & $\sum$
\\ \hline m\"ogliche Pkt. & \aeins & \azwei & \adrei & \avier & \afuenf & \asechs   & \asieben & \aacht  & \aneun & \azehn & \aelf & \azwoelf & \adreizehn
\\ \hline erhaltene Pkt.  & \phantom{11} & \phantom{11} & \phantom{11} & \phantom{11} & \phantom{11}& \phantom{11} & \phantom{11} & \phantom{11}
& \phantom{11} & \phantom{11} & \phantom{11} & & \\ \hline \end{tabular} \end{center} \vspace{2ex} }

\newcommand{\punktetabelledreizehn}
{\vspace{3ex} \begin{center} \renewcommand{\arraystretch}{1.5} \tabcolsep=0.78ex \begin{tabular}{@{}|c|c|c|c|c|c|c|c|c|c|c|c|c|c|c|c|}
\hline Aufgabe & 1 & 2 & 3 & 4 & 5 & 6 & 7 & 8 & 9 & 10 & 11 & 12 & 13   & $\sum$
\\ \hline m\"ogliche Pkt. & \aeins & \azwei & \adrei & \avier & \afuenf & \asechs   & \asieben & \aacht  & \aneun & \azehn & \aelf & \azwoelf & \adreizehn & \avierzehn
\\ \hline erhaltene Pkt.  & \phantom{11} & \phantom{11} & \phantom{11} & \phantom{11} & \phantom{11}& \phantom{11} & \phantom{11} & \phantom{11} & \phantom{11}
& \phantom{11} & \phantom{11} & \phantom{11} & & \\ \hline \end{tabular} \end{center} \vspace{2ex} }

\newcommand{\punktetabellevierzehn}
{\vspace{3ex} \begin{center} \renewcommand{\arraystretch}{1.5} \tabcolsep=0.78ex \begin{tabular}{@{}|c|c|c|c|c|c|c|c|c|c|c|c|c|c|c|c|c|}
\hline Aufgabe & 1 & 2 & 3 & 4 & 5 & 6 & 7 & 8 & 9 & 10 & 11 & 12 & 13 & 14  & $\sum$
\\ \hline m\"ogliche Pkt. & \aeins & \azwei & \adrei & \avier & \afuenf & \asechs   & \asieben & \aacht  & \aneun & \azehn & \aelf & \azwoelf & \adreizehn & \avierzehn & \afuenfzehn
\\ \hline erhaltene Pkt. & \phantom{11} & \phantom{11} & \phantom{11} & \phantom{11} & \phantom{11} & \phantom{11}& \phantom{11} & \phantom{11} & \phantom{11} & \phantom{11}
& \phantom{11} & \phantom{11} & \phantom{11} & & \\ \hline \end{tabular} \end{center} \vspace{2ex} }

\newcommand{\punktetabellefuenfzehn}
{\vspace{3ex} \begin{center} \renewcommand{\arraystretch}{1.5} \tabcolsep=0.78ex \begin{tabular}{@{}|c|c|c|c|c|c|c|c|c|c|c|c|c|c|c|c|c|c|}
\hline Aufgabe & 1 & 2 & 3 & 4 & 5 & 6 & 7 & 8 & 9 & 10 & 11 & 12 & 13 & 14 & 15 & $\sum$
\\ \hline m\"ogliche Pkt. & \aeins & \azwei & \adrei & \avier & \afuenf & \asechs   & \asieben & \aacht  & \aneun & \azehn & \aelf & \azwoelf & \adreizehn & \avierzehn & \afuenfzehn
& \asechzehn
\\ \hline erhaltene Pkt. & \phantom{11} & \phantom{11} & \phantom{11}
& \phantom{11} & \phantom{11} & \phantom{11} & \phantom{11}& \phantom{11} & \phantom{11} & \phantom{11} & \phantom{11}
& \phantom{11} & \phantom{11} & \phantom{11} & & \\ \hline \end{tabular} \end{center} \vspace{2ex} }

\newcommand{\punktetabellesechzehn}
{\vspace{3ex} \begin{center} \renewcommand{\arraystretch}{1.5} \tabcolsep=0.78ex \begin{tabular}{@{}|c|c|c|c|c|c|c|c|c|c|c|c|c|c|c|c|c|c|}
\hline Aufgabe & 1 & 2 & 3 & 4 & 5 & 6 & 7 & 8 & 9 & 10 & 11 & 12 & 13 & 14 & 15 & 16 & $\sum$
\\ \hline m\"ogliche Pkt. & \aeins & \azwei & \adrei & \avier & \afuenf & \asechs   & \asieben & \aacht  & \aneun & \azehn & \aelf & \azwoelf & \adreizehn & \avierzehn & \afuenfzehn
& \asechzehn   & \asiebzehn
\\ \hline erhaltene Pkt. & \phantom{11} & \phantom{11} & \phantom{11} & \phantom{11}
& \phantom{11} & \phantom{11} & \phantom{11} & \phantom{11}& \phantom{11} & \phantom{11} & \phantom{11} & \phantom{11}
& \phantom{11} & \phantom{11} & \phantom{11} & & \\ \hline \end{tabular} \end{center} \vspace{2ex} }

\newcommand{\punktetabellesiebzehn}
{\vspace{3ex} \begin{center} \renewcommand{\arraystretch}{1.5} \tabcolsep=0.78ex \begin{tabular}{@{}|c|c|c|c|c|c|c|c|c|c|c|c|c|c|c|c|c|c|c|c|}
\hline Aufgabe & 1 & 2 & 3 & 4 & 5 & 6 & 7 & 8 & 9 & 10 & 11 & 12 & 13 & 14 & 15 & 16 & 17 & $\sum$
\\ \hline m\"ogliche Pkt. & \aeins & \azwei & \adrei & \avier & \afuenf & \asechs   & \asieben & \aacht  & \aneun & \azehn & \aelf & \azwoelf & \adreizehn & \avierzehn & \afuenfzehn
& \asechzehn   & \asiebzehn & \aachtzehn
\\ \hline erhaltene Pkt. & \phantom{11} & \phantom{11} & \phantom{11} & \phantom{11}
& \phantom{11} & \phantom{11} & \phantom{11} & \phantom{11} & \phantom{11}& \phantom{11} & \phantom{11} & \phantom{11} & \phantom{11}
& \phantom{11} & \phantom{11} & \phantom{11} & & \\ \hline \end{tabular} \end{center} \vspace{2ex} }

\newcommand{\punktetabelleachtzehn}
{\vspace{3ex} \begin{center} \renewcommand{\arraystretch}{1.5} \tabcolsep=0.78ex \begin{tabular}{@{}|c|c|c|c|c|c|c|c|c|c|c|c|c|c|c|c|c|c|c|c|c|}
\hline Aufgabe & 1 & 2 & 3 & 4 & 5 & 6 & 7 & 8 & 9 & 10 & 11 & 12 & 13 & 14 & 15 & 16 & 17 & 18 & $\sum$
\\ \hline m\"ogliche Pkt. & \aeins & \azwei & \adrei & \avier & \afuenf & \asechs   & \asieben & \aacht  & \aneun & \azehn & \aelf & \azwoelf & \adreizehn & \avierzehn & \afuenfzehn
& \asechzehn   & \asiebzehn & \aachtzehn & \aneunzehn
\\ \hline erhaltene Pkt. & \phantom{11} & \phantom{11} & \phantom{11} & \phantom{11} & \phantom{11}
& \phantom{11} & \phantom{11} & \phantom{11} & \phantom{11} & \phantom{11}& \phantom{11} & \phantom{11} & \phantom{11} & \phantom{11}
& \phantom{11} & \phantom{11} & \phantom{11} & & \\ \hline \end{tabular} \end{center} \vspace{2ex} }

\newcommand{\punktetabelleneunzehn}
{\vspace{3ex} \begin{center} \renewcommand{\arraystretch}{1.5} \tabcolsep=0.78ex \begin{tabular}{@{}|c|c|c|c|c|c|c|c|c|c|c|c|c|c|c|c|c|c|c|c|c|c|}
\hline Aufgabe & 1 & 2 & 3 & 4 & 5 & 6 & 7 & 8 & 9 & 10 & 11 & 12 & 13 & 14 & 15 & 16 & 17 & 18 & 19 & $\sum$
\\ \hline m\"ogliche Pkt. & \aeins & \azwei & \adrei & \avier & \afuenf & \asechs   & \asieben & \aacht  & \aneun & \azehn & \aelf & \azwoelf & \adreizehn & \avierzehn & \afuenfzehn
& \asechzehn   & \asiebzehn & \aachtzehn & \aneunzehn  & \azwanzig
\\ \hline erhaltene Pkt. & \phantom{11} & \phantom{11} & \phantom{11} & \phantom{11} & \phantom{11} & \phantom{11}
& \phantom{11} & \phantom{11} & \phantom{11} & \phantom{11} & \phantom{11}& \phantom{11} & \phantom{11} & \phantom{11} & \phantom{11}
& \phantom{11} & \phantom{11} & \phantom{11} & & \\ \hline \end{tabular} \end{center} \vspace{2ex} }

\newcommand{\punktetabellezwanzig}
{\vspace{3ex} \begin{center} \renewcommand{\arraystretch}{1.5} \tabcolsep=0.78ex \begin{tabular}{@{}|c|c|c|c|c|c|c|c|c|c|c|c|c|c|c|c|c|c|c|c|c|c|c|}
\hline Aufgabe & 1 & 2 & 3 & 4 & 5 & 6 & 7 & 8 & 9 & 10 & 11 & 12 & 13 & 14 & 15 & 16 & 17 & 18 & 19 & 20 & $\sum$
\\ \hline m\"ogliche Pkt. & \aeins & \azwei & \adrei & \avier & \afuenf & \asechs   & \asieben & \aacht  & \aneun & \azehn & \aelf & \azwoelf & \adreizehn & \avierzehn & \afuenfzehn
& \asechzehn   & \asiebzehn & \aachtzehn & \aneunzehn  & \azwanzig & \aeinundzwanzig
\\ \hline erhaltene Pkt. & \phantom{11} & \phantom{11} & \phantom{11} & \phantom{11} & \phantom{11} & \phantom{11} & \phantom{11}
& \phantom{11} & \phantom{11} & \phantom{11} & \phantom{11} & \phantom{11}& \phantom{11} & \phantom{11} & \phantom{11} & \phantom{11}
& \phantom{11} & \phantom{11} & \phantom{11} & & \\ \hline \end{tabular} \end{center} \vspace{2ex} }

\newcommand{\punktetabelleeinundzwanzig}
{\vspace{3ex} \begin{center} \renewcommand{\arraystretch}{1.5} \tabcolsep=0.78ex \begin{tabular}{@{}|c|c|c|c|c|c|c|c|c|c|c|c|c|c|c|c|c|c|c|c|c|c|c|c|}
\hline Aufgabe & 1 & 2 & 3 & 4 & 5 & 6 & 7 & 8 & 9 & 10 & 11 & 12 & 13 & 14 & 15 & 16 & 17 & 18 & 19 & 20 & 21 & $\sum$
\\ \hline m\"ogliche Pkt. & \aeins & \azwei & \adrei & \avier & \afuenf & \asechs   & \asieben & \aacht  & \aneun & \azehn & \aelf & \azwoelf & \adreizehn & \avierzehn & \afuenfzehn
& \asechzehn   & \asiebzehn & \aachtzehn & \aneunzehn  & \azwanzig & \aeinundzwanzig  & \azweiundzwanzig
\\ \hline erhaltene Pkt. & \phantom{11} & \phantom{11} & \phantom{11} & \phantom{11} & \phantom{11} & \phantom{11} & \phantom{11}
& \phantom{11} & \phantom{11} & \phantom{11} & \phantom{11} & \phantom{11}& \phantom{11} & \phantom{11} & \phantom{11} & \phantom{11}
& \phantom{11} & \phantom{11} & \phantom{11} & \phantom{11} & & \\ \hline \end{tabular} \end{center} \vspace{2ex} }

\newcommand{\punktetabellezweiundzwanzig}
{\vspace{3ex} \begin{center} \renewcommand{\arraystretch}{1.5} \tabcolsep=0.78ex \begin{tabular}{@{}|c|c|c|c|c|c|c|c|c|c|c|c|c|c|c|c|c|c|c|c|c|c|c|c|c|}
\hline Aufgabe & 1 & 2 & 3 & 4 & 5 & 6 & 7 & 8 & 9 & 10 & 11 & 12 & 13 & 14 & 15 & 16 & 17 & 18 & 19 & 20 & 21 & 22 & $\sum$
\\ \hline  Punkte & \aeins & \azwei & \adrei & \avier & \afuenf & \asechs   & \asieben & \aacht  & \aneun & \azehn & \aelf & \azwoelf & \adreizehn & \avierzehn & \afuenfzehn
& \asechzehn   & \asiebzehn & \aachtzehn & \aneunzehn  & \azwanzig & \aeinundzwanzig  & \azweiundzwanzig  & \adreiundzwanzig
\\ \hline Erhalten & \phantom{11} & \phantom{11} & \phantom{11} & \phantom{11} & \phantom{11} & \phantom{11} & \phantom{11}
& \phantom{11} & \phantom{11} & \phantom{11} & \phantom{11} & \phantom{11}& \phantom{11} & \phantom{11} & \phantom{11}  & \phantom{11} & \phantom{11}
& \phantom{11} & \phantom{11} & \phantom{11} & \phantom{11} & & \\ \hline \end{tabular} \end{center} \vspace{2ex} }

%[[Projekt:Semantische Vorlagen/Wahrheitstabellen/Latex]]

\newcommand{\wahrheitstabelleeins}[4]{
\begin{center}
\begin{tabular}{|c|c|}
\multicolumn{2}{l}{\textbf{#1}}\\
\hline
$p$ & $#2$\\ \hline
w  &	#3 \\ \hline
f  &	#4\\ \hline
\end{tabular}
\end{center} }

\newcommand{\wahrheitstabelleeinszwei}[7]{
\begin{center}
\begin{tabular}{|c|c|c|}
\multicolumn{3}{l}{\textbf{#1}}\\
\hline
$p$ & $#2$ & $#5$ \\ \hline
w  &	#3 &  #4 \\ \hline
f  &	#6 &  #7 \\ \hline
\end{tabular}
\end{center} }

\newcommand{\wahrheitstabelleeinsdrei}[4]{ \begin{center} \begin{tabular}{|c|c|c|c|} \multicolumn{4}{l}{\textbf{#1}} \\  \hline #2  \\  \hline #3  \\  \hline  #4  \\  \hline
\end{tabular} \end{center} }

%Wahrheitstabellen mit zwei Variablen p und q

\newcommand{\wahrheitstabellezweieins}[6]{ \begin{center} \begin{tabular}{|c|c|c|} \multicolumn{3}{l}{\textbf{#1}} \\  \hline  #2  \\  \hline  #3  \\  \hline  #4  \\  \hline  #5  \\  \hline  #6  \\  \hline \end{tabular} \end{center} }

\newcommand{\wahrheitstabellezweizwei}[6]{ \begin{center} \begin{tabular}{|c|c|c|c|} \multicolumn{4}{l}{\textbf{#1}} \\  \hline #2  \\  \hline #3  \\  \hline #4  \\  \hline #5  \\  \hline #6  \\  \hline\end{tabular} \end{center} }

\newcommand{\wahrheitstabellezweidrei}[6]{ \begin{center} \begin{tabular}{|c|c|c|c|c|} \multicolumn{5}{l}{\textbf{#1}} \\  \hline #2  \\  \hline #3  \\  \hline #4  \\  \hline #5  \\  \hline #6  \\  \hline \end{tabular} \end{center} }

\newcommand{\wahrheitstabellezweivier}[6]{ \begin{center} \begin{tabular}{|c|c|c|c|c|c|} \multicolumn{6}{l}{\textbf{#1}} \\  \hline #2  \\  \hline #3  \\  \hline #4  \\  \hline  #5  \\  \hline  #6  \\  \hline \end{tabular} \end{center} }

\newcommand{\wahrheitstabellezweifuenf}[6]{ \begin{center} \begin{tabular}{|c|c|c|c|c|c|c|} \multicolumn{7}{l}{\textbf{#1}} \\  \hline #2  \\  \hline #3  \\  \hline #4  \\  \hline #5  \\  \hline #6  \\  \hline
\end{tabular} \end{center} }

%Wahrheitstabellen mit drei Variablen p,q,r

\newcommand{\wahrheitstabelledreieins}[6]{ \begin{center} \begin{tabular}{|c|c|c|c||} \multicolumn{4}{l}{\textbf{#1}} \\  \hline #2  \\  \hline #3  \\  \hline #4  \\  \hline #5  \\  \hline #6 \\  \hline
\end{tabular} \end{center} }

%[[Projekt:Semantische Vorlagen/Wertetabellen/Latex]]

%Wertetabellen

\newcommand{\wertetabellevorskip}{\par \bigskip}

\newcommand{\wertetabellenachskip}{\par \bigskip}

\newcommand{\wertetabelleeinsausteilzeilen}[4]{\wertetabellevorskip \begin{tabular}{|c|c|} \hline #1 & #2  \\ \hline  #3 & #4  \\ \hline \end{tabular} \wertetabellenachskip }

\newcommand{\wertetabellezweiausteilzeilen}[4]{\wertetabellevorskip \begin{tabular}{|c|c|c|} \hline #1 & #2  \\ \hline  #3 & #4  \\ \hline \end{tabular} \wertetabellenachskip }

\newcommand{\wertetabelledreiausteilzeilen}[4]{\wertetabellevorskip \begin{tabular}{|c|c|c|c|} \hline #1 & #2  \\ \hline  #3 & #4  \\ \hline \end{tabular} \wertetabellenachskip }

\newcommand{\wertetabellevierausteilzeilen}[4]{\wertetabellevorskip \begin{tabular}{|c|c|c|c|c|} \hline #1 & #2  \\ \hline  #3 & #4  \\ \hline \end{tabular} \wertetabellenachskip }

\newcommand{\wertetabellefuenfausteilzeilen}[4]{\wertetabellevorskip \begin{tabular}{|c|c|c|c|c|c|} \hline #1 & #2  \\ \hline  #3 & #4  \\ \hline \end{tabular} \wertetabellenachskip }

\newcommand{\wertetabellesechsausteilzeilen}[6]{\wertetabellevorskip \begin{tabular}{|c|c|c|c|c|c|c|} \hline #1 & #2 & #3 \\ \hline  #4 & #5 & #6 \\ \hline \end{tabular} \wertetabellenachskip }

\newcommand{\wertetabellesiebenausteilzeilen}[6]{\wertetabellevorskip \begin{tabular}{|c|c|c|c|c|c|c|c|} \hline #1 & #2 & #3 \\ \hline  #4 & #5 & #6 \\ \hline \end{tabular} \wertetabellenachskip }

\newcommand{\wertetabelleachtausteilzeilen}[6]{\wertetabellevorskip \begin{tabular}{|c|c|c|c|c|c|c|c|c|} \hline #1 & #2 & #3 \\ \hline  #4 & #5 & #6 \\ \hline \end{tabular} \wertetabellenachskip }

\newcommand{\wertetabelleneunausteilzeilen}[6]{\wertetabellevorskip \begin{tabular}{|c|c|c|c|c|c|c|c|c|c|} \hline #1 & #2 & #3 \\ \hline  #4 & #5 & #6 \\ \hline \end{tabular} \wertetabellenachskip }

\newcommand{\wertetabellezehnausteilzeilen}[6]{\wertetabellevorskip \begin{tabular}{|c|c|c|c|c|c|c|c|c|c|c|} \hline #1 & #2 & #3 \\ \hline  #4 & #5 & #6 \\ \hline \end{tabular} \wertetabellenachskip }

\newcommand{\wertetabelleelfausteilzeilen}[8]{\wertetabellevorskip \begin{tabular}{|c|c|c|c|c|c|c|c|c|c|c|c|} \hline #1 & #2 & #3 & #4 \\ \hline #5 & #6 & #7 & #8 \\ \hline \end{tabular} \wertetabellenachskip }

%[[Projekt:Semantische Vorlagen/Verknüpfungstabellen/Latex]]

%Verknüpfungstabellen

\newcommand{\verknuepfungstabelleeins}[2]{\begin{center}
\begin{tabular}{|c||c|}
\hline #1 \\  \hline \hline #2 \\
\hline
\end{tabular}
\end{center} }

\newcommand{\verknuepfungstabellezwei}[3]{\begin{center}
\begin{tabular}{|c||c|c|}
\hline #1 \\  \hline \hline #2 \\ \hline #3 \\
\hline
\end{tabular}
\end{center} }

\newcommand{\verknuepfungstabelledrei}[4]{\begin{center}
\begin{tabular}{|c||c|c|c|}
\hline #1 \\  \hline \hline #2 \\ \hline #3 \\ \hline #4 \\
\hline
\end{tabular}
\end{center} }

\newcommand{\verknuepfungstabellevier}[5]{\begin{center}
\begin{tabular}{|c||c|c|c|c|}
\hline #1 \\  \hline \hline #2 \\ \hline #3 \\ \hline #4 \\  \hline  #5 \\
\hline
\end{tabular}
\end{center} }

\newcommand{\verknuepfungstabellefuenf}[6]{\begin{center}
\begin{tabular}{|c||c|c|c|c|c|}
\hline #1 \\  \hline \hline #2 \\ \hline #3 \\ \hline #4 \\  \hline  #5 \\ \hline  #6 \\
\hline
\end{tabular}
\end{center} }

\newcommand{\verknuepfungstabellesechs}[7]{\begin{center}
\begin{tabular}{|c||c|c|c|c|c|c|}
\hline #1 \\  \hline \hline #2 \\ \hline #3 \\ \hline #4 \\  \hline  #5 \\ \hline  #6 \\  \hline #7 \\
\hline
\end{tabular}
\end{center} }

\newcommand{\verknuepfungstabellesieben}[8]{\begin{center}
\begin{tabular}{|c||c|c|c|c|c|c|c|c|}
\hline #1 \\  \hline \hline #2 \\ \hline #3 \\ \hline #4 \\  \hline  #5 \\ \hline  #6 \\  \hline #7 \\  \hline  #8 \\
\hline
\end{tabular}
\end{center} }

\newcommand{\verknuepfungstabelleacht}[9]{\begin{center}
\begin{tabular}{|c||c|c|c|c|c|c|c|c|}
\hline #1 \\  \hline \hline #2 \\ \hline #3 \\ \hline #4 \\  \hline  #5 \\ \hline  #6 \\  \hline #7 \\  \hline  #8 \\ \hline  #9 \\
\hline
\end{tabular}
\end{center} }

%[[Projekt:Semantische Vorlagen/Initialisierung für Daten/Latex]]

\newcommand{\aeins}{    }

\newcommand{\azwei}{     }

\newcommand{\adrei}{    }

\newcommand{\avier}{    }

\newcommand{\afuenf}{    }

\newcommand{\asechs}{    }

\newcommand{\asieben}{    }

\newcommand{\aacht}{    }

\newcommand{\aneun}{    }

\newcommand{\azehn}{    }

\newcommand{\aelf}{    }

\newcommand{\azwoelf}{    }

\newcommand{\adreizehn}{     }

\newcommand{\avierzehn}{    }

\newcommand{\afuenfzehn}{    }

\newcommand{\asechzehn}{    }

\newcommand{\asiebzehn}{    }

\newcommand{\aachtzehn}{    }

\newcommand{\aneunzehn}{    }

\newcommand{\azwanzig}{    }

\newcommand{\aeinundzwanzig}{    }

\newcommand{\azweiundzwanzig}{    }

\newcommand{\adreiundzwanzig}{    }

\newcommand{\avierundzwanzig}{    }

\newcommand{\afuenfundzwanzig}{    }

\newcommand{\asechsundzwanzig}{    }

\newcommand{\asiebenundzwanzig}{    }

\newcommand{\aachtundzwanzig}{    }

\newcommand{\aneunundzwanzig}{    }

\newcommand{\leitzeilenull}{}

\newcommand{\leitzeileeins}{}

\newcommand{\leitzeilezwei}{}

\newcommand{\leitzeiledrei}{}

\newcommand{\leitzeilevier}{}

\newcommand{\leitzeilefuenf}{}

\newcommand{\leitzeilesechs}{}

\newcommand{\leitzeilesieben}{}

\newcommand{\leitzeileacht}{}

\newcommand{\leitzeileneun}{}

\newcommand{\leitzeilezehn}{}

\newcommand{\leitzeileelf}{}

\newcommand{\leitzeilezwoelf}{}

\newcommand{\leitspaltenull}{}

\newcommand{\leitspalteeins}{}

\newcommand{\leitspaltezwei}{}

\newcommand{\leitspaltedrei}{}

\newcommand{\leitspaltevier}{}

\newcommand{\leitspaltefuenf}{}

\newcommand{\leitspaltesechs}{}

\newcommand{\leitspaltesieben}{}

\newcommand{\leitspalteacht}{}

\newcommand{\leitspalteneun}{}

\newcommand{\leitspaltezehn}{}

\newcommand{\leitspalteelf}{}

\newcommand{\leitspaltezwoelf}{}

\newcommand{\leitspaltedreizehn}{}

\newcommand{\leitspaltevierzehn}{}

\newcommand{\leitspaltefuenfzehn}{}

\newcommand{\leitspaltesechzehn}{}

\newcommand{\leitspaltesiebzehn}{}

\newcommand{\leitspalteachtzehn}{}

\newcommand{\leitspalteneunzehn}{}

\newcommand{\leitspaltezwanzig}{}

\newcommand{\aeinsxeins}{}

\newcommand{\aeinsxzwei}{}

\newcommand{\aeinsxdrei}{}

\newcommand{\aeinsxvier}{}

\newcommand{\aeinsxfuenf}{}

\newcommand{\aeinsxsechs}{}

\newcommand{\aeinsxsieben}{}

\newcommand{\aeinsxacht}{}

\newcommand{\aeinsxneun}{}

\newcommand{\aeinsxzehn}{}

\newcommand{\aeinsxelf}{}

\newcommand{\aeinsxzwoelf}{}

\newcommand{\azweixeins}{}

\newcommand{\azweixzwei}{}

\newcommand{\azweixdrei}{}

\newcommand{\azweixvier}{}

\newcommand{\azweixfuenf}{}

\newcommand{\azweixsechs}{}

\newcommand{\azweixsieben}{}

\newcommand{\azweixacht}{}

\newcommand{\azweixneun}{}

\newcommand{\azweixzehn}{}

\newcommand{\azweixelf}{}

\newcommand{\azweixzwoelf}{}

\newcommand{\adreixeins}{}

\newcommand{\adreixzwei}{}

\newcommand{\adreixdrei}{}

\newcommand{\adreixvier}{}

\newcommand{\adreixfuenf}{}

\newcommand{\adreixsechs}{}

\newcommand{\adreixsieben}{}

\newcommand{\adreixacht}{}

\newcommand{\adreixneun}{}

\newcommand{\adreixzehn}{}

\newcommand{\adreixelf}{}

\newcommand{\adreixzwoelf}{}

\newcommand{\avierxeins}{}

\newcommand{\avierxzwei}{}

\newcommand{\avierxdrei}{}

\newcommand{\avierxvier}{}

\newcommand{\avierxfuenf}{}

\newcommand{\avierxsechs}{}

\newcommand{\avierxsieben}{}

\newcommand{\avierxacht}{}

\newcommand{\avierxneun}{}

\newcommand{\avierxzehn}{}

\newcommand{\avierxelf}{}

\newcommand{\avierxzwoelf}{}

\newcommand{\afuenfxeins}{}

\newcommand{\afuenfxzwei}{}

\newcommand{\afuenfxdrei}{}

\newcommand{\afuenfxvier}{}

\newcommand{\afuenfxfuenf}{}

\newcommand{\afuenfxsechs}{}

\newcommand{\afuenfxsieben}{}

\newcommand{\afuenfxacht}{}

\newcommand{\afuenfxneun}{}

\newcommand{\afuenfxzehn}{}

\newcommand{\afuenfxelf}{}

\newcommand{\afuenfxzwoelf}{}

\newcommand{\asechsxeins}{}

\newcommand{\asechsxzwei}{}

\newcommand{\asechsxdrei}{}

\newcommand{\asechsxvier}{}

\newcommand{\asechsxfuenf}{}

\newcommand{\asechsxsechs}{}

\newcommand{\asechsxsieben}{}

\newcommand{\asechsxacht}{}

\newcommand{\asechsxneun}{}

\newcommand{\asechsxzehn}{}

\newcommand{\asechsxelf}{}

\newcommand{\asechsxzwoelf}{}

\newcommand{\asiebenxeins}{}

\newcommand{\asiebenxzwei}{}

\newcommand{\asiebenxdrei}{}

\newcommand{\asiebenxvier}{}

\newcommand{\asiebenxfuenf}{}

\newcommand{\asiebenxsechs}{}

\newcommand{\asiebenxsieben}{}

\newcommand{\asiebenxacht}{}

\newcommand{\asiebenxneun}{}

\newcommand{\asiebenxzehn}{}

\newcommand{\asiebenxelf}{}

\newcommand{\asiebenxzwoelf}{}

\newcommand{\aachtxeins}{}

\newcommand{\aachtxzwei}{}

\newcommand{\aachtxdrei}{}

\newcommand{\aachtxvier}{}

\newcommand{\aachtxfuenf}{}

\newcommand{\aachtxsechs}{}

\newcommand{\aachtxsieben}{}

\newcommand{\aachtxacht}{}

\newcommand{\aachtxneun}{}

\newcommand{\aachtxzehn}{}

\newcommand{\aachtxelf}{}

\newcommand{\aachtxzwoelf}{}

\newcommand{\aneunxeins}{}

\newcommand{\aneunxzwei}{}

\newcommand{\aneunxdrei}{}

\newcommand{\aneunxvier}{}

\newcommand{\aneunxfuenf}{}

\newcommand{\aneunxsechs}{}

\newcommand{\aneunxsieben}{}

\newcommand{\aneunxacht}{}

\newcommand{\aneunxneun}{}

\newcommand{\aneunxzehn}{}

\newcommand{\aneunxelf}{}

\newcommand{\aneunxzwoelf}{}

\newcommand{\azehnxeins}{}

\newcommand{\azehnxzwei}{}

\newcommand{\azehnxdrei}{}

\newcommand{\azehnxvier}{}

\newcommand{\azehnxfuenf}{}

\newcommand{\azehnxsechs}{}

\newcommand{\azehnxsieben}{}

\newcommand{\azehnxacht}{}

\newcommand{\azehnxneun}{}

\newcommand{\azehnxzehn}{}

\newcommand{\azehnxelf}{}

\newcommand{\azehnxzwoelf}{}

\newcommand{\aelfxeins}{}

\newcommand{\aelfxzwei}{}

\newcommand{\aelfxdrei}{}

\newcommand{\aelfxvier}{}

\newcommand{\aelfxfuenf}{}

\newcommand{\aelfxsechs}{}

\newcommand{\aelfxsieben}{}

\newcommand{\aelfxacht}{}

\newcommand{\aelfxneun}{}

\newcommand{\aelfxzehn}{}

\newcommand{\aelfxelf}{}

\newcommand{\aelfxzwoelf}{}

\newcommand{\azwoelfxeins}{}

\newcommand{\azwoelfxzwei}{}

\newcommand{\azwoelfxdrei}{}

\newcommand{\azwoelfxvier}{}

\newcommand{\azwoelfxfuenf}{}

\newcommand{\azwoelfxsechs}{}

\newcommand{\azwoelfxsieben}{}

\newcommand{\azwoelfxacht}{}

\newcommand{\azwoelfxneun}{}

\newcommand{\azwoelfxzehn}{}

\newcommand{\azwoelfxelf}{}

\newcommand{\azwoelfxzwoelf}{}

\newcommand{\adreizehnxeins}{}

\newcommand{\adreizehnxzwei}{}

\newcommand{\adreizehnxdrei}{}

\newcommand{\adreizehnxvier}{}

\newcommand{\adreizehnxfuenf}{}

\newcommand{\adreizehnxsechs}{}

\newcommand{\adreizehnxsieben}{}

\newcommand{\adreizehnxacht}{}

\newcommand{\adreizehnxneun}{}

\newcommand{\adreizehnxzehn}{}

\newcommand{\adreizehnxelf}{}

\newcommand{\adreizehnxzwoelf}{}

\newcommand{\avierzehnxeins}{}

\newcommand{\avierzehnxzwei}{}

\newcommand{\avierzehnxdrei}{}

\newcommand{\avierzehnxvier}{}

\newcommand{\avierzehnxfuenf}{}

\newcommand{\avierzehnxsechs}{}

\newcommand{\avierzehnxsieben}{}

\newcommand{\avierzehnxacht}{}

\newcommand{\avierzehnxneun}{}

\newcommand{\avierzehnxzehn}{}

\newcommand{\avierzehnxelf}{}

\newcommand{\avierzehnxzwoelf}{}

\newcommand{\afuenfzehnxeins}{}

\newcommand{\afuenfzehnxzwei}{}

\newcommand{\afuenfzehnxdrei}{}

\newcommand{\afuenfzehnxvier}{}

\newcommand{\afuenfzehnxfuenf}{}

\newcommand{\afuenfzehnxsechs}{}

\newcommand{\afuenfzehnxsieben}{}

\newcommand{\afuenfzehnxacht}{}

\newcommand{\afuenfzehnxneun}{}

\newcommand{\afuenfzehnxzehn}{}

\newcommand{\afuenfzehnxelf}{}

\newcommand{\afuenfzehnxzwoelf}{}

\newcommand{\asechzehnxeins}{}

\newcommand{\asechzehnxzwei}{}

\newcommand{\asechzehnxdrei}{}

\newcommand{\asechzehnxvier}{}

\newcommand{\asechzehnxfuenf}{}

\newcommand{\asechzehnxsechs}{}

\newcommand{\asechzehnxsieben}{}

\newcommand{\asechzehnxacht}{}

\newcommand{\asechzehnxneun}{}

\newcommand{\asechzehnxzehn}{}

\newcommand{\asechzehnxelf}{}

\newcommand{\asechzehnxzwoelf}{}

\newcommand{\asiebzehnxeins}{}

\newcommand{\asiebzehnxzwei}{}

\newcommand{\asiebzehnxdrei}{}

\newcommand{\asiebzehnxvier}{}

\newcommand{\asiebzehnxfuenf}{}

\newcommand{\asiebzehnxsechs}{}

\newcommand{\asiebzehnxsieben}{}

\newcommand{\asiebzehnxacht}{}

\newcommand{\asiebzehnxneun}{}

\newcommand{\asiebzehnxzehn}{}

\newcommand{\asiebzehnxelf}{}

\newcommand{\asiebzehnxzwoelf}{}

\newcommand{\aachtzehnxeins}{}

\newcommand{\aachtzehnxzwei}{}

\newcommand{\aachtzehnxdrei}{}

\newcommand{\aachtzehnxvier}{}

\newcommand{\aachtzehnxfuenf}{}

\newcommand{\aachtzehnxsechs}{}

\newcommand{\aachtzehnxsieben}{}

\newcommand{\aachtzehnxacht}{}

\newcommand{\aachtzehnxneun}{}

\newcommand{\aachtzehnxzehn}{}

\newcommand{\aachtzehnxelf}{}

\newcommand{\aachtzehnxzwoelf}{}

%[[Projekt:Semantische Vorlagen/Tabellen mit benannten Parametern/Latex]]

%Tabellen mit zwei Zeilen

\newcommand{\tabelleleitzweixzwei}{

\begin{center}

\begin{tabular}{|c||c|c|}

\hline  \leitzeilenull  &  \leitzeileeins  &  \leitzeilezwei   \\

\hline \hline  \leitspalteeins  &  $ \aeinsxeins $  &  $ \aeinsxzwei $ \\

\hline  \leitspaltezwei  &  $ \azweixeins $  &  $ \azweixzwei $   \\

\hline

\end{tabular}

\end{center} }

%Tabellen mit drei Zeilen

\newcommand{\tabelleleitdreixdrei}{

\begin{center}

\begin{tabular}{|c||c|c|c||}

\hline  \leitzeilenull  &  \leitzeileeins  &  \leitzeilezwei  &  \leitzeiledrei        \\

\hline \hline  \leitspalteeins  &  $ \aeinsxeins $  &  $ \aeinsxzwei $  &  $ \aeinsxdrei $      \\

\hline  \leitspaltezwei  &  $ \azweixeins $  &  $ \azweixzwei $  &  $ \azweixdrei $       \\

\hline  \leitspaltedrei  &  $ \adreixeins $  &  $ \adreixzwei $  &  $ \adreixdrei $       \\

\hline

\end{tabular}

\end{center} }

\newcommand{\tabelleleitdreixvier}{

\begin{center}

\begin{tabular}{|c||c|c|c|c||}

\hline  \leitzeilenull  &  \leitzeileeins  &  \leitzeilezwei  &  \leitzeiledrei  &  \leitzeilevier      \\

\hline \hline  \leitspalteeins  &  $ \aeinsxeins $  &  $ \aeinsxzwei $  &  $ \aeinsxdrei $  &  $ \aeinsxvier $   \\

\hline  \leitspaltezwei  &  $ \azweixeins $  &  $ \azweixzwei $  &  $ \azweixdrei $  &  $ \azweixvier $     \\

\hline  \leitspaltedrei  &  $ \adreixeins $  &  $ \adreixzwei $  &  $ \adreixdrei $  &  $ \adreixvier $     \\

\hline

\end{tabular}

\end{center} }

\newcommand{\tabelleleitdreixsechs}{

\begin{center}

\begin{tabular}{|c||c|c|c|c|c|c|}

\hline  \leitzeilenull  &  \leitzeileeins  &  \leitzeilezwei  &  \leitzeiledrei  &  \leitzeilevier  &  \leitzeilefuenf   &  \leitzeilesechs    \\

\hline \hline  \leitspalteeins  &  $ \aeinsxeins $  &  $ \aeinsxzwei $  &  $ \aeinsxdrei $  &  $ \aeinsxvier $  &  $ \aeinsxfuenf $  &  $ \aeinsxsechs $  \\

\hline  \leitspaltezwei  &  $ \azweixeins $  &  $ \azweixzwei $  &  $ \azweixdrei $  &  $ \azweixvier $  &  $ \azweixfuenf $   &  $ \azweixsechs $  \\

\hline  \leitspaltedrei  &  $ \adreixeins $  &  $ \adreixzwei $  &  $ \adreixdrei $  &  $ \adreixvier $  &  $ \adreixfuenf $   &  $ \adreixsechs $   \\

\hline

\end{tabular}

\end{center} }

\newcommand{\tabelleleitdreixelf}{

\begin{center}

\begin{tabular}{|c||c|c|c|c|c|c|c|c|c|c|c|}

\hline  \leitzeilenull  &  \leitzeileeins  &  \leitzeilezwei  &  \leitzeiledrei  &  \leitzeilevier  &  \leitzeilefuenf   &  \leitzeilesechs  &  \leitzeilesieben  &  \leitzeileacht  &  \leitzeileneun  &  \leitzeilezehn &  \leitzeileelf    \\

\hline \hline  \leitspalteeins  &  $ \aeinsxeins $  &  $ \aeinsxzwei $  &  $ \aeinsxdrei $  &  $ \aeinsxvier $  &  $ \aeinsxfuenf $  &  $ \aeinsxsechs $  &  $ \aeinsxsieben $  &  $ \aeinsxacht $  &  $ \aeinsxneun $ &  $ \aeinsxzehn $  &  $ \aeinsxelf $  \\

\hline  \leitspaltezwei  &  $ \azweixeins $  &  $ \azweixzwei $  &  $ \azweixdrei $  &  $ \azweixvier $  &  $ \azweixfuenf $   &  $ \azweixsechs $  &  $ \azweixsieben $  &  $ \azweixacht $  &  $ \azweixneun $  &  $ \azweixzehn $ &  $ \azweixelf $     \\

\hline  \leitspaltedrei  &  $ \adreixeins $  &  $ \adreixzwei $  &  $ \adreixdrei $  &  $ \adreixvier $  &  $ \adreixfuenf $   &  $ \adreixsechs $  &  $ \adreixsieben $  &  $ \adreixacht $  &  $ \adreixneun $  &  $ \adreixzehn $     &  $ \adreixelf $   \\

\hline

\end{tabular}

\end{center} }

%Tabellen mit vier Zeilen

\newcommand{\tabelleleitvierxzwei}{

\begin{center}

\begin{tabular}{|c||c|c|}

\hline  \leitzeilenull  &  \leitzeileeins  &  \leitzeilezwei   \\

\hline \hline  \leitspalteeins  &  $ \aeinsxeins $  &  $ \aeinsxzwei $    \\

\hline  \leitspaltezwei  &  $ \azweixeins $  &  $ \azweixzwei $    \\

\hline  \leitspaltedrei  &  $ \adreixeins $  &  $ \adreixzwei $    \\

\hline  \leitspaltevier  &  $ \avierxeins $  &  $ \avierxzwei $    \\

\hline

\end{tabular}

\end{center} }

\newcommand{\tabelleleitvierxdrei}{

\begin{center}

\begin{tabular}{|c||c|c|c|}

\hline  \leitzeilenull  &  \leitzeileeins  &  \leitzeilezwei  &  \leitzeiledrei     \\

\hline \hline  \leitspalteeins  &  $ \aeinsxeins $  &  $ \aeinsxzwei $  &  $ \aeinsxdrei $     \\

\hline  \leitspaltezwei  &  $ \azweixeins $  &  $ \azweixzwei $  &  $ \azweixdrei $      \\

\hline  \leitspaltedrei  &  $ \adreixeins $  &  $ \adreixzwei $  &  $ \adreixdrei $    \\

\hline  \leitspaltevier  &  $ \avierxeins $  &  $ \avierxzwei $  &  $ \avierxdrei $    \\

\hline

\end{tabular}

\end{center} }

\newcommand{\tabelleleitvierxvier}{

\begin{center}

\begin{tabular}{|c||c|c|c|c|}

\hline  \leitzeilenull  &  \leitzeileeins  &  \leitzeilezwei  &  \leitzeiledrei  &  \leitzeilevier    \\

\hline \hline  \leitspalteeins  &  $ \aeinsxeins $  &  $ \aeinsxzwei $  &  $ \aeinsxdrei $  &  $ \aeinsxvier $   \\

\hline  \leitspaltezwei  &  $ \azweixeins $  &  $ \azweixzwei $  &  $ \azweixdrei $  &  $ \azweixvier $     \\

\hline  \leitspaltedrei  &  $ \adreixeins $  &  $ \adreixzwei $  &  $ \adreixdrei $  &  $ \adreixvier $   \\

\hline  \leitspaltevier  &  $ \avierxeins $  &  $ \avierxzwei $  &  $ \avierxdrei $  &  $ \avierxvier $   \\

\hline

\end{tabular}

\end{center} }

\newcommand{\tabelleleittextvierxzwei}{

	\begin{center}

		\begin{tabular}{|c||c|c|}

			\hline  \leitzeilenull  &  \leitzeileeins  &  \leitzeilezwei   \\

			\hline \hline  \leitspalteeins  &   \aeinsxeins   &   \aeinsxzwei    \\

			\hline  \leitspaltezwei  &  \azweixeins   &   \azweixzwei     \\

			\hline  \leitspaltedrei  &  \adreixeins   &   \adreixzwei     \\

			\hline  \leitspaltevier  &  \avierxeins   &   \avierxzwei     \\

			\hline

		\end{tabular}

	\end{center} }

\newcommand{\tabelleleitvierxacht}{

\begin{center}

\begin{tabular}{|c||c|c|c|c|c|c|c|c|}

\hline  \leitzeilenull  &  \leitzeileeins  &  \leitzeilezwei  &  \leitzeiledrei  &  \leitzeilevier  &  \leitzeilefuenf  &  \leitzeilesechs  &  \leitzeilesieben  &  \leitzeileacht  \\

\hline \hline  \leitspalteeins  &  $ \aeinsxeins $  &  $ \aeinsxzwei $  &  $ \aeinsxdrei $  &  $ \aeinsxvier $  &  $ \aeinsxfuenf $  &  $ \aeinsxsechs $  &  $ \aeinsxsieben $  &  $ \aeinsxacht $  \\

\hline  \leitspaltezwei  &  $ \azweixeins $  &  $ \azweixzwei $  &  $ \azweixdrei $  &  $ \azweixvier $   &  $ \azweixfuenf $  &  $ \azweixsechs $  &  $ \azweixsieben $  &  $ \azweixacht $   \\

\hline  \leitspaltedrei  &  $ \adreixeins $  &  $ \adreixzwei $  &  $ \adreixdrei $  &  $ \adreixvier $ &  $ \adreixfuenf $  &  $ \adreixsechs $  &  $ \adreixsieben $  &  $ \adreixacht $   \\

\hline  \leitspaltevier  &  $ \avierxeins $  &  $ \avierxzwei $  &  $ \avierxdrei $  &  $ \avierxvier $  &  $ \avierxfuenf $  &  $ \avierxsechs $  &  $ \avierxsieben $  &  $ \avierxacht $  \\

\hline

\end{tabular}

\end{center} }

%Tabellen mit fünf Zeilen

\newcommand{\tabelleleitfuenfxfuenf}{

\begin{center}

\begin{tabular}{|c||c|c|c|c|c|}

\hline  \leitzeilenull  &  \leitzeileeins  &  \leitzeilezwei  &  \leitzeiledrei  &  \leitzeilevier  &  \leitzeilefuenf   \\

\hline \hline  \leitspalteeins  &  $ \aeinsxeins $  &  $ \aeinsxzwei $  &  $ \aeinsxdrei $  &  $ \aeinsxvier $  &  $ \aeinsxfuenf $   \\

\hline  \leitspaltezwei  &  $ \azweixeins $  &  $ \azweixzwei $  &  $ \azweixdrei $  &  $ \azweixvier $  &  $ \azweixfuenf $     \\

\hline  \leitspaltedrei  &  $ \adreixeins $  &  $ \adreixzwei $  &  $ \adreixdrei $  &  $ \adreixvier $  &  $ \adreixfuenf $      \\

\hline  \leitspaltevier  &  $ \avierxeins $  &  $ \avierxzwei $  &  $ \avierxdrei $  &  $ \avierxvier $  &  $ \avierxfuenf $     \\

\hline  \leitspaltefuenf  &  $ \afuenfxeins $  &  $ \afuenfxzwei $  &  $ \afuenfxdrei $  &  $ \afuenfxvier $  &  $ \afuenfxfuenf $    \\

\hline

\end{tabular}

\end{center} }

%Tabellen mit sechs Zeilen

\newcommand{\tabelleleitsechsxdrei}{

\begin{center}

\begin{tabular}{|c||c|c|c|}

\hline  \leitzeilenull  &  \leitzeileeins  &  \leitzeilezwei  &  \leitzeiledrei    \\

\hline \hline  \leitspalteeins  &  $ \aeinsxeins $  &  $ \aeinsxzwei $  &  $ \aeinsxdrei $    \\

\hline  \leitspaltezwei  &  $ \azweixeins $  &  $ \azweixzwei $  &  $ \azweixdrei $      \\

\hline  \leitspaltedrei  &  $ \adreixeins $  &  $ \adreixzwei $  &  $ \adreixdrei $     \\

\hline  \leitspaltevier  &  $ \avierxeins $  &  $ \avierxzwei $  &  $ \avierxdrei $      \\

\hline  \leitspaltefuenf  &  $ \afuenfxeins $  &  $ \afuenfxzwei $  &  $ \afuenfxdrei $     \\

\hline  \leitspaltesechs  &  $ \asechsxeins $  &  $ \asechsxzwei $  &  $ \asechsxdrei $      \\

\hline

\end{tabular}

\end{center} }

\newcommand{\tabelleleitsechsxfuenf}{

\begin{center}

\begin{tabular}{|c||c|c|c|c|c|}

\hline  \leitzeilenull  &  \leitzeileeins  &  \leitzeilezwei  &  \leitzeiledrei  &  \leitzeilevier  &  \leitzeilefuenf   \\

\hline \hline  \leitspalteeins  &  $ \aeinsxeins $  &  $ \aeinsxzwei $  &  $ \aeinsxdrei $  &  $ \aeinsxvier $  &  $ \aeinsxfuenf $   \\

\hline  \leitspaltezwei  &  $ \azweixeins $  &  $ \azweixzwei $  &  $ \azweixdrei $  &  $ \azweixvier $  &  $ \azweixfuenf $     \\

\hline  \leitspaltedrei  &  $ \adreixeins $  &  $ \adreixzwei $  &  $ \adreixdrei $  &  $ \adreixvier $  &  $ \adreixfuenf $      \\

\hline  \leitspaltevier  &  $ \avierxeins $  &  $ \avierxzwei $  &  $ \avierxdrei $  &  $ \avierxvier $  &  $ \avierxfuenf $     \\

\hline  \leitspaltefuenf  &  $ \afuenfxeins $  &  $ \afuenfxzwei $  &  $ \afuenfxdrei $  &  $ \afuenfxvier $  &  $ \afuenfxfuenf $    \\

\hline  \leitspaltesechs  &  $ \asechsxeins $  &  $ \asechsxzwei $  &  $ \asechsxdrei $  &  $ \asechsxvier $  &  $ \asechsxfuenf $    \\

\hline

\end{tabular}

\end{center} }

\newcommand{\tabelleleitsechsxsechs}{

\begin{center}

\begin{tabular}{|c||c|c|c|c|c|c|}

\hline  \leitzeilenull  &  \leitzeileeins  &  \leitzeilezwei  &  \leitzeiledrei  &  \leitzeilevier  &  \leitzeilefuenf  &  \leitzeilesechs   \\

\hline \hline  \leitspalteeins  &  $ \aeinsxeins $  &  $ \aeinsxzwei $  &  $ \aeinsxdrei $  &  $ \aeinsxvier $  &  $ \aeinsxfuenf $  &  $ \aeinsxsechs $ \\

\hline  \leitspaltezwei  &  $ \azweixeins $  &  $ \azweixzwei $  &  $ \azweixdrei $  &  $ \azweixvier $  &  $ \azweixfuenf $  &  $ \azweixsechs $   \\

\hline  \leitspaltedrei  &  $ \adreixeins $  &  $ \adreixzwei $  &  $ \adreixdrei $  &  $ \adreixvier $  &  $ \adreixfuenf $   &  $ \adreixsechs $   \\

\hline  \leitspaltevier  &  $ \avierxeins $  &  $ \avierxzwei $  &  $ \avierxdrei $  &  $ \avierxvier $  &  $ \avierxfuenf $   &  $ \avierxsechs $  \\

\hline  \leitspaltefuenf  &  $ \afuenfxeins $  &  $ \afuenfxzwei $  &  $ \afuenfxdrei $  &  $ \afuenfxvier $  &  $ \afuenfxfuenf $ &  $ \afuenfxsechs $   \\

\hline  \leitspaltesechs  &  $ \asechsxeins $  &  $ \asechsxzwei $  &  $ \asechsxdrei $  &  $ \asechsxvier $  &  $ \asechsxfuenf $  &  $ \asechsxsechs $  \\

\hline

\end{tabular}

\end{center} }

%Tabellen mit sieben Zeilen

\newcommand{\tabelleleitsiebenxsieben}{

\begin{center}

\begin{tabular}{|c||c|c|c|c|c|c|c|}

\hline  \leitzeilenull  &  \leitzeileeins  &  \leitzeilezwei  &  \leitzeiledrei  &  \leitzeilevier  &  \leitzeilefuenf  &  \leitzeilesechs &  \leitzeilesieben  \\

\hline \hline  \leitspalteeins  &  $ \aeinsxeins $  &  $ \aeinsxzwei $  &  $ \aeinsxdrei $  &  $ \aeinsxvier $  &  $ \aeinsxfuenf $  &  $ \aeinsxsechs $ &  $ \aeinsxsieben $ \\

\hline  \leitspaltezwei  &  $ \azweixeins $  &  $ \azweixzwei $  &  $ \azweixdrei $  &  $ \azweixvier $  &  $ \azweixfuenf $  &  $ \azweixsechs $  &  $ \azweixsieben $ \\

\hline  \leitspaltedrei  &  $ \adreixeins $  &  $ \adreixzwei $  &  $ \adreixdrei $  &  $ \adreixvier $  &  $ \adreixfuenf $   &  $ \adreixsechs $ &  $ \adreixsieben $  \\

\hline  \leitspaltevier  &  $ \avierxeins $  &  $ \avierxzwei $  &  $ \avierxdrei $  &  $ \avierxvier $  &  $ \avierxfuenf $   &  $ \avierxsechs $ &  $ \avierxsieben $ \\

\hline  \leitspaltefuenf  &  $ \afuenfxeins $  &  $ \afuenfxzwei $  &  $ \afuenfxdrei $  &  $ \afuenfxvier $  &  $ \afuenfxfuenf $ &  $ \afuenfxsechs $ &  $ \afuenfxsieben $  \\

\hline  \leitspaltesechs  &  $ \asechsxeins $  &  $ \asechsxzwei $  &  $ \asechsxdrei $  &  $ \asechsxvier $  &  $ \asechsxfuenf $  &  $ \asechsxsechs $ &  $ \asechsxsieben $ \\

\hline  \leitspaltesieben  &  $ \asiebenxeins $  &  $ \asiebenxzwei $  &  $ \asiebenxdrei $  &  $ \asiebenxvier $  &  $ \asiebenxfuenf $  &  $ \asiebenxsechs $ &  $ \asiebenxsieben $ \\

\hline

\end{tabular}

\end{center} }

%Tabellen mit acht Zeilen

\newcommand{\tabelleleitachtxdrei}{

\begin{center}

\begin{tabular}{|c||c|c|c|}

\hline  \leitzeilenull  &  \leitzeileeins  &  \leitzeilezwei  &  \leitzeiledrei     \\

\hline \hline  \leitspalteeins  &  $ \aeinsxeins $  &  $ \aeinsxzwei $  &  $ \aeinsxdrei $    \\

\hline  \leitspaltezwei  &  $ \azweixeins $  &  $ \azweixzwei $  &  $ \azweixdrei $     \\

\hline  \leitspaltedrei  &  $ \adreixeins $  &  $ \adreixzwei $  &  $ \adreixdrei $    \\

\hline  \leitspaltevier  &  $ \avierxeins $  &  $ \avierxzwei $  &  $ \avierxdrei $      \\

\hline  \leitspaltefuenf  &  $ \afuenfxeins $  &  $ \afuenfxzwei $  &  $ \afuenfxdrei $     \\

\hline  \leitspaltesechs  &  $ \asechsxeins $  &  $ \asechsxzwei $  &  $ \asechsxdrei $      \\

\hline  \leitspaltesieben  &  $ \asiebenxeins $  &  $ \asiebenxzwei $  &  $ \asiebenxdrei $     \\

\hline  \leitspalteacht  &  $ \aachtxeins $  &  $ \aachtxzwei $  &  $ \aachtxdrei $      \\

\hline

\end{tabular}

\end{center} }

%Tabellen mit zehn Zeilen

\newcommand{\tabelleleitzehnxdrei}{

\begin{center}

\begin{tabular}{|c||c|c|c|}

\hline  \leitzeilenull  &  \leitzeileeins  &  \leitzeilezwei  &  \leitzeiledrei     \\

\hline \hline  \leitspalteeins  &  $ \aeinsxeins $  &  $ \aeinsxzwei $  &  $ \aeinsxdrei $    \\

\hline  \leitspaltezwei  &  $ \azweixeins $  &  $ \azweixzwei $  &  $ \azweixdrei $     \\

\hline  \leitspaltedrei  &  $ \adreixeins $  &  $ \adreixzwei $  &  $ \adreixdrei $    \\

\hline  \leitspaltevier  &  $ \avierxeins $  &  $ \avierxzwei $  &  $ \avierxdrei $      \\

\hline  \leitspaltefuenf  &  $ \afuenfxeins $  &  $ \afuenfxzwei $  &  $ \afuenfxdrei $     \\

\hline  \leitspaltesechs  &  $ \asechsxeins $  &  $ \asechsxzwei $  &  $ \asechsxdrei $      \\

\hline  \leitspaltesieben  &  $ \asiebenxeins $  &  $ \asiebenxzwei $  &  $ \asiebenxdrei $     \\

\hline  \leitspalteacht  &  $ \aachtxeins $  &  $ \aachtxzwei $  &  $ \aachtxdrei $      \\

\hline  \leitspalteneun  &  $ \aneunxeins $  &  $ \aneunxzwei $  &  $ \aneunxdrei $      \\

\hline  \leitspaltezehn  &  $ \azehnxeins $  &  $ \azehnxzwei $  &  $ \azehnxdrei $      \\

\hline

\end{tabular}

\end{center} }

%Tabellen mit elf Zeilen

\newcommand{\tabelleleitelfxvier}{

\begin{center}

\begin{tabular}{|c||c|c|c|c|}

\hline  \leitzeilenull  &  \leitzeileeins  &  \leitzeilezwei  &  \leitzeiledrei  &  \leitzeilevier  \\

\hline \hline  \leitspalteeins  &  $ \aeinsxeins $  &  $ \aeinsxzwei $  &  $ \aeinsxdrei $  &  $ \aeinsxvier $  \\

\hline  \leitspaltezwei  &  $ \azweixeins $  &  $ \azweixzwei $  &  $ \azweixdrei $  &  $ \azweixvier $   \\

\hline  \leitspaltedrei  &  $ \adreixeins $  &  $ \adreixzwei $  &  $ \adreixdrei $  &  $ \adreixvier $   \\

\hline  \leitspaltevier  &  $ \avierxeins $  &  $ \avierxzwei $  &  $ \avierxdrei $  &  $ \avierxvier $   \\

\hline  \leitspaltefuenf  &  $ \afuenfxeins $  &  $ \afuenfxzwei $  &  $ \afuenfxdrei $  &  $ \afuenfxvier $   \\

\hline  \leitspaltesechs  &  $ \asechsxeins $  &  $ \asechsxzwei $  &  $ \asechsxdrei $  &  $ \asechsxvier $   \\

\hline  \leitspaltesieben  &  $ \asiebenxeins $  &  $ \asiebenxzwei $  &  $ \asiebenxdrei $  &  $ \asiebenxvier $    \\

\hline  \leitspalteacht  &  $ \aachtxeins $  &  $ \aachtxzwei $  &  $ \aachtxdrei $  &  $ \aachtxvier $   \\

\hline  \leitspalteneun  &  $ \aneunxeins $  &  $ \aneunxzwei $  &  $ \aneunxdrei $  &  $ \aneunxvier $   \\

\hline  \leitspaltezehn  &  $ \azehnxeins $  &  $ \azehnxzwei $  &  $ \azehnxdrei $  &  $ \azehnxvier $   \\

\hline  \leitspalteelf  &  $ \aelfxeins $  &  $ \aelfxzwei $  &  $ \aelfxdrei $  &  $ \aelfxvier $   \\

\hline

\end{tabular}

\end{center} }

%Tabellen mit zwölf Zeilen

\newcommand{\tabelleleitzwoelfxsieben}{

\begin{center}

\begin{tabular}{|c||c|c|c|c|c|c|}

\hline  \leitzeilenull  &  \leitzeileeins  &  \leitzeilezwei  &  \leitzeiledrei  &  \leitzeilevier &  \leitzeilefuenf  &  \leitzeilesechs &  \leitzeilesieben \\

\hline \hline  \leitspalteeins  &  $ \aeinsxeins $  &  $ \aeinsxzwei $  &  $ \aeinsxdrei $  &  $ \aeinsxvier $  &  $ \aeinsxfuenf $  &  $ \aeinsxsechs $  &  $ \aeinsxsieben $  \\

\hline  \leitspaltezwei  &  $ \azweixeins $  &  $ \azweixzwei $  &  $ \azweixdrei $  &  $ \azweixvier $  &  $ \azweixfuenf $  &  $ \azweixsechs $  &  $ \azweixsieben $   \\

\hline  \leitspaltedrei  &  $ \adreixeins $  &  $ \adreixzwei $  &  $ \adreixdrei $  &  $ \adreixvier $  &  $ \adreixfuenf $  &  $ \adreixsechs $  &  $ \adreixsieben $   \\

\hline  \leitspaltevier  &  $ \avierxeins $  &  $ \avierxzwei $  &  $ \avierxdrei $  &  $ \avierxvier $   &  $ \avierxfuenf $  &  $ \avierxsechs $  &  $ \avierxsieben $  \\

\hline  \leitspaltefuenf  &  $ \afuenfxeins $  &  $ \afuenfxzwei $  &  $ \afuenfxdrei $  &  $ \afuenfxvier $ &  $ \afuenfxfuenf $  &  $ \afuenfxsechs $  &  $ \afuenfxsieben $  \\

\hline  \leitspaltesechs  &  $ \asechsxeins $  &  $ \asechsxzwei $  &  $ \asechsxdrei $  &  $ \asechsxvier $  &  $ \asechsxfuenf $  &  $ \asechsxsechs $  &  $ \asechsxsieben $  \\

\hline  \leitspaltesieben  &  $ \asiebenxeins $  &  $ \asiebenxzwei $  &  $ \asiebenxdrei $  &  $ \asiebenxvier $  &  $ \asiebenxfuenf $  &  $ \asiebenxsechs $  &  $ \asiebenxsieben $   \\

\hline  \leitspalteacht  &  $ \aachtxeins $  &  $ \aachtxzwei $  &  $ \aachtxdrei $  &  $ \aachtxvier $  &  $ \aachtxfuenf $  &  $ \aachtxsechs $  &  $ \aachtxsieben $ \\

\hline  \leitspalteneun  &  $ \aneunxeins $  &  $ \aneunxzwei $  &  $ \aneunxdrei $  &  $ \aneunxvier $  &  $ \aneunxfuenf $  &  $ \aneunxsechs $  &  $ \aneunxsieben $  \\

\hline  \leitspaltezehn  &  $ \azehnxeins $  &  $ \azehnxzwei $  &  $ \azehnxdrei $  &  $ \azehnxvier $ &  $ \azehnxfuenf $  &  $ \azehnxsechs $  &  $ \azehnxsieben $   \\

\hline  \leitspalteelf  &  $ \aelfxeins $  &  $ \aelfxzwei $  &  $ \aelfxdrei $  &  $ \aelfxvier $  &  $ \aelfxfuenf $  &  $ \aelfxsechs $  &  $ \aelfxsieben $  \\

\hline  \leitspaltezwoelf  &  $ \azwoelfxeins $  &  $ \azwoelfxzwei $  &  $ \azwoelfxdrei $  &  $ \zwoelfxvier $  &  $ \azwoelfxfuenf $  &  $ \azwoelfxsechs $  &  $ \zwoelfxsieben $ \\

\hline

\end{tabular}

\end{center} }

%Tabellen mit sechzehn Zeilen

\newcommand{\tabelleleitsechzehnxzwei}{

\begin{center}

\begin{tabular}{|c||c|c|}

\hline  \leitzeilenull  &  \leitzeileeins  &  \leitzeilezwei   \\

\hline \hline  \leitspalteeins  &  $ \aeinsxeins $  &  $ \aeinsxzwei $ \\

\hline  \leitspaltezwei  &  $ \azweixeins $  &  $ \azweixzwei $  \\

\hline  \leitspaltedrei  &  $ \adreixeins $  &  $ \adreixzwei $  \\

\hline  \leitspaltevier  &  $ \avierxeins $  &  $ \avierxzwei $  \\

\hline  \leitspaltefuenf  &  $ \afuenfxeins $  &  $ \afuenfxzwei $  \\

\hline  \leitspaltesechs  &  $ \asechsxeins $  &  $ \asechsxzwei $  \\

\hline  \leitspaltesieben  &  $ \asiebenxeins $  &  $ \asiebenxzwei $  \\

\hline  \leitspalteacht  &  $ \aachtxeins $  &  $ \aachtxzwei $  \\

\hline  \leitspalteneun  &  $ \aneunxeins $  &  $ \aneunxzwei $   \\

\hline  \leitspaltezehn  &  $ \azehnxeins $  &  $ \azehnxzwei $  \\

\hline  \leitspalteelf  &  $ \aelfxeins $  &  $ \aelfxzwei $   \\

\hline  \leitspaltezwoelf  &  $ \azwoelfxeins $  &  $ \azwoelfxzwei $   \\

\hline  \leitspaltedreizehn  &  $ \adreizehnxeins $  &  $ \adreizehnxzwei $   \\

\hline  \leitspaltevierzehn  &  $ \avierzehnxeins $  &  $ \avierzehnxzwei $  \\

\hline  \leitspaltefuenfzehn  &  $ \afuenfzehnxeins $  &  $ \afuenfzehnxzwei $   \\

\hline  \leitspaltesechzehn  &  $ \asechzehnxeins $  &  $ \asechzehnxzwei $   \\

\hline

\end{tabular}

\end{center} }

\newcommand{\tabelleleitsechzehnxsechs}{

\begin{center}

\begin{tabular}{|c||c|c|c|c|c|c|}

\hline  \leitzeilenull  &  \leitzeileeins  &  \leitzeilezwei  &  \leitzeiledrei  &  \leitzeilevier  &  \leitzeilefuenf  &  \leitzeilesechs  \\

\hline \hline  \leitspalteeins  &  $ \aeinsxeins $  &  $ \aeinsxzwei $  &  $ \aeinsxdrei $  &  $ \aeinsxvier $  &  $ \aeinsxfuenf $  &  $ \aeinsxsechs $  \\

\hline  \leitspaltezwei  &  $ \azweixeins $  &  $ \azweixzwei $  &  $ \azweixdrei $  &  $ \azweixvier $  &  $ \azweixfuenf $  &  $ \azweixsechs $  \\

\hline  \leitspaltedrei  &  $ \adreixeins $  &  $ \adreixzwei $  &  $ \adreixdrei $  &  $ \adreixvier $  &  $ \adreixfuenf $  &  $ \adreixsechs $  \\

\hline  \leitspaltevier  &  $ \avierxeins $  &  $ \avierxzwei $  &  $ \avierxdrei $  &  $ \avierxvier $  &  $ \avierxfuenf $  &  $ \avierxsechs $  \\

\hline  \leitspaltefuenf  &  $ \afuenfxeins $  &  $ \afuenfxzwei $  &  $ \afuenfxdrei $  &  $ \afuenfxvier $  &  $ \afuenfxfuenf $  &  $ \afuenfxsechs $  \\

\hline  \leitspaltesechs  &  $ \asechsxeins $  &  $ \asechsxzwei $  &  $ \asechsxdrei $  &  $ \asechsxvier $  &  $ \asechsxfuenf $  &  $ \asechsxsechs $  \\

\hline  \leitspaltesieben  &  $ \asiebenxeins $  &  $ \asiebenxzwei $  &  $ \asiebenxdrei $  &  $ \asiebenxvier $  &  $ \asiebenxfuenf $  &  $ \asiebenxsechs $  \\

\hline  \leitspalteacht  &  $ \aachtxeins $  &  $ \aachtxzwei $  &  $ \aachtxdrei $  &  $ \aachtxvier $  &  $ \aachtxfuenf $  &  $ \aachtxsechs $  \\

\hline  \leitspalteneun  &  $ \aneunxeins $  &  $ \aneunxzwei $  &  $ \aneunxdrei $  &  $ \aneunxvier $  &  $ \aneunxfuenf $  &  $ \aneunxsechs $  \\

\hline  \leitspaltezehn  &  $ \azehnxeins $  &  $ \azehnxzwei $  &  $ \azehnxdrei $  &  $ \azehnxvier $  &  $ \azehnxfuenf $  &  $ \azehnxsechs $  \\

\hline  \leitspalteelf  &  $ \aelfxeins $  &  $ \aelfxzwei $  &  $ \aelfxdrei $  &  $ \aelfxvier $  &  $ \aelfxfuenf $  &  $ \aelfxsechs $  \\

\hline  \leitspaltezwoelf  &  $ \azwoelfxeins $  &  $ \azwoelfxzwei $  &  $ \azwoelfxdrei $  &  $ \azwoelfxvier $  &  $ \azwoelfxfuenf $  &  $ \azwoelfxsechs $  \\

\hline  \leitspaltedreizehn  &  $ \adreizehnxeins $  &  $ \adreizehnxzwei $  &  $ \adreizehnxdrei $  &  $ \adreizehnxvier $  &  $ \adreizehnxfuenf $  &  $ \adreizehnxsechs $  \\

\hline  \leitspaltevierzehn  &  $ \avierzehnxeins $  &  $ \avierzehnxzwei $  &  $ \avierzehnxdrei $  &  $ \avierzehnxvier $  &  $ \avierzehnxfuenf $  &  $ \avierzehnxsechs $  \\

\hline  \leitspaltefuenfzehn  &  $ \afuenfzehnxeins $  &  $ \afuenfzehnxzwei $  &  $ \afuenfzehnxdrei $  &  $ \afuenfzehnxvier $  &  $ \afuenfzehnxfuenf $  &  $ \afuenfzehnxsechs $  \\

\hline  \leitspaltesechzehn  &  $ \asechzehnxeins $  &  $ \asechzehnxzwei $  &  $ \asechzehnxdrei $  &  $ \asechzehnxvier $  &  $ \asechzehnxfuenf $  &  $ \asechzehnxsechs $  \\

\hline

\end{tabular}

\end{center} }

%Tabellen mit achtzehn Zeilen

\newcommand{\tabelleleitachtzehnxneun}{

\begin{center}

\begin{tabular}{|c||c|c|c|c|c|c|c|c|c|}

\hline  \leitzeilenull  &  \leitzeileeins  &  \leitzeilezwei  &  \leitzeiledrei  &  \leitzeilevier  &  \leitzeilefuenf  &  \leitzeilesechs &  \leitzeilesieben  &  \leitzeileacht  &  \leitzeileneun \\

\hline \hline  \leitspalteeins  &  $ \aeinsxeins $  &  $ \aeinsxzwei $  &  $ \aeinsxdrei $  &  $ \aeinsxvier $  &  $ \aeinsxfuenf $  &  $ \aeinsxsechs $ &  $ \aeinsxsieben $  &  $ \aeinsxacht $  &  $ \aeinsxneun $  \\

\hline  \leitspaltezwei  &  $ \azweixeins $  &  $ \azweixzwei $  &  $ \azweixdrei $  &  $ \azweixvier $  &  $ \azweixfuenf $  &  $ \azweixsechs $ &  $ \azweixsieben $  &  $ \azweixacht $  &  $ \azweixneun $  \\

\hline  \leitspaltedrei  &  $ \adreixeins $  &  $ \adreixzwei $  &  $ \adreixdrei $  &  $ \adreixvier $  &  $ \adreixfuenf $  &  $ \adreixsechs $ &  $ \adreixsieben $  &  $ \adreixacht $  &  $ \adreixneun $  \\

\hline  \leitspaltevier  &  $ \avierxeins $  &  $ \avierxzwei $  &  $ \avierxdrei $  &  $ \avierxvier $  &  $ \avierxfuenf $  &  $ \avierxsechs $  &  $ \avierxsieben $  &  $ \avierxacht $  &  $ \avierxneun $ \\

\hline  \leitspaltefuenf  &  $ \afuenfxeins $  &  $ \afuenfxzwei $  &  $ \afuenfxdrei $  &  $ \afuenfxvier $  &  $ \afuenfxfuenf $  &  $ \afuenfxsechs $ &  $ \afuenfxsieben $  &  $ \afuenfxacht $  &  $ \afuenfxneun $  \\

\hline  \leitspaltesechs  &  $ \asechsxeins $  &  $ \asechsxzwei $  &  $ \asechsxdrei $  &  $ \asechsxvier $  &  $ \asechsxfuenf $  &  $ \asechsxsechs $ &  $ \asechsxsieben $  &  $ \asechsxacht $  &  $ \asechsxneun $  \\

\hline  \leitspaltesieben  &  $ \asiebenxeins $  &  $ \asiebenxzwei $  &  $ \asiebenxdrei $  &  $ \asiebenxvier $  &  $ \asiebenxfuenf $  &  $ \asiebenxsechs $ &  $ \asiebenxsieben $  &  $ \asiebenxacht $  &  $ \asiebenxneun $  \\

\hline  \leitspalteacht  &  $ \aachtxeins $  &  $ \aachtxzwei $  &  $ \aachtxdrei $  &  $ \aachtxvier $  &  $ \aachtxfuenf $  &  $ \aachtxsechs $ &  $ \aachtxsieben $  &  $ \aachtxacht $  &  $ \aachtxneun $  \\

\hline  \leitspalteneun  &  $ \aneunxeins $  &  $ \aneunxzwei $  &  $ \aneunxdrei $  &  $ \aneunxvier $  &  $ \aneunxfuenf $  &  $ \aneunxsechs $ &  $ \aneunxsieben $  &  $ \aneunxacht $  &  $ \aneunxneun $  \\

\hline  \leitspaltezehn  &  $ \azehnxeins $  &  $ \azehnxzwei $  &  $ \azehnxdrei $  &  $ \azehnxvier $  &  $ \azehnxfuenf $  &  $ \azehnxsechs $ &  $ \azehnxsieben $  &  $ \azehnxacht $  &  $ \azehnxneun $  \\

\hline  \leitspalteelf  &  $ \aelfxeins $  &  $ \aelfxzwei $  &  $ \aelfxdrei $  &  $ \aelfxvier $  &  $ \aelfxfuenf $  &  $ \aelfxsechs $ &  $ \aelfxsieben $  &  $ \aelfxacht $  &  $ \aelfxneun $  \\

\hline  \leitspaltezwoelf  &  $ \azwoelfxeins $  &  $ \azwoelfxzwei $  &  $ \azwoelfxdrei $  &  $ \azwoelfxvier $  &  $ \azwoelfxfuenf $  &  $ \azwoelfxsechs $ &  $ \azwoelfxsieben $  &  $ \azwoelfxacht $  &  $ \azwoelfxneun $  \\

\hline  \leitspaltedreizehn  &  $ \adreizehnxeins $  &  $ \adreizehnxzwei $  &  $ \adreizehnxdrei $  &  $ \adreizehnxvier $  &  $ \adreizehnxfuenf $  &  $ \adreizehnxsechs $ &  $ \adreizehnxsieben $  &  $ \adreizehnxacht $  &  $ \adreizehnxneun $  \\

\hline  \leitspaltevierzehn  &  $ \avierzehnxeins $  &  $ \avierzehnxzwei $  &  $ \avierzehnxdrei $  &  $ \avierzehnxvier $  &  $ \avierzehnxfuenf $  &  $ \avierzehnxsechs $  &  $ \avierzehnxsieben $  &  $ \avierzehnxacht $  &  $ \avierzehnxneun $ \\

\hline  \leitspaltefuenfzehn  &  $ \afuenfzehnxeins $  &  $ \afuenfzehnxzwei $  &  $ \afuenfzehnxdrei $  &  $ \afuenfzehnxvier $  &  $ \afuenfzehnxfuenf $  &  $ \afuenfzehnxsechs $ &  $ \afuenfzehnxsieben $  &  $ \afuenfzehnxacht $  &  $ \afuenfzehnxneun $  \\

\hline  \leitspaltesechzehn  &  $ \asechzehnxeins $  &  $ \asechzehnxzwei $  &  $ \asechzehnxdrei $  &  $ \asechzehnxvier $  &  $ \asechzehnxfuenf $  &  $ \asechzehnxsechs $ &  $ \asechzehnxsieben $  &  $ \asechzehnxacht $  &  $ \asechzehnxneun $  \\

\hline  \leitspaltesiebzehn  &  $ \asiebzehnxeins $  &  $ \asiebzehnxzwei $  &  $ \asiebzehnxdrei $  &  $ \asiebzehnxvier $  &  $ \asiebzehnxfuenf $  &  $ \asiebzehnxsechs $ &  $ \asiebzehnxsieben $  &  $ \asiebzehnxacht $  &  $ \asiebzehnxneun $  \\

\hline  \leitspalteachtzehn  &  $ \aachtzehnxeins $  &  $ \aachtzehnxzwei $  &  $ \aachtzehnxdrei $  &  $ \aachtzehnxvier $  &  $ \aachtzehnxfuenf $  &  $ \aachtzehnxsechs $ &  $ \aachtzehnxsieben $  &  $ \aachtzehnxacht $  &  $ \aachtzehnxneun $  \\

\hline

\end{tabular}

\end{center} }

\newcommand{\bruchhilfealign}{\cr & & }

\newcommand{\bruchhilfedisp}{\]  \[ }

\newcommand{\bruchhilfe}{ \- }

%[[Projekt:Semantische Vorlagen/Notlösungen/Latex]]

%
%Bei verschachtelten ifthen-Befehlen
%kann es Probleme geben, daher
%

\newcommand{\faktuebergangpos}[1]{ \faktuebergangskip  {\hspace{-0,55cm} } {#1} }

%Ende des Grundvorspannes

%Globale Strukturierungen

\newcommand{\seitenueberschrift}[1]{\seitenueberschriftvorskip
\begin{center} \large{ #1} \end{center} \seitenueberschriftnachskip}

\newcommand{\zwischenueberschrift}[1]{\section*{#1} }

\newcommand{\zwischenueberschriftaufgaben}[1]{ \zwischenueberschriftvorskip \begin{center} {\bf #1} \end{center} \zwischenueberschriftnachskip}

\newcommand{\unterueberschrift}[1]{ \unterueberschriftvorskip \hspace{0.5cm} {\em #1}  \unterueberschriftnachskip}

\renewcommand{\definitionbenennung}[1]{}

%Grund:Bei Definitionen im Normaltext
%nicht nötig, da Begriff früh kommt

\renewcommand{\beispielbenennung}[1]{}

\renewcommand{\bemerkungbenennung}[1]{}

%Umbrüche bei großer Schrift (entfällt)

%[[Projekt:Semantische Vorlagen/Mathematischer Modus/Umbrüche/Latex]]

\newcommand{\mathdisplaybruch}{}

\newcommand{\mathdisplaybruchinklammer}{}

\newcommand{\mathl}[2]{$#1$#2}

\newcommand{\mathlk}[2]{$#1$#2}

%Nicht durch semantische Vorlage produzierte Befehle

\newcommand{\mathdisplaybruchbruch}{ \]  \[ }

%Gestaltung der vertikalen und horizontalen Abstände

%[[Projekt:Semantische Vorlagen/Latex/Vorlesungsskript/Skips]]

%Abstände bei Überschriften

\newcommand{\seitenueberschriftvorskip}{\par \bigskip \bigskip \bigskip }

\newcommand{\seitenueberschriftnachskip}{\par \bigskip}

\newcommand{\zwischenueberschriftvorskip}{\par \bigskip}

\newcommand{\zwischenueberschriftnachskip}{\par \smallskip}

\newcommand{\unterueberschriftvorskip}{\par \bigskip}

\newcommand{\unterueberschriftnachskip}{\par \smallskip}

%Aufzählungen und Listen

\newcommand{\listskip}{\smallskip}

%Abstände bei mathematischen Strukturen

\newcommand{\aufgabevorskip}{\bigskip}

\newcommand{\aufgabepunktskip}{\par}

\newcommand{\aufgabenachskip}{}

\newcommand{\aufgabeloesungskip}{\par \bigskip}

\newcommand{\beispielvorskip}{}

\newcommand{\beispielnachskip}{}

\newcommand{\beispielbenennungnachskip}{}

\newcommand{\bemerkungvorskip}{}

\newcommand{\bemerkungnachskip}{}

\newcommand{\bemerkungbenennungnachskip}{}

\newcommand{\definitionvorskip}{}

\newcommand{\definitionnachskip}{}

\newcommand{\definitionbenennungnachskip}{}

\newcommand{\faktvorskip}{}

\newcommand{\faktnachskip}{}

\newcommand{\faktbenennungvorskip}{}

\newcommand{\faktbenennungnachskip}{}

\newcommand{\faktsituationskip}{}

\newcommand{\faktvoraussetzungskip}{}

\newcommand{\faktuebergangskip}{}

\newcommand{\faktfolgerungskip}{}

\newcommand{\faktzusatzskip}{}

%Sonstiges

\newcommand{\deckblattabstand}{\vspace{2cm}}

\newcommand{\gesamtskriptnewpage}{ }

\newcommand{\einzeltextnewpage}{ }

%[[Projekt:Semantische Vorlagen/Latex/Vorlesungsskript/Horizontale Abstände]]

\newcommand{\leerzeichen}{ }

%Seitengestaltung

%[[Projekt:Semantische Vorlagen/Latex/Vorlesungsskript/Seitengestaltung]]

%Einzelne Abstände

% page geometry supplied by reader wrapper
% page geometry supplied by reader wrapper

% page geometry supplied by reader wrapper
% page geometry supplied by reader wrapper

% page geometry supplied by reader wrapper
% page geometry supplied by reader wrapper

\setlength{\parindent}{0cm}

\setlength{\parskip}{1ex}

%Bildgestaltung

%[[Projekt:Semantische Vorlagen/Latex/Vorlesungsskript/Bildgestaltung]]

\newcommand{\bild}[1]{\vspace{4mm} #1}

\newcommand{\bildtext}[1]{\begin{center} {\small #1 } \end{center} }

%Lokale Einstellungen

%\input{Lokalername}

\newcommand{\bildeinlesung}[1]{../authority/media/#1}

\renewcommand{\bildeinlesung}[2]{../authority/media/#1.#2}

\newcommand{\bildeinlesungpng}[2]{../authority/media/#1.png}

\newcommand{\bildeinlesungPNG}[2]{../authority/media/#1.png}

\newcommand{\bildeinlesungjpg}[2]{../authority/media/#1.jpg}

\newcommand{\bildeinlesungJPG}[2]{../authority/media/#1.jpg}

\newcommand{\bildeinlesungjpeg}[2]{../authority/media/#1.jpg}

\newcommand{\bildeinlesungsvg}[2]{generated/media/#1.png}

\newcommand{\bildeinlesunggif}[2]{../authority/media/#1.png}

\newcommand{\bildeinlesungGIF}[2]{../authority/media/#1.png}

\newcommand{\bildeinlesungxcf}[2]{../authority/media/#1.png}

%Klausurgestaltung

%[[Projekt:Semantische Vorlagen/Klausurgestaltung/Latex]]

\newcommand{\fachbereich}{Fachbereich}

\newcommand{\dozent}{Dozent}

\newcommand{\klausurdatum}{Datum}

\newcommand{\klausurgebiet}{Mathematik}

\newcommand{\klausurtyp}{Klausur}

\newcommand{\klausurmodus}{Dauer: Zehn Minuten Einlesezeit und zwei volle Stunden Bearbeitungszeit. Es sind keine Hilfsmittel erlaubt. Alle Antworten sind zu begr\"unden.

Es gibt insgesamt 64 Punkte. Es gilt die Sockelregelung, d.h. die Bewertung pro Aufgabe(nteil) beginnt in der Regel bei der halben Punktzahl. Zum Bestehen braucht man $16$ Punkte, ab $32$ Punkten gibt es eine Eins.

Tragen Sie auf dem Deckblatt und jedem weiteren Blatt Ihren Namen und Ihre
Matrikelnummer leserlich ein. Viel Erfolg!}

\newcommand{\klausurname}{
\renewcommand{\arraystretch}{1.5}
\begin{tabular}{@{}p{3.2cm}p{10cm}}
Name, Vorname: &
..................................................................................\\
Matrikelnummer: &
..................................................................................\\
\end{tabular}
}

\newcommand{\klausurvorspann}[5]{\textbf{\fachbereich}\hfill \textbf{\klausurdatum}
\\
\textbf{\dozent} \hfill
\vspace{4ex}

\begin{center}
{\bfseries{\large \klausurgebiet \vspace{2ex}

\klausurtyp} }
\end{center}

\vspace{2ex}

\klausurmodus

\vspace{2ex}

\klausurname
}

\newcommand{\klausurnote}{\begin{tabular}{@{}p{3.2cm}p{10cm}} Note: & \end{tabular} }

\newcommand{\klausurtaeuschungwarnung}{ \vspace{2ex}{Mir ist bewu\ss t, dass bei einem schweren T\"auschungsversuch s\"amtliche Pr\"u\-fungsvorleis\-tungen in diesem Kurs verfallen. Ein schwerer T\"auschungsversuch liegt insbesondere beim Einsatz von Literatur und von elektronischen Hilfsmitteln vor. \vspace{1ex} } }

\newcommand{\klausureinverstaendnis}{
\vspace{2ex}
Ich erkl\"are mich durch meine Unterschrift einverstanden, dass mein Klausurergebnis mit meiner Matrikelnummer im Internet bekanntgegeben wird.
\vspace{1ex}

\begin{tabular}{@{}p{3.2cm}p{10cm}}
Unterschrift: &
.....................................................................................
\end{tabular}
}

%Variante für Testklausur

\newcommand{\testklausurvorspann}[5]{\textbf{\fachbereich}\hfill \textbf{\klausurdatum} \\ \textbf{\dozent} \hfill \vspace{4ex}

\begin{center} {\bfseries{\large \klausurgebiet \vspace{2ex}

\klausurtyp} } \end{center}

\vspace{2ex}

\testklausurmodus

\vspace{3ex}

\klausurname }

\newcommand{\testklausurmodus}{Dauer: Zehn Minuten Einlesezeit und zwei volle Stunden Bearbeitungszeit.

Es sind keine Hilfsmittel erlaubt.

Alle Antworten sind zu begr\"unden.

Es gibt insgesamt 64 Punkte. Es gilt die Sockelregelung, d.h. die Bewertung pro Aufgabe(nteil) beginnt in der Regel bei der halben Punktzahl.

Zur Orientierung: Zum Bestehen braucht man $16$ Punkte, ab $32$ Punkten gibt es eine Eins.

Tragen Sie auf dem Deckblatt und jedem weiteren Blatt Ihren Namen und Ihre Matrikelnummer leserlich ein.

Viel Erfolg!}

%\renewcommand{\inputaufgabegibtloesung}[4]{\aufgabevorskip \begin{aufgabe} \punkte {#1} \aufgabepunktskip #2 #3 #4\end{aufgabe} \aufgabenachskip}

\newcommand{\klausurmitloesungenvorspann}[5]{\textbf{\fachbereich}\hfill \textbf{\klausurdatum} \\ \textbf{\dozent} \hfill \vspace{4ex}

\begin{center} {\bfseries{\large \klausurgebiet \vspace{2ex}

\klausurtyp \hspace{0mm} mit L\"osungen} } \end{center}

}

%Lizenzierung und Bildverzeichnis

%[[Projekt:Semantische Vorlagen/Latex/Vorlesungsskript/Lizenzierung]]

\newcommand{\Lizenztext}{Teks sumber dibuat oleh Holger Brenner di Wikiversity berbahasa Jerman dan digunakan di sini berdasarkan CC BY-SA 4.0. Media mengikuti lisensi per berkas.}

\newcommand{\Abbildungslizenzierung}{ Erl\"auterung: Die in diesem Text verwendeten Bilder stammen aus Commons (also von http://commons.wikimedia.org) und haben eine Lizenz, die die Verwendung hier erlaubt. Die Bilder werden mit ihren Dateinamen auf Commons angef\"uhrt zusammen mit ihrem Autor bzw. Hochlader und der Lizenz. }

%Konfiguration des Bildverzeichnisses

\newcommand{\bildlizenz}[6]{
\ifthenelse {\equal{#2}{}}
{\addcontentsline{lof}{figure}{Sumber = #1, Kreator = Pengguna #3 di #4, Lisensi = #5 \bildlizenzskip  }}
{ \ifthenelse {\equal{#3}{}}
{ \addcontentsline{lof}{figure}{ Sumber = #1, Kreator = #2, Lisensi = #5 \bildlizenzskip }}
{ \addcontentsline{lof}{figure}{ Sumber = #1, Kreator = #2 (diunggah oleh Pengguna #3 di #4), Lisensi = #5 \bildlizenzskip }} } }

\newcommand{\bildlizenzskip}{ }

%Englische Variante

\newcommand{\ignoriereeng}[1]{}

\ignoriereeng{

%[[Projekt:Semantische Vorlagen/Latex/Englische Zusätze]]

\newtheorem{corollary}[fakt]{Corollary}

\newtheorem{Corollary}[fakt]{Corollary}

\theoremstyle{definition}

\newtheorem{example}[fakt]{Example}

\newtheorem{remark}[fakt]{Remark}

\newcommand{\inputremark}[2] {\bemerkungvorskip \begin{remark} \bemerkungbenennung{#1} \bemerkungbenennungnachskip #2 \end{remark} \bemerkungnachskip}

\newcommand{\inputexample}[2] {\beispielvorskip \begin{example} \beispielbenennung{#1} \beispielbenennungnachskip #2 \end{example} \beispielnachskip}

\newcommand{\inputfaktproof}[5] {\faktvorskip \begin{#2}

%\label{#1} (Absetzen, sonst kann das folgende dahinter rutschen)

\faktbenennung{#3} \faktbenennungnachskip #4 \end{#2} \faktnachskip \beweis{#5}}

\renewcommand{\proofname}{\hspace{-0.64cm}{\it Proof}}

}

%Variante für Artikel

\newcommand{\ignoriereart}[1]{}

\ignoriereart{
\newcommand{\titelueberschrift}[1]{\title[#1] { #1} \maketitle }

\renewcommand{\seitenueberschrift}[1]{ \section {#1} }

\renewcommand{\zwischenueberschrift}[1]{ \subsection {#1} }
}

% Indonesian aliases used by the independent translated reader.
\newtheorem{Teorema}[fakt]{Teorema}
\renewcommand{\proofname}{Bukti}

\newtheorem{Proposisi}[fakt]{Proposisi}
\newtheorem{Lema}[fakt]{Lema}
% Keep each source image and caption together when it fits as a single block.
\renewcommand{\bild}[1]{%
\par\vspace{4mm}%
\noindent\begin{minipage}{\linewidth}#1\end{minipage}%
\par}
% Keep source filenames with underscores safe in the persisted list of figures.
\renewcommand{\bildlizenz}[6]{%
\ifthenelse{\equal{#2}{}}%
{\addcontentsline{lof}{figure}{Sumber = \protect\nolinkurl{#1}, Kreator = Pengguna #3 di #4, Lisensi = #5 \bildlizenzskip}}%
{\ifthenelse{\equal{#3}{}}%
{\addcontentsline{lof}{figure}{Sumber = \protect\nolinkurl{#1}, Kreator = #2, Lisensi = #5 \bildlizenzskip}}%
{\addcontentsline{lof}{figure}{Sumber = \protect\nolinkurl{#1}, Kreator = #2 (diunggah oleh Pengguna #3 di #4), Lisensi = #5 \bildlizenzskip}}}}
% The frozen MediaWiki TeX uses underscore-normalized PNG names while some
% canonical Commons filenames retain spaces. The bounded cumulative build
% creates exact generated aliases and routes PNG macros only through that
% generated directory; canonical source files remain untouched.
\renewcommand{\bildeinlesungpng}[2]{generated/media/#1.png}
\renewcommand{\bildeinlesungPNG}[2]{generated/media/#1.png}
% The frozen source numbers facts and exercises by lecture/worksheet section.
% Book chapters are reader scaffolding and must not enter source identifiers.
\renewcommand{\thefakt}{\arabic{section}.\arabic{fakt}}

\addto\captionsindonesian{%
  \renewcommand{\contentsname}{Daftar Isi}%
  \renewcommand{\figurename}{Gambar}%
  \renewcommand{\listfigurename}{Daftar Gambar}%
}
\hypersetup{pdflang={id-ID},bookmarksopen=true}
\emergencystretch=1em
\ifdefined\pdfinfoomitdate\pdfinfoomitdate=1\fi
\ifdefined\pdftrailerid\pdftrailerid{}\fi
\ifdefined\pdfsuppressptexinfo\pdfsuppressptexinfo=15\fi

\title{Geometri Diferensial dan Manifold Mulus\\\large Pembaca kumulatif hingga Unit 22}
\author{Holger Brenner\\Terjemahan Bahasa Indonesia independen}
\date{Sumber: \textit{Differentialgeometrie (Osnabrück 2023)}}

\begin{document}
\hypersetup{pageanchor=false}
\maketitle
\hypersetup{pageanchor=true}
\frontmatter
\chapter*{Tentang edisi ini}
Pembaca kumulatif ini menerjemahkan Kuliah 1--22 dan Lembar Kerja 1--22 dari kursus Holger Brenner di Wikiversity berbahasa Jerman. Teks sumber digunakan berdasarkan CC BY-SA 4.0. Terjemahan ini merupakan karya independen dan bukan edisi resmi atau dukungan dari penulis maupun Wikiversity. Setiap gambar mengikuti status hak atau lisensi berkasnya sendiri, yang dicatat pada daftar gambar dan manifest media.

Sumber permanen dan hash revisi dicatat dalam berkas kontrol edisi ini. Rumus, urutan, latihan, penanda solusi, dan atribusi media dipertahankan; ID mesin hanya merupakan lapisan tambahan. Bagian solusi hanya memuat solusi yang benar-benar disediakan oleh sumber.

Proses terjemahan dan produksi edisi dibantu oleh OpenAI Codex gpt-5.6-sol, Ultra, di bawah arahan pengguna. Kredit penulis, sumber, dan kontributor manusia tetap dipertahankan sebagaimana dicatat dalam edisi dan manifestnya.
\tableofcontents

\mainmatter
\part{Unit 1}
\chapter[Kuliah 1]{Kuliah 1: Analisis dan Geometri}
\input{generated/lecture01.id.build.tex}

\chapter{Lembar Kerja 1}
\setcounter{section}{1}

  \zwischenueberschrift{Soal latihan}

\inputaufgabegibtloesung
{}
{

Misalkan

\mavergleichskette
{\vergleichskette
{ G
}
{ \subseteq }{ \R^n
}
{  }{
}
{    }{
}
{   }{
}
}
{}{}{}
\definitionsverweis {terbuka}{}{}
dan
\maabbdisp {\varphi} { G } { \R^m
} {}
sebuah
\definitionsverweis {pemetaan terdiferensialkan secara kontinu}{}{,}
yang di titik

\mavergleichskette
{\vergleichskette
{ P
}
{ \in }{ G
}
{  }{
}
{    }{
}
{   }{
}
}
{}{}{}
mempunyai
\definitionsverweis {diferensial total}{}{}
yang surjektif. Misalkan

\mavergleichskette
{\vergleichskette
{ v
}
{ \in }{ T_PZ
}
{  }{
}
{    }{
}
{   }{
}
}
{}{}{}
sebuah vektor di
\definitionsverweis {ruang tangen pada serat}{}{}
$Z$ dari $\varphi$ yang melalui $P$. Buktikan bahwa terdapat
\definitionsverweis {kurva terdiferensialkan secara kontinu}{}{}
\maabbdisp {\gamma} { ]- \epsilon, \epsilon[ } { Z
} {}
\zusatzklammer {untuk suatu

\mavergleichskettek
{\vergleichskettek
{ \epsilon
}
{ > }{ 0
}
{  }{
}
{    }{
}
{   }{
}
}
{}{}{}
yang sesuai} {} {}
dengan

\mavergleichskette
{\vergleichskette
{ \gamma(0)
}
{ = }{ P
}
{  }{
}
{    }{
}
{   }{
}
}
{}{}{}
dan

\mavergleichskettedisp
{\vergleichskette
{ \gamma'(0)
}
{ =} { v
}
{ } {
}
{ } {
}
{ } {
}
}
{}{}{.}

}
{} {}

\inputaufgabe
{}
{

Misalkan

\mavergleichskette
{\vergleichskette
{ W
}
{ \subseteq }{ \R^n
}
{  }{
}
{    }{
}
{   }{
}
}
{}{}{}
\definitionsverweis {terbuka}{}{,}
\maabb {h} { W } { \R
} {}
sebuah
\definitionsverweis {fungsi terdiferensialkan secara kontinu}{}{}
dan

\mavergleichskette
{\vergleichskette
{ Y
}
{ = }{ h^{-1}(0)
}
{  }{
}
{    }{
}
{   }{
}
}
{}{}{}
merupakan
\definitionsverweis {serat}{}{}
di atas $0$, dengan $h$
\definitionsverweis {reguler}{}{}
di setiap titik pada $Y$. Misalkan
\maabbdisp {g} { W } {\R
} {}
sebuah fungsi terdiferensialkan secara kontinu yang tidak mempunyai titik nol pada $Y$. Buktikan bahwa $Y$ juga dapat diperoleh sebagai serat dari

\mavergleichskettedisp
{\vergleichskette
{ \tilde{h}
}
{ =} { g h
}
{ } {
}
{ } {
}
{ } {
}
}
{}{}{,}
bahwa
\definitionsverweis {gradien}{}{}
dari $h$ dan $\tilde{h}$ merupakan kelipatan skalar satu sama lain, dan bahwa
\definitionsverweis {ruang-ruang tangen}{}{}
hanya bergantung pada $Y$.

}
{} {}

\inputaufgabe
{}
{

Misalkan

\mavergleichskette
{\vergleichskette
{ W
}
{ \subseteq }{ \R^n
}
{  }{
}
{    }{
}
{   }{
}
}
{}{}{}
\definitionsverweis {terbuka}{}{,}
\maabb {h} { W } { \R
} {}
sebuah
\definitionsverweis {fungsi terdiferensialkan secara kontinu}{}{}
dan

\mavergleichskette
{\vergleichskette
{ Y
}
{ = }{ h^{-1}(c)
}
{  }{
}
{    }{
}
{   }{
}
}
{}{}{}
merupakan
\definitionsverweis {serat}{}{}
di atas

\mavergleichskette
{\vergleichskette
{ c
}
{ \in }{ \R
}
{  }{
}
{    }{
}
{   }{
}
}
{}{}{,}
dengan $h$
\definitionsverweis {reguler}{}{}
di setiap titik pada $Y$. Misalkan
\maabbdisp {L} {\R^n } { \R^n
} {}
sebuah
\definitionsverweis {isomorfisme linear}{}{,}

\mavergleichskette
{\vergleichskette
{ W'
}
{ = }{ L^{-1}(W)
}
{  }{
}
{    }{
}
{   }{
}
}
{}{}{}
dan

\mavergleichskettedisp
{\vergleichskette
{ Z
}
{ =} { L^{-1}(Y)
}
{ =} { (h \circ L)^{-1}(c)
}
{ } {
}
{ } {
}
}
{}{}{.}
Buktikan bahwa konsep
\definitionsverweis {ruang tangen}{}{,}
medan vektor tangensial, dan solusi persamaan diferensial tangensial yang terkait pada $Y$ dan $Z$ secara umum saling bersesuaian
\zusatzklammer {yakni dapat dipetakan satu sama lain dengan bantuan $L$} {} {.}
Bagaimana dengan gradiennya?

}
{} {}

\inputaufgabe
{}
{

Misalkan

\mavergleichskette
{\vergleichskette
{ Y
}
{ = }{ h^{-1}(c)
}
{  }{
}
{    }{
}
{   }{
}
}
{}{}{}
sebuah
\definitionsverweis {hipermuka terdiferensialkan}{}{}
dan misalkan
\maabbdisp {f} { U} { \R
} {}
sebuah fungsi terdiferensialkan secara kontinu yang didefinisikan pada suatu lingkungan terbuka

\mavergleichskette
{\vergleichskette
{ Y
}
{ \subseteq }{ U
}
{  }{
}
{    }{
}
{   }{
}
}
{}{}{}
Andaikan $f$ di

\mavergleichskette
{\vergleichskette
{ P
}
{ \in }{ Y
}
{  }{
}
{    }{
}
{   }{
}
}
{}{}{}
mempunyai
\definitionsverweis {ekstremum lokal}{}{}
pada $Y$. Buktikan bahwa
\definitionsverweis {komponen tangensial}{}{}
dari $\operatorname{Grad} \,  f  (  P )$ sama dengan $0$.

}
{} {}

\inputaufgabe
{}
{

Misalkan

\mavergleichskette
{\vergleichskette
{ Y
}
{ \subseteq }{ \R^3
}
{  }{
}
{    }{
}
{   }{
}
}
{}{}{}
adalah permukaan bola satuan. Tentukan
\definitionsverweis {dekomposisi ortogonal}{}{}
vektor

\mavergleichskette
{\vergleichskette
{ \begin{pmatrix}  6  \\ 7 \\  -3   \end{pmatrix}
}
{ \in }{ \R^3
}
{  }{
}
{    }{
}
{   }{
}
}
{}{}{}
menjadi komponen tangensial dan komponen ortogonal di titik

\mavergleichskette
{\vergleichskette
{ \begin{pmatrix}  -  { \frac{   1  }{  2 } }  \\ { \frac{   1  }{  2 } } \\  { \frac{   1  }{  \sqrt{2}  } }   \end{pmatrix}
}
{ \in }{ Y
}
{  }{
}
{    }{
}
{   }{
}
}
{}{}{.}

}
{} {}

\inputaufgabe
{}
{

Misalkan

\mavergleichskette
{\vergleichskette
{ Y
}
{ \subseteq }{ \R^3
}
{  }{
}
{    }{
}
{   }{
}
}
{}{}{}
adalah
\definitionsverweis {permukaan bola satuan}{}{.}
Misalkan
\maabbdisp {\gamma} {[0, 2 \pi]} { Y
} {}
merupakan gerak melingkar beraturan yang mengelilingi permukaan bola pada ketinggian ${ \frac{   1  }{  2 } }$.
\aufzaehlungzwei {Parametrisasikan $\gamma$.
} {Tentukan
\definitionsverweis {percepatan tangensial}{}{}
dari $\gamma$ pada setiap saat

\mavergleichskette
{\vergleichskette
{ t
}
{ \in }{ [0, 2 \pi]
}
{  }{
}
{    }{
}
{   }{
}
}
{}{}{.}
}

}
{} {}

\inputaufgabe
{}
{

Misalkan $Y$ adalah
\definitionsverweis {grafik}{}{}
dari
\maabbeledisp {f} {\R^2} { \R
} {(x,y)} { x^2+y^3
} {,}
dan misalkan
\maabbdisp {\gamma} {\R} { Y
} {}
adalah lintasan pada grafik tersebut yang berada di atas kurva dasar
\mathl{t \mapsto (t,t)}{.} Tentukan, untuk setiap

\mavergleichskette
{\vergleichskette
{ t
}
{ \in }{ \R
}
{  }{
}
{    }{
}
{   }{
}
}
{}{}{,}
\definitionsverweis {percepatan tangensial}{}{}
dari $\gamma$ di $\gamma(t)$.

}
{} {}

\inputaufgabe
{}
{

Misalkan

\mavergleichskette
{\vergleichskette
{ Y
}
{ \subseteq }{ \R^n
}
{  }{
}
{    }{
}
{   }{
}
}
{}{}{}
sebuah hiperbidang, yaitu serat dari fungsi linear

\mavergleichskettedisp
{\vergleichskette
{ h(x_1 , \ldots , x_n)
}
{ =} { a_1x_1 + \cdots + a_nx_n
}
{ } {
}
{ } {
}
{ } {
}
}
{}{}{}
di atas

\mavergleichskette
{\vergleichskette
{ c
}
{ \in }{ \R
}
{  }{
}
{    }{
}
{   }{
}
}
{}{}{}
dengan

\mavergleichskette
{\vergleichskette
{ a_i
}
{ \neq }{ 0
}
{  }{
}
{    }{
}
{   }{
}
}
{}{}{}
untuk setidaknya satu indeks $i$. Tafsirkan pengertian
\definitionsverweis {diferensial total}{}{,}
\definitionsverweis {gradien}{}{,}
\definitionsverweis {ruang tangen}{}{,}
dekomposisi menjadi komponen tangensial dan normal,
\definitionsverweis {percepatan tangensial}{}{,}
serta
\definitionsverweis {kurva geodesik}{}{}
dalam kasus khusus ini.

}
{} {}

\inputaufgabe
{}
{

Misalkan

\mavergleichskette
{\vergleichskette
{ Y
}
{ = }{ h^{-1}(c)
}
{  }{
}
{    }{
}
{   }{
}
}
{}{}{}
sebuah
\definitionsverweis {hipermuka terdiferensialkan}{}{,}
yang memuat garis
\mathl{P+ \R v}{.} Buktikan bahwa garis yang dilalui secara seragam tersebut merupakan
\definitionsverweis {kurva geodesik}{}{}
pada $Y$.

}
{} {}

Dua titik

\mathbed {P,Q \in K} {}
 {P \neq Q} {}
 {} {} {} {,}
disebut \stichwort {antipodal} {,} jika garis yang menghubungkannya melalui pusat bola.

\inputaufgabe
{}
{

Misalkan

\mavergleichskette
{\vergleichskette
{ K
}
{ \subseteq }{ \R^3
}
{  }{
}
{    }{
}
{   }{
}
}
{}{}{}
sebuah permukaan bola. Buktikan bahwa setiap dua
\definitionsverweis {titik yang tidak antipodal}{}{}

\mathbed {P,Q \in K} {}
 {P \neq Q} {}
 {} {} {} {,}
terletak pada tepat satu
\definitionsverweis {lingkaran besar}{}{}
dari $K$.

}
{} {}

\bild{ \begin{center}
\includegraphics[width=5.5cm]{ \bildeinlesungsvg {Planned flight map of the Oiseau Blanc} {svg} }
\end{center}
\bildtext {Mengapa pesawat terbang menempuh lintasan melengkung?} }

\bildlizenz { Planned flight map of the Oiseau Blanc.svg } {} {Pethrus} {Commons} {CC BY-SA 3.0} {}

\inputaufgabe
{}
{

Sebuah pesawat akan terbang dari Osnabrück menuju suatu tempat tujuan di belahan Bumi selatan. Mungkinkah rutenya lebih pendek jika pesawat itu mula-mula terbang ke arah utara?

}
{} {}

\inputaufgabe
{}
{

Lucy Sonnenschein merasa bahwa hari-hari terlalu singkat. Karena itu, setiap hari dan selama masih terang, ia memutuskan terbang ke arah barat melintasi satu zona waktu agar setiap harinya berlangsung $25$, bukan $24$, jam.
\aufzaehlungfuenf{Dapatkah Lucy Sonnenschein meningkatkan proporsi jam terang dalam hidupnya dengan strategi ini?
}{Dapatkah Lucy Sonnenschein memperpanjang hidupnya dengan strategi ini?
}{Sekarang Lucy mempunyai teman di seluruh dunia. Setelah berapa hari ia bertemu dengan mereka lagi?
}{Apa yang menyebabkan Lucy mengalami kehilangan waktu yang mengimbangi tambahan waktu hariannya?
}{Apakah hal ini berkaitan dengan teori relativitas?
}

}
{} {}

\inputaufgabe
{}
{

Berikan contoh kurva dua kali terdiferensialkan yang bergerak pada suatu
\definitionsverweis {lingkaran besar}{}{}
dari permukaan bola satuan, tetapi bukan
\definitionsverweis {kurva geodesik}{}{.}

}
{} {}

\inputaufgabe
{}
{

Sebuah \stichwort {segitiga geodesik} {} pada permukaan bola

\mavergleichskettedisp
{\vergleichskette
{ S^2
}
{ =} { \left\{ (x,y,z) \in \R^3  \mid   x^2+y^2+z^2 = 1        \right\}
}
{ } {
}
{ } {
}
{ } {
}
}
{}{}{}
terdiri atas tiga titik berbeda yang tidak terletak pada satu lingkaran besar,

\mavergleichskette
{\vergleichskette
{ A,B,C
}
{ \in }{ S^2
}
{  }{
}
{    }{
}
{   }{
}
}
{}{}{,}
dengan setiap pasangan titik dihubungkan oleh busur yang lebih pendek dari suatu
\definitionsverweis {lingkaran besar}{}{}
\zusatzklammer {setiap pasangan titik tidak antipodal} {} {.}
Buktikan bahwa jumlah sudut dalam segitiga itu lebih besar daripada $\pi$.

}
{} {}

  \zwischenueberschrift{Soal untuk dikumpulkan}

\inputaufgabe
{3}
{

Misalkan $Y$ adalah hipermuka yang diberikan oleh syarat

\mavergleichskettedisp
{\vergleichskette
{ 2x^2 +3y^2+5z^2
}
{ =} { 10
}
{ } {
}
{ } {
}
{ } {
}
}
{}{}{}
di $\R^3$. Tentukan dekomposisi ortogonal vektor

\mavergleichskette
{\vergleichskette
{ \begin{pmatrix}  4  \\ 2 \\  -5   \end{pmatrix}
}
{ \in }{ \R^3
}
{  }{
}
{    }{
}
{   }{
}
}
{}{}{}
menjadi komponen tangensial dan komponen ortogonal di titik

\mavergleichskette
{\vergleichskette
{ \begin{pmatrix}  -1 \\ -1\\  1   \end{pmatrix}
}
{ \in }{ Y
}
{  }{
}
{    }{
}
{   }{
}
}
{}{}{.}

}
{} {}

\inputaufgabe
{4}
{

Misalkan $Y$ adalah
\definitionsverweis {grafik}{}{}
dari
\maabbeledisp {f} {\R^2} { \R
} {(x,y)} { x^2+y^2
} {,}
dan misalkan
\maabbdisp {\gamma} {\R} { Y
} {}
adalah lintasan pada grafik tersebut yang berada di atas kurva dasar
\mathl{x \mapsto (x,1)}{.} Tentukan, untuk setiap

\mavergleichskette
{\vergleichskette
{ x
}
{ \in }{ \R
}
{  }{
}
{    }{
}
{   }{
}
}
{}{}{,}
\definitionsverweis {percepatan tangensial}{}{}
dari $\gamma$ di $\gamma(x)$.

}
{} {}

\inputaufgabe
{4}
{

Pada 11 April 2023 pukul 17.45, secara tiba-tiba dan mengejutkan semua orang, rotasi Bumi pada porosnya dihentikan. Pada saat yang sama, hambatan udara dan semua kerugian akibat gesekan ditiadakan. Semua manusia, hewan, dan benda yang tidak terpasang tetap atau tidak sempat berpegangan akan meneruskan gerak sesaatnya sesuai hukum inersia. Gravitasi tetap bekerja sehingga, untungnya, mereka tidak terlepas ke angkasa. Kapan penghapus papan tulis dari ruang kuliah di Osnabrück
\zusatzklammer {penghapus itu terbang keluar melalui jendela} {} {}
untuk pertama kalinya kembali ke posisi awalnya?

}
{} {}

\inputaufgabe
{4 (3+1)}
{

Misalkan

\mavergleichskette
{\vergleichskette
{ W
}
{ \subseteq }{ \R^n
}
{  }{
}
{    }{
}
{   }{
}
}
{}{}{}
\definitionsverweis {terbuka}{}{,}
\maabb {h} { W } { \R
} {}
sebuah
\definitionsverweis {fungsi terdiferensialkan secara kontinu}{}{}
dan

\mavergleichskette
{\vergleichskette
{ Y
}
{ = }{ h^{-1}(c)
}
{  }{
}
{    }{
}
{   }{
}
}
{}{}{}
merupakan
\definitionsverweis {serat}{}{}
di atas

\mavergleichskette
{\vergleichskette
{ c
}
{ \in }{ \R
}
{  }{
}
{    }{
}
{   }{
}
}
{}{}{,}
dengan $h$
\definitionsverweis {reguler}{}{}
di setiap titik pada $Y$. Misalkan
\maabbdisp {\gamma} { I } { Y
} {}
sebuah kurva dua kali
\definitionsverweis {terdiferensialkan secara kontinu}{}{.}
\aufzaehlungzwei {Buktikan bahwa jika $\gamma$ merupakan
\definitionsverweis {kurva geodesik}{}{,}
maka
\definitionsverweis {norma}{}{}
dari $\gamma'$ adalah konstan.
} {Berikan contoh kurva dengan norma kecepatan konstan yang bukan kurva geodesik.
}

}
{} {}

\inputaufgabe
{3}
{

Misalkan
\maabb {h,g} {\R^n } { \R
} {}
adalah
\definitionsverweis {fungsi-fungsi terdiferensialkan secara kontinu}{}{}
dan

\mavergleichskette
{\vergleichskette
{ Y
}
{ = }{ h^{-1}(c)
}
{  }{
}
{    }{
}
{   }{
}
}
{}{}{}
serta

\mavergleichskette
{\vergleichskette
{ Z
}
{ = }{ g^{-1}(d)
}
{  }{
}
{    }{
}
{   }{
}
}
{}{}{}
adalah
\definitionsverweis {hipermuka terdiferensialkan}{}{}
yang terkait, dengan $h$ di setiap titik pada $Y$ dan $g$ di setiap titik pada $Z$ bersifat
\definitionsverweis {reguler}{}{.}
Misalkan

\mavergleichskette
{\vergleichskette
{ T
}
{ = }{ Y \cap Z
}
{  }{
}
{    }{
}
{   }{
}
}
{}{}{}
dan misalkan
\maabbdisp {\gamma} { I } { T
} {}
sebuah kurva dua kali terdiferensialkan secara kontinu yang, jika dipandang sebagai kurva pada $Y$, merupakan
\definitionsverweis {kurva geodesik}{}{.}
Apakah $\gamma$ juga merupakan kurva geodesik pada $Z$?

}
{} {}


\section*{Solusi yang disediakan oleh sumber}
\addcontentsline{toc}{section}{Solusi yang disediakan oleh sumber}
% Authority: German Wikiversity pageid 131353, revid 1111802, 2026-08-15T16:41:12Z.
% Exact UTF-8 witness SHA-256: b07aac01b6f803481fab9b5331e19480fb9bf0058597159b6152c65c7bc955ce.

\zwischenueberschrift{Solusi untuk Soal 1}

Berlaku

\mavergleichskettedisp
{\vergleichskette
{ v
}
{ \in }{ T_PZ
}
{ = }{ \ker \left((D\varphi)_P\right)
}
{ } {
}
{ } {
}
}
{}{}{.}
Menurut
\definitionsverweis {teorema fungsi implisit}{}{,}
terdapat himpunan terbuka

\mathbed {P \in W} {}
 {W \subseteq G} {}
 {} {} {} {,}
himpunan terbuka

\mavergleichskette
{\vergleichskette
{ V
}
{ \subseteq }{ \R^{n-m}
}
{ } {
}
{ } {
}
{ } {
}
}
{}{}{,}
dan pemetaan terdiferensialkan secara kontinu
\maabbdisp {\psi} { V } { W } {}
sedemikian sehingga

\mavergleichskette
{\vergleichskette
{ \psi(V)
}
{ \subseteq }{ Z \cap W
}
{ } {
}
{ } {
}
{ } {
}
}
{}{}{}
dan $\psi$ menginduksi
\definitionsverweis {bijeksi}{}{}
\maabbdisp {\psi} { V } { Z \cap W } {}.
Misalkan

\mavergleichskette
{\vergleichskette
{ Q
}
{ \in }{ V
}
{ } {
}
{ } {
}
{ } {
}
}
{}{}{}
adalah titik yang memenuhi

\mavergleichskettedisp
{\vergleichskette
{ \psi(Q)
}
{ =} { P
}
{ } {
}
{ } {
}
{ } {
}
}
{}{}{.}
Pemetaan $\psi$ di setiap titik

\mavergleichskette
{\vergleichskette
{ Q
}
{ \in }{ V
}
{ } {
}
{ } {
}
{ } {
}
}
{}{}{}
bersifat
\definitionsverweis {reguler}{}{}
dan
\definitionsverweis {diferensial total}{}{}
dari $\psi$ memenuhi

\mavergleichskettedisp
{\vergleichskette
{ (D\varphi)_P \circ (D\psi)_Q
}
{ =} { 0
}
{ } {
}
{ } {
}
{ } {
}
}
{}{}{.}
Dengan demikian,

\mavergleichskettedisp
{\vergleichskette
{ \operatorname{im} \left((D\psi)_Q\right)
}
{ =} { \ker \left((D\varphi)_P\right)
}
{ =} { T_PZ
}
{ } {
}
{ } {
}
}
{}{}{.}
Karena $\psi$ reguler di $Q$, pemetaan
\maabbdisp {(D\psi)_Q} { \R^{n-m} } { \R^n } {}
bersifat injektif, sedangkan pemetaan
\maabbdisp {(D\psi)_Q} { \R^{n-m} } { T_PZ } {}
bersifat bijektif. Misalkan

\mavergleichskette
{\vergleichskette
{ u
}
{ \in }{ \R^{n-m}
}
{ } {
}
{ } {
}
{ } {
}
}
{}{}{}
adalah prapeta dari $v$, dan misalkan
\maabbeledisp {\theta} { ]-\epsilon,\epsilon[ } { \R^{n-m} } {t} {Q+tu} {,}
dengan $\epsilon$ dipilih cukup kecil agar seluruh citra terletak di dalam $V$. Maka
\maabbdisp {\gamma=\psi\circ\theta} { ]-\epsilon,\epsilon[ } { \R^n } {}
mempunyai sifat

\mavergleichskettedisp
{\vergleichskette
{ \gamma(0)
}
{ =} { \psi(\theta(0))
}
{ =} { \psi(Q)
}
{ =} { P
}
{ } {
}
}
{}{}{}
dan

\mavergleichskettedisp
{\vergleichskette
{ \gamma'(0)
}
{ =} { (D\psi)_Q(\theta'(0))
}
{ =} { (D\psi)_Q(u)
}
{ =} { v
}
{ } {
}
}
{}{}{.}


\part{Unit 2}
\chapter[Kuliah 2]{Kuliah 2: Permukaan Putar, Medan Normal, dan Pemetaan Gauss}
\input{generated/lecture02.id.build.tex}

\chapter{Lembar Kerja 2}
\setcounter{section}{2}

  \zwischenueberschrift{Soal latihan}

\inputaufgabegibtloesung
{}
{

Misalkan
\maabbdisp {f,g} {\R} { \R^n
} {}
\definitionsverweis {kurva-kurva terdiferensialkan}{}{.}
Hitung
\definitionsverweis {turunan}{}{}
fungsi
\maabbeledisp {} {\R} {\R
} { t } { \left\langle f(t) , g(t)  \right\rangle
} {.}
Nyatakan hasilnya menggunakan hasil kali dalam.

}
{} {}

\inputaufgabegibtloesung
{}
{

Misalkan
\maabb {h} { I } { V
} {}
sebuah kurva yang terdiferensialkan dua kali dalam
\definitionsverweis {ruang vektor Euklides}{}{}
$V$. Buktikan bahwa jika

\mavergleichskette
{\vergleichskette
{h'(t)
}
{ \neq }{ 0
}
{  }{
}
{    }{
}
{   }{
}
}
{}{}{,}
maka berlaku persamaan

\mavergleichskettedisp
{\vergleichskette
{ \Vert { h'(t)} \Vert'
}
{ =} { { \frac{   \left\langle h'(t) , h^{\prime \prime} (t)  \right\rangle  }{  \left\langle h'(t) , h'(t)  \right\rangle  } } \Vert { h'(t)} \Vert
}
{ } {
}
{ } {
}
{ } {
}
}
{}{}{.}

}
{} {}

Sebuah
\definitionsverweis {kurva terdiferensialkan}{}{}
\maabbeledisp {f} { [a,b] } {\R^n
} { t } {f(t) = \left( f_1(t)  , \,   \ldots  , \,  f_n(t)                 \right)
} {,}
disebut
\definitionswort {berparametrisasi panjang busur}{,}
jika

\mavergleichskette
{\vergleichskette
{ f_1'(t)^2 + \cdots +  f_n'(t)^2
}
{ = }{ 1
}
{  }{
}
{    }{
}
{   }{
}
}
{}{}{}
untuk setiap $t$.

\inputaufgabe
{}
{

Misalkan
\maabb {\gamma} { I } {\R^n
} {}
sebuah kurva yang dua kali
\definitionsverweis {terdiferensialkan secara kontinu}{}{}
dan
\definitionsverweis {berparametrisasi panjang busur}{}{.}
Buktikan bahwa
\mathkor {} {\gamma'(t)} {dan} {\gamma^{\prime \prime}(t)} {}
saling tegak lurus.

}
{} {}

\inputaufgabe
{}
{

Misalkan

\mavergleichskette
{\vergleichskette
{ P
}
{ \in }{ \R^n
}
{  }{
}
{    }{
}
{   }{
}
}
{}{}{}
sebuah titik dan misalkan

\mavergleichskette
{\vergleichskette
{ I
}
{ = }{ {]{-1},1[}
}
{  }{
}
{    }{
}
{   }{
}
}
{}{}{.}
Kita tinjau himpunan

\mavergleichskettedisp
{\vergleichskette
{ M
}
{ =} { \left\{ f:I \rightarrow \R^n   \mid  f \text{ terdiferensialkan} , \,  f(0) = P       \right\}
}
{ } {
}
{ } {
}
{ } {
}
}
{}{}{.}
Dua
\definitionsverweis {kurva}{}{}

\mavergleichskette
{\vergleichskette
{ f,g
}
{ \in }{ M
}
{  }{
}
{    }{
}
{   }{
}
}
{}{}{}
disebut \stichwort {ekuivalen secara tangensial} {,} jika

\mavergleichskettedisp
{\vergleichskette
{ f'(0)
}
{ =} { g'(0)
}
{ } {
}
{ } {
}
{ } {
}
}
{}{}{.}

a) Buktikan bahwa ini merupakan suatu
\definitionsverweis {relasi ekuivalensi}{}{.}

b) Temukan wakil paling sederhana bagi setiap
\definitionsverweis {kelas ekuivalensi}{}{.}

c) Berikan satu wakil lain untuk setiap kelas.

d) Deskripsikan himpunan kelas-kelas ekuivalensi tersebut
\zusatzklammer {yakni
\definitionsverweis {himpunan hasil bagi}{}{}} {} {.}

}
{} {}

\inputaufgabe
{}
{

Dengan syarat
\zusatzklammer {reguler di setiap titik} {} {,}
berikan sebuah contoh
\definitionsverweis {hipermuka terdiferensialkan}{}{,}
yang tidak
\definitionsverweis {terhubung}{}{.}

}
{} {}

\inputaufgabe
{}
{

Misalkan

\mavergleichskette
{\vergleichskette
{ V
}
{ \subseteq }{ \R^{n-1}
}
{  }{
}
{    }{
}
{   }{
}
}
{}{}{}
sebuah
\definitionsverweis {himpunan terbuka}{}{}
dan misalkan
\maabbdisp {f} { V } {\R
} {}
sebuah fungsi yang
\definitionsverweis {terdiferensialkan secara kontinu}{}{.}
Kita tinjau
\definitionsverweis {grafik}{}{}
$Y$ dari $f$ sebagai hipermuka terdiferensialkan di $\R^n$ dalam pengertian
Contoh 1.2.
Tentukan
\definitionsverweis {medan-medan normal satuan}{}{}
dalam kasus ini.

}
{} {}

\inputaufgabegibtloesung
{}
{

Misalkan
\mathkor {} {r} {dan} {R} {}
bilangan real positif yang memenuhi

\mavergleichskette
{\vergleichskette
{ 0
}
{ < }{ r
}
{ < }{ R
}
{    }{
}
{   }{
}
}
{}{}{.}
Tentukan suatu
\definitionsverweis {medan normal satuan}{}{}
untuk torus
\zusatzklammer {tertanam} {} {}

\mavergleichskettedisp
{\vergleichskette
{ T
}
{ =} { \left\{ (x,y,z) \in \R^3  \mid   \left(  \sqrt{x^2+y^2} -R  \right)^2 +z^2  = r^2         \right\}
}
{ } {
}
{ } {
}
{ } {
}
}
{}{}{.}

}
{} {}

\inputaufgabe
{}
{

Misalkan
\maabbdisp {f} { I } {\R_+
} {}
sebuah
\definitionsverweis {fungsi terdiferensialkan}{}{}
dan misalkan $Y$ adalah permukaan putar yang terkait dengan grafik $f$, dipandang sebagai hipermuka terdiferensialkan di $\R^3$ dalam pengertian
Lemma 2.1.
Tentukan
\definitionsverweis {medan-medan normal satuan}{}{}
dalam kasus ini.

}
{} {}

\inputaufgabe
{}
{

Buktikan bahwa
\definitionsverweis {pemetaan Gauss}{}{}
pada sebuah hiperbidang bersifat konstan. Apakah pernyataan sebaliknya juga berlaku?

}
{} {}

\inputaufgabe
{}
{

Buktikan bahwa
\definitionsverweis {pemetaan Gauss}{}{}
pada sebuah permukaan bola di $\R^n$ yang berpusat di titik asal, untuk setiap pilihan orientasi, diinduksi oleh homoteti skalar pada $\R^n$ (dengan faktor positif untuk orientasi keluar dan faktor negatif untuk orientasi masuk).

}
{} {}

\inputaufgabe
{}
{

Misalkan

\mavergleichskette
{\vergleichskette
{ Y
}
{ \subseteq }{ \R^n
}
{  }{
}
{    }{
}
{   }{
}
}
{}{}{}
sebuah
\definitionsverweis {hipermuka terdiferensialkan}{}{}
dengan
\definitionsverweis {medan normal satuan}{}{}
$N$ dan medan normal satuan yang berlawanan, $-N$. Buktikan bahwa kedua
\definitionsverweis {pemetaan Gauss}{}{}
yang terkait saling berhubungan melalui
\definitionsverweis {pemetaan antipodal}{}{}
\maabbeledisp {} {S^{n-1}} {S^{n-1}
} { P } { -P
} {.}

}
{} {}

\inputaufgabegibtloesung
{}
{

Misalkan

\mavergleichskette
{\vergleichskette
{ \alpha,\beta
}
{ \in }{ \R_+
}
{  }{
}
{    }{
}
{   }{
}
}
{}{}{}
dan misalkan

\mavergleichskettedisp
{\vergleichskette
{ Y
}
{ =} { \left\{ (x,y)\in\R^2  \mid   \alpha x^2+\beta y^2 = 1        \right\}
}
{ } {
}
{ } {
}
{ } {
}
}
{}{}{.}
Orientasikan $Y$ dengan medan normal satuan yang diperoleh dari gradien keluar fungsi
$ (x,y) \longmapsto \alpha x^2+\beta y^2 $.
\aufzaehlungzwei {Tentukan
\definitionsverweis {pemetaan Gauss}{}{}
untuk $Y$.
} { Berikan secara eksplisit suatu
\definitionsverweis {pemetaan invers}{}{}
untuk pemetaan Gauss pada $Y$.
}

}
{} {}

\inputaufgabegibtloesung
{}
{

Misalkan

\mavergleichskette
{\vergleichskette
{ \alpha,\beta
}
{ \in }{ \R_+
}
{  }{
}
{    }{
}
{   }{
}
}
{}{}{}
dan misalkan

\mavergleichskettedisp
{\vergleichskette
{ Y
}
{ =} { \left\{ (x,y)\in\R^2  \mid   \alpha x^2+\beta y^2 = 1        \right\}
}
{ } {
}
{ } {
}
{ } {
}
}
{}{}{.}
Buktikan bahwa
\definitionsverweis {pemetaan Gauss}{}{}
untuk $Y$ bersifat bijektif.

}
{} {}

\inputaufgabe
{}
{

Berikan sebuah contoh
\definitionsverweis {hipermuka terdiferensialkan}{}{}

\mavergleichskette
{\vergleichskette
{ Y
}
{ \subseteq }{ \R^2
}
{  }{
}
{    }{
}
{   }{
}
}
{}{}{}
sedemikian sehingga
\definitionsverweis {pemetaan Gauss}{}{}
bersifat surjektif dan terdapat titik-titik pada $S^1$ yang dicapai tak berhingga kali.

}
{} {}

\inputaufgabe
{}
{

Misalkan

\mavergleichskette
{\vergleichskette
{ W
}
{ \subseteq }{ \R^n
}
{  }{
}
{    }{
}
{   }{
}
}
{}{}{}
\definitionsverweis {terbuka}{}{,}
\maabb {h} { W } { \R
} {}
sebuah
\definitionsverweis {fungsi yang terdiferensialkan secara kontinu}{}{}
dan

\mavergleichskette
{\vergleichskette
{ Y
}
{ = }{ h^{-1}(c)
}
{  }{
}
{    }{
}
{   }{
}
}
{}{}{}
adalah
\definitionsverweis {serat}{}{}
di atas

\mavergleichskette
{\vergleichskette
{ c
}
{ \in }{ \R
}
{  }{
}
{    }{
}
{   }{
}
}
{}{}{,}
dengan $h$
\definitionsverweis {reguler}{}{}
di setiap titik pada $Y$. Misalkan
\maabbdisp {L} {\R^n } { \R^n
} {}
sebuah
\definitionsverweis {isometri linear}{}{,}

\mavergleichskette
{\vergleichskette
{ W'
}
{ = }{ L^{-1}(W)
}
{  }{
}
{    }{
}
{   }{
}
}
{}{}{}
dan

\mavergleichskettedisp
{\vergleichskette
{ Z
}
{ =} { L^{-1}(Y)
}
{ =} { (h \circ L)^{-1}(c)
}
{ } {
}
{ } {
}
}
{}{}{.}
Buktikan bahwa konsep-konsep kurva geodesik, medan normal, medan normal satuan, dan pemetaan Gauss pada $Y$ dan $Z$ saling bersesuaian
\zusatzklammer {yakni dapat dipetakan satu sama lain dengan bantuan $L$} {} {.}

}
{} {}

  \zwischenueberschrift{Soal untuk dikumpulkan}

\inputaufgabe
{3}
{

Tentukan citra parabola standar berorientasi ke atas di bawah
\definitionsverweis {pemetaan Gauss}{}{.}

}
{} {}

\inputaufgabe
{4}
{

Misalkan

\mavergleichskette
{\vergleichskette
{ W
}
{ = }{ \R^3 \setminus \{ (0,0,0) \}
}
{  }{
}
{    }{
}
{   }{
}
}
{}{}{}
dan misalkan

\mavergleichskette
{\vergleichskette
{ Y
}
{ \subseteq }{ W
}
{  }{
}
{    }{
}
{   }{
}
}
{}{}{}
adalah
\definitionsverweis {himpunan nol}{}{}
dari

\mavergleichskettedisp
{\vergleichskette
{ h(x,y,z)
}
{ =} { x^2+y^2-z^2
}
{ } {
}
{ } {
}
{ } {
}
}
{}{}{.}
Tentukan citra $Y$
\zusatzklammer {dengan orientasi yang ditentukan oleh gradien $h$} {} {}
di bawah
\definitionsverweis {pemetaan Gauss}{}{.}

}
{} {}

\inputaufgabe
{5}
{

Misalkan

\mavergleichskette
{\vergleichskette
{ \alpha,\beta, \gamma
}
{ \in }{ \R_+
}
{  }{
}
{    }{
}
{   }{
}
}
{}{}{}
dan misalkan

\mavergleichskettedisp
{\vergleichskette
{ Y
}
{ =} { \left\{ (x,y,z) \in \R^3  \mid   \alpha x^2+ \beta y^2 + \gamma z^2 = 1        \right\}
}
{ } {
}
{ } {
}
{ } {
}
}
{}{}{.}
Buktikan bahwa
\definitionsverweis {pemetaan Gauss}{}{}
untuk $Y$ bersifat bijektif.

}
{} {}

\inputaufgabe
{3}
{

Misalkan

\mavergleichskette
{\vergleichskette
{ W
}
{ \subseteq }{ \R^n
}
{  }{
}
{    }{
}
{   }{
}
}
{}{}{}
\definitionsverweis {terbuka}{}{,}
\maabb {h} { W } { \R
} {}
sebuah
\definitionsverweis {fungsi yang terdiferensialkan secara kontinu}{}{}
dan

\mavergleichskette
{\vergleichskette
{ Y
}
{ = }{ h^{-1}(c)
}
{  }{
}
{    }{
}
{   }{
}
}
{}{}{}
adalah
\definitionsverweis {serat}{}{}
di atas

\mavergleichskette
{\vergleichskette
{ c
}
{ \in }{ \R
}
{  }{
}
{    }{
}
{   }{
}
}
{}{}{,}
dengan $h$
\definitionsverweis {reguler}{}{}
di setiap titik pada $Y$. Misalkan
\maabbdisp {L} {\R^n } { \R^n
} {}
sebuah
\definitionsverweis {isomorfisme linear}{}{,}

\mavergleichskette
{\vergleichskette
{ W'
}
{ = }{ L^{-1}(W)
}
{  }{
}
{    }{
}
{   }{
}
}
{}{}{}
dan

\mavergleichskettedisp
{\vergleichskette
{ Z
}
{ =} { L^{-1}(Y)
}
{ =} { (h \circ L)^{-1}(c)
}
{ } {
}
{ } {
}
}
{}{}{.}
Buktikan bahwa konsep-konsep kurva geodesik, medan normal, medan normal satuan, dan pemetaan Gauss pada $Y$ dan $Z$ secara umum tidak saling bersesuaian
\zusatzklammer {berbeda dengan
Soal 2.14,
yang melibatkan suatu isometri} {} {.}

}
{} {}


\section*{Solusi yang disediakan oleh sumber}
\addcontentsline{toc}{section}{Solusi yang disediakan oleh sumber}
\subsection*{Solusi Soal 2.1}
\input{generated/worksheet02_exercise01_solution.id.build.tex}
\subsection*{Solusi Soal 2.2}
% Authority: German Wikiversity pageid 120435, revid 1089268, 2026-05-31T10:02:12Z.
% Exact UTF-8 wikitext witness SHA-256: 0bdd1cde6d1d0dc74757aa824634d10b6b4bf413efb4431538b7d9bac2b6f06c.
% Expanded German TeX source SHA-256: a92e556f50d2216192fc61703eea4bae3d233ab632c8eedb19bd35dec2ed89b0.

\textit{Catatan edisi: solusi sumber memakai basis ortogonal seolah-olah
ortonormal dan menuliskan turunan koordinat yang keliru. Argumen berikut
mengganti langkah tersebut.}

Karena $h'(t)\neq 0$, norma kecepatannya tak nol. Turunkan kuadrat normanya:

\[
\frac{d}{dt}\Vert h'(t)\Vert^2
=2\Vert h'(t)\Vert\bigl(\Vert h'(t)\Vert\bigr)'
=2\left\langle h'(t),h''(t)\right\rangle .
\]

Membagi dengan $2\Vert h'(t)\Vert$ memberikan

\[
\bigl(\Vert h'(t)\Vert\bigr)'
=\frac{\left\langle h'(t),h''(t)\right\rangle}{\Vert h'(t)\Vert}
=\frac{\left\langle h'(t),h''(t)\right\rangle}
       {\left\langle h'(t),h'(t)\right\rangle}\Vert h'(t)\Vert .
\]

\subsection*{Solusi Soal 2.7}
\begingroup\scriptsize
\input{generated/worksheet02_exercise07_solution.id.build.tex}
\endgroup
\subsection*{Solusi Soal 2.12}
\begingroup\small
\input{generated/worksheet02_exercise12_solution.id.build.tex}
\endgroup
\subsection*{Solusi Soal 2.13}
\input{generated/worksheet02_exercise13_solution.id.build.tex}

\part{Unit 3}
\chapter[Kuliah 3]{Kuliah 3: Kelengkungan Kurva Berparameter Panjang Busur}
\input{generated/lecture03.id.build.tex}

\chapter{Lembar Kerja 3}
\setcounter{section}{3}

  \zwischenueberschrift{Soal latihan}

\inputaufgabe
{}
{

Misalkan
\maabb {\gamma} { I } {\R^2
} {}
sebuah
\definitionsverweis {kurva berparametrisasi panjang busur}{}{}
yang dua kali terdiferensialkan secara kontinu, dan misalkan
\maabbeledisp {\delta} { -I } {\R^2
} { s } { \delta(s) = \gamma( -s)
} {,}
yaitu kurva yang ditelusuri dalam arah sebaliknya. Buktikan bahwa
\definitionsverweis {kelengkungan}{}{}
$\delta$ di $s$ merupakan negatif dari kelengkungan
\mathl{\gamma}{} di $-s$.

}
{} {}

\inputaufgabe
{}
{

Misalkan
\maabb {\gamma} { I } {\R^2
} {}
sebuah
\definitionsverweis {kurva berparametrisasi panjang busur}{}{}
yang dua kali terdiferensialkan secara kontinu, dan
\maabbdisp {L} {\R^2} {\R^2
} {}
sebuah
\definitionsverweis {rotasi}{}{.}
Buktikan bahwa
\definitionsverweis {kelengkungan}{}{}
$\gamma$ sama dengan kelengkungan
\mathl{L\circ \gamma}{}.

}
{} {}

\inputaufgabe
{}
{

Buktikan
Lema 3.6
dalam kasus dengan orientasi nonstandar.

}
{} {}

\inputaufgabe
{}
{

Misalkan
\maabb {\gamma} { I } {\R^2
} {}
sebuah
\definitionsverweis {kurva berparametrisasi panjang busur}{}{}
yang dua kali terdiferensialkan secara kontinu, dan
\maabbdisp {L} {\R^2} {\R^2
} {}
sebuah
\definitionsverweis {homoteti skalar}{}{}
dengan faktor homoteti

\mavergleichskette
{\vergleichskette
{ s
}
{ \neq }{ 0
}
{  }{
}
{    }{
}
{   }{
}
}
{}{}{.}
Bagaimana hubungan antara
\definitionsverweis {kelengkungan}{}{}
$\gamma$ dan kelengkungan
\mathl{L\circ \gamma}{?}

}
{} {}

\inputaufgabe
{}
{

Misalkan
\maabbeledisp {\gamma} {\R} { \R^2
} { t } { \begin{pmatrix}  \cos \left(  \omega t \right)  \\ \sin  \left(  \omega t \right)    \end{pmatrix}
} {,}
dengan

\mavergleichskette
{\vergleichskette
{ \omega
}
{ > }{ 0
}
{  }{
}
{    }{
}
{   }{
}
}
{}{}{.}
Buktikan bahwa terdapat tak berhingga banyak gerak melingkar dengan norma kecepatan konstan yang memiliki nilai dan turunan pertama yang sama dengan $\gamma$ pada

\mavergleichskette
{\vergleichskette
{ t
}
{ = }{ 0
}
{  }{
}
{    }{
}
{   }{
}
}
{}{}{.}

}
{} {}

Soal berikut penting dalam pembuatan komidi putar. Kesenangan menaiki komidi putar berasal dari percepatan, bukan dari kecepatan.

\inputaufgabe
{}
{

Misalkan
\maabbeledisp {\gamma} {\R} { \R^2
} { t } { \begin{pmatrix}  \cos \left(  \omega t \right)  \\ \sin  \left(  \omega t \right)    \end{pmatrix}
} {,}
dengan

\mavergleichskette
{\vergleichskette
{ \omega
}
{ > }{ 0
}
{  }{
}
{    }{
}
{   }{
}
}
{}{}{.}
Buktikan bahwa terdapat tak berhingga banyak gerak melingkar dengan norma kecepatan konstan yang berimpit dengan $\gamma$ pada

\mavergleichskette
{\vergleichskette
{ t
}
{ = }{ 0
}
{  }{
}
{    }{
}
{   }{
}
}
{}{}{}
dan juga memiliki percepatan yang sama di sana.

}
{} {}

\inputaufgabegibtloesung
{}
{

Misalkan
\maabbeledisp {\gamma} {\R} { \R^2
} { t } { \begin{pmatrix}  \cos \left(  \omega t \right)  \\ \sin  \left(  \omega t \right)    \end{pmatrix}
} {,}
dengan

\mavergleichskette
{\vergleichskette
{ \omega
}
{ > }{ 0
}
{  }{
}
{    }{
}
{   }{
}
}
{}{}{.}
Buktikan bahwa tidak ada gerak melingkar lain dengan norma kecepatan konstan yang memiliki nilai serta turunan pertama dan kedua yang sama dengan $\gamma$ pada

\mavergleichskette
{\vergleichskette
{ t
}
{ = }{ 0
}
{  }{
}
{    }{
}
{   }{
}
}
{}{}{.}

}
{} {}

\inputaufgabe
{}
{

Misalkan

\mavergleichskette
{\vergleichskette
{ W
}
{ \subseteq }{ \R^2
}
{  }{
}
{    }{
}
{   }{
}
}
{}{}{}
\definitionsverweis {terbuka}{}{,}
\maabb {h} { W } { \R
} {}
sebuah fungsi yang dua kali
\definitionsverweis {terdiferensialkan secara kontinu}{}{,}
dan

\mavergleichskette
{\vergleichskette
{ Y
}
{ = }{ h^{-1}(c)
}
{  }{
}
{    }{
}
{   }{
}
}
{}{}{}
adalah
\definitionsverweis {serat}{}{}
di atas

\mavergleichskette
{\vergleichskette
{ c
}
{ \in }{ \R
}
{  }{
}
{    }{
}
{   }{
}
}
{}{}{,}
dengan $h$
\definitionsverweis {reguler}{}{}
di setiap titik pada $Y$. Buktikan bahwa
\maabbdisp {\gamma} { I } { Y
} {}
merupakan
\definitionsverweis {kurva geodesik}{}{}
jika dan hanya jika $\gamma$
\definitionsverweis {berparametrisasi panjang busur}{}{.}

}
{} {}

\inputaufgabe
{}
{

\bild{ \begin{center}
\includegraphics[width=5.5cm]{ \bildeinlesungsvg {Euler spiral} {svg} }
\end{center}
\bildtext {} }

\bildlizenz { Euler spiral.svg } {} {AdiJapan} {Commons} {CC-by-sa 3.0} {}

Kita tinjau kurva terdiferensialkan
\zusatzklammer {yang disebut \stichwort {klotoid} {}} {} {}
\maabbeledisp {} {\R} {\R^2
} { t } { \gamma(t)
} {,}
dengan

\mavergleichskettedisp
{\vergleichskette
{ \gamma(t)
}
{ =} { \sqrt{\pi} \int_0^t  \begin{pmatrix}  \cos  { \frac{   \pi s^2 }{  2 } }  \\ \sin  { \frac{   \pi s^2 }{  2 } }     \end{pmatrix}  ds
}
{ } {
}
{ } {
}
{ } {
}
}
{}{}{.}
Buktikan bahwa
\definitionsverweis {kelengkungan}{}{}
kurva ini memenuhi persamaan

\mavergleichskettedisp
{\vergleichskette
{ \kappa(t)
}
{ =} { \sqrt{\pi} t
}
{ } {
}
{ } {
}
{ } {
}
}
{}{}{.}

}
{} {}

\inputaufgabe
{}
{

Misalkan
\maabb {f} { I } {\R
} {}
sebuah fungsi yang dua kali
\definitionsverweis {terdiferensialkan secara kontinu}{}{.}
Tentukan
\definitionsverweis {kelengkungan}{}{}
dan lingkaran kelengkungan dari grafik yang terkait.

}
{} {}

\inputaufgabe
{}
{

Misalkan
\maabb {f} {\R} {\R
} {}
dua kali
\definitionsverweis {terdiferensialkan secara kontinu}{}{}
dan

\mavergleichskettedisp
{\vergleichskette
{ \alpha(t)
}
{ \defeq} { \begin{pmatrix}  \cos f(t)  \\ \sin f(t)   \end{pmatrix}
}
{ } {
}
{ } {
}
{ } {
}
}
{}{}{.}
Tentukan
\definitionsverweis {kelengkungan}{}{}
$\alpha$ dan
\definitionsverweis {lingkaran kelengkungan}{}{}
untuk

\mavergleichskette
{\vergleichskette
{ t
}
{ \in }{ \R
}
{  }{
}
{    }{
}
{   }{
}
}
{}{}{}
dengan syarat $f'(t) \neq 0$, menggunakan rumus-rumus dalam
Lema 3.10.

}
{} {}

\inputaufgabe
{}
{

Tentukan
\definitionsverweis {kelengkungan}{}{}
\definitionsverweis {grafik}{}{}
\maabbdisp {\sin} {\R} {\R
} {}
untuk

\mavergleichskette
{\vergleichskette
{ t
}
{ = }{ { \frac{   \pi }{  2 } }
}
{  }{
}
{    }{
}
{   }{
}
}
{}{}{.}

}
{} {}

\inputaufgabe
{}
{

Tentukan
\definitionsverweis {kelengkungan}{}{}
spiral logaritmik
\maabbeledisp {\gamma} {\R} { \R^2
} { t } { \gamma(t) = e^t \left(  \cos  t   , \,   \sin  t                   \right)
} {,}
untuk setiap

\mavergleichskette
{\vergleichskette
{ t
}
{ \in }{ \R
}
{  }{
}
{    }{
}
{   }{
}
}
{}{}{.}

}
{} {}

Sebuah fungsi
\maabb {f} { I } {\R
} {}
pada interval terbuka

\mavergleichskette
{\vergleichskette
{ I
}
{ \subseteq }{ \R
}
{  }{
}
{    }{
}
{   }{
}
}
{}{}{}
disebut
\definitionswort {analitik}{,}
jika pada setiap titik

\mavergleichskette
{\vergleichskette
{ a
}
{ \in }{ I
}
{  }{
}
{    }{
}
{   }{
}
}
{}{}{}
fungsi tersebut dapat dinyatakan oleh suatu
\definitionsverweis {deret pangkat}{}{.}

Sebuah kurva disebut analitik jika setiap fungsi komponennya analitik.

\inputaufgabe
{}
{

Misalkan
\maabb {\kappa} { I } {\R
} {}
sebuah
\definitionsverweis {fungsi analitik}{}{.}
Buktikan bahwa terdapat kurva analitik
\maabb {\gamma} { I } {\R^2
} {}
yang
\definitionsverweis {kelengkungan}{}{}
di titik $\gamma(t)$ sama dengan $\kappa(t)$.

}
{} {}

\inputaufgabe
{}
{

\bild{ \begin{center}
\includegraphics[width=5.5cm]{ \bildeinlesungsvg {Evolute-parab} {svg} }
\end{center}
\bildtext {} }

\bildlizenz { Evolute-parab.svg } {} {Ag2gaeh} {Commons} {CC-by-sa 4.0} {}

Tentukan
\definitionsverweis {evolut}{}{}
kurva
\maabbeledisp {} {\R} {\R^2
} { t } { \begin{pmatrix}  t  \\ t^2    \end{pmatrix}
} {.}

}
{} {}

\inputaufgabegibtloesung
{}
{

Tentukan
\definitionsverweis {kelengkungan}{}{}
di setiap titik lingkaran satuan yang diberikan secara implisit oleh

\mavergleichskettedisp
{\vergleichskette
{ h(x,y)
}
{ =} { x^2+y^2
}
{ =} { 1
}
{ } {
}
{ } {
}
}
{}{}{}
menggunakan
Lema 3.11
dan medan gradien $h$.

}
{} {}

  \zwischenueberschrift{Soal untuk dikumpulkan}

\inputaufgabe
{2}
{

Tentukan
\definitionsverweis {kelengkungan}{}{}
\definitionsverweis {grafik}{}{}
fungsi eksponensial
\maabbdisp {\exp} {\R} {\R
} {}
untuk setiap

\mavergleichskette
{\vergleichskette
{ t
}
{ \in }{ \R
}
{  }{
}
{    }{
}
{   }{
}
}
{}{}{.}

}
{} {}

\inputaufgabe
{2}
{

Tentukan
\definitionsverweis {kelengkungan}{}{}
spiral Archimedes
\maabbeledisp {\gamma} {\R} { \R^2
} { t } { \gamma(t) = t \left(  \cos  t   , \,   \sin  t                   \right)
} {,}
untuk setiap

\mavergleichskette
{\vergleichskette
{ t
}
{ \in }{ \R
}
{  }{
}
{    }{
}
{   }{
}
}
{}{}{.}

}
{} {}

\inputaufgabe
{4}
{

Berikan sebuah contoh kurva berparametrisasi panjang busur
\maabbdisp {\gamma} {\R} {\R^2
} {}
yang dua kali terdiferensialkan secara kontinu, sedemikian sehingga untuk dua waktu

\mavergleichskette
{\vergleichskette
{ t_0
}
{ \neq }{ t_1
}
{  }{
}
{    }{
}
{   }{
}
}
{}{}{}
berlaku kesamaan

\mavergleichskette
{\vergleichskette
{ \gamma (t_0)
}
{ = }{ \gamma(t_1)
}
{  }{
}
{    }{
}
{   }{
}
}
{}{}{}
dan
\definitionsverweis {lingkaran-lingkaran kelengkungan}{}{}
pada
\mathkor {} {t_0} {dan} {t_1} {}
tidak sama.

}
{} {}

\inputaufgabe
{4}
{

Tentukan
\definitionsverweis {evolut}{}{}
kurva
\maabbeledisp {} {\R} {\R^2
} { t } { \begin{pmatrix}  2 \cos  t  \\ 3 \sin  t    \end{pmatrix}
} {.}
Di titik-titik manakah turunan evolut mempunyai titik nol?

}
{} {}

\inputaufgabe
{4}
{

Tentukan
\definitionsverweis {kelengkungan}{}{}
lingkaran satuan yang diberikan secara implisit oleh

\mavergleichskettedisp
{\vergleichskette
{ h(x,y)
}
{ =} { x^4+y^4
}
{ =} { 1
}
{ } {
}
{ } {
}
}
{}{}{}
menggunakan
Lema 3.11
dan medan gradien $h$ di titik-titik berikut.
\aufzaehlungzwei {
\mavergleichskette
{\vergleichskette
{ P
}
{ = }{ \begin{pmatrix}  1  \\ 0   \end{pmatrix}
}
{  }{
}
{    }{
}
{   }{
}
}
{}{}{,}
} {
\mavergleichskette
{\vergleichskette
{ Q
}
{ = }{ \begin{pmatrix}  { \frac{   1  }{  \sqrt[4]{2}  } }  \\ { \frac{   1  }{  \sqrt[4]{2}  } }     \end{pmatrix}
}
{  }{
}
{    }{
}
{   }{
}
}
{}{}{.}
}

}
{} {}


\section*{Solusi yang disediakan oleh sumber}
\addcontentsline{toc}{section}{Solusi yang disediakan oleh sumber}
\subsection*{Solusi Soal 3.7}
% Authority: German Wikiversity pageid 151168, revid 1095913, 2026-06-15T07:25:24Z.
% Exact UTF-8 wikitext witness SHA-256: c771b7c1ed173ddc05c8c4219c3c34e7e8099fe42079740b292d991b2ce3fb25.
% Expanded German TeX source SHA-256: 7e3437274bf4f79b6a3fa719876b09fdbcee4e10523a215e9e6f4ecb566cdb85.

Nilai-nilainya adalah

\mavergleichskettedisp
{\vergleichskette
{ \gamma(0)
}
{ =} { \begin{pmatrix}  1  \\ 0   \end{pmatrix}
}
{ } {
}
{ } {
}
{ } {
}
}
{}{}{,}

\mavergleichskettedisp
{\vergleichskette
{ \gamma'(0)
}
{ =} { \begin{pmatrix}  0  \\ \omega   \end{pmatrix}
}
{ } {
}
{ } {
}
{ } {
}
}
{}{}{}
dan

\mavergleichskettedisp
{\vergleichskette
{ \gamma^{\prime \prime} (0)
}
{ =} { \begin{pmatrix}  - \omega^2  \\ 0   \end{pmatrix}
}
{ } {
}
{ } {
}
{ } {
}
}
{}{}{.}
Suatu gerak melingkar yang memiliki nilai serta turunan pertama dan kedua yang sama dengan gerak yang diberikan pada

\mavergleichskette
{\vergleichskette
{ t
}
{ = }{ 0
}
{  }{
}
{    }{
}
{   }{
}
}
{}{}{}
harus, khususnya, mempunyai garis singgung yang sama di titik ini. Karena itu, pusat lingkaran harus terletak pada sumbu $x$. Selain itu, pusat tersebut harus berada di sebelah kiri
\mathl{\begin{pmatrix}  1  \\ 0   \end{pmatrix}}{},
karena percepatan mengarah ke kiri. Oleh sebab itu, kita dapat menuliskan gerak melingkar beraturan yang dicari, dengan pusat
\mathl{\begin{pmatrix}  x  \\ 0   \end{pmatrix}}{}
\zusatzklammer {
\mavergleichskettek
{\vergleichskettek
{ x
}
{ < }{  1
}
{  }{
}
{    }{
}
{   }{
}
}
{}{}{}} {} {}
dan jari-jari
\mathl{1-x}{}, sebagai

\mavergleichskettedisp
{\vergleichskette
{  \delta (t)
}
{ =} {  \begin{pmatrix}  x  \\ 0   \end{pmatrix} + (1-x) \begin{pmatrix}  \cos \left(  u t \right)  \\ \sin  \left(  u t \right)     \end{pmatrix}
}
{ } {
}
{ } {
}
{ } {
}
}
{}{}{.}
Kita memperoleh

\mavergleichskettedisp
{\vergleichskette
{ \delta(0)
}
{ =} { \gamma(0)
}
{ } {
}
{ } {
}
{ } {
}
}
{}{}{.}
Selanjutnya,

\mavergleichskettedisp
{\vergleichskette
{ \delta'( 0)
}
{ =} {  (1-x)  \begin{pmatrix}  0  \\ u     \end{pmatrix}
}
{ } {
}
{ } {
}
{ } {
}
}
{}{}{}
dan

\mavergleichskettedisp
{\vergleichskette
{ \delta^{\prime \prime}( 0)
}
{ =} {  (1-x)  \begin{pmatrix}  -u^2  \\ 0     \end{pmatrix}
}
{ } {
}
{ } {
}
{ } {
}
}
{}{}{.}
Jadi, agar syarat-syarat yang diminta terpenuhi, harus berlaku

\mavergleichskettedisp
{\vergleichskette
{ \omega
}
{ =} { (1-x) u
}
{ } {
}
{ } {
}
{ } {
}
}
{}{}{}
dan

\mavergleichskettedisp
{\vergleichskette
{ \omega^2
}
{ =} { (1-x) u^2
}
{ } {
}
{ } {
}
{ } {
}
}
{}{}{.}
Dari sini diperoleh

\mavergleichskettedisp
{\vergleichskette
{ \omega^2
}
{ =} {  \omega u
}
{ } {
}
{ } {
}
{ } {
}
}
{}{}{}
dan dengan demikian

\mavergleichskettedisp
{\vergleichskette
{ u
}
{ =} { \omega
}
{ } {
}
{ } {
}
{ } {
}
}
{}{}{}
serta

\mavergleichskettedisp
{\vergleichskette
{ x
}
{ =} { 0
}
{ } {
}
{ } {
}
{ } {
}
}
{}{}{.}

\subsection*{Solusi Soal 3.16}
\input{generated/worksheet03_exercise16_solution.id.build.tex}

\part{Unit 4}
\chapter[Kuliah 4]{Kuliah 4: Pemetaan Weingarten}
\setcounter{section}{4}

  \zwischenueberschrift{Pemetaan Weingarten}

Misalkan

\mavergleichskette
{\vergleichskette
{ Y
}
{ \subseteq }{ \R^n
}
{  }{
}
{    }{
}
{   }{
}
}
{}{}{}
adalah suatu
\definitionsverweis {hipermuka terdiferensialkan}{}{.}
Untuk suatu titik

\mavergleichskette
{\vergleichskette
{ P
}
{ \in }{ Y
}
{  }{
}
{    }{
}
{   }{
}
}
{}{}{}
kita mendefinisikan suatu pemetaan linear
\maabbdisp {L_P} { T_PY } { T_PY
} {}
pada ruang tangen $Y$ di P. Misalkan $N$ adalah suatu medan normal satuan terdiferensialkan pada $Y$ yang didefinisikan pada suatu lingkungan terbuka di sekitar $Y$. Gagasan pokoknya adalah merepresentasikan suatu vektor tangen

\mavergleichskette
{\vergleichskette
{ v
}
{ \in }{ T_PY
}
{  }{
}
{    }{
}
{   }{
}
}
{}{}{}
dengan suatu kurva terdiferensialkan
\maabbdisp {\gamma} { I } { Y
} {}
sehingga

\mavergleichskettedisp
{\vergleichskette
{ \gamma'(0)
}
{ =} { v
}
{ } {
}
{ } {
}
{ } {
}
}
{}{}{.}
Medan normal satuan itu kemudian mendefinisikan pemetaan
\maabbdisp {N \circ \gamma} { I } { \R^n
} {}
sepanjang $I$. Perubahan infinitesimal medan normal satuan sepanjang $\gamma$ diukur oleh limit

\mathdisp {\operatorname{lim}_{   t  \rightarrow  0   } \,  { \frac{   N(\gamma(t))-N(P)  }{  t  } }} {  }

jika limit ini ada. Menurut
Lemma 43.4 (Analysis (Osnabrück 2021-2023)),
limit ini sama dengan turunan berarah $\left(  D_{ v }  N   \right) \left(   P   \right)$ dan, pada gilirannya, sama dengan
\mathl{\left(  D N   \right)_{ P } \left(   v   \right)}{,} yaitu
\definitionsverweis {diferensial total}{}{}
yang dievaluasi pada vektor $v$. Akan ditunjukkan bahwa pemetaan
\maabbeledisp {} { T_PY } { \R^n
} {v } { \left(  D_{v}  N   \right) \left(   P   \right)
} {,}
ini mengambil nilai di $T_PY$.

\inputdefinition
{}
{

Misalkan

\mavergleichskette
{\vergleichskette
{ W
}
{ \subseteq }{ \R^n
}
{  }{
}
{    }{
}
{   }{
}
}
{}{}{}
\definitionsverweis {terbuka}{}{,}
\maabb {h} { W } { \R
} {}
suatu
\definitionsverweis {fungsi yang dua kali terdiferensialkan secara kontinu}{}{}
dan

\mavergleichskette
{\vergleichskette
{ Y
}
{ = }{ h^{-1}(c)
}
{  }{
}
{    }{
}
{   }{
}
}
{}{}{}
adalah
\definitionsverweis {serat}{}{}
di atas

\mavergleichskette
{\vergleichskette
{ c
}
{ \in }{ \R
}
{  }{
}
{    }{
}
{   }{
}
}
{}{}{,}
dengan $h$
\definitionsverweis {reguler}{}{}
pada setiap titik di $Y$. Misalkan $N$ adalah suatu
\definitionsverweis {medan normal satuan}{}{}
dan misalkan

\mavergleichskette
{\vergleichskette
{ P
}
{ \in }{ Y
}
{  }{
}
{    }{
}
{   }{
}
}
{}{}{.}
Pemetaan
\maabbeledisp {L_P} {T_P Y} { T_PY
} { v } { - \left(  D_{v}  N   \right) \left(   P  \right)
} {,}
disebut
\definitionswort {pemetaan Weingarten}{}
pada $T_PY$.

}

Perhatikan bahwa pemetaan Weingarten bergantung pada medan normal satuan, meskipun hal ini tidak selalu dinyatakan secara eksplisit. Jika $Y$ diberikan sebagai serat dari $h$, biasanya digunakan medan gradien ternormalisasi yang bersesuaian dengan $h$.


\inputfaktbeweis
{Hyperfläche/Weingartenabbildung/Tangentialraum/Fakt}
{Lemma}
{}
{

\faktsituation {Misalkan

\mavergleichskette
{\vergleichskette
{ W
}
{ \subseteq }{ \R^n
}
{  }{
}
{    }{
}
{   }{
}
}
{}{}{}
\definitionsverweis {terbuka}{}{,}
\maabb {h} { W } { \R
} {}
suatu
\definitionsverweis {fungsi yang dua kali terdiferensialkan secara kontinu}{}{}
dan

\mavergleichskette
{\vergleichskette
{ Y
}
{ = }{ h^{-1}(c)
}
{  }{
}
{    }{
}
{   }{
}
}
{}{}{}
adalah
\definitionsverweis {serat}{}{}
di atas

\mavergleichskette
{\vergleichskette
{ c
}
{ \in }{ \R
}
{  }{
}
{    }{
}
{   }{
}
}
{}{}{,}
dengan $h$
\definitionsverweis {reguler}{}{}
pada setiap titik di $Y$. Misalkan

\mavergleichskette
{\vergleichskette
{ P
}
{ \in }{ Y
}
{  }{
}
{    }{
}
{   }{
}
}
{}{}{.}}
\faktfolgerung {Maka
\definitionsverweis {pemetaan Weingarten}{}{}
$L_P$ adalah suatu endomorfisme linear dari
\definitionsverweis {ruang tangen}{}{}
$T_PY$.}
Catatan edisi: sumber memakai huruf kecil p pada ruang tangen ini, padahal titik yang dikuantifikasi adalah P.
\faktzusatz {}
\faktzusatz {}

}
{

Medan

\mavergleichskette
{\vergleichskette
{ N(P)
}
{ = }{ { \frac{
\operatorname{Grad} \,  h  (  P )  }{  \Vert {
\operatorname{Grad} \,  h  (  P ) } \Vert  } }
}
{  }{
}
{    }{
}
{   }{
}
}
{}{}{}
adalah suatu medan vektor yang terdiferensialkan secara kontinu pada suatu lingkungan terbuka

\mavergleichskette
{\vergleichskette
{ Y
}
{ \subseteq }{ U
}
{  }{
}
{    }{
}
{   }{
}
}
{}{}{}
tempat medan tersebut didefinisikan. Oleh karena itu, menurut
Proposition 46.1 (Analysis (Osnabrück 2021-2023)),

\mavergleichskettedisp
{\vergleichskette
{ \left(  D_{v}  N   \right) \left(   P  \right)
}
{ =} { \left(  D N   \right)_{ P } \left(  v  \right)
}
{ } {
}
{ } {
}
{ } {
}
}
{}{}{}
bersifat linear terhadap arah $v$. Dari syarat normal satuan diperoleh

\mavergleichskette
{\vergleichskette
{ \left\langle  N(P)  ,  N(P)   \right\rangle
}
{ = }{ 1
}
{  }{
}
{    }{
}
{   }{
}
}
{}{}{}
untuk setiap

\mavergleichskette
{\vergleichskette
{ P
}
{ \in }{ U
}
{  }{
}
{    }{
}
{   }{
}
}
{}{}{}
dan, dengan menggunakan
Soal 43.14 (Analysis (Osnabrück 2021-2023)), diperoleh

\mavergleichskettedisp
{\vergleichskette
{ 0
}
{ =} { \left(  D_{v}  \left\langle  N  ,  N  \right\rangle   \right) \left(   P  \right)
}
{ =} { 2 \left\langle  N(P)  ,  \left(  D_{v}  N   \right) \left(   P  \right)   \right\rangle
}
{ } {
}
{ } {
}
}
{}{}{.}
Jadi,
\mathl{\left(  D_{v}  N   \right) \left(   P  \right)}{} tegak lurus terhadap $N(P)$ dan, untuk

\mavergleichskette
{\vergleichskette
{ P
}
{ \in }{ Y
}
{  }{
}
{    }{
}
{   }{
}
}
{}{}{}
vektor tersebut termasuk dalam ruang tangen $T_PY$.

}

Dengan demikian, pemetaan Weingarten $L_P$ adalah negatif dari pembatasan diferensial total medan normal satuan pada ruang tangen $T_PY$, dengan kodomain ruang tangen $T_PY$. Jika diferensial total dihitung melalui
\definitionsverweis {matriks Jacobi}{}{}
dari $-N$, kita harus menentukan cara matriks itu bekerja pada suatu basis ruang tangen untuk memperoleh representasi matriks dari pemetaan Weingarten.



\inputfaktbeweis
{Ebene Kurve/Weingartenabbildung/Krümmung/Fakt}
{Lemma}
{}
{

\faktsituation {Misalkan

\mavergleichskette
{\vergleichskette
{ W
}
{ \subseteq }{ \R^2
}
{  }{
}
{    }{
}
{   }{
}
}
{}{}{}
\definitionsverweis {terbuka}{}{,}
\maabb {h} { W } { \R
} {}
suatu
\definitionsverweis {fungsi yang dua kali terdiferensialkan secara kontinu}{}{}
dan

\mavergleichskette
{\vergleichskette
{ C
}
{ = }{ h^{-1}(c)
}
{  }{
}
{    }{
}
{   }{
}
}
{}{}{}
adalah serat di atas

\mavergleichskette
{\vergleichskette
{ c
}
{ \in }{ \R
}
{  }{
}
{    }{
}
{   }{
}
}
{}{}{,}
dengan $h$
\definitionsverweis {reguler}{}{}
pada setiap titik di $C$, dan misalkan

\mavergleichskette
{\vergleichskette
{ P
}
{ \in }{ C
}
{  }{
}
{    }{
}
{   }{
}
}
{}{}{.}}
\faktfolgerung {Maka
\definitionsverweis {pemetaan Weingarten}{}{}
di $P$ adalah perkalian dengan
\definitionsverweis {kelengkungan}{}{}
$\kappa(P)$ dari $C$ di $P$.}
\faktzusatz {}
\faktzusatz {}

}
{

Pernyataan ini merupakan perumusan ulang dari
Lemma 3.11.

}

Pernyataan di atas tidak secara eksplisit merujuk pada orientasi; orientasi yang dimaksud harus turut diperhitungkan.

\inputbeispiel{}
{

Misalkan

\mavergleichskettedisp
{\vergleichskette
{ h(x,y,z)
}
{ =} { x^2+y^2+z^2
}
{ } {
}
{ } {
}
{ } {
}
}
{}{}{}
dan perhatikan serat $Y$ dari $h$ di atas $r^2$, yakni permukaan bola berjari-jari

\mavergleichskette
{\vergleichskette
{ r
}
{ > }{ 0
}
{  }{
}
{    }{
}
{   }{
}
}
{}{}{}
yang berpusat di titik asal. Misalkan

\mavergleichskette
{\vergleichskette
{ P
}
{ = }{ \left(  x_0   , \,   y_0  , \,  z_0                 \right)
}
{ \in }{ Y
}
{    }{
}
{   }{
}
}
{}{}{.}
Dengan suatu isometri, titik ini dapat dipetakan ke

\mavergleichskette
{\vergleichskette
{ Q
}
{ = }{ \left(  r   , \,   0  , \,   0                 \right)
}
{  }{
}
{    }{
}
{   }{
}
}
{}{}{}
tanpa mengubah
\definitionsverweis {pemetaan Weingarten}{}{}.
Suatu basis ruang tangennya ialah
\mathkor {} {\begin{pmatrix}  0  \\ 1 \\  0   \end{pmatrix}} {dan} {\begin{pmatrix}  0  \\ 0\\  1   \end{pmatrix}} {.}
Medan normal satuan $N$ yang mengarah ke dalam adalah
\mathl{- { \frac{   1  }{  2r } } \begin{pmatrix}  2x \\ 2y\\  2z   \end{pmatrix}}{} dan karena itu, untuk

\mavergleichskettedisp
{\vergleichskette
{ v
}
{ =} { \begin{pmatrix}  0  \\ b \\ c   \end{pmatrix}
}
{ } {
}
{ } {
}
{ } {
}
}
{}{}{}

\mavergleichskettedisp
{\vergleichskette
{ \left(  D_{v}  N   \right) \left(   Q   \right)
}
{ =} { b { \frac{   \partial   N   }{  \partial   y  } } (Q) +  c { \frac{   \partial   N   }{  \partial  z  } }  (Q)
}
{ =} { -b { \frac{   1  }{  2r } } \begin{pmatrix}  0  \\ 2 \\  0   \end{pmatrix} - c { \frac{   1  }{  2r } } \begin{pmatrix}  0  \\ 0\\  2   \end{pmatrix}
}
{ =} { - { \frac{   1  }{ r } } \begin{pmatrix}  0  \\ b \\ c   \end{pmatrix}
}
{ } {
}
}
{}{}{.}
Jadi, pemetaan Weingarten adalah perkalian dengan kebalikan ${ \frac{   1  }{ r } }$ dari jari-jari.

}

\inputbeispiel{}
{

Kita melanjutkan
Contoh 2.5.
Matriks Jacobi dari medan normal satuan

\mavergleichskettedisphandlinks
{\vergleichskettedisphandlinks
{ N( \begin{pmatrix}  x  \\ y \\ z   \end{pmatrix} )
}
{ =} { \left(  x^2+y^2+z^2  \right)^{ - { \frac{   1  }{  2 } } } \begin{pmatrix}  x  \\ y \\  -z   \end{pmatrix}
}
{ } {
}
{ } {
}
{ } {
}
}
{}{}{}
adalah

\mathdisp {\left(  x^2+y^2+z^2  \right)^{- { \frac{   3  }{  2 } } } \begin{pmatrix}  y^2+z^2  &   -xy &  -xz \\  -xy &  x^2+z^2 &  -yz \\ xz &  yz &  -x^2-y^2  \end{pmatrix}} { . }

\definitionsverweis {Pemetaan Weingarten}{}{}
adalah negatif dari pembatasan pemetaan ini pada ruang tangen permukaan tersebut; menurut
Lemma 4.2,
hasilnya kembali berada di ruang tangen. Kita mengasumsikan

\mavergleichskette
{\vergleichskette
{ x
}
{ \neq }{ 0
}
{  }{
}
{    }{
}
{   }{
}
}
{}{}{}
dan menggunakan basis
\mathkor {} {\begin{pmatrix}  y  \\ -x\\  0   \end{pmatrix}} {dan} {\begin{pmatrix}  z  \\ 0 \\  x   \end{pmatrix}} {}
dari ruang tangen. Kita mempunyai

\mavergleichskettealignhandlinks
{\vergleichskettealignhandlinks
{ \begin{pmatrix}  y^2+z^2  &   -xy &  -xz \\  -xy &  x^2+z^2 &  -yz \\ xz &  yz &  -x^2-y^2  \end{pmatrix} \begin{pmatrix}  y  \\ -x\\  0   \end{pmatrix}
}
{ =} { \begin{pmatrix}  y \left(  y^2+z^2  \right) +x^2y \\ -xy^2 -x \left(  x^2+z^2  \right) \\  0   \end{pmatrix}
}
{ =} { \left(  x^2+y^2+z^2  \right) \begin{pmatrix}  y  \\ -x\\  0   \end{pmatrix}
}
{ } {
}
{ } {
}
}
{}
{}{,}
sehingga langsung diperoleh vektor eigen $\begin{pmatrix}  y  \\ -x\\  0   \end{pmatrix}$ dengan nilai eigen
\mathl{- \left(  x^2+y^2+z^2  \right)^{ - { \frac{   1  }{  2 } } }}{}.
Selain itu,


\mavergleichskettealigndrucklinks
{\vergleichskettealigndrucklinks
{ \begin{pmatrix}  y^2+z^2  &   -xy &  -xz \\  -xy &  x^2+z^2 &  -yz \\ xz &  yz &  -x^2-y^2  \end{pmatrix} \begin{pmatrix}  z  \\ 0 \\  x   \end{pmatrix}
}
{ =} { \begin{pmatrix} z \left(  y^2+z^2  \right) -x^2z \\ -2xyz\\  xz^2-x \left(  x^2+y^2  \right)   \end{pmatrix}
}
{ =} { 2yz \begin{pmatrix}  y  \\ -x\\  0   \end{pmatrix} + \left(  z^2-x^2-y^2  \right) \begin{pmatrix}  z  \\ 0 \\  x   \end{pmatrix}
}
{ =} { 2yz \begin{pmatrix}  y  \\ -x\\  0   \end{pmatrix}- \begin{pmatrix}  z  \\ 0 \\  x   \end{pmatrix}
}
{ } {
}
}
{}{}{.}

Dengan demikian, terhadap basis ini, pemetaan Weingarten direpresentasikan oleh matriks

\mathdisp {\begin{pmatrix}  - \left(  x^2+y^2+z^2  \right)^{ - { \frac{   1  }{  2 } } }   &   - 2yz \left(  x^2+y^2+z^2  \right)^{ - { \frac{   3  }{  2 } } }    \\  0  &  \left(  x^2+y^2+z^2  \right)^{ - { \frac{   3  }{  2 } } }   \end{pmatrix}} {  }

Untuk memperoleh vektor eigen kedua dalam ruang tangen,
\zusatzklammer {dengan memperhatikan
Korolari 4.8} {} {}
cara termudah ialah menentukan suatu vektor yang tegak lurus terhadap vektor eigen pertama dan vektor normal. Hasilnya adalah vektor
\mathl{\begin{pmatrix}  xz \\ yz\\  x^2+y^2   \end{pmatrix}}{,} dan memang,
\zusatzklammer {tanpa faktor skalar di depan} {} {}


\mavergleichskettealigndrucklinks
{\vergleichskettealigndrucklinks
{ \begin{pmatrix}  y^2+z^2  &   -xy &  -xz \\  -xy &  x^2+z^2 &  -yz \\ xz &  yz &  -x^2-y^2  \end{pmatrix} \begin{pmatrix}  xz \\ yz\\  x^2+y^2   \end{pmatrix}
}
{ =} { \begin{pmatrix}  xz \left(  y^2+z^2  \right) -xy^2z -xz \left(  x^2+y^2  \right)  \\ - x^2yz +yz \left(  x^2+z^2  \right) -yz \left(  x^2+y^2  \right) \\  x^2z^2 +y^2z^2 - \left(  x^2+y^2  \right)^2   \end{pmatrix}
}
{ =} { \left(  - x^2-y^2+z^2  \right) \begin{pmatrix}  xz \\ yz\\  x^2+y^2   \end{pmatrix}
}
{ =} { - \begin{pmatrix}  xz \\ yz\\  x^2+y^2   \end{pmatrix}
}
{ } {
}
}
{}
{}{.}

Jadi, nilai eigennya adalah
\mathl{\left(  x^2+y^2+z^2  \right)^{- { \frac{   3  }{  2 } } }}{}
\zusatzklammer {hal ini juga dapat langsung dibaca dari matriks segitiga yang merepresentasikannya} {} {.}

}




\inputfaktbeweis
{Hyperfläche/Weingartenabbildung/Zweite Ableitung/Skalarprodukt/Fakt}
{Lemma}
{}
{

\faktsituation {Misalkan

\mavergleichskette
{\vergleichskette
{ W
}
{ \subseteq }{ \R^n
}
{  }{
}
{    }{
}
{   }{
}
}
{}{}{}
\definitionsverweis {terbuka}{}{,}
\maabb {h} { W } { \R
} {}
suatu
\definitionsverweis {fungsi yang dua kali terdiferensialkan secara kontinu}{}{}
dan

\mavergleichskette
{\vergleichskette
{ Y
}
{ = }{ h^{-1}(c)
}
{  }{
}
{    }{
}
{   }{
}
}
{}{}{}
adalah
\definitionsverweis {serat}{}{}
di atas

\mavergleichskette
{\vergleichskette
{ c
}
{ \in }{ \R
}
{  }{
}
{    }{
}
{   }{
}
}
{}{}{,}
dengan $h$
\definitionsverweis {reguler}{}{}
pada setiap titik di $Y$. Misalkan $N$ adalah suatu
\definitionsverweis {medan normal satuan}{}{}
dan misalkan

\mavergleichskette
{\vergleichskette
{ P
}
{ \in }{ Y
}
{  }{
}
{    }{
}
{   }{
}
}
{}{}{.}
Misalkan $\gamma$ adalah suatu realisasi yang terdiferensialkan dua kali pada $Y$ dari vektor tangen

\mavergleichskette
{\vergleichskette
{ v
}
{ \in }{ T_PY
}
{  }{
}
{    }{
}
{   }{
}
}
{}{}{.}}
\faktfolgerung {Maka

\mavergleichskettedisp
{\vergleichskette
{ \left\langle  \gamma^{\prime \prime} (0)  ,  N(P)  \right\rangle
}
{ =} { \left\langle  L_P(v) , v  \right\rangle
}
{ } {
}
{ } {
}
{ } {
}
}
{}{}{.}}
\faktzusatz {}
\faktzusatz {}

}
{

Misalkan
\maabbeledisp {\gamma} { I } { Y
} { t } { \gamma(t)
} {,}
terdiferensialkan dua kali, dengan

\mavergleichskette
{\vergleichskette
{ \gamma(0)
}
{ = }{ P
}
{  }{
}
{    }{
}
{   }{
}
}
{}{}{}
dan

\mavergleichskette
{\vergleichskette
{ \gamma'(0)
}
{ = }{ v
}
{ \in }{ T_PY
}
{    }{
}
{   }{
}
}
{}{}{.}
Kita mempunyai

\mavergleichskettedisp
{\vergleichskette
{ \left\langle  \gamma'(t)  ,  N(\gamma(t))   \right\rangle
}
{ =} { 0
}
{ } {
}
{ } {
}
{ } {
}
}
{}{}{}
untuk semua

\mavergleichskette
{\vergleichskette
{ t
}
{ \in }{ I
}
{  }{
}
{    }{
}
{   }{
}
}
{}{}{}
sehingga, dengan
Soal 43.14 (Analysis (Osnabrück 2021-2023)), diperoleh

\mavergleichskettealign
{\vergleichskettealign
{ 0
}
{ =} { \left\langle  \gamma'(t)  ,  N(\gamma(t))   \right\rangle' (0)
}
{ =} { \left\langle  \gamma^{\prime \prime}(0)  ,  N(\gamma(0))  \right\rangle + \left\langle  \gamma'(0)  ,  \left(  D_{ \gamma'(0) }  N   \right) \left(   \gamma(0)   \right)   \right\rangle
}
{ =} { \left\langle  \gamma^{\prime \prime}(0) , N(P)  \right\rangle - \left\langle  L_P(v)  ,  v   \right\rangle
}
{ } {
}
}
{}
{}{.}

}

Dalam pernyataan di atas tidak diperlukan penetapan tambahan mengenai orientasi, sebab medan normal satuan muncul pada kedua ruas, dan pada ruas kanan muncul melalui pemetaan Weingarten. Perhatikan pula bahwa $\gamma$ tidak harus memiliki kecepatan konstan. Dengan demikian, terdapat banyak realisasi dan percepatan $\gamma^{\prime \prime} (0)$ tidak ditentukan secara tunggal. Namun, nilai
\mathl{\left\langle  \gamma^{\prime \prime}(0)  ,  N(P)  \right\rangle}{,} yang menggambarkan komponen normal percepatan, tidak berubah karena $\gamma$ bergerak pada hipermuka tersebut.
Catatan edisi: sumber kehilangan argumen N(P) dalam hasil kali dalam pada kalimat sebelumnya; identitas yang baru saja dibuktikan menentukannya secara unik.



\inputfaktbeweis
{Hyperfläche/Weingartenabbildung/Selbstadjungiert/Fakt}
{Satz}
{}
{

\faktsituation {Misalkan

\mavergleichskette
{\vergleichskette
{ W
}
{ \subseteq }{ \R^n
}
{  }{
}
{    }{
}
{   }{
}
}
{}{}{}
\definitionsverweis {terbuka}{}{,}
\maabb {h} { W } { \R
} {}
suatu
\definitionsverweis {fungsi yang dua kali terdiferensialkan secara kontinu}{}{}
dan

\mavergleichskette
{\vergleichskette
{ Y
}
{ = }{ h^{-1}(c)
}
{  }{
}
{    }{
}
{   }{
}
}
{}{}{}
adalah serat di atas

\mavergleichskette
{\vergleichskette
{ c
}
{ \in }{ \R
}
{  }{
}
{    }{
}
{   }{
}
}
{}{}{,}
dengan $h$
\definitionsverweis {reguler}{}{}
pada setiap titik di $Y$.}
\faktfolgerung {Maka
\definitionsverweis {pemetaan Weingarten}{}{}
$L_P$ pada setiap titik

\mavergleichskette
{\vergleichskette
{ P
}
{ \in }{ Y
}
{  }{
}
{    }{
}
{   }{
}
}
{}{}{}
bersifat
\definitionsverweis {adjoin-diri (self-adjoint)}{}{.}}
\faktzusatz {}
\faktzusatz {}

}
{

Untuk vektor

\mavergleichskette
{\vergleichskette
{ v,w
}
{ \in }{ T_PY
}
{  }{
}
{    }{
}
{   }{
}
}
{}{}{}
kita harus menunjukkan bahwa

\mavergleichskettedisp
{\vergleichskette
{ \left\langle  v  ,  L_P(w)  \right\rangle
}
{ =} { \left\langle L_P(v) ,  w  \right\rangle
}
{ } {
}
{ } {
}
{ } {
}
}
{}{}{}
tersebut. Dengan

\mavergleichskette
{\vergleichskette
{ N(P)
}
{ = }{ { \frac{
\operatorname{Grad} \,  h  (  P )  }{  \Vert {
\operatorname{Grad} \,  h  (  P ) } \Vert  } }
}
{  }{
}
{    }{
}
{   }{
}
}
{}{}{}
menurut
Lemma 45.10 (Analysis (Osnabrück 2021-2023)), kita mempunyai

\mavergleichskettealignhandlinks
{\vergleichskettealignhandlinks
{ \left\langle  v  ,  L_P(w)  \right\rangle
}
{ =} { -\left\langle  v  ,  \left(  D_{w}  N   \right) \left(   P   \right)   \right\rangle
}
{ =} { -\left\langle  v  ,  \left(  D_{w}  { \frac{
\operatorname{Grad} \,  h  }{  \Vert {
\operatorname{Grad} \,  h } \Vert  } }   \right) \left(   P   \right)     \right\rangle
}
{ =} { -\left\langle  v  ,  \left(  D_{w}  { \frac{   1  }{  \Vert {
\operatorname{Grad} \,  h } \Vert  } }   \right) \left(   P   \right) \cdot
\operatorname{Grad} \,  h  (  P ) + { \frac{   1  }{  \Vert {
\operatorname{Grad} \,  h  (  P ) } \Vert  } } \cdot  \left(  D_{w}
\operatorname{Grad} \,  h   \right) \left(   P   \right)   \right\rangle
}
{ =} { - { \frac{   1  }{  \Vert {
\operatorname{Grad} \,  h  (  P ) } \Vert  } } \left\langle  v  ,  \left(  D_{w}
\operatorname{Grad} \,  h   \right) \left(   P   \right)    \right\rangle
}
}
{}
{}{,}
karena suku pertama tegak lurus terhadap vektor tangen $v$. Dalam fungsi-fungsi koordinat,

\mavergleichskettealignhandlinks
{\vergleichskettealignhandlinks
{ \left(  D_{w}
\operatorname{Grad} \,  h   \right) \left(   P   \right)
}
{ =} { \left(  D_{w}  \left(  { \frac{   \partial   h   }{  \partial   x_1   } }   , \,   \ldots  , \,   { \frac{   \partial   h   }{  \partial   x_n   } }                 \right)    \right) \left(   P   \right)
}
{ =} { \sum_{j = 1}^n  w_j  { \frac{   \partial    }{  \partial   x_j   } }  \left(  { \frac{   \partial   h   }{  \partial   x_1   } }   , \,   \ldots  , \,   { \frac{   \partial   h   }{  \partial   x_n   } }                 \right) (P)
}
{ =} { \left(  \sum_{j = 1}^n  w_j  { \frac{   \partial    }{  \partial   x_j   } }  { \frac{   \partial   h   }{  \partial   x_1   } } (P)   , \,   \ldots  , \,   \sum_{j = 1}^n  w_j  { \frac{   \partial    }{  \partial   x_j   } }  { \frac{   \partial   h   }{  \partial   x_n   } } (P)                 \right)
}
{ } {
}
}
{}
{}{.}
Dengan demikian, ungkapan di atas sama dengan

\mavergleichskettealign
{\vergleichskettealign
{ \left\langle  v  ,  L_P(w)  \right\rangle
}
{ =} { - { \frac{   1  }{  \Vert {
\operatorname{Grad} \,  h  (  P ) } \Vert  } } \left\langle  v  ,  \left(  D_{w}
\operatorname{Grad} \,  h   \right) \left(   P   \right)   \right\rangle
}
{ =} { -  { \frac{   1  }{  \Vert {
\operatorname{Grad} \,  h  (  P ) } \Vert  } }  \sum_{i,j = 1}^n v_i  w_j  { \frac{   \partial    }{  \partial   x_j   } } { \frac{   \partial   h   }{  \partial   x_i   } }  (P)
}
{ } {
}
{ } {
}
}
{}
{}{.}
Menurut
Teorema 44.10 (Analysis (Osnabrück 2021-2023)),
urutan turunan parsial dapat dipertukarkan, sehingga peran
\mathkor {} {v} {dan} {w} {}
juga dapat dipertukarkan.

}



\inputfaktbeweis
{Hyperfläche/Weingartenabbildung/Selbstadjungiert/Eigenwerte/Fakt}
{Korollar}
{}
{

\faktsituation {Misalkan

\mavergleichskette
{\vergleichskette
{ W
}
{ \subseteq }{ \R^n
}
{  }{
}
{    }{
}
{   }{
}
}
{}{}{}
\definitionsverweis {terbuka}{}{,}
\maabb {h} { W } { \R
} {}
suatu
\definitionsverweis {fungsi yang dua kali terdiferensialkan secara kontinu}{}{}
dan

\mavergleichskette
{\vergleichskette
{ Y
}
{ = }{ h^{-1}(c)
}
{  }{
}
{    }{
}
{   }{
}
}
{}{}{}
adalah serat di atas

\mavergleichskette
{\vergleichskette
{ c
}
{ \in }{ \R
}
{  }{
}
{    }{
}
{   }{
}
}
{}{}{,}
dengan $h$
\definitionsverweis {reguler}{}{}
pada setiap titik di $Y$.}
\faktfolgerung {Maka
\definitionsverweis {pemetaan Weingarten}{}{}
$L_P$ pada setiap titik

\mavergleichskette
{\vergleichskette
{ P
}
{ \in }{ Y
}
{  }{
}
{    }{
}
{   }{
}
}
{}{}{}
\definitionsverweis {dapat didiagonalkan}{}{}
dengan
\definitionsverweis {nilai-nilai eigen}{}{}
real dan ruang-ruang eigen yang saling ortogonal.}
\faktzusatz {}
\faktzusatz {}

}
{

Hal ini mengikuti dari
Teorema 4.7
dan
Teorema 41.11 (Lineare Algebra (Osnabrück 2024-2025)).

}



\inputfaktbeweis
{Differenzierbare Funktion/Graph/Hyperfläche/Weingartenabbildung/Fakt}
{Lemma}
{}
{

\faktsituation {Misalkan

\mavergleichskette
{\vergleichskette
{ V
}
{ \subseteq }{ \R^{n-1}
}
{  }{
}
{    }{
}
{   }{
}
}
{}{}{}
suatu
\definitionsverweis {himpunan terbuka}{}{}
dan misalkan

\mavergleichskette
{\vergleichskette
{ Y
}
{ \subseteq }{ V \times \R
}
{  }{
}
{    }{
}
{   }{
}
}
{}{}{}
adalah
\definitionsverweis {grafik}{}{}
dari suatu fungsi yang
\definitionsverweis {dua kali terdiferensialkan secara kontinu}{}{}
\maabbdisp {f} { V } {\R
} {.}}
\faktfolgerung {Maka
\definitionsverweis {pemetaan Weingarten}{}{}
$L_P$ pada suatu titik

\mavergleichskette
{\vergleichskette
{ P
}
{ = }{ (x_1 , \ldots , x_{n-1}, f(x_1 , \ldots , x_{n-1}))
}
{ \in }{ Y
}
{    }{
}
{   }{
}
}
{}{}{}
\zusatzklammer {terhadap medan normal satuan yang mengarah ke atas} {} {}
dapat direpresentasikan sebagai berikut. Tuliskan \mathl{g=\operatorname{Grad} f}{} sebagai vektor kolom dan
\mathl{G=I_{n-1}+gg^{\mathsf T}}{,}
yakni matriks metrik dalam basis grafik. Matriks pemetaan Weingarten dalam basis tersebut adalah
\mathl{G^{-1}{ \frac{   1  }{  \sqrt{ 1+ \left(  { \frac{   \partial   f   }{  \partial   x_1   } }  \right)^2 + \cdots + \left(  { \frac{   \partial   f   }{  \partial   x_{n-1}   } }  \right)^2 }  } } \operatorname{Hess}_{   } \,  f}{}, dengan $\operatorname{Hess}_{   } \,  f$ menyatakan
\definitionsverweis {matriks Hesse}{}{}
dari $f$, apabila vektor-vektor dasar

\mavergleichskette
{\vergleichskette
{ v
}
{ \in }{ \R^{n-1}
}
{  }{
}
{    }{
}
{   }{
}
}
{}{}{}
diidentifikasi dengan vektor-vektor tangen dari $T_PY$ dalam pengertian
Contoh 1.2.}
\faktzusatz {}
\faktzusatz {}

}
{

Kita melanjutkan
Contoh 1.2.
Medan normal satuan yang mengarah ke atas diberikan oleh

\mavergleichskettedisphandlinks
{\vergleichskettedisphandlinks
{ N(x_1  , \ldots , x_{n-1}, f(x_1 , \ldots , x_{n-1}))
}
{ =} { { \frac{   1  }{  \sqrt{ 1+ \left(  { \frac{   \partial   f   }{  \partial   x_1   } }  \right)^2 + \cdots + \left(  { \frac{   \partial   f   }{  \partial   x_{n-1}   } }  \right)^2 }  } } \begin{pmatrix}  -{ \frac{   \partial   f   }{  \partial   x_{1}   } }  \\\vdots\\  -{ \frac{   \partial   f   }{  \partial   x_{n-1}   } }\\ 1   \end{pmatrix}
}
{ } {
}
{ } {
}
{ } {
}
}
{}{}{}
yang akan digunakan. Kita meninjau lintasan

\mavergleichskettedisp
{\vergleichskette
{ \gamma(t)
}
{ =} { \begin{pmatrix}  x_1+tv_1  \\ \vdots \\  x_{n-1} +tv_{n-1}\\f \begin{pmatrix}  x_1+tv_1  \\ \vdots \\  x_{n-1} +tv_{n-1}   \end{pmatrix}   \end{pmatrix}
}
{ } {
}
{ } {
}
{ } {
}
}
{}{}{}
pada grafik yang bersesuaian dengan vektor dasar $v$. Turunan keduanya adalah

\mavergleichskettedisp
{\vergleichskette
{ \gamma^{\prime \prime} (0)
}
{ =} { \begin{pmatrix}  0  \\ \vdots \\  0\\ \sum_{1 \leq i, j \leq n-1}  { \frac{   \partial    }{  \partial   x_i   } }  { \frac{   \partial   f   }{  \partial   x_j   } } v_iv_j   \end{pmatrix}
}
{ } {
}
{ } {
}
{ } {
}
}
{}{}{.}
Menurut
Lemma 4.6,

\mavergleichskettealigndrucklinks
{\vergleichskettealigndrucklinks
{ \left\langle  L_P( \gamma'(0) ) ,  \gamma'(0)  \right\rangle
}
{ =} { \left\langle  \gamma^{\prime \prime} (0) ,  N(x_1 , \ldots , x_{n-1}, f(x_1  , \ldots , x_{n-1}))   \right\rangle
}
{ =} { { \frac{   1  }{  \sqrt{ 1 + \left(  { \frac{   \partial   f   }{  \partial   x_1   } }  \right)^2 + \cdots + \left(  { \frac{   \partial   f  }{  \partial   x_{n-1}   } }  \right)^2 }  } } \left\langle  \begin{pmatrix}  0  \\\vdots\\  0\\ \sum_{ 1 \leq i, j \leq n-1 }  { \frac{   \partial    }{  \partial   x_i   } }  { \frac{   \partial   f   }{  \partial   x_j   } } v_iv_j   \end{pmatrix}  ,  \begin{pmatrix}  -{ \frac{   \partial   f   }{  \partial   x_{1}   } }  \\\vdots\\  -{ \frac{   \partial   f   }{  \partial   x_{n-1}   } }\\ 1   \end{pmatrix}   \right\rangle
}
{ =} { { \frac{   \sum_{ 1 \leq i, j \leq n-1 } { \frac{   \partial    }{  \partial   x_i   } }  { \frac{   \partial   f   }{  \partial   x_j   } } v_i v_j     }{  \sqrt{ 1+  \left(  { \frac{   \partial   f   }{  \partial   x_1   } }  \right)^2 + \cdots + \left(  { \frac{   \partial   f  }{  \partial   x_{n-1}   } }  \right)^2 }  } }
}
{ =} { { \frac{   1  }{  \sqrt{ 1+ \left(  { \frac{   \partial   f   }{  \partial   x_1   } }  \right)^2 + \cdots + \left(  { \frac{   \partial   f   }{  \partial   x_{n-1}   } }  \right)^2 }  } }  \left( v_1   , \,   \ldots  , \,   v_{n-1}                 \right)  \operatorname{Hess}_{   } \,  f  \begin{pmatrix} v_1  \\ \vdots \\  v_{n-1}   \end{pmatrix}
}
}
{}{}{.}
Artinya, menurut
Teorema 4.7,
\definitionsverweis {bentuk bilinear}{}{}
simetris

\mathl{\left\langle  L_P(-) ,  -  \right\rangle}{} pada ruang tangen $T_PY$ bersesuaian dengan bentuk bilinear yang diberikan oleh matriks Hesse yang diskalakan
\zusatzklammer {oleh faktor di depan} {} {}
pada $\R^{n-1}$, apabila vektor yang sama disubstitusikan pada kedua tempat.
Menurut
Lemma 38.10 (Lineare Algebra (Osnabrück 2024-2025)),
kedua bentuk bilinear tersebut berimpit untuk sembarang pasangan argumen. Akan tetapi, ini menentukan bentuk fundamental kedua, bukan langsung matriks operator dalam basis grafik. Jika \mathl{A}{} adalah matriks $L_P$ dalam basis grafik dan \mathl{G=I_{n-1}+gg^{\mathsf T}}{} adalah matriks metriknya, maka matriks bentuk bilinear di atas adalah \mathl{GA}{.} Oleh karena itu,

\mathdisp {GA={ \frac{   1  }{  \sqrt{1+\lVert g\rVert^2}  } }\operatorname{Hess}_{   }\, f} { , }

sehingga

\mathdisp {A=G^{-1}{ \frac{   1  }{  \sqrt{1+\lVert g\rVert^2}  } }\operatorname{Hess}_{   }\, f} { . }

Catatan edisi: sumber menghilangkan faktor invers metrik G dan dengan demikian menyamakan matriks bentuk fundamental kedua dengan matriks pemetaan Weingarten. Rumus dan simpulan di atas memperbaiki pembedaan tersebut.

}


\chapter{Lembar Kerja 4}
\setcounter{section}{4}

  \zwischenueberschrift{Soal latihan}

\inputaufgabe
{}
{

Misalkan

\mavergleichskette
{\vergleichskette
{ h(x_1 , \ldots , x_n)
}
{ = }{ \sum_{i = 1}^n a_ix_i
}
{  }{
}
{    }{
}
{   }{
}
}
{}{}{}
adalah suatu
\definitionsverweis {bentuk linear}{}{}
$\neq 0$ dan

\mavergleichskette
{\vergleichskette
{ Y
}
{ = }{ h^{-1}(c)
}
{  }{
}
{    }{
}
{   }{
}
}
{}{}{}
adalah hipermuka yang bersesuaian. Buktikan bahwa
\definitionsverweis {pemetaan Weingarten}{}{}
merupakan pemetaan nol.

}
{} {}

\inputaufgabe
{}
{

Misalkan

\mavergleichskettedisp
{\vergleichskette
{ Y
}
{ =} { \left\{  \left(  x   , \,   y  , \,  z                 \right)   \mid   x^2+y^2 = z^2 , \,  \left(  x   , \,   y  , \,  z                 \right)  \neq  \left(  0   , \,   0  , \,   0                 \right)       \right\}
}
{ \subseteq} { W
}
{ =} { \R^3 \setminus \{  \left(  0   , \,   0  , \,   0                 \right) \}
}
{ } {
}
}
{}{}{}
adalah kerucut standar, dan misalkan

\mavergleichskette
{\vergleichskette
{ P
}
{ = }{ \left(  x   , \,   y  , \,  z                 \right)
}
{ \in }{ Y
}
{    }{
}
{   }{
}
}
{}{}{.}
Misalkan

\mavergleichskettedisp
{\vergleichskette
{ \gamma(t)
}
{ =} { \left(  x\cos t-y\sin t   , \,   x\sin t+y\cos t  , \,  z                 \right)
}
{ } {
}
{ } {
}
{ } {
}
}
{}{}{.}

Catatan edisi: kurva pada sumber tidak selalu melalui P dan tidak selalu tetap berada pada kerucut. Rumus di atas menggantinya dengan rotasi standar pada bidang xy.
Tentukan

\mathdisp {\operatorname{lim}_{   t  \rightarrow  0   } \,  { \frac{  N( \gamma(t) ) -  N( \gamma(0) )  }{ t } }} { , }

dengan $N$ medan normal satuan yang diperoleh dari gradien $x^2+y^2-z^2$.

}
{} {}

\inputaufgabe
{}
{

Misalkan

\mavergleichskettedisp
{\vergleichskette
{ Y
}
{ =} { \left\{  \left(  x   , \,   y  , \,  z                 \right)   \mid   x^2+y^2 =1         \right\}
}
{ \subseteq} { \R^3
}
{ } {
}
{ } {
}
}
{}{}{}
adalah silinder standar, dan misalkan

\mavergleichskette
{\vergleichskette
{ P
}
{ = }{ \left(  1   , \,   0  , \,  z                 \right)
}
{ \in }{ Y
}
{    }{
}
{   }{
}
}
{}{}{.}
Misalkan

\mavergleichskettedisp
{\vergleichskette
{ \gamma(t)
}
{ =} { \left(  \cos  t   , \,   \sin  t  , \,  z                 \right)
}
{ } {
}
{ } {
}
{ } {
}
}
{}{}{.}
Tentukan

\mathdisp {\operatorname{lim}_{   t  \rightarrow  0   } \,  { \frac{  N( \gamma(t) ) -  N( \gamma(0) )  }{ t } }} { , }

dengan $N$ medan normal satuan yang diperoleh dari gradien $x^2+y^2-1$.

}
{} {}

\inputaufgabe
{}
{

Misalkan

\mavergleichskettedisp
{\vergleichskette
{ Y
}
{ =} { \left\{  \left(  x   , \,   y  , \,  z                 \right)   \mid   x^2+y^2 =1         \right\}
}
{ \subseteq} { \R^3
}
{ } {
}
{ } {
}
}
{}{}{}
dan misalkan

\mavergleichskette
{\vergleichskette
{ P
}
{ \in }{ Y
}
{  }{
}
{    }{
}
{   }{}
}
{}{}{}
adalah suatu titik. Buktikan bahwa
\definitionsverweis {pemetaan Weingarten}{}{}
$L_P$ tidak bijektif, tetapi juga bukan pemetaan nol.

}
{} {}

\inputaufgabe
{}
{

Buktikan bahwa suatu
\definitionsverweis {permukaan terdiferensialkan}{}{}

\mavergleichskette
{\vergleichskette
{ Y
}
{ \subseteq }{ \R^3
}
{  }{
}
{    }{
}
{   }{
}
}
{}{}{}
dapat memuat suatu garis

\mavergleichskette
{\vergleichskette
{ G
}
{ \subseteq }{ Y
}
{  }{
}
{    }{
}
{   }{
}
}
{}{}{}
yang vektor arahnya, jika dipandang sebagai vektor di
\definitionsverweis {ruang singgung}{}{}
$T_PY$ untuk suatu titik

\mavergleichskette
{\vergleichskette
{ P
}
{ \in }{ G
}
{  }{
}
{    }{
}
{   }{
}
}
{}{}{,}
tidak harus termasuk dalam
\definitionsverweis {kernel}{}{}
dari
\definitionsverweis {pemetaan Weingarten}{}{}
$L_P$.

}
{} {}

\inputaufgabe
{}
{

Misalkan

\mavergleichskette
{\vergleichskette
{ W
}
{ \subseteq }{ \R^n
}
{  }{
}
{    }{
}
{   }{
}
}
{}{}{}
\definitionsverweis {terbuka}{}{,}
\maabb {h} { W } { \R
} {}
adalah suatu
\definitionsverweis {fungsi yang terdiferensialkan secara kontinu sebanyak dua kali}{}{}
dan

\mavergleichskette
{\vergleichskette
{ Y
}
{ = }{ h^{-1}(c)
}
{  }{
}
{    }{
}
{   }{
}
}
{}{}{}
adalah
\definitionsverweis {serat}{}{}
di atas

\mavergleichskette
{\vergleichskette
{ c
}
{ \in }{ \R
}
{  }{
}
{    }{
}
{   }{
}
}
{}{}{,}
dengan $h$
\definitionsverweis {reguler}{}{}
di setiap titik pada $Y$.
Catatan edisi: sumber hanya mengasumsikan h kelas C1. Di sini hipotesis diperkuat menjadi h kelas C2 agar medan normal gradien dan pemetaan Weingarten yang ditanyakan terdefinisi sebagaimana pada kuliah.
Kita juga memandang $h$ sebagai pemetaan $\tilde{h}$ pada
\mathl{W \times \R^m}{}{,}
dengan $m$ variabel terakhir tidak muncul secara eksplisit dalam $\tilde{h}$. Misalkan

\mavergleichskettedisp
{\vergleichskette
{ Z
}
{ =} { \tilde{h}^{-1}(c)
}
{ =} { Y \times \R^m
}
{ \subseteq} { W \times \R^{m}
}
{ } {
}
}
{}{}{.}
Misalkan

\mavergleichskette
{\vergleichskette
{ Q
}
{ \in }{ Z
}
{  }{
}
{    }{
}
{   }{
}
}
{}{}{}
adalah suatu titik yang terletak di atas

\mavergleichskette
{\vergleichskette
{ P
}
{ \in }{ Y
}
{  }{
}
{    }{
}
{   }{
}
}
{}{}{.}
Bagaimana hubungan antara
\definitionsverweis {pemetaan-pemetaan Weingarten}{}{}
\mathkor {} {L_Q} {dan} {L_P} {?}

}
{} {}

\inputaufgabegibtloesung
{}
{

Misalkan
\mathkor {} {r} {dan} {R} {}
adalah bilangan real positif dengan

\mavergleichskette
{\vergleichskette
{ 0
}
{ < }{ r
}
{ < }{ R
}
{    }{
}
{   }{
}
}
{}{}{,}
dan kita tinjau torus
\zusatzklammer {tertanam} {} {}

\mavergleichskettedisp
{\vergleichskette
{ T
}
{ =} { \left\{ (x,y,z) \in \R^3  \mid   \left(  \sqrt{x^2+y^2} -R  \right)^2 +z^2  = r^2         \right\}
}
{ } {
}
{ } {
}
{ } {
}
}
{}{}{.}
\aufzaehlungdrei{Tentukan suatu
\definitionsverweis {medan normal satuan}{}{}
$N$ pada $T$.
}{Pada titik
\mathl{\left( R-r  , \,   0  , \,   0                 \right)}{} dan untuk vektor

\mavergleichskette
{\vergleichskette
{ v
}
{ = }{ \begin{pmatrix}  0  \\ 1 \\  0   \end{pmatrix}
}
{  }{
}
{    }{
}
{   }{
}
}
{}{}{}
tentukan
\definitionsverweis {turunan berarah}{}{}

\mathl{D_vN}{.}
}{Pada titik
\mathl{\left( R-r  , \,   0  , \,   0                 \right)}{} dan untuk vektor

\mavergleichskette
{\vergleichskette
{ w
}
{ = }{ \begin{pmatrix}  0  \\ 0\\  1   \end{pmatrix}
}
{  }{
}
{    }{
}
{   }{
}
}
{}{}{}
tentukan
\definitionsverweis {turunan berarah}{}{}

\mathl{D_w N}{.}
}

}
{} {}

\inputaufgabe
{}
{

Untuk grafik $Y$ dari fungsi

\mavergleichskettedisp
{\vergleichskette
{ f(x,y)
}
{ =} { xy
}
{ } {
}
{ } {
}
{ } {
}
}
{}{}{}
tentukan
\definitionsverweis {pemetaan Weingarten}{}{}
$L_P$ pada setiap titik

\mavergleichskettedisp
{\vergleichskette
{ P
}
{ =} { (x,y,f(x,y) )
}
{ } {
}
{ } {
}
{ } {
}
}
{}{}{.}
Tentukan nilai eigen, vektor eigen, dan suatu matriks diagonal untuk $L_P$.

}
{} {}

\inputaufgabe
{}
{

Misalkan

\mavergleichskette
{\vergleichskette
{ W
}
{ \subseteq }{ \R^n
}
{  }{
}
{    }{
}
{   }{
}
}
{}{}{}
\definitionsverweis {terbuka}{}{,}
\maabb {h} { W } { \R
} {}
adalah suatu
\definitionsverweis {fungsi yang terdiferensialkan secara kontinu sebanyak dua kali}{}{}
dan

\mavergleichskette
{\vergleichskette
{ Y
}
{ = }{ h^{-1}(c)
}
{  }{
}
{    }{
}
{   }{
}
}
{}{}{}
adalah
\definitionsverweis {serat}{}{}
di atas

\mavergleichskette
{\vergleichskette
{ c
}
{ \in }{ \R
}
{  }{
}
{    }{
}
{   }{
}
}
{}{}{,}
dengan $h$
\definitionsverweis {reguler}{}{}
di setiap titik pada $Y$.
Catatan edisi: sumber hanya mengasumsikan h kelas C1. Di sini hipotesis diperkuat menjadi h kelas C2 agar pemetaan Weingarten yang dibandingkan terdefinisi sebagaimana pada kuliah.
Misalkan
\maabbdisp {L} {\R^n } { \R^n
} {}
adalah suatu
\definitionsverweis {isometri linear}{}{,}

\mavergleichskette
{\vergleichskette
{ W'
}
{ = }{ L^{-1}(W)
}
{  }{
}
{    }{
}
{   }{
}
}
{}{}{}
dan

\mavergleichskettedisp
{\vergleichskette
{ Z
}
{ =} { L^{-1}(Y)
}
{ =} { (h \circ L)^{-1}(c)
}
{ } {
}
{ } {
}
}
{}{}{.}
Buktikan bahwa
\definitionsverweis {pemetaan Weingarten}{}{}
$L_P$ untuk

\mavergleichskette
{\vergleichskette
{ P
}
{ \in }{ Y
}
{  }{
}
{    }{
}
{   }{
}
}
{}{}{}
bersesuaian dengan pemetaan Weingarten untuk

\mavergleichskette
{\vergleichskette
{ Q
}
{ = }{ L^{-1}(P)
}
{ \in }{ Z
}
{    }{
}
{   }{
}
}
{}{}{,}
jika melalui $L$ ruang singgung $T_PY$ dan $T_QZ$ dihubungkan satu sama lain.

}
{} {}

\inputaufgabegibtloesung
{}
{

Untuk permukaan $Y$ yang diberikan oleh

\mavergleichskettedisp
{\vergleichskette
{ x^2+y^2+2z^2
}
{ =} { 4
}
{ } {
}
{ } {
}
{ } {
}
}
{}{}{}
dan titik

\mavergleichskettedisp
{\vergleichskette
{ P
}
{ =} { \left(  1  , \,   1  , \,   1                 \right)
}
{ \in} { Y
}
{ } {
}
{ } {
}
}
{}{}{}
tentukan suatu matriks diagonal untuk
\definitionsverweis {pemetaan Weingarten}{}{}
$L_P$.

}
{} {}

\inputaufgabe
{}
{

Misalkan

\mavergleichskette
{\vergleichskette
{ W
}
{ \subseteq }{ \R^n
}
{  }{
}
{    }{
}
{   }{
}
}
{}{}{}
\definitionsverweis {terbuka}{}{,}
\maabb {h} { W } { \R
} {}
adalah suatu
\definitionsverweis {fungsi yang terdiferensialkan secara kontinu sebanyak dua kali}{}{}
dan

\mavergleichskette
{\vergleichskette
{ Y
}
{ = }{ h^{-1}(c)
}
{  }{
}
{    }{
}
{   }{
}
}
{}{}{}
adalah
\definitionsverweis {serat}{}{}
di atas

\mavergleichskette
{\vergleichskette
{ c
}
{ \in }{ \R
}
{  }{
}
{    }{
}
{   }{
}
}
{}{}{,}
dengan $h$
\definitionsverweis {reguler}{}{}
di setiap titik pada $Y$.
Bagaimana
\definitionsverweis {pemetaan Weingarten}{}{}
dan
\definitionsverweis {ruang-ruang eigennya}{}{}
berubah jika kita beralih dari suatu
\definitionsverweis {orientasi}{}{}
pada $Y$ ke orientasi yang berlawanan?

}
{} {}

  \zwischenueberschrift{Soal untuk dikumpulkan}

\inputaufgabe
{4}
{

Dalam
Contoh 4.5,
tentukan matriks representasi
\definitionsverweis {pemetaan Weingarten}{}{}
untuk titik-titik berbentuk
\mathl{\begin{pmatrix}  0  \\ y \\ z   \end{pmatrix}}{} terhadap basis
\mathkor {} {\begin{pmatrix}  y  \\ 0 \\  0   \end{pmatrix}} {dan} {\begin{pmatrix}  0  \\ z \\  y   \end{pmatrix}} {}
dari ruang singgung. Tentukan nilai-nilai eigen dan suatu basis yang terdiri atas vektor-vektor eigen.

}
{} {}

\inputaufgabe
{5}
{

Untuk permukaan $Y$ yang diberikan oleh

\mavergleichskettedisp
{\vergleichskette
{ x^2+2y^2+5z^2
}
{ =} { 11
}
{ } {
}
{ } {
}
{ } {
}
}
{}{}{}
dan titik

\mavergleichskettedisp
{\vergleichskette
{ P
}
{ =} { \left(  2  , \,   1  , \,   1                 \right)
}
{ \in} { Y
}
{ } {
}
{ } {
}
}
{}{}{}
tentukan suatu matriks diagonal untuk
\definitionsverweis {pemetaan Weingarten}{}{}
$L_P$.

}
{} {}

\inputaufgabe
{4}
{

Untuk hipermuka yang diberikan oleh

\mavergleichskettedisp
{\vergleichskette
{ x^2+2y^2+3z^2+4w^2
}
{ =} { 10
}
{ } {
}
{ } {
}
{ } {
}
}
{}{}{}

\mavergleichskette
{\vergleichskette
{ Y
}
{ \subseteq }{ \R^4
}
{  }{
}
{    }{
}
{   }{
}
}
{}{}{}
dan titik

\mavergleichskettedisp
{\vergleichskette
{ P
}
{ =} { \left(  1  , \,   1  , \,   1 , \,   1                \right)
}
{ \in} { Y
}
{ } {
}
{ } {
}
}
{}{}{}
tentukan suatu matriks untuk
\definitionsverweis {pemetaan Weingarten}{}{}
$L_P$ terhadap suatu
\definitionsverweis {basis}{}{}
yang sesuai dari
\definitionsverweis {ruang singgung}{}{}
$T_PY$.

}
{} {}

\inputaufgabe
{6}
{

Untuk grafik $Y$ dari fungsi

\mavergleichskettedisp
{\vergleichskette
{ f(x,y)
}
{ =} { x^3y^5
}
{ } {
}
{ } {
}
{ } {
}
}
{}{}{}
tentukan
\definitionsverweis {pemetaan Weingarten}{}{}
$L_P$ pada titik-titik

\mavergleichskette
{\vergleichskette
{ P_1
}
{ = }{ (0,0,0)
}
{  }{
}
{    }{
}
{   }{
}
}
{}{}{,}

\mavergleichskette
{\vergleichskette
{ P_2
}
{ = }{ (0,1,0)
}
{  }{
}
{    }{
}
{   }{
}
}
{}{}{,}

\mavergleichskette
{\vergleichskette
{ P_3
}
{ = }{ (2,1,8)
}
{  }{
}
{    }{
}
{   }{
}
}
{}{}{.}
Untuk masing-masing titik, tentukan nilai eigen, vektor eigen, dan suatu matriks diagonal untuk $L_P$.

}
{} {}


\section*{Solusi yang disediakan oleh sumber}
\addcontentsline{toc}{section}{Solusi yang disediakan oleh sumber}
\subsection*{Solusi Soal 4.7}
\input{generated/worksheet04_exercise07_solution.id.build.tex}
\subsection*{Solusi Soal 4.10}
\input{generated/worksheet04_exercise10_solution.id.build.tex}

\part{Unit 5}
\chapter[Kuliah 5]{Kuliah 5: Kelengkungan Utama}
\input{generated/lecture05.id.build.tex}

\chapter{Lembar Kerja 5}
\setcounter{section}{5}

  \zwischenueberschrift{Soal latihan}

\inputaufgabegibtloesung
{}
{

Misalkan

\mavergleichskettedisp
{\vergleichskette
{ f(x,y)
}
{ =} { ax^2+by^2+cxy
}
{ } {
}
{ } {
}
{ } {
}
}
{}{}{.}
Tentukan
\definitionsverweis {kelengkungan utama}{}{}
dan
\definitionsverweis {arah kelengkungan utama}{}{}
pada
\definitionsverweis {grafik}{}{}
fungsi $f$ di titik nol.

}
{} {}

\inputaufgabe
{}
{

Tentukan
\definitionsverweis {kelengkungan utama}{}{}
dan
\definitionsverweis {arah kelengkungan utama}{}{}
pada setiap titik
\definitionsverweis {grafik}{}{}
fungsi
\maabbeledisp {f} {\R^{n-1}} {\R
} { \left(  x_1   , \,   \ldots  , \,   x_{n-1}                 \right) } { a_1x_1^2  + \cdots + a_{n-1} x_{n-1}^2
} {.}

}
{} {}

\inputaufgabe
{}
{

Misalkan

\mavergleichskettedisp
{\vergleichskette
{ f(x,y,z)
}
{ =} { ax^2 +by^2 +c z^2 +dxy
}
{ } {
}
{ } {
}
{ } {
}
}
{}{}{.}
Hitung
\definitionsverweis {polinom karakteristik}{}{}
dari
\definitionsverweis {pemetaan Weingarten}{}{}
pada
\definitionsverweis {grafik}{}{}
fungsi $f$ di titik nol. Hitung pula
\definitionsverweis {kelengkungan utama}{}{}
dan
\definitionsverweis {arah kelengkungan utama}{}{.}

}
{} {}

\inputaufgabe
{}
{

Misalkan

\mavergleichskettedisp
{\vergleichskette
{ f(x,y,z)
}
{ =} { ax^2 +by^2 +c z^2 +dxy + e yz
}
{ } {
}
{ } {
}
{ } {
}
}
{}{}{.}
Hitung
\definitionsverweis {polinom karakteristik}{}{}
dari
\definitionsverweis {pemetaan Weingarten}{}{}
pada
\definitionsverweis {grafik}{}{}
fungsi $f$ di titik nol. Dalam kasus ini, seberapa sulitkah menentukan
\definitionsverweis {kelengkungan utama}{}{?}

}
{} {}

\inputaufgabe
{}
{

Misalkan

\mavergleichskettedisp
{\vergleichskette
{ f(x,y)
}
{ =} { x^2+y^3
}
{ } {
}
{ } {
}
{ } {
}
}
{}{}{.}
Tentukan
\definitionsverweis {kelengkungan utama}{}{,}
\definitionsverweis {arah kelengkungan utama}{}{}
dan
\definitionsverweis {kelengkungan Gauss}{}{}
pada
\definitionsverweis {grafik}{}{}
fungsi $f$ di setiap titik

\mavergleichskette
{\vergleichskette
{ P
}
{ = }{ (Q,f(Q))
}
{  }{
}
{    }{
}
{   }{
}
}
{}{}{,}

\mavergleichskette
{\vergleichskette
{ Q
}
{ \in }{ \R^2
}
{  }{
}
{    }{
}
{   }{
}
}
{}{}{.}

}
{} {}

\inputaufgabe
{}
{

Misalkan

\mavergleichskette
{\vergleichskette
{ V
}
{ \subseteq }{ \R^{n-1}
}
{  }{
}
{    }{
}
{   }{
}
}
{}{}{}
\definitionsverweis {terbuka}{}{}
dan misalkan
\maabb {f} { V } {\R
} {}
adalah suatu
\definitionsverweis {fungsi yang dua kali terdiferensialkan secara kontinu}{}{.}
Buktikan bahwa
\definitionsverweis {kelengkungan utama}{}{}
dan
\definitionsverweis {arah kelengkungan utama}{}{}
pada
\definitionsverweis {grafik}{}{}
fungsi $f$ di suatu titik

\mavergleichskette
{\vergleichskette
{ P
}
{ = }{ (Q,f(Q))
}
{  }{
}
{    }{
}
{   }{
}
}
{}{}{}
hanya bergantung pada turunan-turunan parsial pertama dan kedua dari $f$ di $Q$.

}
{} {}

\inputaufgabe
{}
{

Misalkan

\mavergleichskette
{\vergleichskette
{ W
}
{ \subseteq }{ \R^n
}
{  }{
}
{    }{
}
{   }{
}
}
{}{}{}
\definitionsverweis {terbuka}{}{,}
\maabb {h} { W } { \R
} {}
adalah suatu
\definitionsverweis {fungsi yang dua kali terdiferensialkan secara kontinu}{}{}
dan

\mavergleichskette
{\vergleichskette
{ Y
}
{ = }{ h^{-1}(c)
}
{  }{
}
{    }{
}
{   }{
}
}
{}{}{}
adalah
\definitionsverweis {serat}{}{}
di atas

\mavergleichskette
{\vergleichskette
{ c
}
{ \in }{ \R
}
{  }{
}
{    }{
}
{   }{
}
}
{}{}{,}
dengan $h$
\definitionsverweis {reguler}{}{}
di setiap titik pada $Y$.
Catatan edisi: sumber hanya mengasumsikan $h$ kelas $C^1$, padahal kelengkungan utama dan arah utamanya memerlukan pemetaan Weingarten. Hipotesis di atas diperkuat menjadi $C^2$ agar semua objek yang ditanyakan terdefinisi sebagaimana pada kuliah.
Kita juga memandang $h$ sebagai suatu pemetaan $\tilde{h}$ pada
\mathl{W \times \R^m}{},
dengan $m$ variabel terakhir tidak muncul secara eksplisit dalam $\tilde{h}$. Misalkan

\mavergleichskettedisp
{\vergleichskette
{ Z
}
{ =} { \tilde{h}^{-1}(c)
}
{ =} { Y \times \R^m
}
{ \subseteq} { W \times \R^{m}
}
{ } {
}
}
{}{}{.}
Misalkan

\mavergleichskette
{\vergleichskette
{ Q
}
{ \in }{ Z
}
{  }{
}
{    }{
}
{   }{
}
}
{}{}{}
adalah suatu titik yang terletak di atas

\mavergleichskette
{\vergleichskette
{ P
}
{ \in }{ Y
}
{  }{
}
{    }{
}
{   }{
}
}
{}{}{.}
Bagaimana hubungan antara
\definitionsverweis {kelengkungan utama}{}{}
dan
\definitionsverweis {arah kelengkungan utama}{}{}
dari $Z$ di $Q$ dengan yang dari $Y$ di $P$?

}
{} {}

Misalkan

\mavergleichskette
{\vergleichskette
{ W
}
{ \subseteq }{ \R^n
}
{  }{
}
{    }{
}
{   }{
}
}
{}{}{}
\definitionsverweis {terbuka}{}{,}
\maabb {h} { W } { \R
} {}
adalah suatu
\definitionsverweis {fungsi yang dua kali terdiferensialkan secara kontinu}{}{}
dan

\mavergleichskette
{\vergleichskette
{ Y
}
{ = }{ h^{-1}(c)
}
{  }{
}
{    }{
}
{   }{
}
}
{}{}{}
adalah
\definitionsverweis {serat}{}{}
di atas

\mavergleichskette
{\vergleichskette
{ c
}
{ \in }{ \R
}
{  }{
}
{    }{
}
{   }{
}
}
{}{}{,}
dengan $h$
\definitionsverweis {reguler}{}{}
di setiap titik pada $Y$. Untuk suatu titik

\mavergleichskette
{\vergleichskette
{ P
}
{ \in }{ Y
}
{  }{
}
{    }{
}
{   }{
}
}
{}{}{}
\definitionsverweis {determinan}{}{}
dari
\definitionsverweis {pemetaan Weingarten}{}{}
$L_P$ disebut
\definitionswort {kelengkungan Gauss--Kronecker}{}
dari $Y$ di $P$.

Kelengkungan Gauss--Kronecker merupakan perumuman langsung dari kelengkungan Gauss.

\inputaufgabe
{}
{

Misalkan

\mavergleichskette
{\vergleichskette
{ W
}
{ \subseteq }{ \R^n
}
{  }{
}
{    }{
}
{   }{
}
}
{}{}{}
\definitionsverweis {terbuka}{}{,}
\maabb {h} { W } { \R
} {}
adalah suatu
\definitionsverweis {fungsi yang dua kali terdiferensialkan secara kontinu}{}{}
dan

\mavergleichskette
{\vergleichskette
{ Y
}
{ = }{ h^{-1}(c)
}
{  }{
}
{    }{
}
{   }{
}
}
{}{}{}
adalah
\definitionsverweis {serat}{}{}
di atas

\mavergleichskette
{\vergleichskette
{ c
}
{ \in }{ \R
}
{  }{
}
{    }{
}
{   }{
}
}
{}{}{,}
dengan $h$
\definitionsverweis {reguler}{}{}
di setiap titik pada $Y$. Misalkan

\mavergleichskette
{\vergleichskette
{ P
}
{ \in }{ Y
}
{  }{
}
{    }{
}
{   }{
}
}
{}{}{.}
Buktikan bahwa hasil kali
\definitionsverweis {kelengkungan utama}{}{}
dari $Y$ di $P$ sama dengan
\definitionsverweis {kelengkungan Gauss--Kronecker}{}{}
dari $Y$ di $P$.

}
{} {}

\inputaufgabe
{}
{

Misalkan

\mavergleichskettedisp
{\vergleichskette
{ h(x,y,z,w)
}
{ =} { x^2+y^3+z^4+w^5
}
{ } {
}
{ } {
}
{ } {
}
}
{}{}{}
dan

\mavergleichskette
{\vergleichskette
{ Y
}
{ = }{ h^{-1}(1)
}
{  }{
}
{    }{
}
{   }{
}
}
{}{}{.}
Untuk titik

\mavergleichskettedisp
{\vergleichskette
{ P
}
{ =} { \left(  1   , \,   -1  , \,   1 , \,   0                \right)
}
{ \in} { Y
}
{ } {
}
{ } {
}
}
{}{}{}
tentukan
\definitionsverweis {pemetaan Weingarten}{}{}
$L_P$,
\definitionsverweis {polinom karakteristik}{}{}
darinya, serta
\definitionsverweis {kelengkungan Gauss--Kronecker}{}{}
dari $Y$ di $P$.

}
{} {}

\inputaufgabe
{}
{

Misalkan

\mavergleichskettedisp
{\vergleichskette
{ h(x,y,z)
}
{ =} { x^2+y^3+z^5
}
{ } {
}
{ } {
}
{ } {
}
}
{}{}{}
pada

\mavergleichskette
{\vergleichskette
{ W
}
{ = }{ \R^3 \setminus \{0\}
}
{  }{
}
{    }{
}
{   }{
}
}
{}{}{}
dan

\mavergleichskette
{\vergleichskette
{ Y
}
{ = }{ h^{-1}(0)
}
{  }{
}
{    }{
}
{   }{
}
}
{}{}{.}
Tentukan
\definitionsverweis {kelengkungan normal}{}{}
di titik

\mavergleichskette
{\vergleichskette
{ P
}
{ = }{ \begin{pmatrix}  1  \\ -1\\  0   \end{pmatrix}
}
{  }{
}
{    }{
}
{   }{
}
}
{}{}{}
dalam arah vektor singgung ternormalisasi
\mathl{{ \frac{   1  }{  \sqrt{29}  } }  \begin{pmatrix}  3  \\ -2\\  4   \end{pmatrix}}{.}

}
{} {}

  \zwischenueberschrift{Soal untuk dikumpulkan}

\inputaufgabe
{4}
{

Misalkan

\mavergleichskettedisp
{\vergleichskette
{ f(x,y)
}
{ =} { x^3 +x^2y+3xy^2-y^4
}
{ } {
}
{ } {
}
{ } {
}
}
{}{}{.}
Tentukan
\definitionsverweis {kelengkungan utama}{}{,}
\definitionsverweis {arah kelengkungan utama}{}{}
dan
\definitionsverweis {kelengkungan Gauss}{}{}
pada
\definitionsverweis {grafik}{}{}
fungsi $f$ di setiap titik

\mavergleichskette
{\vergleichskette
{ P
}
{ = }{ (Q,f(Q))
}
{  }{
}
{    }{
}
{   }{
}
}
{}{}{,}

\mavergleichskette
{\vergleichskette
{ Q
}
{ \in }{ \R^2
}
{  }{
}
{    }{
}
{   }{
}
}
{}{}{.}

}
{} {}

\inputaufgabe
{4}
{

Misalkan

\mavergleichskettedisp
{\vergleichskette
{ f(x,y,z)
}
{ =} { -7x^2 +3y^2 -5z^2+ 6xy
}
{ } {
}
{ } {
}
{ } {
}
}
{}{}{.}
Tentukan
\definitionsverweis {kelengkungan utama}{}{}
dan
\definitionsverweis {arah kelengkungan utama}{}{}
pada
\definitionsverweis {grafik}{}{}
fungsi $f$ di titik nol.

}
{} {}

\inputaufgabe
{6}
{

Misalkan

\mavergleichskettedisp
{\vergleichskette
{ h(x,y,z,w)
}
{ =} { x^2+yz+z^3+xyzw
}
{ } {
}
{ } {
}
{ } {
}
}
{}{}{}
dan

\mavergleichskette
{\vergleichskette
{ Y
}
{ = }{ h^{-1}(2)
}
{  }{
}
{    }{
}
{   }{
}
}
{}{}{.}
Untuk titik

\mavergleichskettedisp
{\vergleichskette
{ P
}
{ =} { \left(  1   , \,   1  , \,   1 , \,   -1                \right)
}
{ \in} { Y
}
{ } {
}
{ } {
}
}
{}{}{}
tentukan
\definitionsverweis {pemetaan Weingarten}{}{}
$L_P$,
\definitionsverweis {polinom karakteristik}{}{}
darinya, serta
\definitionsverweis {kelengkungan Gauss--Kronecker}{}{}
dari $Y$ di $P$.

}
{} {Petunjuk: Selesaikan persamaan terhadap suatu variabel yang sesuai dan perlakukan hipermuka tersebut sebagai grafik dalam variabel-variabel lainnya.}

\inputaufgabe
{6 (2+2+2)}
{

Misalkan

\mavergleichskettedisp
{\vergleichskette
{ h(x,y,z)
}
{ =} { x^3+y^3+z^3
}
{ } {
}
{ } {
}
{ } {
}
}
{}{}{}
pada

\mavergleichskette
{\vergleichskette
{ W
}
{ = }{ \R^3 \setminus \{0\}
}
{  }{
}
{    }{
}
{   }{
}
}
{}{}{}
dan

\mavergleichskette
{\vergleichskette
{ Y
}
{ = }{ h^{-1}(0)
}
{  }{
}
{    }{
}
{   }{
}
}
{}{}{.}
Misalkan

\mavergleichskette
{\vergleichskette
{ P
}
{ = }{ \begin{pmatrix}  1  \\ 1\\  - \sqrt[3]{2}   \end{pmatrix}
}
{ \in }{ Y
}
{    }{
}
{   }{
}
}
{}{}{}
dan

\mavergleichskettedisp
{\vergleichskette
{ v
}
{ =} { \begin{pmatrix}  1  \\ -1\\  0   \end{pmatrix}
}
{ \in} { T_PY
}
{ } {
}
{ } {
}
}
{}{}{.}
\aufzaehlungdrei{Tentukan suatu persamaan bidang untuk
\definitionsverweis {bidang normal}{}{}
yang bersesuaian dengan $v$.
}{Tentukan
\definitionsverweis {kelengkungan normal}{}{}
di $P$ dalam arah vektor singgung ternormalisasi
\mathl{{ \frac{   v  }{  \Vert { v } \Vert  } }}{.}
}{Tentukan suatu kurva berparametrisasi panjang busur
\maabbdisp {\gamma} { I } { Y
} {}
dengan

\mavergleichskette
{\vergleichskette
{ \gamma (0)
}
{ = }{ P
}
{  }{
}
{    }{
}
{   }{
}
}
{}{}{,}

\mavergleichskette
{\vergleichskette
{ \gamma' (0)
}
{ = }{ { \frac{   v  }{  \Vert { v } \Vert  } }
}
{  }{
}
{    }{
}
{   }{
}
}
{}{}{.}
Verifikasikan berlakunya
Teorema 5.14
dalam situasi ini.
}

Catatan edisi: pada butir ketiga, orientasikan bidang normal dengan menetapkan pasangan $(v/\lVert v\rVert,N(P))$ sebagai basis positif, sehingga normal satuan kurva irisan di dalam bidang yang bersesuaian dengan arah singgung $v/\lVert v\rVert$ adalah $N(P)$. Tanpa konvensi ini, verifikasi kelengkungan bertanda hanya menentukan kesamaan nilai mutlak.

}
{} {}

\inputaufgabe
{2}
{

Buatlah sketsa suatu permukaan terdiferensialkan $Y$ di ruang dan suatu bidang normal $E$ yang irisannya dengan permukaan tersebut menghasilkan suatu kurva yang tidak reguler pada setidaknya satu titik.

}
{} {}


\section*{Solusi yang disediakan oleh sumber}
\addcontentsline{toc}{section}{Solusi yang disediakan oleh sumber}
\subsection*{Solusi Soal 5.1}
\input{generated/worksheet05_exercise01_solution.id.build.tex}

\part{Unit 6}
\chapter[Kuliah 6]{Kuliah 6: Turunan Kovarian, Medan Vektor Paralel, dan Transpor Paralel}
\input{generated/lecture06.id.build.tex}

\chapter{Lembar Kerja 6}
\input{generated/worksheet06.id.build.tex}

\section*{Solusi yang disediakan oleh sumber}
\addcontentsline{toc}{section}{Solusi yang disediakan oleh sumber}
\subsection*{Solusi Soal 6.2}
\aufzaehlungdrei{Matriks Jacobi tersebut adalah

\mathdisp {\begin{pmatrix}  0   &   1  &  0  \\  -1 &  0  &  0 \\ 0  &  0 &  0  \end{pmatrix}} { . }

}{Sebagai medan normal, kita ambil $\begin{pmatrix}  x  \\ y \\ z   \end{pmatrix}$. Di

\mavergleichskette
{\vergleichskette
{ P
}
{ = }{ \begin{pmatrix}  0  \\ 1 \\  0   \end{pmatrix}
}
{  }{
}
{    }{
}
{   }{
}
}
{}{}{,}
vektor-vektor
\mathkor {} {\begin{pmatrix}  1  \\ 0 \\  0   \end{pmatrix}} {dan} {\begin{pmatrix}  0  \\ 0\\  1   \end{pmatrix}} {}
membentuk suatu basis ruang singgung. Kita mempunyai

\mavergleichskettedisp
{\vergleichskette
{ \begin{pmatrix}  0   &   1  &  0  \\  -1 &  0  &  0 \\ 0  &  0 &  0  \end{pmatrix} \begin{pmatrix}  1  \\ 0 \\  0   \end{pmatrix}
}
{ =} { \begin{pmatrix}  0  \\ -1\\  0   \end{pmatrix}
}
{ } {
}
{ } {
}
{ } {
}
}
{}{}{}
dan

\mavergleichskettedisp
{\vergleichskette
{  \begin{pmatrix}  0   &   1  &  0  \\  -1 &  0  &  0 \\ 0  &  0 &  0  \end{pmatrix} \begin{pmatrix}  0  \\ 0\\  1   \end{pmatrix}
}
{ =} { \begin{pmatrix}  0  \\ 0\\  0   \end{pmatrix}
}
{ } {
}
{ } {
}
{ } {
}
}
{}{}{.}
Di sini,
\mathl{\begin{pmatrix}  0  \\ -1\\  0   \end{pmatrix}}{} merupakan kelipatan vektor normal, sehingga proyeksi ortogonalnya ke ruang singgung juga merupakan vektor nol. Dengan demikian, pemetaan
\mathl{v \mapsto \nabla_vF}{} adalah pemetaan nol.
}{Di

\mavergleichskette
{\vergleichskette
{ Q
}
{ = }{ \begin{pmatrix}  0  \\ 0\\  1   \end{pmatrix}
}
{  }{
}
{    }{
}
{   }{
}
}
{}{}{,}
vektor-vektor
\mathkor {} {\begin{pmatrix}  1  \\ 0 \\  0   \end{pmatrix}} {dan} {\begin{pmatrix}  0  \\ 1 \\  0   \end{pmatrix}} {}
membentuk suatu basis ruang singgung. Kita mempunyai

\mavergleichskettedisp
{\vergleichskette
{ \begin{pmatrix}  0   &   1  &  0  \\  -1 &  0  &  0 \\ 0  &  0 &  0  \end{pmatrix} \begin{pmatrix}  1  \\ 0 \\  0   \end{pmatrix}
}
{ =} { \begin{pmatrix}  0  \\ -1\\  0   \end{pmatrix}
}
{ } {
}
{ } {
}
{ } {
}
}
{}{}{}
dan

\mavergleichskettedisp
{\vergleichskette
{ \begin{pmatrix}  0   &   1  &  0  \\  -1 &  0  &  0 \\ 0  &  0 &  0  \end{pmatrix} \begin{pmatrix}  0  \\ 1 \\  0   \end{pmatrix}
}
{ =} { \begin{pmatrix}  1  \\ 0 \\  0   \end{pmatrix}
}
{ } {
}
{ } {
}
{ } {
}
}
{}{}{.}
Kedua vektor citra ini sudah tangensial pada $Y$ di $Q$. Suatu matriks yang mendeskripsikan
\mathl{v \mapsto \nabla_vF}{} adalah

\mathdisp {\begin{pmatrix}  0   &   1  \\  -1  &  0 \end{pmatrix}} { . }

\emph{Catatan edisi: Tanda pada kedua kolom diperbaiki agar matriks ini sesuai dengan citra kedua vektor basis yang dihitung tepat di atasnya.}

}

\subsection*{Solusi Soal 6.6}
Kita mempunyai

\mavergleichskettealign
{\vergleichskettealign
{ (N (\gamma(t)) \times F(t))'
}
{ =} { (N (\gamma(t)))' \times F(t) +  N (\gamma(t)) \times F'(t)
}
{ =} { \left(  D N   \right)_{ \gamma(t)} \left(   \gamma'(t)   \right)  \times F(t) +  N (\gamma(t)) \times F'(t)
}
{ } {
}
{ } {
}
}
{}
{}{.}
Kita perlu membuktikan bahwa vektor ini selalu tegak lurus terhadap ruang singgung $T_{\gamma(t)}Y$. Karena $F'(t)$ tegak lurus terhadap ruang singgung, turunan tersebut merupakan kelipatan skalar vektor normal $N(\gamma(t))$; oleh karena itu, suku kedua sama dengan $0$ menurut
Fakta (3).
Pada suku di sebelah kiri, $F(t)$ bersifat tangensial menurut asumsi, sedangkan $\left(  D N   \right)_{ \gamma(t)} \left(   \gamma'(t)   \right)$ bersifat tangensial menurut
Fakta.
Oleh karena itu, menurut
Fakta (6),
suku di sebelah kiri tegak lurus terhadap ruang singgung. Jadi, medan vektor tersebut secara keseluruhan bersifat paralel.

\subsection*{Solusi Soal 6.9}
\input{generated/worksheet06_exercise09_solution.id.build.tex}

\part{Unit 7}
\chapter[Kuliah 7]{Kuliah 7: Konsep sebuah manifold}
\setcounter{section}{7}

  \zwischenueberschrift{Konsep sebuah manifold}

Dalam geometri diferensial, gagasan tentang manifold berada di pusat pembahasan. Sebagai contoh, kita meninjau Bumi
\zusatzklammer {permukaannya} {} {,}
yang dalam sejarah ilmu pengetahuan lama dianggap sebagai sebuah cakram, dan ada alasan yang baik untuk itu. Secara lokal Bumi tampak seperti sebuah bidang. Hal ini juga tercermin dalam peta yang dibuat tentangnya. Sebuah peta adalah sebuah \anfuehrung{lembaran}{,} yang titik-titiknya berada dalam bijeksi dengan suatu bagian permukaan Bumi. Khususnya untuk bagian-bagian kecil hal ini tidak bermasalah, tetapi pada peta yang harus menggambarkan bagian besar atau bahkan seluruh Bumi segera muncul pertanyaan tentang apa yang dipertahankan oleh peta dan apa yang tidak: pertanyaan tentang pelestarian panjang, pelestarian luas, pelestarian sudut, titik-titik yang hilang atau muncul berkali-kali, pertanyaan perpanjangan, pertanyaan kelengkungan, dan seterusnya.

\bild{ \begin{center}
\includegraphics[width=5.5cm]{ \bildeinlesungpng {Stereographic_projection_in_3D} {png} }
\end{center}
\bildtext {Proyeksi stereografis ketika bidang tidak diletakkan melalui khatulistiwa, melainkan melalui Kutub Selatan.} }

\bildlizenz { Stereographic projection in 3D.png } {} {Mark.Howison} {en.Wikipedia} {PD} {}

Mula-mula kita membahas \stichwort {proyeksi stereografis} {} dari permukaan bola.

\inputbeispiel{
}
{

Kita meninjau permukaan bola

\mavergleichskettedisp
{\vergleichskette
{ K
}
{ =} { \left\{ (x,y,z)  \mid   x^2+y^2+z^2 = 1        \right\}
}
{ } {
}
{ } {
}
{ } {
}
}
{}{}{}
dan menyebut titik

\mavergleichskette
{\vergleichskette
{ N
}
{ = }{ (0,0,1)
}
{  }{
}
{    }{
}
{   }{
}
}
{}{}{}
sebagai Kutub Utara dan titik

\mavergleichskette
{\vergleichskette
{ S
}
{ = }{ (0,0,-1)
}
{  }{
}
{    }{
}
{   }{
}
}
{}{}{}
sebagai Kutub Selatan. Sebuah titik

\mathbed {P=(x,y,z) \in K} {}
 {P \neq N} {}
 {} {} {} {,}
bersama Kutub Utara menentukan sebuah garis tunggal di ruang, yang diparameterkan oleh

\mathdisp {(tx,ty ,1+t(z-1))=(0,0,1)  + t ((x,y,z) -(0,0,1)), \, t \in \R} { , }

Vektor \mathl{(x,y,z) -(0,0,1)}{,} yang menentukan garis ini, tidak sejajar dengan bidang $x-y$; artinya garis tersebut mempunyai tepat satu titik potong dengan bidang ini. Titik itu diperoleh untuk parameter

\mavergleichskettedisp
{\vergleichskette
{ t
}
{ =} { { \frac{   1  }{  1-z } }
}
{ } {
}
{ } {
}
{ } {
}
}
{}{}{,}
yaitu titik

\mathdisp {\left(  { \frac{   x  }{  1-z  } }, { \frac{   y  }{  1-z  } } ,0  \right)} { . }

Kita rangkum konstruksi ini sebagai pemetaan
\maabbeledisp {\alpha} {K \setminus \{N\}} { \R^2
} {(x,y,z)} { \left(  { \frac{   x  }{  1-z  } }, { \frac{   y  }{  1-z  } }  \right)
} {}
Pemetaan ini secara intuitif jelas merupakan bijeksi, dan hal tersebut juga mudah dihitung dari rumus-rumusnya. Pemetaan invers diperoleh dengan menghubungkan sebuah titik \mathl{(u,v,0)}{} di bidang dengan Kutub Utara, lalu menghitung titik tembusnya $\neq N$ dengan permukaan bola. Hal ini menghasilkan syarat

\mavergleichskettealignhandlinks
{\vergleichskettealignhandlinks
{ \Vert {(0,0,1) +a((u,v,0) - (0,0,1) ) } \Vert^2
}
{ =} { \Vert {(au,av,1-a) } \Vert^2
}
{ =} { a^2u^2+a^2v^2+(1-a)^2
}
{ =} { 1
}
{ } {
}
}
{}
{}{,}
yang berujung pada

\mavergleichskettedisp
{\vergleichskette
{ a \left(  au^2+av^2 -2 +a  \right)
}
{ =} { 0
}
{ } {
}
{ } {
}
{ } {
}
}
{}{}{}
Solusinya adalah

\mavergleichskette
{\vergleichskette
{ a
}
{ = }{ 0
}
{  }{
}
{    }{
}
{   }{
}
}
{}{}{}
yang bersesuaian dengan Kutub Utara dan tidak kita minati, sehingga kita memperoleh

\mavergleichskette
{\vergleichskette
{ a
}
{ = }{ { \frac{   2  }{ u^2+v^2+1 } }
}
{  }{
}
{    }{
}
{   }{
}
}
{}{}{}
dan dengan demikian pemetaan
\maabbeledisp {} { \R^2 } { K \setminus \{N\}
} { (u,v) } { \left(  { \frac{   2u }{ u^2+v^2+1 } }, { \frac{   2v }{ u^2+v^2+1 } } ,1- { \frac{   2  }{ u^2+v^2+1 } }  \right)
} {.}
Jadi bidang riil adalah \definitionsverweis {homeomorf}{}{} dengan permukaan bola yang \anfuehrung{dilubangi}{} pada satu titik.

Pertimbangan yang sama dapat dilakukan untuk \mathl{K \setminus \{S\}}{}. Hal ini menghasilkan pemetaan
\maabbeledisp {\beta} { K \setminus \{S\} } { \R^2
} { (x,y,z) } { \left(  { \frac{   x  }{ z+1 } }, { \frac{   y  }{ z+1 } }  \right)
} {}
dengan pemetaan invers
\maabbeledisp {} { \R^2 } { K \setminus \{S\}
} { (s,t) } { \left(  { \frac{   2s }{ s^2+t^2+1 } }, { \frac{   2t }{ s^2+t^2+1 } }, -1+ { \frac{   2  }{ s^2+t^2+1 } }  \right)
} {.}
Kutub Selatan dipetakan ke pusat bidang Euklides oleh pemetaan pertama, sedangkan Kutub Utara dipetakan ke pusat bidang Euklides oleh pemetaan kedua. Kita menyebut kedua pemetaan, maupun pemetaan inversnya, sebagai peta. Kedua peta tersebut bersama-sama menutupi seluruh permukaan bola. Karena merupakan homeomorfisme, keduanya mempertahankan sifat-sifat topologis utama bola. Keduanya tidak cocok untuk geografi Bumi karena peta membutuhkan seluruh bidang dan sangat mendistorsi panjang.

Kedua peta itu sama baiknya. Titik-titik \zusatzklammer {dan secara umum gambar-gambar lain} {} {} pada satu peta dapat dengan mudah diubah ke peta yang lain. Namun, perlu diperhatikan bahwa kedua kutub hanya muncul pada salah satu peta. Titik-titik dari himpunan \mathl{K \setminus \{N,S\}}{} terdapat pada kedua peta, dan melalui keduanya berada dalam bijeksi dengan bidang yang dilubangi di pusatnya, \mathl{\R^2 \setminus \{0\}}{.} Pemetaan \stichwort {transisi} {} \zusatzklammer {atau \stichwort {pergantian peta} {}} {} {} diberikan oleh

\mavergleichskette
{\vergleichskette
{ \psi
}
{ = }{ \beta \circ \left(  \alpha^{-1}\!\mid_{\R^2 \setminus \{0\} }  \right)
}
{  }{
}
{    }{
}
{   }{
}
}
{}{}{}
Selanjutnya,

\mavergleichskettealignhandlinks
{\vergleichskettealignhandlinks
{ \psi(u,v)
}
{ =} { \beta \left(  { \frac{   2u }{ u^2+v^2+1 } }, { \frac{   2v }{ u^2+v^2+1 } } ,1- { \frac{   2  }{ u^2+v^2+1 } }  \right)
}
{ =} { \left(  { \frac{   { \frac{   2u }{ u^2+v^2+1 } }  }{  2 - { \frac{   2  }{ u^2+v^2+1 } }  } } , { \frac{   { \frac{   2v }{ u^2+v^2+1 } }  }{  2 - { \frac{   2  }{ u^2+v^2+1 } }  } }  \right)
}
{ =} { \left(  { \frac{   { \frac{   u  }{ u^2+v^2+1 } }  }{  { \frac{  (u^2+v^2+1) -1  }{ u^2+v^2+1 } }  } } , { \frac{   { \frac{   v  }{ u^2+v^2+1 } }  }{  { \frac{  (u^2+v^2+1) -1  }{ u^2+v^2+1 } }  } }  \right)
}
{ =} { \left(  { \frac{   u  }{ u^2+v^2  } } , { \frac{   v  }{ u^2+v^2 } }  \right)
}
}
{}
{}{.}
Pemetaan transisi ini tidak hanya menginduksi homeomorfisme antara \mathl{\R^2 \setminus \{0\}}{} dengan \mathl{\R^2 \setminus \{0\}}{\zusatzfussnote {Sebaiknya di sini tidak dikatakan \anfuehrung{dengan dirinya sendiri}{} karena kedua bidang peta sebaiknya dibayangkan sebagai bidang yang independen. Hubungan di antara keduanya muncul semata-mata karena keduanya menggambarkan permukaan bola yang sama} {.} {,}} sebagaimana langsung mengikuti dari fakta bahwa kedua pemetaan peta \mathkor {} {\alpha} {dan} {\beta} {} adalah homeomorfisme, melainkan bahkan sebuah \definitionsverweis {difeomorfisme}{}{.} Ini langsung terlihat dari rumusnya; tetapi tidak masuk akal mengatakan bahwa pemetaan peta adalah difeomorfisme, karena permukaan bola bukan himpunan terbuka di dalam $\R^3$. Yang masih kurang adalah sebuah \anfuehrung{struktur diferensiabel}{} pada permukaan ini agar kita dapat berbicara tentang difeomorfisme.

}
Sebuah \stichwort {manifold} {} adalah suatu bangun geometri yang secara \anfuehrung{lokal}{} tampak seperti ruang Euklides $\R^n$. Kita menganggap bangun ini sebagai sebuah
\definitionsverweis {ruang topologis}{}{}
dan memperjelas kata lokal dengan menyatakan bahwa terdapat tutupan oleh himpunan-himpunan terbuka yang homeomorfik dengan himpunan terbuka di $\R^n$. Walaupun selanjutnya kita bekerja dengan ruang topologis, kandungan intuitif pembahasan tidak berkurang jika ruang topologis dibayangkan sekadar sebagai ruang metrik.

  \zwischenueberschrift{Manifold diferensiabel}

\inputdefinition
{}
{

Sebuah
\definitionsverweis {ruang topologis}{}{}
\definitionsverweis {Hausdorff}{}{}
$M$ disebut \definitionswort {manifold topologis}{} berdimensi \definitionswort {dimensi}{} $n$, jika terdapat suatu
\definitionsverweis {tutupan terbuka}{}{}

\mavergleichskette
{\vergleichskette
{ M
}
{ = }{  \bigcup_{i \in I} U_i
}
{  }{
}
{    }{
}
{   }{
}
}
{}{}{}
sedemikian sehingga setiap $U_i$
\definitionsverweis {homeomorf}{}{}
dengan suatu
\definitionsverweis {himpunan bagian terbuka}{}{}
dari $\R^n$.

}

Untuk setiap titik

\mavergleichskette
{\vergleichskette
{P
}
{ \in }{ M
}
{  }{
}
{    }{
}
{   }{
}
}
{}{}{}
terdapat lingkungan terbuka

\mavergleichskette
{\vergleichskette
{P
}
{ \in }{ U
}
{ \subseteq }{ M
}
{    }{
}
{   }{
}
}
{}{}{,}
yang homeomorfik dengan suatu himpunan terbuka

\mavergleichskette
{\vergleichskette
{V
}
{ \subseteq }{ \R^n
}
{  }{
}
{    }{
}
{   }{
}
}
{}{}{}.
Misalkan
\maabbdisp {\varphi} {U} {V
} {}
sebuah homeomorfisme dan misalkan

\mavergleichskette
{\vergleichskette
{ Q
}
{ = }{ \varphi(P)
}
{  }{
}
{    }{
}
{   }{
}
}
{}{}{.}
Maka lingkungan bola terbuka

\mavergleichskette
{\vergleichskette
{Q
}
{ \in }{ U \left(  Q, \epsilon  \right)
}
{ \subseteq }{ V
}
{    }{
}
{   }{
}
}
{}{}{}
berpadanan dengan lingkungan terbuka

\mavergleichskette
{\vergleichskette
{ U'
}
{ = }{  \varphi^{-1}(U \left(  Q, \epsilon  \right))
}
{  }{
}
{    }{
}
{   }{
}
}
{}{}{}
dengan

\mavergleichskette
{\vergleichskette
{P
}
{ \in }{ U'
}
{ \subseteq }{ U
}
{    }{
}
{   }{
}
}
{}{}{,}
yang menurut konstruksi homeomorfik dengan bola terbuka. Karena itu manifold topologis juga dapat dicirikan sebagai ruang Hausdorff topologis yang \stichwort {secara lokal Euklides} {}.

\inputdefinition
{}
{

Misalkan $M$ adalah sebuah
\definitionsverweis {manifold topologis}{}{.}
Setiap
\definitionsverweis {homeomorfisme}{}{}
\maabbdisp {\varphi} { U } { V
} {,}
dengan

\mavergleichskette
{\vergleichskette
{ U
}
{ \subseteq }{ M
}
{  }{
}
{    }{
}
{   }{
}
}
{}{}{}
dan

\mavergleichskette
{\vergleichskette
{ V
}
{ \subseteq }{ \R^n
}
{  }{
}
{    }{
}
{   }{
}
}
{}{}{}
\definitionsverweis {terbuka}{}{}
adalah sebuah
\zusatzklammer {topologis} {} {}
\definitionswort {peta}{} untuk $M$.

}

Himpunan terbuka

\mavergleichskette
{\vergleichskette
{ U
}
{ \subseteq }{ M
}
{  }{
}
{    }{
}
{   }{
}
}
{}{}{}
kadang-kadang disebut \stichwort {daerah peta} {}, sedangkan

\mavergleichskette
{\vergleichskette
{ V
}
{ \subseteq }{ \R^n
}
{  }{
}
{    }{
}
{   }{
}
}
{}{}{}
kadang-kadang disebut \stichwort {citra peta} {.} Untuk sebuah peta
\maabbdisp {\varphi} {U} {V
} {}
dan himpunan bagian terbuka

\mavergleichskette
{\vergleichskette
{ U'
}
{ \subseteq }{ U
}
{  }{
}
{    }{
}
{   }{
}
}
{}{}{}
pemetaan terinduksi
\maabbdisp {\varphi \vert_{U'}} {U'} {\varphi(U')
} {}
juga merupakan peta. Kadang-kadang pemetaan invers juga disebut peta. Selain peta, kita juga berbicara tentang \stichwort {sistem koordinat lokal} {.} Melalui peta
\maabb {\varphi} {U} {V
} {}
koordinat pada

\mavergleichskette
{\vergleichskette
{V
}
{ \subseteq }{ \R^n
}
{  }{
}
{    }{
}
{   }{
}
}
{}{}{}
dipindahkan ke $U$. Koordinat ke-$j$
\zusatzklammer {yakni proyeksi ke-$j$} {} {}
\maabb {x_j} {V} {\R
} {}
menginduksi \zusatzklammer {koordinat lokal} {} {-}fungsi
\maabbdisp {x_j \circ \varphi} {U} {\R
} {}
\zusatzklammer {yang sering kembali dilambangkan dengan $x_j$} {} {,}
dan sebuah titik

\mavergleichskette
{\vergleichskette
{ Q
}
{ = }{ \left(  x_1   , \,   , \ldots ,  , \,   x_n                 \right)
}
{ \in }{ V
}
{    }{
}
{   }{
}
}
{}{}{}
berpadanan dengan titik

\mavergleichskette
{\vergleichskette
{P
}
{ = }{\varphi^{-1}(Q)
}
{  }{
}
{    }{
}
{   }{
}
}
{}{}{.}

\bild{ \begin{center}
\includegraphics[width=5.5cm]{ \bildeinlesungpng {Manifold_zahyou3} {png} }
\end{center}
\bildtext {} }

% TeX-safe display form for the moving list-of-figures argument.  The exact
% Unicode creator string remains in source/unit_media.json and the rights
% manifest; this ASCII rendering avoids an engine-dependent Unicode failure.
\bildlizenz { Manifold zahyou3.png } {} {132ninme} {ja. Wikipedia} {CC-by-sa 3.0} {}

\inputdefinition
{}
{

Misalkan $M$ adalah sebuah
\definitionsverweis {manifold topologis}{}{,}
misalkan

\mavergleichskette
{\vergleichskette
{  U_1, U_2
}
{ \subseteq }{ M
}
{  }{
}
{    }{
}
{   }{
}
}
{}{}{}
merupakan
\definitionsverweis {himpunan bagian terbuka}{}{}
dan
\maabb {\alpha_1} { U_1 } { V_1
} {} dan \maabb {\alpha_2} { U_2 } { V_2
} {}
merupakan
\definitionsverweis {peta}{}{}
\zusatzklammer {dengan

\mavergleichskettek
{\vergleichskettek
{  V_1,V_2
}
{ \subseteq }{\R^n
}
{  }{
}
{    }{
}
{   }{
}
}
{}{}{}
terbuka} {} {.}
Maka
\definitionsverweis {pemetaan}{}{}
\maabbdisp {\alpha_2  \circ \alpha_1^{-1}} { \alpha_1(U_1 \cap U_2) } { \alpha_2(U_1 \cap U_2)
} {} disebut \definitionswort {pemetaan transisi}{} untuk peta-peta tersebut.

}

Irisan \mathl{U_1 \cap U_2}{} adalah himpunan bagian terbuka tempat kedua peta didefinisikan dan tempat keduanya dapat dibandingkan. Secara lebih teliti, pada definisi tersebut kita harus berbicara tentang pembatasan $\alpha_1^{-1}$ pada himpunan bagian terbuka

\mavergleichskette
{\vergleichskette
{ \alpha_1(U_1 \cap U_2)
}
{ \subseteq }{ V_1
}
{  }{
}
{    }{
}
{   }{
}
}
{}{}{}.

\inputdefinition
{}
{

Misalkan

\mavergleichskette
{\vergleichskette
{ n
}
{ \in }{ \N
}
{  }{
}
{    }{
}
{   }{
}
}
{}{}{}
dan

\mavergleichskette
{\vergleichskette
{ k
}
{ \in }{ \overline{ \N }_+
}
{  }{
}
{    }{
}
{   }{
}
}
{}{}{.}
Sebuah
\definitionsverweis {ruang topologis}{}{}
\definitionsverweis {Hausdorff}{}{}
$M$ bersama suatu
\definitionsverweis {tutupan terbuka}{}{}

\mavergleichskette
{\vergleichskette
{ M
}
{ = }{ \bigcup_{i \in I} U_i
}
{  }{
}
{    }{
}
{   }{
}
}
{}{}{}
dan
\definitionsverweis {peta}{}{}
\maabbdisp {\alpha_i} {U_i } { V_i
} {}
dengan

\mavergleichskette
{\vergleichskette
{ V_i
}
{ \subseteq }{ \R^n
}
{  }{
}
{    }{
}
{   }{
}
}
{}{}{}
terbuka, sedemikian sehingga
\definitionsverweis {pemetaan transisi}{}{}
\maabbdisp {\alpha_j \circ (\alpha_i)^{-1}} {V_i \cap \alpha_i(U_i \cap U_j) } { V_j \cap \alpha_j(U_i \cap U_j)
} {}
merupakan
$C^k$-\definitionsverweis {difeomorfisme}{}{}
untuk semua

\mavergleichskette
{\vergleichskette
{ i,j
}
{ \in }{ I
}
{  }{
}
{    }{
}
{   }{
}
}
{}{}{}
dinamakan
\definitionswortpraemath {C^k}{ manifold }{}
atau \definitionswort {manifold diferensiabel}{}
\zusatzklammer {berdimensi \definitionswort {dimensi}{} $n$ dan berderajat diferensiabilitas $k$} {} {.} Kumpulan peta

\mathbed {(U_i,\alpha_i)} {}
 {i \in I} {}
 {} {} {} {,}
dinamakan juga
\definitionswortpraemath {C^k}{ atlas }{}
manifold tersebut.

}

Manifold diferensiabel khususnya merupakan sebuah
\definitionsverweis {manifold topologis}{}{.}
Menurut definisi kita, atlas merupakan bagian integral dari konsep manifold. Namun, yang lebih penting daripada atlas adalah \stichwort {struktur diferensiabel} {} yang ditentukan oleh atlas tersebut. Hal ini akan menjadi lebih jelas setelah kita memiliki konsep pemetaan diferensiabel antara dua manifold dan dapat berbicara tentang manifold yang difeomorf.

\inputdefinition
{}
{

Misalkan $M$ sebuah
\definitionsverweis {manifold diferensiabel}{}{.}
Sebuah
\definitionsverweis {himpunan bagian terbuka}{}{}

\mavergleichskette
{\vergleichskette
{ U
}
{ \subseteq }{ M
}
{  }{
}
{    }{
}
{   }{
}
}
{}{}{,}
yang dilengkapi dengan
\definitionsverweis {peta}{}{}
\definitionsverweis {terbatas}{}{}
disebut \definitionswort {submanifold terbuka}{.}

}

\inputbeispiel{}
{

Setiap himpunan bagian terbuka

\mavergleichskette
{\vergleichskette
{ V
}
{ \subseteq }{ \R^n
}
{  }{
}
{    }{
}
{   }{
}
}
{}{}{}
merupakan
$C^\infty$-\definitionsverweis {manifold}{}{,}
jika kita menggunakan
\definitionsverweis {identitas}{}{}
\maabb {\operatorname{Id}} { V } { V
} {}
sebagai
\definitionsverweis {peta}{}{}.
Satu-satunya
\definitionsverweis {pemetaan transisi}{}{}
adalah identitas itu sendiri, yang merupakan sebuah $C^\infty$-\definitionsverweis {difeomorfisme}{}{}.
Ini adalah manifold dengan sebuah
\definitionsverweis {atlas}{}{,}
yang terdiri atas satu peta. Namun kita juga dapat menggunakan atlas yang terdiri atas semua himpunan bagian terbuka

\mavergleichskette
{\vergleichskette
{ U
}
{ \subseteq }{ V
}
{  }{
}
{    }{
}
{   }{
}
}
{}{}{}
beserta peta identitas terkait $\varphi_U$. Pemetaan transisinya adalah identitas pada \mathl{U_1 \cap U_2}{.}

}

Kita telah membahas ruang metrik yang
\definitionsverweis {terhubung}{}{}
dalam konteks teorema nilai antara. Definisi yang sama juga kita gunakan untuk ruang topologis.

\inputdefinition
{}
{

Sebuah
\definitionsverweis {ruang topologis}{}{}
$X$ disebut \definitionswort {terhubung}{,} jika di dalam $X$ tepat ada dua himpunan
\zusatzklammer {yaitu $\emptyset$ dan seluruh ruang \mathlk{X \neq \emptyset}{}} {} {,}
yang sekaligus
\definitionsverweis {terbuka}{}{}
dan
\definitionsverweis {tertutup}{}{}.

}

Sering kali kita hanya tertarik pada manifold yang terhubung, terutama karena pada kasus tak terhubung komponen-komponen \anfuehrung{keterhubungan}{} dapat diteliti secara terpisah. Sekarang kita membahas secara singkat manifold berdimensi rendah.

\inputbeispiel{}
{

Pada sebuah
\definitionsverweis {manifold}{}{}
$M$ berdimensi nol, untuk setiap titik

\mavergleichskette
{\vergleichskette
{ P
}
{ \in }{ M
}
{  }{
}
{    }{
}
{   }{
}
}
{}{}{}
terdapat lingkungan terbuka

\mavergleichskette
{\vergleichskette
{ P
}
{ \in }{ U
}
{  }{
}
{    }{
}
{   }{
}
}
{}{}{,}
yang
\definitionsverweis {homeomorf}{}{}
dengan suatu himpunan terbuka dari

\mavergleichskette
{\vergleichskette
{ \R^0
}
{ = }{ \{0\}
}
{  }{
}
{    }{
}
{   }{
}
}
{}{}{}.
Artinya, himpunan satu titik \mathl{\{P\}}{} harus terbuka, sehingga $M$ harus membawa
\definitionsverweis {topologi diskret}{}{},
yaitu setiap himpunan bagiannya terbuka. Karena itu satu-satunya manifold berdimensi nol yang
\definitionsverweis {terhubung}{}{}
adalah himpunan satu titik.

}

\bild{ \begin{center}
\includegraphics[width=5.5cm]{ \bildeinlesungsvg {Circle_-_black_simple} {svg} }
\end{center}
\bildtext {Sebuah garis lingkaran adalah manifold kompak berdimensi satu} }

\bildlizenz { Circle - black simple.svg } {} {Dakdada} {Commons} {PD} {}

\inputbeispiel{}
{

Pada manifold berdimensi satu
\definitionsverweis {manifold}{}{}
mula-mula terdapat himpunan bagian terbuka dari $\R^1$. Himpunan-himpunan ini merupakan gabungan interval terbuka, dan \definitionsverweis {terhubung}{}{,} tepat ketika merupakan sebuah interval terbuka. Setiap interval terbuka, baik terbatas maupun tak terbatas, adalah
\definitionsverweis {homeomorf}{}{}
dan juga
\definitionsverweis {difeomorf}{}{}
dengan
\definitionsverweis {interval satuan terbuka}{}{}
\mathl{]0,1[}{} dan dengan bilangan-bilangan riil $\R$ sendiri. Interval tertutup

\mathbed {[a,b]} {dengan}
 {a < b} {}
 {} {} {} {}
bukan manifold, karena untuk titik-titik batasnya
\zusatzklammer {ujung-ujung interval} {} {}
tidak ada lingkungan terbuka yang homeomorfik dengan interval terbuka
\zusatzklammer {meskipun interval itu disebut \stichwort {manifold dengan batas} {}} {} {.}

Selain itu ada \stichwort {lingkaran} {}
\zusatzklammer {atau \stichwort {sfera} {}} {} {}
$S^1$ sebagai manifold terhubung berdimensi satu lainnya. Berlaku

\mavergleichskettedisp
{\vergleichskette
{ S^1
}
{ =} { \left\{ (x,y) \in \R^2  \mid   x^2+y^2 = 1         \right\}
}
{ } {
}
{ } {
}
{ } {
}
}
{}{}{.}
Untuk setiap titik

\mavergleichskette
{\vergleichskette
{ P
}
{ \in }{ S^1
}
{  }{
}
{    }{
}
{   }{
}
}
{}{}{}
berlaku
\mathl{S^1 \setminus \{P \}}{} homeomorfik dengan $\R$
\zusatzklammer {melalui proyeksi stereografis} {} {.}
Lingkaran tidak homeomorfik dengan $\R$, sebab lingkaran
\definitionsverweis {kompak}{}{\zusatzfussnote {Sebelumnya kita mendefinisikan kekompakan hanya untuk himpunan bagian di dalam $\R^n$; sebentar lagi akan terlihat bahwa ini merupakan konsep absolut yang tidak bergantung pada penyematan. Jadi, tidak mungkin menyematkan $\R$ secara homeomorfik ke dalam $\R^n$ sedemikian sehingga citranya kompak} {.} {}}
sedangkan bilangan riil tidak kompak. Selain \mathkor {} {S^1} {dan} {\R} {}, tidak ada manifold terhubung berdimensi satu lainnya yang memiliki
\definitionsverweis {basis topologi terhitung}{}{}
\zusatzklammer {hal ini dinyatakan tanpa bukti di sini} {} {.}

}
Mulai dari dimensi dua, tanpa syarat tambahan yang kuat, mustahil membuat gambaran menyeluruh tentang semua manifold.


\chapter{Lembar Kerja 7}
\setcounter{section}{7}

  \zwischenueberschrift{Soal latihan}

\inputaufgabe
{}
{

Tunjukkan bahwa untuk dua titik
\mathl{A,B \neq N,S}{} pada bola satuan $S^2$, selisih jarak yang diambil dalam dua peta standar stereografis, yaitu
\mathkor {} {d( \alpha_1(A),\alpha_1(B))} {dan} {d( \alpha_2(A),\alpha_2(B))} {,}
dapat dibuat sekecil dan sebesar apa pun.

}
{} {}
Soal di atas menunjukkan bahwa secara umum kita tidak dapat memperoleh pengertian jarak yang bermakna pada suatu manifold melalui peta. Metrik alami pada bola satuan diperoleh dari metrik terinduksi pada $\R^3$ atau dari jarak geodesik; lihat
Aufgabe 7.16.

\inputaufgabe
{}
{

Tunjukkan bahwa suatu
\definitionsverweis {interval setengah terbuka}{}{}
\mathl{[a,b[}{} bukan
\definitionsverweis {manifold topologis}{}{}.

}
{} {}

\inputaufgabe
{}
{

Misalkan
\mathl{I=[0,1[}{} adalah
\zusatzklammer {setengah terbuka ke atas} {} {}
interval satuan dan $S^1$ adalah
\definitionsverweis {lingkaran satuan}{}{.}
Tunjukkan bahwa terdapat suatu
\definitionsverweis {pemetaan bijektif}{}{}
\definitionsverweis {kontinu}{}{}
\maabbdisp {f} { I } { S^1
} {}
tetapi
\mathkor {} {I} {dan} {S^1} {}
tidak
\definitionsverweis {homeomorfik}{}{}
.

}
{} {}

\inputaufgabegibtloesung
{}
{

Tunjukkan bahwa ketiga manifold berdimensi satu
\mathlistdisp {S^1 = \left\{ (x,y) \in \R^2  \mid   \betrag { (x,y) } = 1         \right\}} {} {\R} {dan} {\R \setminus \{0\}} {}
tidak homeomorfik satu sama lain secara berpasangan.

}
{} {}

\inputaufgabe
{}
{

Misalkan $M$ adalah
\definitionsverweis {grafik}{}{}
dari fungsi

\mavergleichskettedisp
{\vergleichskette
{ f(x)
}
{ =} {\begin{cases} \sin \frac{1}{x} \text{ untuk } x \neq 0\, , \\ 0 \text{ selain itu }\, , \end{cases}
}
{ } {
}
{ } {
}
{ } {
}
}
{}{}{}
yang diberikan oleh pemetaan
\maabb {f} {\R} {\R
} {.}
Tunjukkan bahwa $M$ bukan
\definitionsverweis {manifold topologis}{}{}.

}
{} {}

\inputaufgabe
{}
{

Tunjukkan bahwa sebuah
\definitionsverweis {bola terbuka}{}{}

\mavergleichskette
{\vergleichskette
{ U \left(   P ,r  \right)
}
{ \subseteq }{ \R^m
}
{  }{
}
{    }{
}
{   }{
}
}
{}{}{}
adalah $C^\infty$-\definitionsverweis {difeomorfik}{}{}
dengan $\R^m$.

}
{} {}

\inputaufgabegibtloesung
{}
{

Misalkan

\mathdisp {S= \left\{ P \in \R^3  \mid   \Vert { P } \Vert = 1         \right\}} {  }

adalah bola satuan. Untuk
\mathl{v=(a,b,c) \neq 0}{} berlaku

\mathdisp {E_v = \left\{ (x,y,z) \in \R^3  \mid  ax+by+cz = 0         \right\}} {  }

sebuah bidang melalui titik nol, yang mendefinisikan sebuah lingkaran besar
\zusatzklammer {sebuah \anfuehrung{khatulistiwa}{}} {} {}
dan dua hemisfer terbuka pada $S$.

a) Jelaskan, untuk
\mathl{v=(a,b,c) \neq 0}{}, lingkaran besar terkait dan kedua hemisfer dengan persamaan atau pertidaksamaan.

b) Tunjukkan bahwa $S$ tidak dapat ditutupi oleh tiga hemisfer terbuka.

}
{} {}

\inputaufgabe
{}
{

Tunjukkan bahwa permukaan silinder terbuka
\mathl{S^1 \times {]0,1[}}{} homeomorfik dengan
\mathl{S^1 \times \R}{,} dengan bidang yang ditusuk
\mathl{\R^2 \setminus \{(0,0)\}}{} dan dengan
\mathl{S^2 \setminus \{N,S\}}{}
\definitionsverweis {homeomorfik}{}{.}

}
{} {}

\inputaufgabe
{}
{

Misalkan
\mathl{]a,b[}{} sebuah
\definitionsverweis {interval terbuka}{}{}
dan
\maabbdisp {f} {]a,b[} {\R_{\geq 0}
} {}
sebuah
\definitionsverweis {fungsi kontinu}{}{.}
Misalkan $M$ adalah permukaan luar benda hasil rotasi terkait. Tunjukkan bahwa ketika

\mavergleichskette
{\vergleichskette
{ f
}
{ > }{ 0
}
{  }{
}
{    }{
}
{   }{
}
}
{}{}{}
ia merupakan manifold yang homeomorfik dengan selubung silinder terbuka. Tunjukkan lebih lanjut bahwa tidak ada manifold jika $f$ memiliki baik titik nol maupun nilai positif.

}
{} {}

\inputaufgabe
{}
{

Tunjukkan bahwa sebuah
\definitionsverweis {manifold topologis}{}{}
\definitionsverweis {terhubung}{}{}
tepat jika dan hanya jika
\definitionsverweis {terhubung lintasan}{}{}
.

}
{} {}

\inputaufgabe
{}
{

Misalkan $M$ sebuah
\definitionsverweis {manifold topologis}{}{}
dan misalkan

\mathl{M= \bigcup_{i \in I} U_i}{} sebuah
\definitionsverweis {tutupan terbuka}{}{}
dengan
\definitionsverweis {peta}{}{}
\maabbdisp {\alpha_i} {U_i } { V_i
} {}
dengan
\mathl{V_i \subseteq \R^n}{.} Untuk
\mathl{i,j \in I}{} definisikan
\maabbdisp {\varphi_{ij} = \alpha_j \circ (\alpha_i)^{-1}} {V_i \cap \alpha_i(U_i \cap U_j)  } { V_j \cap \alpha_j(U_i \cap U_j)
} {}
yang disebut
\definitionsverweis {pemetaan transisi}{}{.}
Tunjukkan bahwa untuk
\mathl{i,j,k \in I}{} apa yang disebut \stichwort {syarat koklus} {}
\zusatzklammer {pada himpunan bagian mana} {?} {}

\mavergleichskettedisp
{\vergleichskette
{ \varphi_{ij}
}
{ =} {\varphi_{kj} \circ \varphi_{ik}
}
{ } {
}
{ } {
}
{ } {
}
}
{}{}{}
berlaku.

}
{} {}

\inputaufgabe
{}
{

Periksa syarat koklus
\zusatzklammer {lihat
Aufgabe 7.11} {} {}
untuk bola satuan $M=S^2$ dan tiga proyeksi stereografis dari Kutub Utara, Kutub Selatan, dan
\mathl{P=(1,0,0)}{.}

}
{} {}

\inputaufgabegibtloesung
{}
{

Buat sebuah animasi yang menampilkan objek-objek geometri dari
Aufgabe 7.17.

}
{} {}

\inputaufgabe
{}
{

Perhatikan kertas kado. Dengan cara apa kertas tersebut dapat dipotong dan/atau direkatkan sehingga menghasilkan manifold berdimensi dua? Apakah tepi kertas harus menjadi bagian darinya atau tidak? Manifold mana yang dihasilkan yang terhubung, mana yang kompak? Bagaimana pita Möbius dibuat? Kemungkinan apa yang ada jika kita mengizinkan sejumlah hingga titik pengecualian, tempat tidak ada struktur manifold?

Terapkan teori ini untuk merancang kemasan yang seorisinal mungkin. Ikatlah kemasan tersebut dengan manifold berdimensi satu yang sesuai.

}
{} {}

  \zwischenueberschrift{Soal untuk dikumpulkan}

\inputaufgabe
{3}
{

Tentukan
\definitionsverweis {citra}{}{} dari
\definitionsverweis {lingkaran-lingkaran besar}{}{}
melalui kedua kutub pada
\definitionsverweis {bola satuan}{}{} di bawah
\definitionsverweis {proyeksi stereografis}{}{}
dari Kutub Utara.

}
{} {}

\inputaufgabe
{5}
{

Tunjukkan bahwa pada
\definitionsverweis {bola satuan}{}{}

\mavergleichskette
{\vergleichskette
{ K
}
{ \subset }{ \R^3
}
{  }{
}
{    }{
}
{   }{
}
}
{}{}{}
melalui korespondensi berikut ditentukan sebuah
\definitionsverweis {metrik}{}{} . Untuk

\mavergleichskette
{\vergleichskette
{ P,Q
}
{ \in }{ K
}
{  }{
}
{    }{
}
{   }{
}
}
{}{}{}
\mathl{d(P,Q)}{} adalah panjang lintasan penghubung
\zusatzklammer {yang lebih pendek} {} {} dari $P$ ke $Q$ pada
\definitionsverweis {lingkaran besar}{}{}
\zusatzklammer {perhatikan juga kasus

\mavergleichskette
{\vergleichskette
{ P
}
{ = }{ Q
}
{  }{
}
{    }{
}
{   }{
}
}
{}{}{}
dan $P,Q$
\definitionsverweis {antipodal}{}{}} {} {.}

}
{} {}

\inputaufgabe
{8}
{

Kita tetapkan dua titik $N=(0,0,1)$ dan $P=(1,0,0)$ pada
\definitionsverweis {bola satuan}{}{} $K$. Misalkan $G$ adalah garis penghubung dan $H$ adalah bidang tegak lurus terhadap $G$ melalui $N$. Pada $H$, buat lingkaran satuan berparameter $E$ dengan $N$ sebagai pusat. Tentukan, untuk $S \in E$, panjang lintasan
\zusatzklammer {yang lebih pendek} {} {} dari $N$ ke $P$ pada lingkaran yang ditentukan oleh irisan $K$ dengan bidang yang melalui
\mathkor {} {N,P} {dan} {S} {},
tersebut.

}
{} {}

\inputaufgabe
{6}
{

Berikan suatu
\definitionsverweis {pemetaan injektif}{}{}
\definitionsverweis {kontinu}{}{}
\maabbdisp {\varphi} {\R} {S^2
} {,}
yang
\zusatzklammer {sebagai pemetaan ke $\R^3$} {} {}
\definitionsverweis {dapat direktifikasi}{}{}
memiliki
\definitionsverweis {panjang}{}{}
tak hingga, dan memenuhi
\mathkor {} {\operatorname{lim}_{   t  \rightarrow  \infty  } \, \varphi (t) = N} {dan} {\operatorname{lim}_{   t  \rightarrow  - \infty  } \, \varphi (t) = S} {}.

}
{} {}

\inputaufgabe
{4}
{

Tunjukkan bahwa sumbu-sumbu koordinat bukan
\definitionsverweis {manifold topologis}{}{}.

}
{} {}


\section*{Solusi yang disediakan oleh sumber}
\addcontentsline{toc}{section}{Solusi yang disediakan oleh sumber}
\subsection*{Solusi Soal 7.4}
Sebagai himpunan bagian tertutup dan terbatas dari $\R^2$, lingkaran satuan kompak
\zusatzklammer {Satz von Heine-Borel} {} {.} Garis riil $\R$ dan $\R \setminus \{0\}$, sebaliknya, tidak kompak karena keduanya merupakan himpunan bagian tak terbatas dari $\R$. Jadi
\mathl{S^1 \not \cong \R,\, \R \setminus \{0\}}{.}

Garis riil terhubung, sebagaimana mengikuti dari
Zwischenwertsatz.
Sebaliknya, $\R \setminus \{0\}$ tidak terhubung, karena kita dapat menuliskannya sebagai

\mavergleichskettedisp
{\vergleichskette
{\R_+
}
{ =} { (\R \setminus \{0\}) \cap [0, \infty]
}
{ =} { (\R \setminus \{0\}) \cap \, ]0, \infty]
}
{ } {
}
{ } {
}
}
{}{}{}
sehingga diperoleh suatu himpunan bagian tak kosong yang sekaligus terbuka dan tertutup. Jadi
\mathl{\R \,\not \cong \R \setminus \{0\}}{.}

\subsection*{Solusi Soal 7.7}
a) Lingkaran besar adalah

\mathdisp {G= \left\{ (x,y,z) \in S  \mid  ax+by+cz = 0        \right\}} {  }

dan kedua hemisfer terbuka adalah

\mathdisp {U_+ = \left\{ (x,y,z) \in S  \mid  ax+by+cz > 0        \right\}} {  }

dan

\mathdisp {U_- = \left\{ (x,y,z) \in S  \mid  ax+by+cz < 0        \right\}} { . }

b) Untuk setiap hemisfer terbuka $U$ dan lingkaran besar $G$ yang terkait berlaku
\mathl{U \cap G = \emptyset}{.} Jarak maksimum antara dua titik
\mathl{P,Q \in S}{} adalah
\mathl{2}{,} dan hal ini terjadi tepat ketika kedua titik tersebut saling berhadapan
\zusatzklammer {antipodal} {} {}, yaitu ketika garis penghubungnya melalui pusat bola
\zusatzklammer {yakni pada
\mathl{Q=-P}{}} {} {}. Pasangan antipodal seperti itu tidak terletak pada satu hemisfer terbuka, sebab untuk
\mathl{P=(x,y,z)}{} dan
\mathl{P \in U_+}{} berlaku
\mathl{ax+by+cz>0}{} dan karena itu
\mathl{a(-x) +b(-y)+c(-z) <0}{} berlaku, sehingga
\mathl{-P \in U_-}{.}

Sekarang anggaplah bahwa terdapat suatu tutupan terbuka

\mathdisp {S \subseteq U_1 \cup U_2 \cup U_3} {  }

dengan tiga hemisfer terbuka
\mathl{U_1,U_2,U_3}{}
\zusatzklammer {bidang dan lingkaran besar yang bersesuaian masing-masing dilambangkan dengan
\mathkor {} {E_i} {atau} {G_i} {}} {} {.}
Karena
\mathl{U_1 \cap G_1 = \emptyset}{} diperoleh

\mathdisp {G_1 \subseteq U_2 \cup U_3} { . }

Irisan
\mathl{G_1 \cap E_2}{} memuat setidaknya dua titik antipodal
\mathkor {} {P} {dan} {-P} {}
dan
\mathl{P,-P \not\in U_2}{.} Karena, berdasarkan pengamatan sebelumnya, $U_3$ tidak memuat pasangan titik antipodal, salah satu dari kedua titik ini juga tidak termasuk dalam $U_3$, sehingga kita memperoleh kontradiksi.

\subsection*{Solusi Soal 7.13}
Animasi pertama memperlihatkan secara bertahap konstruksi objek-objek individual. Animasi lainnya memperlihatkan variasi S dan perubahan panjang lintasan yang bersesuaian.

Dari Lukas Freudenberg.
\href{https://commons.wikimedia.org/wiki/File:Aufgabe75.22.1.gif}{Konstruktion der Objekte}
\href{https://commons.wikimedia.org/wiki/File:Aufgabe75.22.2.gif}{Variation von S}


\part{Unit 8}
\chapter[Kuliah 8]{Kuliah 8: Teorema Pemetaan Implisit dan Manifold}
\input{generated/lecture08.id.build.tex}

\chapter{Lembar Kerja 8}
\setcounter{section}{8}

  \zwischenueberschrift{Soal latihan}

\inputaufgabe
{}
{

Misalkan
\mathkor {} {M} {dan} {N} {}
adalah
\definitionsverweis {manifold diferensiabel}{}{}
dan
\maabbdisp {\varphi} { M } { N
} {}
sebuah pemetaan. Misalkan
\mathl{M= \bigcup_{i \in I} U_i}{} suatu
\definitionsverweis {tutupan terbuka}{}{}
dari $M$. Tunjukkan bahwa $\varphi$
\definitionsverweis {diferensiabel}{}{}
jika dan hanya jika semua pembatasan
\mathl{\varphi_i= \varphi \vert_{U_i }}{} diferensiabel.

}
{} {}

\inputaufgabe
{}
{

Misalkan $M$ sebuah
\definitionsverweis {manifold diferensiabel}{}{}
dan
\maabbdisp {\alpha} { U } { V
} {}
sebuah peta
\zusatzklammer {jadi
\mathl{U \subseteq M}{} dan
\mathl{V \subseteq \R^n}{} terbuka} {} {.}
Tunjukkan bahwa $\alpha$ adalah sebuah
\definitionsverweis {difeomorfisme}{}{}.

}
{} {}

\inputaufgabe
{}
{

Misalkan $M$ sebuah
\definitionsverweis {manifold diferensiabel}{}{} dan
\maabbdisp {f,g} { M } {\R
} {}
\definitionsverweis {fungsi-fungsi diferensiabel}{}{} pada $M$. Buktikan pernyataan-pernyataan berikut. \aufzaehlungvier{Pemetaan
\maabbeledisp {f \times g} { M } {\R^2
} { x } {(f(x),g(x))
} {,} diferensiabel.
}{ $f+g$
\definitionsverweis {diferensiabel}{}{.}
}{ $f \cdot g$
\definitionsverweis {diferensiabel}{}{.}
}{Jika $f$ tidak memiliki titik nol, maka
\mathl{f^{-1}}{} juga diferensiabel.
}

}
{} {}

\inputaufgabe
{}
{

Misalkan
\mathl{S^2 \subseteq \R^3}{} adalah
\definitionsverweis {sfera}{}{.}
Dengan menggunakan
\definitionsverweis {peta-peta stereografis}{}{,}
tunjukkan bahwa pembatasan koordinat ruang
\mathl{x,y,z}{} pada sfera merupakan
\definitionsverweis {fungsi-fungsi diferensiabel}{}{}.

}
{} {}

\inputaufgabe
{}
{

Misalkan
\mathkor {} {M} {dan} {N} {}
adalah
\definitionsverweis {manifold diferensiabel}{}{}
dan
\maabbdisp {\varphi} { M } { N
} {}
sebuah
\definitionsverweis {pemetaan diferensiabel}{}{.}
Tunjukkan bahwa pemetaan ini menginduksi sebuah
\definitionsverweis {homomorfisme gelanggang}{}{}
\maabbeledisp {\varphi^*} {C^1(N,\R)} {C^1(M,\R)
} { f } {f \circ \varphi
} {,}
melalui penarikan balik.

}
{} {}

\inputaufgabe
{}
{

Tunjukkan bahwa untuk
\mathl{m \leq n}{}, penyematan subruang
\mathl{\R^m}{} ke dalam $\R^n$ yang diberikan oleh
\mathl{(x_1 , \ldots , x_m) \mapsto (x_1 , \ldots , x_m,0 , \ldots , 0)}{}
\definitionsverweis {diferensiabel hingga orde berapa pun}{}{}.

}
{} {}

\inputaufgabe
{}
{

Tunjukkan bahwa permukaan silinder terbuka
\mathl{S^1 \times {]0,1[}}{},
\mathl{S^1 \times \R}{,} bidang yang ditusuk
\mathl{\R^2 \setminus \{(0,0)\}}{}, dan
\mathl{S^2 \setminus \{N,S\}}{}
saling \definitionsverweis {difeomorfik}{}{}.

}
{} {}

Soal berikut menggunakan definisi di bawah ini.

Sebuah fungsi
\maabbdisp {f} { {\mathbb K}^n } { {\mathbb K}
} {}
disebut
\definitionswort {homogen berderajat}{}
$d$, jika untuk setiap titik
\mathl{P \in {\mathbb K}^n}{} dan setiap
\mathl{\lambda \in {\mathbb K}}{} berlaku hubungan

\mavergleichskettedisp
{\vergleichskette
{f(\lambda P)
}
{ =} { \lambda^d f(P)
}
{ } {
}
{ } {
}
{ } {
}
}
{}{}{}
untuk semua pilihan tersebut.

\inputaufgabe
{}
{

Misalkan
\maabbdisp {f} {\R^n } {\R
} {}
sebuah
\definitionsverweis {fungsi homogen}{}{} yang
\definitionsverweis {terdiferensialkan secara kontinu}{}{,}
dan yang pada
\definitionsverweis {serat}{}{}
$F$ di atas
\mathl{a \neq 0}{}
bersifat \definitionsverweis {reguler}{}{.}
Tunjukkan bahwa setiap serat di atas
\mathl{b \neq 0,\ b=\lambda^d a}{}
untuk suatu skalar real tak nol,
\definitionsverweis {difeomorfik}{}{} dengan $F$ dan merupakan sebuah
\definitionsverweis {manifold}{}{}.

Catatan edisi: Syarat penskalaan real ditambahkan. Klaim sumber untuk semua nilai tak nol gagal pada derajat genap ketika tanda nilai serat berbeda; misalnya, fungsi kuadrat norma mempunyai serat positif berbentuk bola tetapi serat negatif kosong.

}
{} {}

\inputaufgabe
{}
{

Misalkan
\mathl{]a,b[}{} sebuah
\definitionsverweis {interval terbuka}{}{}
dan
\maabbdisp {f} {]a,b[} {\R_{+}
} {}
sebuah
\definitionsverweis {fungsi yang terdiferensialkan secara kontinu}{}{.}
Misalkan $M$ adalah permukaan benda putar terkait. Tunjukkan bahwa himpunan ini merupakan manifold yang difeomorfik dengan silinder terbuka.

}
{} {}

\inputaufgabe
{}
{

Tunjukkan bahwa permukaan elipsoid dan permukaan bola satuan
$C^\infty$-\definitionsverweis {difeomorfik}{}{}.

}
{} {}

\inputaufgabegibtloesung
{}
{

Berikan sebuah contoh
\definitionsverweis {manifold diferensiabel}{}{}
\definitionsverweis {berdimensi dua}{}{} yang
\definitionsverweis {terhubung}{}{},
sebutlah $M$, dan sebuah titik

\mavergleichskette
{\vergleichskette
{ P
}
{ \in }{ M
}
{  }{
}
{    }{
}
{   }{
}
}
{}{}{}
sedemikian sehingga
\mathkor {} {M} {dan} {M \setminus \{P\}} {}
\definitionsverweis {difeomorfik}{}{}
satu sama lain.

}
{} {}

\inputaufgabe
{}
{

Tunjukkan bahwa
\definitionsverweis {pemetaan antipodal}{}{}
\maabbeledisp {} {S^n } { S^n
} {(x_1 , \ldots , x_{n+1})} {(-x_1 , \ldots , -x_{n+1})
} {,}
adalah sebuah
\definitionsverweis {difeomorfisme}{}{} \definitionsverweis {tanpa titik tetap}{}{}
yang inversnya adalah dirinya sendiri.

}
{} {}

\inputaufgabegibtloesung
{}
{

Tunjukkan bahwa sebuah
\definitionsverweis {manifold diferensiabel}{}{}
$M$ berdimensi

\mavergleichskette
{\vergleichskette
{ n
}
{ \geq }{ 1
}
{  }{
}
{    }{
}
{   }{
}
}
{}{}{}
memiliki tak hingga banyak
\definitionsverweis {difeomorfisme}{}{}
\maabbdisp {\varphi} { M } { M
} {}
yang berbeda.

}
{} {}

Soal-soal berikut dimaksudkan untuk menjelaskan mengapa manifold dirumuskan dengan tutupan terbuka.

Misalkan
\mathl{\left(  x_n  \right)_{n \in  \N  }}{} sebuah
\definitionsverweis {barisan}{}{}
dalam suatu
\definitionsverweis {ruang topologis}{}{}
$X$. Kita mengatakan bahwa, untuk suatu

\mavergleichskette
{\vergleichskette
{ x
}
{ \in }{ X
}
{  }{
}
{    }{
}
{   }{
}
}
{}{}{}
barisan tersebut \definitionswort {konvergen ke titik itu}{,} jika sifat berikut terpenuhi.

Untuk setiap
\definitionsverweis {lingkungan terbuka}{}{}

\mavergleichskette
{\vergleichskette
{ U
}
{ \subseteq }{ X
}
{  }{
}
{    }{
}
{   }{
}
}
{}{}{}
yang memuat $x$, terdapat suatu

\mavergleichskette
{\vergleichskette
{ n_0
}
{ \in }{ \N
}
{  }{
}
{    }{
}
{   }{
}
}
{}{}{}
sedemikian sehingga untuk semua

\mavergleichskette
{\vergleichskette
{ n
}
{ \geq }{ n_0
}
{  }{
}
{    }{
}
{   }{
}
}
{}{}{}
suku barisan $x_n$ termasuk dalam $U$.

Dalam hal ini $x$ disebut \definitionswort {nilai batas}{} atau \definitionswort {limit}{} barisan tersebut. Kita juga menuliskannya sebagai

\mavergleichskettedisp
{\vergleichskette
{  \lim_{n \rightarrow \infty}  x_n
}
{ =} { x
}
{ } {
}
{ } {
}
{ } {
}
}
{}{}{.}
Jika barisan memiliki limit, kita juga mengatakan bahwa barisan itu \definitionswort {konvergen}{}
\zusatzklammer {tanpa menyebut nilai limit tertentu} {} {,}
dan jika tidak, bahwa barisan itu \definitionswort {divergen}{.}

\inputaufgabe
{}
{

Misalkan $M$ sebuah
\definitionsverweis {manifold topologis}{}{}
dan
\mathl{\left(   x_n  \right)_{n \in \N }}{} sebuah
\definitionsverweis {barisan}{}{}
dalam $M$. Tunjukkan bahwa barisan tersebut
\definitionsverweis {konvergen}{}{,}
jika dan hanya jika terdapat sebuah
\definitionsverweis {domain peta}{}{}
dari $M$ yang memuat hampir semua suku barisan dan sedemikian sehingga barisan citra yang bersesuaian konvergen di dalam citra peta.

}
{} {}

\inputaufgabe
{}
{

Misalkan $M$ sebuah
\definitionsverweis {manifold topologis}{}{}
dan
\maabbdisp {\alpha_i} {U_i } { V_i
} {}
sebuah keluarga
\definitionsverweis {peta}{}{}
dengan
\definitionsverweis {pemetaan transisi}{}{}
\maabbdisp {\varphi_{ij} = \alpha_j \circ (\alpha_i)^{-1}} {V_i \cap \alpha_i(U_i \cap U_j) } { V_j \cap \alpha_j(U_i \cap U_j)
} {.}
Tunjukkan bahwa manifold $M$ dapat direkonstruksi dari keluarga

\mathbed {V_i} {}
 {i \in I} {}
 {} {} {} {,}
himpunan-himpunan bagian
\mathl{V_{ij} \subseteq V_i}{} dan pemetaan transisi
\maabbdisp {\varphi_{ij}} {V_{ij}} { V_{ji}
} {}
yang bersesuaian.

a) Pada

\mavergleichskettedisp
{\vergleichskette
{N
}
{ \defeq} { \biguplus_{i \in I} V_i
}
{ } {
}
{ } {
}
{ } {
}
}
{}{}{}
perhatikan
\definitionsverweis {relasi ekuivalensi}{}{,}
yang menyamakan dua titik
\mathl{P\in V_i}{} dan
\mathl{Q \in V_j}{} jika
\mathl{\varphi_{ij}}{} memetakan titik pertama ke titik kedua.

b) Lengkapi himpunan hasil bagi
\mathl{N/ \sim}{} dengan sebuah topologi yang sesuai.

c) Definisikan peta-peta pada
\mathl{N/ \sim}{}.

d) Tunjukkan bahwa
\mathkor {} {M} {dan} {N/ \sim} {}
\definitionsverweis {homeomorfik}{}{}.

}
{} {}

Prosedur konstruksi manifold yang dijelaskan dalam soal sebelumnya berlaku untuk suatu keluarga himpunan bagian terbuka di dalam $\R^n$ dengan pemetaan transisi yang memenuhi syarat koklus dari
Soal 7.11.
Namun, ruang topologis yang dihasilkan tidak serta-merta merupakan ruang Hausdorff. Konstruksi ini disebut \stichwort {perekatan terbuka} {} ruang-ruang.

\inputaufgabe
{}
{

Kita mengambil dua salinan garis riil, yaitu

\mavergleichskette
{\vergleichskette
{G_1
}
{ = }{\R
}
{  }{
}
{    }{
}
{   }{
}
}
{}{}{}
dan

\mavergleichskette
{\vergleichskette
{G_2
}
{ = }{\R
}
{  }{
}
{    }{
}
{   }{
}
}
{}{}{}
beserta pemetaan perekatan
\maabbeledisp {\varphi} {G_1 \setminus \{0\} } { G_2 \setminus \{0\}
} { x } { x^{-1}
} {.}
Misalkan $M$ ruang topologis yang dihasilkan menurut konstruksi yang dijelaskan dalam
Soal 8.15.
Tunjukkan bahwa $M$
\definitionsverweis {homeomorfik}{}{}
dengan
$1$-\definitionsverweis {sfera}{}{}.

}
{} {}

\inputaufgabe
{}
{

Kita mengambil dua salinan garis riil, yaitu

\mavergleichskette
{\vergleichskette
{G_1
}
{ = }{\R
}
{  }{
}
{    }{
}
{   }{
}
}
{}{}{}
dan

\mavergleichskette
{\vergleichskette
{G_2
}
{ = }{\R
}
{  }{
}
{    }{
}
{   }{
}
}
{}{}{}
beserta pemetaan perekatan
\maabbdisp {\operatorname{Id}} {G_1 \setminus \{0\} } { G_2 \setminus \{0\}
} {.}
Misalkan $M$ ruang topologis yang dihasilkan menurut konstruksi yang dijelaskan dalam
Soal 8.15.
Tunjukkan bahwa $M$ bukan
\definitionsverweis {manifold}{}{}.

}
{} {}

  \zwischenueberschrift{Soal untuk dikumpulkan}

Dalam soal berikut, tafsirkan ${\mathbb C}$ sebagai $\R^2$.

\inputaufgabe
{4}
{

Perhatikan pemetaan
\maabbeledisp {} { {\mathbb C}^2} { {\mathbb C}
} {(z,w)} { zw
} {.}
Untuk titik manakah
\mathl{u \in {\mathbb C}}{} serat di atas $u$ merupakan manifold? Dalam setiap kasus, berikan deskripsi kelas difeomorfisme yang sesederhana mungkin.

}
{} {}

\inputaufgabe
{6}
{

Misalkan diberikan dua titik
\mathkor {} {P} {dan} {Q} {}
pada
\definitionsverweis {permukaan bola satuan}{}{.}
Tunjukkan bahwa terdapat sebuah
\definitionsverweis {difeomorfisme}{}{}
dari permukaan bola itu ke dirinya sendiri yang memetakan $P$ ke $Q$.

}
{} {}

\inputaufgabe
{4 (1+1+2)}
{

Misalkan $M$ sebuah
\definitionsverweis {manifold diferensiabel}{}{.}
Untuk setiap himpunan bagian terbuka

\mavergleichskette
{\vergleichskette
{ U
}
{ \subseteq }{ M
}
{  }{
}
{    }{
}
{   }{
}
}
{}{}{}
kita perhatikan himpunan
\mathl{C^1(U,\R)}{} yang terdiri atas
\definitionsverweis {fungsi-fungsi diferensiabel}{}{}
pada $U$. Misalkan

\mavergleichskette
{\vergleichskette
{ M
}
{ = }{ \bigcup_{i \in I} U_i
}
{  }{
}
{    }{
}
{   }{
}
}
{}{}{}
sebuah
\definitionsverweis {tutupan terbuka}{}{.}
\aufzaehlungdreiabc{Tunjukkan bahwa jika

\mavergleichskette
{\vergleichskette
{ V
}
{ \subseteq }{ U
}
{  }{
}
{    }{
}
{   }{
}
}
{}{}{}
terbuka dan

\mavergleichskette
{\vergleichskette
{ f
}
{ \in }{ C^1(U,\R)
}
{  }{
}
{    }{
}
{   }{
}
}
{}{}{}
maka
\definitionsverweis {pembatasan}{}{}
$f \vert_V$ juga termasuk dalam
\mathl{C^1(V,\R)}{}.
}{Misalkan

\mavergleichskette
{\vergleichskette
{ f
}
{ \in }{ C^1(M,\R)
}
{  }{
}
{    }{
}
{   }{
}
}
{}{}{.}
Tunjukkan bahwa

\mavergleichskette
{\vergleichskette
{ f
}
{ = }{ 0
}
{  }{
}
{    }{
}
{   }{
}
}
{}{}{}
jika dan hanya jika semua pembatasan

\mavergleichskette
{\vergleichskette
{ f \vert_{U_i }
}
{ = }{ 0
}
{  }{
}
{    }{
}
{   }{
}
}
{}{}{}
berlaku.
}{Misalkan diberikan suatu keluarga

\mavergleichskette
{\vergleichskette
{ f_i
}
{ \in }{ C^1(U_i,\R)
}
{  }{
}
{    }{
}
{   }{
}
}
{}{}{}
fungsi-fungsi yang memenuhi \anfuehrung{syarat kecocokan}{}

\mavergleichskettedisp
{\vergleichskette
{ f_i \vert_{U_i \cap U_j }
}
{ =} { f_j \vert_{U_i \cap U_j }
}
{ } {
}
{ } {
}
{ } {
}
}
{}{}{}
untuk semua $i,j$. Tunjukkan bahwa terdapat suatu

\mavergleichskette
{\vergleichskette
{ f
}
{ \in }{ C^1(M,\R)
}
{  }{
}
{    }{
}
{   }{
}
}
{}{}{}
yang memenuhi

\mavergleichskette
{\vergleichskette
{ f \vert_{U_i }
}
{ = }{ f_i
}
{  }{
}
{    }{
}
{   }{
}
}
{}{}{}
untuk semua $i$.
}

}
{} {}

\inputaufgabe
{5}
{

Misalkan $M$ sebuah
\definitionsverweis {manifold diferensiabel}{}{,} yang memiliki setidaknya dua titik. Tunjukkan bahwa terdapat
\definitionsverweis {fungsi-fungsi diferensiabel}{}{}
\maabbdisp {f,g} { M } {\R
} {}
dengan
\mathl{f,g \neq 0}{,} tetapi
\mathl{fg=0}{.}

}
{} {}


\section*{Solusi yang disediakan oleh sumber}
\addcontentsline{toc}{section}{Solusi yang disediakan oleh sumber}
\subsection*{Solusi Soal 8.11}
Misalkan

\mavergleichskettedisp
{\vergleichskette
{M
}
{ =} { \R^2 \setminus \N_+
}
{ =} { \R^2 \setminus \{ (1,0) , (2,0), (3,0), \ldots  \}
}
{ } {
}
{ } {
}
}
{}{}{,}
yakni dari $\R^2$ kita mengeluarkan bilangan-bilangan asli positif yang terletak pada sumbu $x$. Himpunan ini terbuka di dalam $\R^2$ dan karena itu merupakan manifold diferensiabel. Manifold ini terhubung lintasan: misalnya, setiap titik di $M$ dengan
\mathl{y \geq 0}{} dapat dihubungkan oleh lintasan lurus dengan
\mathl{(0,1)}{}, dan setiap titik di $M$ dengan
\mathl{y \leq 0}{} dapat dihubungkan oleh lintasan lurus dengan
\mathl{(0,-1)}{}; kedua titik tersebut juga dapat dihubungkan oleh lintasan lurus. Sekarang misalkan

\mavergleichskettedisp
{\vergleichskette
{P
}
{ =} { (0,0)
}
{ } {
}
{ } {
}
{ } {
}
}
{}{}{}
dan

\mavergleichskettedisp
{\vergleichskette
{M'
}
{ =} { M \setminus \{P\}
}
{ } {
}
{ } {
}
{ } {
}
}
{}{}{.}
Pemetaan
\maabbeledisp {} {\R^2} {\R^2
} {(x,y)} {(x-1,y)
} {,}
merupakan sebuah difeomorfisme
\zusatzklammer {sebagai translasi} {} {},
dan citra $M$ tepat sama dengan $M'$. Oleh karena itu,
\mathkor {} {M} {dan} {M'=M \setminus \{P\}} {}
difeomorfik satu sama lain.

\subsection*{Solusi Soal 8.13}
\input{generated/worksheet08_exercise13_solution.id.build.tex}

\part{Unit 9}
\chapter[Kuliah 9]{Kuliah 9: Ruang Singgung Manifold Diferensiabel}
\input{generated/lecture09.id.build.tex}

\chapter{Lembar Kerja 9}
\input{generated/worksheet09.id.build.tex}

\part{Unit 10}
\chapter[Kuliah 10]{Kuliah 10: Submanifold Tertutup}
\input{generated/lecture10.id.build.tex}

\chapter{Lembar Kerja 10}
\input{generated/worksheet10.id.build.tex}

\section*{Solusi yang disediakan oleh sumber}
\addcontentsline{toc}{section}{Solusi yang disediakan oleh sumber}
\subsection*{Solusi Soal 10.9}
\input{generated/worksheet10_exercise09_solution.id.build.tex}
\subsection*{Solusi Soal 10.10}
\input{generated/worksheet10_exercise10_solution.id.build.tex}
\subsection*{Solusi Soal 10.15}
Dengan regularitas $C^2$ yang dinyatakan dalam soal edisi ini, kita klaim

\mavergleichskettealignhandlinks
{\vergleichskettealignhandlinks
{ TY
}
{ =} { \left\{  (P,u)   \mid   h(P) = c , \,  \left(  D h   \right)_{ P } \left(  u  \right) = 0       \right\}
}
{ =} { \left\{  (P,u_1 , \ldots , u_n)   \mid   h(P) = c , \,  \sum_{i = 1}^n u_i { \frac{   \partial   h   }{  \partial   x_i   } } (P) = 0       \right\}
}
{ \subseteq} { W \times \R^n
}
{ } {
}
}
{}
{}{}
dengan proyeksi alami ke $Y$.

Definisikan pemetaan
\maabbeledisp {\Phi} {W\times\R^n} {\R^2
} {(P,u)} {\left(h(P),\left(Dh\right)_P(u)\right)
} {.}
Karena $h$ berkelas $C^2$, pemetaan $\Phi$ berkelas $C^1$, dan

\mathdisp {TY=\Phi^{-1}(c,0)} { . }

Kita tunjukkan bahwa $(c,0)$ merupakan nilai reguler. Untuk
\mathl{(P,u)\in TY}{}, diferensial $\Phi$ diberikan oleh

\mathdisp {(D\Phi)_{(P,u)}(v,w)=\left((Dh)_P(v),\,(D^2h)_P(v,u)+(Dh)_P(w)\right)} { . }

Karena $h$ reguler pada $Y$, bentuk linear $(Dh)_P$ surjektif. Untuk sebarang
\mathl{(r,s)\in\R^2}{}, pilih $v$ dengan $(Dh)_P(v)=r$, lalu pilih $w$ dengan

\mathdisp {(Dh)_P(w)=s-(D^2h)_P(v,u)} { . }

Dengan pilihan ini, $(D\Phi)_{(P,u)}(v,w)=(r,s)$, sehingga diferensial $\Phi$ surjektif di setiap titik serat. Teorema pemetaan implisit kini menunjukkan bahwa $TY$ merupakan submanifold tertutup dari $W\times\R^n$ berdimensi $2n-2$. Ketertutupannya juga langsung mengikuti dari kontinuitas $\Phi$ dan ketertutupannya titik $(c,0)$ di $\R^2$.

Untuk mencocokkan realisasi tertanam ini dengan struktur bundel, bagi setiap
\mathl{P\in Y}{}
pilih suatu parametrisasi lokal
\maabbdisp {\beta} {V} {U\subseteq Y
} {,}
dengan $V\subseteq\R^{n-1}$ terbuka. Pemetaan
\maabbeledisp {} {V \times \R^{n-1} } { U \times \R^n
} {(Q,z)} { ( \beta(Q), \left(  D\beta  \right)_{ Q } \left(  z  \right) )
} {,}
memberikan trivialisasi lokal yang citranya tepat terdiri atas pasangan titik dan vektor singgung pada serat. Jadi topologi serta proyeksi realisasi tertanam ini sama dengan topologi dan proyeksi bundel singgung yang telah didefinisikan.

Catatan edisi: Solusi sumber mengganti syarat $h(P)=c$ dengan $h(P)=0$, menulis turunan parsial dari fungsi $f$ yang tidak didefinisikan, dan memakai $\beta(q)$ alih-alih $\beta(Q)$. Lebih mendasar, asumsi $C^1$ sumber hanya membuat turunan pertama kontinu dan tidak cukup untuk menyimpulkan bahwa $TY$ adalah submanifold diferensiabel dari $W\times\R^n$. Edisi ini memperkuat asumsi soal menjadi $C^2$ dan memberikan pemeriksaan nilai reguler yang hilang.

\subsection*{Solusi Soal 10.25}
Misalkan $N=(0,0,1)$ dan gunakan invers proyeksi stereografis
\maabbeledisp {\sigma} {\R^2} {S^2\setminus\{N\}
} {(x,y)} {\left(\frac{2x}{1+x^2+y^2},\frac{2y}{1+x^2+y^2},\frac{x^2+y^2-1}{1+x^2+y^2}\right)
} {.}
Pada $S^2\setminus\{N\}$, dorong medan vektor konstan $e_1$ ke depan dan definisikan

\mathdisp {X(\sigma(x,y))=(D\sigma)_{(x,y)}(e_1)} { . }

Secara eksplisit,

\mathdisp {X(\sigma(x,y))=\frac{1}{(1+x^2+y^2)^2}\left(2(1-x^2+y^2),-4xy,4x\right)} { . }

Karena $\sigma$ merupakan difeomorfisme, diferensialnya injektif; jadi medan ini tidak mempunyai titik nol pada $S^2\setminus\{N\}$. Selain itu,

\mathdisp {\left\|X(\sigma(x,y))\right\|=\frac{2}{1+x^2+y^2}} { , }

yang menuju $0$ ketika $x^2+y^2\to\infty$, yakni ketika $\sigma(x,y)\to N$. Oleh sebab itu, medan tersebut dapat diperluas secara kontinu dengan menetapkan

\mavergleichskettedisp
{\vergleichskette
{X(N)
}
{ =} {0
}
{ \in }{T_NS^2
}
{ } {
}
{ } {
}
}
{}{}{.}
Dalam realisasi $TS^2\subseteq S^2\times\R^3$, rumus norma di atas langsung membuktikan kontinuitas pada $N$. Dengan demikian, $X$ adalah medan vektor kontinu pada $S^2$ dan satu-satunya titik nolnya ialah $N$.

Catatan edisi: Solusi sumber menyatakan medan $e_1/(x^2+y^2)$ pada seluruh $\R^2$, padahal rumus itu tidak terdefinisi di titik asal, lalu menyimpulkan kontinuitas di kutub hanya dari komponen koordinat. Konstruksi di atas memakai medan konstan yang terdefinisi di seluruh bidang dan memeriksa norma dorong-majunya di dalam realisasi ambien, sehingga perluasan di kutub benar-benar kontinu.


\part{Unit 11}
\chapter[Kuliah 11]{Kuliah 11: Produk Manifold}
\setcounter{section}{11}

  \zwischenueberschrift{Produk manifold}

\inputdefinition
{}
{

Misalkan
\mathkor {} {M} {dan} {N} {}
dua
\definitionsverweis {manifold diferensiabel}{}{}
dengan
\definitionsverweis {atlas-atlas}{}{}
\mathkor {} {(U_i,U_i',\alpha_i, i \in I)} {dan} {(V_j,V_j',\beta_j, j \in J)} {.}
Maka
\definitionsverweis {ruang produk}{}{}

\mathl{M \times N}{} dengan
\definitionsverweis {peta-peta}{}{}
\maabbdisp {\alpha_i \times \beta_j} {U_i \times V_j } { U_i' \times V_j'
} {}
\zusatzklammer {dengan

\mavergleichskettek
{\vergleichskettek
{  (i,j)
}
{ \in }{ I \times J
}
{  }{
}
{    }{
}
{   }{
}
}
{}{}{}
dan

\mavergleichskettek
{\vergleichskettek
{ U_i' \times V_j'
}
{ \subseteq }{ \R^m \times \R^n
}
{  }{
}
{    }{
}
{   }{
}
}
{}{}{}} {} {}
disebut \definitionswort {produk manifold}{}
\mathkor {} {M} {dan} {N} {.}

}

Objek ini memang merupakan suatu manifold; lihat
Soal 11.1.
Manifold produk berbentuk
\mathl{M \times \R}{} disebut juga \stichwort {silinder} {} di atas $M$. Jika

\mavergleichskette
{\vergleichskette
{ M
}
{ \subseteq }{ \R^\ell
}
{  }{
}
{    }{
}
{   }{
}
}
{}{}{}
dan

\mavergleichskette
{\vergleichskette
{ N
}
{ \subseteq }{ \R^k
}
{  }{
}
{    }{
}
{   }{
}
}
{}{}{}
merupakan submanifold-submanifold tertutup, maka manifold produk
\mathl{M \times N}{} merupakan submanifold tertutup dari

\mavergleichskette
{\vergleichskette
{ \R^\ell \times \R^k
}
{ = }{ \R^{\ell + k}
}
{  }{
}
{    }{
}
{   }{
}
}
{}{}{,}
lihat
Soal 11.2.



\inputfaktbeweis
{Produkt von Mannigfaltigkeiten/Abbildungseigenschaften/Fakt}
{Lema}
{}
{

\faktsituation {Misalkan
\mathkor {} {M} {dan} {N} {}
\definitionsverweis {manifold-manifold diferensiabel}{}{}
dan
\mathl{M \times N}{} merupakan
\definitionsverweis {produk}{}{} keduanya.}
\faktuebergang {Maka berlaku sifat-sifat berikut.}
\faktfolgerung {\aufzaehlungdrei{Kedua
\definitionsverweis {proyeksi}{}{}
\maabbeledisp {p_1} {M \times N} {M
} {(x,y)} { x
} {,}
dan
\maabbeledisp {p_2} {M \times N} {N
} {(x,y)} { y
} {,}
merupakan
\definitionsverweis {pemetaan diferensiabel}{}{.}
}{
\definitionsverweis {Ruang singgung}{}{}
di suatu titik

\mavergleichskette
{\vergleichskette
{ R
}
{ = }{ (P,Q)
}
{  }{
}
{    }{
}
{   }{
}
}
{}{}{}
secara kanonik teridentifikasi dengan

\mavergleichskette
{\vergleichskette
{ T_R(M \times N)
}
{ = }{ T_PM \times T_QN
}
{  }{
}
{    }{
}
{   }{
}
}
{}{}{.}
}{Misalkan $L$ suatu manifold diferensiabel lain. Maka suatu pemetaan
\maabbeledisp {\varphi \times \psi} { L } { M \times N
} { u } { (\varphi(u), \psi(u))
} {,}
diferensiabel jika dan hanya jika kedua
\definitionsverweis {pemetaan komponennya}{}{}
\mathkor {} {\varphi} {dan} {\psi} {}
diferensiabel.
}}
\faktzusatz {}
\faktzusatz {}

}
{

\teilbeweis {}{}{}
{(1). Dengan beralih ke peta, kita dapat mengasumsikan bahwa
\mathkor {} {M} {dan} {N} {}
merupakan
\definitionsverweis {himpunan-himpunan bagian terbuka}{}{}
masing-masing di
\mathkor {} {\R^m} {dan} {\R^n} {.}
Dalam hal ini, pemetaan tersebut merupakan suatu
\definitionsverweis {pembatasan}{}{}
dari
\definitionsverweis {proyeksi linear}{}{}
\maabb {} {\R^m \times \R^n } {\R^m
} {,}
yang
menurut Proposisi 45.3 (Analysis (Osnabrück 2021-2023))
\definitionsverweis {terdiferensialkan secara kontinu}{}{.}}
{}
\teilbeweis {}{}{}
{(2). Proyeksi-proyeksi diferensiabel
\maabb {p_1} {M \times N} {M
} {}
dan
\maabb {p_2} {M \times N} {N
} {}
menghasilkan
\definitionsverweis {pemetaan-pemetaan singgung}{}{}
linear
\maabb {T_R(p_1)} {T_R(M \times N)} {T_PM
} {}
dan
\maabb {T_R(p_2)} {T_R(M \times N)} {T_QN
} {}
sehingga secara keseluruhan menghasilkan pemetaan linear
\maabbdisp {T_R(p_1) \times T_R(p_2)} { T_R(M \times N)} {T_PM \times T_QN} {.}
Untuk membuktikan bijektivitas, kita dapat beralih ke peta dan mengasumsikan bahwa

\mavergleichskette
{\vergleichskette
{ M
}
{ \subseteq }{ \R^m
}
{  }{
}
{    }{
}
{   }{
}
}
{}{}{}
dan

\mavergleichskette
{\vergleichskette
{ N
}
{ \subseteq }{ \R^n
}
{  }{
}
{    }{
}
{   }{
}
}
{}{}{}
merupakan himpunan-himpunan bagian terbuka. Pemetaan ini kemudian menjadi bijeksi
\maabbdisp {p_1 \times p_2} {\R^{m+n} } { \R^m \times \R^n
} {.}}
{}
\teilbeweis {}{}{}
{(3). Untuk suatu titik tetap

\mavergleichskette
{\vergleichskette
{ A
}
{ \in }{ L
}
{  }{
}
{    }{
}
{   }{
}
}
{}{}{}
dengan menggunakan lingkungan peta dari $A$ dan dari
\mathkor {} {\varphi(A)} {dan} {\psi(A)} {}
kita dapat mereduksi keadaan ke kasus ketika ketiga manifold merupakan himpunan-himpunan terbuka dalam ruang-ruang Euklides. Jika kedua pemetaan terdiferensialkan secara kontinu, maka menurut
Soal 45.9 (Analysis (Osnabrück 2021-2023))
pemetaan keseluruhannya terdiferensialkan secara kontinu. Arah sebaliknya jelas.}
{}

}

\bild{ \begin{center}
\includegraphics[width=5.5cm]{ \bildeinlesungpng {Toroidal_coord} {png} }
\end{center}
\bildtext {} }

\bildlizenz { Toroidal coord.png } {} {DaveBurke} {Commons} {CC-by 2.5} {}

\inputbeispiel{}
{

\definitionsverweis {Produk}{}{}
dari
\definitionsverweis {lingkaran}{}{}
dengan dirinya sendiri, yaitu

\mavergleichskette
{\vergleichskette
{ M
}
{ = }{ S^1 \times S^1
}
{  }{
}
{    }{
}
{   }{
}
}
{}{}{,}
disebut \stichwort {torus} {.} Ini merupakan manifold berdimensi dua. Karena

\mavergleichskette
{\vergleichskette
{ S^1
}
{ \subset }{ \R^2
}
{  }{
}
{    }{
}
{   }{
}
}
{}{}{}
merupakan
\definitionsverweis {submanifold tertutup}{}{,}
torus dapat direalisasikan sebagai submanifold tertutup di

\mavergleichskette
{\vergleichskette
{ \R^2 \times \R^2
}
{ = }{ \R^4
}
{  }{
}
{    }{
}
{   }{
}
}
{}{}{.}
Namun, torus juga dapat direalisasikan sebagai submanifold tertutup di $\R^3$. Untuk itu, misalkan
\mathkor {} {r} {dan} {R} {}
bilangan-bilangan real positif dengan

\mavergleichskette
{\vergleichskette
{ 0
}
{ < }{ r
}
{ < }{ R
}
{    }{
}
{   }{
}
}
{}{}{.}
Maka himpunan

\mathdisp {\left\{ (x,y,z) \in \R^3  \mid   \left(  \sqrt{x^2+y^2} -R  \right)^2 +z^2 = r^2         \right\}} {  }

merupakan suatu torus. Dalam realisasi ini, torus merupakan permukaan
\zusatzklammer {yang digelembungkan} {} {}
\anfuehrung{ban dalam sepeda}{,} dengan \anfuehrung{jari-jari roda}{} sama dengan $R$ dan \anfuehrung{jari-jari tabung}{} sama dengan $r$
\zusatzklammer {roda terletak pada bidang \mathlk{x-y}{-}} {} {.}
Setelah setiap faktor lingkaran diidentifikasi melalui
\mathl{S^1 \cong \R/(2\pi\mathbb{Z})}{,}
hubungan dengan produk
\mathl{S^1 \times S^1}{} diperoleh dengan memetakan pasangan kelas sudut
\mathl{([ \varphi ],[ \psi ])}{} ke titik
\mathl{( (R +r \cos  \psi) \cos \varphi, (R+ r \cos  \psi ) \sin \varphi ,  r \sin  \psi )}{}.

}

  \zwischenueberschrift{Braket Lie medan vektor}

\inputdefinition
{}
{

Misalkan $M$ suatu
\definitionsverweis {manifold diferensiabel}{}{.}
Suatu pemetaan
\maabbdisp {D} {C^1(M,\R) } { C^0(M,\R)
} {}
disebut
\definitionswort {operator diferensial orde pertama}{,}
jika memenuhi sifat-sifat berikut.
\aufzaehlungzwei { $D$ bersifat
$\R$-\definitionsverweis {linear}{}{.}
} {Berlaku

\mavergleichskette
{\vergleichskette
{ D(fg)
}
{ = }{ fD(g)+gD(f)
}
{  }{
}
{    }{
}
{   }{
}
}
{}{}{.}
}

}

Operator diferensial orde pertama disebut juga \stichwort {derivasi} {.}

Untuk suatu himpunan terbuka

\mavergleichskette
{\vergleichskette
{ U
}
{ \subseteq }{ \R^n
}
{  }{
}
{    }{
}
{   }{
}
}
{}{}{,}
turunan-turunan parsial $\partial_i$ dan $\partial_j$, sebagai operasi pada fungsi-fungsi yang dua kali terdiferensialkan secara kontinu, dapat dipertukarkan; jadi

\mavergleichskettedisp
{\vergleichskette
{ \partial_i \circ \partial_j
}
{ =} { \partial_j \circ \partial_i
}
{ } {
}
{ } {
}
{ } {
}
}
{}{}{.}
Suatu kombinasi linear

\mavergleichskettedisp
{\vergleichskette
{ V
}
{ =} { \sum_{i = 1}^n f_i \partial_i
}
{ } {
}
{ } {
}
{ } {
}
}
{}{}{}
dengan koefisien-koefisien kontinu
\mathl{f_i \in C^0(U,\R)}{}
dapat dipandang sebagai
\definitionsverweis {medan vektor}{}{}
\maabbeledisp {} { U } { TU = U \times \R^n
} { P } { (P, \sum_{i = 1}^n f_i(P) \partial_i ) = (P,f_1(P) , \ldots , f_n(P) )
} {,}
dan juga sebagai operator turunan
\maabbeledisp {} {C^1(U,\R)} {C^0(U,\R)
} {h } { \sum_{i = 1}^n f_i  \partial_i  h
} {,}
\zusatzklammer {atau dari $C^\infty(U,\R)$ ke $C^\infty(U,\R)$ apabila
\mathl{f_i \in C^\infty(U,\R)}{} untuk setiap $i$} {} {.}
Dalam pengertian ini, medan vektor dan operator diferensial orde pertama saling bersesuaian. Lebih lanjut, operator-operator diferensial dapat dikomposisikan satu sama lain, dan ternyata urutannya penting.



\inputfaktbeweis
{Offene Menge/R^n/Vektorfelder/Hintereinanderschaltung/Explizit/Fakt}
{Lema}
{}
{

\faktsituation {Misalkan

\mavergleichskette
{\vergleichskette
{ U
}
{ \subseteq }{ \R^n
}
{  }{
}
{    }{
}
{   }{
}
}
{}{}{}
\definitionsverweis {terbuka}{}{}
dan misalkan $f_1 , \ldots , f_n, g_1 , \ldots , g_n,h$ merupakan fungsi-fungsi pada $U$ yang dua kali
\definitionsverweis {terdiferensialkan secara kontinu}{}{.}}
\faktfolgerung {Maka berlaku

\mavergleichskettedisphandlinks
{\vergleichskettedisphandlinks
{ \left(  \sum_{i = 1}^n f_i \partial_i  \right) \left(  \left(  \sum_{j = 1}^n g_j \partial_j   \right) (h)  \right)
}
{ =} { \sum_{j = 1}^n  \left(  \sum_{i = 1}^n  f_i \partial_i g_j  \right) \partial_j h + \sum_{j = 1}^n \sum_{i = 1}^n    f_i g_j \partial_i \partial_j h
}
{ } {
}
{ } {
}
{ } {
}
}
{}{}{.}}
\faktzusatz {}
\faktzusatz {}

}
{

Berlaku

\mavergleichskettealignhandlinks
{\vergleichskettealignhandlinks
{ \left(  \sum_{i = 1}^n f_i \partial_i  \right) \left(  \left(  \sum_{j = 1}^n g_j \partial_j   \right) (h)  \right)
}
{ =} { \left(  \sum_{i = 1}^n f_i \partial_i  \right) \left(  \sum_{j = 1}^n g_j \partial_j (h)  \right)
}
{ =} { \sum_{i = 1}^n f_i \partial_i \left(  \sum_{j = 1}^n g_j \partial_j (h)  \right)
}
{ =} { \sum_{i = 1}^n f_i \left(  \sum_{j = 1}^n \left(  \partial_i (g_j) \partial_j (h) +  g_j \partial_i \partial_j (h)  \right)  \right)
}
{ =} { \sum_{j = 1}^n \left(  \sum_{i = 1}^n f_i  \partial_i g_j  \right) \partial_j h +  \sum_{j = 1}^n \sum_{i = 1}^n    f_i g_j \partial_i \partial_j h
}
}
{}
{}{.}

}

Dari deskripsi eksplisit ini tampak bahwa kedua komposisi pada umumnya dapat berbeda; yakni, dapat terjadi

\mavergleichskettedisp
{\vergleichskette
{ \left(  \sum_{i = 1}^n f_i \partial_i  \right) \circ  \left(  \sum_{j = 1}^n g_j \partial_j   \right)
}
{ \neq} { \left(  \sum_{j = 1}^n g_j \partial_j   \right) \circ  \left(  \sum_{i = 1}^n f_i \partial_i  \right)
}
{ } {
}
{ } {
}
{ } {
}
}
{}{}{,}
karena meskipun suku kedua selalu sama, suku pertamanya pada umumnya dapat berbeda. Hal ini juga mempunyai konsekuensi berikut.


\inputfaktbeweis
{Offene Menge/R^n/Vektorfelder/Lie-Klammer/Vektorfeld/Fakt}
{Lema}
{}
{

\faktsituation {Misalkan

\mavergleichskette
{\vergleichskette
{ U
}
{ \subseteq }{ \R^n
}
{  }{
}
{    }{
}
{   }{
}
}
{}{}{}
\definitionsverweis {terbuka}{}{}
dan misalkan $f_1 , \ldots , f_n, g_1 , \ldots , g_n$ merupakan fungsi-fungsi pada $U$ yang dua kali
\definitionsverweis {terdiferensialkan secara kontinu}{}{}
dengan operator-operator diferensial terkait

\mavergleichskette
{\vergleichskette
{ D
}
{ = }{ \sum_{i = 1}^n f_i \partial_i
}
{  }{
}
{    }{
}
{   }{
}
}
{}{}{}
dan

\mavergleichskette
{\vergleichskette
{ E
}
{ = }{ \sum_{j = 1}^n g_j \partial_j
}
{  }{
}
{    }{
}
{   }{
}
}
{}{}{}}
\faktfolgerung {Maka berlaku

\mavergleichskettedisp
{\vergleichskette
{ D \circ E -E \circ D
}
{ =} { \sum_{j = 1}^n \left(  \sum_{i = 1}^n \left(  f_i  \partial_i g_j - g_i \partial_i f_j  \right)  \right)  \partial_j
}
{ } {
}
{ } {
}
{ } {
}
}
{}{}{.}
Secara khusus, ungkapan ini sendiri bersesuaian dengan suatu operator diferensial orde $1$, dan dengan demikian dengan suatu medan vektor.}
\faktzusatz {}
\faktzusatz {}

}
{

Menurut
Lema 11.5
berlaku

\mavergleichskettealignhandlinks
{\vergleichskettealignhandlinks
{D \circ E-E \circ D
}
{ =} { \left(  \sum_{i = 1}^n f_i \partial_i  \right) \circ  \left(  \sum_{j = 1}^n g_j \partial_j   \right) -  \left(  \sum_{j = 1}^n g_j \partial_j   \right) \circ  \left(  \sum_{i = 1}^n f_i \partial_i  \right)
}
{ =} { \sum_{j = 1}^n  \left(  \sum_{i = 1}^n  \left(  f_i  \partial_i g_j - g_i \partial_i f_j  \right)  \right)  \partial_j
}
{ } {
}
{ } {
}
}
{}
{}{.}

}

Pada suatu manifold, operator diferensial yang bersesuaian dengan medan vektor $V$ kita nyatakan dengan $D_V$. Untuk suatu fungsi diferensiabel
\maabb {f} { M } {\R
} {}
operator $D_V$ didefinisikan secara titik-demi-titik oleh

\mavergleichskette
{\vergleichskette
{ \left( D_V f \right)(P)
}
{ = }{ d f_P\left(V_P\right)
}
{ \in }{ \R
}
{    }{
}
{   }{
}
}
{}{}{,}
dengan $V_P=V(P)$ dan dengan identifikasi kanonik
\mathl{T_{f(P)}\R \cong \R}{.}
Dengan kata lain, untuk suatu titik

\mavergleichskette
{\vergleichskette
{ P
}
{ \in }{ M
}
{  }{
}
{    }{
}
{   }{
}
}
{}{}{}
pemetaan singgung penuh memenuhi

\mavergleichskettedisp
{\vergleichskette
{ \left( T(f) \circ V \right)(P)
}
{ =} { \left( f(P), \left(D_V f\right)(P) \right)
}
{ } {
}
{ } {
}
{ } {
}
}
{}{}{.}
Pernyataan yang bersesuaian berlaku pula untuk fungsi diferensiabel yang didefinisikan pada suatu himpunan bagian terbuka dari $M$.

\inputdefinition
{}
{

Misalkan
\mathkor {} {V} {dan} {W} {}
merupakan
\definitionsverweis {medan-medan vektor}{}{}
yang terdiferensialkan secara kontinu pada suatu
$C^2$-\definitionsverweis {manifold diferensiabel}{}{}
dengan operator-operator diferensial terkait $D_V$ dan $D_W$. Medan vektor yang bersesuaian dengan komposisi
\mathl{D_V \circ D_W -D_W \circ D_V}{} disebut
\definitionswort {braket Lie}{}
dari kedua medan vektor tersebut. Braket ini dinotasikan dengan $[V,W]$.

}

Urutan operator yang diberi tanda minus dalam definisi merupakan suatu konvensi.


\inputfaktbeweis
{Mannigfaltigkeit/Vektorfelder/Lie-Klammer/Eigenschaften/Fakt}
{Lema}
{}
{

\faktsituation {Misalkan $M$ suatu
$C^2$-\definitionsverweis {manifold diferensiabel}{}{.}}
\faktuebergang {Maka
\definitionsverweis {braket Lie}{}{}
dari
\definitionsverweis {medan-medan vektor}{}{}
memenuhi sifat-sifat berikut.}
\faktfolgerung {\aufzaehlungdrei{Berlaku

\mavergleichskette
{\vergleichskette
{ [V,W]
}
{ = }{ -[W,V]
}
{  }{
}
{    }{
}
{   }{
}
}
{}{}{.}
}{Berlaku

\mavergleichskette
{\vergleichskette
{ [V_1+V_2,W]
}
{ = }{ [V_1,W] + [V_2,W]
}
{  }{
}
{    }{
}
{   }{
}
}
{}{}{.}
}{Untuk suatu fungsi diferensiabel $f$ pada $M$ berlaku

\mavergleichskettedisp
{\vergleichskette
{ [fV,W]
}
{ =} { f [V,W] - \left(  D_Wf  \right)  V
}
{ } {
}
{ } {
}
{ } {
}
}
{}{}{.}
}}
\faktzusatz {}
\faktzusatz {}

}
{

(1) mengikuti dari definisi, dan (2) mengikuti dari hubungan

\mavergleichskettedisp
{\vergleichskette
{ D_{V_1 +V_2 }
}
{ =} { D_{V_1 } + D_{V_2 }
}
{ } {
}
{ } {
}
{ } {
}
}
{}{}{}
antara operator diferensial dan medan vektor. (3). Karena pernyataan tersebut bersifat lokal, kita dapat menetapkan

\mavergleichskettedisp
{\vergleichskette
{ V
}
{ =} { \sum_{i = 1}^n f_i \partial_i
}
{ } {
}
{ } {
}
{ } {
}
}
{}{}{}
dan

\mavergleichskette
{\vergleichskette
{ W
}
{ = }{ \sum_{j = 1}^n g_j \partial_j
}
{  }{
}
{    }{
}
{   }{
}
}
{}{}{.}
Maka

\mavergleichskettedisp
{\vergleichskette
{ fV
}
{ =} { f \sum_{i = 1}^n f_i \partial_i
}
{ =} { \sum_{i = 1}^n ff_i \partial_i
}
{ } {
}
{ } {
}
}
{}{}{}
dan menurut
Lema 11.6
berlaku

\mavergleichskettealignhandlinks
{\vergleichskettealignhandlinks
{ \, [fV,W]
}
{ =} { \sum_{j = 1}^n \left(  \sum_{i = 1}^n \left(  f f_i \partial_i g_j - g_i \partial_i (f f_j) \right)  \right) \partial_j
}
{ =} { \sum_{j = 1}^n \left(  \sum_{i = 1}^n \left(  f f_i \partial_i g_j - g_i \left(  f \partial_i ( f_j) +f_j \partial_i (f)  \right)  \right)  \right)  \partial_j
}
{ =} { f \sum_{j = 1}^n \left(  \sum_{i = 1}^n \left(  f_i \partial_i g_j - g_i \partial_i ( f_j)  \right)  \right) \partial_j - \sum_{j = 1}^n \left(  \sum_{i = 1}^n g_i f_j \partial_i (f)  \right)  \partial_j
}
{ =} { f[V,W] - \sum_{j = 1}^n \left(  \sum_{i = 1}^n g_i \partial_i (f)  \right) f_j \partial_j
}
}
{
\vergleichskettefortsetzungalign
{ =} { f[V,W] - \left(  \sum_{i = 1}^n g_i \partial_i (f)  \right) \left(  \sum_{j = 1}^n f_j \partial_j  \right)
}
{ =} { f[V,W] - D_W(f) V
}
{ } {}
{ } {}
}
{}{.}

}


\chapter{Lembar Kerja 11}
\input{generated/worksheet11.id.build.tex}

\section*{Solusi yang disediakan oleh sumber}
\addcontentsline{toc}{section}{Solusi yang disediakan oleh sumber}
\subsection*{Solusi Soal 11.10}
Tuliskan $\rho=\sqrt{x^2+y^2}$. Fungsi
$F(x,y,z)=(\rho-R)^2+z^2$ tidak diferensiabel ketika $\rho=0$, yaitu pada seluruh sumbu $z$. Pada himpunan terbuka $U=\{\rho>0\}$ berlaku
\mathdisp {dF=2(\rho-R)\left(\frac{x}{\rho}\,dx+\frac{y}{\rho}\,dy\right)+2z\,dz} { . }
Jadi, lokus kritis pembatasan $F\vert_U$ adalah lingkaran yang diberikan oleh $\rho=R$ dan $z=0$; semua titik lain di $U$ merupakan titik reguler. Titik-titik pada sumbu $z$ bukanlah titik kritis, sebab $F$ sama sekali tidak diferensiabel di sana.

Untuk $s<0$, seratnya kosong. Untuk $s=0$, seratnya adalah lingkaran kritis tersebut. Jika $0<s<R^2$, seratnya berupa torus cincin yang mulus. Pada $s=R^2$, lubangnya menutup dan diperoleh torus tanduk yang bertemu sumbu $z$ di titik asal, tempat $F$ tidak diferensiabel. Jika $s>R^2$, seratnya berupa permukaan berbentuk torus gelendong yang bertemu sumbu $z$ di dua titik; pada kedua titik itu $F$ juga tidak diferensiabel.

Animasi di samping memperlihatkan, sebagai contoh untuk $R=2$, serat-serat bagi berbagai nilai $s$ dalam soal ini. Peralihan terjadi pada $s=4=R^2$, atau secara ekuivalen $\sqrt{s}=R$.
\href{https://commons.wikimedia.org/wiki/File:Aufgabe79.27.gif}{Animasi yang memperlihatkan serat-serat fungsi torus pada Soal 11.10 untuk berbagai nilai s.}
Oleh Lukas Freudenberg.

\subsection*{Solusi Soal 11.14}
\input{generated/worksheet11_exercise14_solution.id.build.tex}

\part{Unit 12}
\chapter[Kuliah 12]{Kuliah 12: Bundel Vektor}
\input{generated/lecture12.id.build.tex}

\chapter{Lembar Kerja 12}
\input{generated/worksheet12.id.build.tex}

\section*{Solusi yang disediakan oleh sumber}
\addcontentsline{toc}{section}{Solusi yang disediakan oleh sumber}
\subsection*{Solusi Soal 12.11}
Misalkan $Y$ merupakan gabungan saling lepas dari semua $U_i$. Pada $Y$ kita definisikan sebuah
\definitionsverweis {relasi ekuivalensi}{}{}
$\sim$ dengan menyatakan titik-titik

\mavergleichskette
{\vergleichskette
{ x_i
}
{ \in }{ U_i
}
{  }{
}
{    }{
}
{   }{
}
}
{}{}{}
dan

\mavergleichskette
{\vergleichskette
{ x_j
}
{ \in }{ U_j
}
{  }{
}
{    }{
}
{   }{
}
}
{}{}{}
ekuivalen jika

\mavergleichskette
{\vergleichskette
{ x_i
}
{ \in }{ U_{ij}
}
{  }{
}
{    }{
}
{   }{
}
}
{}{}{,}

\mavergleichskette
{\vergleichskette
{ x_j
}
{ \in }{ U_{ji}
}
{  }{
}
{    }{
}
{   }{
}
}
{}{}{}
dan

\mavergleichskette
{\vergleichskette
{ \varphi_{ji} (x_i)
}
{ = }{ x_j
}
{  }{
}
{    }{
}
{   }{
}
}
{}{}{}.
Sifat-sifat relasi ekuivalensi dijamin oleh syarat koklus; lihat
soal tentang relasi ekuivalensi.
Kita tetapkan

\mavergleichskettedisp
{\vergleichskette
{X
}
{ \defeq} { Y/ \sim
}
{ } {
}
{ } {
}
{ } {
}
}
{}{}{}
dan melengkapi $X$ dengan
\definitionsverweis {topologi hasil bagi}{}{.}
Komposisi-komposisi

\mathdisp {U_i \longrightarrow Y \longrightarrow X} {  }

adalah $\psi_i$, dan $V_i$ merupakan citra pemetaan-pemetaan ini. Jadi diperoleh homeomorfisme-homeomorfisme
\maabb {\psi_i} {U_i } { V_i
} {}.
Untuk

\mavergleichskette
{\vergleichskette
{ x
}
{ \in }{ U_i
}
{  }{
}
{    }{
}
{   }{
}
}
{}{}{}

\mavergleichskettedisp
{\vergleichskette
{ \psi_i (x)
}
{ \in} { V_j
}
{ } {
}
{ } {
}
{ } {
}
}
{}{}{}
berlaku jika dan hanya jika

\mavergleichskette
{\vergleichskette
{ x
}
{ \in }{ U_{ij}
}
{  }{
}
{    }{
}
{   }{
}
}
{}{}{},
sebab tepat dalam kasus ini $x$ diidentifikasi dengan
\mathl{\varphi_{ji}(x)}{}.
Oleh karena itu,

\mavergleichskette
{\vergleichskette
{ \psi_i(U_{ij})
}
{ = }{ V_i \cap V_j
}
{  }{
}
{    }{
}
{   }{
}
}
{}{}{.}
Komutativitas diagram

\mathdisp {\begin{matrix} U_{ij}   & \stackrel{ \varphi_{ji} }{\longrightarrow} &  U_{ji}   & \\ & \!\!\! \!\! \psi_i \searrow & \downarrow \psi_j \!\!\! \!\! & \\ & &  V_i \cap V_j   & \!\!\!\! \!\!\!\!  \\ \end{matrix}} {  }

mengikuti dengan cara yang sama.

\subsection*{Solusi Soal 12.12}
\input{generated/worksheet12_exercise12_solution.id.build.tex}

\part{Unit 13}
\chapter[Kuliah 13]{Kuliah 13: Konstruksi Bundel Vektor}
\input{generated/lecture13.id.build.tex}

\chapter{Lembar Kerja 13}
\setcounter{section}{13}

  \zwischenueberschrift{Soal latihan}

\definitionswort {Produk Kronecker}{}
dari matriks-matriks

\mavergleichskettedisp
{\vergleichskette
{ A
}
{ =} { \left(  a_{ij}  \right)_{1 \leq i \leq m,\, 1 \leq j \leq n}
}
{ } {
}
{ } {
}
{ } {
}
}
{}{}{}
dan

\mavergleichskettedisp
{\vergleichskette
{ B
}
{ =} { \left(  b_{k \ell }  \right)_{1 \leq k \leq p,\, 1 \leq \ell \leq r}
}
{ } {
}
{ } {
}
{ } {
}
}
{}{}{}
diberikan oleh

\mathdisp {\left(  a_{ij} \cdot b_{k \ell}  \right)_{1 \leq i \leq m,\,     1 \leq  k \leq p;\,   1 \leq j \leq n,     \, 1 \leq \ell \leq r   }} {  }

\inputaufgabegibtloesung
{}
{

Hitunglah
\definitionsverweis {produk Kronecker}{}{}
dari kedua matriks
\mathkor {} {\begin{pmatrix}  3   &   -4  \\  5  &  -2 \end{pmatrix}} {dan} {\begin{pmatrix}  -2  &   7   \\  6  &  3 \end{pmatrix}} {.}

}
{} {}

\inputaufgabe
{}
{

Misalkan $K$ sebuah
\definitionsverweis {lapangan}{}{}
dan

\mavergleichskettedisp
{\vergleichskette
{ A
}
{ =} { \left(  a_{ij}  \right)_{1 \leq i \leq m,\, 1 \leq j \leq n}
}
{ } {
}
{ } {
}
{ } {
}
}
{}{}{}
serta

\mavergleichskettedisp
{\vergleichskette
{ B
}
{ =} { \left(  b_{k \ell }  \right)_{1 \leq  k \leq p,\, 1 \leq \ell \leq r}
}
{ } {
}
{ } {
}
{ } {
}
}
{}{}{}
\definitionsverweis {matriks-matriks}{}{}
dengan
\definitionsverweis {pemetaan-pemetaan linear}{}{}
yang bersesuaian
\maabb {A} {K^n  } { K^m
} {}
dan
\maabb {B} {K^r } { K^p
} {.}
Tunjukkan bahwa
\definitionsverweis {produk tensor}{}{}
dari pemetaan-pemetaan linear tersebut, terhadap basis-basis

\mathbed {e_j \otimes e_\ell} {}
 {1 \leq j \leq n,\, 1 \leq \ell \leq r} {}
 {} {} {} {,}
dari
\mathl{K^n \otimes K^r}{} dan

\mathbed {e_i \otimes e_k} {}
 {1 \leq i \leq m,\, 1 \leq  k \leq p} {}
 {} {} {} {,}
dari
\mathl{K^m \otimes K^p}{,}
dideskripsikan oleh
\definitionsverweis {produk Kronecker}{}{}
dari $A$ dan $B$.

}
{} {}

\inputaufgabe
{}
{

Tunjukkan bahwa
\definitionsverweis {produk tensor}{}{}
dari
\definitionsverweis {pita Möbius}{}{}
dengan dirinya sendiri merupakan sebuah bundel garis trivial.

}
{} {}

\inputaufgabe
{}
{

Misalkan
\mathkor {} {M_1} {dan} {M_2} {}
\definitionsverweis {manifold-manifold berorientasi}{}{.}
Tunjukkan bahwa
\definitionsverweis {produk}{}{}

\mathl{M_1 \times M_2}{} merupakan sebuah manifold berorientasi
\zusatzklammer {dengan orientasi yang bergantung pada urutan pada \mathlk{\{1,2\}}{}} {} {.}

}
{} {}

\inputaufgabe
{}
{

Tunjukkan bahwa sfera berdimensi $1$, yaitu $S^1$, merupakan sebuah
\definitionsverweis {manifold diferensiabel}{}{}
yang
\definitionsverweis {dapat diorientasikan}{}{.}

}
{} {}

Misalkan
\mathkor {} {M} {dan} {N} {}
\definitionsverweis {manifold-manifold berorientasi}{}{}
dan
\maabbdisp {\varphi} { M } { N
} {}
sebuah
\definitionsverweis {pemetaan diferensiabel}{}{.}
Pemetaan $\varphi$ disebut \definitionswort {mempertahankan orientasi}{,} jika untuk setiap titik

\mavergleichskette
{\vergleichskette
{ P
}
{ \in }{ M
}
{  }{
}
{    }{
}
{   }{
}
}
{}{}{}
\definitionsverweis {pemetaan singgung}{}{}
\maabbdisp {T\varphi} { T_PM } { T_{\varphi(P)} N
} {}
bijektif dan
\definitionsverweis {mempertahankan orientasi}{}{.}

\inputaufgabe
{}
{

Tunjukkan bahwa
\definitionsverweis {pemetaan antipodal}{}{}
\maabbeledisp {\varphi} {S^1} {S^1
} { P } { -P
} {,}
\definitionsverweis {mempertahankan orientasi}{}{.}

}
{} {}

\inputaufgabe
{}
{

Misalkan $M$ sebuah
\definitionsverweis {manifold berorientasi}{}{.}
Tunjukkan bahwa pemetaan pertukaran
\maabbeledisp {\varphi} {M \times M} {M \times M
} {(P,Q)} {(Q,P)
} {,}
terhadap masing-masing
\definitionsverweis {orientasi produk}{}{,}
tidak harus
\definitionsverweis {mempertahankan orientasi}{}{.}

}
{} {}

\inputaufgabe
{}
{

Misalkan $X$ sebuah
\definitionsverweis {ruang topologis}{}{}
yang hanya terdiri atas
\definitionsverweis {berhingga banyak}{}{}
elemen. Tunjukkan bahwa $X$
\definitionsverweis {kompak}{}{.}

}
{} {}

\inputaufgabe
{}
{

Misalkan $X$ sebuah
\definitionsverweis {ruang topologis}{}{}
dan misalkan

\mavergleichskette
{\vergleichskette
{Y_1 , \ldots , Y_n
}
{ \subseteq }{ X
}
{  }{
}
{    }{
}
{   }{
}
}
{}{}{}
\definitionsverweis {himpunan-himpunan bagian kompak}{}{.}
Tunjukkan bahwa
\definitionsverweis {gabungan}{}{}

\mavergleichskette
{\vergleichskette
{ Y
}
{ = }{ \bigcup_{i = 1}^n Y_i
}
{  }{
}
{    }{
}
{   }{
}
}
{}{}{}
juga kompak.

}
{} {}

\inputaufgabegibtloesung
{}
{

Misalkan $X$ sebuah
\definitionsverweis {ruang kompak}{}{}
dan misalkan

\mavergleichskette
{\vergleichskette
{Y
}
{ \subseteq }{ X
}
{  }{
}
{    }{
}
{   }{
}
}
{}{}{}
sebuah
\definitionsverweis {himpunan bagian tertutup}{}{,}
yang dilengkapi dengan
\definitionsverweis {topologi terinduksi}{}{.}
Tunjukkan bahwa $Y$ juga kompak.

}
{} {}

\inputaufgabegibtloesung
{}
{

Misalkan $X$ sebuah
\definitionsverweis {ruang topologis}{}{}
\definitionsverweis {kompak}{}{}
yang tak kosong dan misalkan
\maabbdisp {f} { X } {\R
} {}
sebuah
\definitionsverweis {fungsi kontinu}{}{.}
Tunjukkan bahwa terdapat suatu

\mavergleichskette
{\vergleichskette
{ x
}
{ \in }{ X
}
{  }{
}
{    }{
}
{   }{
}
}
{}{}{}
dengan

\mathdisp {f(x) \geq f(x')  \text { untuk semua } x' \in X} {  }
.

}
{} {}

\inputaufgabe
{}
{

Misalkan
\maabbdisp {f} {\R} {S^1
} {}
sebuah
\definitionsverweis {pemetaan kontinu}{}{.}
Tunjukkan bahwa
\definitionsverweis {citra}{}{}
dari $f$
\definitionsverweis {homeomorfik}{}{}
dengan suatu
\definitionsverweis {interval}{}{}
terbuka, setengah terbuka, atau tertutup, atau dengan $S^1$.

}
{} {}

\inputaufgabe
{}
{

Misalkan
\maabbdisp {f} {S^1} {\R
} {}
sebuah
\definitionsverweis {pemetaan kontinu}{}{.}
Tunjukkan bahwa
\definitionsverweis {citra}{}{}
dari $f$
\definitionsverweis {homeomorfik}{}{}
dengan suatu
\definitionsverweis {interval tertutup}{}{.}

}
{} {}

\inputaufgabe
{}
{

Tunjukkan bahwa himpunan
\definitionsverweis {bilangan-bilangan riil}{}{} $\R$ tidak
\definitionsverweis {kompak terhadap tutupan}{}{.}

}
{} {}

\inputaufgabe
{}
{

Kita meninjau
\definitionsverweis {bilangan-bilangan asli}{}{}
$\N$ dan melengkapinya dengan
\definitionsverweis {metrik diskret}{}{.}
Tunjukkan bahwa $\N$
\definitionsverweis {
tertutup}{}{}
dan
\definitionsverweis {terbatas}{}{,}
tetapi tidak
\definitionsverweis {
kompak terhadap tutupan}{}{.}

}
{} {}

\inputaufgabegibtloesung
{}
{

Misalkan $X$ sebuah
\definitionsverweis {ruang metrik}{}{}
\definitionsverweis {kompak}{}{.}
Tunjukkan bahwa $X$
\definitionsverweis {
lengkap}{}{.}

}
{} {}

\inputaufgabe
{}
{

Tunjukkan bahwa terdapat sebuah
\definitionsverweis {keluarga terhitung}{}{}
dari
\definitionsverweis {bola-bola terbuka}{}{}
di $\R^n$ yang membentuk sebuah
\definitionsverweis {basis topologi}{}{.}

}
{} {}

\inputaufgabegibtloesung
{}
{

Tunjukkan bahwa himpunan

\mavergleichskettedisp
{\vergleichskette
{ M
}
{ =} { \left\{ (x,y,z) \in \R^3  \mid   x^2+y^4+z^6 = 1        \right\}
}
{ } {
}
{ } {
}
{ } {
}
}
{}{}{}
merupakan sebuah manifold diferensiabel kompak berdimensi dua.

}
{} {}

\inputaufgabegibtloesung
{}
{

Misalkan $M$ sebuah manifold topologis kompak berdimensi $d$,

\mavergleichskette
{\vergleichskette
{ d
}
{ \geq }{ 1
}
{  }{
}
{    }{
}
{   }{
}
}
{}{}{.}
Tunjukkan bahwa terdapat sebuah himpunan bagian terbuka terbatas

\mavergleichskette
{\vergleichskette
{ U
}
{ \subseteq }{ \R^d
}
{  }{
}
{    }{
}
{   }{
}
}
{}{}{}
dan sebuah pemetaan kontinu surjektif
\maabbdisp {\varphi} { U } { M
} {}.

}
{} {}

  \zwischenueberschrift{Soal untuk dikumpulkan}

\inputaufgabe
{4}
{

Tunjukkan bahwa sfera berdimensi $2$, yaitu $S^2$, merupakan sebuah
\definitionsverweis {manifold diferensiabel}{}{}
yang
\definitionsverweis {dapat diorientasikan}{}{.}

}
{} {}

\inputaufgabegibtloesung
{4}
{

Tunjukkan bahwa
\definitionsverweis {pemetaan antipodal}{}{}
\maabbeledisp {} {S^2} {S^2
} {(x,y,z)} {(-x,-y,-z)
} {,}
tidak
\definitionsverweis {mempertahankan orientasi}{}{.}

}
{} {}

\inputaufgabegibtloesung
{4}
{

Misalkan $X$ sebuah
\definitionsverweis {ruang Hausdorff}{}{}
dan misalkan

\mavergleichskette
{\vergleichskette
{Y
}
{ \subseteq }{ X
}
{  }{
}
{    }{
}
{   }{
}
}
{}{}{}
sebuah himpunan bagian yang dilengkapi dengan
\definitionsverweis {topologi terinduksi}{}{.}
Misalkan $Y$
\definitionsverweis {kompak}{}{.}
Tunjukkan bahwa $Y$
\definitionsverweis {tertutup}{}{}
di $X$.

}
{} {}

\inputaufgabe
{4}
{

Misalkan
\mathkor {} {X} {dan} {Y} {}
\definitionsverweis {ruang-ruang topologis}{}{}
dan misalkan
\maabbdisp {\varphi} { X } { Y
} {}
sebuah
\definitionsverweis {pemetaan kontinu}{}{.}
Misalkan $X$
\definitionsverweis {kompak}{}{.}
Tunjukkan bahwa
\definitionsverweis {citra}{}{}

\mavergleichskette
{\vergleichskette
{ \varphi(X)
}
{ \subseteq }{ Y
}
{  }{
}
{    }{
}
{   }{
}
}
{}{}{}
juga kompak.

}
{} {}

\inputaufgabe
{4}
{

Misalkan $V$ sebuah
\definitionsverweis {ruang vektor Euklides}{}{}
ber-
\definitionsverweis {dimensi}{}{}
$n$ dan misalkan
\definitionsverweis {produk}{}{}

\mavergleichskette
{\vergleichskette
{ V^n
}
{ = }{ V \times \cdots \times V
}
{  }{
}
{    }{
}
{   }{
}
}
{}{}{}
dilengkapi dengan
\definitionsverweis {topologi produk}{}{.}
Misalkan $I$ sebuah
\definitionsverweis {interval riil}{}{}
dan
\maabbdisp {\varphi} { I } { V^n
} {}
sebuah
\definitionsverweis {pemetaan kontinu}{}{}
dengan sifat bahwa

\mavergleichskettedisp
{\vergleichskette
{ \varphi(t)
}
{ =} { (\varphi_1(t), \varphi_2(t) , \ldots , \varphi_n(t))
}
{ } {
}
{ } {
}
{ } {
}
}
{}{}{}
untuk setiap

\mavergleichskette
{\vergleichskette
{ t
}
{ \in }{ I
}
{  }{
}
{    }{
}
{   }{
}
}
{}{}{}
merupakan sebuah
\definitionsverweis {basis}{}{}
dari $V$. Tunjukkan bahwa semua basis

\mathbed {\varphi(t)} {}
 {t \in I} {}
 {} {} {} {,}
merepresentasikan
\definitionsverweis {orientasi}{}{}
yang sama pada $V$.

}
{} {}


\section*{Solusi yang disediakan oleh sumber}
\addcontentsline{toc}{section}{Solusi yang disediakan oleh sumber}
\subsection*{Solusi Soal 13.1}
\input{generated/worksheet13_exercise01_solution.id.build.tex}
\subsection*{Solusi Soal 13.10}
Misalkan

\mavergleichskettedisp
{\vergleichskette
{ Y
}
{ =} { \bigcup_{i \in I} V_i
}
{ } {
}
{ } {
}
{ } {
}
}
{}{}{}
sebuah tutupan oleh himpunan-himpunan bagian terbuka

\mavergleichskette
{\vergleichskette
{ V_i
}
{ \subseteq }{ Y
}
{  }{
}
{    }{
}
{   }{
}
}
{}{}{.}
Ini berarti terdapat himpunan-himpunan terbuka

\mavergleichskette
{\vergleichskette
{ U_i
}
{ \subseteq }{ X
}
{  }{
}
{    }{
}
{   }{
}
}
{}{}{}
sedemikian sehingga

\mavergleichskette
{\vergleichskette
{ V_i
}
{ = }{ Y \cap U_i
}
{  }{
}
{    }{
}
{   }{
}
}
{}{}{.}
Karena itu,

\mavergleichskettedisp
{\vergleichskette
{ Y
}
{ \subseteq} { \bigcup_{i \in I} U_i
}
{ } {
}
{ } {
}
{ } {
}
}
{}{}{.}
Karena $Y$ tertutup di $X$, maka $X \setminus Y$ terbuka, sehingga

\mavergleichskettedisp
{\vergleichskette
{ X
}
{ =} { Y \cup (X \setminus Y)
}
{ =} { \bigcup_{i \in I} U_i  \cup  (X \setminus Y)
}
{ } {
}
{ } {
}
}
{}{}{}
merupakan sebuah tutupan terbuka dari $X$. Karena $X$ kompak, terdapat sebuah subtutupan hingga

\mavergleichskettedisp
{\vergleichskette
{ X
}
{ =} { \bigcup_{i \in J} U_i  \cup  (X \setminus Y)
}
{ } {
}
{ } {
}
{ } {}
}
{}{}{.}
Hal ini pada gilirannya berarti

\mavergleichskettedisp
{\vergleichskette
{ Y
}
{ =} { \bigcup_{i \in J} V_i
}
{ =} { Y \cap \bigcup_{i \in J} U_i
}
{ } {
}
{ } {}
}
{}{}{.}

\subsection*{Solusi Soal 13.11}
Berdasarkan
Fakta,

\mavergleichskette
{\vergleichskette
{ f(X)
}
{ \subseteq }{ \R
}
{  }{
}
{    }{
}
{   }{
}
}
{}{}{}
kompak, sehingga menurut
Fakta
himpunan ini
\definitionsverweis {tertutup}{}{}
dan
\definitionsverweis {terbatas}{}{.}
Secara khusus,

\mavergleichskette
{\vergleichskette
{ f(X)
}
{ \subseteq }{ (-\infty,M]
}
{  }{
}
{    }{
}
{   }{
}
}
{}{}{}
untuk suatu bilangan riil $M$. Karena

\mavergleichskette
{\vergleichskette
{ X
}
{ \neq }{ \emptyset
}
{  }{
}
{    }{
}
{   }{
}
}
{}{}{}
menurut
Fakta,
$f(X)$ memiliki sebuah
\definitionsverweis {supremum}{}{}
$s$ di $\R$. Karena himpunan itu tertutup, menurut Fakta,
supremum tersebut termasuk dalam $f(X)$, sehingga merupakan maksimum dari $f(X)$. Oleh karena itu, terdapat pula suatu

\mavergleichskette
{\vergleichskette
{ x
}
{ \in }{ X
}
{  }{
}
{    }{
}
{   }{
}
}
{}{}{}
dengan

\mavergleichskette
{\vergleichskette
{ f(x)
}
{ = }{ s
}
{  }{
}
{    }{
}
{   }{
}
}
{}{}{.}

\subsection*{Solusi Soal 13.16}
\input{generated/worksheet13_exercise16_solution.id.build.tex}
\subsection*{Solusi Soal 13.18}
\input{generated/worksheet13_exercise18_solution.id.build.tex}
\subsection*{Solusi Soal 13.19}
Untuk setiap titik
\mathl{P \in M}{} kita pilih sebuah lingkungan peta terbuka
\mathl{P \in U_P}{} dengan sebuah peta
\maabbdisp {\alpha_P} {U_P } {V_P
} {}
dan
\mathl{V_P \subseteq \R^d}{.} Dengan beralih ke sebuah lingkungan bola dari
\mathl{\alpha_P(P) \in V_P}{} beserta prabayangnya, kita dapat mengandaikan bahwa $V_P$ adalah bola-bola terbuka dengan jari-jari paling besar
\mathl{{ \frac{   1 }{  3 } }}{}. Keluarga

\mathbed {U_P} {}
 {P\in M} {}
 {} {} {} {,}
menutupi manifold tersebut. Karena $M$ kompak, terdapat berhingga banyak titik
\mathl{P_1 , \ldots , P_n}{} sedemikian sehingga

\mathbed {U_i=U_{P_i}} {}
 {i=1 , \ldots , n} {}
 {} {} {} {,}
juga menutupi manifold tersebut. Kita menempatkan bola-bola terbuka
\mathl{V_i=V_{P_i}}{} di dalam
\zusatzklammer {\anfuehrung{$\R^d$ baru}{}} {} {}
dengan titik-titik pusat

\mathdisp {(1,0 , \ldots , 0), (2,0 , \ldots , 0) , \ldots , (n,0 , \ldots , 0)} { . }

Berdasarkan syarat jari-jari, bola-bola ini saling lepas. Kita meninjau himpunan terbuka terbatas
\mathl{U= \bigcup_{i=1}^n V_i}{.} Peta-peta $\alpha_i$ menghasilkan pemetaan-pemetaan kontinu

\mathdisp {\varphi_i :V_i \stackrel{ \left(  \alpha_i  \right)^{-1} }{\longrightarrow} U_i \longrightarrow M} { . }

Karena himpunan-himpunan tersebut saling lepas, dengan menetapkan
\mathl{\varphi(z)= \varphi_i(z)}{,} jika
\mathl{z \in V_i}{} berlaku, kita memperoleh sebuah pemetaan kontinu
\maabbdisp {\varphi} { U } {M
} {.}
Karena sifat tutupan tersebut, pemetaan ini surjektif.

\subsection*{Solusi Soal 13.21}
Untuk pemetaan antipodal, kita memperoleh matriks
:$
\ dA = \begin{pmatrix}
-1 & 0 & 0 \\
0 & -1 & 0 \\
0 & 0 & -1 \end{pmatrix}
$ $ \in $ $\R^{3x3}$.

Jadi, determinan $dA$ tepat sama dengan -1, yakni
 $ \det(dA) = -1 <0$.
Dengan orientasi standar $S^2$ sebagai batas bola satuan, pemetaan antipodal tersebut membalik orientasi dan dengan demikian tidak mempertahankan orientasi.

Secara umum, derajat pemetaan antipodal pada $S^n$ adalah $(-1)^{n+1}$; karena itu, pemetaan tersebut membalik orientasi tepat ketika dimensinya genap.

\subsection*{Solusi Soal 13.22}
Kita tunjukkan bahwa untuk setiap titik

\mavergleichskette
{\vergleichskette
{ x
}
{ \in }{ U
}
{ = }{ X \setminus Y
}
{    }{
}
{   }{
}
}
{}{}{}
terdapat sebuah lingkungan terbuka

\mavergleichskette
{\vergleichskette
{ x
}
{ \in }{ V
}
{ \subseteq }{ U
}
{    }{
}
{   }{
}
}
{}{}{}.
Gabungan semua lingkungan tersebut kemudian sama dengan $U$. Jadi, misalkan

\mavergleichskette
{\vergleichskette
{ x
}
{ \notin }{ Y
}
{  }{
}
{    }{
}
{   }{
}
}
{}{}{}
diberikan. Untuk setiap titik

\mavergleichskette
{\vergleichskette
{ y
}
{ \in }{ Y
}
{  }{
}
{    }{
}
{   }{
}
}
{}{}{}
terdapat himpunan-himpunan terbuka

\mavergleichskette
{\vergleichskette
{ x
}
{ \in }{ V_y
}
{  }{
}
{    }{
}
{   }{
}
}
{}{}{}
dan

\mavergleichskette
{\vergleichskette
{ y
}
{ \in }{ W_y
}
{  }{
}
{    }{
}
{   }{
}
}
{}{}{,}
yang saling lepas. Himpunan-himpunan terbuka

\mathbed {W_y} {}
 {y \in Y} {}
 {} {} {} {,}
menutupi $Y$. Berdasarkan kekompakan, terdapat sebuah subtutupan hingga, katakanlah

\mavergleichskettedisp
{\vergleichskette
{ Y
}
{ \subseteq} { \bigcup_{i = 1}^n  W_i
}
{ } {
}
{ } {
}
{ } {
}
}
{}{}{.}
Maka

\mavergleichskette
{\vergleichskette
{ V
}
{ = }{ \bigcap_{i = 1}^n  V_i
}
{  }{
}
{    }{
}
{   }{
}
}
{}{}{}
terbuka sebagai irisan berhingga dari himpunan-himpunan terbuka, dan karena itu merupakan sebuah lingkungan terbuka dari $x$ yang saling lepas dengan $Y$.


\part{Unit 14}
\chapter[Kuliah 14]{Kuliah 14: Bentuk Diferensial pada Manifold}
\setcounter{section}{14}

  \zwischenueberschrift{Bentuk diferensial pada manifold}

\inputdefinition
{}
{

Misalkan $M$ suatu
\definitionsverweis {manifold diferensiabel}{}{.}
Suatu
\definitionswortpraemath {k}{ bentuk diferensial }{}
\zusatzklammer {atau $k$-\definitionswort {bentuk}{} atau \definitionswort {bentuk berderajat}{} $k$} {} {}
adalah suatu
\definitionsverweis {penampang}{}{}
dalam
\definitionsverweis {produk eksterior}{}{}
ke-$k$ dari
\definitionsverweis {bundel kotangen}{}{,}
yaitu suatu pemetaan
\maabbeledisp {\omega} { M } { \bigwedge^k T^*M
} { P } { \omega(P)
} {,}
dengan

\mavergleichskette
{\vergleichskette
{ \omega(P)
}
{ \in }{ \bigwedge^k T_P^*M
}
{  }{
}
{    }{
}
{   }{
}
}
{}{}{.}

}

Kita melambangkan himpunan semua $k$-bentuk pada $M$ dengan

\mathdisp {{ \mathcal E }^{ k } ( M  )} { . }

\inputbemerkung
{}
{

Jadi, suatu $k$-\definitionsverweis {bentuk}{}{}
memasangkan kepada setiap titik $P$ pada manifold sebuah elemen dari
\mathl{\bigwedge^k T^*_PM}{}. Berdasarkan
Korolari 57.9 (Aljabar Linear (Osnabrück 2024-2025))
dan Teorema 58.4 (Aljabar Linear (Osnabrück 2024-2025)),
hal ini sama dengan suatu
\definitionsverweis {pemetaan multilinear}{}{}
\definitionsverweis {selang-seling}{}{}
\maabbdisp {} {T_PM \times \cdots \times T_PM } { \R
} {.}
Pemetaan yang bersesuaian ini juga kita lambangkan dengan
\mathl{\omega(P)}{;} untuk $k$
\definitionsverweis {vektor singgung}{}{}

\mavergleichskette
{\vergleichskette
{ v_1 , \ldots , v_k
}
{ \in }{ T_PM
}
{  }{
}
{    }{
}
{   }{
}
}
{}{}{}
dengan demikian,

\mathdisp {\omega(P)(v_1 , \ldots , v_k)} {  }

adalah suatu bilangan riil. Di sini muncul dua jenis argumen yang sama sekali berbeda: di satu pihak titik pada manifold, dan di pihak lain elemen-elemen ruang singgung di titik tersebut. Ketergantungan pada vektor-vektor singgung relatif sederhana karena pemetaannya multilinear selang-seling, sedangkan ketergantungan pada titik manifold dapat serumit apa pun. Karena produk-produk eksterior bundel kotangen sendiri merupakan manifold menurut
Soal 14.2,
kita dapat langsung membicarakan bentuk diferensial yang kontinu atau
\zusatzklammer {jika $M$ merupakan manifold $C^2$} {} {}
diferensiabel.

Untuk

\mavergleichskette
{\vergleichskette
{ k
}
{ = }{ 0
}
{  }{
}
{    }{
}
{   }{
}
}
{}{}{}
ruang kotangen hanya muncul secara formal; suatu $0$-bentuk tidak lain adalah sebuah fungsi
\maabb {f} { M } {\R
} {.}
Suatu $1$-bentuk
\zusatzklammer {juga disebut
\definitionswortenp{bentuk Pfaff}{}} {} {}
memasangkan suatu bilangan riil kepada setiap titik $P$ dan setiap vektor singgung pada titik tersebut. Untuk

\mavergleichskette
{\vergleichskette
{k
}
{ > }{ n
}
{ = }{ \operatorname{dim} \, M
}
{    }{
}
{   }{
}
}
{}{}{}
produk eksterior ke-$k$ adalah ruang nol sehingga sama sekali tidak ada bentuk taktrivial berderajat ini. Kasus yang sangat penting adalah

\mavergleichskette
{\vergleichskette
{ k
}
{ = }{ n
}
{ = }{ \operatorname{dim} \, M
}
{    }{
}
{   }{
}
}
{}{}{.}
Dalam hal ini, produk eksterior ke-$n$ memiliki rank $1$
\zusatzklammer {artinya, dimensinya $1$ di setiap titik} {} {}
dan suatu penampang di dalamnya secara lokal dideskripsikan oleh satu fungsi saja. Gambaran yang berguna ialah bahwa, untuk $n$ vektor singgung, bilangan
\mathl{\omega(P)(v_1 , \ldots , v_n)}{} menyatakan volume
\zusatzklammer {\anfuehrung{berorientasi}{}} {} {}
dari paralelotop yang direntang oleh vektor-vektor tersebut di ruang singgung. Gambaran ini juga membantu untuk $k$ yang lebih kecil: dengan
\mathl{\omega(P) (v_1 , \ldots , v_k)}{} kita dapat menghitung volume berdimensi $k$ dari paralelotop yang dibangkitkan oleh $k$ vektor singgung. Gambaran ini akan dibuat tepat ketika kita mengintegralkan pada suatu
\definitionsverweis {submanifold tertutup}{}{}
berdimensi $k$.

}



\inputfaktbeweis
{Mannigfaltigkeit/Differentialform/Elementare Eigenschaften/Fakt}
{Lemma}
{}
{

\faktsituation {Misalkan $M$ suatu
\definitionsverweis {manifold diferensiabel}{}{}
dan, untuk

\mavergleichskette
{\vergleichskette
{ k
}
{ \in }{ \N
}
{  }{
}
{    }{
}
{   }{
}
}
{}{}{},
misalkan
\mathl{{ \mathcal E }^{ k } (  M   )}{} himpunan semua
$k$-\definitionsverweis {bentuk}{}{}
pada $M$.}
\faktuebergang {Maka berlaku sifat-sifat berikut.}
\faktfolgerung {\aufzaehlungfuenf{Dengan operasi-operasi alaminya,
\mathl{{ \mathcal E }^{ k } (  M   )}{} membentuk
\definitionsverweis {ruang vektor riil}{}{.}
}{Untuk suatu bentuk diferensial

\mavergleichskette
{\vergleichskette
{ \omega
}
{ \in }{
{ \mathcal E }^{ k } (  M   )
}
{  }{
}
{    }{
}
{   }{
}
}
{}{}{}
dan suatu fungsi
\maabbdisp {f} { M } {\R
} {},

\mavergleichskette
{\vergleichskette
{ f \omega
}
{ \in }{
{ \mathcal E }^{ k } (  M   )
}
{  }{
}
{    }{
}
{   }{
}
}
{}{}{,}
dengan $f \omega$ didefinisikan oleh

\mavergleichskettedisp
{\vergleichskette
{ (f \omega)(P)
}
{ \defeq} { f(P) \omega (P)
}
{ } {
}
{ } {
}
{ } {
}
}
{}{}{}.
}{Setiap
\definitionsverweis {fungsi diferensiabel}{}{}
$C^1$
\maabbdisp {f} { M } {\R
} {}
melalui
\definitionsverweis {pemetaan singgung}{}{}

\mathl{Tf}{} mendefinisikan suatu
$1$-\definitionsverweis {bentuk diferensial}{}{}
\maabbeledisp {df} { M } { T^*M
} { P } { T_Pf
} {,}
dengan ruang singgung dari $\R$ di
\mathl{f(P)}{} diidentifikasi dengan $\R$. Hal ini menghasilkan suatu pemetaan
\maabbeledisp {d} {C^1(M,\R)} {
{ \mathcal E }^{  1 } (  M   )
} { f } {df
} {.}
}{Jika

\mavergleichskette
{\vergleichskette
{ M
}
{ \subseteq }{ \R^m
}
{  }{
}
{    }{
}
{   }{
}
}
{}{}{}
merupakan himpunan bagian terbuka, maka di bawah identifikasi

\mavergleichskettedisp
{\vergleichskette
{ TM
}
{ \cong} { M \times \R^m
}
{ } {
}
{ } {
}
{ } {
}
}
{}{}{}
pemetaan pada (3) sama dengan
\maabbeledisp {} { M } { M \times ( \R^m)^*
} { P } { (P, \left(  D f   \right)_{ P } \left(   -  \right) )
} {.}
}{Pemetaan $d$ pada (3) adalah
$\R$-\definitionsverweis {linear}{}{.}
}}
\faktzusatz {}
\faktzusatz {}

}
{

\teilbeweis {}{}{}
{(1) dan (2) langsung mengikuti dari peninjauan titik demi titik.}
{} \teilbeweis {}{}{}
{(3). Untuk setiap titik

\mavergleichskette
{\vergleichskette
{ P
}
{ \in }{ M
}
{  }{
}
{    }{
}
{   }{
}
}
{}{}{},

\maabbdisp {T_Pf} { T_PM } { T_{f(P)} \R \cong \R
} {}
merupakan suatu
\definitionsverweis {pemetaan linear}{}{}
menurut Lemma 9.10  (3),
dan dengan demikian merupakan elemen dari
\mathl{T_P^* M}{,} yang kita lambangkan dengan
\mathl{(df)_P}{}. Oleh karena itu, pemasangan
\mathl{P \mapsto (df)_P}{} merupakan suatu
\definitionsverweis {bentuk diferensial}{}{.}}
{}
\teilbeweis {}{}{}
{(4) mengikuti dari
Lemma 9.10  (1).}
{}
\teilbeweis {}{}{}
{(5). Pemetaan pada (3) ditentukan untuk setiap titik

\mavergleichskette
{\vergleichskette
{ P
}
{ \in }{ M
}
{  }{
}
{    }{
}
{   }{
}
}
{}{}{}
pada setiap lingkungan terbuka. Karena itu, kita boleh menganggap bahwa

\mavergleichskette
{\vergleichskette
{ M
}
{ \subseteq }{ \R^m
}
{  }{
}
{    }{
}
{   }{
}
}
{}{}{}
merupakan himpunan terbuka, sehingga pernyataannya mengikuti dari (4) dan
Proposisi 45.6 (Analisis (Osnabrück 2021-2023)).}
{}

}

Untuk suatu himpunan terbuka

\mavergleichskette
{\vergleichskette
{ V
}
{ \subseteq }{ \R^n
}
{  }{
}
{    }{
}
{   }{
}
}
{}{}{},
tersedia fungsi-fungsi koordinat
\maabbdisp {x_j} {V} {\R
} {},
yang, jika diberikan suatu peta, dapat dipindahkan ke suatu himpunan bagian terbuka dari manifold. Di setiap titik

\mavergleichskette
{\vergleichskette
{ Q
}
{ \in }{ V
}
{  }{
}
{    }{
}
{   }{
}
}
{}{}{},
bentuk-bentuk

\mathbed {dx_j} {}
 {j=1 , \ldots , n} {}
 {} {} {} {,}
membentuk basis ruang kotangen di $Q$. Basis ini tidak lain adalah
\definitionsverweis {basis dual}{}{}
dari basis standar di ruang ambien $\R^n$, yang digunakan sebagai ruang singgung di seluruh $V$. Untuk suatu himpunan bagian beranggotakan $k$

\mavergleichskettedisp
{\vergleichskette
{ J
}
{ =} { \{j_1 , \ldots , j_k\}
}
{ \subseteq} { \{1 , \ldots , n\}
}
{ } {
}
{ } {
}
}
{}{}{},
kita tetapkan

\mavergleichskettedisp
{\vergleichskette
{  dx_J
}
{ =} {  dx_{j_1} \wedge \ldots \wedge dx_{j_k}
}
{ } {
}
{ } {
}
{ } {
}
}
{}{}{,}
yang merupakan $k$-bentuk yang sangat sederhana pada $V$. Untuk setiap titik

\mavergleichskette
{\vergleichskette
{ Q
}
{ \in }{ V
}
{  }{
}
{    }{
}
{   }{
}
}
{}{}{},
berlaku

\mavergleichskettedisp
{\vergleichskette
{ dx_J(Q)
}
{ =} { (dx_{j_1} \wedge \ldots \wedge dx_{j_k} )(Q)
}
{ =} { dx_{j_1}(Q) \wedge \ldots \wedge dx_{j_k} (Q)
}
{ =} { e_{j_1}^* \wedge \ldots \wedge e_{j_k}^*
}
{ } {
}
}
{}{}{.}
Aksi bentuk ini pada

\mavergleichskette
{\vergleichskette
{  v_1 \wedge \ldots \wedge v_k
}
{ \in }{ \bigwedge^k T_QV
}
{  }{
}
{    }{
}
{   }{
}
}
{}{}{}
, menurut
Teorema 58.4 (Aljabar Linear (Osnabrück 2024-2025))
, diberikan oleh

\mavergleichskettedisp
{\vergleichskette
{ \left(  e_{j_1}^* \wedge \ldots \wedge e_{j_k}^*  \right)  (v_1 \wedge \ldots \wedge v_k)
}
{ =} { \det  \left(  e_{j_i}^*(v_\ell)  \right)_{i \ell}
}
{ =} { \det  \left(  (v_\ell)_{j_i}  \right)_{i \ell}
}
{ } {
}
{ } {
}
}
{}{}{.}
Menurut
Teorema 58.1 (Aljabar Linear (Osnabrück 2024-2025)),
nilai-nilai bentuk diferensial
\zusatzklammer {dengan \mathlk{j_1 <   \ldots < j_k}{}} {} {}

\mavergleichskettedisp
{\vergleichskette
{ dx_J
}
{ =} { dx_{j_1} \wedge \ldots \wedge dx_{j_k}
}
{ } {
}
{ } {
}
{ } {
}
}
{}{}{}
membentuk suatu
\definitionsverweis {basis}{}{}
dari
\mathl{\bigwedge^k T_Q^*V}{}
untuk setiap titik $Q$. Oleh karena itu, setiap
\definitionsverweis {bentuk diferensial}{}{} pada $V$ yang berderajat $k$

\mavergleichskette
{\vergleichskette
{ \omega
}
{ \in }{
{ \mathcal E }^{ k } ( V  )
}
{  }{
}
{    }{
}
{   }{
}
}
{}{}{}
dapat ditulis secara tunggal sebagai

\mavergleichskettedisp
{\vergleichskette
{ \omega
}
{ =} { \sum_{J,\, { \# \left(  J \right) } = k} f_J dx_J
}
{ } {
}
{ } {
}
{ } {
}
}
{}{}{}
dengan fungsi-fungsi yang ditentukan secara tunggal
\maabbdisp {f_J} {V} {\R
} {.}



\inputfaktbeweis
{Mannigfaltigkeit/Differentialform/Lokale Beschreibung/Fakt}
{Lemma}
{}
{

\faktsituation {Misalkan $M$ suatu
\definitionsverweis {manifold diferensiabel}{}{}
dan

\mavergleichskette
{\vergleichskette
{ U
}
{ \subseteq }{ M
}
{  }{
}
{    }{
}
{   }{
}
}
{}{}{}
suatu
\definitionsverweis {himpunan bagian terbuka}{}{}
dengan peta
\maabbdisp {\alpha} { U } { V
} {},
dan misalkan

\mavergleichskette
{\vergleichskette
{ V
}
{ \subseteq }{ \R^n
}
{  }{
}
{    }{
}
{   }{
}
}
{}{}{}
terbuka. Misalkan
\maabbdisp {x_j} { U } {\R
} {}
fungsi-fungsi koordinat yang bersesuaian,

\mavergleichskette
{\vergleichskette
{ 1
}
{ \leq }{ j
}
{ \leq }{ n
}
{    }{
}
{   }{
}
}
{}{}{.}}
\faktfolgerung {Maka setiap
\definitionsverweis {bentuk diferensial}{}{} pada $U$ yang berderajat $k$

\mavergleichskette
{\vergleichskette
{ \omega
}
{ \in }{
{ \mathcal E }^{ k } (  U   )
}
{  }{
}
{    }{
}
{   }{
}
}
{}{}{}
dapat ditulis secara tunggal sebagai

\mavergleichskettedisp
{\vergleichskette
{ \omega
}
{ =} { \sum_{J,\, { \# \left(   J  \right) } = k}  f_J dx_J
}
{ } {
}
{ } {
}
{ } {
}
}
{}{}{}
dengan fungsi-fungsi yang ditentukan secara tunggal
\maabbdisp {f_J} { U } {\R
} {.}}
\faktzusatz {}
\faktzusatz {}

}
{

Hal ini mengikuti dari
Teorema 58.1 (Aljabar Linear (Osnabrück 2024-2025)).

}



\inputfaktbeweis
{Mannigfaltigkeit/Pfaffsche Differentialform zu Funktion/Beschreibung mit partiellen Ableitungen/Fakt}
{Korollar}
{}
{

\faktsituation {Misalkan $M$ suatu
\definitionsverweis {manifold diferensiabel}{}{}
dan

\mavergleichskette
{\vergleichskette
{U
}
{ \subseteq }{ M
}
{  }{
}
{    }{
}
{   }{
}
}
{}{}{}
suatu
\definitionsverweis {himpunan bagian terbuka}{}{}
dengan peta
\maabbdisp {\alpha} { U } { V
} {},
dan misalkan

\mavergleichskette
{\vergleichskette
{V
}
{ \subseteq }{\R^n
}
{  }{
}
{    }{
}
{   }{
}
}
{}{}{}
terbuka. Misalkan
\maabbdisp {x_j} { U } {\R
} {}
fungsi-fungsi koordinat yang bersesuaian,

\mavergleichskette
{\vergleichskette
{ 1
}
{ \leq }{j
}
{ \leq }{n
}
{    }{
}
{   }{
}
}
{}{}{.}
Misalkan
\maabbdisp {f} { U } {\R
} {}
suatu
\definitionsverweis {fungsi diferensiabel}{}{.}}
\faktfolgerung {Maka $1$-\definitionsverweis {bentuk diferensial}{}{}
$df$ yang bersesuaian mempunyai representasi\zusatzfussnote {Turunan-turunan
\mathl{{ \frac{   \partial  f  }{  \partial  x_j  } }}{} diperkenalkan dalam kuliah kedelapan} {.} {}

\mavergleichskettedisp
{\vergleichskette
{ df
}
{ =} { \sum_{j = 1}^n { \frac{   \partial   f   }{  \partial   x_j   } } dx_j
}
{ } {
}
{ } {
}
{ } {
}
}
{}{}{.}}
\faktzusatz {}
\faktzusatz {}

}
{

Kita dapat langsung menganggap bahwa semuanya berlangsung pada himpunan terbuka

\mavergleichskette
{\vergleichskette
{V
}
{ \subseteq }{\R^n
}
{  }{
}
{    }{
}
{   }{
}
}
{}{}{}.
Untuk setiap titik

\mavergleichskette
{\vergleichskette
{Q
}
{ \in }{ V
}
{  }{
}
{    }{
}
{   }{
}
}
{}{}{},
berlaku kesamaan berikut antara
\definitionsverweis {bentuk-bentuk linear}{}{}
pada $\R^n$,

\mavergleichskettealign
{\vergleichskettealign
{(df)_Q
}
{ =} { \left(D f \right)_{ Q }
}
{ =} { \left(  { \frac{   \partial   f   }{  \partial   x_1   } } ( Q)   , \,   \ldots  , \,   { \frac{   \partial   f   }{  \partial   x_n   } } ( Q)                 \right)
}
{ =} { { \frac{   \partial   f   }{  \partial   x_1   } } ( Q) dx_1 + \cdots + { \frac{   \partial   f   }{  \partial   x_n   } } ( Q) dx_n
}
{ } {
}
}
{}
{}{.}

}

  \zwischenueberschrift{Penarikan balik bentuk diferensial}

\inputdefinition
{}
{

Misalkan
\mathkor {} {L} {dan} {M} {}
\definitionsverweis {dua manifold diferensiabel}{}{}
dan misalkan
\maabbdisp {\varphi} { L } { M
} {}
suatu
\definitionsverweis {pemetaan diferensiabel}{}{.}
Misalkan $\omega$ suatu $k$-\definitionsverweis {bentuk diferensial}{}{}
pada $M$. Maka pemetaan yang menentukan suatu $k$-bentuk pada $L$ diberikan oleh

\mathdisp {(P ,v_1 , \ldots , v_k) \longmapsto   \omega(\varphi(P), T_P(\varphi)(v_1) , \ldots ,  T_P(\varphi)(v_k) )} {  }

Pemetaan ini merupakan
\definitionsverweis {pemetaan selang-seling}{}{}.
Bentuk yang bersesuaian dengannya disebut, terhadap $\varphi$, \definitionswort {bentuk tarik balik}{} berderajat $k$. Bentuk ini dilambangkan dengan

\mathdisp {\varphi^* \omega} {  }.

}



\inputfaktbeweis
{Mannigfaltigkeit/Differenzierbare Abbildung/Zurückziehen von Differentialformen/Elementare Eigenschaften/Fakt}
{Lemma}
{}
{

\faktsituation {Misalkan
\mathkor {} {L} {dan} {M} {}
\definitionsverweis {dua manifold diferensiabel}{}{}
dan
\maabbdisp {\varphi} { L } { M
} {}
suatu
\definitionsverweis {pemetaan diferensiabel}{}{.}}
\faktuebergang {Maka
\definitionsverweis {penarikan balik bentuk diferensial}{}{}
memenuhi sifat-sifat berikut.}
\faktfolgerung {\aufzaehlungvier{Untuk suatu fungsi

\mavergleichskette
{\vergleichskette
{ f
}
{ \in }{ \operatorname{Map} (M,\R)
}
{ = }{
{ \mathcal E }^{  0 } (  M   )
}
{    }{
}
{   }{
}
}
{}{}{},
berlaku

\mavergleichskette
{\vergleichskette
{ \varphi^*(f)
}
{ = }{ f \circ \varphi
}
{  }{
}
{    }{
}
{   }{
}
}
{}{}{.}
}{Pemetaan-pemetaan
\maabbeledisp {\varphi^*} {
{ \mathcal E }^{ k } (  M   ) } {
{ \mathcal E }^{ k } (  L   )
} { \omega} { \varphi^* \omega
} {,}
adalah
$\R$-\definitionsverweis {linear}{}{.}
}{Jika

\mavergleichskette
{\vergleichskette
{ L
}
{ \subseteq }{ M
}
{  }{
}
{    }{
}
{   }{
}
}
{}{}{}
merupakan suatu
\definitionsverweis {submanifold terbuka}{}{,}
dan pemetaan di atas adalah inklusi, maka penarikan balik suatu bentuk diferensial $\omega$ tidak lain adalah pembatasan
\mathl{\omega \vert_L}{} pada himpunan bagian ini.
}{Misalkan $N$ suatu manifold diferensiabel lain dan
\maabbdisp {\psi} { M } { N
} {}
suatu pemetaan diferensiabel lain. Maka berlaku

\mavergleichskettedisp
{\vergleichskette
{ ( \psi \circ \varphi )^*(\omega)
}
{ =} { \varphi^*(\psi^*(\omega))
}
{ } {
}
{ } {
}
{ } {
}
}
{}{}{}
untuk setiap bentuk diferensial

\mavergleichskette
{\vergleichskette
{ \omega
}
{ \in }{
{ \mathcal E }^{ k } (  N   )
}
{  }{
}
{    }{
}
{   }{
}
}
{}{}{.}
}}
\faktzusatz {}
\faktzusatz {}

}
{

\teilbeweis {}{}{}
{(1) langsung mengikuti dari
Definisi 14.6.}
{}
\teilbeweis {}{}{}
{(2). Untuk bentuk-bentuk diferensial
\mathkor {} {\omega_1} {dan} {\omega_2} {}
dan skalar-skalar

\mavergleichskette
{\vergleichskette
{ a,b
}
{ \in }{\R
}
{  }{
}
{    }{
}
{   }{
}
}
{}{}{},
kita harus menunjukkan bahwa

\mavergleichskette
{\vergleichskette
{ \varphi^*(a \omega_1+ b \omega_2)
}
{ = }{ a \varphi^* \omega_1 + b \varphi^* \omega_2
}
{  }{
}
{    }{
}
{   }{
}
}
{}{}{}.
Kesamaan bentuk diferensial semacam ini berarti bahwa kesamaan tersebut berlaku di setiap titik

\mavergleichskette
{\vergleichskette
{ P
}
{ \in }{ L
}
{  }{
}
{    }{
}
{   }{
}
}
{}{}{}
dan untuk setiap $k$-tupel vektor singgung

\mavergleichskette
{\vergleichskette
{ v_1 , \ldots , v_k
}
{ \in }{ T_PL
}
{  }{
}
{    }{
}
{   }{
}
}
{}{}{}.
Karena itu, pernyataannya mengikuti dari

\mavergleichskettealignhandlinks
{\vergleichskettealignhandlinks
{ \left(  \varphi^*(a \omega_1 + b \omega_2)  \right) \left(  P, v_1 , \ldots , v_k  \right)
}
{ =} { (a \omega_1 + b \omega_2) \left(  \varphi(P), T_P(\varphi)(v_1) , \ldots ,  T_P(\varphi)(v_k)  \right)
}
{ =} { a \omega_1 \left(  \varphi(P), T_P(\varphi)(v_1) , \ldots ,  T_P(\varphi)(v_k)  \right)
}
{ \,\,\,\, \,} {+ b \omega_2 \left(  \varphi(P), T_P(\varphi)(v_1) , \ldots ,  T_P(\varphi)(v_k)  \right)
}
{ =} { \left(  a  \varphi^* \omega_1  \right)(P, v_1 , \ldots , v_k)  + \left(  b \varphi^* \omega_2  \right) (P, v_1 , \ldots , v_k)
}
}
{}
{}{.}}
{}
\teilbeweis {}{}{}
{(3) langsung mengikuti dari definisi.}
{}
\teilbeweis {}{}{}
{(4). Misalkan

\mavergleichskette
{\vergleichskette
{ Q
}
{ \in }{ L
}
{  }{
}
{    }{
}
{   }{
}
}
{}{}{,}

\mavergleichskette
{\vergleichskette
{ u_1 , \ldots , u_k
}
{ \in }{ T_QL
}
{  }{
}
{    }{
}
{   }{
}
}
{}{}{},
dan $\omega$ suatu $k$-bentuk pada $N$. Maka, dengan menggunakan
Lemma 9.10  (4),
berlaku

\mavergleichskettealigndrucklinks
{\vergleichskettealigndrucklinks
{(( \psi \circ \varphi)^* (\omega))(Q, u_1 , \ldots , u_k)
}
{ =} { \omega \left(  ( \psi \circ \varphi) (Q), T_Q( \psi \circ \varphi)(u_1) , \ldots , T_Q(  \psi \circ \varphi) (u_k)  \right)
}
{ =} { \omega \left(  \psi ( \varphi (Q)), \left(  T_{\varphi(Q)}\psi  \right)  ((T_Q\varphi)(u_1)) , \ldots , \left(  T_{\varphi(Q)} \psi  \right) ((T_Q\varphi)(u_k))  \right)
}
{ =} { \left(  \psi^*(\omega)  \right) \left(  \varphi(Q),T_Q(\varphi)(u_1) , \ldots , T_Q(\varphi)(u_k)  \right)
}
{ =} { \left(  \varphi^* \left(  \psi^*(\omega)  \right)   \right) (Q,u_1 , \ldots , u_k)
}
}
{}{}{,}
dan inilah pernyataannya.}
{}

}



\inputfaktbeweis
{Offene Mengen/Differenzierbare Abbildung/Zurückziehen von Differentialformen/In Koordinaten/Fakt}
{Lemma}
{}
{

\faktsituation {Misalkan
\mathkor {} {U \subseteq \R^n} {dan} {V \subseteq \R^m} {}
\definitionsverweis {himpunan-himpunan bagian terbuka}{}{,}
dengan koordinat masing-masing dilambangkan
\mathl{x_1 , \ldots , x_n}{} dan
\mathl{y_1 , \ldots , y_m}{}. Misalkan
\maabbdisp {\varphi} { U } { V
} {}
suatu
\definitionsverweis {pemetaan diferensiabel}{}{}
dan misalkan $\omega$ suatu
$k$-\definitionsverweis {bentuk diferensial}{}{}
pada $V$ dengan representasi

\mavergleichskettedisp
{\vergleichskette
{ \omega
}
{ =} { \sum_{I,\, { \# \left(   I  \right) } = k}  f_I  dy_I
}
{ } {
}
{ } {
}
{ } {
}
}
{}{}{,}
dengan
\maabb {f_I} { V } {\R
} {}
merupakan
\definitionsverweis {fungsi-fungsi}{}{.}}
\faktfolgerung {Maka
\definitionsverweis {bentuk tarik balik}{}{}
mempunyai representasi

\mavergleichskettealignhandlinks
{\vergleichskettealignhandlinks
{ \varphi^* \omega
}
{ =} { \sum_{I,\, { \# \left(   I  \right) } = k}  \left(  f_I \circ \varphi  \right)  \left( \sum_{J,\, { \# \left(   J  \right) } = k} \det  \left(  \left(  { \frac{   \partial  \varphi_i   }{  \partial   x_j   } }  \right)_{ i \in I, j \in J}  \right)  dx_J \right)
}
{ =} { \sum_{J,\, { \# \left(   J  \right) } = k}   \left( \sum_{I,\, { \# \left(   I  \right) } = k} \left(  f_I \circ \varphi  \right)  \det  \left(  \left(  { \frac{   \partial  \varphi_i   }{  \partial   x_j   } }  \right)_{ i \in I, j \in J}  \right)  \right)  dx_J
}
{ } {
}
{ } {
}
}
{}{}{.}}
\faktzusatz {}
\faktzusatz {}

}
{

Kesamaan kedua didasarkan pada pengaturan ulang yang sederhana. Berdasarkan
Lemma 14.7 (2),
kita dapat membatasi diri pada kasus

\mavergleichskette
{\vergleichskette
{ \omega
}
{ = }{ f_Idy_I
}
{  }{
}
{    }{
}
{   }{
}
}
{}{}{}.
Kita tetapkan

\mavergleichskette
{\vergleichskette
{ f
}
{ = }{ f_I
}
{  }{
}
{    }{
}
{   }{
}
}
{}{}{}
dan boleh menganggap

\mavergleichskette
{\vergleichskette
{ I
}
{ = }{ \{ 1 , \ldots , k \}
}
{  }{
}
{    }{
}
{   }{
}
}
{}{}{}.
Kita menunjukkan kesamaan antara kedua $k$-bentuk pada $U$ dengan menunjukkan bahwa keduanya menghasilkan nilai yang sama untuk setiap titik

\mavergleichskette
{\vergleichskette
{ P
}
{ \in }{ U
}
{  }{
}
{    }{
}
{   }{
}
}
{}{}{}
dan setiap produk eksterior
\mathl{e_{j_1 } \wedge \ldots \wedge e_{j_k }}{} dengan

\mavergleichskette
{\vergleichskette
{ 1
}
{ \leq }{ j_1
}
{ < }{ \ldots
}
{ <   }{ j_k
}
{ \leq  }{ n
}
}
{}{}{}.
Di satu pihak,

\mavergleichskettealigndrucklinks
{\vergleichskettealigndrucklinks
{ \varphi^*( f dy_1 \wedge \ldots \wedge dy_k)(P,e_{j_1 } \wedge \ldots \wedge e_{j_k }  )
}
{ =} { ( f dy_1 \wedge \ldots \wedge dy_k)(\varphi(P),T_P(\varphi)(e_{j_1 }) \wedge \ldots \wedge T_P(\varphi)(e_{j_k }))
}
{ =} { f(\varphi(P)) (dy_1 \wedge \ldots \wedge dy_k)( \begin{pmatrix}  { \frac{   \partial  \varphi_1   }{  \partial   x_{j_1 }  } } (P )  \\\vdots\\  { \frac{   \partial  \varphi_m   }{  \partial   x_{j_1 }  } } (P )   \end{pmatrix} \wedge \ldots \wedge \begin{pmatrix}  { \frac{   \partial  \varphi_1   }{  \partial   x_{j_k }  } } (P )  \\\vdots\\  { \frac{   \partial  \varphi_m   }{  \partial   x_{j_k }  } } (P )   \end{pmatrix} )
}
{ =} { f(\varphi(P)) \cdot \det  \left(  (dy_i) \begin{pmatrix}  { \frac{   \partial  \varphi_1   }{  \partial   x_{j_\ell}  } } (P )  \\\vdots\\  { \frac{   \partial  \varphi_m   }{  \partial   x_{j_\ell}  } } (P )   \end{pmatrix}  \right)_{1 \leq i, \ell \leq k}
}
{ =} { f(\varphi(P)) \cdot \det  \left(  { \frac{   \partial  \varphi_i   }{  \partial   x_{j_\ell}  } } (P )  \right)_{1 \leq i, \ell \leq k}
}
}
{}{}{.}
Di pihak lain, jika jumlah tersebut diterapkan pada
\mathl{e_{j_1 } \wedge \ldots \wedge e_{j_k }}{},
maka

\mavergleichskette
{\vergleichskette
{ dx_J (e_{j_1 } \wedge \ldots \wedge e_{j_k })
}
{ = }{ 0
}
{  }{
}
{    }{
}
{   }{
}
}
{}{}{}
kecuali apabila

\mavergleichskette
{\vergleichskette
{ J
}
{ = }{ \{j_1 , \ldots , j_k\}
}
{  }{
}
{    }{
}
{   }{
}
}
{}{}{,}
dan dalam kasus itu nilainya $1$, sehingga diperoleh nilai yang sama.

}



\inputfaktbeweis
{Offene Mengen/Differenzierbare Abbildung/Zurückziehen von Volumenform auf gleichdimensionalem Raum/In Koordinaten/Fakt}
{Korollar}
{}
{

\faktsituation {Misalkan

\mavergleichskette
{\vergleichskette
{ U,V
}
{ \subseteq }{ \R^n
}
{  }{
}
{    }{
}
{   }{
}
}
{}{}{}
\definitionsverweis {himpunan-himpunan bagian terbuka}{}{,}
dengan koordinat masing-masing dilambangkan
\mathl{x_1 , \ldots , x_n}{} dan
\mathl{y_1 , \ldots , y_n}{}. Misalkan
\maabbdisp {\varphi} { U } { V
} {}
suatu
\definitionsverweis {pemetaan diferensiabel}{}{}
dan misalkan $\omega$ suatu
$n$-\definitionsverweis {bentuk diferensial}{}{}
pada $V$ dengan representasi

\mavergleichskettedisp
{\vergleichskette
{ \omega
}
{ =} { f dy_1 \wedge \ldots \wedge dy_n
}
{ } {
}
{ } {
}
{ } {
}
}
{}{}{.}}
\faktfolgerung {Maka
\definitionsverweis {bentuk tarik balik}{}{}
mempunyai representasi

\mavergleichskettedisp
{\vergleichskette
{ \varphi^* \omega
}
{ =} { (f \circ \varphi) \cdot \det  \left(  \left(  { \frac{   \partial  \varphi_i   }{  \partial   x_j   } }  \right)_{1 \leq  i, j \leq n}  \right) dx_1 \wedge \ldots \wedge dx_n
}
{ } {
}
{ } {
}
{ } {}
}
{}{}{.}}
\faktzusatz {}
\faktzusatz {}

}
{

Hal ini langsung mengikuti dari
Lemma 14.8.

}
Dengan demikian, fungsi-fungsi yang mendeskripsikan suatu bentuk diferensial mempunyai perilaku transformasi yang sama dengan kerapatan berorientasi pada suatu peta. Untuk kerapatan ukuran positif, faktor determinannya harus diganti dengan nilai mutlak determinan.



\inputfaktbeweis
{Differentialform/Lokal/Zurückziehen unter partiell konstanter Abbildung/Fakt}
{Korollar}
{}
{

\faktsituation {Misalkan

\mavergleichskette
{\vergleichskette
{ U
}
{ \subseteq }{ \R^n
}
{  }{
}
{    }{
}
{   }{
}
}
{}{}{}
dan

\mavergleichskette
{\vergleichskette
{ V
}
{ \subseteq }{ \R^m
}
{  }{
}
{    }{
}
{   }{
}
}
{}{}{}
\definitionsverweis {himpunan-himpunan bagian terbuka}{}{,}
dengan koordinat masing-masing dilambangkan
\mathl{x_1 , \ldots , x_n}{} dan
\mathl{y_1 , \ldots , y_m}{}. Misalkan
\maabbdisp {\varphi} { U } { V
} {}
suatu
\definitionsverweis {pemetaan diferensiabel}{}{}
dengan
\mathl{\varphi_{i_0 }}{} konstan untuk suatu

\mavergleichskette
{\vergleichskette
{ i_0
}
{ \in }{ \{1 , \ldots , m\}
}
{  }{
}
{    }{
}
{   }{
}
}
{}{}{},
dan misalkan $\omega$ suatu
$k$-\definitionsverweis {bentuk diferensial}{}{}
pada $V$ dengan representasi

\mavergleichskettedisp
{\vergleichskette
{ \omega
}
{ =} { f dy_I
}
{ } {
}
{ } {
}
{ } {
}
}
{}{}{}
dengan

\mavergleichskette
{\vergleichskette
{ i_0
}
{ \in }{ I
}
{  }{
}
{    }{
}
{   }{
}
}
{}{}{.}}
\faktfolgerung {Maka

\mavergleichskette
{\vergleichskette
{ \varphi^*\omega
}
{ = }{ 0
}
{  }{
}
{    }{
}
{   }{
}
}
{}{}{.}}
\faktzusatz {}
\faktzusatz {}

}
{

Menurut
Lemma 14.8,
berlaku

\mavergleichskettedisp
{\vergleichskette
{ \varphi^* \omega
}
{ =} { \sum_{J,\, { \# \left(   J  \right) } = k} (f \circ \varphi) \cdot \det  \left(  \left(  { \frac{   \partial  \varphi_i   }{  \partial   x_j   } }  \right) _{ i \in I, j \in J}  \right) dx_J
}
{ } {
}
{ } {
}
{ } {}
}
{}{}{.}
Karena

\mavergleichskettedisp
{\vergleichskette
{ { \frac{   \partial  \varphi_{i_0 }  }{  \partial   x_j   } }
}
{ =} { 0
}
{ } {
}
{ } {
}
{ } {
}
}
{}{}{}
untuk semua

\mavergleichskette
{\vergleichskette
{ j
}
{ \in }{ J
}
{  }{
}
{    }{
}
{   }{
}
}
{}{}{,}
untuk setiap $J$, matriks tersebut mempunyai satu baris bernilai $0$, sehingga semua determinannya bernilai $0$.

}


\chapter{Lembar Kerja 14}
\setcounter{section}{14}

  \zwischenueberschrift{Soal latihan}

\inputaufgabe
{}
{

Misalkan $M$ sebuah
\definitionsverweis {manifold diferensiabel}{}{}
dengan
\definitionsverweis {bundel kotangen}{}{}

\mathl{T^*M}{.} Tunjukkan bahwa pada
\mathl{\bigwedge^k T^*M}{} untuk setiap $k$ dapat didefinisikan suatu
\definitionsverweis {topologi}{}{}
sedemikian sehingga, untuk setiap peta
\maabb {\alpha} { U } { V
} {}
pemetaan
\maabbdisp {} {\bigwedge^k T^*U} { \bigwedge^k T^*V \cong V \times \bigwedge^k \R^{n*}
} {}
merupakan sebuah
\definitionsverweis {homeomorfisme}{}{.}

}
{} {}
Dengan demikian, kita dapat berbicara tentang bentuk diferensial kontinu maupun terukur.

\inputaufgabe
{}
{

Misalkan $M$ sebuah
$C^2$-\definitionsverweis {manifold diferensiabel}{}{}
dengan
\definitionsverweis {bundel kotangen}{}{}

\mathl{T^*M}{.} Tunjukkan bahwa
\mathl{\bigwedge^k T^*M}{} merupakan sebuah manifold diferensiabel untuk setiap $k$.

}
{} {}
Dengan demikian, kita juga dapat berbicara tentang bentuk diferensial yang terdiferensialkan secara kontinu. Secara lebih umum, jika $M$ sebuah
$C^m$-\definitionsverweis {manifold}{}{}
dan

\mavergleichskette
{\vergleichskette
{  \ell
}
{ < }{ m
}
{  }{
}
{    }{
}
{   }{
}
}
{}{}{}
maka
\mathl{{ \mathcal E }^{ k }_{\ell } ( M  )}{} menyatakan himpunan semua bentuk diferensial yang terdiferensialkan secara kontinu hingga orde $\ell$ dan berderajat $k$.

\inputaufgabe
{}
{

Misalkan $M$ sebuah
\definitionsverweis {manifold diferensiabel}{}{}
dan
\mathl{{ \mathcal E }^{ k } (  M   )}{} himpunan semua
$k$-\definitionsverweis {bentuk}{}{}
pada $M$. Tunjukkan bahwa ${ \mathcal E }^{ k } (  M   )$ merupakan sebuah
$R$-\definitionsverweis {modul}{}{}
untuk

\mavergleichskette
{\vergleichskette
{R
}
{ = }{ C^0(M,\R)
}
{  }{
}
{    }{
}
{   }{
}
}
{}{}{.}

}
{} {}

\inputaufgabe
{}
{

Misalkan $M$ sebuah
\definitionsverweis {manifold diferensiabel}{}{,}

\mavergleichskette
{\vergleichskette
{ P
}
{ \in }{ M
}
{  }{
}
{    }{
}
{   }{
}
}
{}{}{}
sebuah titik, dan

\mavergleichskette
{\vergleichskette
{ f
}
{ \in }{ C^1(M,\R)
}
{  }{
}
{    }{
}
{   }{
}
}
{}{}{}
sebuah
\definitionsverweis {fungsi yang terdiferensialkan secara kontinu}{}{.}
Misalkan

\mavergleichskette
{\vergleichskette
{ v
}
{ \in }{ T_PM
}
{  }{
}
{    }{
}
{   }{
}
}
{}{}{}
sebuah
\definitionsverweis {vektor singgung}{}{,}
yang direpresentasikan oleh suatu lintasan diferensiabel
\maabbdisp {\gamma} {{]{-\delta},\delta [}} {M
} {}
dengan

\mavergleichskette
{\vergleichskette
{ \gamma(0)
}
{ = }{ P
}
{  }{
}
{    }{
}
{   }{
}
}
{}{}{.}
Tunjukkan kesamaan

\mavergleichskettedisp
{\vergleichskette
{(df)(P,v)
}
{ =} { (f \circ \gamma)'(0)
}
{ } {
}
{ } {
}
{ } {
}
}
{}{}{.}

}
{} {}

\inputaufgabegibtloesung
{}
{

Misalkan $M$ sebuah
\definitionsverweis {manifold diferensiabel}{}{}
dan
\maabbdisp {f} { M } {\R
} {}
sebuah
\definitionsverweis {fungsi yang terdiferensialkan secara kontinu}{}{.}
Pada titik

\mavergleichskette
{\vergleichskette
{ P
}
{ \in }{ M
}
{  }{
}
{    }{
}
{   }{
}
}
{}{}{}
fungsi tersebut mempunyai suatu
\definitionsverweis {ekstremum lokal}{}{.}
Tunjukkan bahwa

\mavergleichskettedisp
{\vergleichskette
{ (df)_P
}
{ =} { 0
}
{ } {
}
{ } {
}
{ } {
}
}
{}{}{.}

}
{} {}

\inputaufgabegibtloesung
{}
{

Misalkan $M$ sebuah
\definitionsverweis {manifold diferensiabel}{}{}
\definitionsverweis {kompak}{}{}
dan
\maabbdisp {f} { M } {\R
} {}
sebuah
\definitionsverweis {fungsi yang terdiferensialkan secara kontinu}{}{.}
Tunjukkan bahwa terdapat sedikitnya dua titik pada $M$
\zusatzklammer {dengan syarat $M$ mempunyai sedikitnya dua titik} {} {}
di mana
$1$-\definitionsverweis {bentuk diferensial}{}{}
$df$ bernilai nol.

}
{} {}

\inputaufgabe
{}
{

Misalkan
\mathl{i: M\subseteq \R^n}{} sebuah
\definitionsverweis {submanifold tertutup}{}{.}
Tunjukkan bahwa untuk setiap
\definitionsverweis {fungsi diferensiabel}{}{}
\maabbdisp {f} {\R^n } {\R
} {}
berlaku hubungan

\mavergleichskettedisp
{\vergleichskette
{ i^* (df)
}
{ =} { d(f \circ i )
}
{ } {
}
{ } {
}
{ } {
}
}
{}{}{.}

}
{} {}

\inputaufgabe
{}
{

Misalkan
\maabbeledisp {f} {\R} {\R
} { t } {f(t)
} {,}
sebuah
\definitionsverweis {fungsi}{}{}
\definitionsverweis {yang terdiferensialkan secara kontinu}{}{,}
dan misalkan

\mavergleichskette
{\vergleichskette
{ \omega
}
{ = }{ g(s)ds
}
{  }{
}
{    }{
}
{   }{
}
}
{}{}{}
sebuah
$1$-\definitionsverweis {bentuk diferensial}{}{}
pada $\R$. Tentukan
\mathl{f^* \omega}{.}

}
{} {}

\inputaufgabegibtloesung
{}
{

Kita meninjau pemetaan
\maabbeledisp {\varphi} {\R^3} {\R^2
} {(x,y,z)} {(xyz^2,xy-z^3)
} {.}
Hitunglah matriks pemetaan
\maabbdisp {\bigwedge^2 T_P (\varphi)} { \bigwedge^2 T_P \R^3 } { \bigwedge^2 T_{\varphi(P)}  \R^2
} {}
di titik

\mavergleichskette
{\vergleichskette
{ P
}
{ = }{ (1,3,5)
}
{  }{
}
{    }{
}
{   }{
}
}
{}{}{}
terhadap suatu basis yang sesuai.

}
{} {}

\inputaufgabe
{}
{

Kita meninjau
\definitionsverweis {pemetaan terdiferensialkan kontinu}{}{}
\maabbeledisp {\varphi} {\R^2 \setminus \{(0,0)\}} {\R^2 \setminus \{(0,0)\}
} {(u,v)} {(u^2,v^3-u)
} {,}
dan
$2$-\definitionsverweis {bentuk diferensial}{}{}

\mavergleichskettedisp
{\vergleichskette
{ \omega
}
{ =} { { \frac{   1  }{  x^2+y^2 } } dx \wedge dy
}
{ } {
}
{ } {
}
{ } {
}
}
{}{}{.}
Tentukan
\definitionsverweis {bentuk diferensial tarik balik}{}{}

\mathl{\varphi^* \omega}{.}

}
{} {}

\inputaufgabegibtloesung
{}
{

Kita meninjau
$1$-\definitionsverweis {bentuk diferensial}{}{}

\mavergleichskettedisp
{\vergleichskette
{ \omega
}
{ =} { -vdu+udv
}
{ } {
}
{ } {
}
{ } {
}
}
{}{}{}
pada
$1$-\definitionsverweis {sfera}{}{}

\mavergleichskette
{\vergleichskette
{ S^1
}
{ \subset }{ \R^2
}
{  }{
}
{    }{
}
{   }{
}
}
{}{}{,}
dengan koordinat-koordinat pada ruang ambien $\R^2$ dinotasikan sebagai
\mathkor {} {u} {dan} {v} {.}
Tentukan
\mathl{\pi^* \omega}{} di bawah
\definitionsverweis {pemetaan terdiferensialkan kontinu}{}{}
\maabbeledisp {\pi} {\R^2 \setminus \{0\}} { S^{1}
} {(x , y) } { { \frac{   1  }{  \sqrt{x^2 + y^2}  } } (x , y)
} {.}

}
{} {}

\inputaufgabegibtloesung
{}
{

Hitunglah bentuk diferensial tarik balik
\mathl{\varphi^* \tau}{} untuk

\mavergleichskettedisp
{\vergleichskette
{ \tau
}
{ =} { dx \wedge dy \wedge dz - w dx \wedge dy \wedge dw + \cos (xy)  dx \wedge dz \wedge dw-yw dy \wedge dz \wedge dw
}
{ } {
}
{ } {
}
{ } {
}
}
{}{}{}
di bawah pemetaan
\maabbeledisp {\varphi} {\R^3} {\R^4
} {(r,s,t)} { (r^2s,t, \sin  r ,e^{st}) = (x,y,z,w)
} {.}

}
{} {}

\inputaufgabegibtloesung
{}
{

Misalkan
\maabbdisp {\psi} { L } { M
} {}
sebuah
\definitionsverweis {pemetaan diferensiabel}{}{}
antara
\definitionsverweis {dua manifold diferensiabel}{}{,}
misalkan
\maabbdisp {f} { M } {\R
} {}
sebuah fungsi pada $M$, dan $\omega$ sebuah
$k$-\definitionsverweis {bentuk}{}{}
pada $M$. Tunjukkan bahwa

\mavergleichskettedisp
{\vergleichskette
{ \psi^*(f \omega)
}
{ =} { \psi^*(f)  \psi^* \omega
}
{ } {
}
{ } {
}
{ } {
}
}
{}{}{.}

}
{} {}

\inputaufgabegibtloesung
{}
{

Misalkan

\mavergleichskette
{\vergleichskette
{ W
}
{ \subseteq }{ \R^m
}
{  }{
}
{    }{
}
{   }{
}
}
{}{}{}
dan

\mavergleichskette
{\vergleichskette
{ U
}
{ \subseteq }{ \R^n
}
{  }{
}
{    }{
}
{   }{
}
}
{}{}{}
merupakan
\definitionsverweis {himpunan-himpunan bagian terbuka}{}{}
dan misalkan
\maabbdisp {\psi} { W } { U
} {}
sebuah
\definitionsverweis {pemetaan}{}{} yang
\definitionsverweis {terdiferensialkan secara kontinu}{}{.}
Misalkan
\maabbdisp {f} { U } {\R
} {}
sebuah fungsi yang terdiferensialkan secara kontinu. Simpulkan dari
aturan rantai
bahwa

\mavergleichskettedisp
{\vergleichskette
{ d (\psi^*f)
}
{ =} { \psi^*(df)
}
{ } {
}
{ } {
}
{ } {
}
}
{}{}{,}
dengan $\psi^*$ menyatakan
\definitionsverweis {penarikan balik bentuk diferensial}{}{.}

}
{} {}

  \zwischenueberschrift{Soal untuk dikumpulkan}

\inputaufgabe
{6}
{

Misalkan $M$ sebuah
\definitionsverweis {manifold diferensiabel}{}{}
dengan
\definitionsverweis {bundel kotangen}{}{}

\mathl{T^*M}{.} Misalkan $\omega$ sebuah
$k$-\definitionsverweis {bentuk diferensial}{}{,} yaitu suatu pemetaan
\maabbdisp {\omega} { M } {\bigwedge^k T^*M
} {}
dengan
\mathl{\omega(P) \in \bigwedge^k T^*_P M}{} untuk setiap
\mathl{P \in M}{,} dengan produk eksterior ini dilengkapi topologi alaminya
\zusatzklammer {lihat
Soal 14.1} {} {.}
Tunjukkan bahwa pernyataan-pernyataan berikut ekuivalen:
\aufzaehlungdrei{ $\omega$
\definitionsverweis {
kontinu}{}{.}
}{Untuk setiap
\definitionsverweis {peta}{}{}
\maabb {\alpha} { U } { V
} {}
dengan
\mathl{V \subseteq \R^n}{} dan representasi lokal
\mathl{\alpha_* \omega =  \sum_{J,\, { \# \left(   J  \right) } =k} f_J dx_J}{,}
fungsi-fungsi $f_J$ kontinu.
}{Terdapat suatu tutupan terbuka
\mathl{M= \bigcup_{i \in I} U_i}{} oleh
\definitionsverweis {domain-domain peta}{}{} $U_i$ sedemikian sehingga, dalam representasi lokal
\mathl{\alpha_{i*} \omega =  \sum_{J,\, { \# \left(   J  \right) } =k} f_{i J} dx_J}{,}
fungsi-fungsi $f_{i J}$ kontinu.
}

}
{} {}

\inputaufgabe
{4}
{

Misalkan
\maabbdisp {\varphi} { L } { M
} {}
sebuah
\definitionsverweis {pemetaan diferensiabel}{}{}
antara dua
\definitionsverweis {manifold diferensiabel}{}{}
\mathkor {} {L} {dan} {M} {.}
Misalkan
\mathkor {} {\omega \in
{ \mathcal E }^{ k } (  M   )} {dan} {\tau \in
{ \mathcal E }^{  \ell } (  M   )} {}
\definitionsverweis {bentuk-bentuk diferensial}{}{} pada $M$. Tunjukkan kesamaan

\mavergleichskettedisp
{\vergleichskette
{ \varphi^* ( \omega \wedge \tau)
}
{ =} {\varphi^* \omega \wedge \varphi^* \tau
}
{ } {
}
{ } {
}
{ } {
}
}
{}{}{.}

}
{} {}

\inputaufgabe
{4}
{

Kita meninjau
\definitionsverweis {pemetaan terdiferensialkan kontinu}{}{}
\maabbeledisp {\varphi} {\R^3 \setminus \{(0,0,0)\}} {N = \left\{ (x,y,z) \in \R^3  \mid   z \neq 0        \right\}
} {(u,v,w)} {(uvw,u^2-vw^5,u^2+v^2+w^2)
} {,}
dan
$2$-\definitionsverweis {bentuk diferensial}{}{}

\mathdisp {\omega =  z^2dx \wedge dy+{ \frac{   xy }{ z } }dx \wedge dz+(xe^y-z)dy \wedge dz} {  }

pada $N$. Tentukan
\definitionsverweis {bentuk diferensial tarik balik}{}{}

\mathl{\varphi^* \omega}{.}

}
{} {}

\inputaufgabe
{6}
{

Kita meninjau
\definitionsverweis {permukaan bola satuan}{}{}

\mavergleichskette
{\vergleichskette
{ S^2
}
{ \subseteq }{ \R^3
}
{  }{
}
{    }{
}
{   }{
}
}
{}{}{,}
dengan koordinat-koordinat pada $\R^3$ dinotasikan sebagai $x,y,z$. Untuk titik-titik manakah

\mavergleichskette
{\vergleichskette
{ P
}
{ \in }{ S^2
}
{  }{
}
{    }{
}
{   }{
}
}
{}{}{}
pembatasan
\mathkor {} {dx} {dan} {dy} {}
pada $S^2$ membentuk suatu
\definitionsverweis {basis}{}{}
dari
\definitionsverweis {ruang kotangen}{}{}

\mathl{T^*_PS^2}{.}

}
{} {}


\section*{Solusi yang disediakan oleh sumber}
\addcontentsline{toc}{section}{Solusi yang disediakan oleh sumber}
\subsection*{Solusi Soal 14.5}
\input{generated/worksheet14_exercise05_solution.id.build.tex}
\subsection*{Solusi Soal 14.6}
Jika $f$ konstan, maka

\mavergleichskette
{\vergleichskette
{ df
}
{ = }{ 0
}
{  }{
}
{    }{
}
{   }{
}
}
{}{}{}
dan pernyataan tersebut jelas. Jadi, misalkan $f$ tidak konstan. Berdasarkan
Fakta,
fungsi kontinu tersebut mencapai maksimum global dan minimum global, katakanlah masing-masing di titik
\mathkor {} {P} {dan} {Q} {.}
Berdasarkan
soal,
maka

\mavergleichskettedisp
{\vergleichskette
{ (df)_P
}
{ =} { (df)_Q
}
{ =} { 0
}
{ } {
}
{ } {
}
}
{}{}{.}

\subsection*{Solusi Soal 14.9}
Matriks Jacobi dari $\varphi$ secara umum adalah

\mathdisp {\begin{pmatrix}  yz^2 &  xz^2 &  2xyz \\  y  &  x  &  -3z^2 \end{pmatrix}} { . }

Oleh karena itu, di titik

\mavergleichskette
{\vergleichskette
{ P
}
{ = }{ (1,3,5)
}
{  }{
}
{    }{
}
{   }{
}
}
{}{}{}
matriks Jacobi tersebut adalah

\mathdisp {\begin{pmatrix}  75 &  25 &  30 \\  3 &  1 &  -75  \end{pmatrix}} {  }
.

Matriks ini mendeskripsikan pemetaan linear
\maabbdisp {L} {T_P\R^3 \cong \R^3} { T_{\varphi(P)} \R^2 \cong \R^2
} {}
terhadap basis-basis standar. Kita menentukan representasi matriks untuk pemetaan
\maabbdisp {\bigwedge^2 L} {\bigwedge^2 \R^3 } { \bigwedge^2 \R^2
} {}
terhadap basis-basis
$e_1 \wedge e_2,\, e_1 \wedge e_3$ , $e_2 \wedge e_3$
\zusatzklammer {di ruas kiri} {} {}
dan
\mathl{e_1 \wedge e_2}{}
\zusatzklammer {di ruas kanan} {} {.}
Untuk itu, kita perlu menghitung citra produk-produk eksterior tersebut. Kita peroleh

\mavergleichskettealign
{\vergleichskettealign
{ (\bigwedge^2 L ) (e_1 \wedge e_2)
}
{ =} { L(e_1) \wedge L(e_2)
}
{ =} { (75e_1 + 3e_2) \wedge (  25 e_1 + 1 e_2 )
}
{ =} { 75 e_1 \wedge e_2 + 75 e_2 \wedge  e_1
}
{ =} { 75 e_1 \wedge e_2 - 75 e_1 \wedge  e_2
}
}
{
\vergleichskettefortsetzungalign
{ =} { 0
}
{ } {}
{ } {}
{ } {}
}
{}{,}

\mavergleichskettealign
{\vergleichskettealign
{ (\bigwedge^2 L ) (e_1 \wedge e_3)
}
{ =} { L(e_1) \wedge L(e_3)
}
{ =} { (75e_1 + 3e_2) \wedge (  30 e_1 - 75 e_2 )
}
{ =} { - 75^2 e_1 \wedge e_2 + 90 e_2 \wedge  e_1
}
{ =} { (-5625 - 90) e_1 \wedge e_2
}
}
{
\vergleichskettefortsetzungalign
{ =} { -5715 e_1 \wedge e_2
}
{ } {}
{ } {}
{ } {}
}
{}{,}
dan

\mavergleichskettealign
{\vergleichskettealign
{ (\bigwedge^2 L ) (e_2 \wedge e_3)
}
{ =} { L(e_2) \wedge L(e_3)
}
{ =} { (  25 e_1 + 1 e_2 )   \wedge ( 30 e_1 - 75 e_2)
}
{ =} { - 1875 e_1 \wedge e_2 + 30 e_2 \wedge  e_1
}
{ =} { -  1905 e_1 \wedge  e_2
}
}
{}
{}{.}
Jadi, matriks yang mendeskripsikan pemetaan tersebut adalah

\mathdisp {\left(  0,-5715,-1905                     \right)} { . }

\subsection*{Solusi Soal 14.11}
\input{generated/worksheet14_exercise11_solution.id.build.tex}
\subsection*{Solusi Soal 14.12}
\input{generated/worksheet14_exercise12_solution.id.build.tex}
\subsection*{Solusi Soal 14.13}
\input{generated/worksheet14_exercise13_solution.id.build.tex}
\subsection*{Solusi Soal 14.14}
\input{generated/worksheet14_exercise14_solution.id.build.tex}

\part{Unit 15}
\chapter[Kuliah 15]{Kuliah 15: Integrasi pada Manifold}
\input{generated/lecture15.id.build.tex}

\chapter{Lembar Kerja 15}
\input{generated/worksheet15.id.build.tex}

\section*{Solusi yang disediakan oleh sumber}
\addcontentsline{toc}{section}{Solusi yang disediakan oleh sumber}
\subsection*{Solusi Soal 15.1}
\input{generated/worksheet15_exercise01_solution.id.build.tex}
\subsection*{Solusi Soal 15.11}
Kita peroleh

\mavergleichskettealign
{\vergleichskettealign
{ \gamma^* \omega
}
{ =} { (-t^2)^3 d(-t^2) - (t^3-1) (t+2)d (t^3-1)  - t^2(t+2)^2d(t+2)
}
{ =} { 2 t^7 dt     - 3t^2 (t^4 +2t^3 -t -2) dt - t^2(  t^2  +4 t +4) dt
}
{ =} {  (2t^7 -3t^6 -6t^5   +3t^3 +6t^2 -t^4  -4t^3 -4t^2       )dt
}
{ =} { (2t^7 -3t^6 -6t^5   -t^4  -t^3 +2 t^2       )dt
}
}
{}
{}{.}
Oleh karena itu,

\mavergleichskettealign
{\vergleichskettealign
{ \int_\gamma \omega
}
{ =} { \int_{-1}^0    (2t^7 -3t^6 -6t^5   -t^4  -t^3 +2 t^2       )dt
}
{ =} { ({ \frac{   1 }{  4 } }  t^8 - { \frac{   3 }{  7 } }  t^7 - t^6- { \frac{   1 }{  5 } } t^5 - { \frac{   1 }{  4 } } t^4 + { \frac{   2 }{  3 } } t^3) \vert_{  -1 }^{  0 }
}
{ =} { -( { \frac{   1 }{  4 } }  + { \frac{   3 }{  7 } }  -1 + { \frac{   1 }{  5 } }  - { \frac{   1 }{  4 } }  - { \frac{   2 }{  3 } } )
}
{ =} { - { \frac{   3 }{  7 } }  +1 - { \frac{   1 }{  5 } } + { \frac{   2 }{  3 } }
}
}
{
\vergleichskettefortsetzungalign
{ =} { { \frac{   -45 +105-21+70 }{  105 } }
}
{ =} { { \frac{   109 }{  105 } }
}
{ } {}
{ } {}
}
{}{.}

\subsection*{Solusi Soal 15.12}
\input{generated/worksheet15_exercise12_solution.id.build.tex}
\subsection*{Solusi Soal 15.13}
\input{generated/worksheet15_exercise13_solution.id.build.tex}

\part{Unit 16}
\chapter[Kuliah 16]{Kuliah 16: Manifold Riemann}
\input{generated/lecture16.id.build.tex}

\chapter{Lembar Kerja 16}
\setcounter{section}{16}

  \zwischenueberschrift{Soal latihan}

\inputaufgabegibtloesung
{}
{

Kita meninjau parabola standar

\mavergleichskette
{\vergleichskette
{ Y
}
{ \subseteq }{ \R^2
}
{  }{
}
{    }{
}
{   }{
}
}
{}{}{}
dengan struktur Riemann terinduksi dan parametrisasi
\maabbeledisp {\varphi} {\R} { Y
} { t } { \left(  t   , \,  t^2                   \right)
} {,}
yang kita pandang sebagai suatu peta
\zusatzklammer {invers} {} {.}
Tentukan
\definitionsverweis {fungsi fundamental}{}{}
Riemann

\mavergleichskette
{\vergleichskette
{ g(t)
}
{ = }{ g_{11} (t)
}
{  }{
}
{    }{
}
{   }{
}
}
{}{}{.}

}
{} {}

\inputaufgabe
{}
{

Misalkan $M$ suatu
\definitionsverweis {manifold diferensiabel}{}{}
sedemikian sehingga
\definitionsverweis {bundel tangen}{}{}
$TM$ trivial. Tunjukkan bahwa melalui suatu trivialisasi diferensiabel

\mavergleichskette
{\vergleichskette
{ TM
}
{ \cong }{ M \times \R^n
}
{  }{
}
{    }{
}
{   }{
}
}
{}{}{}
\zusatzklammer {dengan bantuan
\definitionsverweis {hasil kali dalam standar}{}{}
pada $\R^n$} {} {}
$M$ dapat dijadikan suatu
\definitionsverweis {manifold Riemann}{}{.}

}
{} {}

\inputaufgabe
{}
{

Kita meninjau lingkaran

\mavergleichskette
{\vergleichskette
{ S^1
}
{ \subseteq }{ \R^2
}
{  }{
}
{    }{
}
{   }{
}
}
{}{}{}
dengan
\definitionsverweis {struktur Riemann terinduksi}{}{.}
Trivialisasi
\definitionsverweis {bundel tangen}{}{}
\maabbeledisp {} { S^1 \times \R } { T S^1
} {( a,b, u )} { (a,b, -ub,ua )
} {,}
melalui hasil kali dalam standar pada $\R$, juga menghasilkan suatu struktur Riemann pada $S^1$. Tunjukkan bahwa kedua struktur Riemann tersebut berimpit.

}
{} {}

\inputaufgabe
{}
{

Misalkan
\mathkor {} {L} {dan} {M} {}
merupakan
\definitionsverweis {manifold Riemann}{}{.}
Tunjukkan bahwa $L \times M$ secara alami juga merupakan suatu manifold Riemann.

}
{} {}

\inputaufgabe
{}
{

Kita meninjau suatu himpunan terbuka
\mathl{V \subseteq \R^n}{} sebagai
\definitionsverweis {manifold Riemann}{}{.} Apa
\definitionsverweis {bentuk volume kanonik}{}{} pada $V$?

}
{} {}

\inputaufgabe
{}
{

Misalkan $M$ suatu
\definitionsverweis {manifold Riemann}{}{.}
Tunjukkan bahwa
\definitionsverweis {pemetaan}{}{}
\maabbeledisp {} {
{ \mathcal V } ( M )  } {
{ \mathcal E }^{  1 } (  M   )
} { F } { \omega_F
} {,}
dengan

\mavergleichskettedisp
{\vergleichskette
{ \left(  \omega_F(P)  \right) (v)
}
{ \defeq} { \left\langle F(P) , v  \right\rangle_P
}
{ } {
}
{ } {
}
{ } {
}
}
{}{}{,}
bersifat
\definitionsverweis {linear}{}{.}

}
{} {}

\inputaufgabe
{}
{

Kita meninjau suatu himpunan terbuka
\mathl{V \subseteq \R^n}{} sebagai
\definitionsverweis {manifold Riemann}{}{.} Apa yang dinyatakan oleh korespondensi antara
\definitionsverweis {medan vektor}{}{}
dan
$1$-\definitionsverweis {bentuk diferensial}{}{}
yang diuraikan dalam
Lema 16.3
untuk situasi ini?

}
{} {}

\inputaufgabe
{}
{

Berikan alasan untuk
Catatan 16.4.

}
{} {}

\inputaufgabe
{}
{

Misalkan $M$ suatu
\definitionsverweis {manifold Riemann}{}{}
\definitionsverweis {terorientasi}{}{.} Tunjukkan bahwa
\definitionsverweis {bentuk volume kanonik}{}{} $\omega$ ditentukan oleh syarat bahwa pada setiap titik, nilainya pada setiap basis ortonormal yang merepresentasikan orientasi adalah $1$.

}
{} {}

\inputaufgabe
{}
{

Misalkan $V$ suatu ruang vektor riil terorientasi berdimensi $n$ dan $\lambda$ suatu
\definitionsverweis {ukuran yang invarian terhadap translasi}{}{}
pada $V$. Tunjukkan bahwa pemetaan
\maabbeledisp {} {V \times \cdots \times V} {\R
} {(v_1 , \ldots , v_n) } { \pm \lambda (P(v_1 , \ldots , v_n))
} {,}
dengan tanda positif dipilih jika vektor tersebut merepresentasikan orientasi, merupakan suatu pemetaan multilinear selang-seling.

}
{} {}

\inputaufgabe
{}
{

Tunjukkan bahwa pada suatu
\definitionsverweis {manifold Riemann}{}{},
\definitionsverweis {pemetaan koordinat}{}{}
\maabbdisp {\alpha} { U } { V
} {}
secara umum tidak menginduksi suatu
\definitionsverweis {isometri}{}{}
\maabbdisp {T_P(\alpha)} {T_PU} {T_{\alpha(P)} V
} {}
\zusatzklammer {jika
\mathl{T_PU}{} dilengkapi dengan
\mathl{\left\langle  - ,  -  \right\rangle_P}{} dan
\mathl{T_{\alpha(P)} V= \R^n}{} dilengkapi dengan hasil kali dalam standar} {} {.}

}
{} {}

\inputaufgabegibtloesung
{}
{

Kita meninjau grafik $M$ dari pemetaan
\maabbeledisp {\varphi} {\R^2} {\R
} {(u,v)} { u^2+uv-v^3
} {,} sebagai suatu submanifold tertutup berdimensi dua dari $\R^3$, yaitu

\mavergleichskettedisp
{\vergleichskette
{ M
}
{ =} { \left\{ (u,v,u^2+uv-v^3)  \mid  (u,v) \in \R^2         \right\}
}
{ } {
}
{ } {
}
{ } {
}
}
{}{}{,}
dengan metrik Riemann yang diinduksi dari $\R^3$. Pemetaan
\maabbeledisp {\psi} {\R^2} { M
} {(u,v)} { (u,v,u^2+uv-v^3)
} {,}
merupakan difeomorfisme yang terkait.

a) Tentukan diferensial total dari $\psi$ serta vektor citra
\mathl{T_P(\psi) (e_1)}{} dan
\mathl{T_P(\psi) (e_2)}{} dalam
\mathl{T_{\psi(P)}M}{.}

b) Untuk setiap titik berbentuk

\mavergleichskette
{\vergleichskette
{ P
}
{ = }{ (u,0)
}
{  }{
}
{    }{
}
{   }{
}
}
{}{}{}
tentukan luas jajaran genjang yang direntang oleh
\mathl{T_P(\psi) (e_1)}{} dan
\mathl{T_P(\psi) (e_2)}{} dalam
\mathl{T_{\psi(P)}M}{.}

c) Untuk setiap titik berbentuk

\mavergleichskette
{\vergleichskette
{ P
}
{ = }{ (0,v)
}
{  }{
}
{    }{
}
{   }{
}
}
{}{}{}
tentukan luas jajaran genjang yang direntang oleh
\mathl{T_P(\psi) (e_1)}{} dan
\mathl{T_P(\psi) (e_2)}{} dalam
\mathl{T_{\psi(P)}M}{.}

}
{} {}

\inputaufgabe
{}
{

Misalkan
\maabbdisp {\varphi} {\R^n } {\R
} {}
\zusatzklammer {dengan \mathlk{m=n -1 \geq 0}{}} {} {}
suatu
\definitionsverweis {fungsi yang terdiferensialkan secara kontinu}{}{,}
yang pada setiap titik dari
\definitionsverweis {serat}{}{}
$M$ di atas
\mathl{0 \in \R}{}, bersifat
\definitionsverweis {reguler}{}{.}
Kita memandang $M$ sebagai suatu manifold Riemann terorientasi. Tunjukkan bahwa antara bentuk volume $\tau$ dari
Korolari 15.6
dan
\definitionsverweis {bentuk volume kanonik}{}{}
$\omega$ berlaku hubungan

\mavergleichskettedisp
{\vergleichskette
{\tau(P,v_1 , \ldots , v_m)
}
{ =} { \pm \Vert {
\operatorname{Grad} \, \varphi ( P ) } \Vert  \omega(P, v_1 , \ldots , v_m)
}
{ } {
}
{ } {
}
{ } {
}
}
{}{}{.}

}
{} {}

\inputaufgabe
{}
{

Misalkan
\maabbdisp {\varphi} {\R^n } {\R^\ell
} {}
\zusatzklammer {dengan \mathlk{m=n - \ell \geq 0}{}} {} {}
suatu
\definitionsverweis {pemetaan yang terdiferensialkan secara kontinu}{}{,}
yang pada setiap titik dari
\definitionsverweis {serat}{}{}
$M$ di atas
\mathl{0 \in \R^\ell}{}, bersifat
\definitionsverweis {reguler}{}{.}
Kita memandang $M$ sebagai suatu manifold Riemann terorientasi. Andaikan bahwa gradien

\mathdisp {\operatorname{Grad} \, \varphi_1  (  P ) , \ldots ,
\operatorname{Grad} \, \varphi_\ell ( P )} {  }

saling tegak lurus pada setiap titik
\mathl{P\in M}{.} Tunjukkan bahwa antara bentuk volume $\tau$ dari
Korolari 15.6
dan
\definitionsverweis {bentuk volume kanonik}{}{}
$\omega$ berlaku hubungan

\mavergleichskettedisp
{\vergleichskette
{\tau(P,v_1 , \ldots , v_m)
}
{ =} { \pm \Vert {
\operatorname{Grad} \, \varphi_1  (  P ) } \Vert  \cdots \Vert {
\operatorname{Grad} \, \varphi_\ell ( P ) } \Vert \omega(P, v_1 , \ldots , v_m)
}
{ } {
}
{ } {
}
{ } {
}
}
{}{}{.}

}
{} {}

\inputaufgabe
{}
{

Tunjukkan bahwa
\definitionsverweis {bentuk luas kanonik}{}{}
pada
\definitionsverweis {permukaan bola satuan}{}{}

\mavergleichskette
{\vergleichskette
{S^2
}
{ \subset }{ \R^3
}
{  }{
}
{    }{
}
{   }{
}
}
{}{}{}
adalah

\mathdisp {x dy \wedge dz  -y dx \wedge dz +  z dx \wedge dy} {  }

dengan $x,y,z$ menyatakan koordinat pada $\R^3$.

}
{} {}

Suatu
\definitionsverweis {bundel vektor}{}{}
\maabb {p} { E } { X
} {}
di atas suatu
\definitionsverweis {manifold diferensiabel}{}{}
$X$ disebut
\definitionswort {bundel vektor Riemann}{,}
jika pada setiap serat $E_x$ diberikan suatu
\definitionsverweis {hasil kali dalam}{}{}
dengan sifat bahwa terdapat
\definitionsverweis {peta}{}{}
yang mentrivialkan
\maabbdisp {\alpha} { U } { V \subseteq \R^n
} {}
bersama
\maabbdisp {\theta} {E\vert_U \cong U \times \R^r } { V \times \R^r
} {}
sedemikian sehingga fungsi

\mavergleichskettedisp
{\vergleichskette
{ h_{ij} (Q )
}
{ \defeq} { \left\langle  \theta^{-1}(Q,e_i)  , \theta^{-1}(Q,e_j)    \right\rangle_{\alpha^{-1}(Q)}
}
{ } {
}
{ } {
}
{ } {
}
}
{}{}{}
yang terdiferensialkan secara kontinu.

Jadi, pada suatu
\definitionsverweis {manifold Riemann}{}{},
\definitionsverweis {bundel tangen}{}{}
merupakan suatu bundel vektor Riemann.

\inputaufgabe
{}
{

Misalkan
\maabb {p} { E } { M
} {}
suatu
\definitionsverweis {bundel vektor Riemann}{}{}
di atas suatu
\definitionsverweis {manifold diferensiabel}{}{}
$M$, dan misalkan
\maabb {\varphi} { L } { M
} {}
suatu
\definitionsverweis {pemetaan diferensiabel}{}{}
dengan
\definitionsverweis {bundel vektor tarik balik}{}{}
\maabb {} {\varphi^*E} {L
} {.}
Tunjukkan bahwa $\varphi^*E$ secara alami juga merupakan suatu bundel vektor Riemann.

}
{} {}

  \zwischenueberschrift{Soal untuk dikumpulkan}

Pada suatu manifold Riemann $M$, untuk suatu vektor singgung

\mavergleichskette
{\vergleichskette
{ v
}
{ \in }{ T_PM
}
{  }{
}
{    }{
}
{   }{
}
}
{}{}{}
normanya didefinisikan oleh

\mavergleichskette
{\vergleichskette
{ \Vert {v} \Vert
}
{ = }{ \sqrt{ \left\langle v , v  \right\rangle_P }
}
{  }{
}
{    }{
}
{   }{
}
}
{}{}{.}

\inputaufgabe
{4}
{

Misalkan $M$ suatu
\definitionsverweis {manifold Riemann}{}{.}
Tunjukkan bahwa pemetaan
\maabbeledisp {} {TM} {\R
} { v } { \Vert { v } \Vert
} {,}
bersifat
\definitionsverweis {kontinu}{}{.}

}
{} {}

\inputaufgabe
{3}
{

Tunjukkan bahwa
\mathl{\R \times \R_+}{} dengan bentuk bilinear yang diberikan oleh
\definitionsverweis {bentuk Hesse}{}{} dari
\definitionsverweis {fungsi}{}{}
\maabbeledisp {f} {\R \times \R_+} {\R
} {(x,y)} { x^2+y^4
} {,}
merupakan suatu
\definitionsverweis {manifold Riemann}{}{.}

}
{} {}

\inputaufgabe
{4}
{

Untuk setiap titik
\mathl{P=(x,y,z)}{} pada
\definitionsverweis {permukaan bola satuan}{}{}
$K$, berikan suatu
\definitionsverweis {basis ortonormal}{}{}
dalam
\mathl{T_PK \subset \R^3}{}
\zusatzklammer {terhadap
\definitionsverweis {struktur Riemann terinduksi}{}{}} {} {.}

}
{} {}

\inputaufgabe
{6}
{

Dalam $\R^3$, diberikan
\definitionsverweis {elipsoid}{}{}

\mathdisp {E= \left\{ (x,y,z)  \mid   x^2+y^2+3z^2 \leq 5        \right\}} {  }
dan bidang

\mathdisp {M= \left\{ (x,y,z)  \mid   7x-3y-2z =  2         \right\}} {  }.
Hitunglah luas irisan
\mathl{M \cap E}{.}

}
{} {}

\inputaufgabe
{6}
{

Buatlah suatu grafik komputer yang memvisualisasikan situasi yang diuraikan dalam
Catatan 16.4
dengan menggunakan suatu permukaan di $\R^3$.

}
{} {}


\section*{Solusi yang disediakan oleh sumber}
\addcontentsline{toc}{section}{Solusi yang disediakan oleh sumber}
\subsection*{Solusi Soal 16.1}
\input{generated/worksheet16_exercise01_solution.id.build.tex}
\subsection*{Solusi Soal 16.12}
\input{generated/worksheet16_exercise12_solution.id.build.tex}

\part{Unit 17}
\chapter[Kuliah 17]{Kuliah 17: Perhitungan pada Manifold Riemann}
\setcounter{section}{17}

Dalam kuliah ini kita melanjutkan teori manifold Riemann dan, khususnya, menghitung beberapa luas permukaan.

  \zwischenueberschrift{Perhitungan pada manifold Riemann}



\inputfaktbeweis
{Graph einer Funktion/Riemannsche Untermannigfaltigkeit des R^n/Volumenform/Fakt}
{Korolari}
{}
{

\faktsituation {Misalkan

\mavergleichskette
{\vergleichskette
{W
}
{ \subseteq }{\R^n
}
{  }{
}
{    }{
}
{   }{
}
}
{}{}{}
suatu
\definitionsverweis {himpunan bagian terbuka}{}{}
dan
\maabbdisp {\psi} { W } {\R
} {}
suatu
\definitionsverweis {fungsi terdiferensialkan}{}{.}
Misalkan

\mavergleichskettedisphandlinks
{\vergleichskettedisphandlinks
{M
}
{ =} { \left\{  \left(  x_1 , \ldots , x_n , \, \psi(x_1 , \ldots , x_n) \right)   \mid  (x_1 , \ldots , x_n) \in W         \right\}
}
{ \subseteq} { W \times \R
}
{ } {
}
{ } {
}
}
{}{}{}
merupakan
\definitionsverweis {grafik}{}{}
dari $\psi$.}
\faktfolgerung {Maka $M$ merupakan suatu
\definitionsverweis {manifold Riemann}{}{} yang \definitionsverweis {berorientasi}{}{,}
dan untuk
\definitionsverweis {bentuk volume kanonik}{}{}
$\omega$ pada $M$ berlaku

\mavergleichskettedisp
{\vergleichskette
{ (\operatorname{Id} \times \psi)^*(\omega)
}
{ =} { \left(  1+ \sum_{i = 1}^n \left(  { \frac{   \partial   \psi  }{  \partial   x_i   } }  \right)^2  \right)^{1/2} dx_1 \wedge \ldots \wedge dx_n
}
{ } {
}
{ } {
}
{ } {
}
}
{}{}{.}}
\faktzusatz {}
\faktzusatz {}

}
{

Pemetaan
\maabbeledisp {\varphi = \operatorname{id} \times \psi} { W } { W \times \R
} { x } {(x, \psi(x))
} {,}
merupakan suatu
\definitionsverweis {difeomorfisme}{}{}
antara $W$ dan grafik $M$. Grafik tersebut merupakan suatu
\definitionsverweis {submanifold tertutup}{}{}
dari
\mathl{W \times \R}{} sehingga membawa
\definitionsverweis {struktur Riemann terinduksi}{}{}
dan
\zusatzklammer {karena orientasi $W$ dipindahkan ke $M$} {} {}
suatu
\definitionsverweis {bentuk volume kanonik}{}{}
$\omega$. Pada situasi ini kita dapat menerapkan
Teorema 16.8.
\definitionsverweis {Turunan parsial}{}{}
$\varphi$ terhadap variabel ke-$i$ adalah

\mavergleichskettedisp
{\vergleichskette
{ { \frac{   \partial  \varphi  }{  \partial   x_i   } }
}
{ =} { \begin{pmatrix}  0  \\\vdots\\  0\\ 1\\  0\\ \vdots\\  0\\  { \frac{   \partial   \psi  }{  \partial   x_i   } }   \end{pmatrix}
}
{ } {
}
{ } {
}
{ } {
}
}
{}{}{.}
Misalkan

\mavergleichskette
{\vergleichskette
{Q
}
{ \in }{ W
}
{  }{
}
{    }{
}
{   }{
}
}
{}{}{}
suatu titik yang selanjutnya kita substitusikan ke dalam fungsi-fungsi tersebut, sehingga di seluruh perhitungan kita bekerja dengan bilangan real. Hasil kali dalam yang membentuk entri-entri
\mathl{b_{ij}}{} dari matriks $B$
\zusatzklammer {yang determinannya perlu kita ambil akarnya} {} {}
ialah

\mavergleichskettedisp
{\vergleichskette
{ b_{ij}
}
{ =} { \left\langle  { \frac{   \partial  \varphi  }{  \partial   x_i   } }(Q)  ,  { \frac{   \partial  \varphi  }{  \partial   x_j   } }(Q)   \right\rangle
}
{ } {}
{ } {}
{ } {}
}
{}{}{}
Kita tulis

\mavergleichskette
{\vergleichskette
{B
}
{ = }{ E_n+A
}
{  }{
}
{    }{
}
{   }{
}
}
{}{}{}
dengan

\mavergleichskette
{\vergleichskette
{ A
}
{ = }{ \left(  a_{ij}  \right)_{1 \leq i,j \leq n}
}
{  }{
}
{    }{
}
{   }{
}
}
{}{}{.}
Dengan

\mavergleichskette
{\vergleichskette
{ c_i
}
{ = }{ { \frac{   \partial   \psi  }{  \partial   x_i   } } ( Q)
}
{  }{
}
{    }{
}
{   }{
}
}
{}{}{}
kita dapat menuliskan

\mavergleichskette
{\vergleichskette
{ a_{ij}
}
{ = }{ c_ic_j
}
{  }{
}
{    }{
}
{   }{
}
}
{}{}{}
dan seluruh matriks $A$ sebagai

\mavergleichskettedisp
{\vergleichskette
{A
}
{ =} { \begin{pmatrix} c_1  \\\vdots\\ c_n   \end{pmatrix} \left( c_1   , \,   \ldots  , \,  c_n                 \right)
}
{ } {
}
{ } {
}
{ } {
}
}
{}{}{}
Dengan demikian, $A$ mendeskripsikan suatu pemetaan linear dari $\R^n$ ke $\R^n$ yang
\definitionsverweis {terfaktorkan}{}{} melalui $\R$,
sehingga mempunyai suatu
\definitionsverweis {kernel}{}{}
yang berdimensi sekurang-kurangnya $(n-1)$. Sebut kernel itu $K$. Jika dimensinya $n$, maka

\mavergleichskette
{\vergleichskette
{A
}
{ = }{ 0
}
{  }{
}
{    }{
}
{   }{
}
}
{}{}{}
dan $B$ adalah identitas, sehingga pernyataannya benar. Jadi, misalkan

\mavergleichskette
{\vergleichskette
{A
}
{ \neq }{ 0
}
{  }{
}
{    }{
}
{   }{
}
}
{}{}{.}
Maka

\mavergleichskette
{\vergleichskette
{v
}
{ = }{ \begin{pmatrix} c_1  \\\vdots\\ c_n   \end{pmatrix}
}
{  }{
}
{    }{
}
{   }{
}
}
{}{}{}
merupakan suatu
\definitionsverweis {vektor eigen}{}{}
dari $A$ untuk
\definitionsverweis {nilai eigen}{}{}

\mavergleichskette
{\vergleichskette
{ c_1^2 + \cdots + c_n^2
}
{ \neq }{ 0
}
{  }{
}
{    }{
}
{   }{
}
}
{}{}{.}
Vektor ini merupakan vektor eigen $B$ untuk nilai eigen
\mathl{1+c_1^2 + \cdots + c_n^2}{,} sedangkan $K$ merupakan
\definitionsverweis {ruang eigen}{}{}
bagi $B$ untuk nilai eigen $1$ dan berdimensi
\mathl{(n-1)}{.} Secara keseluruhan, $B$
\definitionsverweis {dapat didiagonalkan}{}{}
dan
\definitionsverweis {determinan}{}{}
matriks tersebut adalah hasil kali nilai-nilai eigennya, yaitu
\mathl{1+c_1^2 + \cdots + c_n^2}{.}

}

Dengan pendekatan ini kita, misalnya, dapat menghitung luas permukaan sfer satuan; lihat
Soal 17.2.



\inputfaktbeweis
{Flächenstück im Raum/Einbettung/Flächenform/Fakt}
{Korolari}
{}
{

\faktsituation {Misalkan

\mavergleichskette
{\vergleichskette
{ M
}
{ \subseteq }{ G
}
{  }{
}
{    }{
}
{   }{
}
}
{}{}{}
suatu
\definitionsverweis {submanifold tertutup terbenam}{}{\zusatzfussnote {Permukaan adalah manifold dua dimensi} {.} {}}
berdimensi dua dan berorientasi di dalam suatu
\definitionsverweis {himpunan terbuka}{}{}

\mavergleichskette
{\vergleichskette
{ G
}
{ \subseteq }{ \R^3
}
{  }{
}
{    }{
}
{   }{
}
}
{}{}{,}
yang dilengkapi dengan struktur Riemann terinduksi dan
\definitionsverweis {bentuk luas kanonik}{}{}
$\omega$. Misalkan

\mavergleichskette
{\vergleichskette
{ W
}
{ \subseteq }{ \R^2
}
{  }{
}
{    }{
}
{   }{
}
}
{}{}{}
terbuka dan misalkan
\maabbdisp {\varphi} { W } { U
} {}
suatu
\definitionsverweis {difeomorfisme}{}{}
yang mempertahankan orientasi ke himpunan terbuka

\mavergleichskette
{\vergleichskette
{ U
}
{ \subseteq }{ M
}
{  }{
}
{    }{
}
{   }{
}
}
{}{}{.}
Koordinat-koordinat pada $\R^2$ adalah
\mathkor {} {u} {dan} {v} {}
dan kita tetapkan
\zusatzfussnote {Notasi ini sudah digunakan oleh Carl Friedrich Gauss} {.} {}

\mathdisp {E = \left\langle  \begin{pmatrix}  { \frac{   \partial  \varphi_1   }{  \partial  u  } }  \\ { \frac{   \partial  \varphi_2   }{  \partial  u  } } \\  { \frac{   \partial  \varphi_3   }{  \partial  u  } }   \end{pmatrix}  ,  \begin{pmatrix}  { \frac{   \partial  \varphi_1   }{  \partial  u  } }  \\ { \frac{   \partial  \varphi_2   }{  \partial  u  } } \\  { \frac{   \partial  \varphi_3   }{  \partial  u  } }   \end{pmatrix}   \right\rangle ,\, F = \left\langle  \begin{pmatrix}  { \frac{   \partial  \varphi_1   }{  \partial  u  } }  \\ { \frac{   \partial  \varphi_2   }{  \partial  u  } } \\  { \frac{   \partial  \varphi_3   }{  \partial  u  } }   \end{pmatrix}   ,  \begin{pmatrix}  { \frac{   \partial  \varphi_1   }{  \partial  v  } }  \\ { \frac{   \partial  \varphi_2   }{  \partial  v  } } \\  { \frac{   \partial  \varphi_3   }{  \partial  v  } }   \end{pmatrix}   \right\rangle \text{ dan } G = \left\langle  \begin{pmatrix}  { \frac{   \partial  \varphi_1   }{  \partial  v  } }  \\ { \frac{   \partial  \varphi_2   }{  \partial  v  } } \\  { \frac{   \partial  \varphi_3   }{  \partial  v  } }   \end{pmatrix}   ,  \begin{pmatrix}  { \frac{   \partial  \varphi_1   }{  \partial  v  } }  \\ { \frac{   \partial  \varphi_2   }{  \partial  v  } } \\  { \frac{   \partial  \varphi_3   }{  \partial  v  } }   \end{pmatrix}   \right\rangle} { . }
}
\faktfolgerung {Maka pada $W$ berlaku

\mavergleichskettealign
{\vergleichskettealign
{ \varphi^*\left(  \omega \vert_U  \right)
}
{ =} { \sqrt{EG-F^2}  du \wedge dv
}
{ } {
}
{ } {
}
{ } {
}
}
{}
{}{.}}
\faktzusatz {}
\faktzusatz {}

}
{

Hal ini langsung mengikuti dari
Teorema 16.8.

}

\inputbemerkung
{}
{

Misalkan

\mavergleichskettedisp
{\vergleichskette
{ M
}
{ =} { \left\{ (u,v, \psi(u,v))  \mid  (u,v) \in W         \right\}
}
{ \subseteq} { W \times \R
}
{ } {
}
{ } {
}
}
{}{}{}
merupakan
\definitionsverweis {grafik}{}{}
dari
\maabbdisp {\psi} { W } {\R
} {,}
dengan

\mavergleichskette
{\vergleichskette
{W
}
{ \subseteq }{ \R^2
}
{  }{
}
{    }{
}
{   }{
}
}
{}{}{}
suatu himpunan bagian terbuka. Dalam kasus ini,
Korolari 17.1
dan
Korolari 17.2
berkaitan sebagai berikut. Turunan-turunan parsialnya adalah
\mathkor {} {\begin{pmatrix}  1  \\ 0 \\  { \frac{   \partial \psi }{  \partial u } }   \end{pmatrix}} {dan} {\begin{pmatrix}  0  \\ 1 \\  { \frac{   \partial \psi }{  \partial v } }   \end{pmatrix}} {.}
Oleh karena itu,

\mavergleichskette
{\vergleichskette
{ E
}
{ = }{ 1 + \left(  { \frac{   \partial \psi }{  \partial u } }  \right)^2
}
{  }{
}
{    }{
}
{   }{
}
}
{}{}{,}

\mavergleichskette
{\vergleichskette
{ F
}
{ = }{ { \frac{   \partial \psi }{  \partial u } } { \frac{   \partial \psi }{  \partial v } }
}
{  }{
}
{    }{
}
{   }{
}
}
{}{}{,}
dan

\mavergleichskette
{\vergleichskette
{ G
}
{ = }{ 1 + \left(  { \frac{   \partial \psi }{  \partial v } }  \right)^2
}
{  }{
}
{    }{
}
{   }{
}
}
{}{}{.}
Dengan demikian,

\mavergleichskettedisphandlinks
{\vergleichskettedisphandlinks
{EG-F^2
}
{ =} { \left(  1 + \left(  { \frac{   \partial \psi }{  \partial u } }  \right)^2  \right) \left(  1 + \left(  { \frac{   \partial \psi }{  \partial v } }  \right)^2  \right) - \left(  { \frac{   \partial \psi }{  \partial u } }  \right)^2 \left(  { \frac{   \partial \psi }{  \partial v } }  \right)^2
}
{ =} { 1  + \left(  { \frac{   \partial \psi }{  \partial u } }  \right)^2  + \left(  { \frac{   \partial \psi }{  \partial v } }  \right)^2
}
{ } {
}
{ } {
}
}
{}{}{.}

}

\inputbemerkung
{}
{

Kita melanjutkan dengan notasi dari
Korolari 17.2.
Jika himpunan bagian terbuka yang dicakup oleh
\mathkor {} {W} {dan} {\varphi} {}

\mavergleichskette
{\vergleichskette
{ U
}
{ \subseteq }{ M
}
{  }{
}
{    }{
}
{   }{
}
}
{}{}{}
mempunyai sifat bahwa
\definitionsverweis {komplemen}{}{}

\mathl{M \setminus U}{} merupakan suatu
\definitionsverweis {himpunan nol}{}{}
terhadap
\definitionsverweis {ukuran kanonik}{}{}
pada $M$, maka luas permukaan $M$ dapat dihitung hanya dengan memakai rumus untuk
\mathl{\varphi^*\left(  \omega \vert_U  \right)}{.} Hal ini, misalnya, berlaku jika
\mathl{M \setminus U}{} merupakan suatu
\definitionsverweis {submanifold tertutup}{}{}
dari $M$ yang berdimensi
\mathl{\leq 1}{;} lihat
Soal 15.14.
Dalam perhitungan, himpunan nol sering diabaikan secara tersirat.

}

  \zwischenueberschrift{Permukaan putar}

Misalkan diberikan suatu
\definitionsverweis {kurva terdiferensialkan}{}{}
\maabbeledisp {\gamma} {]a,b[} {\R^2
} {t} {(x(t),y(t))
} {,}
dengan

\mavergleichskette
{\vergleichskette
{ y(t)
}
{ \geq }{ 0
}
{  }{
}
{    }{
}
{   }{
}
}
{}{}{.}
Kita tertarik pada
\definitionsverweis {permukaan putar}{}{}
yang bersesuaian, yakni himpunan bagian

\mathdisp {\left\{  (x(t), y(t) \cos \alpha , y(t) \sin \alpha)   \mid   t \in ]a,b[  , \,  \alpha \in [0,2\pi]       \right\}} {  }

dari $\R^3$ yang diperoleh dengan memutar lintasan kurva tersebut mengelilingi sumbu $x$. Kita juga mensyaratkan bahwa $\gamma$ menghasilkan suatu difeomorfisme ke citranya

\mavergleichskette
{\vergleichskette
{ M
}
{ = }{ \gamma \left(  ]a,b[  \right)
}
{  }{
}
{    }{
}
{   }{
}
}
{}{}{}
dan bahwa $M$ merupakan submanifold tertutup dalam suatu himpunan terbuka

\mavergleichskettedisp
{\vergleichskette
{ G
}
{ \subseteq} {\R \times \R_+
}
{ } {
}
{ } {
}
{ } {
}
}
{}{}{}

\zusatzklammer {jadi juga disyaratkan bahwa komponen kedua $\gamma$ bernilai positif di mana-mana} {} {.}
Permukaan putarnya kemudian merupakan suatu submanifold tertutup dua dimensi dari $\R^3$ yang tidak berpotongan dengan sumbu $x$, sehingga kita memperoleh suatu manifold Riemann. Luas permukaannya dapat dihitung sebagai berikut.



\inputfaktbeweis
{Rotationsfläche zu ebener Kurve/Flächenberechnung/Fakt}
{Teorema}
{}
{

\faktsituation {Misalkan
\maabbeledisp {\gamma} {]a,b[} {\R^2
} { t } {(x(t),y(t))
} {,}
suatu
\definitionsverweis {kurva terdiferensialkan}{}{}
dengan

\mavergleichskette
{\vergleichskette
{ y(t)
}
{ > }{ 0
}
{  }{
}
{    }{
}
{   }{
}
}
{}{}{,}}
\faktvoraussetzung {yang menginduksi suatu
\definitionsverweis {difeomorfisme}{}{}
ke

\mavergleichskette
{\vergleichskette
{M
}
{ = }{ \gamma(]a,b[)
}
{  }{
}
{    }{
}
{   }{
}
}
{}{}{,}
dengan

\mavergleichskette
{\vergleichskette
{M
}
{ \subseteq }{ G
}
{  }{
}
{    }{
}
{   }{
}
}
{}{}{}
suatu
\definitionsverweis {submanifold tertutup}{}{}
satu dimensi dalam suatu himpunan terbuka

\mavergleichskette
{\vergleichskette
{G
}
{ \subseteq }{\R \times \R_+
}
{  }{
}
{    }{
}
{   }{
}
}
{}{}{.}}
\faktfolgerung {Maka
\definitionsverweis {permukaan putar}{}{}
yang bersesuaian merupakan suatu submanifold tertutup dari
\mathl{\R^3}{} yang tidak berpotongan dengan sumbu $x$, dan luas permukaannya sama dengan

\mathdisp {2 \pi \int_{  a  }^{  b  } \sqrt{ (x'(t))^2+ (y'(t))^2 } \cdot y(t) \, d t} { . }
}
\faktzusatz {}
\faktzusatz {}

}
{

Misalkan $S$ permukaan putar tersebut, yang merupakan submanifold tertutup dua dimensi dalam suatu himpunan terbuka di $\R^3$. Kita menerapkan
Korolari 17.2
pada parametrisasi
\maabbeledisp {} { ]a,b[ \times ]0, 2 \pi[} {S
} { (t, \alpha) } { (x(t),  y(t) \cos \alpha  , y(t) \sin \alpha  )
} {,}
Turunan-turunan parsialnya adalah

\mathdisp {\begin{pmatrix}  x'(t) \\ y'(t) \cos \alpha \\  y'(t) \sin \alpha   \end{pmatrix} \text{   dan } \begin{pmatrix}  0  \\ - y(t) \sin \alpha \\  y(t) \cos \alpha   \end{pmatrix}} {  }

dan karena itu

\mathdisp {E(t, \alpha) = (x'(t))^2 + (y'(t))^2,\, F(t, \alpha)=0,\, G(t, \alpha) = (y(t))^2} { . }

Dengan demikian, luas permukaannya sama dengan

\mavergleichskettedisphandlinks
{\vergleichskettedisphandlinks
{ \int_{  a  }^{  b  } \int_{  0  }^{  2 \pi } \sqrt{(x'(t))^2+(y'(t))^2} \cdot   y(t) \, d \alpha \, d t
}
{ =} { 2 \pi \int_{  a  }^{  b  } \sqrt{(x'(t))^2+(y'(t))^2} \cdot   y(t) \, d t
}
{ } {
}
{ } {
}
{ } {
}
}
{}{}{.}

}

  \zwischenueberschrift{Kartografi}

Kartografi
\zusatzklammer {abstrak} {} {}
mempelajari peta untuk permukaan suatu sfer.

\bild{ \begin{center}
\includegraphics[width=5.5cm]{ \bildeinlesungjpg {Cilinderprojectie-constructie} {jpg} }
\end{center}
\bildtext {} }

\bildlizenz { Cilinderprojectie-constructie.jpg } {} {KoenB} {Commons} {PD} {}

\inputbeispiel{}
{

Kita meninjau pemetaan
\maabbeledisp {\varphi} { [0, 2 \pi] \times [-1,1] } {\R^3
} {(u,v)} { \left(  \sqrt{1-v^2} \cos  u   , \,   \sqrt{1-v^2} \sin  u  , \,   v                 \right)
} {,}
yang
\definitionsverweis {citranya}{}{}
terletak pada \stichwort {sfer satuan} {.} Pemetaan ini kontinu pada persegi panjang tertutup tersebut. Pemetaan ini dapat dibayangkan dengan mula-mula membentuk persegi panjang itu menjadi permukaan tabung, kemudian memproyeksikan lingkaran-lingkaran tabung ke lingkaran-lingkaran horizontal pada sebuah sfer yang mempunyai tinggi sama. Pada interior persegi panjang, pemetaan ini
\definitionsverweis {terdiferensialkan}{}{}
dengan
\definitionsverweis {turunan parsial}{}{}

\mathdisp {{ \frac{   \partial  \varphi  }{  \partial  u  } } = \begin{pmatrix}  - \sqrt{1-v^2} \sin  u  \\ \sqrt{1-v^2} \cos  u \\  0   \end{pmatrix} \text{   dan } { \frac{   \partial  \varphi  }{  \partial  v  } } = \begin{pmatrix}  { \frac{   -v }{  \sqrt{1-v^2}  } } \cos  u   \\ { \frac{   -v }{  \sqrt{1-v^2}  } } \sin  u \\  1   \end{pmatrix}} { . }

\definitionsverweis {Pembatasan}{}{}
pemetaan ini pada persegi panjang terbuka bersifat
\definitionsverweis {injektif}{}{,}
dan citranya adalah sfer satuan kecuali sebuah busur setengah lingkaran bujur. Dengan demikian, kita dapat menghitung luas permukaan sfer memakai koordinat-koordinat ini. Dengan notasi yang digunakan dalam
Korolari 17.2,
kita memperoleh

\mavergleichskettedisp
{\vergleichskette
{E
}
{ =} { \left(  1-v^2  \right) \sin^{ 2  }  u  + \left(  1-v^2  \right) \cos^{ 2  }  u
}
{ =} { 1-v^2
}
{ } {
}
{ } {
}
}
{}{}{,}

\mavergleichskettedisp
{\vergleichskette
{F
}
{ =} { v \sin  u\cos  u  - v \sin  u\cos  u
}
{ =} { 0
}
{ } {
}
{ } {
}
}
{}{}{}
dan

\mavergleichskettedisp
{\vergleichskette
{G
}
{ =} { { \frac{  v^2 }{  1-v^2 } } \cos^{ 2  }  u + { \frac{  v^2 }{  1-v^2 } } \sin^{ 2  }  u +1
}
{ =} { { \frac{  v^2 }{  1-v^2 } }  +1
}
{ =} { { \frac{   1  }{  1-v^2 } }
}
{ } {
}
}
{}{}{.}
Oleh karena itu,

\mavergleichskettedisp
{\vergleichskette
{ \sqrt{EG-F^2 }
}
{ =} { 1
}
{ } {
}
{ } {
}
{ } {
}
}
{}{}{,}
artinya, pemetaan kartografis ini \stichwort {mempertahankan luas} {\zusatzfussnote {Akan tetapi, pemetaan ini tidak mempertahankan panjang. Ruas horizontal pada persegi panjang sangat dipendekkan ketika mendekati kutub. Sebagai gantinya, ruas vertikal semakin diregangkan ketika mendekati kutub, dan kedua gejala ini saling menetralkan} {.} {,}}
sehingga luas permukaan sfer sama dengan

\mavergleichskettealign
{\vergleichskettealign
{A
}
{ =} {
\int_{  ]0, 2 \pi[ \times ]- 1 ,  1 [   }   1  du \wedge dv
}
{ =} { \int_{  [0, 2 \pi] \times [- 1 ,  1 ]   }   1     \,  d  \lambda^2
}
{ =} { 2 \cdot 2 \pi
}
{ =} { 4 \pi
}
}
{}
{}{.}

}

\inputbeispiel{}
{

Kita meninjau pemetaan
\maabbeledisp {\varphi} { [0, 2 \pi] \times \R} {\R^3
} {(u,v)} { { \frac{   1  }{  \sqrt{1+v^2}  } }  ( \cos  u  , \sin  u  , v )
} {,}
yang
\definitionsverweis {citranya}{}{}
terletak pada \stichwort {sfer satuan} {.} Pemetaan ini dapat dibayangkan dengan mula-mula membentuk persegi panjang
\zusatzklammer {yang tak terbatas dalam satu arah} {} {}
\mathl{[0, 2 \pi] \times \R}{} menjadi selimut tabung tak hingga di atas lingkaran satuan, lalu memproyeksikan setiap titik selimut tabung itu ke sfer melalui garis yang menghubungkannya dengan pusat sfer. Di bawah pemetaan ini, semua titik permukaan sfer tercapai kecuali kutub utara dan kutub selatan. Selain itu, pemetaan tersebut bersifat injektif jika titik-titik ujung interval dibuang
\zusatzklammer {maka setengah lingkaran bujur tidak termasuk dalam citra} {} {.}
Pemetaan ini
\definitionsverweis {terdiferensialkan}{}{}
dengan
\definitionsverweis {turunan parsial}{}{}

\mathdisp {{ \frac{   \partial  \varphi  }{  \partial  u  } } = \begin{pmatrix}  { \frac{   - \sin  u  }{  \sqrt{1+v^2}  } }   \\ { \frac{   \cos  u  }{  \sqrt{1+v^2}  } }  \\  0   \end{pmatrix} = { \frac{   1  }{  \sqrt{1+v^2}  } } \begin{pmatrix}  - \sin  u  \\ \cos  u  \\  0   \end{pmatrix} \text{   dan } { \frac{   \partial  \varphi  }{  \partial  v  } } = \begin{pmatrix}  { \frac{   -v \cos  u  }{  \sqrt{1+v^2}^3  } }    \\ { \frac{   -v \sin  u  }{  \sqrt{1+v^2}^3  } }    \\  { \frac{   \sqrt{1+v^2 } - v^2 \left(  1+v^2  \right)^{- 1/2}   }{  1+v^2  } }   \end{pmatrix} = { \frac{   1  }{  \sqrt{1+v^2}^3  } } \begin{pmatrix}  -v \cos  u  \\ -v \sin  u  \\  1   \end{pmatrix}} { . }

Kita dapat menghitung luas permukaan sfer memakai koordinat-koordinat ini. Dengan notasi yang digunakan dalam
Korolari 17.2,
kita memperoleh

\mavergleichskettedisp
{\vergleichskette
{E
}
{ =} { { \frac{   1  }{  1+v^2 } }
}
{ } {
}
{ } {
}
{ } {
}
}
{}{}{,}

\mavergleichskettedisp
{\vergleichskette
{F
}
{ =} { 0
}
{ } {
}
{ } {
}
{ } {
}
}
{}{}{}
dan

\mavergleichskettedisp
{\vergleichskette
{G
}
{ =} { { \frac{   1  }{  \left(  1+v^2  \right)^3 } } \left(  v^2 \cos^{ 2  }  u + v^2 \sin^{ 2  }  u +1  \right)
}
{ =} { { \frac{   1  }{  \left(  1+v^2  \right)^2 } }
}
{ } {
}
{ } {
}
}
{}{}{.}
Oleh karena itu,

\mavergleichskettedisp
{\vergleichskette
{ \sqrt{EG-F^2 }
}
{ =} { \sqrt{ { \frac{   1  }{  \left(  1+v^2  \right)^3  } }  }
}
{ =} { { \frac{   1  }{  \sqrt{1+v^2}^3  } }
}
{ } {
}
{ } {
}
}
{}{}{.}
Dengan menggunakan
Teorema 12.10 (Teori Ukuran dan Integrasi (Osnabrück 2022-2023)),
luas permukaan sfer dengan demikian sama dengan

\mavergleichskettealign
{\vergleichskettealign
{A
}
{ =} {
\int_{  ]0, 2 \pi[ \times \R   }   { \frac{   1  }{  \sqrt{1+v^2}^3  } }  du \wedge dv
}
{ =} { \int_{ ]0, 2 \pi[ \times  \R   }   { \frac{   1  }{  \sqrt{1+v^2}^3  } }    \,  d  \lambda^2
}
{ =} { 2 \pi \int_{  - \infty }^{ + \infty } { \frac{   1  }{  \sqrt{1+v^2}^3  } } \, d  v
}
{ } {
}
}
{}
{}{.}
Menurut
Contoh 31.7 (Analisis (Osnabrück 2021-2023)),
integral tersebut sama dengan $2$, sehingga diperoleh luas permukaan
\mathl{4\pi}{.}

}
\stichwort {Proyeksi Mercator} {} berangkat dari proyeksi yang disebut terakhir dan mereparametrisasi koordinat vertikal sebagai $v=\sinh w,\;w\in\R$. Dengan demikian, koordinat lintang yang terbatas diwakili oleh koordinat vertikal tak terbatas, dan peta yang dihasilkan mempertahankan sudut.

\inputbeispiel{}
{

Kita meninjau pemetaan
\maabbeledisp {\varphi} {\R^2} {\R^3
} {(u,v)} {(\cos  u \cos  v  , \cos  u \sin  v  , \sin  u )
} {,}
yang
\definitionsverweis {citranya}{}{}
terletak pada \stichwort {sfer satuan} {.} Dalam istilah geografi, $u$ menyatakan \stichwort {lingkaran lintang} {} dan $v$ menyatakan \stichwort {lingkaran bujur} {} dari titik yang bersesuaian pada bumi satuan
\zusatzklammer {dalam \stichwort {koordinat geosentrik} {;} koordinat yang dipakai dalam geografi sedikit berbeda karena bumi bukan benar-benar sebuah sfer} {} {.}
Pemetaan ini
\definitionsverweis {terdiferensialkan}{}{}
dengan
\definitionsverweis {turunan parsial}{}{}

\mathdisp {{ \frac{   \partial  \varphi  }{  \partial  u  } } = \begin{pmatrix}  - \sin  u \cos  v  \\ -\sin  u \sin  v \\  \cos  u   \end{pmatrix} \text{   dan } { \frac{   \partial  \varphi  }{  \partial  v  } } = \begin{pmatrix}  - \cos  u \sin  v  \\ \cos  u \cos  v \\  0   \end{pmatrix}} { . }

\definitionsverweis {Pembatasan}{}{}
pemetaan ini pada persegi panjang terbuka

\mathdisp {{]{- { \frac{   \pi }{  2 } }}, { \frac{   \pi }{  2 } } [} \times {]{-\pi}, \pi[}} {  }

bersifat
\definitionsverweis {injektif}{}{,}
dan citranya adalah sfer satuan kecuali satu lingkaran bujur. Dengan demikian, kita dapat menghitung luas permukaan sfer memakai koordinat-koordinat ini. Dengan notasi yang digunakan dalam
Korolari 17.2,
kita memperoleh

\mavergleichskettedisp
{\vergleichskette
{E
}
{ =} { \sin^{ 2  }  u \cos^{ 2  }  v + \sin^{ 2  }  u \sin^{ 2  }  v + \cos^{ 2  }  u
}
{ =} { \sin^{ 2  }  u + \cos^{ 2  }  u
}
{ =} { 1
}
{ } {
}
}
{}{}{,}

\mavergleichskettedisp
{\vergleichskette
{F
}
{ =} { \sin  u\cos  u \sin  v \cos  v  -\sin  u\cos  u \sin  v \cos  v
}
{ =} { 0
}
{ } {
}
{ } {
}
}
{}{}{}
dan

\mavergleichskettedisp
{\vergleichskette
{G
}
{ =} { \cos^{ 2  }  u \sin^{ 2  }  v + \cos^{ 2  }  u \cos^{ 2  }  v
}
{ =} { \cos^{ 2  }  u
}
{ } {
}
{ } {
}
}
{}{}{.}
Oleh karena itu,

\mavergleichskettedisp
{\vergleichskette
{ \sqrt{EG-F^2 }
}
{ =} { \sqrt{ 1 \cdot \cos^{ 2  }  u  }
}
{ =} { \cos  u
}
{ } {
}
{ } {
}
}
{}{}{.}
Jadi, menurut
teorema Fubini,
luas permukaan sfer sama dengan

\mavergleichskettealign
{\vergleichskettealign
{A
}
{ =} {
\int_{  ]- { \frac{   \pi }{  2 } } , { \frac{   \pi }{  2 } } [ \times ]-\pi, \pi[   }   \cos  u  \,du \wedge dv
}
{ =} { \int_{  [- { \frac{   \pi }{  2 } }, { \frac{   \pi }{  2 } }] \times [-\pi, \pi]   }   \cos  u    \,  d  \lambda^2
}
{ =} { \int_{  -\pi }^{  \pi } \int_{  - { \frac{   \pi }{  2 } }  }^{  { \frac{   \pi }{  2 } }  } \cos  u \, d  u \, d v
}
{ =} { \int_{  -\pi }^{  \pi } 2 \, d v
}
}
{
\vergleichskettefortsetzungalign
{ =} { 4 \pi
}
{ } {}
{ } {}
{ } {}
}
{}{.}

}


\chapter{Lembar Kerja 17}
\setcounter{section}{17}

  \zwischenueberschrift{Soal latihan}

\inputaufgabe
{}
{

Misalkan
\maabbeledisp {\varphi} {\R^n } {\R
} { \left(  x_1   , \,   \ldots  , \,   x_n                 \right) } {a_1x_1 + \cdots + a_nx_n
} {,}
suatu
\definitionsverweis {bentuk linear}{}{.}
Misalkan $M$ adalah
\definitionsverweis {grafik}{}{}
fungsi ini, yang kita pandang sebagai
\definitionsverweis {manifold Riemann}{}{.}
Tunjukkan bahwa terdapat hubungan konstan antara volume himpunan-himpunan bagian yang saling bersesuaian di $\R^n$ dan pada grafik tersebut.

}
{} {}

\inputaufgabegibtloesung
{}
{

Hitung luas permukaan bola dengan menggunakan
Korolari 17.1.

}
{} {}

\inputaufgabe
{}
{

Bahas
\definitionsverweis {permukaan putar}{}{} $S$ yang diperoleh dengan memutar

\mathdisp {M= \left\{ ( \sin  { \frac{   1  }{  y } } ,y)  \mid   y> 0        \right\}} {  }

mengelilingi sumbu $x$, yaitu $A$. Apakah $S$ merupakan
\definitionsverweis {submanifold tertutup}{}{} dari
\mathl{\R^3 \setminus A}{?} Apakah himpunan $S$
\definitionsverweis {tertutup}{}{} di $\R^3$? Apakah
\definitionsverweis {penutupan}{}{} $S$ di $\R^3$ merupakan suatu manifold?

}
{} {}

\inputaufgabegibtloesung
{}
{

Tunjukkan bahwa luas permukaan putar yang diperoleh dengan memutar grafik

\mathdisp {\Gamma = \left\{ (x,e^x)  \mid   x \leq 0        \right\}} {  }

mengelilingi sumbu $x$ lebih kecil daripada $10$.

}
{} {}

\inputaufgabe
{}
{

Pastikan bahwa
\definitionsverweis {pemetaan-pemetaan}{}{}
yang diberikan dalam
Contoh 17.6,
Contoh 17.7,
dan
Contoh 17.8
memiliki
\definitionsverweis {citra}{}{}
yang terletak pada
\definitionsverweis {sfer satuan}{}{}
dan bersifat surjektif kecuali pada suatu
\definitionsverweis {himpunan nol}{}{.}

}
{} {}

\inputaufgabe
{}
{

Tentukan
\zusatzklammer {yang terdefinisi secara parsial} {} {}
\definitionsverweis {pemetaan-pemetaan invers}{}{}
dari
\definitionsverweis {pemetaan-pemetaan}{}{}
yang diberikan dalam
Contoh 17.6,
Contoh 17.7,
dan
Contoh 17.8.

}
{} {}

\inputaufgabe
{}
{

Tunjukkan bahwa
\definitionsverweis {lingkaran bujur}{}{}
dan
\definitionsverweis {lingkaran lintang}{}{}
pada bola bumi saling
\definitionsverweis {tegak lurus}{}{.}

}
{} {}

\inputaufgabe
{}
{

Berapakah panjang
\definitionsverweis {lingkaran lintang}{}{} pada lintang $30^\circ$ di Bumi
\zusatzklammer {gunakan $6370$ km sebagai jari-jari bumi} {} {?}

}
{} {}

\inputaufgabe
{}
{

Berapa persen permukaan Bumi yang pada suatu saat disinari Matahari
\zusatzklammer {karena jaraknya sangat jauh, sinar Matahari dianggap saling sejajar} {} {?}

}
{} {}

\inputaufgabe
{}
{

Tinjau elipsoid

\mathdisp {M= \left\{ (x,y,z)\in \R^3  \mid   \ \frac{x^2}{a^2}+\frac{y^2}{b^2}+\frac{z^2}{c^2} =1         \right\}} { , }

untuk $a,b,c\in\mathbb{R}_{+}$. Hitung luas permukaan M untuk
\mathl{a= b}{.} Apa yang terjadi ketika
\mathl{a\rightarrow c}{?}

}
{} {}

\inputaufgabe
{}
{

Tentukan
\definitionsverweis {infimum}{}{}
dan
\definitionsverweis {supremum}{}{}
dari
\definitionsverweis {panjang}{}{}
citra
\definitionsverweis {lingkaran-lingkaran besar}{}{}
pada peta yang diuraikan dalam
Contoh 17.6.

}
{} {}

\inputaufgabe
{}
{

Hubungkan
Korolari 17.1
dengan
Contoh 15.10
dengan menggunakan
Latihan 16.10.

}
{} {}

\inputaufgabe
{}
{

Misalkan
\maabbdisp {\alpha} {S^2 \setminus \{N\}} { \R^2
} {}
adalah
\definitionsverweis {proyeksi stereografis}{}{.}

a) Tentukan
\mathl{\alpha_* \omega}{} untuk
\definitionsverweis {bentuk luas kanonik}{}{}
pada $S^2$.

b) Gunakan
\mathl{\alpha_* \omega}{} untuk menghitung luas permukaan bola.

}
{} {}

  \zwischenueberschrift{Soal untuk dikumpulkan}

\inputaufgabe
{5}
{

Kita meninjau
\definitionsverweis {grafik}{}{}
$M$ dari fungsi
\maabbeledisp {\psi} {\R^2} {\R
} {(x,y)} { y+x^2
} {,}
sebagai
\definitionsverweis {manifold Riemann}{}{.}
Hitung luas permukaan grafik di atas persegi
\mathl{[-1,1] \times [-1,1]}{.}

}
{} {}

\inputaufgabe
{5}
{

Misalkan

\mathdisp {M= \left\{ (x,x^2)  \mid   x \in \R        \right\}  \subseteq \R^2} {  }

adalah
\definitionsverweis {parabola}{}{,}
yaitu
\definitionsverweis {grafik}{}{}
dari
\definitionsverweis {fungsi}{}{}
\maabbeledisp {} {\R} {\R
} { x } { x^2
} {.}
Tunjukkan bahwa
\definitionsverweis {permukaan putar}{}{}
terkait yang diperoleh dengan memutarnya mengelilingi sumbu $x$ bukan suatu
\definitionsverweis {manifold}{}{.}

}
{} {}

\inputaufgabe
{4}
{

Nyatakan suatu
\definitionsverweis {permukaan bola}{}{} sebagai
\definitionsverweis {permukaan putar}{}{} dan gunakan representasi ini untuk menghitung luas permukaan bola tersebut.

}
{} {}

\inputaufgabe
{4}
{

Nyatakan suatu
\definitionsverweis {torus}{}{} sebagai
\definitionsverweis {permukaan putar}{}{} dan gunakan representasi ini untuk menghitung luas permukaannya.

}
{} {}

\inputaufgabe
{6}
{

Tentukan \anfuehrung{jarak}{} antara Osnabrück dan Bangalore
\zusatzklammer {gunakan $6370$ km sebagai jari-jari bumi} {} {}
dalam kedua pengertian berikut.

a) Menyusuri permukaan bumi
\zusatzklammer {jarak garis udara} {} {.}

b) Menembus bumi
\zusatzklammer {jarak garis tikus tanah} {} {.}

}
{} {}

\inputaufgabe
{6}
{

Berapakah panjang citra
\definitionsverweis {lingkaran lintang}{}{} pada lintang $30^\circ$
pada
\definitionsverweis {peta-peta}{}{}
yang diuraikan dalam
Contoh 17.6,
Contoh 17.7,
dan
Contoh 17.8
\zusatzklammer {gunakan $6370$ km sebagai jari-jari bumi} {} {?}

}
{} {}


\section*{Solusi yang disediakan oleh sumber}
\addcontentsline{toc}{section}{Solusi yang disediakan oleh sumber}
\subsection*{Solusi Soal 17.2}
\input{generated/worksheet17_exercise02_solution.id.build.tex}
\subsection*{Solusi Soal 17.4}
\input{generated/worksheet17_exercise04_solution.id.build.tex}

\part{Unit 18}
\chapter[Kuliah 18]{Kuliah 18: Isometri dan Permukaan Hiperbolik}
\setcounter{section}{18}

  \zwischenueberschrift{Rumus transformasi untuk bentuk volume}



\inputfaktbeweis
{Differenzierbare Mannigfaltigkeit/Positive Volumenform/Diffeomorphismus/Rückzug/Maß/Fakt}
{Lema}
{}
{

\faktsituation {Misalkan $M$ suatu
\definitionsverweis {manifold diferensiabel}{}{}
dengan
\definitionsverweis {basis topologi terhitung}{}{}
dan misalkan $\omega$ suatu
\definitionsverweis {bentuk volume positif}{}{}
pada $M$. Misalkan
\maabbdisp {\varphi} { L } { M
} {}
suatu
\definitionsverweis {difeomorfisme}{}{}
dengan manifold $L$, yang diberi orientasi tarikan balik melalui pemetaan tersebut, dan

\mavergleichskette
{\vergleichskette
{ S
}
{ \subseteq }{ L
}
{  }{
}
{    }{
}
{   }{
}
}
{}{}{}
suatu
\definitionsverweis {himpunan bagian terukur}{}{.}}
\faktfolgerung {Maka

\mavergleichskettedisp
{\vergleichskette
{ \int_S \varphi^* \omega
}
{ =} { \int_{ \varphi(S)} \omega
}
{ } {
}
{ } {
}
{ } {
}
}
{}{}{.}}
\faktzusatz {}
\faktzusatz {}

}
{

Menurut
Lema 15.2,
kita dapat mengasumsikan bahwa
\mathkor {} {S} {dan} {\varphi(S)} {}
masing-masing seluruhnya terletak di dalam sebuah peta berorientasi positif pada
\mathkor {} {L} {dan} {M} {;}
katakanlah

\mavergleichskette
{\vergleichskette
{ S
}
{ \subseteq }{ U
}
{  }{
}
{    }{
}
{   }{
}
}
{}{}{}
dengan peta
\maabbdisp {\alpha} { U } { U' \subseteq \R^n
} {}
dan

\mavergleichskette
{\vergleichskette
{ \varphi(S)
}
{ \subseteq }{ V
}
{  }{
}
{    }{
}
{   }{
}
}
{}{}{}
dengan peta
\maabbdisp {\beta} { V } { V' \subseteq \R^n
} {.}
Dengan memperkecil peta-peta itu lebih lanjut, kita dapat mengasumsikan bahwa
\mathkor {} {U} {dan} {V} {}
saling difeomorfik melalui $\varphi$. Dengan demikian, terdapat diagram komutatif difeomorfisme

\mathdisp {\begin{matrix}  U  & \stackrel{ \varphi }{\longrightarrow} &  V  &  \\ \!\!\!\!\! \alpha \downarrow & & \downarrow \beta \!\!\!\!\!  & \\  U'  & \stackrel{ \beta \circ \varphi \circ \alpha^{-1} }{\longrightarrow} &  V' & \! \! \! \!\!\!\!\!   \\ \end{matrix}} {  }

yang menginduksi diagram komutatif himpunan terukur

\mathdisp {\begin{matrix}  S  & \stackrel{ \varphi }{\longrightarrow} &  \varphi (S)  &  \\ \!\!\!\!\! \alpha \downarrow & & \downarrow \beta \!\!\!\!\!  & \\  \alpha (S) & \stackrel{ \beta \circ \varphi \circ \alpha^{-1} }{\longrightarrow} &  \beta ( \varphi (S) ) & \! \! \! \!\!\!\!\!   \\ \end{matrix}} {  }

Bagi bentuk volume $\omega$ pada $V$, berlaku

\mavergleichskettedisp
{\vergleichskette
{ \alpha^{ {-1}*} ( \varphi^* \omega)
}
{ =} { \left(  \beta \circ \varphi \circ \alpha^{-1}  \right)^*  ( \beta^{ {-1}*} \omega)
}
{ } {
}
{ } {
}
{ } {
}
}
{}{}{.}
menurut
Lema 14.7 (4).
Bentuk volume $\beta^{ {-1}*} \omega$ pada $V'$ mempunyai representasi
\mathl{g dy_1  \wedge \ldots \wedge dy_n}{} dengan suatu fungsi terukur
\maabb {g} {V'} {\R
} {}
dan, menurut definisi, $\int_{\varphi(S)} \omega$ sama dengan $\int_{\beta( \varphi(S))} g d \lambda^n$. Demikian pula,
\mathl{\alpha^{ {-1}*} ( \varphi^* \omega)}{} pada $U'$ mempunyai representasi berbentuk
\mathl{f dx_1 \wedge \ldots \wedge dx_n}{.} Bentuk volume pada $U'$ ini berimpit dengan penarikan balik
\mathl{g dy_1  \wedge \ldots \wedge dy_n}{} dan, menurut
Korolari 14.9,
penarikan balik tersebut sama dengan

\mathdisp {(g \circ \psi) \cdot \det  \left(  \left(  { \frac{   \partial   \psi_i   }{  \partial   x_j   } }  \right)_{1 \leq  i, j \leq n}  \right) dx_1 \wedge \ldots \wedge dx_n} {  }

\zusatzklammer {dengan

\mavergleichskettek
{\vergleichskettek
{ \psi
}
{ = }{ \beta \circ \varphi \circ \alpha^{-1}
}
{  }{
}
{    }{
}
{   }{
}
}
{}{}{}} {} {.}
Dengan demikian, pernyataan tersebut mengikuti dari
Teorema 14.3 (Teori Ukuran dan Integrasi (Osnabrück 2022-2023)).

}

  \zwischenueberschrift{Isometri antara manifold Riemann}

\inputdefinition
{}
{

Misalkan
\mathkor {} {L} {dan} {M} {}
merupakan
\definitionsverweis {manifold Riemann}{}{.}
Suatu
\definitionsverweis {pemetaan terdiferensialkan}{}{}
\maabb {\varphi} { L } { M
} {}
disebut
\definitionswort {isometri lokal}{,}
jika untuk setiap titik

\mavergleichskette
{\vergleichskette
{ P
}
{ \in }{ L
}
{  }{
}
{    }{
}
{   }{
}
}
{}{}{}
\definitionsverweis {pemetaan singgung}{}{}
\maabbdisp {T_P \varphi} {T_PL} { T_{\varphi(P) } M
} {}
merupakan suatu
\definitionsverweis {isometri}{}{}
terhadap hasil kali dalam yang diberikan.

}

Pemetaan
\maabbeledisp {} {\R} { \R^2
} { t } { \begin{pmatrix}  \cos  t  \\ \sin  t    \end{pmatrix}
} {,}
sudah menunjukkan bahwa suatu isometri lokal tidak harus injektif.

\inputdefinition
{}
{

Misalkan
\mathkor {} {L} {dan} {M} {}
merupakan
\definitionsverweis {manifold Riemann}{}{.}
Suatu
\definitionsverweis {pemetaan terdiferensialkan}{}{}
\maabb {\varphi} { L } { M
} {}
disebut
\definitionswort {isometri}{,}
jika pemetaan tersebut merupakan suatu
\definitionsverweis {difeomorfisme}{}{}
dan secara lokal merupakan isometri.

}

Manifold-manifold Riemann yang isometrik dianggap sama dalam geometri Riemann. Secara khusus, panjang, sudut, dan volume dipertahankan oleh isometri; lihat
Lema 18.6.
Akan tetapi, sama sekali tidak mudah menentukan sifat mana dari suatu manifold Riemann yang tertanam di dalam $\R^n$ yang hanya bergantung pada struktur Riemann dan bukan pada penanamannya. Sifat yang hanya bergantung pada struktur Riemann juga disebut intrinsik
\zusatzklammer {sedangkan sifat yang bergantung pada penanaman disebut ekstrinsik} {} {.}
Untuk sifat intrinsik, gambaran yang sesuai ialah menanyakan apakah sifat itu dapat dikenali dan dideskripsikan oleh suatu makhluk yang hanya boleh bergerak pada manifold tersebut dan tidak dapat meninggalkannya.

\inputbeispiel{}
{

Misalkan

\mavergleichskette
{\vergleichskette
{ M
}
{ \subseteq }{ \R^n
}
{  }{
}
{    }{
}
{   }{
}
}
{}{}{}
suatu submanifold Riemann tertutup. Maka setiap
\definitionsverweis {isometri linear}{}{}
dari $\R^n$ menginduksi suatu
\definitionsverweis {isometri}{}{}
dari $M$ ke $\varphi(M)$; simbol pada kodomain menyatakan isometri linear yang bersangkutan.

}

\inputbeispiel{}
{

Misalkan
\maabb {\gamma} {]a,b[} { \R^2
} {}
suatu kurva injektif
\definitionsverweis {terdiferensialkan}{}{}
dan
\definitionsverweis {berparametrisasi panjang busur}{}{}
dengan kurva citra

\mavergleichskettedisp
{\vergleichskette
{ C
}
{ \subseteq} { V
}
{ \subseteq} { \R^2
}
{ } {
}
{ } {
}
}
{}{}{,}
yang tertutup di dalam himpunan bagian terbuka $V$. Maka pemetaan
\maabbdisp {\varphi = \gamma \times
\operatorname{Id}_{ \R }} { ]a,b[ \times \R} { C \times \R
} {}
merupakan suatu
\definitionsverweis {isometri}{}{,}
dengan
\mathl{]a,b[ \times \R}{} membawa struktur Euklides alami dan $C \times \R$ membawa struktur submanifold Riemann

\mavergleichskette
{\vergleichskette
{ C \times \R
}
{ \subseteq }{ V \times \R
}
{ \subseteq }{ \R^3
}
{    }{
}
{   }{
}
}
{}{}{}
yang terinduksi. Pada dasarnya, pemetaan ini berarti bahwa selembar kertas dibentangkan mengikuti sebuah kurva dasar. Serat-serat vertikal tidak berubah, sedangkan serat-serat horizontal merupakan salinan kurva dasar. Contoh tak trivial yang paling sederhana ialah membengkokkan lembar tersebut menjadi selimut tabung bercelah. Secara intuitif jelas bahwa panjang kurva pada permukaan dan luas permukaan tidak berubah. Kita dapat menunjukkan bahwa pemetaan ini merupakan isometri sebagai berikut. Untuk itu, misalkan

\mavergleichskette
{\vergleichskette
{ P
}
{ = }{ (s,t)
}
{ \in }{ ]a,b[ \times \R
}
{    }{
}
{   }{
}
}
{}{}{.}
Maka

\mavergleichskettedisp
{\vergleichskette
{ \left(  D\varphi  \right)_{ P } \left(  e_1   \right)
}
{ =} { \begin{pmatrix}  \gamma_1'(s) \\ \gamma_2'(s)\\  0   \end{pmatrix}
}
{ } {
}
{ } {
}
{ } {
}
}
{}{}{}
dan

\mavergleichskettedisp
{\vergleichskette
{ \left(  D\varphi  \right)_{ P } \left(  e_2   \right)
}
{ =} { \begin{pmatrix}  0  \\ 0\\  1   \end{pmatrix}
}
{ } {
}
{ } {
}
{ } {
}
}
{}{}{,}
jadi kedua vektor citra dari vektor-vektor standar mempunyai panjang $1$ dan saling tegak lurus. Dengan demikian,
\maabbdisp {} {T_P(  ]a,b[ \times \R)} { T_{\varphi(P)} (C \times \R)
} {}
merupakan suatu isometri.

}



\inputfaktbeweis
{Riemannsche Mannigfaltigkeit/Isometrie/Grundlegende Eigenschaften/Fakt}
{Lema}
{}
{

\faktsituation {Misalkan $L,M$ merupakan
\definitionsverweis {manifold Riemann}{}{}
yang
\definitionsverweis {berorientasi}{}{}
dan misalkan
\maabb {\varphi} { L } { M
} {}
suatu
\definitionsverweis {isometri}{}{}
yang
\definitionsverweis {mempertahankan orientasi}{}{.}}
\faktuebergang {Maka pernyataan-pernyataan berikut berlaku.}
\faktfolgerung {\aufzaehlungvier{Bentuk volume kanonik
\definitionsverweis {ditarik balik}{}{}
dari
\definitionsverweis {bentuk volume kanonik}{}{}
pada $M$ menjadi bentuk volume kanonik pada $L$.
}{Untuk setiap kurva terdiferensialkan
\maabbdisp {\gamma} {[a,b]} { L
} {}
\definitionsverweis {panjang kurva}{}{}
$\gamma$ sama dengan panjang kurva $\varphi \circ \gamma$.
}{ $\varphi$
\definitionsverweis {mempertahankan ukuran}{}{.}
}{ $\varphi$
\definitionsverweis {konformal}{}{.}
}}
\faktzusatz {}
\faktzusatz {}

}
{

\aufzaehlungvier{Misalkan $\omega$ bentuk volume kanonik pada $M$ dan

\mavergleichskette
{\vergleichskette
{ P
}
{ \in }{ L
}
{  }{
}
{    }{
}
{   }{
}
}
{}{}{.}
Misalkan

\mavergleichskette
{\vergleichskette
{ u_1 , \ldots , u_n
}
{ \in }{ T_PL
}
{  }{
}
{    }{
}
{   }{
}
}
{}{}{}
suatu
\definitionsverweis {basis ortonormal}{}{}
dari $T_PL$ yang merepresentasikan orientasi $L$. Maka

\mavergleichskettedisp
{\vergleichskette
{ \left(  \varphi^* \omega  \right) (P, u_1 , \ldots , u_n)
}
{ =} { \omega( \varphi(P), T_P\varphi(u_1) , \ldots , T_P\varphi(u_n))
}
{ } {
}
{ } {
}
{ } {
}
}
{}{}{.}
Karena sifat isometri,
\mathl{T_P\varphi(u_1) , \ldots , T_P\varphi(u_n)}{} merupakan suatu basis ortonormal dari $T_{\varphi(P)} M$. Karena pemetaan tersebut mempertahankan orientasi, basis itu merepresentasikan orientasi pada $M$, sehingga nilainya sama dengan $1$. Sifat ini mengarakterisasi bentuk volume kanonik pada $L$.
}{Panjang kurva $\gamma$ adalah integral dari
\mathl{\Vert { \gamma'(t) } \Vert}{,} dengan norma diambil dalam
\mathl{T_{\gamma(t)}L}{.} Demikian pula, panjang kurva $\varphi \circ \gamma$ adalah integral dari
\mathl{\Vert { \left(  \varphi \circ \gamma  \right)'(t) } \Vert}{,} yang sama karena sifat isometri.
}{Untuk suatu himpunan bagian

\mavergleichskette
{\vergleichskette
{ S
}
{ \subseteq }{ L
}
{  }{
}
{    }{
}
{   }{
}
}
{}{}{}
ukuran yang berasal dari bentuk volume kanonik pada $L$, jika diterapkan pada $S$, menurut bagian (1) dan
Lema 18.1,
sama dengan

\mavergleichskettedisp
{\vergleichskette
{ \int_S \omega_L
}
{ =} { \int_S \varphi^*(\omega_M)
}
{ =} { \int_{\varphi(S)} \omega_M
}
{ } {
}
{ } {
}
}
{}{}{.}
}{Hal ini langsung mengikuti dari sifat isometri.
}

}



\inputfaktbeweis
{Riemannsche Mannigfaltigkeit/Karte/Isometrie/Fakt}
{Lema}
{}
{

\faktsituation {Misalkan $M$ suatu
\definitionsverweis {manifold Riemann}{}{}
dan misalkan
\maabbdisp {\alpha} { U } { V \subseteq \R^n
} {}
suatu
\definitionsverweis {peta}{}{}
pada $M$.}
\faktfolgerung {Maka $\alpha$ merupakan suatu
\definitionsverweis {isometri}{}{}
antara
\mathkor {} {U} {dan} {V} {,}
jika $V$ dilengkapi dengan
\definitionsverweis {hasil kali dalam}{}{}
yang ditentukan oleh
\definitionsverweis {fungsi-fungsi fundamental metrik}{}{}
$g_{ij}$.}
\faktzusatz {}
\faktzusatz {}

}
{

Hal ini langsung mengikuti dari definisi $g_{ij}$ melalui

\mavergleichskettedisp
{\vergleichskette
{ g_{ij}(Q)
}
{ =} { \left\langle  T(\alpha^{-1})(e_i)  ,  T(\alpha^{-1}) (e_j)    \right\rangle_{ \alpha^{-1}(Q) }
}
{ } {
}
{ } {
}
{ } {
}
}
{}{}{.}

}

  \zwischenueberschrift{Permukaan hiperbolik}

Kita membahas tiga model untuk apa yang disebut permukaan hiperbolik
\zusatzklammer {atau bidang hiperbolik} {} {}
dan memberikan
\definitionsverweis {isometri}{}{}
di antara model-model tersebut. Dalam
Contoh 29.10,
kita akan melihat bahwa ini merupakan permukaan dengan
\definitionsverweis {kelengkungan seksional}{}{}
konstan $-1$. Model berikut disebut setengah bidang Poincaré.

\inputbeispiel{}
{

Misalkan

\mavergleichskettedisp
{\vergleichskette
{ H = {\mathbb H}
}
{ =} { \R \times \R_+
}
{ =} { \left\{ (x,y)  \mid   y > 0        \right\}
}
{ } {
}
{ } {
}
}
{}{}{}
merupakan setengah bidang atas yang dilengkapi dengan
\definitionsverweis {metrik Riemann}{}{,}
yang diberikan oleh

\mavergleichskettedisp
{\vergleichskette
{ g_{11}
}
{ =} { g_{22}
}
{ =} { { \frac{   1  }{  y^2 } }
}
{ } {
}
{ } {
}
}
{}{}{}
dan

\mavergleichskette
{\vergleichskette
{ g_{12}
}
{ = }{ 0
}
{  }{
}
{    }{
}
{   }{
}
}
{}{}{.}
Jadi, hasil kali dalam standar diskalakan dengan fungsi ${ \frac{   1  }{  y^2 } }$. Norma suatu vektor singgung pada titik
\mathl{(x,y)}{} adalah norma standar yang diskalakan dengan faktor ${ \frac{   1  }{  \betrag {  y  } } }$. Dengan demikian, bentuk luas diberikan oleh kerapatan ${ \frac{   1  }{  y^2 } }$.

}

Model berikut disebut \stichwort {cakram hiperbolik} {.}

\inputbeispiel{}
{

Pada cakram terbuka

\mavergleichskettedisp
{\vergleichskette
{ V
}
{ =} { \left\{ (x,y)  \mid   x^2+y^2 < 1        \right\}
}
{ } {
}
{ } {
}
{ } {
}
}
{}{}{}
kita mendefinisikan suatu
\definitionsverweis {struktur Riemann}{}{}
melalui bentuk bilinear

\mavergleichskettedisp
{\vergleichskette
{ g(Q) (u,v)
}
{ =} { { \frac{   4  }{  \left(  1-  \Vert { Q } \Vert^2   \right)^2  } }   \left\langle  u  , v  \right\rangle
}
{ } {
}
{ } {
}
{ } {
}
}
{}{}{,}
jadi hasil kali dalam standar diskalakan dengan fungsi
\mathl{{ \frac{   4  }{  \left(  1 - \Vert { Q } \Vert^2   \right)^2  } }}{,} dan matriks fundamental pertamanya adalah

\mathdisp {{ \frac{   4  }{  \left(  1- \Vert { Q } \Vert^2   \right)^2  } } \begin{pmatrix}  1   &   0   \\  0 &  1 \end{pmatrix}} { . }

}

\bild{ \begin{center}
\includegraphics[width=5.5cm]{ \bildeinlesunggif {Poincarehalfplaneconform} {gif} }
\end{center}
\bildtext {Animasi transformasi konformal setengah bidang Poincaré menjadi cakram satuan melalui inversi terhadap lingkaran yang berpusat di titik bergerak I. Edisi PDF menampilkan bingkai pertama yang statis; animasi asli dipertahankan untuk edisi HTML dan unduhan.} }

\bildlizenz { Poincarehalfplaneconform.gif } {} {HB} {Commons} {CC-by-sa 3.0} {}



\inputfaktbeweis
{Hyperbolische Fläche/Kreisscheibe/Halbebene/Isometrie/Fakt}
{Lema}
{}
{

\faktsituation {Pemetaan
\maabbeledisp {\psi} { V } { H
} { w } { { \frac{   \left(  w +1  \right)  { \mathrm i}  }{  1 - w } }
} {,}}
\faktfolgerung {mendefinisikan suatu
\definitionsverweis {isometri}{}{}
antara
\definitionsverweis {cakram hiperbolik}{}{}
dan
\definitionsverweis {setengah bidang Poincaré}{}{.}}
\faktzusatz {}
\faktzusatz {}

}
{

Menurut
Soal 18.10,
pemetaan ini bijektif dan terdiferensialkan kompleks dengan invers yang terdiferensialkan kompleks, sehingga khususnya terdapat suatu difeomorfisme antara manifold-manifold diferensiabel. Menurut
Soal 18.12,
pemetaan ini dalam koordinat real diberikan oleh

\mavergleichskettedisp
{\vergleichskette
{ \psi \begin{pmatrix}  a  \\b   \end{pmatrix}
}
{ =} { { \frac{   1  }{  -2a+1+a^2+b^2 } }  \begin{pmatrix}  -2b \\ 1-a^2-b^2   \end{pmatrix}
}
{ } {
}
{ } {
}
{ } {
}
}
{}{}{}
Menurut
Soal 18.13,
turunan-turunan parsialnya adalah

\mavergleichskettedisp
{\vergleichskette
{ { \frac{   \partial   \psi  }{  \partial  a  } }
}
{ =} { { \frac{   1  }{  \left(  -2a+1+a^2+b^2  \right)^2  } } \begin{pmatrix}  4ab -4b  \\ 2a^2 -4a+2 -2 b^2     \end{pmatrix}
}
{ } {
}
{ } {}
{ } {}
}
{}{}{}
dan

\mavergleichskettedisp
{\vergleichskette
{ { \frac{   \partial   \psi  }{  \partial  b  } }
}
{ =} { { \frac{   1  }{  \left(  -2a+1+a^2+b^2  \right)^2  } } \begin{pmatrix}  4a  -2-2a^2+2b^2   \\ 4ab -4b     \end{pmatrix}
}
{ } {
}
{ } {
}
{ } {}
}
{}{}{.}
Kita harus menunjukkan bahwa

\mavergleichskettedisp
{\vergleichskette
{ \left\langle e_i  , e_j   \right\rangle_{V, (a,b)}
}
{ =} { \left\langle  \left(  D \psi  \right)_{(a,b)} \left(  e_i   \right)   ,  \left(  D \psi  \right)_{(a,b)} \left(  e_j   \right)   \right\rangle_{H, \psi (a,b)}
}
{ } {
}
{ } {
}
{ } {
}
}
{}{}{.}
Karena kedua turunan parsial saling tegak lurus dan karena kedua metrik Riemann mewarisi ortogonalitas hasil kali dalam standar, hal ini hanya perlu ditunjukkan untuk

\mavergleichskette
{\vergleichskette
{ i
}
{ = }{ j
}
{  }{
}
{    }{
}
{   }{
}
}
{}{}{.}
Jadi, kita harus menunjukkan bahwa diferensial total mempertahankan panjang. Kita memperoleh

\mavergleichskettealigndrucklinks
{\vergleichskettealigndrucklinks
{ \Vert { { \frac{   \partial   \psi  }{  \partial   a   } } ( \begin{pmatrix}  a  \\b   \end{pmatrix} ) } \Vert_{H, \psi(a,b) }
}
{ =} { \Vert { { \frac{   1  }{  \left(  -2a+1+a^2+b^2  \right)^2  } } \begin{pmatrix}  4ab -4b  \\ 2a^2 -4a+2 -2 b^2    \end{pmatrix} } \Vert_{H, \psi(a,b) }
}
{ =} { { \frac{   \left(  -2a+1+a^2+b^2  \right)  }{  \left(  1-a^2-b^2  \right)  } } \Vert { { \frac{   2  }{  \left(  -2a+1+a^2+b^2  \right)^2  } } \begin{pmatrix}  2(a-1)b  \\ \left(  a-1  \right)^2 - b^2    \end{pmatrix} } \Vert
}
{ =} { { \frac{   2  }{  \left(  1-a^2-b^2  \right) \left(  \left(  a-1  \right)^2 +b^2  \right)  } } \sqrt{ 4(a-1)^2b^2 + \left(  a-1  \right)^4 + b^4 -2\left(  a-1  \right)^2  b^2 }
}
{ =} { { \frac{   2  }{  \left(  1-a^2-b^2  \right) \left(  \left(  a-1  \right)^2 +b^2  \right)  } } \left(  \left(  a-1  \right)^2 + b^2   \right)
}
}
{
\vergleichskettefortsetzungalign
{ =} { { \frac{   2  }{  \left(  1- \left(  a^2+b^2  \right)  \right)  } }
}
{ =} { \Vert {e_1 } \Vert_{V, (a,b)}
}
{ } {}
{ } {}
}{}{.}
Dengan cara yang sama, pernyataan tersebut diperoleh untuk turunan parsial kedua.

}

Metrik Euklides standar pada $\R^n$ menginduksi suatu metrik Riemann pada setiap submanifold tertutup. Akan tetapi, terdapat pula situasi ketika $\R^n$ dilengkapi dengan suatu
\definitionsverweis {bentuk bilinear}{}{}
tak terdegenerasi yang berbeda, yaitu tidak
\definitionsverweis {definit positif}{}{,}
namun pembatasannya pada suatu submanifold tertutup tetap menghasilkan manifold Riemann. Untuk contoh berikut, yaitu hiperboloid dua lembar, lihat juga Kurs:Lineare Algebra (Osnabrück 2017-2018)/Teil II/Vorlesung 40.

\inputbeispiel{}
{

\bild{ \begin{center}
\includegraphics[width=5.5cm]{ \bildeinlesungpng {Hyperboloid2} {png} }
\end{center}
\bildtext {Yang dibahas ialah lembar atas hiperboloid dua lembar,\\
yang dilengkapi dengan bentuk bilinear terinduksi dari bentuk Minkowski.} }

\bildlizenz { Hyperboloid2.png } {} {RokerHRO} {Commons} {gemeinfrei} {}

Kita melengkapi $\R^3$ dengan
\definitionsverweis {bentuk Minkowski standar}{}{,}
yakni bentuk bilinear yang bersesuaian dengan bentuk kuadratik $x^2+y^2-z^2$, dan kita meninjau himpunan bagian

\mavergleichskettedisp
{\vergleichskette
{ Y
}
{ =} { \left\{ (x,y,z)  \mid   x^2 +y^2-z^2 = -1  , \,  z > 0       \right\}
}
{ } {
}
{ } {
}
{ } {
}
}
{}{}{.}
Jika $\R^3$ dengan bentuk Minkowski dipandang sebagai model tiga dimensi untuk relativitas khusus, maka ini merupakan himpunan vektor pengamat berarah masa depan. Untuk suatu titik

\mavergleichskette
{\vergleichskette
{ P
}
{ = }{ (x,y,z)
}
{ \in }{ Y
}
{    }{
}
{   }{
}
}
{}{}{}
menurut
Lema 40.4 (Aljabar Linear (Osnabrück 2024-2025)),
pembatasan bentuk tersebut pada ruang singgung bersifat definit positif. Vektor-vektor
\mathkor {} {\begin{pmatrix}  z  \\ 0\\  x   \end{pmatrix}} {dan} {\begin{pmatrix}  0  \\ z \\  y   \end{pmatrix}} {}
membentuk suatu basis dari
\definitionsverweis {ruang singgung}{}{}
$T_PY$. Untuk sembarang vektor singgung

\mavergleichskettedisp
{\vergleichskette
{ v
}
{ =} { a \begin{pmatrix}  z  \\ 0\\  x   \end{pmatrix} + b \begin{pmatrix}  0  \\ z\\  y   \end{pmatrix}
}
{ =} { \begin{pmatrix} az \\ bz\\ ax+by   \end{pmatrix}
}
{ } {
}
{ } {
}
}
{}{}{}
berlaku

\mavergleichskettealignhandlinks
{\vergleichskettealignhandlinks
{ \left\langle  v  , v  \right\rangle_{Y,P}
}
{ =} { {  \begin{pmatrix} az \\ bz\\ ax+by \end{pmatrix} ^{ \text{tr} } } \begin{pmatrix}  1   &   0  &  0 \\  0  &  1  &  0  \\ 0 &  0  &  -1 \end{pmatrix} \begin{pmatrix} az \\ bz\\ ax+by \end{pmatrix}
}
{ =} { {  \begin{pmatrix} az \\ bz\\ ax+by \end{pmatrix} ^{ \text{tr} } }\begin{pmatrix} az \\ bz\\  -ax-by \end{pmatrix}
}
{ =} { a^2z^2+b^2z^2-\left(  ax+by  \right)^2
}
{ =} { a^2+b^2+\left(  ay-bx  \right)^2
}
}
{}
{}{,}
yang positif setiap kali kedua koefisien tidak sekaligus nol.

}



\inputfaktbeweis
{Hyperbolische Fläche/Kreisscheibe/Zweischaliges Hyperboloid/Isometrie/Fakt}
{Lema}
{}
{

\faktsituation {Pemetaan
\maabbeledisp {\theta} { V } { Y
} { \begin{pmatrix}  a  \\b   \end{pmatrix} } { \begin{pmatrix}  { \frac{   -2a }{ a^2+b^2-1  } }  \\ { \frac{   -2b }{ a^2+b^2-1  } } \\  { \frac{   -a^2-b^2-1  }{ a^2+b^2-1  } }   \end{pmatrix}
} {,}}
\faktfolgerung {mendefinisikan suatu
\definitionsverweis {isometri}{}{}
antara
\definitionsverweis {cakram hiperbolik}{}{}
$V$ dan lembar atas
\definitionsverweis {hiperboloid dua lembar}{}{}
dengan
\definitionsverweis {metrik Minkowski}{}{}
terinduksi.}
\faktzusatz {}
\faktzusatz {}

}
{

Turunan parsial dari $\theta$ adalah

\mavergleichskettealignhandlinks
{\vergleichskettealignhandlinks
{ { \frac{   \partial  \theta  }{  \partial   a   } } ( \begin{pmatrix}  a  \\b   \end{pmatrix} )
}
{ =} { { \frac{   1  }{  \left(  a^2+b^2-1  \right)^2  } } \begin{pmatrix}  -2 \left(  a^2+b^2-1  \right)+ 2a (2a)  \\ 4ab\\  -2a \left(  a^2+b^2-1  \right) +2a  \left(  a^2+b^2+1  \right)   \end{pmatrix}
}
{ =} { { \frac{   1  }{  \left(  a^2+b^2-1  \right)^2  } } \begin{pmatrix}  2a^2-2b^2+2  \\ 4ab\\  4a   \end{pmatrix}
}
{ } {
}
{ } {
}
}
{}
{}{}
dan

\mavergleichskettedisp
{\vergleichskette
{ { \frac{   \partial  \theta  }{  \partial   b   } } ( \begin{pmatrix}  a  \\b   \end{pmatrix} )
}
{ =} { { \frac{   1  }{  \left(  a^2+b^2-1  \right)^2  } } \begin{pmatrix}  4ab \\ - 2a^2+2b^2+2 \\  4b   \end{pmatrix}
}
{ } {
}
{ } {
}
{ } {}
}
{}{}{.}
Kita memperoleh

\mavergleichskettealignhandlinks
{\vergleichskettealignhandlinks
{ \left\langle  \begin{pmatrix}  2a^2-2b^2+2  \\ 4ab\\  4a   \end{pmatrix}   ,  \begin{pmatrix}  4ab \\ - 2a^2+2b^2+2 \\  4b   \end{pmatrix}     \right\rangle_{Y}
}
{ =} { 4 ab \left(  2a^2-2b^2+2   \right) +  4 ab \left(  - 2a^2 +2b^2+2   \right)      -16ab
}
{ =} { 0
}
{ } {
}
{ } {
}
}
{}
{}{,}

\mavergleichskettealignhandlinks
{\vergleichskettealignhandlinks
{ \left\langle  \begin{pmatrix}  2a^2-2b^2+2  \\ 4ab\\  4a   \end{pmatrix}   ,  \begin{pmatrix}  2a^2-2b^2+2  \\ 4ab\\  4a   \end{pmatrix}      \right\rangle_{Y}
}
{ =} { \left(  2a^2-2b^2+2   \right)^2  +  16 a^2b^2 -16a^2
}
{ =} { 4  \left(  \left(  a^2-b^2+1   \right)^2  +  4 a^2b^2 -4a^2  \right)
}
{ =} { 4 \left(  a^4+b^4+1 -2a^2+2a^2b^2 -2b^2  \right)
}
{ =} { 4  \left(  a^2 +b^2-1  \right)^2
}
}
{}
{}{,}
dan

\mavergleichskettedisphandlinks
{\vergleichskettedisphandlinks
{ \left\langle  \begin{pmatrix}  4ab \\ - 2a^2+2b^2+2 \\  4b   \end{pmatrix}   ,  \begin{pmatrix}  4ab \\ - 2a^2+2b^2+2 \\  4b   \end{pmatrix}     \right\rangle_{Y}
}
{ =} { 4  \left(  a^2 +b^2-1  \right)^2
}
{ } {}
{ } {}
{ } {}
}
{}{}{.}
Dengan memperhitungkan faktor awal ${ \frac{   1  }{  \left(  a^2+b^2-1  \right)^2  } }$, nilai hasil kali dalam pada $Y$ berimpit dengan hasil kali dalam pada cakram.

}


\chapter{Lembar Kerja 18}
\setcounter{section}{18}

  \zwischenueberschrift{Soal latihan}

\inputaufgabe
{}
{

Misalkan $L,M,N$ adalah
\definitionsverweis {manifold-manifold Riemann}{}{.}
Tunjukkan sifat-sifat berikut.
\aufzaehlungdrei{Identitas
\maabbdisp {\operatorname{Id}_{  M  }} { M } { M
} {}
adalah suatu
\definitionsverweis {isometri}{}{.}
}{Jika
\maabb {\varphi} { L } { M
} {}
adalah suatu isometri, maka
\definitionsverweis {pemetaan invers}{}{}
\maabb {\varphi^{-1}} { M } { L
} {}
juga merupakan suatu isometri.
}{Jika
\maabb {\varphi} { L } { M
} {}
dan
\maabb {\psi} { M } { N
} {}
adalah isometri, maka $\psi \circ \varphi$ juga merupakan suatu isometri.
}

}
{} {}

\inputaufgabe
{}
{

Misalkan $M$ adalah suatu
\definitionsverweis {manifold Riemann}{}{.}
Tunjukkan bahwa himpunan semua
\definitionsverweis {isometri}{}{}
pada $M$ membentuk suatu
\definitionsverweis {grup}{}{.}

}
{} {}

\inputaufgabe
{}
{

Misalkan
\maabbeledisp {\varphi} {\R} {S^1 \subseteq \R^2
} { t } { \begin{pmatrix}  \cos  t  \\ \sin  t    \end{pmatrix}
} {.}
Tunjukkan bahwa $\varphi$ adalah suatu
\definitionsverweis {isometri lokal}{}{}
antara
\definitionsverweis {manifold-manifold Riemann}{}{}
jika $\R$ dan $\R^2$ dilengkapi dengan
\definitionsverweis {struktur Euklides}{}{}
dan $S^1$ dilengkapi dengan struktur Riemann terinduksi.

}
{} {}

\inputaufgabe
{}
{

Misalkan
\maabb {\varphi} {\R^n } { \R^n
} {}
adalah suatu
\definitionsverweis {pemetaan linear}{}{,}
dan misalkan $\R^n$ dilengkapi dengan
\definitionsverweis {struktur Euklides}{}{}
beserta
\definitionsverweis {struktur Riemann}{}{}
yang diinduksikannya. Tunjukkan bahwa $\varphi$ merupakan suatu
\definitionsverweis {isometri linear}{}{}
jika dan hanya jika $\varphi$ merupakan suatu
\definitionsverweis {isometri}{}{}
terhadap struktur Riemann tersebut.

}
{} {}

\inputaufgabe
{}
{

Misalkan

\mavergleichskette
{\vergleichskette
{ I
}
{ \subseteq }{ \R
}
{  }{
}
{    }{
}
{   }{
}
}
{}{}{}
adalah suatu
\definitionsverweis {interval terbuka}{}{}
dan
\maabbdisp {\gamma} { I } { \R^n
} {}
adalah suatu
\definitionsverweis {kurva yang terdiferensialkan secara kontinu}{}{}
yang injektif. Kita memandang $I$ maupun $\R^n$ sebagai
\definitionsverweis {manifold Riemann}{}{}
melalui
\definitionsverweis {struktur Euklides}{}{}
masing-masing. Andaikan bahwa

\mavergleichskette
{\vergleichskette
{ \gamma(I)
}
{ \subseteq }{ U
}
{ \subseteq }{ \R^n
}
{    }{
}
{   }{
}
}
{}{}{}
dengan $U$
\definitionsverweis {terbuka}{}{}
merupakan suatu
\definitionsverweis {submanifold tertutup}{}{.}
Tunjukkan bahwa pemetaan

\mavergleichskettedisp
{\vergleichskette
{ \gamma
}
{ :} { I
}
{ \longrightarrow} { \gamma(I) \subseteq \R^n
}
{ } {
}
{ } {
}
}
{}{}{}
merupakan suatu
\definitionsverweis {isometri}{}{}
jika dan hanya jika $\gamma$
\definitionsverweis {berparametrisasi panjang busur}{}{.}

}
{} {}

\inputaufgabe
{}
{

Misalkan $L_1,L_2,M_1,M_2$ adalah
\definitionsverweis {manifold-manifold Riemann}{}{}
dan misalkan
\maabb {\varphi_1} {L_1 } { M_1
} {}
serta
\maabb {\varphi_2} {L_2 } { M_2
} {}
adalah
\definitionsverweis {isometri}{}{.}
Tunjukkan bahwa
\maabbdisp {\varphi_1 \times \varphi_2} {L_1 \times L_2 } { M_1 \times M_2
} {}
juga merupakan suatu isometri.

}
{} {}

\inputaufgabe
{}
{

Kita meninjau
\definitionsverweis {difeomorfisme}{}{}
\maabbeledisp {\varphi} {S^1 \times \R_+} { \Complex \setminus \{0\}  =  \R^2 \setminus \{(0,0)\}
} {(s,t) } { st
} {,}
dengan silinder $S^1 \times \R_+$ dilengkapi
\definitionsverweis {struktur produk Riemann}{}{.}
Uraikan struktur Riemann pada
\mathl{\R^2 \setminus \{(0,0)\}}{}
yang membuat $\varphi$ menjadi suatu
\definitionsverweis {isometri}{}{.}

}
{} {}

\inputaufgabegibtloesung
{}
{

Misalkan
\mathkor {} {r} {dan} {R} {}
adalah bilangan real positif dengan

\mavergleichskette
{\vergleichskette
{ 0
}
{ < }{ r
}
{ < }{ R
}
{    }{
}
{   }{
}
}
{}{}{,}
dan kita meninjau torus
\zusatzklammer {tertanam} {} {}

\mavergleichskettedisp
{\vergleichskette
{ T
}
{ =} { \left\{ (x,y,z) \in \R^3  \mid   \left(  \sqrt{x^2+y^2} -R  \right)^2 +z^2  = r^2         \right\}
}
{ \subseteq} { \R^3
}
{ } {
}
{ } {
}
}
{}{}{}
dengan
\definitionsverweis {struktur Riemann terinduksi}{}{.}
Tunjukkan bahwa untuk setiap

\mavergleichskette
{\vergleichskette
{ \alpha
}
{ \in }{ [0,2 \pi [
}
{  }{
}
{    }{
}
{   }{
}
}
{}{}{}
pemetaan
\maabbeledisp {\Psi_\alpha} {\R^3} {\R^3
} { \left(  x   , \,   y  , \,  z                 \right) } { \left(  x \cos \alpha - y \sin  \alpha   , \,   x \sin \alpha + y \cos  \alpha  , \,  z                 \right)
} {,}
menginduksi suatu
\definitionsverweis {isometri}{}{}
dari $T$ ke dirinya sendiri.

}
{} {}

\inputaufgabe
{}
{

Hitung panjang lintasan linear dari
\mathl{\begin{pmatrix}  0  \\ 1   \end{pmatrix}}{} ke
\mathl{\begin{pmatrix}  1  \\ 1   \end{pmatrix}}{} dalam
\definitionsverweis {setengah bidang Poincaré}{}{.}

}
{} {}

\inputaufgabe
{}
{

Hitung luas persegi panjang dalam
\definitionsverweis {setengah bidang Poincaré}{}{}
yang diberikan oleh titik-titik sudut
\mathl{(1,1),\, (5,1),\, (1,3),\, (5,3)}{}{.}

}
{} {}

\inputaufgabegibtloesung
{}
{

Tunjukkan bahwa
\definitionsverweis {pemetaan-pemetaan}{}{}
\maabbeledisp {\varphi} { {\mathbb H} } { U \left(  0,1  \right)
} { z} { { \frac{  z-  { \mathrm i}  }{ z + { \mathrm i}  } }
} {,}
dan
\maabbeledisp {\psi} {  U \left(  0,1  \right)  } { {\mathbb H}
} {w} { { \frac{   (w+1)  { \mathrm i}  }{ 1-w  } }
} {,}
memberikan suatu
\definitionsverweis {pemetaan biholomorfik}{}{}
antara
\definitionsverweis {setengah bidang atas}{}{}
${\mathbb H}$ dan cakram satuan terbuka $U \left(  0,1  \right)$.

}
{} {}

\inputaufgabe
{}
{

Tentukan
\definitionsverweis {turunan}{}{}
dari
\definitionsverweis {pemetaan-pemetaan}{}{}
\maabbeledisp {\varphi} { {\mathbb H} } { U \left(  0,1  \right)
} { z} { { \frac{  z-  { \mathrm i}  }{ z + { \mathrm i}  } }
} {,}
dan
\maabbeledisp {\psi} {  U \left(  0,1  \right)  } { {\mathbb H}
} {w} { { \frac{   (w+1)  { \mathrm i}  }{ 1-w  } }
} {.}

}
{} {}

\inputaufgabegibtloesung
{}
{

Uraikan
\definitionsverweis {pemetaan}{}{}
\maabbeledisp {\psi} { \Complex \setminus \{1\} } { \Complex
} {w} { { \frac{   (w+1)  { \mathrm i}  }{ 1-w  } }
} {}
dalam koordinat real.

}
{} {}

\inputaufgabegibtloesung
{}
{

Tentukan
\definitionsverweis {turunan parsial}{}{}
dari
\definitionsverweis {pemetaan}{}{}
\maabbeledisp {\psi} { \R^2 \setminus \{  (1,0 )\} } { \R^2
} {  \begin{pmatrix} a \\b   \end{pmatrix} } {  { \frac{  1 }{ -2a+1+a^2+b^2 } }  \begin{pmatrix} -2b \\1-a^2-b^2   \end{pmatrix}
} {.}

}
{} {}

\inputaufgabe
{}
{

Tentukan titik-titik potong garis yang melalui
\mathl{\begin{pmatrix}  0  \\ 0\\  -1   \end{pmatrix}}{} dan suatu titik
\mathl{\begin{pmatrix}  a  \\ b \\  0   \end{pmatrix}}{} dengan

\mavergleichskette
{\vergleichskette
{ a^2+b^2
}
{ < }{ 1
}
{  }{
}
{    }{
}
{   }{
}
}
{}{}{}
dengan hiperboloid dua lembar; lihat
Contoh 18.11.

}
{} {}

Misalkan
\mathkor {} {L} {dan} {M} {}
adalah dua
$C^k$-\definitionsverweis {manifold}{}{.}
Suatu
\definitionsverweis {pemetaan kontinu}{}{}
\maabbdisp {\varphi} { L } { M
} {}
disebut suatu
\definitionswortpraemath {C^k}{ difeomorfisme lokal }{,}
jika untuk setiap titik

\mavergleichskette
{\vergleichskette
{ P
}
{ \in }{ L
}
{  }{
}
{    }{
}
{   }{
}
}
{}{}{}
terdapat suatu
\definitionsverweis {lingkungan terbuka}{}{}

\mavergleichskette
{\vergleichskette
{ P
}
{ \in }{ U
}
{ \subseteq }{ L
}
{    }{
}
{   }{
}
}
{}{}{}
sedemikian sehingga $\varphi  \vert_U$ menginduksi suatu
\definitionsverweis {difeomorfisme}{}{}
antara $U$ dan $\varphi(U)$.

  \zwischenueberschrift{Soal untuk dikumpulkan}

\inputaufgabe
{}
{

Kita meninjau permukaan elipsoid

\mavergleichskettedisp
{\vergleichskette
{ E
}
{ =} { \left\{ (x,y,z)  \mid   2x^2+3y^2+5z^2 = 1        \right\}
}
{ \subset} { \R^3
}
{ } {
}
{ } {
}
}
{}{}{}
dengan
\definitionsverweis {struktur Riemann}{}{}
yang diinduksi oleh ruang sekitarnya. Tentukan
\definitionsverweis {isometri linear}{}{}
pada $\R^3$ yang menginduksi suatu
\definitionsverweis {isometri}{}{}
pada $E$.

}
{} {}

\inputaufgabe
{4}
{

Misalkan $L,M$ adalah
\definitionsverweis {manifold terdiferensialkan}{}{,}
misalkan
\maabbdisp {\varphi} { L } { M
} {}
adalah suatu
\definitionsverweis {difeomorfisme lokal}{}{,}
dan misalkan $M$ adalah suatu
\definitionsverweis {manifold Riemann}{}{.}
Tunjukkan bahwa terdapat tepat satu struktur Riemann pada $L$ yang membuat $\varphi$ menjadi suatu
\definitionsverweis {isometri lokal}{}{.}

}
{} {}

\inputaufgabe
{3}
{

Berikan suatu contoh
\definitionsverweis {manifold Riemann}{}{}
$M$ dan suatu
\definitionsverweis {difeomorfisme}{}{}
\maabbdisp {\varphi} { M } { M
} {,}
yang bukan suatu
\definitionsverweis {isometri}{}{.}

}
{} {}

\inputaufgabe
{4}
{

Hitung panjang busur pada setengah lingkaran atas berjari-jari $\sqrt{2}$ dari
\mathl{\begin{pmatrix}  1  \\ 1   \end{pmatrix}}{} ke
\mathl{\begin{pmatrix}  -1 \\ 1   \end{pmatrix}}{} dalam
\definitionsverweis {setengah bidang Poincaré}{}{.}

}
{} {}

\inputaufgabe
{4}
{

Hitung luas cakram berpusat di
\mathl{\begin{pmatrix}  2  \\ 2   \end{pmatrix}}{} dan berjari-jari $1$ dalam
\definitionsverweis {setengah bidang Poincaré}{}{.}

}
{} {}

\inputaufgabe
{4}
{

Uraikan suatu
\definitionsverweis {isometri}{}{}
eksplisit antara
\definitionsverweis {setengah bidang Poincaré}{}{}
dan lembar atas hiperboloid dua lembar.

}
{} {}


\section*{Solusi yang disediakan oleh sumber}
\addcontentsline{toc}{section}{Solusi yang disediakan oleh sumber}
\subsection*{Solusi Soal 18.8}
\input{generated/worksheet18_exercise08_solution.id.build.tex}
\subsection*{Solusi Soal 18.11}
\input{generated/worksheet18_exercise11_solution.id.build.tex}
\subsection*{Solusi Soal 18.13}
Kita tulis

\mavergleichskette
{\vergleichskette
{ w
}
{ = }{ a+b { \mathrm i}
}
{  }{
}
{    }{
}
{   }{
}
}
{}{}{.}
Maka

\mavergleichskettealign
{\vergleichskettealign
{  \psi ( a+b { \mathrm i}   )
}
{ =} { { \frac{   \left(  a+b { \mathrm i} +1  \right) { \mathrm i}      }{  1-  a-b { \mathrm i}      } }
}
{ =} { { \frac{     a { \mathrm i} +{ \mathrm i}  -b  }{  1-  a-b { \mathrm i}      } }
}
{ =} { { \frac{     \left(  a { \mathrm i} +{ \mathrm i}  -b  \right) \left(  1-  a+b { \mathrm i}  \right)  }{  \left(  1-  a-b { \mathrm i}  \right) \left(  1-  a+b { \mathrm i}  \right)      } }
}
{ =} { { \frac{     \left(  a+1 \right) \left(  1-a  \right) { \mathrm i}  -b^2 { \mathrm i} -b \left(  1-  a \right) -b\left(  a+1  \right)  }{  \left(  1-  a \right)^2 +b^2  } }
}
}
{
\vergleichskettefortsetzungalign
{ =} { { \frac{     \left(  1-a ^2-b^2  \right) { \mathrm i}   -2 b  }{  \left(  1-  a \right)^2 +b^2  } }
}
{ } {
}
{ } {}
{ } {}
}
{}{.}
Jadi, dalam koordinat real,

\mavergleichskettedisp
{\vergleichskette
{ \psi \begin{pmatrix}  a  \\b   \end{pmatrix}
}
{ =} { { \frac{   1  }{  -2a+1+a^2+b^2 } } \begin{pmatrix}  -2b \\ 1-a^2-b^2   \end{pmatrix}
}
{ } {
}
{ } {
}
{ } {
}
}
{}{}{.}

\subsection*{Solusi Soal 18.14}
Turunan parsial dari

\mavergleichskettedisp
{\vergleichskette
{ \psi \begin{pmatrix}  a  \\b   \end{pmatrix}
}
{ =} { { \frac{   1 }{  -2a+1+a^2+b^2 } } \begin{pmatrix}  -2b \\ 1-a^2-b^2   \end{pmatrix}
}
{ } {
}
{ } {
}
{ } {
}
}
{}{}{}
adalah

\mavergleichskettealign
{\vergleichskettealign
{ { \frac{   \partial   \psi  }{  \partial  a  } }
}
{ =} { { \frac{   1 }{  \left(  -2a+1+a^2+b^2  \right)^2  } } \begin{pmatrix}  2b \left(  2a-2  \right)  \\ -2a \left(  -2a+1+a^2+b^2  \right) - \left(  1-a^2-b^2  \right) \left(  -2+2a  \right)    \end{pmatrix}
}
{ =} { { \frac{   1 }{  \left(  -2a+1+a^2+b^2  \right)^2  } } \begin{pmatrix}  2b \left(  2a-2  \right)  \\ 4a^2 -4a+2 \left(  1-a^2-b^2  \right)    \end{pmatrix}
}
{ =} { { \frac{   1 }{  \left(  -2a+1+a^2+b^2  \right)^2  } } \begin{pmatrix}  4ab -4b  \\ 2a^2 -4a+2 -2 b^2     \end{pmatrix}
}
{ } {
}
}
{}
{}{}
dan

\mavergleichskettealign
{\vergleichskettealign
{ { \frac{   \partial   \psi  }{  \partial  b  } }
}
{ =} { { \frac{   1 }{  \left(  -2a+1+a^2+b^2  \right)^2  } } \begin{pmatrix}  -2 \left(  -2a+1+a^2+b^2  \right) +4b^2   \\ -2b \left(  -2a+1+a^2+b^2  \right) - 2b \left(  1-a^2-b^2  \right)     \end{pmatrix}
}
{ =} { { \frac{   1 }{  \left(  -2a+1+a^2+b^2  \right)^2  } } \begin{pmatrix}  4a  -2-2a^2+2b^2   \\ 4ab -4b    \end{pmatrix}
}
{ } {
}
{ } {
}
}
{}
{}{.}


\part{Unit 19}
\chapter[Kuliah 19]{Kuliah 19: Pemetaan Weingarten dan Kelengkungan Gauss}
\input{generated/lecture19.id.build.tex}

\chapter{Lembar Kerja 19}
\input{generated/worksheet19.id.build.tex}

\part{Unit 20}
\chapter[Kuliah 20]{Kuliah 20: Turunan Eksterior dan Setengah Ruang Euklides}
\setcounter{section}{20}

  \zwischenueberschrift{Turunan eksterior}

Pada kuliah ini kita akan mengenal suatu objek matematika jenis baru, yang disebut turunan eksterior. Turunan ini merupakan konsep turunan yang mengubah bentuk diferensial berderajat $k$ menjadi bentuk diferensial berderajat
\mathl{k+1}{.} Untuk bentuk diferensial berderajat $0$, yaitu suatu fungsi $f$, turunan eksterior yang bersesuaian hanyalah bentuk diferensial berderajat $1$, yaitu $df$, yang memasangkan pada setiap titik $P$
\zusatzklammer {untuk ruang Euklides} {} {}
diferensial total
\maabbdisp {\left(Df\right)_{P}} { \R^n } { \R
} {}
atau
\zusatzklammer {untuk suatu manifold $M$} {} {}
pemetaan singgung
\maabbdisp {T_P(f)} {T_PM} {\R
} {.}

Dalam kalkulus diferensial satu dimensi, fungsi dan turunannya atau antiturunannya merupakan objek sejenis
\zusatzklammer {hal ini masih berlaku untuk kurva terdiferensialkan} {} {,}
tetapi sejak diperkenalkannya diferensial total bagi suatu fungsi beberapa variabel, turunannya sudah merupakan objek yang secara mendasar berbeda dari fungsi itu sendiri. Turunan berarah tingkat tinggi memang dapat didefinisikan sepanjang arah-arah yang telah ditentukan dan kembali berupa fungsi, tetapi masing-masing hanya menangkap satu aspek parsial dari turunan fungsi, sedangkan diferensial total memuat seluruh informasinya.

Perbedaan mendasar antara fungsi dan turunannya ini juga berkaitan dengan kenyataan bahwa, dalam dimensi lebih tinggi, kita belum membahas pertanyaan sebaliknya: turunan mana yang memiliki antiturunan. Suatu fungsi beberapa variabel tidak dapat memiliki antiturunan dalam pengertian ini; pertanyaan tersebut hanya bermakna bagi suatu bentuk diferensial-$1$
\zusatzklammer {atau medan vektor yang bersesuaian} {} {.}
Teorema Schwarz tentang dapat dipertukarkannya turunan berarah sudah memberikan suatu kriteria perlu yang penting bagi keberadaan antiturunan dari suatu bentuk diferensial-$1$.

Pertanyaan tentang antiturunan atau bentuk primitif memperoleh kerangka alaminya dalam teori turunan eksterior. Selain itu, teori ini memungkinkan teorema Stokes dirumuskan secara ringkas. Turunan eksterior juga dapat digunakan untuk mengarakterisasi sifat-sifat topologis penting suatu manifold, walaupun hal tersebut jauh melampaui kuliah ini.

\inputdefinition
{}
{

Misalkan

\mavergleichskette
{\vergleichskette
{ U
}
{ \subseteq }{ \R^n
}
{  }{
}
{    }{
}
{   }{
}
}
{}{}{}
\definitionsverweis {terbuka}{}{}
dan misalkan

\mavergleichskette
{\vergleichskette
{ \omega
}
{ \in }{
{ \mathcal E }^{ k } (  U   )
}
{  }{
}
{    }{
}
{   }{
}
}
{}{}{}
suatu
$k$-\definitionsverweis {bentuk diferensial}{}{}
yang \definitionsverweis {terdiferensialkan secara kontinu}{}{}
dengan penyajian

\mavergleichskettedisp
{\vergleichskette
{ \omega
}
{ =} { \sum_{I \subseteq \{1 , \ldots , n \},\, { \# \left(   I  \right) } = k } f_I dx_I
}
{ } {
}
{ } {
}
{ } {
}
}
{}{}{}
dengan
\definitionsverweis {fungsi-fungsi yang terdiferensialkan secara kontinu}{}{}
\maabbdisp {f_I} { U } {\R
} {.}
Maka bentuk diferensial berderajat $(k+1)$

\mavergleichskettedisp
{\vergleichskette
{d \omega
}
{ =} { \sum_{I \subseteq \{1 , \ldots , n \},\,  { \# \left(   I  \right) } = k } df_I \wedge dx_I
}
{ =} { \sum_{I \subseteq \{1 , \ldots , n \},\,  { \# \left(   I  \right) } = k } \left(  \sum_{j = 1}^n { \frac{   \partial  f_I  }{  \partial   x_j   } } dx_j  \right)  \wedge dx_I
}
{ } {
}
{ } {
}
}
{}{}{}
disebut \definitionswort {turunan eksterior}{} dari $\omega$.

}

Kadang-kadang kita menyebutnya secara lebih tepat sebagai turunan eksterior ke-$k$. Derajat keterdiferensialan bentuk diferensial tersebut berkurang sebesar $1$, sebagaimana dapat langsung dilihat dari fungsi-fungsi koefisiennya.

Turunan eksterior relevan untuk

\mavergleichskette
{\vergleichskette
{ k
}
{ = }{ 0 , \ldots , n
}
{  }{
}
{    }{
}
{   }{
}
}
{}{}{}
sedangkan untuk

\mavergleichskette
{\vergleichskette
{ k
}
{ \geq }{ n
}
{  }{
}
{    }{
}
{   }{
}
}
{}{}{}
pemetaan ini merupakan pemetaan nol. Jika kita membatasi diri pada bentuk-bentuk diferensial mulus, diperoleh suatu barisan turunan eksterior

\mathdisp {C^\infty(U,\R) =
{ \mathcal E }^{  0  }_{ \infty } (  U   ) \stackrel{d}{\longrightarrow}
{ \mathcal E }^{  1  }_{ \infty } (  U   ) \stackrel{d}{\longrightarrow}
{ \mathcal E }^{  2  }_{ \infty } (  U   ) \stackrel{d}{\longrightarrow} \ldots \stackrel{d}{\longrightarrow}
{ \mathcal E }^{ n-1 }_{ \infty } (  U   ) \stackrel{d}{\longrightarrow}
{ \mathcal E }^{  n  }_{ \infty } (  U   ) \stackrel{d}{\longrightarrow} 0} { . }

Pada posisi pertama hanya terdapat turunan suatu fungsi
\zusatzklammer {satu-satunya himpunan indeks dengan $0$ elemen adalah himpunan kosong} {} {,}
yakni pemetaan
\mathl{f \mapsto df}{.}

Sifat-sifat terpenting turunan eksterior dirangkum sebagai berikut.


\inputfaktbeweis
{Differentialform/Offene Menge/Äußere Ableitung/Eigenschaften/Fakt}
{Lemma}
{}
{

\faktsituation {Misalkan

\mavergleichskette
{\vergleichskette
{ U
}
{ \subseteq }{ \R^n
}
{  }{
}
{    }{
}
{   }{
}
}
{}{}{}
\definitionsverweis {terbuka}{}{,}

\mavergleichskette
{\vergleichskette
{ k
}
{ \in }{ \N
}
{  }{
}
{    }{
}
{   }{
}
}
{}{}{}
dan misalkan
\maabbeledisp {d} {
{ \mathcal E }^{  k  }_{ 1 } (  U   ) } {
{ \mathcal E }^{ k+1 }_{ 0 } (  U   )
} { \omega} { d \omega
} {,}
adalah
\definitionsverweis {turunan eksterior}{}{.}}
\faktuebergang {Maka berlaku sifat-sifat berikut.}
\faktfolgerung {\aufzaehlungfuenf{Turunan eksterior
\maabbdisp {d} {
{ \mathcal E }^{  0  }_{ 1 } (  U   ) } {
{ \mathcal E }^{  1  }_{ 0 } (  U   )
} {,}
adalah
\definitionsverweis {diferensial total}{}{.}
}{Turunan eksterior bersifat
$\R$-\definitionsverweis {linear}{}{.}
}{Untuk

\mavergleichskette
{\vergleichskette
{ \omega
}
{ \in }{
{ \mathcal E }^{  k  }_{ 1 } (  U   )
}
{  }{
}
{    }{
}
{   }{
}
}
{}{}{}
dan

\mavergleichskette
{\vergleichskette
{ \tau
}
{ \in }{
{ \mathcal E }^{  \ell }_{ 1 } (  U   )
}
{  }{
}
{    }{
}
{   }{
}
}
{}{}{}
berlaku aturan hasil kali

\mavergleichskettedisp
{\vergleichskette
{ d( \omega \wedge \tau)
}
{ =} {( d \omega) \wedge \tau  + (-1)^k \omega  \wedge  (d \tau)
}
{ } {
}
{ } {
}
{ } {
}
}
{}{}{.}
Untuk

\mavergleichskette
{\vergleichskette
{k
}
{ = }{ 0
}
{  }{
}
{    }{
}
{   }{
}
}
{}{}{}
hal ini dibaca sebagai

\mavergleichskettedisp
{\vergleichskette
{ d( f \tau)
}
{ =} {( d f) \wedge \tau  +  f  d \tau
}
{ } {
}
{ } {
}
{ } {
}
}
{}{}{}
}{Untuk setiap bentuk diferensial yang dua kali terdiferensialkan secara kontinu, $\omega$,

\mavergleichskette
{\vergleichskette
{ d(d \omega)
}
{ = }{ 0
}
{  }{
}
{    }{
}
{   }{
}
}
{}{}{.}
}{Untuk suatu
\definitionsverweis {pemetaan yang dua kali terdiferensialkan secara kontinu}{}{}
\zusatzklammer {dengan

\mavergleichskettek
{\vergleichskettek
{ W
}
{ \subseteq \R^m }{}
{  }{}
{    }{}
{   }{}
}
{}{}{}
terbuka} {} {}
\maabbdisp {\psi} { W } { U
} {}
dan setiap

\mavergleichskette
{\vergleichskette
{ \omega
}
{ \in }{
{ \mathcal E }^{  k  }_{ 1 } (  U   )
}
{  }{
}
{    }{
}
{   }{
}
}
{}{}{}
berlaku bagi
\definitionsverweis {bentuk diferensial tarik balik}{}{}

\mavergleichskettedisp
{\vergleichskette
{d (\psi^* \omega)
}
{ =} { \psi^* (d \omega)
}
{ } {
}
{ } {
}
{ } {
}
}
{}{}{.}
}}
\faktzusatz {}
\faktzusatz {}

}
{

\teilbeweis {}{}{}
{(1) langsung mengikuti definisi
\zusatzklammer {himpunan kosong adalah satu-satunya himpunan indeks yang relevan} {} {.}}
{}
\teilbeweis {}{}{}
{(2). Linearitas langsung mengikuti definisi,
linearitas
dari
\definitionsverweis {diferensial total}{}{}
dan
\definitionsverweis {multilinearitas}{}{}
dari
\definitionsverweis {produk eksterior}{}{.}}
{}

\teilbeweis {}{}{}
{(3). Misalkan
\mathl{x_1 , \ldots , x_n}{} koordinat-koordinat pada $\R^n$. Berkat linearitas $d$ dan
multilinearitas produk eksterior,
kedua bentuk diferensial dapat ditulis sebagai
\mathkor {} {\omega=f dx_I} {dan} {\tau=g dx_J} {}
dengan himpunan indeks
\mathkor {} {I=\{i_1 , \ldots , i_k\}} {dan} {J=\{ j_1 , \ldots , j_\ell \}} {.}
Maka

\mavergleichskettealignhandlinks
{\vergleichskettealignhandlinks
{ d( \omega \wedge \tau)
}
{ =} { d \left(  fdx_I \wedge gdx_J  \right)
}
{ =} { d \left(  (f g) dx_I \wedge dx_J  \right)
}
{ =} { \sum_{s = 1}^n { \frac{   \partial  (fg)  }{  \partial   x_s  } }  dx_s  \wedge  dx_I  \wedge dx_J
}
{ =} { \sum_{s = 1}^n \left(  g { \frac{   \partial   f   }{  \partial   x_s  } } + f { \frac{   \partial   g   }{  \partial   x_s  } }  \right)  dx_s \wedge  dx_I  \wedge dx_J
}
}
{
\vergleichskettefortsetzungalign
{ =} { \sum_{s = 1}^n  g { \frac{   \partial   f   }{  \partial   x_s  } } dx_s \wedge  dx_I \wedge dx_J + \sum_{s = 1}^n  f { \frac{   \partial   g   }{  \partial   x_s  } }  dx_s \wedge  dx_I  \wedge dx_J
}
{ =} { \sum_{s = 1}^n  { \frac{   \partial   f   }{  \partial   x_s  } } dx_s \wedge  dx_I \wedge gdx_J + \sum_{s = 1}^n  { \frac{   \partial   g   }{  \partial   x_s  } }  dx_s \wedge f dx_I \wedge dx_J
}
{ =} { d(fdx_I ) \wedge g dx_J  + \sum_{s = 1}^n (-1)^k  f dx_I \wedge { \frac{   \partial   g   }{  \partial   x_s  } }  dx_s \wedge dx_J
}
{ =} { d(fdx_I ) \wedge g dx_J + (-1)^k  f dx_I \wedge  d( g dx_J)
}
}
{}{.}}
{}

\teilbeweis {}{}{}
{(4). Untuk suatu bentuk-$1$

\mavergleichskette
{\vergleichskette
{ \omega
}
{ = }{ \sum_{j = 1}^n g_jdx_j
}
{  }{
}
{    }{
}
{   }{
}
}
{}{}{}
dengan menggunakan

\mavergleichskette
{\vergleichskette
{ dx_i \wedge dx_j
}
{ = }{ -dx_j \wedge dx_i
}
{  }{
}
{    }{
}
{   }{
}
}
{}{}{}
diperoleh

\mavergleichskettealign
{\vergleichskettealign
{d \omega
}
{ =} { \sum_{j = 1 }^n  d(g_j dx_j)
}
{ =} { \sum_{j = 1 }^n   \left(  \sum_{i = 1}^n { \frac{   \partial  g_j   }{  \partial   x_i   } } dx_i \wedge dx_j  \right)
}
{ =} { \sum _{1 \leq i <j\leq n} \left(  { \frac{   \partial  g_j   }{  \partial   x_i   } } - { \frac{   \partial  g_i   }{  \partial   x_j   } }  \right) dx_i \wedge dx_j
}
{ } {
}
}
{}
{}{.}
Untuk fungsi $f$ yang dua kali terdiferensialkan secara kontinu,

\mavergleichskette
{\vergleichskette
{ df
}
{ = }{ \sum_{j = 1}^n g_j dx_j
}
{  }{
}
{    }{
}
{   }{
}
}
{}{}{}
dengan
\definitionsverweis {turunan parsial}{}{}

\mavergleichskette
{\vergleichskette
{ g_j
}
{ = }{ { \frac{   \partial   f   }{  \partial   x_j   } }
}
{  }{
}
{    }{
}
{   }{
}
}
{}{}{,}
sehingga

\mavergleichskette
{\vergleichskette
{ d(df)
}
{ = }{ 0
}
{  }{
}
{    }{
}
{   }{
}
}
{}{}{}
menurut
teorema Schwarz.

Untuk bentuk diferensial berderajat $k$, kita mengambil

\mavergleichskette
{\vergleichskette
{ \omega
}
{ = }{ f dx_{i_1 } \wedge \ldots \wedge dx_{i_k }
}
{  }{
}
{    }{
}
{   }{
}
}
{}{}{}
dan memperoleh

\mavergleichskettedisp
{\vergleichskette
{ d(d \omega)
}
{ =} { d (df \wedge  dx_{i_1 } \wedge \ldots \wedge dx_{i_k })
}
{ } {
}
{ } {
}
{ } {
}
}
{}{}{.}
Menurut aturan hasil kali (3), ungkapan ini merupakan jumlah dari
\mathl{k+1}{} produk eksterior; pada setiap produk itu, salah satu \anfuehrung{faktor produk-eksterior}{} berbentuk

\mavergleichskette
{\vergleichskette
{ d(dg)
}
{ = }{ 0
}
{  }{
}
{    }{
}
{   }{
}
}
{}{}{.}}
{}

\teilbeweis {}{}{}
{(5). Kita tulis

\mavergleichskettedisp
{\vergleichskette
{ \psi_i
}
{ =} { x_i \circ \psi
}
{ } {
}
{ } {
}
{ } {
}
}
{}{}{.}
Berkat linearitas turunan eksterior (2) dan
linearitas penarikan balik bentuk diferensial,
kita dapat mengambil

\mavergleichskette
{\vergleichskette
{ \omega
}
{ = }{ f dx_I
}
{  }{
}
{    }{
}
{   }{
}
}
{}{}{}
dengan

\mavergleichskette
{\vergleichskette
{ I
}
{ = }{ \{i_1 , \ldots , i_k\}
}
{  }{
}
{    }{
}
{   }{
}
}
{}{}{.}
Menurut
Soal 14.13
dan
Soal 14.16,
penarikan balik serasi dengan perkalian oleh fungsi skalar dan dengan produk eksterior. Dengan menggunakan aturan hasil kali (3), aturan (4), dan
aturan rantai
\zusatzklammer {dalam pengertian
Soal 14.14} {} {,}
diperoleh

\mavergleichskettealignhandlinks
{\vergleichskettealignhandlinks
{ d \left(  \psi^* \omega  \right)
}
{ =} { d \left(  \psi^* \left(  f dx_{i_1 } \wedge \ldots \wedge dx_{i_k }  \right)   \right)
}
{ =} { d \left(  \psi^*(f)  \psi^* \left(  dx_{i_1 } \wedge \ldots \wedge dx_{i_k }  \right)  \right)
}
{ =} { d \left(  \psi^*(f) \left(  \psi^* dx_{i_1 }  \right) \wedge \ldots \wedge \left(  \psi^* dx_{i_k }  \right)  \right)
}
{ =} { d \left(  \psi^*(f)  d \psi_{i_1 } \wedge \ldots \wedge d\psi_{i_k }  \right)
}
}
{
\vergleichskettefortsetzungalign
{ =} { d \left(  \psi^*(f)  \right) \wedge \left(  d \psi_{i_1 } \wedge \ldots \wedge d\psi_{i_k }   \right) + (\psi^*(f))  d \left(  d \psi_{i_1 } \wedge \ldots \wedge d\psi_{i_k }  \right)
}
{ =} { d \left(  \psi^*(f)  \right)  \wedge \left(  d \psi_{i_1 } \wedge \ldots \wedge d\psi_{i_k }  \right)
}
{ =} { \psi^* (df)  \wedge  d \psi_{i_1 } \wedge \ldots \wedge d\psi_{i_k }
}
{ =} { \psi^* ( df) \wedge \psi^*\left(  d x_{i_1 }  \right) \wedge \ldots \wedge \psi^*\left(  d x_{i_k }  \right)
}
}
{
\vergleichskettefortsetzungalign
{ =} { \psi^* \left(  df \wedge  d x_{i_1 } \wedge \ldots \wedge d x_{i_k }  \right)
}
{ =} { \psi^* \left(  d \left(  f d x_{i_1 } \wedge \ldots \wedge d x_{i_k }  \right)   \right)
}
{ } {}
{ } {}
}{.}}
{}

}

\inputdefinition
{}
{

Misalkan $M$ suatu
\definitionsverweis {manifold diferensiabel}{}{.}
Untuk suatu
\definitionsverweis {bentuk diferensial}{}{}
\definitionsverweis {terdiferensialkan}{}{}

\mavergleichskette
{\vergleichskette
{ \omega
}
{ \in }{
{ \mathcal E }^{  k  }_{ 1 } (  M   )
}
{  }{
}
{    }{
}
{   }{
}
}
{}{}{}
kita definisikan \definitionswort {turunan eksterior}{}

\mavergleichskette
{\vergleichskette
{ d\omega
}
{ \in }{
{ \mathcal E }^{ k+1 }_{ 0 } (  M   )
}
{  }{
}
{    }{
}
{   }{
}
}
{}{}{}
dengan mengacu pada
\definitionsverweis {kasus lokal}{}{}
dan peta-peta
\maabbdisp {\alpha} { U } { V
} {}
\zusatzklammer {dengan
\mavergleichskettek
{\vergleichskettek
{ U
}
{ \subseteq }{ M
}
{  }{
}
{    }{
}
{   }{
}
}
{}{}{}
dan

\mavergleichskettek
{\vergleichskettek
{ V
}
{ \subseteq }{ \R^n
}
{  }{
}
{    }{
}
{   }{
}
}
{}{}{}
terbuka} {} {}
melalui

\mavergleichskettedisp
{\vergleichskette
{ (d \omega){{|}}_U
}
{ =} { d \left(  \omega{{|}}_U  \right)
}
{ =} { \alpha^* \left(  d \left(  \alpha^{-1 *} \left(  \omega {{|}}_U  \right)  \right)  \right)
}
{ } {
}
{ } {
}
}
{}{}{.}

}

Jadi, bentuk diferensial yang dibatasi pada $U$ ditarik balik ke $V$, turunan eksteriornya diambil di sana sesuai aturan lokal, lalu hasilnya ditarik balik ke $U$. Kita perlu memastikan bahwa prosedur ini menghasilkan bentuk diferensial yang terdefinisi dengan baik pada $M$: untuk suatu titik

\mavergleichskette
{\vergleichskette
{ P
}
{ \in }{ M
}
{  }{
}
{    }{
}
{   }{
}
}
{}{}{}
hasil turunan eksterior tidak bergantung pada lingkungan koordinat yang digunakan. Misalkan diberikan dua peta untuk $P$; kita dapat langsung menganggap bahwa keduanya mempunyai domain definisi yang sama, yaitu $U$. Misalkan peta-peta itu
\maabbdisp {\alpha} { U } { V
} {}
dan
\maabbdisp {\beta} { U } { W
} {}
dan kita tetapkan

\mavergleichskette
{\vergleichskette
{\tau
}
{ = }{ \omega  \vert_U
}
{  }{
}
{    }{
}
{   }{
}
}
{}{}{.}
Dengan menerapkan
Lema 20.2 (5)
pada
\mathl{\alpha \circ \beta^{-1}}{} dan
\mathl{\alpha^{-1*} \tau}{,}
kita memperoleh

\mavergleichskettealignhandlinks
{\vergleichskettealignhandlinks
{ \beta^* \left(  d \left(  \beta^{-1*}(\tau)  \right)   \right)
}
{ =} { \left(  \beta \circ \alpha^{-1} \circ \alpha  \right)^* \left(  d \left(  \left(  \alpha^{-1} \circ \alpha  \circ \beta^{-1}  \right)^* (\tau )  \right)   \right)
}
{ =} {\alpha^* \left(  \beta \circ \alpha^{-1}  \right)^* \left(  d \left(  \left(  \alpha  \circ  \beta^{-1}  \right)^*  \alpha^{-1*} (\tau)  \right)   \right)
}
{ =} {\alpha^* \left(  d \left(  \alpha^{-1*}( \tau)  \right)   \right)
}
{ } {}
}
{}
{}{.}

Sifat-sifat dasar di atas juga terbawa ke manifold.


\inputfaktbeweis
{Differentialform/Mannigfaltigkeit/Äußere Ableitung/Eigenschaften/Fakt}
{Satz}
{}
{

\faktsituation {Misalkan $M$ suatu
\definitionsverweis {manifold diferensiabel}{}{,}

\mavergleichskette
{\vergleichskette
{ k
}
{ \in }{ \N
}
{  }{
}
{    }{
}
{   }{
}
}
{}{}{}
dan misalkan
\maabbeledisp {d} {
{ \mathcal E }^{  k  }_{ 1 } (  M   ) } {
{ \mathcal E }^{ k+1 }_{ 0 } (  M   )
} { \omega} { d \omega
} {,}
adalah
\definitionsverweis {turunan eksterior}{}{.}}
\faktuebergang {Maka berlaku sifat-sifat berikut.}
\faktfolgerung {\aufzaehlungfuenf{Turunan eksterior
\maabbdisp {d} {
{ \mathcal E }^{  0  }_{ 1 } (  M   ) } {
{ \mathcal E }^{  1  }_{ 0 } (  M   )
} {,}
adalah
\definitionsverweis {pemetaan singgung}{}{.}
}{Turunan eksterior bersifat
$\R$-\definitionsverweis {linear}{}{.}
}{Untuk

\mavergleichskette
{\vergleichskette
{ \omega
}
{ \in }{
{ \mathcal E }^{  k  }_{ 1 } (  M   )
}
{  }{
}
{    }{
}
{   }{
}
}
{}{}{}
dan

\mavergleichskette
{\vergleichskette
{ \tau
}
{ \in }{
{ \mathcal E }^{  \ell }_{ 1 } (  M   )
}
{  }{
}
{    }{
}
{   }{
}
}
{}{}{}
berlaku aturan hasil kali

\mavergleichskettedisp
{\vergleichskette
{ d( \omega \wedge \tau)
}
{ =} { ( d \omega) \wedge \tau  + (-1)^k \omega \wedge d \tau
}
{ } {
}
{ } {
}
{ } {
}
}
{}{}{}
}{Untuk setiap bentuk diferensial yang dua kali terdiferensialkan secara kontinu, $\omega$,

\mavergleichskette
{\vergleichskette
{ d(d \omega)
}
{ = }{ 0
}
{  }{
}
{    }{
}
{   }{
}
}
{}{}{.}
}{Misalkan $L$ suatu manifold diferensiabel lainnya. Untuk suatu
\definitionsverweis {pemetaan yang dua kali terdiferensialkan secara kontinu}{}{}
\maabbdisp {\psi} { L } { M
} {}
dan setiap

\mavergleichskette
{\vergleichskette
{ \omega
}
{ \in }{
{ \mathcal E }^{  k  }_{ 1 } (  M   )
}
{  }{
}
{    }{
}
{   }{
}
}
{}{}{}
berlaku bagi
\definitionsverweis {bentuk diferensial tarik balik}{}{}

\mavergleichskettedisp
{\vergleichskette
{ d (\psi^* \omega)
}
{ =} { \psi^* (d \omega)
}
{ } {
}
{ } {
}
{ } {
}
}
{}{}{.}
}}
\faktzusatz {}
\faktzusatz {}

}
{

Semua pernyataan ini bersifat lokal, sehingga langsung diperoleh dari
Lema 20.2.

}

\inputdefinition
{}
{

Misalkan $M$ suatu
\definitionsverweis {manifold diferensiabel}{}{.}
Suatu
\definitionsverweis {bentuk diferensial}{}{}
\definitionsverweis {terdiferensialkan}{}{}
$\omega$ pada $M$ disebut \definitionswort {tertutup}{,} jika
\definitionsverweis {turunan eksterior}{}{}

\mavergleichskette
{\vergleichskette
{ d \omega
}
{ = }{ 0
}
{  }{
}
{    }{
}
{   }{
}
}
{}{}{.}

}

\inputdefinition
{}
{

Misalkan $M$ suatu
\definitionsverweis {manifold diferensiabel}{}{.}
Untuk derajat positif, suatu
$k$-\definitionsverweis {bentuk diferensial}{}{}
$\omega$ pada $M$, disebut \definitionswort {eksak}{,} jika terdapat suatu bentuk diferensial yang
\definitionsverweis {terdiferensialkan}{}{}
dan berderajat \mathl{(k-1)}{}, yang kita nyatakan dengan $\sigma$, pada $M$ sedemikian sehingga

\mavergleichskette
{\vergleichskette
{ d \sigma
}
{ = }{ \omega
}
{  }{
}
{    }{
}
{   }{
}
}
{}{}{.}

}

Jadi, bentuk diferensial eksak adalah bentuk diferensial berderajat positif yang memiliki suatu \stichwort {bentuk primitif} {} $\sigma$. Dengan istilah-istilah ini, pernyataan di atas

\mavergleichskette
{\vergleichskette
{ dd
}
{ = }{ 0
}
{  }{
}
{    }{
}
{   }{
}
}
{}{}{}
dapat dirumuskan dengan mengatakan bahwa setiap bentuk eksak adalah tertutup. Jadi, sifat tertutup merupakan syarat perlu bagi keberadaan suatu bentuk primitif. Tanpa bukti, kita catat bahwa untuk bentuk berderajat positif kriteria perlu ini juga cukup pada $\R^n$. Namun, ekuivalensi tersebut sama sekali tidak berlaku pada setiap manifold.

  \zwischenueberschrift{Setengah ruang Euklides}

\inputdefinition
{}
{

Dalam dimensi positif, yang dimaksud dengan \definitionswort {setengah ruang Euklides}{} dengan
\definitionswort {dimensi}{} $n$ adalah himpunan

\mavergleichskettedisp
{\vergleichskette
{H
}
{ =} { \left\{  x \in \R^n   \mid   x_1 \geq 0        \right\}
}
{ } {
}
{ } {
}
{ } {
}
}
{}{}{}
dengan
\definitionsverweis {topologi terinduksi}{}{.}

}
Untuk

\mavergleichskette
{\vergleichskette
{ n
}
{ = }{ 0
}
{  }{
}
{    }{
}
{   }{
}
}
{}{}{}
kita gunakan konvensi bahwa setengah ruang Euklides berupa satu titik; untuk

\mavergleichskette
{\vergleichskette
{ n
}
{ = }{ 1
}
{  }{
}
{    }{
}
{   }{
}
}
{}{}{}
himpunan ini berupa interval
\mathl{[0, \infty)}{;} untuk

\mavergleichskette
{\vergleichskette
{ n
}
{ = }{ 2
}
{  }{
}
{    }{
}
{   }{
}
}
{}{}{}
himpunan ini merupakan suatu \stichwort {setengah bidang} {;} dan untuk

\mavergleichskette
{\vergleichskette
{ n
}
{ = }{ 3
}
{  }{
}
{    }{
}
{   }{
}
}
{}{}{}
himpunan ini merupakan suatu setengah ruang. Jika indeks koordinat $1$ diganti dengan indeks lain, atau \anfuehrung{$\leq$}{} digunakan sebagai ganti \anfuehrung{$\geq$ }{,} objek-objek ini juga disebut setengah ruang. Karena suatu setengah ruang $H$ tertutup dalam $\R^n$, suatu himpunan bagian

\mavergleichskette
{\vergleichskette
{ T
}
{ \subseteq }{ H
}
{  }{
}
{    }{
}
{   }{
}
}
{}{}{}
tertutup di $H$ jika dan hanya jika ia tertutup di $\R^n$. Ekuivalensi ini
\betonung{tidak}{} berlaku untuk himpunan terbuka. Sebagai contoh, seluruh ruang $H$ terbuka di $H$, tetapi tidak terbuka di $\R^n$. Himpunan

\mavergleichskettedisp
{\vergleichskette
{ \partial H
}
{ =} { \left\{  x \in \R^n   \mid   x_1 = 0        \right\}
}
{ } {
}
{ } {
}
{ } {
}
}
{}{}{}
merupakan bagian dari $H$ dan disebut \stichwort {batas} {} dari $H$. Batas ini homeomorfik dengan $\R^{n-1}$
\zusatzklammer {yang untuk $n=0$ harus ditafsirkan sebagai himpunan kosong} {} {.}
Dengan $H_+$ kita nyatakan bagian positif

\mavergleichskette
{\vergleichskette
{ H_+
}
{ = }{ \left\{  x \in \R^n   \mid   x_1 >0        \right\}
}
{  }{
}
{    }{
}
{   }{
}
}
{}{}{,}
yang merupakan himpunan bagian terbuka dari $\R^n$.

Setengah ruang merupakan model standar bagi manifold berbatas yang akan kita perkenalkan sekarang. Konsep ini memperumum konsep manifold. Contoh khas manifold berbatas adalah bola pejal tertutup; batasnya adalah sfer. Suatu titik di interior bola mempunyai lingkungan berupa bola terbuka yang lebih kecil, sehingga di sekitar titik itu bentuknya \anfuehrung{secara lokal}{} seperti $\R^3$. Titik pada batas bola tidak mempunyai lingkungan seperti itu: batasnya hadir dalam setiap lingkungan terbuka titik tersebut. Secara lokal, titik batas demikian tampak seperti titik batas suatu setengah ruang Euklides.

Citra koordinat peta-peta suatu manifold berbatas berupa himpunan-himpunan terbuka di dalam setengah ruang. Karena itu, untuk membahas pemetaan transisi, kita harus dapat berbicara tentang pemetaan diferensiabel yang didefinisikan pada setengah ruang. Definisi berikut memungkinkan hal tersebut.

\inputdefinition
{}
{

Misalkan

\mavergleichskette
{\vergleichskette
{ U
}
{ \subseteq }{ H
}
{  }{
}
{    }{
}
{   }{
}
}
{}{}{}
suatu
\definitionsverweis {himpunan bagian terbuka}{}{}
di dalam
\definitionsverweis {setengah ruang Euklides}{}{}

\mavergleichskette
{\vergleichskette
{ H
}
{ \subseteq }{ \R^n
}
{  }{
}
{    }{
}
{   }{
}
}
{}{}{,}

\mavergleichskette
{\vergleichskette
{ P
}
{ \in }{ U
}
{  }{
}
{    }{
}
{   }{
}
}
{}{}{}
suatu titik, dan misalkan
\maabbdisp {\varphi} { U } {\R^m
} {}
suatu
\definitionsverweis {pemetaan}{}{.}
Maka $\varphi$ disebut \definitionswort {terdiferensialkan secara kontinu}{} di $P$ jika terdapat suatu lingkungan terbuka

\mavergleichskette
{\vergleichskette
{ P
}
{ \in }{ V
}
{ \subseteq }{ \R^n
}
{    }{
}
{   }{
}
}
{}{}{}
dan suatu fungsi
\definitionsverweis {terdiferensialkan secara kontinu}{}{}
\maabbdisp {\tilde{\varphi}} { V } {\R^m
} {}
dengan

\mavergleichskette
{\vergleichskette
{ \tilde{\varphi} {{|}}_{U \cap V}
}
{ = }{\varphi {{|}}_{U \cap V}
}
{  }{
}
{    }{
}
{   }{
}
}
{}{}{.}

}
Jadi, konsep keterdiferensialan yang baru direduksi ke konsep lama. Untuk suatu himpunan terbuka

\mavergleichskette
{\vergleichskette
{ U
}
{ \subseteq }{ H
}
{  }{
}
{    }{
}
{   }{
}
}
{}{}{,}
yang tidak memotong batas $H$, definisi ini ekuivalen dengan definisi bagi suatu himpunan terbuka di $\R^n$.

Melalui strategi mendefinisikan konsep-konsep pada titik batas dengan mensyaratkan adanya lingkungan terbuka dan objek yang diperluas, banyak konsep penting terbawa ke situasi baru yang lebih umum ini; kita tidak akan selalu menguraikannya satu per satu. Sebagai contoh, sudah jelas apa yang dimaksud dengan \stichwort {difeomorfisme} {} antara himpunan-himpunan terbuka di dalam setengah ruang, serta apa yang dimaksud dengan diferensial total suatu pemetaan diferensiabel. Dengan latar ini, definisi manifold berbatas juga tidak mengejutkan.


\chapter{Lembar Kerja 20}
\input{generated/worksheet20.id.build.tex}

\section*{Solusi yang disediakan oleh sumber}
\addcontentsline{toc}{section}{Solusi yang disediakan oleh sumber}
\subsection*{Solusi Soal 20.3}
\input{generated/worksheet20_exercise03_solution.id.build.tex}
\subsection*{Solusi Soal 20.4}
Kita peroleh

\mavergleichskettealign
{\vergleichskettealign
{ d \omega
}
{ =} { xe^{xz} dz \wedge d x \wedge dy - xzdy \wedge dx \wedge dz  -y \sin (xy)  \cos ( \cos (xy) )  dx \wedge dy \wedge dz
}
{ =} {  \left(  xe^{xz} +xz  -y \sin (xy)  \cos ( \cos (xy) )  \right) dx \wedge dy \wedge dz
}
{ } {
}
{ } {}
}
{}
{}{.}

\subsection*{Solusi Soal 20.5}
\input{generated/worksheet20_exercise05_solution.id.build.tex}
\subsection*{Solusi Soal 20.7}
\textit{Catatan edisi: solusi sumber ditulis dengan notasi ASCII yang rusak
dan menyebut $d\omega$ sebagai bentuk yang tertutup. Perhitungan berikut
menormalkan notasi, membuktikan bahwa $\omega$ tertutup, dan menampilkan
bentuk primitifnya secara eksplisit.}

Pertama,

\mavergleichskettealign
{\vergleichskettealign
{d\omega
}
{ =} { -\cos(y)\,dy\wedge dx-\cos(y)\,dx\wedge dy
}
{ =} { \cos(y)\,dx\wedge dy-\cos(y)\,dx\wedge dy
}
{ =} {0
}
{ } {}
}
{}
{}{.}

Jadi, $\omega$ tertutup. Selain itu,

\mavergleichskettedisp
{\vergleichskette
{\omega
}
{ =} {d\left(x^2-x\sin(y)\right)
}
{ } {
}
{ } {
}
{ } {
}
}
{}{}{,}

sehingga $\omega$ eksak.

\subsection*{Solusi Soal 20.14 (fragmen sumber; belum selesai)}
\textit{Catatan edisi: halaman sumber secara eksplisit berstatus belum
selesai. Halaman itu berhenti setelah menampilkan kandidat bentuk
berderajat $n-1$ dan sebuah rantai persamaan kosong. Jika setiap
$\partial_iH_i=h$, turunan kandidat yang ditampilkan sumber justru sama
dengan $-n\omega$, bukan $\omega$; sumber juga belum membuktikan bahwa
pilihan antiturunan pada irisan-irisan koordinat menghasilkan fungsi global
yang sesuai pada himpunan terbuka sembarang. Fragmen berikut dipertahankan
sebagai fragmen solusi sumber, bukan sebagai bukti lengkap.}

Kita dapat menuliskan bentuk tersebut sebagai

\mavergleichskettedisp
{\vergleichskette
{ \omega
}
{ =} { h(x_1 , \ldots , x_n) dx_1 \wedge \ldots \wedge dx_n
}
{ } {
}
{ } {
}
{ } {
}
}
{}{}{}
dengan suatu fungsi kontinu
\maabbdisp {h} { U } {\R
} {}
dengan
\mathl{x_1 , \ldots , x_n}{} menyatakan koordinat-koordinat pada $\R^n$.
Untuk setiap tupel tetap

\mathdisp {(x_1 , \ldots , x_{i-1},x_{i+1} , \ldots , x_n)} { , }

pembatasan fungsi $h$
pada satu komponen interval dari irisan koordinat hanya bergantung pada
$x_i$; kita nyatakan pembatasan ini dengan
\mathl{h_i(x_1 , \ldots , x_{i-1},-,x_{i+1} , \ldots , x_n)}{.}
Karena $h_i$ kontinu, pada komponen interval tersebut terdapat suatu
antiturunan
\mathl{H_i(x_1 , \ldots , x_{i-1},-,x_{i+1} , \ldots , x_n)}{.}
Turunan parsialnya terhadap $x_i$ sama dengan $h_i$. Sumber kemudian
menampilkan kandidat

\mathdisp {\sum_{i = 1}^n (-1)^{i} H_i dx_1 \wedge \ldots \wedge  dx_{i-1} \wedge d x_{i+1} \wedge \ldots \wedge  dx_n} { . }

\textit{Fragmen sumber berakhir di sini. Rantai persamaan kosong pada
halaman sumber tidak dimasukkan sebagai isi matematis.}


\part{Unit 21}
\chapter[Kuliah 21]{Kuliah 21: Manifold Berbatas dan Orientasi Batas}
\input{generated/lecture21.id.build.tex}

\chapter{Lembar Kerja 21}
\input{generated/worksheet21.id.build.tex}

\section*{Solusi yang disediakan oleh sumber}
\addcontentsline{toc}{section}{Solusi yang disediakan oleh sumber}
\subsection*{Solusi Soal 21.7}
\input{generated/worksheet21_exercise07_solution.id.build.tex}
\clearpage
\subsection*{Solusi Soal 21.10}
\input{generated/worksheet21_exercise10_solution.id.build.tex}
\subsection*{Solusi Soal 21.12}
Misalkan
\mathl{P \in N}{} dan
\mathl{P \in U}{} suatu daerah peta dari $M$. Misalkan
\maabbdisp {\alpha} { U } {V
} {}
suatu peta dengan himpunan terbuka
\mathl{V \subseteq H}{,}
\mathl{H=\R_{\geq 0} \times \R^{n-1}}{.} Peta ini menginduksi suatu homeomorfisme
\maabbdisp {} {N \cap U} { \alpha (N \cap U)
} {.}
Peta-peta inilah yang kita gunakan untuk $N$. Sifat difeomorfisme dari pergantian peta untuk $M$ langsung diwariskan oleh pergantian peta untuk $N$. Misalkan
\mathl{P \in N \cap \partial M}{.} Di bawah suatu peta, $P$ kemudian dipetakan ke suatu titik batas dari $H$. Karena peta-peta untuk $N$ merupakan pembatasan dari peta-peta semacam itu, sifat ini tetap berlaku. Sebaliknya, jika
\mathl{P \in \partial N}{}, maka tentu saja, di satu pihak,
\mathl{P \in N \subseteq M}{.} Di pihak lain, di sekitar $P$ terdapat suatu peta untuk $N$ yang diperoleh dengan membatasi suatu peta untuk $M$, dan dalam peta tersebut $P$ dipetakan ke suatu titik batas dari $H$.

\subsection*{Solusi Soal 21.13}
\input{generated/worksheet21_exercise13_solution.id.build.tex}

\part{Unit 22}
\chapter[Kuliah 22]{Kuliah 22: Partisi Satuan}
\setcounter{section}{22}

Sekarang kita membahas suatu teknik bantu analitik penting yang disebut \stichwort {partisi satuan} {.} Kita akan menggunakannya dalam pembuktian pernyataan bahwa manifold yang dapat diorientasikan mempunyai bentuk volume positif, dan dalam pembuktian teorema Stokes. Dalam kuliah ini kita akan mengonstruksi partisi satuan; untuk itu mula-mula kita memerlukan beberapa konsep topologis.

  \zwischenueberschrift{Penghabisan kompak}

\bild{ \begin{center}
\includegraphics[width=5.5cm]{ \bildeinlesungpng {Inner_point} {png} }
\end{center}
\bildtext {Interior terbuka adalah gabungan semua titik interior, yaitu titik-titik dari $T$
\zusatzklammer {pada gambar, $E$} {} {,}
yang memiliki suatu lingkungan terbuka yang seluruhnya termuat di dalam $T$.} }

\bildlizenz { Inner point.png } {} {Zasdfgbnm} {Commons} {PD} {}

\inputdefinition
{}
{

Misalkan $X$ suatu
\definitionsverweis {ruang topologis}{}{}
dan

\mavergleichskette
{\vergleichskette
{T
}
{ \subseteq }{ X
}
{  }{
}
{    }{
}
{   }{
}
}
{}{}{}
suatu himpunan bagian. Maka

\mavergleichskettedisp
{\vergleichskette
{ \overline{  T  }
}
{ =} {  \bigcap_{A \text{ tertutup}, \, T \subseteq A} A
}
{ } {
}
{ } {
}
{ } {
}
}
{}{}{}
disebut \definitionswort {penutupan}{}
\zusatzklammer {atau \definitionswort {penutupan topologis}{}} {} {}
dari $T$.

}
Untuk ruang metrik, kita sebelumnya telah memperkenalkan penutupan sebagai himpunan semua
\definitionsverweis {titik pelekatan}{}{}{;}
lihat khususnya
Soal 35.3 (Analisis (Osnabrück 2021--2023)).

\inputdefinition
{}
{

Misalkan $X$ suatu
\definitionsverweis {ruang topologis}{}{}
dan

\mavergleichskette
{\vergleichskette
{ T
}
{ \subseteq }{ X
}
{  }{
}
{    }{
}
{   }{
}
}
{}{}{}
suatu himpunan bagian. Maka

\mavergleichskettedisp
{\vergleichskette
{ T ^{o}
}
{ =} { \bigcup_{U \text{ terbuka}, \, U \subseteq T}   U
}
{ } {
}
{ } {
}
{ } {
}
}
{}{}{}
disebut \definitionswort {interior terbuka}{}
\zusatzklammer {atau \stichwort {interior} {}} {} {}
dari $T$.

}
Kedua konsep ini berhubungan melalui

\mavergleichskettedisp
{\vergleichskette
{ \overline{ T }
}
{ =} { X \setminus (X \setminus T)^{o}
}
{ } {
}
{ } {
}
{ } {
}
}
{}{}{.}
Konsep batas juga terbawa dari situasi metrik ke ruang topologis sembarang.

\inputdefinition
{
}
{

Misalkan $X$ suatu
\definitionsverweis {ruang topologis}{}{}
dan

\mavergleichskette
{\vergleichskette
{ T
}
{ \subseteq }{ X
}
{  }{
}
{    }{
}
{   }{
}
}
{}{}{}
suatu himpunan bagian. Yang dimaksud dengan
\definitionswort {batas}{}
dari $T$ adalah himpunan

\mavergleichskettedisp
{\vergleichskette
{ \operatorname{Batas} \, ( T )
}
{ =} { \overline{  T  }  \setminus T ^{o}
}
{ } {
}
{ } {
}
{ } {
}
}
{}{}{.}

}
Perhatikan bahwa batas topologis ini merupakan konsep yang berbeda dari batas suatu manifold berbatas. Meskipun demikian, terdapat hubungan yang dibahas dalam
Soal 22.1.

\inputdefinition
{}
{

Misalkan $X$ suatu
\definitionsverweis {ruang topologis}{}{}
dan
\maabbdisp {f} { X } {\R
} {}
suatu
\definitionsverweis {fungsi}{}{.}
Maka
\definitionsverweis {penutupan topologis}{}{}

\mathdisp {\overline{  \left\{  x \in X  \mid  f(x) \neq 0        \right\}   }} {  }

disebut \definitionswort {tumpuan}{} dari $f$.

}

\inputdefinition
{}
{

Misalkan $X$ suatu
\definitionsverweis {ruang topologis}{}{.}
Suatu \definitionswort {penghabisan kompak}{}

\mathbed {A_n} {}
 {n \in \N} {}
 {} {} {} {,}
dari $X$ adalah suatu
\definitionsverweis {barisan}{}{}
\definitionsverweis {himpunan-himpunan bagian kompak}{}{}

\mavergleichskette
{\vergleichskette
{ A_n
}
{ \subseteq }{ X
}
{  }{
}
{    }{
}
{   }{
}
}
{}{}{}
dengan

\mathdisp {A_n \subseteq A_{n+1}^{o} \text{ dan } \bigcup_{n=0}^\infty A_n = X} { . }

}

Ruang $\R^d$ dihabiskan secara kompak oleh bola-bola tertutup

\mathbed {B  \left( 0,n \right)} {}
 {n \in \N} {}
 {} {} {} {.}
Lihat
Soal 22.6.

  \zwischenueberschrift{Partisi satuan}



\inputfaktbeweis
{Mannigfaltigkeit/Abzählbare Basis/Kompakte Ausschöpfung/Fakt}
{Lema}
{}
{

\faktsituation {Misalkan $M$ suatu
\definitionsverweis {manifold}{}{}}
\faktvoraussetzung {dengan suatu
\definitionsverweis {basis terhitung bagi topologinya}{}{.}}
\faktfolgerung {Maka $M$ mempunyai suatu
\definitionsverweis {penghabisan kompak}{}{.}}
\faktzusatz {}
\faktzusatz {}

}
{

Untuk setiap titik

\mavergleichskette
{\vergleichskette
{ P
}
{ \in }{ M
}
{  }{
}
{    }{
}
{   }{
}
}
{}{}{}
terdapat suatu lingkungan koordinat terbuka

\mavergleichskette
{\vergleichskette
{ P
}
{ \in }{ U
}
{  }{
}
{    }{
}
{   }{
}
}
{}{}{,}
\maabbdisp {\alpha} { U } { V
} {}
serta lingkungan-lingkungan bola

\mavergleichskettedisp
{\vergleichskette
{ U \left(  \alpha(P), \epsilon  \right)
}
{ \subseteq} { B  \left( \alpha(P), \epsilon \right)
}
{ \subset} { V
}
{ } {
}
{ } {
}
}
{}{}{.}
Karena pemetaan koordinat tersebut merupakan
\definitionsverweis {homeomorfisme}{}{}
dan bola tertutup bersifat kompak, maka

\mavergleichskette
{\vergleichskette
{ B_P
}
{ = }{ \alpha^{-1} \left(  B  \left( \alpha(P), \epsilon \right)  \right)
}
{  }{
}
{    }{
}
{   }{
}
}
{}{}{}
merupakan suatu
\definitionsverweis {himpunan bagian kompak}{}{}
dari $M$ yang memuat
\definitionsverweis {lingkungan terbuka}{}{}

\mavergleichskette
{\vergleichskette
{ U_P
}
{ = }{ \alpha^{-1} \left(  U \left(  \alpha(P), \epsilon  \right)  \right)
}
{  }{
}
{    }{
}
{   }{
}
}
{}{}{}
dari $P$. Himpunan-himpunan

\mathbed {U_P} {}
 {P \in M} {}
 {} {} {} {,}
membentuk suatu
\definitionsverweis {tutupan terbuka}{}{}
dari $M$, sehingga menurut
Soal 2.8 (Teori Ukuran dan Integrasi (Osnabrück 2022--2023))
terdapat suatu subtutupan terhitung. Nyatakan subtutupan ini dengan

\mathbed {U_n} {}
 {n \in \N} {}
 {} {} {} {,}
\zusatzklammer {dengan $U_n$ termuat di dalam himpunan kompak $B_n$} {} {.}
Sekarang kita definisikan secara
\definitionsverweis {rekursif}{}{}
suatu pemetaan
\definitionsverweis {naik secara monoton}{}{}
\maabbeledisp {} {\N} {\N
} { k } {n_k
} {,}
sedemikian sehingga

\mathbeddisp {A_k = \bigcup_{n=0}^{n_k } B_n} {}
 {k \in \N} {}
 {} {} {} {,}
merupakan suatu
\definitionsverweis {penghabisan kompak}{}{}
dari $M$. Sebagai gabungan berhingga dari himpunan-himpunan kompak, setiap $A_k$ bersifat kompak. Kita mulai dengan

\mavergleichskette
{\vergleichskette
{ n_0
}
{ = }{ 0
}
{  }{
}
{    }{}
{   }{}
}
{}{}{.}
Misalkan $n_k$ sudah dikonstruksi. Himpunan

\mathdisp {A_k \cup B_{n_k +1}} {  }

bersifat
\definitionsverweis {kompak}{}{}
dan karena itu ditutupi oleh berhingga banyak himpunan terbuka
\mathl{\bigcup_{n=0}^{n_{k+1} } U_n}{,} dengan memilih

\mavergleichskette
{\vergleichskette
{ n_{k+1}
}
{ \geq }{ n_k +1
}
{  }{
}
{    }{
}
{   }{
}
}
{}{}{.}
Dengan pilihan ini,

\mavergleichskettedisp
{\vergleichskette
{ A_k
}
{ \subseteq} { \bigcup_{n = 0}^{n_{k+1} } U_n
}
{ \subseteq} { \bigcup_{n = 0}^{n_{k+1} } B_n
}
{ =} { A_{k+1}
}
{ } {
}
}
{}{}{,}
dan barisan ini membentuk suatu penghabisan karena himpunan-himpunan

\mathbed {U_n} {}
 {n \in \N} {}
 {} {} {} {,}
membentuk suatu tutupan terbuka dari $M$.

}

\bild{ \begin{center}
\includegraphics[width=5.5cm]{ \bildeinlesungsvg {Partition_of_unity_illustration} {svg} }
\end{center}
\bildtext {Ilustrasi beberapa fungsi taknegatif yang jumlahnya selalu 1.} }

\bildlizenz { Partition of unity illustration.svg } {} {Oleg Alexandrov} {Commons} {PD} {}

\inputdefinition
{}
{

Misalkan $X$ suatu
\definitionsverweis {ruang topologis}{}{.}
Suatu keluarga fungsi
\maabbdisp {h_j} { X } {\R
} {}
dengan

\mavergleichskette
{\vergleichskette
{ j
}
{ \in }{ J
}
{  }{
}
{    }{
}
{   }{
}
}
{}{}{}
disebut suatu \definitionswort {partisi satuan}{,} jika sifat-sifat berikut berlaku.
\aufzaehlungdrei{Berlaku

\mavergleichskette
{\vergleichskette
{ h_j(X)
}
{ \subseteq }{ [0,1]
}
{  }{
}
{    }{
}
{   }{
}
}
{}{}{}
untuk setiap

\mavergleichskette
{\vergleichskette
{ j
}
{ \in }{ J
}
{  }{
}
{    }{
}
{   }{
}
}
{}{}{.}
}{Setiap titik

\mavergleichskette
{\vergleichskette
{ P
}
{ \in }{ X
}
{  }{
}
{    }{
}
{   }{
}
}
{}{}{}
mempunyai suatu
\definitionsverweis {lingkungan terbuka}{}{}

\mavergleichskette
{\vergleichskette
{ P
}
{ \in }{ U
}
{  }{
}
{    }{
}
{   }{
}
}
{}{}{}
sedemikian sehingga
\definitionsverweis {fungsi-fungsi hasil pembatasan}{}{}

\mathl{h_j {{|}}_{ U }}{} semuanya, kecuali berhingga banyak, merupakan fungsi nol.
}{Berlaku

\mavergleichskette
{\vergleichskette
{  \sum_{j \in J} h_j
}
{ = }{ 1
}
{  }{
}
{    }{
}
{   }{
}
}
{}{}{.}
} Jika semua $h_j$ \definitionsverweis {kontinu}{}{,} partisi tersebut disebut suatu \definitionswort {partisi satuan kontinu}{.}

}

Sifat kedua memastikan bahwa jumlah dalam (3) terdefinisi, sebab untuk setiap titik

\mavergleichskette
{\vergleichskette
{ P
}
{ \in }{ X
}
{  }{
}
{    }{
}
{   }{
}
}
{}{}{}
dan bagi semua indeks

\mavergleichskette
{\vergleichskette
{ j
}
{ \in }{ J
}
{  }{
}
{    }{
}
{   }{
}
}
{}{}{}
kecuali berhingga banyak, berlaku kesamaan

\mavergleichskette
{\vergleichskette
{  h_j(P)
}
{ = }{ 0
}
{  }{
}
{    }{
}
{   }{
}
}
{}{}{.}
Pada suatu manifold, partisi seperti itu disebut \stichwort {diferensiabel} {,} jika semua $h_j$ merupakan
\definitionsverweis {fungsi-fungsi diferensiabel}{}{.}

\inputdefinition
{}
{

Misalkan

\mavergleichskette
{\vergleichskette
{ X
}
{ = }{ \bigcup_{i \in I} W_i
}
{  }{
}
{    }{
}
{   }{
}
}
{}{}{}
suatu
\definitionsverweis {tutupan terbuka}{}{}
dari suatu
\definitionsverweis {ruang topologis}{}{}
$X$. Suatu
\definitionsverweis {partisi satuan}{}{}
\maabbdisp {h_j} { X } {\R
} {}
dengan

\mavergleichskette
{\vergleichskette
{ j
}
{ \in }{ J
}
{  }{
}
{    }{
}
{   }{
}
}
{}{}{}
disebut suatu \definitionswort {partisi satuan subordinat terhadap tutupan}{,} jika untuk setiap

\mavergleichskette
{\vergleichskette
{ j
}
{ \in }{ J
}
{  }{
}
{    }{
}
{   }{
}
}
{}{}{}
terdapat suatu himpunan terbuka
\mathl{W_{i(j)}}{} dari tutupan tersebut sedemikian sehingga
\definitionsverweis {tumpuan}{}{}
dari $h_j$ termuat di dalam
\mathl{W_{i(j)}}{}{.}

}



\inputfaktbeweis
{Mannigfaltigkeit/Abzählbare Basis/Lokal endlicher Atlas/Fakt}
{Lema}
{}
{

\faktsituation {Misalkan $M$ suatu
\definitionsverweis {manifold}{}{}
dengan suatu
\definitionsverweis {basis terhitung}{}{}
bagi topologinya. Misalkan

\mavergleichskette
{\vergleichskette
{ M
}
{ = }{ \bigcup_{i \in I} W _i
}
{  }{
}
{    }{
}
{   }{
}
}
{}{}{}
suatu
\definitionsverweis {tutupan terbuka}{}{}
dari $M$.}
\faktfolgerung {Maka terdapat suatu
\definitionsverweis {atlas kompatibel}{}{}
yang terhitung

\mathbed {\left(  U_j, \alpha_j, V_j  \right)} {}
 {j \in J} {}
 {} {} {} {,}
dengan lingkungan-lingkungan bola

\mavergleichskettedisp
{\vergleichskette
{ U \left(   0 , \delta_j   \right)
}
{ \subset} { B  \left(  0 , \epsilon_j  \right)
}
{ \subset} { V_j
}
{ } {
}
{ } {
}
}
{}{}{}
\zusatzklammer {di sini,

\mavergleichskettek
{\vergleichskettek
{ 0
}
{ \in }{ V_j
}
{ \subseteq }{ \R^n
}
{    }{
}
{   }{
}
}
{}{}{}
dan

\mavergleichskettek
{\vergleichskettek
{ \delta_j
}
{ < }{ \epsilon_j
}
{  }{
}
{    }{
}
{   }{
}
}
{}{}{}} {} {}
sedemikian sehingga untuk setiap

\mavergleichskette
{\vergleichskette
{ j
}
{ \in }{ J
}
{  }{
}
{    }{
}
{   }{
}
}
{}{}{}
terdapat suatu
\mathl{W_{i(j)}}{} dengan

\mavergleichskette
{\vergleichskette
{ U_j
}
{ \subseteq }{ W_{i(j)}
}
{  }{
}
{    }{
}
{   }{
}
}
{}{}{}; selain itu, $M$ ditutupi oleh

\mathbed {\alpha_j^{-1} \left(  U \left(   0 , \delta_j   \right)  \right)} {}
 {j \in J} {}
 {} {} {} {,}
dan setiap titik

\mavergleichskette
{\vergleichskette
{ P
}
{ \in }{ M
}
{  }{
}
{    }{
}
{   }{
}
}
{}{}{}
terletak hanya di dalam berhingga banyak himpunan $U_j$.}
\faktzusatz {}
\faktzusatz {}

}
{

\teilbeweis {}{}{}
{Misalkan tutupan terbuka

\mathbed {W_i} {}
 {i \in I} {}
 {} {} {} {,}
diberikan. Selanjutnya, misalkan

\mathbed {A_n} {}
 {n \in \N} {}
 {} {} {} {,}
suatu
\definitionsverweis {penghabisan kompak}{}{}
dari $M$, yang keberadaannya dijamin oleh
Lema 22.6.
Himpunan-himpunan terbuka
\mathl{A_{n+2}^{o}  \setminus A_{n-1}}{} juga membentuk suatu tutupan terbuka, sebab untuk setiap titik

\mavergleichskette
{\vergleichskette
{ P
}
{ \in }{ M
}
{  }{
}
{    }{
}
{   }{
}
}
{}{}{}
terdapat suatu nilai minimal

\mavergleichskette
{\vergleichskette
{ n
}
{ \in }{ \N
}
{  }{
}
{    }{
}
{   }{
}
}
{}{}{}
dengan

\mavergleichskette
{\vergleichskette
{ P
}
{ \in }{ A_n
}
{  }{
}
{    }{
}
{   }{
}
}
{}{}{}
\zusatzklammer {tetapkan

\mavergleichskettek
{\vergleichskettek
{ A_{-1}
}
{ = }{ \emptyset
}
{  }{
}
{    }{
}
{   }{
}
}
{}{}{}} {} {.}
Untuk $n$ ini berlaku
\mathkor {} {P \not\in A_{n-1}} {dan} {P \in A_n \subseteq A_{n+1} ^{o}} {.}
Dengan meninjau irisan-irisan
\mathl{W_i \cap \left(  A_{n+2}^{o}  \setminus A_{n-1}  \right)}{}{,} kita dapat menganggap bahwa setiap himpunan dari tutupan termuat di dalam suatu
\mathl{A_{n+2}^{o} \setminus A_{n-1}}{}{.}}
{}
\teilbeweis {}{}{}
{Untuk setiap titik

\mavergleichskette
{\vergleichskette
{ P
}
{ \in }{ M
}
{  }{
}
{    }{
}
{   }{
}
}
{}{}{}
terdapat suatu lingkungan koordinat terbuka
\zusatzklammer {yang kompatibel} {} {}

\mavergleichskette
{\vergleichskette
{ P
}
{ \in }{ U_P
}
{  }{
}
{    }{
}
{   }{
}
}
{}{}{,}
yang termuat di dalam salah satu $W_i$ dan yang mempunyai lingkungan-lingkungan bola

\mavergleichskettedisp
{\vergleichskette
{ U \left(   0 , \delta_P   \right)
}
{ \subset} { B  \left(  0 , \epsilon_P  \right)
}
{ \subset} { V_P
}
{ } {
}
{ } {
}
}
{}{}{}
dengan

\mavergleichskette
{\vergleichskette
{ P
}
{ \in }{ \alpha_P^{-1} \left(  U \left(   0 , \delta_P   \right)  \right)
}
{  }{
}
{    }{
}
{   }{
}
}
{}{}{}
dan

\mavergleichskette
{\vergleichskette
{ \delta_P
}
{ < }{ \epsilon_P
}
{  }{
}
{    }{
}
{   }{
}
}
{}{}{.}
Himpunan-himpunan

\mathbed {\alpha_P^{-1} \left(  U \left(   0 , \delta_P   \right)  \right)} {}
 {P \in M} {}
 {} {} {} {,}
kemudian juga membentuk suatu tutupan terbuka dari $M$. Menurut
Soal 2.8 (Teori Ukuran dan Integrasi (Osnabrück 2022--2023))
kita dapat beralih ke suatu subtutupan terhitung. Jadi, kita dapat menganggap bahwa suatu sistem domain peta

\mathbed {U_j} {}
 {j \in \N} {}
 {} {} {} {,}
bersama dengan lingkungan-lingkungan bola

\mavergleichskettedisp
{\vergleichskette
{ U \left(   0 , \delta_j   \right)
}
{ \subset} { B  \left(  0 , \epsilon_j  \right)
}
{ \subset} { V_j
}
{ } {
}
{ } {
}
}
{}{}{}
diberikan sedemikian sehingga

\mathbed {\alpha_j^{-1} \left(  U \left(   0 , \delta_j   \right)  \right)} {}
 {j \in \N} {}
 {} {} {} {,}
juga merupakan suatu tutupan terbuka dari $M$, setiap $U_j$ termuat di dalam suatu
\mathl{W_{i(j)}}{}{,} dan hubungan dengan penghabisan kompak yang diuraikan di atas berlaku.}
{}
\teilbeweis {}{Kita akan mendefinisikan suatu himpunan bagian

\mavergleichskette
{\vergleichskette
{ J
}
{ \subseteq }{ \N
}
{  }{
}
{    }{
}
{   }{
}
}
{}{}{}
sedemikian sehingga keluarga

\mathbed {U_j} {}
 {j \in J} {}
 {} {} {} {,}
juga memenuhi sifat keberhinggaan tersebut.\leerzeichen{}}{}
{Mula-mula, kekompakan $A_0$ memberikan suatu himpunan indeks berhingga
\mathl{J_{-1}}{} sedemikian sehingga himpunan-himpunan

\mathbed {\alpha_j^{-1} \left(  U \left(   0 , \delta_j   \right)  \right)} {}
 {j \in J_{-1}} {}
 {} {} {} {,}
menutupi $A_0$. Untuk

\mavergleichskette
{\vergleichskette
{ n
}
{ \in }{ \N
}
{  }{
}
{    }{
}
{   }{
}
}
{}{}{}
tinjau himpunan kompak
\mathl{A_{n+1} \setminus A_n ^{o}}{.} Himpunan ini ditutupi oleh berhingga banyak himpunan

\mathbed {\alpha_j^{-1} \left(  U \left(   0 , \delta_j   \right)  \right)} {}
 {j \in \N} {}
 {} {} {} {,}
dan untuk itu hanya diperlukan indeks-indeks $j$ sedemikian sehingga $U_j$ termuat di dalam
\mathl{A_{n+2}^{o}  \setminus A_{n-1}}{}{.} Nyatakan himpunan indeks berhingga yang bersangkutan dengan $J_n$, dan tetapkan

\mavergleichskette
{\vergleichskette
{ J
}
{ = }{ J_{-1} \cup \bigcup_{n \in \N} J_n
}
{  }{
}
{    }{
}
{   }{
}
}
{}{}{.}
Maka setiap $A_n$ hanya memotong berhingga banyak himpunan

\mathbed {U_j} {}
 {j \in J} {}
 {} {} {} {.}}
{}

}




\inputfaktbeweis
{Mannigfaltigkeit/Abzählbare Basis/Überdeckung/Partition/Fakt}
{Teorema}
{}
{

\faktsituation {Misalkan $M$ suatu
\definitionsverweis {manifold diferensiabel}{}{}
dengan suatu
\definitionsverweis {basis terhitung}{}{}
bagi topologinya.}
\faktfolgerung {Maka untuk setiap tutupan terbuka terdapat suatu
\definitionsverweis {partisi satuan}{}{}
yang \definitionsverweis {terdiferensialkan secara kontinu}{}{}
dan subordinat terhadap tutupan tersebut.}
\faktzusatz {}
\faktzusatz {}

}
{

Menurut
Lema 22.9
kita dapat menganggap bahwa suatu tutupan terbuka terdiri atas domain-domain peta

\mathbed {U_j} {}
 {j \in J} {}
 {} {} {} {,}
\zusatzklammer {$J$ terhitung} {} {}
dengan
\maabbdisp {\alpha_j} {U_j } { V_j
} {}
dan dengan lingkungan-lingkungan bola

\mavergleichskettedisp
{\vergleichskette
{ U \left(   0 , \delta_j   \right)
}
{ \subset} { B  \left(  0 , \epsilon_j  \right)
}
{ \subset} { V_j
}
{ } {
}
{ } {
}
}
{}{}{}
\zusatzklammer {dengan

\mavergleichskettek
{\vergleichskettek
{ \delta_j
}
{ < }{ \epsilon_j
}
{  }{
}
{    }{
}
{   }{
}
}
{}{}{}} {} {}
sedemikian sehingga himpunan-himpunan
\mathl{\alpha_j^{-1} \left(  U \left(   0 , \delta_j   \right)  \right)}{} juga membentuk suatu tutupan dari $M$ dan setiap titik

\mavergleichskette
{\vergleichskette
{ P
}
{ \in }{ M
}
{  }{
}
{    }{
}
{   }{
}
}
{}{}{}
termuat hanya dalam berhingga banyak himpunan $U_j$ dan, khususnya, hanya dalam berhingga banyak himpunan
\mathl{\alpha_j^{-1} \left(  U \left(   0 , \delta_j   \right)  \right)}{}{.}
Pada $V_j$, tinjau fungsi $g_j$ yang didefinisikan oleh

\mavergleichskettedisp
{\vergleichskette
{ g_j (v)
}
{ =} { \begin{cases}  e^{ - { \frac{   1  }{  \left(  \delta_j^2- \Vert { v } \Vert^2  \right)^2 } } }  \text{ untuk } \Vert { v } \Vert < \delta_j \, , \\ 0 \text{ selainnya} \, , \end{cases}
}
{ } {
}
{ } {
}
{ } {
}
}
{}{}{}
Fungsi ini bernilai positif tepat pada
\mathl{U \left(   0 , \delta_j   \right)}{} dan
\definitionsverweis {tumpuan}{}{}
fungsi ini adalah
\mathl{B  \left(  0 , \delta_j  \right)}{.}
Peninjauan pada kedua himpunan bagian terbuka
\zusatzklammer {yang menutupi $V_j$} {} {}
\mathkor {} {U \left(   0 , \epsilon_j   \right)} {dan} {V_j \setminus B  \left(  0 , \delta_j  \right)} {}
menunjukkan bahwa $g_j$ terdiferensialkan tak hingga kali. Kita mendefinisikan suatu fungsi
\maabbdisp {\tilde{g}_j} { M } {\R
} {}
melalui

\mavergleichskettedisp
{\vergleichskette
{ \tilde{g}_j(x)
}
{ =} { \begin{cases}  g_j( \alpha_j(x)) \text{ untuk }  x \in \alpha_j^{-1}( U \left(   0 , \delta_j   \right) ) \, , \\  0 \text{ selainnya} \, . \end{cases}
}
{ } {
}
{ } {
}
{ } {
}
}
{}{}{}
Fungsi ini
\definitionsverweis {terdiferensialkan secara kontinu}{}{}
pada $M$, sebab \anfuehrung{pita}{}
\mathl{B  \left(  0 , \epsilon_j  \right) \setminus U \left(   0 , \delta_j   \right)}{} memungkinkan suatu transisi mulus. Kita tetapkan

\mavergleichskettedisp
{\vergleichskette
{ \tilde{g}(x)
}
{ \defeq} { \sum_{j \in J} \tilde{g}_j(x)
}
{ } {
}
{ } {
}
{ } {
}
}
{}{}{,}
yang merupakan jumlah berhingga untuk setiap titik, sebab
\definitionsverweis {tumpuan}{}{}
dari $\tilde{g}_j$ termuat di dalam

\mavergleichskettedisp
{\vergleichskette
{ \alpha_j^{-1} \left(  B  \left(  0 , \delta_j  \right)  \right)
}
{ \subset} { U_j
}
{ } {
}
{ } {
}
{ } {
}
}
{}{}{.}
Fungsi ini
\definitionsverweis {terdiferensialkan secara kontinu}{}{}
pada $M$ dan positif di setiap titik, sebab
\mathl{\tilde{g}_j(x)}{} masing-masing bernilai positif pada himpunan yang bersesuaian dari tutupan
\mathl{\alpha_j^{-1} \left(  U \left(   0 , \delta_j   \right)  \right)}{}{.}
Dengan demikian, fungsi-fungsi

\mavergleichskettedisp
{\vergleichskette
{ h_j
}
{ =} { { \frac{   \tilde{g}_j  }{ \tilde{g}  } }
}
{ } {
}
{ } {
}
{ } {
}
}
{}{}{}
membentuk partisi satuan yang dicari.

}

  \zwischenueberschrift{Orientasi pada manifold dan bentuk volume}

Dengan bantuan partisi satuan, sekarang kita dapat membuktikan kebalikan dari
Lema 15.5.



\inputfaktbeweis
{Nullstellenfreie Volumenform/Orientierte (Karten) Mannigfaltigkeit/Äquivalenz/Fakt}
{Teorema}
{}
{

\faktsituation {Misalkan $M$ suatu
\definitionsverweis {manifold diferensiabel
}{}{}
dengan suatu
\definitionsverweis {basis terhitung bagi topologinya}{}{.}}
\faktfolgerung {Maka terdapat suatu
\definitionsverweis {bentuk volume}{}{}
\definitionsverweis {kontinu}{}{}
tanpa titik nol pada $M$ jika dan hanya jika $M$
\definitionsverweis {dapat diorientasikan}{}{.}}
\faktzusatz {Bentuk volume ini juga dapat dipilih agar terdiferensialkan secara kontinu dan positif.}
\faktzusatz {}

}
{

\teilbeweis {}{}{}
{Salah satu arah sudah dibuktikan
dalam Lema 15.5.}
{}
\teilbeweis {}{}{}
{Sebaliknya, misalkan $M$
\definitionsverweis {dapat diorientasikan}{}{}
dan diberikan suatu
\definitionsverweis {atlas berorientasi}{}{}
yang
\definitionsverweis {terhitung}{}{}

\mathbed {\left(  U_i,\alpha_i,V_i  \right)} {}
 {i \in I} {}
 {} {} {} {,}
dari $M$. Di sini,

\mavergleichskette
{\vergleichskette
{ V_i
}
{ \subseteq }{ \R^n
}
{  }{
}
{    }{
}
{   }{
}
}
{}{}{}
merupakan himpunan
\definitionsverweis {terbuka}{}{,}
dan
\definitionsverweis {koordinat-koordinat}{}{}

\mathl{x_1 , \ldots , x_n}{} mendefinisikan suatu
\definitionsverweis {bentuk volume}{}{}
kontinu tanpa titik nol
\zusatzklammer {bahkan terdiferensialkan tak hingga kali} {} {}

\mathl{dx_1 \wedge \ldots \wedge dx_n}{} pada $V_i$. Kita tetapkan

\mavergleichskettedisp
{\vergleichskette
{ \omega_i
}
{ =} { \alpha_i^* \left( dx_1 \wedge \ldots \wedge dx_n \right)
}
{ } {
}
{ } {
}
{ } {
}
}
{}{}{}
dan dengan demikian memperoleh suatu bentuk volume tanpa titik nol pada $U_i$, yang kita perluas dengan $0$ di luar $U_i$\zusatzfussnote {Perluasan ini tentu saja tidak kontinu, tetapi hal itu tidak menjadi masalah untuk pembahasan berikut} {.} {.}

Sekarang, misalkan

\mathbed {h_j} {}
 {j \in J} {}
 {} {} {} {,}
suatu keluarga fungsi yang subordinat terhadap tutupan

\mathbed {U_i} {}
 {i \in I} {}
 {} {} {} {,}
dan terdiferensialkan secara kontinu, serta membentuk suatu
\definitionsverweis {partisi satuan}{}{,}
yang keberadaannya dijamin oleh
Teorema 22.10.
Jadi, khususnya, untuk setiap $j$ terdapat suatu
\mathl{i(j)}{} sedemikian sehingga
\definitionsverweis {tumpuan}{}{}
dari $h_j$ termuat di dalam
\mathl{U_{i(j)}}{}{.}
Oleh karena itu,
\mathl{h_j \omega_{i(j)}}{} merupakan suatu
\definitionsverweis {bentuk diferensial}{}{}
kontinu berderajat $n$ pada $M$. Kita tetapkan

\mavergleichskettedisp
{\vergleichskette
{ \omega
}
{ =} { \sum_{j \in J} h_j \omega_{i(j)}
}
{ } {
}
{ } {
}
{ } {
}
}
{}{}{.}
Di setiap titik

\mavergleichskette
{\vergleichskette
{ P
}
{ \in }{ M
}
{  }{
}
{    }{
}
{   }{
}
}
{}{}{}
jumlah ini berhingga dan karena itu mendefinisikan suatu bentuk diferensial kontinu berderajat $n$ pada $M$. Untuk suatu titik

\mavergleichskette
{\vergleichskette
{ P
}
{ \in }{ M
}
{  }{
}
{    }{
}
{   }{
}
}
{}{}{}
dan suatu
\definitionsverweis {basis}{}{}

\mathl{v_1 , \ldots , v_n}{} dari
\mathl{T_PM}{} yang merepresentasikan orientasi, berlaku

\mavergleichskettedisp
{\vergleichskette
{ \omega(P; v_1 , \ldots , v_n)
}
{ =} { \sum_{j \in J} h_j(P) \omega_{i(j)} \left(  P; v_1 , \ldots , v_n  \right)
}
{ } {
}
{ } {
}
{ } {
}
}
{}{}{.}
Terdapat suatu $j$ dengan

\mavergleichskette
{\vergleichskette
{ h_j(P)
}
{ > }{ 0
}
{  }{
}
{    }{
}
{   }{
}
}
{}{}{,}
dan untuk $j$ ini juga berlaku

\mavergleichskette
{\vergleichskette
{ \omega_{i(j)}(P; v_1 , \ldots , v_n)
}
{ > }{ 0
}
{  }{
}
{    }{
}
{   }{
}
}
{}{}{,}
sebab

\mavergleichskette
{\vergleichskette
{ P
}
{ \in }{ U_{i(j)}
}
{  }{
}
{    }{
}
{   }{
}
}
{}{}{}
berlaku; dengan demikian, bentuk ini positif di setiap titik.}
{}

}


\chapter{Lembar Kerja 22}
\setcounter{section}{22}

  \zwischenueberschrift{Soal latihan}

\inputaufgabe
{}
{

Misalkan $M$ suatu
\definitionsverweis {manifold berbatas}{}{}
dan
\mathl{\partial M}{} merupakan batasnya. Tunjukkan bahwa
\definitionsverweis {batas topologis}{}{}
dari
\mathl{M \setminus \partial M}{} sama dengan $\partial M$.

}
{} {}

\inputaufgabe
{}
{

Tentukan
\definitionsverweis {tumpuan}{}{} fungsi-fungsi berikut dari $\R$ ke $\R$.
\aufzaehlungsieben{Suatu fungsi polinomial.
}{Fungsi sinus.
}{Fungsi eksponensial.
}{Fungsi indikator $e_{ \Z }$.
}{Fungsi indikator $e_{ \Q }$.
}{Fungsi indikator $e_{  [a,b]  }$.
}{Fungsi indikator $e_{ ]a,b[ }$.
}

}
{} {}

\inputaufgabe
{}
{

Misalkan
\maabbdisp {f} { X } {\R
} {}
suatu fungsi pada
\definitionsverweis {ruang topologis}{}{}
$X$. Andaikan terdapat suatu
\definitionsverweis {himpunan bagian terbuka}{}{}

\mavergleichskette
{\vergleichskette
{U
}
{ \subseteq }{ X
}
{  }{
}
{    }{
}
{   }{
}
}
{}{}{,}
yang memuat
\definitionsverweis {tumpuan}{}{}
dari $f$. Tunjukkan bahwa $f$
\definitionsverweis {kontinu}{}{}
jika dan hanya jika pembatasan
\mathl{f \vert_U}{} kontinu.

}
{} {}

\inputaufgabe
{}
{

Misalkan $X$ suatu
\definitionsverweis {ruang topologis}{}{}
dan

\mavergleichskette
{\vergleichskette
{ T
}
{ \subseteq }{ X
}
{  }{
}
{    }{
}
{   }{
}
}
{}{}{}
suatu himpunan bagian. Tunjukkan bahwa
\definitionsverweis {penutupan}{}{}
dari $T$ sama dengan
\definitionsverweis {tumpuan}{}{}
dari
\definitionsverweis {fungsi indikator}{}{}

\mathl{e_{  T  }}{}.

}
{} {}

\inputaufgabe
{}
{

Berikan suatu
\definitionsverweis {penghabisan kompak}{}{} untuk
\definitionsverweis {bilangan riil}{}{} $\R$.

}
{} {}

\inputaufgabegibtloesung
{}
{

Berikan suatu
\definitionsverweis {penghabisan kompak}{}{}
untuk $\R^d$.

}
{} {}

\inputaufgabe
{}
{

Misalkan $X$ suatu
\definitionsverweis {ruang topologis}{}{.}
Untuk tutupan atas $X$ yang hanya terdiri atas $X$, tentukan suatu
\definitionsverweis {partisi satuan subordinat terhadap tutupan tersebut}{}{.}

}
{} {}

\inputaufgabe
{}
{

Misalkan $X$ suatu
\definitionsverweis {ruang topologis}{}{}
dan

\mavergleichskette
{\vergleichskette
{ X
}
{ = }{ \bigcup_{i \in I} U_i
}
{  }{
}
{    }{
}
{   }{
}
}
{}{}{}
suatu
\definitionsverweis {tutupan terbuka}{}{.}
Kita tinjau keluarga
\definitionsverweis {fungsi indikator}{}{}

\mathbeddisp {e_{  P  }} {}
 {P \in X} {}
 {} {} {} {.}
Untuk tutupan yang dimaksud
\zusatzklammer {tutupan ini} {}{}{,}
manakah sifat-sifat suatu
\definitionsverweis {partisi satuan subordinat}{}{}
terhadap tutupan tersebut yang dipenuhi oleh keluarga ini?

}
{} {}

\inputaufgabe
{}
{

Kita tinjau
\definitionsverweis {penghabisan kompak}{}{}

\mathbed {A_n= [-n,n]} {}
 {n \in \N} {}
 {} {} {} {,} dari
\definitionsverweis {bilangan riil}{}{}
beserta tutupan terbuka

\mathbed {W_n= A_{n+1}^{o}  \setminus A_{n-1}} {}
 {n \in \N} {}
 {} {} {} {,}
\zusatzklammer {dengan
\mathl{A_{-1}= \emptyset}{}} {} {.}
Carilah suatu tutupan $\R$ dengan
\definitionsverweis {interval terbuka}{}{,}
yang memenuhi sifat-sifat
dari Lemma 22.9
\zusatzklammer {beserta pembuktiannya} {} {}.

}
{} {}

\inputaufgabe
{}
{

Tunjukkan bahwa tidak terdapat barisan
\definitionsverweis {fungsi kontinu}{}{}
\maabbdisp {f_n} { [0,1] } { [0,1]
} {}
untuk

\mavergleichskette
{\vergleichskette
{ n
}
{ \in }{ \N
}
{  }{
}
{    }{
}
{   }{
}
}
{}{}{,}
yang membentuk suatu
\definitionsverweis {partisi satuan}{}{}
dan memiliki sifat bahwa $0$ termasuk dalam
\definitionsverweis {tumpuan}{}{}
setiap fungsi $f_n$.

}
{} {}

\inputaufgabe
{}
{

Misalkan

\mavergleichskette
{\vergleichskette
{ V
}
{ \subseteq }{ H
}
{  }{
}
{    }{
}
{   }{
}
}
{}{}{}
suatu
\definitionsverweis {himpunan terbuka}{}{}
di dalam
\definitionsverweis {setengah ruang}{}{}

\mavergleichskette
{\vergleichskette
{ H
}
{ \subset }{ \R^n
}
{  }{
}
{    }{
}
{   }{
}
}
{}{}{}
dan misalkan
\maabbdisp {f} { V } {\R
} {}
suatu
\definitionsverweis {fungsi yang terdiferensialkan secara kontinu}{}{.}
Tunjukkan bahwa terdapat suatu himpunan terbuka

\mavergleichskette
{\vergleichskette
{ \tilde{V}
}
{ \subseteq }{ \R^n
}
{  }{
}
{    }{
}
{   }{
}
}
{}{}{}
dan suatu
\definitionsverweis {fungsi yang terdiferensialkan secara kontinu}{}{}
\maabbdisp {\tilde{f}} {\tilde{V}} { \R
} {}
sedemikian sehingga

\mavergleichskette
{\vergleichskette
{ \tilde{V} \cap H
}
{ = }{ V
}
{  }{
}
{    }{
}
{   }{
}
}
{}{}{}
dan

\mavergleichskette
{\vergleichskette
{ \tilde{f} \vert_V
}
{ = }{ f
}
{  }{
}
{    }{
}
{   }{
}
}
{}{}{}
berlaku.

}
{} {}

\inputaufgabe
{}
{

Misalkan
\maabbdisp {f} {\R} {\R
} {}
suatu
\definitionsverweis {fungsi yang terdiferensialkan secara kontinu}{}{}
dengan
\definitionsverweis {tumpuan kompak}{}{.}
Tunjukkan bahwa
\definitionsverweis {turunan}{}{}
$f'$ juga memiliki tumpuan kompak dan bahwa

\mavergleichskette
{\vergleichskette
{ \int_\R f' d \lambda^1
}
{ = }{ 0
}
{  }{
}
{    }{
}
{   }{
}
}
{}{}{}
berlaku.

}
{} {}

\inputaufgabe
{}
{

Misalkan
\maabbdisp {f} {\R_{\geq 0}} {\R
} {}
suatu
\definitionsverweis {fungsi yang terdiferensialkan secara kontinu}{}{}
dengan
\definitionsverweis {tumpuan kompak}{}{.}
Tunjukkan bahwa
\definitionsverweis {turunan}{}{}
$f'$ juga memiliki tumpuan kompak dan bahwa

\mavergleichskette
{\vergleichskette
{ \int_{\R_{\geq 0} } f' d \lambda^1
}
{ = }{ -f(0)
}
{  }{
}
{    }{
}
{   }{
}
}
{}{}{}
berlaku. Tunjukkan bahwa pernyataan ini tidak harus berlaku jika $f$ tidak memiliki tumpuan kompak.

}
{} {}

\inputaufgabe
{}
{

Misalkan $X$ suatu
\definitionsverweis {manifold diferensiabel}{}{}
dengan suatu
\definitionsverweis {basis terhitung bagi topologinya}{}{.}
Tunjukkan bahwa terdapat suatu
\definitionsverweis {struktur Riemann}{}{}
pada $X$.

}
{} {}

\inputaufgabe
{}
{

Misalkan $X$ suatu
\definitionsverweis {manifold diferensiabel}{}{}
dengan suatu
\definitionsverweis {basis terhitung bagi topologinya}{}{}
dan misalkan
\maabb {p} { E } { X
} {}
suatu
\definitionsverweis {bundel vektor diferensiabel}{}{}
di atas $X$. Tunjukkan bahwa $E$ dapat dilengkapi dengan struktur
\definitionsverweis {bundel vektor Riemann}{}{.}

}
{} {}

  \zwischenueberschrift{Soal untuk dikumpulkan}

\inputaufgabe
{3}
{

Misalkan

\mathbed {A_n} {}
 {n \in \N} {}
 {} {} {} {,}
suatu
\definitionsverweis {penghabisan kompak}{}{} dari
\definitionsverweis {ruang topologis}{}{} $X$. Dengan konvensi
\mathl{A_{-1}= \emptyset}{,}
tunjukkan bahwa relasi

\mathdisp {A_{n+1} \setminus A_n ^{o}  \subseteq A_{n+2}^{o} \setminus A_{n-1}} {  }
 berlaku.

}
{} {}

\inputaufgabe
{4}
{

Untuk
\definitionsverweis {tutupan terbuka}{}{}

\mavergleichskettedisp
{\vergleichskette
{ \R
}
{ =} { \bigcup_{n \in \Z} ]n,n+3[
}
{ } {
}
{ } {
}
{ } {
}
}
{}{}{}
berikan suatu
\definitionsverweis {partisi satuan}{}{}
kontinu yang subordinat terhadap tutupan tersebut.

}
{} {}

\inputaufgabe
{6}
{

Kita tinjau
\definitionsverweis {penghabisan kompak}{}{}

\mathbeddisp {A_n = B  \left(  0 ,n \right)} {}
 {n \in \N} {}
 {} {} {} {,}
dari $\R^2$ beserta
\definitionsverweis {tutupan terbuka}{}{}

\mathbeddisp {W_n= A_{n+1}^{o} \setminus A_{n-1}} {}
 {n \in \N} {}
 {} {} {} {,}
\zusatzklammer {dengan
\mathl{A_{-1}= \emptyset}{}} {} {.}
Carilah suatu tutupan $\R^2$ dengan
\definitionsverweis {cakram terbuka}{}{,}
yang memenuhi sifat-sifat
dari Lemma 22.9
\zusatzklammer {beserta pembuktiannya} {} {}.

}
{} {}

\inputaufgabe
{3}
{

Berikan suatu contoh
\definitionsverweis {ruang topologis}{}{,}
yang tidak memiliki
\definitionsverweis {penghabisan kompak}{}{.}

}
{} {}


\section*{Solusi yang disediakan oleh sumber}
\addcontentsline{toc}{section}{Solusi yang disediakan oleh sumber}
\subsection*{Solusi Soal 22.6}
\input{generated/worksheet22_exercise06_solution.id.build.tex}

% O011-COMPLETE-EXTENSION-BEGIN
\part{Kelanjutan Edisi Lengkap}
\chapter*{Catatan provenans bagian tambahan}
\addcontentsline{toc}{chapter}{Catatan provenans bagian tambahan}
Unit 23--29 melanjutkan sumber Holger Brenner dari Wikiversity berbahasa Jerman. Hanya solusi inti yang benar-benar disediakan oleh sumber yang dimuat pada bagian unit. Jembatan Grup Lie dan de Rham serta enam perbaikan kekosongan ujian merupakan materi asli edisi ini dan tidak dinisbahkan kepada Brenner atau Wikiversity. Formulir solusi resmi mempertahankan kekosongan solusi sumber; perbaikan asli disajikan kemudian sebagai bagian yang terpisah dan berlabel jelas.
\newtheorem{Korolari}[fakt]{Korolari}
\newtheorem{Akibat}[fakt]{Akibat}

\part{Unit 23}
\chapter[Kuliah 23]{Kuliah 23: Teorema Stokes dan Teorema Titik Tetap Brouwer}
\setcounter{section}{23}

\bild{ \begin{center}
\includegraphics[width=5.5cm]{ \bildeinlesungjpg {SS-stokes} {jpg} }
\end{center}
\bildtext {George Stokes (1819 -1903)} }

\bildlizenz { SS-stokes.jpg } {} {Kelson} {Commons} {PD} {}

Teorema Stokes termasuk salah satu teorema terpenting dalam matematika. Teorema ini membentuk hubungan langsung antara integral suatu bentuk diferensial di atas batas manifold berbatas dan integral turunan eksterior bentuk tersebut di atas seluruh manifold. Dengan demikian, teorema ini merupakan perumuman yang sangat luas dari teorema dasar kalkulus, yang menyatakan bahwa integral tentu suatu fungsi yang didefinisikan pada sebuah interval dapat dinyatakan melalui antiturunannya hanya dengan menggunakan nilai-nilai pada ujung interval.

  \zwischenueberschrift{Teorema Stokes versi balok}

Sebelum kita merumuskan dan membuktikan teorema Stokes secara umum, kita berikan dahulu versi baloknya. Dalam versi ini, domain definisi bentuk diferensial adalah sebuah balok yang batasnya terdiri atas muka-mukanya. Agar objek geometris ini menjadi manifold berbatas, kita harus membuang \anfuehrung{rusuk-rusuknya}{} dari balok tersebut. Namun, rusuk-rusuk itu merupakan himpunan nol di dalam setiap muka
\zusatzklammer {dan demikian pula muka-muka merupakan himpunan nol di dalam keseluruhan balok} {} {,}
sehingga himpunan-himpunan bagian tersebut dapat diabaikan dalam pengintegralan.



\inputfaktbeweis
{Satz von Stokes/Quaderversion/Fakt}
{Teorema}
{}
{

\faktsituation {Misalkan

\mavergleichskette
{\vergleichskette
{ Q
}
{ \subseteq }{ \R^n
}
{  }{
}
{    }{
}
{   }{
}
}
{}{}{}
sebuah balok berdimensi $n$ yang sejajar sumbu
\zusatzklammer {beserta muka-mukanya, tetapi tanpa rusuk-rusuknya} {} {}\zusatzfussnote {Syarat-syarat ini menjamin bahwa yang diperoleh adalah manifold berbatas, dengan batas berupa gabungan saling lepas dari muka-muka terbuka. Namun, dalam pembuktian kita juga akan menggunakan balok tertutup} {.} {}
dengan
\definitionsverweis {batas}{}{}

\mathl{\partial Q}{,} dan misalkan $\omega$ didefinisikan pada $Q$ sebagai suatu
\definitionsverweis {bentuk diferensial}{}{}
yang
\definitionsverweis {terdiferensialkan secara kontinu}{}{}
dan berderajat $(n-1)$.}
\faktfolgerung {Maka

\mavergleichskettedisp
{\vergleichskette
{
\int_{  Q  }  d \omega
}
{ =} {
\int_{  \partial Q  }   \omega
}
{ } {
}
{ } {
}
{ } {
}
}
{}{}{.}}
\faktzusatz {}
\faktzusatz {}

}
{

Karena kedua ruas persamaan ini linear terhadap $\omega$, kita dapat menganggap bahwa $\omega$ berbentuk

\mavergleichskettedisp
{\vergleichskette
{ \omega
}
{ =} { f dx_2 \wedge \ldots \wedge dx_n
}
{ } {
}
{ } {
}
{ } {
}
}
{}{}{}
dengan suatu fungsi yang terdiferensialkan secara kontinu dan didefinisikan pada suatu lingkungan terbuka dari $Q$, yaitu $f$. Integral pada ruas kiri dan ruas kanan adalah
\definitionsverweis {integral Lebesgue}{}{}
dari fungsi-fungsi kontinu pada himpunan bagian
\mathkor {} {\R^n} {dan} {\R^{n-1}} {.}
Karena itu, pada kedua ruas kita dapat beralih ke
\definitionsverweis {penutupan topologis}{}{}
sebab hal tersebut hanya mengubah daerah integrasi yang bersangkutan sebesar suatu himpunan nol, sehingga tidak mengubah integralnya.

Kita tuliskan balok tertutup itu sebagai

\mavergleichskettedisp
{\vergleichskette
{ \overline{Q}
}
{ =} { [a,b] \times \tilde{Q}
}
{ } {
}
{ } {
}
{ } {
}
}
{}{}{.}
Kita menerapkan
Korolari 14.10
pada setiap muka $S$, kecuali
\mathkor {} {a \times \tilde{Q}} {dan} {b \times \tilde{Q}} {,}
dan memperoleh pada muka tersebut

\mavergleichskettedisp
{\vergleichskette
{
\int_{ S  }   \omega
}
{ =} {
\int_{  S  }  f dx_2 \wedge \ldots \wedge dx_n
}
{ =} {
\int_{ S  }   0
}
{ =} { 0
}
{ } {
}
}
{}{}{,}
karena pada setiap muka ini salah satu variabel
\mathl{x_2 , \ldots , x_n}{} konstan. Berdasarkan
teorema Fubini
dan
teorema dasar kalkulus
\zusatzklammer {yang diterapkan untuk setiap \mathlk{(x_2 , \ldots , x_n)}{} yang tetap} {} {,}
berlaku

\mavergleichskettealignhandlinks
{\vergleichskettealignhandlinks
{
\int_{ Q  }   d\omega
}
{ =} {
\int_{  Q  }   df \wedge  dx_2 \wedge \ldots \wedge dx_n
}
{ =} {
\int_{  Q  }   \left(  \sum_{j = 1}^n { \frac{   \partial   f   }{  \partial   x_j   } } dx_j  \right) \wedge  dx_2 \wedge \ldots \wedge dx_n
}
{ =} {
\int_{  Q  }   { \frac{   \partial   f   }{  \partial   x_1   } } dx_1 \wedge  dx_2 \wedge \ldots \wedge dx_n
}
{ =} {
\int_{ \tilde{Q}   }   \left(  \int_{[a,b] } { \frac{   \partial   f   }{  \partial   x_1   } } dx_1  \right)  dx_2 \wedge \ldots \wedge dx_n
}
}
{
\vergleichskettefortsetzungalign
{ =} {
\int_{  \tilde{Q}   }   \left(  f(b,x_2 , \ldots , x_n ) - f(a,x_2 , \ldots , x_n )   \right) dx_2 \wedge \ldots \wedge dx_n
}
{ =} {
\int_{ b \times \tilde{Q}   }   f dx_2 \wedge \ldots \wedge dx_n  -
\int_{  a \times \tilde{Q}   }  f dx_2 \wedge \ldots \wedge dx_n
}
{ =} { \sum_{S \text{ muka berorientasi dari } Q}
\int_{  S  }  f dx_2 \wedge \ldots \wedge dx_n
}
{ =} {
\int_{  \partial Q  }   f dx_2 \wedge \ldots \wedge dx_n
}
}
{
\vergleichskettefortsetzungalign
{ =} {
\int_{  \partial Q  }   \omega
}
{ } {}
{ } {}
{ } {}
}{.}

}

  \zwischenueberschrift{Teorema Stokes}



\inputfaktbeweis
{Mannigfaltigkeit mit Rand/Satz von Stokes/Fakt}
{Teorema}
{}
{

\faktsituation {Misalkan $M$ suatu
\definitionsverweis {manifold diferensiabel berbatas}{}{}
yang
\definitionsverweis {berdimensi}{}{} $n$,
\definitionsverweis {berorientasi}{}{},
memiliki batas $\partial M$, serta memiliki
\definitionsverweis {basis terhitung bagi topologinya}{}{.}
Misalkan pula $\omega$ suatu
\definitionsverweis {bentuk diferensial}{}{}
berderajat $(n-1)$ yang
\definitionsverweis {terdiferensialkan secara kontinu}{}{}
dan memiliki
\definitionsverweis {tumpuan}{}{}
\definitionsverweis {kompak}{}{}\zusatzfussnote {Tumpuan suatu bentuk diferensial adalah penutupan topologis dari himpunan titik tempat bentuk tersebut \mathlk{\neq 0}{}} {.} {} pada $M$.}
\faktfolgerung {Maka

\mavergleichskettedisp
{\vergleichskette
{
\int_{ M  }  d\omega
}
{ =} {
\int_{  \partial M  }   \omega
}
{ } {
}
{ } {
}
{ } {
}
}
{}{}{.}}
\faktzusatz {}
\faktzusatz {}

}
{

\teilbeweis {}{}{}
{Misalkan

\mathbed {U_i} {}
 {i \in I} {}
 {} {} {} {,}
suatu
\definitionsverweis {tutupan terbuka}{}{}
dari $M$ oleh
\definitionsverweis {peta-peta berorientasi}{}{,}
dan misalkan

\mathbed {h_j} {}
 {j \in J} {}
 {} {} {} {,}
suatu
\definitionsverweis {partisi satuan yang terdiferensialkan secara kontinu dan subordinat terhadap tutupan ini}{}{,}
yang ada menurut
Teorema 22.10.
Untuk setiap

\mavergleichskette
{\vergleichskette
{ P
}
{ \in }{ M
}
{  }{
}
{    }{
}
{   }{
}
}
{}{}{}
terdapat suatu lingkungan terbuka

\mavergleichskette
{\vergleichskette
{ P
}
{ \in }{ V_P
}
{ \subseteq }{ M
}
{    }{
}
{   }{
}
}
{}{}{}
sedemikian rupa sehingga

\mavergleichskette
{\vergleichskette
{ h_j \vert_{V_P }
}
{ = }{ 0
}
{  }{
}
{    }{
}
{   }{
}
}
{}{}{}
selain untuk sejumlah berhingga pengecualian. Misalkan $Y$ adalah tumpuan $\omega$. Tutupan

\mavergleichskette
{\vergleichskette
{ Y
}
{ \subseteq }{ \bigcup_{P \in M} V_P
}
{  }{
}
{    }{
}
{   }{
}
}
{}{}{}
memiliki suatu subtutupan berhingga karena
\definitionsverweis {kekompakan}{}{}
yang diasumsikan; sebutlah

\mavergleichskettedisp
{\vergleichskette
{ Y
}
{ \subseteq} { V_{1} \cup \ldots \cup V_{r}
}
{ =} { V
}
{ } {
}
{ } {
}
}
{}{}{.}
Karena itu, hanya sejumlah berhingga dari $h_j$ pada $V$ yang berbeda dari $0$. Kita definisikan

\mavergleichskette
{\vergleichskette
{ \omega_j
}
{ = }{ h_j \omega
}
{  }{
}
{    }{
}
{   }{
}
}
{}{}{;}
bentuk-bentuk diferensial ini juga terdiferensialkan secara kontinu. Tumpuan $h_j$ adalah suatu
\definitionsverweis {himpunan bagian tertutup}{}{}
dari $M$ yang termuat dalam
\mathl{U_{i(j)}}{.} Oleh karena itu, tumpuan $\omega_j$ termuat dalam
\mathl{Y \cap U_{i(j)}}{} dan dengan sendirinya kompak menurut
Soal 13.10.
Berlaku

\mavergleichskettedisp
{\vergleichskette
{ \omega
}
{ =} { \sum_{j \in J} \omega_j
}
{ } {
}
{ } {
}
{ } {
}
}
{}{}{,}
dan hanya sejumlah berhingga dari bentuk-bentuk diferensial
\mathl{\omega_j}{} tersebut yang berbeda dari $0$, sebab

\mavergleichskette
{\vergleichskette
{ \omega_j \vert_{M \setminus Y}
}
{ = }{ 0
}
{  }{
}
{    }{
}
{   }{
}
}
{}{}{}
untuk semua $j$, sedangkan

\mavergleichskettedisp
{\vergleichskette
{ \omega_j \vert_V
}
{ =} { 0
}
{ } {
}
{ } {
}
{ } {
}
}
{}{}{}
untuk semua $j$ selain sejumlah berhingga pengecualian. Berdasarkan
aditivitas integral bentuk diferensial
dan
aditivitas turunan eksterior,
pernyataan tersebut dapat dibuktikan secara terpisah untuk setiap $\omega_j$. Jadi, kita dapat menganggap bahwa diberikan suatu bentuk diferensial berderajat $(n-1)$ yang terdiferensialkan secara kontinu, mempunyai tumpuan kompak, dan tumpuan itu seluruhnya termuat dalam suatu lingkungan peta $U$.}
{}
\teilbeweis {}{}{}
{Terdapat suatu
\definitionsverweis {difeomorfisme}{}{}
\maabb {\alpha} { U } { V
} {}
dengan

\mavergleichskette
{\vergleichskette
{ V
}
{ \subseteq }{ H
}
{  }{
}
{    }{
}
{   }{
}
}
{}{}{}
terbuka, yang sekaligus menginduksi suatu difeomorfisme antara batas-batas

\mavergleichskette
{\vergleichskette
{ \partial U
}
{ = }{ \partial M \cap U
}
{  }{
}
{    }{
}
{   }{
}
}
{}{}{}
dan

\mavergleichskette
{\vergleichskette
{ \partial V
}
{ = }{ \partial H \cap V
}
{  }{
}
{    }{
}
{   }{
}
}
{}{}{.}
Dalam hal ini berlaku

\mavergleichskettedisp
{\vergleichskette
{
\int_{  \partial U  }   \omega
}
{ =} {
\int_{  \partial V  }  \alpha^{-1 *}\omega
}
{ } {
}
{ } {
}
{ } {
}
}
{}{}{}
dan

\mavergleichskettedisp
{\vergleichskette
{
\int_{ U  }  d \omega
}
{ =} {
\int_{ V  }  \alpha^{-1 *} (  d \omega)
}
{ =} {
\int_{  V  }  d \left(  \alpha^{-1 *}\omega  \right)
}
{ } {
}
{ } {
}
}
{}{}{}
menurut
Lema 20.2 (5).}
{\leerzeichen{}Dengan demikian, kita dapat bertolak dari suatu bentuk diferensial yang terdiferensialkan secara kontinu, didefinisikan pada

\mavergleichskette
{\vergleichskette
{ V
}
{ \subseteq }{ H
}
{  }{
}
{    }{
}
{   }{
}
}
{}{}{,}
serta mempunyai tumpuan kompak. Di luar tumpuannya, bentuk ini dapat kita perluas ke seluruh $H$ sebagai bentuk nol.}
\teilbeweis {}{}{}
{Karena tumpuannya kompak, terdapat suatu balok yang cukup besar

\mavergleichskette
{\vergleichskette
{ Q
}
{ \subseteq }{ H
}
{  }{
}
{    }{
}
{   }{
}
}
{}{}{,}
yang salah satu mukanya, $S$, terletak pada $\partial H$ dan yang hanya berpotongan dengan tumpuan $\omega$ di $S$. Pada semua muka lain dari $Q$, $\omega$
\zusatzklammer {dan karena itu juga $d\omega$} {} {}
adalah bentuk nol. Oleh sebab itu, pada satu pihak berlaku

\mavergleichskettedisp
{\vergleichskette
{
\int_{ H  }  d\omega
}
{ =} {
\int_{ Q  }  d\omega
}
{ } {
}
{ } {
}
{ } {
}
}
{}{}{}
dan pada pihak lain

\mavergleichskettedisp
{\vergleichskette
{
\int_{  \partial H  }   \omega
}
{ =} {
\int_{  S  }   \omega
}
{ =} {
\int_{  \partial Q  }   \omega
}
{ } {
}
{ } {
}
}
{}{}{.}
Dengan demikian, pernyataan tersebut mengikuti dari
teorema Stokes versi balok.}
{}

}



\inputfaktbeweis
{Mannigfaltigkeiten ohne Rand/Satz von Stokes/Fakt}
{Korolari}
{}
{

\faktsituation {Misalkan $M$ suatu manifold
\definitionsverweis {berdimensi}{}{} $n$ yang
\definitionsverweis {berorientasi}{}{} dan
\definitionsverweis {diferensiabel}{}{}
\zusatzklammer {tanpa batas} {} {}
serta memiliki
\definitionsverweis {basis terhitung bagi topologinya}{}{.}
Misalkan pula $\omega$ suatu
\definitionsverweis {bentuk diferensial}{}{}
berderajat $(n-1)$ yang
\definitionsverweis {terdiferensialkan secara kontinu}{}{}
dan memiliki
\definitionsverweis {tumpuan}{}{}
\definitionsverweis {kompak}{}{} pada $M$.}
\faktfolgerung {Maka

\mavergleichskettedisp
{\vergleichskette
{
\int_{ M  }  d\omega
}
{ =} { 0
}
{ } {
}
{ } {
}
{ } {
}
}
{}{}{.}}
\faktzusatz {}
\faktzusatz {}

}
{

Hal ini segera mengikuti dari
Teorema 23.2
karena

\mavergleichskette
{\vergleichskette
{ \partial M
}
{ = }{ \emptyset
}
{  }{
}
{    }{
}
{   }{
}
}
{}{}{.}

}



\inputfaktbeweis
{Mannigfaltigkeiten mit Rand/Satz von Stokes/Differentialform ist null auf Rand/Fakt}
{Korolari}
{}
{

\faktsituation {Misalkan $M$ suatu
\definitionsverweis {manifold diferensiabel berbatas}{}{}
yang
\definitionsverweis {berdimensi}{}{} $n$,
\definitionsverweis {berorientasi}{}{},
dan memiliki
\definitionsverweis {basis terhitung bagi topologinya}{}{.}
Misalkan pula $\omega$ suatu
\definitionsverweis {bentuk diferensial}{}{}
berderajat $(n-1)$ yang
\definitionsverweis {terdiferensialkan secara kontinu}{}{}
dan memiliki
\definitionsverweis {tumpuan}{}{}
\definitionsverweis {kompak}{}{}
pada $M$,}
\faktvoraussetzung {serta pada batas $\partial M$ bernilai tetap $0$.}
\faktfolgerung {Maka

\mavergleichskettedisp
{\vergleichskette
{
\int_{ M  }  d\omega
}
{ =} { 0
}
{ } {
}
{ } {
}
{ } {
}
}
{}{}{.}}
\faktzusatz {}
\faktzusatz {}

}
{

Hal ini segera mengikuti dari
Teorema 23.2.

}

Hal penting dalam pernyataan di atas adalah bahwa $\omega$ bernilai $0$ pada batas; tidak cukup bahwa turunan eksterior $d\omega$ bernilai $0$ pada batas, sebagaimana sudah ditunjukkan oleh situasi berdimensi satu.

Ada banyak cara untuk merealisasikan bentuk volume

\mavergleichskette
{\vergleichskette
{ \tau
}
{ = }{ dx_1 \wedge \ldots \wedge dx_n
}
{  }{
}
{    }{
}
{   }{
}
}
{}{}{}
dari $\R^n$ sebagai turunan eksterior suatu bentuk berderajat
\mathl{(n-1)}{,} misalnya dengan

\mavergleichskette
{\vergleichskette
{  \omega
}
{ = }{ x_1dx_2 \wedge \ldots \wedge dx_n
}
{  }{
}
{    }{
}
{   }{
}
}
{}{}{.}
Dengan demikian, perhitungan volume suatu benda berbatas dapat dikembalikan pada perhitungan integral di atas batasnya. Dalam kasus bidang, pernyataan ini juga disebut \stichwort {teorema Green} {.}



\inputfaktbeweis
{Integration auf ebener Mannigfaltigkeit mit Rand/Satz von Green/Fakt}
{Teorema}
{}
{

\faktsituation {Misalkan

\mavergleichskette
{\vergleichskette
{M
}
{ \subseteq }{ \R^2
}
{  }{
}
{    }{
}
{   }{
}
}
{}{}{}
suatu manifold
\definitionsverweis {kompak}{}{}
\definitionsverweis {dengan batas}{}{}\zusatzfussnote {Bidang real sekitar hanya berperan dalam menetapkan koordinat dan bentuk diferensial pada $M$} {.} {}
$\partial M$, dan misalkan
\maabbdisp {f,g} { M } {\R
} {}
\definitionsverweis {fungsi-fungsi yang terdiferensialkan secara kontinu}{}{.}}
\faktfolgerung {Maka

\mavergleichskettealignhandlinks
{\vergleichskettealignhandlinks
{
\int_{  M  }   \left(  - { \frac{   \partial   f   }{  \partial   y   } } ( x,y) + { \frac{   \partial   g   }{  \partial   x   } } ( x,y)  \right) dx \wedge dy
}
{ =} {
\int_{  \partial M  }  f(x,y)dx +g(x,y)dy
}
{ } {
}
{ } {
}
{ } {
}
}
{}
{}{}}
\faktzusatz {}
\faktzusatz {}

}
{

Hal ini mengikuti dari
Teorema 23.2,
yang diterapkan pada suatu bentuk-$1$ yang
\definitionsverweis {terdiferensialkan secara kontinu}{}{},
yakni suatu
\definitionsverweis {bentuk diferensial}{}{}

\mavergleichskette
{\vergleichskette
{ \omega
}
{ = }{ f(x,y)dx + g(x,y)dy
}
{  }{
}
{    }{
}
{   }{
}
}
{}{}{.}

}



\inputfaktbeweis
{Integration auf ebener Mannigfaltigkeit mit Rand/Satz von Green/Flächenversion/Fakt}
{Korolari}
{}
{

\faktsituation {Misalkan

\mavergleichskette
{\vergleichskette
{ M
}
{ \subseteq }{ \R^2
}
{  }{
}
{    }{
}
{   }{
}
}
{}{}{}
suatu manifold
\definitionsverweis {kompak}{}{}
\definitionsverweis {dengan batas}{}{}

\mathl{\partial M}{.}}
\faktfolgerung {Maka luas $M$ sama dengan

\mavergleichskettedisp
{\vergleichskette
{ \lambda^2(M)
}
{ =} {
\int_{  \partial M  }   x dy
}
{ =} { -
\int_{  \partial M  }   y dx
}
{ } {
}
{ } {
}
}
{}{}{.}}
\faktzusatz {}
\faktzusatz {}

}
{

Hal ini mengikuti karena

\mavergleichskette
{\vergleichskette
{ \lambda^2(M)
}
{ = }{
\int_{ M  }  dx \wedge dy
}
{  }{
}
{    }{
}
{   }{
}
}
{}{}{}
dari
Teorema 23.5,
yang diterapkan pada
\mathkor {} {f=0,\, g=x} {atau} {f=-y,\, g=0} {.}

}

\inputbemerkung
{}
{

Teorema Green juga dapat diterapkan pada daerah kompak

\mavergleichskette
{\vergleichskette
{ M
}
{ \subseteq }{ \R^2
}
{  }{
}
{    }{
}
{   }{
}
}
{}{}{}
yang batas
\zusatzklammer {topologisnya} {} {}
terhubung dan tersusun atas sejumlah berhingga kurva mulus yang tidak harus mulus pada titik-titik sambungannya
\zusatzklammer {misalnya sebuah segitiga} {} {.}
Integral batas kemudian merupakan jumlah integral lintasan di atas potongan-potongan kurva mulus tersebut. Situasi yang sedikit lebih umum ini dapat dikembalikan pada
Teorema 23.5
dengan salah satu dari dua cara: melicinkan titik-titik sambungan untuk memperoleh batas mulus, sambil membuat penyimpangan integral sekecil yang diinginkan; atau melicinkan bentuk diferensial hingga bernilai $0$ dalam lingkungan-lingkungan kecil dari titik-titik sambungan.

}

\inputbeispiel{}
{

Salah satu kasus khusus
teorema Stokes
adalah apa yang disebut \stichwort {teorema divergensi} {} atau \stichwort {teorema Gauss} {.} Untuk suatu manifold kompak berdimensi tiga dengan batas

\mavergleichskette
{\vergleichskette
{ M
}
{ \subset }{ U
}
{ \subseteq }{ \R^3
}
{    }{
}
{   }{
}
}
{}{}{}
\zusatzklammer {$U$ suatu himpunan bagian terbuka} {} {}
dan suatu bentuk diferensial-$2$ $\omega$ pada $U$, teorema ini menyatakan persamaan

\mavergleichskettedisp
{\vergleichskette
{ \int_{\partial M} \omega
}
{ =} { \int_M d \omega
}
{ } {
}
{ } {
}
{ } {
}
}
{}{}{.}
Teorema ini berkaitan dengan situasi fisis suatu aliran. Bentuk $\omega$, untuk suatu titik

\mavergleichskette
{\vergleichskette
{ P
}
{ \in }{ U
}
{  }{
}
{    }{
}
{   }{
}
}
{}{}{}
dan sebuah jajaran genjang infinitesimal pada titik itu, menyatakan berapa banyak fluida yang mengalir menembus potongan tersebut per satuan waktu. Dalam deskripsi ekuivalen situasi ini dengan suatu medan vektor,
\mathl{F(P)}{} menyatakan arah aliran beserta kekuatannya. Turunan $d \omega$ adalah suatu bentuk volume yang disebut \stichwort {divergensi} {.} Pada suatu titik, bentuk ini menyatakan pertambahan infinitesimal bahan aliran
\zusatzklammer {\stichwort {sumber} {} atau \stichwort {serapan} {}} {} {}
di titik tersebut. Syarat

\mavergleichskette
{\vergleichskette
{d \omega
}
{ = }{ 0
}
{  }{
}
{    }{
}
{   }{
}
}
{}{}{}
sering dipenuhi dan berarti bahwa fluida tidak mengalami penambahan bahan di $U$. Teorema Gauss kemudian menyatakan bahwa aliran total yang menembus batas
\zusatzklammer {dengan orientasi menentukan aliran mana yang dianggap mengarah ke luar atau ke dalam} {} {}
sama dengan $0$. Jadi, apa pun yang mengalir masuk ke $M$ pada suatu tempat akan mengalir keluar lagi di tempat lain.

}

  \zwischenueberschrift{Teorema titik tetap Brouwer}



\inputfaktbeweis
{Mannigfaltigkeit mit Rand/Stokes/Nichtexistenz von Retraktionen/Fakt}
{Teorema}
{}
{

\faktsituation {Misalkan $M$ suatu
\definitionsverweis {manifold diferensiabel berbatas}{}{}
yang
\definitionsverweis {kompak}{}{},
\definitionsverweis {berorientasi}{}{},
memiliki batas
\mathl{\partial M}{,} dan memiliki
\definitionsverweis {basis terhitung bagi topologinya}{}{.}}
\faktfolgerung {Maka tidak ada
\definitionsverweis {pemetaan yang terdiferensialkan secara kontinu}{}{}
\maabbdisp {\varphi} { M } { \partial M
} {,}
yang
\definitionsverweis {pembatasannya}{}{}
pada $\partial M$ adalah identitas.}
\faktzusatz {}
\faktzusatz {}

}
{

Menurut
Teorema 21.8,
batas
\mathl{\partial M}{} adalah suatu
\definitionsverweis {manifold diferensiabel}{}{}
\definitionsverweis {berorientasi}{}{}
\zusatzklammer {tanpa batas} {} {.}
Oleh karena itu, menurut
Teorema 22.11,
terdapat suatu
\definitionsverweis {bentuk volume positif}{}{}
yang
\definitionsverweis {terdiferensialkan secara kontinu}{}{},
yaitu $\tau$, pada
\mathl{\partial M}{.} Berlaku

\mavergleichskette
{\vergleichskette
{
\int_{  \partial M  }  \tau
}
{ > }{ 0
}
{  }{
}
{    }{
}
{   }{
}
}
{}{}{.}
\definitionsverweis {Turunan eksterior}{}{}
dari bentuk volume $\tau$ adalah $0$.
Andaikan terdapat suatu
\definitionsverweis {pemetaan yang terdiferensialkan secara kontinu}{}{}
\maabbdisp {\varphi} { M } { \partial M
} {}
dengan

\mavergleichskette
{\vergleichskette
{ \varphi \vert_{\partial M}
}
{ = }{ \operatorname{Id}_{\partial M}
}
{  }{
}
{    }{
}
{   }{
}
}
{}{}{.}
Maka
\definitionsverweis {bentuk tarik balik}{}{}

\mathl{\varphi^* \tau}{} adalah suatu
\definitionsverweis {bentuk diferensial}{}{}
berderajat $(n-1)$ pada $M$, yang pembatasannya pada batas sama dengan $\tau$. Dengan menggunakan
Teorema 23.2
dan
Teorema 20.4 (5),
kita memperoleh

\mavergleichskettealign
{\vergleichskettealign
{
\int_{  \partial M  }   \tau
}
{ =} {
\int_{  \partial  M  }   \varphi^* \tau
}
{ =} {
\int_{  M  }   d(\varphi^* \tau)
}
{ =} {
\int_{  M  }   \varphi^* (d\tau)
}
{ =} {
\int_{  M  }   \varphi^* (0)
}
}
{
\vergleichskettefortsetzungalign
{ =} { 0
}
{ } {}
{ } {}
{ } {}
}
{}{.}
Ini sebuah kontradiksi.

}

Pernyataan ini juga dirumuskan dengan mengatakan bahwa tidak ada
\zusatzklammer {yang terdiferensialkan secara kontinu} {} {} \stichwort {retraksi} {} ke batas.

\bild{ \begin{center}
\includegraphics[width=5.5cm]{ \bildeinlesungjpg {Théorème-de-Brouwer-(cond-2)} {jpg} }
\end{center}
\bildtext {Ilustrasi cakram dengan lintasan-lintasan berarah yang berpilin menuju titik pusat.} }

\bildlizenz { Théorème-de-Brouwer-(cond-2).jpg } {} {Jean-Luc W} {Commons} {CC-by-sa 3.0} {}

\bild{ \begin{center}
\includegraphics[width=5.5cm]{ \bildeinlesungjpg {Théorème-de-Brouwer-(cond-1)} {jpg} }
\end{center}
\bildtext {Ilustrasi lintasan-lintasan melingkar berarah pada anulus dan pada pita persegi panjang padanannya.} }

\bildlizenz { Théorème-de-Brouwer-(cond-1).jpg } {} {Jean-Luc W} {Commons} {CC-by-sa 3.0} {}

Teorema berikut disebut \stichwort {teorema titik tetap Brouwer} {.}


\inputfaktbeweis
{Brouwersche Fixpunktsatz/Stetig differenzierbar/Retraktion/Fakt}
{Teorema}
{}
{

\faktsituation {Misalkan
\maabbdisp {\psi} { B  \left(  0 ,r \right) } { B  \left(  0 ,r \right)
} {}
suatu
\definitionsverweis {pemetaan yang terdiferensialkan secara kontinu}{}{}
dari
\definitionsverweis {bola tertutup}{}{}
di $\R^n$ ke dirinya sendiri.}
\faktfolgerung {Maka $\psi$ memiliki suatu
\definitionsverweis {titik tetap}{}{.}}
\faktzusatz {}
\faktzusatz {}

}
{

Untuk menyederhanakan notasi, misalkan

\mavergleichskette
{\vergleichskette
{ r
}
{ = }{ 1
}
{  }{
}
{    }{
}
{   }{
}
}
{}{}{.}
Andaikan ada suatu pemetaan $\psi$ yang terdiferensialkan secara kontinu dan tidak memiliki titik tetap. Maka selalu berlaku

\mavergleichskettedisp
{\vergleichskette
{ x
}
{ \neq} { \psi(x)
}
{ } {
}
{ } {
}
{ } {
}
}
{}{}{,}
sehingga kedua titik itu menentukan sebuah garis. Gagasannya adalah menggunakan garis ini untuk mengambil salah satu
\zusatzklammer {dari kedua} {} {}
titik potong dengan sfer sebagai titik citra suatu retraksi ke batas. Dengan fungsi bantu

\mavergleichskettedisp
{\vergleichskette
{ h(x)
}
{ =} { { \frac{   \psi(x)-x }{  \Vert { \psi (x)-x} \Vert  } }
}
{ } {
}
{ } {
}
{ } {
}
}
{}{}{,}
kita mendefinisikan suatu pemetaan
\maabbdisp {\varphi} { B  \left(  0 , 1 \right) } { \R^n
} {}
melalui

\mavergleichskettedisphandlinks
{\vergleichskettedisphandlinks
{ \varphi(x)
}
{ =} { x + \left(  - \left\langle  x  , h(x)  \right\rangle + \sqrt{1 + \left\langle  x  , h(x)  \right\rangle^2 - \Vert { x} \Vert^2  }  \right) \cdot h(x)
}
{ } {
}
{ } {
}
{ } {
}
}
{}{}{.}
Ekspresi di bawah akar kuadrat bernilai positif. Hal ini jelas jika

\mavergleichskette
{\vergleichskette
{ \Vert { x } \Vert
}
{ < }{ 1
}
{  }{
}
{    }{
}
{   }{
}
}
{}{}{,}
sedangkan jika

\mavergleichskette
{\vergleichskette
{ \Vert { x } \Vert
}
{ = }{ 1
}
{  }{
}
{    }{
}
{   }{
}
}
{}{}{,}
terdapat suatu titik pada sfer yang garis penghubungnya dengan titik bola
\mathl{\psi(x)}{} tidak tegak lurus terhadap $x$
\zusatzklammer {ruang singgung afin pada suatu titik sfer hanya berpotongan dengan bola di satu titik} {} {,}
sehingga

\mavergleichskettedisp
{\vergleichskette
{ \left\langle  x  , h(x)  \right\rangle
}
{ \neq} { 0
}
{ } {
}
{ } {
}
{ } {
}
}
{}{}{.}
Karena akar kuadrat dan nilai mutlak
\definitionsverweis {terdiferensialkan secara kontinu}{}{}
di luar titik nol, baik
\mathkor {} {h(x)} {maupun} {\varphi(x)} {}
merupakan pemetaan yang terdiferensialkan secara kontinu. Menurut
Soal 23.23,
pemetaan $\varphi$ memetakan bola ke sfer, dan pembatasannya pada sfer adalah identitas. Dengan demikian, terdapat suatu
\definitionsverweis {retraksi}{}{}
yang terdiferensialkan secara kontinu dari bola pejal tertutup ke batasnya, yang mustahil menurut
Teorema 23.9.

}


\chapter{Lembar Kerja 23}
\setcounter{section}{23}

  \zwischenueberschrift{Soal latihan}

\inputaufgabe
{}
{

Bahaslah
rumus Newton--Leibniz
sebagai kasus khusus dari
teorema Stokes.

}
{} {}

\inputaufgabe
{}
{

Buktikan secara langsung
versi balok dari teorema Stokes
untuk balok
\mathl{[a_1,b_1] \times \cdots \times [a_n,b_n]}{} dan suatu
$n-1$-\definitionsverweis {bentuk diferensial}{}{}
$\omega$ berbentuk

\mavergleichskettedisp
{\vergleichskette
{ \omega
}
{ =} { x_1^{m_1 } \cdots x_n^{m_n }  dx_2 \wedge \ldots \wedge dx_n
}
{ } {
}
{ } {
}
{ } {
}
}
{}{}{.}

}
{} {}

\inputaufgabe
{}
{

Misalkan $M$ suatu
$n$-\definitionsverweis {dimensional}{}{}
\definitionsverweis {berorientasi}{}{}
\definitionsverweis {manifold diferensiabel berbatas}{}{}
dengan
\definitionsverweis {basis topologi terhitung}{}{},
dan misalkan $\omega$ suatu $(n-1)$-\definitionsverweis {bentuk diferensial}{}{}
yang
\definitionsverweis {terdiferensialkan secara kontinu}{}{},
\definitionsverweis {tertutup}{}{}, dan mempunyai
\definitionsverweis {tumpuan kompak}{}{}
pada $M$. Tunjukkan bahwa

\mavergleichskettedisp
{\vergleichskette
{ \int_{\partial M} \omega
}
{ =} { 0
}
{ } {
}
{ } {
}
{ } {
}
}
{}{}{.}

}
{} {}

\inputaufgabe
{}
{

Misalkan $M$ suatu
\definitionsverweis {manifold diferensiabel}{}{}
$n$-\definitionsverweis {dimensional}{}{},
\definitionsverweis {kompak}{}{}, dan
\definitionsverweis {berorientasi}{}{}
\zusatzklammer {tanpa batas} {} {},
dengan
\definitionsverweis {basis topologi terhitung}{}{},
dan misalkan $\omega$ suatu
$(n-1)$-\definitionsverweis {bentuk diferensial}{}{}
yang
\definitionsverweis {terdiferensialkan secara kontinu}{}{}
pada $M$. Tunjukkan bahwa

\mavergleichskettedisp
{\vergleichskette
{
\int_{ M  }  d\omega
}
{ =} { 0
}
{ } {
}
{ } {
}
{ } {
}
}
{}{}{.}

}
{Apa arti pernyataan ini untuk $S^1$? Bagaimana pernyataan tersebut dapat dibuktikan dalam kasus ini melalui suatu integral lintasan?} {}

\inputaufgabe
{}
{

Misalkan $M$ suatu
\definitionsverweis {manifold diferensiabel}{}{}
\definitionsverweis {kompak}{}{} dan
\definitionsverweis {berorientasi}{}{}
\zusatzklammer {tanpa batas} {} {},
dengan
\definitionsverweis {basis topologi terhitung}{}{},
dan misalkan $\tau$ suatu
\definitionsverweis {bentuk volume positif}{}{}
pada $M$. Tunjukkan bahwa $\tau$ tidak
\definitionsverweis {eksak}{}{.}

}
{Bagaimana keadaannya tanpa asumsi kekompakan?} {}

\inputaufgabegibtloesung
{}
{

Misalkan $M$ suatu
\definitionsverweis {manifold diferensiabel}{}{}
dan $\omega$ suatu
\definitionsverweis {eksak}{}{} \definitionsverweis {bentuk diferensial}{}{}
pada $M$. Misalkan $S$ suatu
\definitionsverweis {manifold diferensiabel}{}{}
\definitionsverweis {kompak}{}{} dan
\definitionsverweis {berorientasi}{}{}
\zusatzklammer {tanpa batas} {} {},
dengan
\definitionsverweis {basis topologi terhitung}{}{},
serta misalkan
\maabbdisp {\varphi} { S } { M
} {}
suatu pemetaan yang
\definitionsverweis {terdiferensialkan secara kontinu}{}{.}
Tunjukkan bahwa

\mavergleichskettedisp
{\vergleichskette
{
\int_{  S  }   \varphi^*\omega
}
{ =} { 0
}
{ } {
}
{ } {
}
{ } {
}
}
{}{}{.}

}
{} {}

\inputaufgabe
{}
{

Kita tinjau
\definitionsverweis {pemetaan yang terdiferensialkan secara kontinu}{}{}
\zusatzklammer {\mathlk{n \geq 2}{}} {} {}
\maabbeledisp {\pi} {\R^n \setminus \{0\}} { S^{n-1}
} {(x_1 , \ldots , x_n) } { { \frac{   1  }{  \sqrt{x_1^2 + \cdots + x_n^2}  } } (x_1 , \ldots , x_n)
} {.}
Misalkan $\omega$ adalah
\definitionsverweis {bentuk volume kanonik}{}{}
pada $S^{n-1}$. Tunjukkan bahwa
\mathl{\pi^* \omega}{} pada $\R^n \setminus \{0\}$ merupakan
\definitionsverweis {tertutup}{}{},
tetapi tidak
\definitionsverweis {eksak}{}{}
$n-1$-\definitionsverweis {bentuk diferensial}{}{.}

}
{} {}

\inputaufgabe
{}
{

Misalkan

\mavergleichskette
{\vergleichskette
{ M
}
{ = }{ B  \left(  0 , 1 \right) \setminus \{(0,0)\}
}
{ \subset }{ \R^2
}
{    }{
}
{   }{
}
}
{}{}{}
dengan batas

\mavergleichskette
{\vergleichskette
{ \partial M
}
{ = }{ S^1
}
{  }{
}
{    }{
}
{   }{
}
}
{}{}{.}
Kita tinjau kedua
$1$-\definitionsverweis {bentuk diferensial}{}{}
yang
\definitionsverweis {terdiferensialkan secara kontinu}{}{}

\mathdisp {\omega= ydx-xdy \text{ dan } \tau= { \frac{   1  }{  x^2+y^2 } } \left(  ydx-xdy  \right)} {  }

pada $M$. Tunjukkan bahwa pembatasan kedua bentuk tersebut pada batas saling berimpit, dan khususnya

\mavergleichskette
{\vergleichskette
{ \int_{\partial M} \omega
}
{ =} { \int_{\partial M} \tau
}
{  }{
}
{    }{
}
{   }{
}
}
{}{}{}
berlaku. Bandingkan
\mathkor {} {\int_M d \omega} {dan} {\int_M d \tau} {.}

}
{} {}

\inputaufgabe
{}
{

Yakinkan diri bahwa
teorema Green
tidak menyatakan bahwa luas suatu daerah berbatas di $\R^2$ hanya bergantung pada panjang batasnya.

}
{} {}

\inputaufgabe
{}
{

Misalkan $D$ segitiga dengan titik sudut
$(0,2),\, (1,-1)$, dan $(-2,-1)$,
serta

\mavergleichskette
{\vergleichskette
{ \tau
}
{ = }{ x^2ydx \wedge dy
}
{  }{
}
{    }{
}
{   }{
}
}
{}{}{}
suatu
$2$-\definitionsverweis {bentuk diferensial}{}{} pada $D$. Temukan suatu
\definitionsverweis {bentuk primitif}{}{}
untuk $\tau$, lalu gunakan bentuk tersebut untuk menghitung
\mathl{\int_{ D  }  \tau}{} melalui suatu integral pada batas segitiga.

}
{} {}

\inputaufgabe
{}
{

Verifikasikan
teorema Green
melalui perhitungan eksplisit untuk persegi satuan

\mavergleichskette
{\vergleichskette
{ T
}
{ = }{ [0,1] \times [0,1]
}
{  }{
}
{    }{
}
{   }{
}
}
{}{}{}
dan
\definitionsverweis {bentuk-bentuk diferensial}{}{}

\mavergleichskettedisp
{\vergleichskette
{ \omega
}
{ =} { x^ay^b dx + x^cy^d dy
}
{ } {
}
{ } {
}
{ } {
}
}
{}{}{}
dengan

\mavergleichskette
{\vergleichskette
{ a,b,c,d
}
{ \in }{ \N
}
{  }{
}
{    }{
}
{   }{
}
}
{}{}{.}

}
{} {}

\inputaufgabe
{}
{

Verifikasikan
teorema Green
melalui perhitungan eksplisit untuk himpunan
\mathl{T=[-2,2] \times [-2,2] \setminus U \left(   0 , 1  \right)}{}
\zusatzklammer {yakni persegi berpusat di asal dengan panjang sisi $4$, tanpa lingkaran satuan terbuka} {} {}
dan
\definitionsverweis {bentuk diferensial}{}{}

\mavergleichskette
{\vergleichskette
{ \omega
}
{ = }{(x-y)dx + xy dy
}
{  }{
}
{    }{
}
{   }{
}
}
{}{}{.}

}
{} {}

\inputaufgabegibtloesung
{}
{

Buktikan teorema Green untuk segitiga $D$ dengan titik-titik sudut

\mavergleichskette
{\vergleichskette
{ P_1,P_2,P_3
}
{ \in }{ \R^2
}
{  }{
}
{    }{
}
{   }{
}
}
{}{}{}
dan untuk bentuk diferensial
\mathl{xdy}{.}

}
{} {}

\inputaufgabe
{}
{

Misalkan
\maabbdisp {f} { [a,b] } { \R_{+}
} {}
suatu
\definitionsverweis {fungsi yang terdiferensialkan secara kontinu}{}{.}
Kita pandang
\definitionsverweis {subgraf}{}{}
sebagai suatu
\definitionsverweis {manifold berbatas}{}{},
dengan batas yang terdiri atas grafik, interval dasar, dan kedua sisi tegak, tetapi tanpa keempat titik sudutnya. Hitung luas subgraf dengan menggunakan kedua
\definitionsverweis {bentuk diferensial}{}{}

\mathl{xdy}{} dan
\mathl{ydx}{}
pada batas.

}
{} {}

\inputaufgabe
{}
{

Tentukan volume
\definitionsverweis {bola satuan tertutup}{}{}
tiga dimensi melalui suatu integral permukaan yang sesuai pada
\definitionsverweis {permukaan bola satuan}{}{}
$S^2$.

}
{} {}

\inputaufgabegibtloesung
{}
{

Kita tinjau
$2$-\definitionsverweis {bentuk diferensial}{}{}

\mavergleichskettedisp
{\vergleichskette
{ \omega
}
{ =} { x dy \wedge dz -y dx \wedge dz + z dx \wedge dy
}
{ } {
}
{ } {
}
{ } {
}
}
{}{}{}
pada bola satuan
\mathl{B  \left(  0 , 1  \right)}{.}
\aufzaehlungdreiabc{Tunjukkan bahwa $d \omega$ adalah tiga kali bentuk volume standar pada bola tersebut.
}{Tunjukkan bahwa
\mathl{\omega\vert_{S^2}}{} adalah bentuk luas standar pada permukaan bola satuan.
}{Hitung luas permukaan bola dari
volume bola
dengan menggunakan
teorema Stokes.
}

}
{} {}

\inputaufgabegibtloesung
{}
{

Misalkan $\omega$ suatu
$n-1$-\definitionsverweis {bentuk diferensial}{}{}
yang terdiferensialkan secara kontinu pada suatu
\definitionsverweis {himpunan terbuka}{}{}

\mavergleichskette
{\vergleichskette
{ U
}
{ \subseteq }{ \R^n
}
{  }{
}
{    }{
}
{   }{
}
}
{}{}{},
dan misalkan
\definitionsverweis {turunan eksterior}{}{}
diberikan oleh

\mavergleichskettedisp
{\vergleichskette
{d \omega
}
{ =} { fdx_1 \wedge \ldots \wedge dx_n
}
{ } {
}
{ } {
}
{ } {
}
}
{}{}{}
dengan suatu fungsi
\maabb {f} { U } {\R
} {.}
Tunjukkan untuk

\mavergleichskette
{\vergleichskette
{ P
}
{ \in }{ U
}
{  }{
}
{    }{
}
{   }{
}
}
{}{}{}
kesamaan

\mavergleichskettedisp
{\vergleichskette
{ f(P)
}
{ =} { { \frac{   1  }{  \beta_n  } } \cdot \operatorname{lim}_{   \epsilon \rightarrow  0  } \, { \frac{   1  }{  \epsilon^n } } \int_{S^{n-1}(P, \epsilon)} \omega
}
{ } {
}
{ } {
}
{ } {
}
}
{}{}{,}
dengan $\beta_n$ menyatakan
volume
bola satuan $n$-dimensional.

}
{} {}

Untuk suatu kasus khusus dari pernyataan di atas, lihat
Soal 58.23 (Analisis (Osnabrück 2021--2023)).

\inputaufgabe
{}
{

Misalkan $M$ suatu
\definitionsverweis {manifold diferensiabel}{}{}
dengan batas tak kosong. Tunjukkan bahwa terdapat suatu
\definitionsverweis {pemetaan diferensiabel}{}{}
\maabb {} { M } { \partial M
} {}
ada.

}
{} {}

\inputaufgabe
{}
{

Misalkan

\mavergleichskette
{\vergleichskette
{ H
}
{ \subseteq }{ \R^n
}
{  }{
}
{    }{
}
{   }{
}
}
{}{}{}
\zusatzklammer {
\mavergleichskette
{\vergleichskette
{ n
}
{ \geq }{ 1
}
{  }{
}
{    }{
}
{   }{
}
}
{}{}{}} {} {}
suatu
\definitionsverweis {setengah ruang}{}{.}
Tunjukkan bahwa terdapat suatu
\definitionsverweis {pemetaan diferensiabel}{}{}
\maabbdisp {\varphi} { H } { \partial H
} {}
yang
\definitionsverweis {pembatasannya}{}{}
pada
\mathl{\partial H}{} adalah
\definitionsverweis {identitas}{}{.}

}
{Bagaimana keadaannya untuk

\mavergleichskette
{\vergleichskette
{ n
}
{ = }{ 0
}
{  }{
}
{    }{
}
{   }{
}
}
{}{}{}?} {}

\inputaufgabe
{}
{

Kita tinjau
\definitionsverweis {manifold berbatas}{}{}

\mavergleichskette
{\vergleichskette
{ M
}
{ = }{ \R^n \setminus  U \left(   0 , 1  \right)
}
{  }{
}
{    }{
}
{   }{
}
}
{}{}{.}
Tunjukkan bahwa terdapat suatu pemetaan yang terdiferensialkan secara kontinu dari $M$ ke batas
\mathl{S \left(   0 ,  1  \right)}{}
yang merupakan identitas pada batas tersebut.

}
{} {}

Untuk soal berikut, Anda dapat menggunakan
Soal 21.8.

\inputaufgabe
{}
{

Kita tinjau
\definitionsverweis {manifold berbatas}{}{}

\mavergleichskettedisp
{\vergleichskette
{ M
}
{ =} { B  \left(  0 , 1 \right) \setminus \{  \left(  1   , \,   0                   \right) \}
}
{ \subseteq} { \R^2
}
{ } {
}
{ } {
}
}
{}{}{.}
Tunjukkan bahwa terdapat suatu
\definitionsverweis {pemetaan yang terdiferensialkan secara kontinu}{}{}
\maabbdisp {\varphi} { M } { \partial M
} {}
yang merupakan identitas pada batas.

}
{} {}

\inputaufgabe
{}
{

Tunjukkan bahwa pada suatu
\definitionsverweis {anulus}{}{}
terdapat pemetaan
\definitionsverweis {bijektif}{}{} yang
\definitionsverweis {terdiferensialkan secara kontinu}{}{}
dan tidak mempunyai
\definitionsverweis {titik tetap}{}{.}

}
{} {}

\inputaufgabe
{}
{

Tunjukkan bahwa pemetaan

\mavergleichskettedisp
{\vergleichskette
{ \varphi(x)
}
{ =} { x + \left( - \left\langle  x  , h(x)  \right\rangle + \sqrt{1 + \left\langle  x  , h(x)  \right\rangle^2 - \Vert { x} \Vert^2  }\right) \cdot h(x)
}
{ } {
}
{ } {
}
{ } {
}
}
{}{}{}
\zusatzklammer {dengan

\mavergleichskette
{\vergleichskette
{ h(x)
}
{ = }{ { \frac{   \psi(x)-x }{  \Vert { \psi (x)-x} \Vert  } }
}
{  }{
}
{    }{
}
{   }{
}
}
{}{}{}} {} {}
yang digunakan dalam pembuktian
teorema titik tetap Brouwer
memetakan bola satuan ke permukaan bola satuan.

}
{} {}

  \zwischenueberschrift{Tugas untuk dikumpulkan}

\inputaufgabe
{4}
{

Misalkan $D$ segitiga dengan titik sudut
$(0,2),\, (1,-1)$, dan $(-2,-1)$,
serta
\mathl{\tau =( 3x^2y^5-x \sin  y ) dx \wedge dy}{} suatu
$2$-\definitionsverweis {bentuk diferensial}{}{} pada $D$. Temukan suatu
\definitionsverweis {bentuk primitif}{}{} untuk $\tau$, lalu gunakan bentuk tersebut untuk menghitung
\mathl{\int_{ D  }  \tau}{} melalui suatu integral pada batas segitiga.

}
{} {}

\inputaufgabe
{6}
{

Kita tinjau kubus

\mavergleichskettedisp
{\vergleichskette
{ Q
}
{ =} { [-1,1]^3
}
{ \subseteq} { \R^3
}
{ } {
}
{ } {
}
}
{}{}{}
dan
$2$-\definitionsverweis {bentuk diferensial}{}{}

\mavergleichskettedisp
{\vergleichskette
{ \omega
}
{ =} { dx \wedge dy +y dx \wedge dz +x^2y^2z^2 dy \wedge dz
}
{ } {
}
{ } {
}
{ } {
}
}
{}{}{.}
Hitung
\mathl{d \omega}{} dan kedua integral
\mathkor {} {\int_{  \partial Q  }   \omega} {dan} {\int_{  Q  }  d \omega} {}
\zusatzklammer {secara terpisah} {} {.}

}
{} {}

\inputaufgabe
{4}
{

Hitung luas
\definitionsverweis {cakram satuan tertutup}{}{}
melalui suatu
\definitionsverweis {integral lintasan}{}{}
yang sesuai.

}
{} {}

\inputaufgabe
{2}
{

Tunjukkan bahwa pada suatu
\definitionsverweis {torus}{}{}
terdapat pemetaan
\definitionsverweis {bijektif}{}{} yang
\definitionsverweis {terdiferensialkan secara kontinu}{}{}
dan tidak mempunyai
\definitionsverweis {titik tetap}{}{.}

}
{} {}

\inputaufgabe
{3}
{

Misalkan $B$ adalah
\definitionsverweis {cakram satuan tertutup}{}{}
dan $K$ busur setengah lingkaran atas. Tunjukkan bahwa terdapat suatu
\definitionsverweis {pemetaan diferensiabel}{}{}
\maabbdisp {\varphi} { B } { K
} {}
yang
\definitionsverweis {pembatasannya}{}{}
pada
\mathl{K}{} adalah
\definitionsverweis {identitas}{}{.}

}
{} {}

\inputaufgabe
{3}
{

Misalkan

\mavergleichskette
{\vergleichskette
{ v,w
}
{ \in }{ \R^n
}
{  }{
}
{    }{
}
{   }{
}
}
{}{}{}
dengan
\mathkor {} {\Vert { v } \Vert \leq 1} {dan} {\Vert { w } \Vert = 1} {.}
Tunjukkan bahwa terdapat

\mavergleichskette
{\vergleichskette
{ a
}
{ \in }{ \R
}
{  }{
}
{    }{
}
{   }{
}
}
{}{}{}
sedemikian sehingga

\mavergleichskette
{\vergleichskette
{ \Vert {v+aw} \Vert
}
{ = }{ 1
}
{  }{
}
{    }{
}
{   }{
}
}
{}{}{.}

}
{} {}

\inputaufgabe
{4}
{

Kita tinjau
\definitionsverweis {manifold berbatas}{}{}

\mavergleichskettedisp
{\vergleichskette
{ M
}
{ =} { B  \left(  0 , 1 \right) \setminus \{  \left(  1   , \,   0                   \right) \}
}
{ \subseteq} { \R^2
}
{ } {
}
{ } {
}
}
{}{}{.}
Tunjukkan bahwa terdapat suatu pemetaan bebas titik tetap
\definitionsverweis {yang terdiferensialkan secara kontinu}{}{}
\maabbdisp {\varphi} { M } { M
} {}
ada.

}
{} {}


\section*{Solusi yang disediakan oleh sumber}
\addcontentsline{toc}{section}{Solusi yang disediakan oleh sumber}
\subsection*{Solusi Soal 23.6}
Karena menurut
Fakta (5)
penarikan balik bentuk kompatibel dengan turunan eksterior, maka
\mathl{\varphi^*\omega}{} juga eksak. Misalkan $\tau$ suatu bentuk diferensiabel pada $S$ dengan

\mavergleichskettedisp
{\vergleichskette
{d \tau
}
{ =} { \varphi^* \omega
}
{ } {
}
{ } {
}
{ } {
}
}
{}{}{.}
Menurut
teorema Stokes,

\mavergleichskettedisp
{\vergleichskette
{ \int_S \varphi^* \omega
}
{ =} { \int_S d \tau
}
{ =} { \int_{ \partial S } \tau
}
{ =} { \int_\emptyset \tau
}
{ =} { 0
}
}
{}{}{.}

\subsection*{Solusi Soal 23.13}
Misalkan

\mathdisp {P_1 = \begin{pmatrix} a_1  \\b_1   \end{pmatrix} ,\, P_2 = \begin{pmatrix} a_2  \\b_2   \end{pmatrix} \text{ dan } P_3 = \begin{pmatrix} a_3  \\b_3   \end{pmatrix}} {  }

adalah ketiga titik sudut. Karena
\mathl{d (xdy )= dx \wedge dy}{}, integral bentuk luas ini pada $D$ sama dengan luas segitiga tersebut. Dari $P_1$, segitiga itu direntang oleh kedua vektor
\mathkor {} {\begin{pmatrix} a_2-a_1  \\b_2-b_1   \end{pmatrix}} {dan} {\begin{pmatrix} a_3-a_1  \\b_3-b_1   \end{pmatrix}} {.}
Jadi, menurut rumus determinan untuk sebuah paralelotop, luasnya sama dengan

\mavergleichskettedisp
{\vergleichskette
{ { \frac{   1 }{  2 } } \betrag { (a_2-a_1)(b_3-b_1) -(b_2-b_1)  (a_3-a_1)   }
}
{ =} { { \frac{   1 }{  2 } } \betrag {  a_2b_3-a_2b_1 -a_1b_3 -a_3b_2 +a_1b_2 +a_3b_1   }
}
{ } {
}
{ } {
}
{ } {
}
}
{}{}{.}
Jika kedua vektor tersebut merepresentasikan orientasi standar, sebagaimana akan kita asumsikan mulai sekarang, tanda nilai mutlak dapat dihilangkan.

Sekarang kita hitung integral lintasan dari
\mathl{xdy}{} di sepanjang batas segitiga yang ditelusuri berlawanan arah jarum jam. Lintasan bergerak dari $P_1$ ke $P_2$, lalu ke $P_3$, dan kembali ke $P_1$
\zusatzklammer {ini sesuai dengan arah berlawanan jarum jam untuk orientasi yang telah ditetapkan} {} {.}
Lintasan-lintasan linear ini adalah
\zusatzklammer {masing-masing didefinisikan pada interval satuan} {} {}

\mathdisp {\gamma_1(t)=P_1+t(P_2-P_1) = \begin{pmatrix} a_1  \\b_1   \end{pmatrix} + t \begin{pmatrix} a_2-a_1  \\b_2-b_1   \end{pmatrix}} { , }

\mathdisp {\gamma_2(t) = P_2+t(P_3-P_2) = \begin{pmatrix} a_2  \\b_2   \end{pmatrix} + t \begin{pmatrix} a_3-a_2  \\b_3-b_2   \end{pmatrix}} {  }

dan

\mathdisp {\gamma_3(t) = P_3+t(P_1-P_3)= \begin{pmatrix} a_3  \\b_3   \end{pmatrix} + t \begin{pmatrix} a_1-a_3  \\b_1-b_3   \end{pmatrix}} { . }

Berlaku

\mavergleichskettealign
{\vergleichskettealign
{ \int_{\gamma_1}  xdy
}
{ =} { \int_0^1  (a_1 +t (a_2-a_1)) (b_2-b_1) dt
}
{ =} { (  a_1 (b_2-b_1)t + { \frac{   1 }{  2 } } (a_2-a_1)(b_2-b_1) t^2) \vert_{  0 }^{  1 }
}
{ =} { a_1 (b_2-b_1)  + { \frac{   1 }{  2 } } (a_2-a_1)(b_2-b_1)
}
{ =} { { \frac{   1 }{  2 } } (a_2+a_1)(b_2-b_1)
}
}
{}
{}{.}
Demikian pula,

\mavergleichskettedisp
{\vergleichskette
{ \int_{\gamma_2}  xdy
}
{ =} { { \frac{   1 }{  2 } } (a_3+a_2)(b_3-b_2)
}
{ } {
}
{ } {
}
{ } {
}
}
{}{}{}
dan

\mavergleichskettedisp
{\vergleichskette
{ \int_{\gamma_3}  xdy
}
{ =} { { \frac{   1 }{  2 } } (a_1+a_3)(b_1-b_3)
}
{ } {
}
{ } {
}
{ } {
}
}
{}{}{.}
Jumlah ketiga integral lintasan ini adalah separuh dari

\mavergleichskettealigndrucklinks
{\vergleichskettealigndrucklinks
{(a_2+a_1)(b_2-b_1) + (a_3+a_2)(b_3-b_2) + (a_1+a_3)(b_1-b_3)
}
{ =} { a_2b_2 -a_2b_1+a_1b_2-a_1b_1 +a_3b_3-a_3b_2 +a_2b_3 -a_2b_2+a_1b_1-a_1b_3+a_3b_1-a_3b_3
}
{ =} { -a_2b_1+a_1b_2 -a_3b_2 +a_2b_3 -a_1b_3+a_3b_1
}
{ } {
}
{ } {
}
}
{}{}{,}
sehingga kedua integral itu berimpit.

\subsection*{Solusi Soal 23.16}
\input{generated/worksheet23_exercise16_solution.id.build.tex}
\subsection*{Solusi Soal 23.17}
Kita terapkan
teorema Stokes
pada bola-bola
\mathl{B^n(P, \epsilon)}{} dengan batas
\mathl{S^{n-1}(P, \epsilon)}{}
\zusatzklammer {berpusat di $P$ dan berjari-jari $\epsilon$} {} {},
sehingga diperoleh

\mavergleichskettealign
{\vergleichskettealign
{ \operatorname{lim}_{   \epsilon \rightarrow  0  } \, { \frac{   1 }{  \epsilon^n } }     \int_{S^{n-1}(P, \epsilon)} \omega
}
{ =} { \operatorname{lim}_{   \epsilon \rightarrow  0  } \, { \frac{   1 }{  \epsilon^n } }     \int_{ \partial B^{n}(P, \epsilon)}   \omega
}
{ =} { \operatorname{lim}_{   \epsilon \rightarrow  0  } \, { \frac{   1 }{  \epsilon^n } }     \int_{B^{n}(P, \epsilon)} d \omega
}
{ =} { \operatorname{lim}_{   \epsilon \rightarrow  0  } \, { \frac{   1 }{  \epsilon^n } }     \int_{B^{n}(P, \epsilon)} f d \lambda^n
}
{ =} {\beta_n \cdot f(P)
}
}
{}
{}{,}
di mana pada kesamaan terakhir kita menggunakan kekontinuan $f$.


\part{Unit 24}
\chapter[Kuliah 24]{Kuliah 24: Koneksi dan Turunan Vertikal}
\setcounter{section}{24}

Misalkan

\mavergleichskette
{\vergleichskette
{ Y
}
{ \subseteq }{ \R^n
}
{  }{
}
{    }{
}
{   }{
}
}
{}{}{ }
adalah suatu hipermuka terdiferensialkan. Bagi setiap titik

\mavergleichskette
{\vergleichskette
{ P
}
{ \in }{ Y
}
{  }{
}
{    }{
}
{   }{
}
}
{}{}{,}
\definitionsverweis {ruang singgung}{}{}
merupakan subruang vektor

\mavergleichskette
{\vergleichskette
{ T_PY
}
{ \subseteq }{ \R^n
}
{  }{
}
{    }{
}
{   }{
}
}
{}{}{}
di ruang ambien. Dengan menggunakan hasil kali dalam standar pada $\R^n$, kita memperoleh
\definitionsverweis {proyeksi ortogonal}{}{}
\maabbdisp {\pi_P} {\R^n } { T_PY
} {,}
yang memungkinkan data di ruang ambien ditafsirkan sebagai data tangensial pada hipermuka. Dari konstruksi ini muncul berbagai konsep penting pada $Y$, seperti
\definitionsverweis {percepatan tangensial}{}{,}
\definitionsverweis {kurva geodesik}{}{,}
\definitionsverweis {kelengkungan}{}{,}
\definitionsverweis {turunan kovarian suatu medan vektor}{}{,}
\definitionsverweis {medan vektor paralel}{}{}
serta
\definitionsverweis {transpor paralel}{}{.}
Di sisi lain, hasil kali dalam standar tersebut juga menginduksi
\definitionsverweis {struktur Riemann}{}{}
pada $Y$. Pertanyaannya kemudian ialah apakah konsep-konsep tadi dapat dibangun secara intrinsik hanya dari $Y$ dan struktur Riemann-nya
\zusatzklammer {misalnya untuk
\definitionsverweis {permukaan hiperbolik}{}{}} {} {.}
Untuk menjawabnya, diperlukan perangkat baru, yaitu koneksi pada bundel vektor, terutama pada bundel tangen suatu
\definitionsverweis {manifold diferensiabel}{}{.}
Pada manifold Riemann, struktur ini menghasilkan secara kanonik \stichwort {koneksi Levi-Civita} {}
\zusatzklammer {lihat
Teorema 26.12} {} {,}
yang memungkinkan seluruh konsep di atas dirumuskan secara intrinsik.

  \zwischenueberschrift{Koneksi}

Misalkan $X$ adalah
\definitionsverweis {manifold diferensiabel}{}{}
riil, dan
\maabbdisp {p} { E } { X
} {}
adalah bundel vektor di atas $X$ yang ruang totalnya juga diperlakukan sebagai manifold diferensiabel. Adapun
\definitionsverweis {pemetaan singgung}{}{}
\maabb {T(p)} {TE} {TX
} {,}
berada dalam diagram komutatif

\mathdisp {\begin{matrix} TE & \stackrel{ T(p) }{\longrightarrow} & TX &  \\ \downarrow & & \downarrow  & \\  E  & \stackrel{ p }{\longrightarrow} &  X  & \! \! \! \!\!\!\!\!   \\ \end{matrix}} {  }

dan, bagi setiap titik

\mavergleichskette
{\vergleichskette
{ e
}
{ \in }{ E
}
{  }{
}
{    }{
}
{   }{
}
}
{}{}{}
dengan titik dasar

\mavergleichskette
{\vergleichskette
{ x
}
{ = }{ p(e)
}
{  }{
}
{    }{
}
{   }{
}
}
{}{}{,}
\definitionsverweis {pemetaan singgung}{}{}
\definitionsverweis {linear}{}{}
\maabbdisp {T_e(p)} {T_e E} { T_{x}X
} {}
bersifat surjektif. Salah satu alasannya ialah bahwa secara lokal, melalui setiap titik

\mavergleichskette
{\vergleichskette
{ e
}
{ \in }{ E
}
{  }{
}
{    }{
}
{   }{
}
}
{}{}{,}
terdapat suatu penampang bundel vektor $E$. Agar domain dan kodomain memiliki ruang dasar yang sama, kita memandang $T(p)$ sebagai
\definitionsverweis {pemetaan bundel}{}{}
\maabbdisp {Tp} {TE} { p^* TX
} {}
antara bundel vektor di atas $E$. Dalam notasi ini,

\mavergleichskettedisp
{\vergleichskette
{ p^*TX
}
{ =} { E \oplus TX
}
{ =} { E \times_X TX
}
{ } {
}
{ } {
}
}
{}{}{}
merupakan
\definitionsverweis {jumlah langsung}{}{}
bundel vektor di atas $X$; lihat
Soal 12.24.
Di sini, $p^*TX$ dipandang sebagai bundel vektor di atas $E$. Serat $p^*TX$ di atas titik

\mavergleichskette
{\vergleichskette
{ e
}
{ \in }{ E
}
{  }{
}
{    }{
}
{   }{
}
}
{}{}{}
dengan titik dasar

\mavergleichskette
{\vergleichskette
{ x
}
{ = }{ p(e)
}
{  }{
}
{    }{
}
{   }{
}
}
{}{}{}
hanyalah $T_{x}X$, sedangkan pemetaan pada serat tersebut tetaplah
\maabbdisp {T_e (p)} { T_e E } { T_{x} X
} {,}
sehingga di dalam $p^*TX$ terdapat satu salinan $T_{x} X$ untuk setiap

\mavergleichskette
{\vergleichskette
{ e
}
{ \in }{ E
}
{  }{
}
{    }{
}
{   }{
}
}
{}{}{.}

\inputbemerkung
{}
{

Misalkan
\maabb {p} { E } { X
} {}
adalah
\definitionsverweis {bundel vektor}{}{}
diferensiabel di atas
\definitionsverweis {manifold diferensiabel}{}{}
$X$. Ambil himpunan terbuka

\mavergleichskette
{\vergleichskette
{ U
}
{ \subseteq }{ X
}
{  }{
}
{    }{
}
{   }{
}
}
{}{}{,}
tempat $E$ dapat ditrivialkan, sehingga

\mavergleichskette
{\vergleichskette
{ E  \vert_U
}
{ \cong }{ U \times W
}
{  }{
}
{    }{
}
{   }{
}
}
{}{}{}
dengan ruang vektor atas $\R$ yang kita namakan $W$. Dengan demikian,

\mavergleichskettedisp
{\vergleichskette
{ T  \left(  E \vert_U  \right)
}
{ =} { T(U \times  W )
}
{ =} { T U \times T W
}
{ =} { TU \times W \times W
}
{ } {
}
}
{}{}{,}
Sementara itu,
\definitionsverweis {pemetaan singgung}{}{}
yang berkaitan dengan
\maabb {p_U} {E \vert_U } { U
} {}
merupakan proyeksi ke $TU$. Untuk menuliskannya sebagai homomorfisme bundel di atas $U\times W$, kodomainnya diidentifikasi sebagai

\mavergleichskettedisp
{\vergleichskette
{ p^*TU
}
{ =} { TU \times W
}
{ =} { TU \times_U (U \times W)
}
{ } {
}
{ } {
}
}
{}{}{}
dan pemetaannya ialah
\maabbdisp {} {T  \left(  E \vert_U  \right)  = TU \times W \times W  } { TU \times W
} {}
yang mengambil proyeksi ke $TU$ sekaligus proyeksi kedua ke $W$.

}

\inputdefinition
{}
{

Misalkan
\maabb {p} { E } { X
} {}
adalah
\definitionsverweis {bundel vektor diferensiabel}{}{}
di atas
\definitionsverweis {manifold diferensiabel}{}{}
$X$.
\definitionsverweis {Bundel kernel}{}{}
dari
\definitionsverweis {homomorfisme bundel}{}{}
\definitionsverweis {surjektif}{}{}
\maabbdisp {T(p)} {TE} {p^*TX
} {}
di atas $E$ disebut \definitionswort {bundel vertikal}{.} Notasinya adalah $V$.

}

Menurut Soal 12.19, bundel vertikal adalah bundel vektor di atas $E$; lebih khusus lagi, ia merupakan
\definitionsverweis {subbundel vektor}{}{}
dari $TE$.



\inputfaktbeweis
{Mannigfaltigkeit/Vektorbündel/R/Vertikalbündel/Eigenschaften/Fakt}
{Lema}
{}
{

\faktsituation {Misalkan
\maabb {p} { E } { X
} {}
adalah
\definitionsverweis {bundel vektor diferensiabel}{}{}
di atas
$C^2$-\definitionsverweis {manifold diferensiabel}{}{}
$X$.}
\faktuebergang {Pernyataan berikut berlaku.}
\faktfolgerung {\aufzaehlungdrei{Terdapat
\definitionsverweis {barisan eksak pendek}{}{}

\mathdisp {0 \, \stackrel{  }{\longrightarrow} \,  V   \, \stackrel{  }{ \longrightarrow}   \,  TE  \, \stackrel{  }{\longrightarrow} \, p^* TX \,\stackrel{  } {\longrightarrow}  \, 0} {  }

yang terdiri dari
\definitionsverweis {homomorfisme}{}{}
bundel vektor di atas $E$.
}{Untuk
\definitionsverweis {bundel vertikal}{}{}
terdapat
\definitionsverweis {isomorfisme bundel}{}{}

\mavergleichskette
{\vergleichskette
{ V \cong p^*E
}
{ = }{ E \oplus E
}
{  }{
}
{    }{
}
{   }{
}
}
{}{}{.}
}{Bagi setiap titik

\mavergleichskette
{\vergleichskette
{ e
}
{ \in }{ E
}
{  }{
}
{    }{
}
{   }{
}
}
{}{}{,}
diperoleh barisan eksak pendek ruang vektor berikut:

\mathdisp {0 \, \stackrel{  }{\longrightarrow} \, E_{p(e)}   \, \stackrel{  }{ \longrightarrow}   \,  T_eE  \, \stackrel{  }{\longrightarrow} \, T_{p(e)}X \,\stackrel{  } {\longrightarrow}  \, 0} {  }

}}
\faktzusatz {}
\faktzusatz {}

}
{

\aufzaehlungzwei {Pernyataan ini langsung mengikuti definisi.
} {Terdapat homomorfisme injektif bundel vektor
\maabbdisp {} {p^*E = E \oplus E} { TE
} {}
di atas $E$ yang memetakan
\mathl{(P,u)}{,}

\mavergleichskette
{\vergleichskette
{ u
}
{ \in }{ E_P
}
{  }{
}
{    }{
}
{   }{
}
}
{}{}{,}
ke vektor singgung yang diwakili oleh kurva diferensiabel
\zusatzklammer {vertikal} {} {}
\maabbeledisp {} {\R} {E
} { t } {(P,tu)
} {,}
ini. Selanjutnya, oleh pemetaan
\maabbdisp {} {TE} { p^*TX
} {}
vektor singgung tersebut dikirim ke $0$, sebab kurva wakilnya dikirim ke $P$ oleh
\maabb {p} { E } { X
} {}
.
Oleh karena itu,

\mavergleichskette
{\vergleichskette
{ E
}
{ \subseteq }{ V
}
{  }{
}
{    }{
}
{   }{
}
}
{}{}{,}
dan perbandingan rank menunjukkan bahwa inklusi ini sebenarnya suatu isomorfisme.
}

}

Kasus yang sering dijumpai adalah ketika $E$ merupakan bundel tangen dari $X$. Dalam kasus ini diperoleh barisan eksak pendek

\mathdisp {0 \, \stackrel{  }{\longrightarrow} \, p^* TX  \, \stackrel{  }{ \longrightarrow}   \, TTX  \, \stackrel{  }{\longrightarrow} \, p^* TX \,\stackrel{  } {\longrightarrow}  \, 0} {  }

Oleh sebab itu, kedua salinan bundel tangen di ruas kiri dan kanan harus selalu dibedakan dengan cermat.

\inputbemerkung
{}
{

Misalkan
\maabbdisp {f} {\R^n } { \R
} {}
adalah fungsi yang dua kali
\definitionsverweis {terdiferensialkan secara kontinu}{}{,}
dan

\mavergleichskette
{\vergleichskette
{ Y
}
{ = }{ f^{-1} (0)
}
{  }{
}
{    }{
}
{   }{
}
}
{}{}{}
adalah
\definitionsverweis {hipermuka terdiferensialkan}{}{}
terkait yang
\definitionsverweis {reguler}{}{}
pada setiap titik.
\definitionsverweis {Bundel tangen}{}{}
$TY$ dapat pula dinyatakan sebagai serat bersama dari fungsi-fungsi, yakni

\mavergleichskettealignhandlinks
{\vergleichskettealignhandlinks
{ TY
}
{ =} { \left\{  (P,u)   \mid   f(P) = 0 , \,  \left(  D f   \right)_{ P } \left(  u  \right) = 0       \right\}
}
{ =} { \left\{  (P,u_1 , \ldots , u_n)   \mid   f(P) = 0 , \,  \sum_{i = 1}^n u_i { \frac{   \partial   f   }{  \partial   x_i   } } ( P) = 0       \right\}
}
{ \subseteq} { \R^n \times \R^n
}
{ } {
}
}
{}
{}{,}
sebagaimana dalam
Soal 10.15.
Dalam rumusan ini, $f$ dipandang sebagai fungsi pada $\R^{2n}$ yang hanya bergantung pada $n$ variabel pertama, dan kita definisikan

\mavergleichskettedisp
{\vergleichskette
{ g(P,u)
}
{ =} { \sum_{i = 1}^n u_i { \frac{   \partial   f   }{  \partial   x_i   } } ( P)
}
{ } {
}
{ } {
}
{ } {
}
}
{}{}{.}
Bundel tangen kedua $TTY$, yaitu bundel tangen dari $TY$, selanjutnya dapat dideskripsikan dengan
\definitionsverweis {diferensial total}{}{}
dari pemetaan
\maabbdisp {\varphi = (f,g)} {\R^{2n}} { \R^2
} {}
sebagai berikut:

\mavergleichskettealignhandlinks
{\vergleichskettealignhandlinks
{ TTY
}
{ =} { \left\{  (P,u, v,w)  \mid   (P,u) \in TY , \,  \left(  D\varphi  \right)_{(P,u)} \left(  (v,w)  \right) = 0       \right\}
}
{ \subseteq} { \R^{4n}
}
{ } {
}
{ } {}
}
{}
{}{.}
\definitionsverweis {Matriks Jacobi}{}{,}
yang merepresentasikan diferensial total $D\varphi$, ialah

\mathdisp {\begin{pmatrix}   { \frac{   \partial   f   }{  \partial   x_1   } } ( P)   &   \ldots &  { \frac{   \partial   f   }{  \partial   x_n   } } ( P)  &   0 &  \ldots &  0  \\   \sum_{i = 1}^n u_i   { \frac{   \partial    }{  \partial   x_1   } } { \frac{   \partial   f   }{  \partial   x_i   } } ( P)  &  \ldots  &   \sum_{i = 1}^n u_i   { \frac{   \partial    }{  \partial   x_n   } } { \frac{   \partial   f   }{  \partial   x_i   } } ( P)   &  { \frac{   \partial   f   }{  \partial   x_1   } } ( P)  &  \ldots &  { \frac{   \partial   f   }{  \partial   x_n   } } ( P)    \end{pmatrix}} { . }

Di samping dua syarat yang mendefinisikan $TY$ dan hanya menyangkut
\mathkor {} {P} {dan} {u} {,}
diperoleh dua syarat tambahan yang melibatkan seluruh variabel, yaitu

\mavergleichskettedisp
{\vergleichskette
{ \sum_{j = 1}^n v_j { \frac{   \partial   f   }{  \partial   x_j   } } ( P)
}
{ =} { 0
}
{ } {
}
{ } {
}
{ } {
}
}
{}{}{}
dan

\mavergleichskettedisp
{\vergleichskette
{ \sum_{j = 1}^n v_j  \left(  \sum_{i = 1}^n u_i { \frac{   \partial    }{  \partial   x_i   } }  { \frac{   \partial   f   }{  \partial   x_j   } } ( P)  \right)      + \sum_{k = 1}^n w_k { \frac{   \partial   f   }{  \partial   x_k   } } ( P)
}
{ =} { 0
}
{ } {
}
{ } {
}
{ } {
}
}
{}{}{.}

}

\inputbemerkung
{}
{

Misalkan
\maabbdisp {f} {\R^n } { \R
} {}
adalah fungsi yang dua kali
\definitionsverweis {terdiferensialkan secara kontinu}{}{,}
dan

\mavergleichskette
{\vergleichskette
{ Y
}
{ = }{ f^{-1} (0)
}
{  }{
}
{    }{
}
{   }{
}
}
{}{}{}
adalah
\definitionsverweis {hipermuka terdiferensialkan}{}{}
terkait yang
\definitionsverweis {reguler}{}{}
pada setiap titik. Kita gunakan deskripsi bundel tangen kedua $TTY$ dari
Catatan 24.4
beserta notasi di sana. Bundel tangen tarik balik $\pi^*TY$ terhadap
\maabbdisp {\pi} {TY} {Y
} {,}
yang juga merupakan
\definitionsverweis {bundel vertikal}{}{,}
adalah

\mavergleichskettealign
{\vergleichskettealign
{ V
}
{ \cong} { \pi^*TY
}
{ =} { TY \times_Y TY
}
{ =} { \left\{ (P,u,z)  \mid  f(P) = 0 , \,  \sum_{i = 1}^n u_i { \frac{   \partial   f   }{  \partial   x_i   } } ( P) = 0   , \,  \sum_{k = 1}^n z_k { \frac{   \partial   f   }{  \partial   x_k   } } ( P) = 0    \right\}
}
{ } {
}
}
{}
{}{.}
Sekarang kita tentukan bentuk pemetaan singgung
\maabb {T(\pi)} {TTY } { \pi^* TY
} {}
serta pembenaman bundel vertikal ke dalam $TTY$. Berdasarkan
Soal 24.2,
diperoleh
\maabbeledisp {T(\pi)} {TTY} { \pi^* TY
} {(P,u,v,w)} { (P,u,v)
} {.}
Kernel pemetaan ini adalah

\mavergleichskettedisp
{\vergleichskette
{ V
}
{ =} { \left\{  (P,u,0,w)   \mid   (P,u,0,w) \in TTY          \right\}
}
{ \subset} { TTY
}
{ } {
}
{ } {
}
}
{}{}{,}
dan langsung terlihat bahwa
\mathl{(P,u,w)}{} memenuhi syarat bagi
\mathl{\pi^*TY}{}.
Oleh sebab itu, kedua homomorfisme bundel dalam barisan eksak pendek

\mathdisp {0 \, \stackrel{  }{\longrightarrow} \,  \pi^*TY  \, \stackrel{  }{ \longrightarrow}   \, TTY  \, \stackrel{  }{\longrightarrow} \,  \pi^*TY \,\stackrel{  } {\longrightarrow}  \, 0} {  }

dapat ditulis dalam koordinat sebagai

\mathl{(P,u,w) \mapsto  (P,u,0,w)}{} dan
\mathl{(P,u,v,w) \mapsto (P,u,v)}{}.

}

Jika $E$ merupakan bundel trivial, yaitu

\mavergleichskette
{\vergleichskette
{ E
}
{ = }{ X \times \R^r
}
{  }{
}
{    }{
}
{   }{
}
}
{}{}{,}
maka terdapat penguraian langsung
\zusatzklammer {di atas $E$} {} {}

\mavergleichskette
{\vergleichskette
{ TE
}
{ = }{ V \oplus p^* TX
}
{  }{
}
{    }{
}
{   }{
}
}
{}{}{.}
Di setiap titik

\mavergleichskette
{\vergleichskette
{ e
}
{ \in }{ E
}
{  }{
}
{    }{
}
{   }{
}
}
{}{}{,}
vektor singgung

\mavergleichskette
{\vergleichskette
{ t
}
{ \in }{ V
}
{  }{
}
{    }{
}
{   }{
}
}
{}{}{}
mewakili arah vertikal, sedangkan

\mavergleichskette
{\vergleichskette
{ t
}
{ \in }{ p^* T_xX
}
{  }{
}
{    }{
}
{   }{
}
}
{}{}{}
mewakili \anfuehrung{arah horizontal}{.}
Pada bundel vektor sembarang, suatu trivialisasi lokal menguraikan ruang singgung
\mathl{T_eE}{} sebagai jumlah langsung ruang vertikal dan ruang horizontal. Namun, subruang horizontal yang dihasilkan sangat bergantung pada trivialisasi yang dipilih. Tidak seperti subruang vertikal
\zusatzklammer {yang bersama-sama membentuk bundel vertikal} {} {,}
subruang horizontal ini belum tentu menyatu menjadi suatu subbundel. Keberadaan subbundel horizontal semacam inilah yang dirumuskan oleh konsep koneksi.

\inputdefinition
{}
{

Misalkan $X$ adalah
\definitionsverweis {manifold diferensiabel}{}{}
dan $E$ adalah
\definitionsverweis {bundel vektor diferensiabel}{}{}
di atas $X$. Suatu
\definitionswort {koneksi}{}
pada $E$ didefinisikan sebagai penguraian jumlah langsung
\definitionsverweis {bundel tangen}{}{}

\mavergleichskette
{\vergleichskette
{ TE
}
{ = }{ V \oplus H
}
{  }{
}
{    }{
}
{   }{
}
}
{}{}{}
menjadi dua
\definitionsverweis {subbundel vektor}{}{}
\mathkor {} {V} {dan} {H} {,}
dengan ketentuan bahwa

\mavergleichskette
{\vergleichskette
{ V
}
{ \subseteq }{ TE
}
{  }{
}
{    }{
}
{   }{
}
}
{}{}{}
merupakan
\definitionsverweis {bundel vertikal}{}{.}
Subbundel

\mavergleichskette
{\vergleichskette
{ H
}
{ \subseteq }{ TE
}
{  }{
}
{    }{
}
{   }{
}
}
{}{}{}
disebut
\definitionswort {bundel horizontal}{}
koneksi tersebut.

}

Koneksi biasanya dilambangkan dengan $\nabla$, khususnya ketika proses diferensiasi yang ditentukannya hendak ditekankan. Jadi, menurut definisi, koneksi hanyalah subbundel dari $TE$ yang komplementer terhadap bundel vertikal. Ada beberapa rumusan lain yang ekuivalen. Memberikan suatu koneksi sama artinya dengan memberikan homomorfisme bundel
\zusatzklammer {disebut \stichwort {proyeksi vertikal} {} koneksi} {} {}
\maabbdisp {\pi_{\text{vert} }} {TE} {V
} {}
dengan

\mavergleichskettedisp
{\vergleichskette
{ \pi_{\text{vert} }  \circ \iota
}
{ =} {
\operatorname{Id}_{  V  }
}
{ } {
}
{ } {
}
{ } {
}
}
{}{}{,}
di mana $\iota$ adalah pembenaman $V$ ke dalam $TE$. Dalam rumusan ini, bundel horizontal ialah kernel $\pi_{\text{vert}}$. Sebaliknya, penguraian jumlah langsung menyediakan proyeksi ke masing-masing suku. Penguraian seperti itu juga disebut \stichwort {pemisahan} {} dari barisan eksak pendek. Setiap koneksi dapat dibatasi ke himpunan terbuka

\mavergleichskette
{\vergleichskette
{ U
}
{ \subseteq }{ X
}
{  }{
}
{    }{
}
{   }{
}
}
{}{}{.}

\inputbeispiel{}
{

Pertimbangkan bundel trivial
\maabbdisp {p} {X \times W = W_X } { X
} {,}
dengan $W$ ruang vektor riil
\definitionsverweis {berdimensi}{}{}
$r$ dan $X$ sebuah
\definitionsverweis {manifold diferensiabel}{}{,}
kita memperoleh barisan eksak pendek

\mathdisp {0 \, \stackrel{  }{\longrightarrow} \, p^*(X \times W)  =  X \times W \times W   \, \stackrel{  }{ \longrightarrow}   \, p^* T(X \times W )  = TX \times W \times W   \, \stackrel{  }{\longrightarrow} \, p^* TX  =  TX \times W   \,\stackrel{  } {\longrightarrow}  \, 0} {  }

bundel vektor di atas
\mathl{X \times W}{.}
Semua produk di sini dipahami sebagai produk Kartesius biasa, setelah dilakukan identifikasi

\mavergleichskette
{\vergleichskette
{ (X \times W) \times_X (X \times W)
}
{ = }{ X \times W \times W
}
{  }{
}
{    }{
}
{   }{
}
}
{}{}{}
di sebelah kiri serta identifikasi

\mavergleichskette
{\vergleichskette
{ TX \times_X (X \times W)
}
{ = }{ TX \times W
}
{  }{
}
{    }{
}
{   }{
}
}
{}{}{}
di sebelah kanan. Pada titik

\mavergleichskette
{\vergleichskette
{ (x,w)
}
{ \in }{ X \times W
}
{  }{
}
{    }{
}
{   }{
}
}
{}{}{,}
barisan eksak pada serat berbentuk

\mathdisp {0 \, \stackrel{  }{\longrightarrow} \,  W   \, \stackrel{  }{ \longrightarrow}   \,  T_xX \times W   \, \stackrel{  }{\longrightarrow} \,  T_xX \,\stackrel{  } {\longrightarrow}  \, 0} { . }

Pemisahan baku dari barisan eksak pendek bundel vektor ini disebut \stichwort {koneksi trivial} {} pada bundel trivial tersebut.

}

\inputbemerkung
{}
{

Misalkan
\maabbdisp {f} {\R^n } { \R
} {}
adalah fungsi yang dua kali
\definitionsverweis {terdiferensialkan secara kontinu}{}{,}
dan

\mavergleichskette
{\vergleichskette
{ Y
}
{ = }{ f^{-1} (0)
}
{  }{
}
{    }{
}
{   }{
}
}
{}{}{}
adalah
\definitionsverweis {hipermuka terdiferensialkan}{}{}
terkait yang
\definitionsverweis {reguler}{}{}
pada setiap titik. Kita kembali memakai deskripsi bundel tangen kedua $TTY$ dari
Catatan 24.4
dan homomorfisme bundel dari
Catatan 24.5.
Melalui proyeksi ortogonal
\maabb {} {\R^n } { T_PY
} {}
diperoleh proyeksi vertikal
\maabbeledisp {} {TTY} { V = p^*TY
} {(P,u,v,w)} { (P,u, \pi_P (w) )
} {,}
pada
\definitionsverweis {bundel vertikal}{}{.}
Memang, $\pi_P(w)$ berada dalam ruang singgung $Y$. Apabila kita mulai dari vektor

\mavergleichskette
{\vergleichskette
{ (P,u,w)
}
{ \in }{ V
}
{  }{
}
{    }{
}
{   }{
}
}
{}{}{,}
komponen terakhir $w$ pada
\mathl{(P,u,0,w)}{} sudah otomatis berada dalam ruang singgung. Karena itu, proyeksi ortogonal mengirim unsur ini kembali ke $(P,u,w)$. Jadi, konstruksi tersebut mendefinisikan
\definitionsverweis {koneksi}{}{}
pada $TY$.

}

  \zwischenueberschrift{Turunan vertikal}

\inputdefinition
{}
{

Misalkan $X$ adalah
\definitionsverweis {manifold diferensiabel}{}{}
dan $E$ adalah
\definitionsverweis {bundel vektor diferensiabel}{}{}
di atas $X$ yang dilengkapi dengan suatu
\definitionsverweis {koneksi}{}{.}
\definitionswort {turunan vertikal}{}
$\nabla$ yang berkaitan dengan koneksi tersebut adalah pemetaan
\maabbeledisp {\nabla} {C^1(E)} { C^0(E \otimes T^* X )
} { s } { \nabla(s) = \pi_{\text{vert} }  \circ T(s)
} {.}

}
Dengan kata lain, setiap penampang diferensiabel $s$ dari $E$
\zusatzklammer {di atas $X$ ataupun di atas himpunan terbuka sembarang

\mavergleichskettek
{\vergleichskettek
{ U
}
{ \subseteq }{ X
}
{  }{
}
{    }{
}
{   }{
}
}
{}{}{}} {} {}
dipetakan ke komposisi

\mathdisp {TX \stackrel{T(s)}{ \longrightarrow } TE \stackrel{ \pi_{\text{vert} } }{\longrightarrow} V \cong p^*E \longrightarrow E} {  }

dengan identifikasi

\mavergleichskette
{\vergleichskette
{ \operatorname{Hom}(TX,E)
}
{ \cong }{ E \otimes T^*X
}
{  }{
}
{    }{
}
{   }{
}
}
{}{}{.}

Jika diberikan pula
\definitionsverweis {medan vektor}{}{}
$F$ pada $X$, yaitu penampang
\maabb {F} { X } { TX
} {}
pada
\definitionsverweis {bundel tangen}{}{}
$TX$, maka diperoleh pemetaan
\maabbeledisp {\nabla_F} {C^1(E)} {C^0(E)
} {s } {\nabla_F(s) = \nabla (s) \circ F
} {,}
yang disebut \stichwort {turunan vertikal} {} dalam arah $F$. Ketergantungan terhadap $F$ hanya bersifat titik demi titik. Sebaliknya, ketergantungan terhadap $s$ bersifat infinitesimal karena melibatkan pemetaan singgung
\mathl{T(s)}{}.



\inputfaktbeweis
{Mannigfaltigkeit/Vektorbündel/R/Zusammenhang/Ableitung bezüglich Vektorfeld/Eigenschaften/Fakt}
{Lema}
{}
{

\faktsituation {Misalkan $X$ adalah
\definitionsverweis {manifold diferensiabel}{}{}
dan $E$ adalah
\definitionsverweis {bundel vektor diferensiabel}{}{}
di atas $X$ yang dilengkapi dengan suatu
\definitionsverweis {koneksi}{}{.}
Tuliskan
\maabbdisp {\nabla_F} {C^1(E)} {C^0(E)
} {}
untuk operator turunan yang bersesuaian dengan
\definitionsverweis {medan vektor}{}{}
$F$.}
\faktuebergang {Sifat-sifat berikut berlaku.}
\faktfolgerung {\aufzaehlungzwei {
\mathl{\left(  \nabla_F(s)  \right)  (P)}{} hanya bergantung pada $F(P)$
\zusatzklammer {serta pada $s$} {} {.}
} {Berlaku

\mavergleichskettedisp
{\vergleichskette
{ \nabla_{g_1F_1+g_2F_2 }
}
{ =} { g_1 \nabla_{F_1 } +g_2 \nabla_{F_2 }
}
{ } {
}
{ } {
}
{ } {
}
}
{}{}{}
untuk medan vektor $F_1,F_2$ dan fungsi
\maabb {g_1,g_2} { X} { \R
} {.}
}}
\faktzusatz {}
\faktzusatz {}

}
{

\aufzaehlungzwei {Tetapkan titik

\mavergleichskette
{\vergleichskette
{ P
}
{ \in }{ X
}
{  }{
}
{    }{
}
{   }{
}
}
{}{}{}
dan penampang

\mavergleichskette
{\vergleichskette
{ s
}
{ \in }{ C^1(U,E)
}
{  }{
}
{    }{
}
{   }{
}
}
{}{}{}
pada himpunan terbuka

\mavergleichskette
{\vergleichskette
{ P
}
{ \in }{ U
}
{ \subseteq }{ X
}
{    }{
}
{   }{
}
}
{}{}{,}
maka

\mavergleichskettedisp
{\vergleichskette
{ \nabla_F (s)(P)
}
{ =} { ( \pi_{\text{vert} } \circ  T_P(s) ) (F(P))
}
{ } {
}
{ } {
}
{ } {
}
}
{}{}{.}
} {Pemetaan singgung dan proyeksi vertikal keduanya linear. Menurut (1), ketergantungan pada $F$ juga linear dan selaras dengan perkalian oleh fungsi.
}

}



\inputfaktbeweis
{Differenzierbare Hyperfläche/Zweites Tangentialbündel/Induzierter Zusammenhang/Vertikale Ableitung/Fakt}
{Lema}
{}
{

\faktsituation {Misalkan

\mavergleichskette
{\vergleichskette
{ Y
}
{ \subseteq }{ W
}
{ \subseteq }{ \R^n
}
{    }{
}
{   }{
}
}
{}{}{,}
dengan $W$
\definitionsverweis {terbuka}{}{,}
dan himpunan pertama adalah
\definitionsverweis {hipermuka terdiferensialkan}{}{.}
Misalkan pula
\maabbdisp {F} { Y} { \R^n
} {}
adalah medan vektor tangensial.}
\faktfolgerung {Maka
\definitionsverweis {turunan vertikal}{}{}
$\nabla_F$ yang ditentukan oleh
\definitionsverweis {koneksi}{}{}
pada
Catatan 24.8
berimpit dengan
\definitionsverweis {turunan kovarian}{}{.}}
\faktzusatz {}
\faktzusatz {}

}
{

Untuk penampang diferensiabel
\maabb {s} { Y } { TY
} {}
turunan vertikal koneksi tersebut searah $F$ diberikan oleh komposisi

\mathdisp {Y \stackrel{F}{\longrightarrow} TY \stackrel{T(s)}{\longrightarrow} TTY  \stackrel{  \pi_{\text{vert} } }{\longrightarrow} TY} { . }

Pemetaan $T(s)$ mengirim titik
\mathl{(P,u)}{} ke
\mathl{(P,u, s(P), \left(  D s   \right)_{ P } \left(  u   \right))}{} dalam notasi
Catatan 24.8,
kemudian proyeksi vertikal mengirimnya ke
\mathl{(P, \pi_P  \left(  \left(  D s   \right)_{ P } \left(  u   \right)  \right) )}{.}
Inilah tepatnya rumus dalam definisi turunan kovarian.

}


\chapter{Lembar Kerja 24}
\input{generated/worksheet24.id.build.tex}

\part{Unit 25}
\chapter[Kuliah 25]{Kuliah 25: Penampang Horizontal dan Koneksi Linear}
\input{generated/lecture25.id.build.tex}

\chapter{Lembar Kerja 25}
\setcounter{section}{25}

  \zwischenueberschrift{Soal latihan}

\inputaufgabegibtloesung
{}
{

Misalkan $X$ suatu
\definitionsverweis {manifold diferensiabel}{}{}
dan
\maabb {p} { E  =  X \times \R } { X
} {}
bundel vektor trivial di atas $X$ yang ber-
\definitionsverweis {rank}{}{}
$1$. Misalkan $\omega$ suatu
$1$-\definitionsverweis {bentuk diferensial}{}{}
pada $X$. Tunjukkan pernyataan-pernyataan berikut.
\aufzaehlungfuenf{Pemetaan
\maabbeledisp {} {TE} { p^*E
} {(P,u,v,w) } { (P, u, \omega(P,v) +w)
} {,}
mendefinisikan suatu
\definitionsverweis {koneksi}{}{.}
}{Untuk fungsi-fungsi yang terdiferensialkan secara kontinu
\maabb {f} { U } {\R
} {}
pada suatu himpunan terbuka

\mavergleichskette
{\vergleichskette
{ U
}
{ \subseteq }{ X
}
{  }{
}
{    }{
}
{   }{
}
}
{}{}{}
turunan vertikal dari koneksi tersebut diberikan oleh

\mavergleichskettedisp
{\vergleichskette
{ \nabla (f) (P)
}
{ =} { \omega (P) + (df)_P
}
{ } {
}
{ } {
}
{ } {
}
}
{}{}{}
\zusatzklammer {di kedua ruas kita mengidentifikasi fungsi bernilai riil $f$ dengan penampang
\maabbele {} { X } { X \times \R
} { x } { (x,f(x))
} {.}} {} {}
}{Suatu fungsi yang terdiferensialkan secara kontinu
\maabb {f} { U } {\R
} {}
merupakan
\definitionsverweis {penampang horizontal}{}{}
di atas $U$ terhadap $\nabla$ jika dan hanya jika $-f$ merupakan
\definitionsverweis {bentuk primitif}{}{}
dari $\omega$ pada $U$.
}{Bentuk $\omega$
\definitionsverweis {tertutup}{}{}
jika dan hanya jika koneksi tersebut
\definitionsverweis {terintegralkan secara lokal}{}{.}
}{Bentuk $\omega$
\definitionsverweis {eksak}{}{}
jika dan hanya jika koneksi tersebut
\definitionsverweis {terintegralkan secara global}{}{.}
}

}
{} {}

\inputaufgabe
{}
{

Kita tinjau bundel vektor trivial
\maabbdisp {p} {E= \R^2 \times \R} { \R^2
} {.}
Deskripsikan suatu
\definitionsverweis {koneksi}{}{}
pada $E$ sedemikian sehingga tidak terdapat
\definitionsverweis {penampang horizontal}{}{.}

}
{} {}

\inputaufgabe
{}
{

Misalkan $X$ suatu
\definitionsverweis {manifold diferensiabel}{}{}
dan
\maabb {p} { E  =  X \times \R } { X
} {}
bundel vektor trivial di atas $X$ yang ber-
\definitionsverweis {rank}{}{}
$1$. Misalkan $\omega$ suatu
$1$-\definitionsverweis {bentuk diferensial}{}{}
pada $X$. Tunjukkan pernyataan-pernyataan berikut.
\aufzaehlungdrei{Pemetaan
\maabbeledisp {} {TE} { p^*E
} {(P,u,v,w) } { (P, u, u\omega(P,v) +w)
} {,}
mendefinisikan suatu
\definitionsverweis {koneksi}{}{.}
}{Untuk fungsi-fungsi yang terdiferensialkan secara kontinu
\maabb {f} { U } {\R
} {}
pada suatu himpunan terbuka

\mavergleichskette
{\vergleichskette
{ U
}
{ \subseteq }{ X
}
{  }{
}
{    }{
}
{   }{
}
}
{}{}{}
turunan vertikal dari koneksi tersebut diberikan oleh

\mavergleichskettedisp
{\vergleichskette
{ \nabla (f) (P)
}
{ =} { f\omega (P) + (df)_P
}
{ } {
}
{ } {
}
{ } {
}
}
{}{}{}
\zusatzklammer {di kedua ruas kita mengidentifikasi fungsi bernilai riil $f$ dengan penampang
\maabbele {} { X } { X \times \R
} { x } { (x,f(x))
} {.}} {} {}
}{Suatu fungsi yang terdiferensialkan secara kontinu dan tidak pernah nol
\maabb {f} { U } {\R
} {}
merupakan
\definitionsverweis {penampang horizontal}{}{}
di atas $U$ terhadap $\nabla$ jika dan hanya jika $-   \ln  \betrag {  f  }$ merupakan
\definitionsverweis {bentuk primitif}{}{}
dari $\omega$ pada $U$.
}

}
{} {}

\inputaufgabe
{}
{

Misalkan $X$ suatu
\definitionsverweis {manifold diferensiabel}{}{}
dan $E$ suatu
\definitionsverweis {bundel vektor diferensiabel}{}{}
di atas $X$ yang dilengkapi dengan suatu
\definitionsverweis {koneksi}{}{.}
Tunjukkan bahwa suatu penampang yang terdiferensialkan secara kontinu
\maabb {s} { U } { E\vert_U
} {}
pada suatu himpunan terbuka

\mavergleichskette
{\vergleichskette
{ U
}
{ \subseteq }{ X
}
{  }{
}
{    }{
}
{   }{
}
}
{}{}{}
bersifat
\definitionsverweis {horizontal}{}{}
jika dan hanya jika
\definitionsverweis {turunan vertikal}{}{}-nya memenuhi

\mavergleichskette
{\vergleichskette
{ \nabla s
}
{ = }{ 0
}
{  }{
}
{    }{
}
{   }{
}
}
{}{}{.}

}
{} {}

\inputaufgabe
{}
{

Misalkan $X$ suatu
\definitionsverweis {manifold diferensiabel}{}{}
dan $E$ suatu
\definitionsverweis {bundel vektor diferensiabel}{}{}
di atas $X$ yang dilengkapi dengan suatu
\definitionsverweis {koneksi}{}{.}
Misalkan $Z$ suatu manifold diferensiabel lain dengan pemetaan diferensiabel
\maabb {\psi} { Z } { X
} {.}
Tunjukkan bahwa suatu penampang yang terdiferensialkan secara kontinu
\maabb {s} { Z } { E
} {}
di atas $\psi$ bersifat
\definitionsverweis {horizontal}{}{}
jika dan hanya jika
\definitionsverweis {turunan vertikal}{}{}-nya memenuhi

\mavergleichskette
{\vergleichskette
{ \nabla s
}
{ = }{ 0
}
{  }{
}
{    }{
}
{   }{
}
}
{}{}{.}

}
{} {}

Untuk soal berikut, gunakan
Soal 25.21.

\inputaufgabe
{}
{

Misalkan
\maabb {p} { E } { X
} {}
suatu
\definitionsverweis {bundel vektor}{}{}
ber-
\definitionsverweis {rank}{}{}
$r$ di atas suatu
\definitionsverweis {manifold diferensiabel}{}{}
$X$, yang dilengkapi dengan suatu
\definitionsverweis {koneksi yang terintegralkan secara global}{}{.}
Tunjukkan bahwa terdapat
\zusatzklammer {yang terdefinisi pada $X$} {} {}
\definitionsverweis {penampang-penampang horizontal}{}{}

\mathl{s_1  , \ldots , s_r}{} di $E$ sedemikian sehingga
\maabbeledisp {} {X \times \R^r} { E
} {(P,v_1  , \ldots , v_r)} { \sum_{i = 1}^r v_i s_i(P)
} {,}
merupakan suatu
\definitionsverweis {isomorfisme bundel vektor}{}{.}

}
{} {}

\inputaufgabegibtloesung
{}
{

Misalkan
\maabb {p} { E } { X
} {}
suatu
\definitionsverweis {bundel vektor}{}{}
ber-
\definitionsverweis {rank}{}{}
$r$ di atas suatu
\definitionsverweis {manifold diferensiabel}{}{}
$X$, yang dilengkapi dengan suatu koneksi linear yang
\definitionsverweis {terintegralkan secara global}{}{.}
Tunjukkan bahwa terdapat
\zusatzklammer {yang terdefinisi pada $X$} {} {}
\definitionsverweis {penampang-penampang horizontal}{}{}

\mathl{s_1 , \ldots , s_r}{} di $E$ sedemikian sehingga, di bawah
\definitionsverweis {isomorfisme bundel vektor}{}{}
\maabbeledisp {} {X \times \R^r} { E
} {(P,v_1 , \ldots , v_r)} { \sum_{i = 1}^r v_i s_i(P)
} {,}
\definitionsverweis {simbol-simbol Christoffel}{}{}
dari koneksi tersebut terhadap penampang basis
\mathl{s_1  , \ldots , s_r}{} semuanya bernilai nol.

}
{} {}

\inputaufgabegibtloesung
{}
{

Misalkan diberikan suatu
\definitionsverweis {medan vektor}{}{}
kontinu yang bergantung pada waktu,
\maabbeledisp {F} {I \times \R^n } { \R^n
} { \left(  t   , \,  u_1  , \,   \ldots , \,  u_n                \right) } { \left(  F_1 \left(  u_1 , \ldots ,  u_n  \right)    , \,   \ldots  , \,   F_n \left(  u_1 , \ldots , u_n  \right)                 \right)
} {,}
beserta
\definitionsverweis {sistem persamaan diferensial biasa}{}{}
yang terkait,

\mavergleichskettedisp
{\vergleichskette
{ u'
}
{ =} { F(t,u)
}
{ } {
}
{ } {
}
{ } {
}
}
{}{}{}
pada suatu
\definitionsverweis {interval terbuka}{}{}

\mavergleichskette
{\vergleichskette
{ I
}
{ \subseteq }{ \R
}
{  }{
}
{    }{
}
{   }{
}
}
{}{}{.}
\aufzaehlungdrei{Tunjukkan bahwa pemetaan
\maabbeledisp {} {I \times \R^n \times \R \times \R^n } { I \times \R^n  \times \R^n
} {(t,u,v,w)} { \left(  t   , \,  u  , \,   - vF \left(  t, u_1  , \ldots , u_n  \right)  +w                 \right)
} {,}
memberikan
\zusatzklammer {proyeksi vertikal dari suatu} {} {}
\definitionsverweis {koneksi}{}{}
pada bundel vektor trivial
\maabbdisp {} {I \times \R^n } { \R^n
} {}
diberikan.
}{Tentukan
\definitionsverweis {turunan vertikal}{}{}
$\nabla_\partial u$ dari suatu
\definitionsverweis {kurva diferensiabel}{}{}
\maabbdisp {u} { I } { \R^n
} {}
\zusatzklammer {yang dipandang sebagai penampang di $I \times \R^n$} {} {.}
}{Tunjukkan bahwa suatu kurva diferensiabel
\maabb {u} { I } { \R^n
} {}
merupakan solusi dari sistem persamaan diferensial jika dan hanya jika $u$ merupakan
\definitionsverweis {penampang horizontal}{}{}
terhadap koneksi tersebut.
}

}
{} {}

\inputaufgabe
{}
{

Misalkan $\nabla$ suatu
\definitionsverweis {koneksi linear}{}{}
pada suatu
\definitionsverweis {bundel vektor diferensiabel}{}{}
$E$ di atas suatu
\definitionsverweis {manifold diferensiabel}{}{}
$X$. Tunjukkan bahwa penampang nol bersifat
\definitionsverweis {horizontal}{}{.}

}
{} {}

\inputaufgabe
{}
{

Misalkan $\nabla$
\definitionsverweis {koneksi trivial}{}{}
pada bundel vektor trivial
\maabb {p} {X \times W } { X
} {}
yang terkait dengan suatu ruang vektor riil
\mathl{W}{} di atas suatu
\definitionsverweis {manifold diferensiabel}{}{}
$X$. Tunjukkan bahwa $\nabla$
\definitionsverweis {linear}{}{.}

}
{} {}

\inputaufgabegibtloesung
{}
{

Misalkan

\mavergleichskette
{\vergleichskette
{ U
}
{ \subseteq }{ \R^n
}
{  }{
}
{    }{
}
{   }{
}
}
{}{}{}
suatu
\definitionsverweis {himpunan bagian terbuka}{}{}
dan misalkan $\nabla$
\definitionsverweis {koneksi trivial}{}{}
pada
\mathl{U \times \R}{.} Tunjukkan bahwa
\definitionsverweis {simbol-simbol Christoffel}{}{}
dan
\definitionsverweis {turunan vertikal}{}{}
dari koneksi tersebut memenuhi sifat-sifat berikut.
\aufzaehlungvier{Simbol-simbol Christoffel memenuhi

\mavergleichskettedisp
{\vergleichskette
{ \Gamma_{ij}^k
}
{ =} { 0
}
{ } {
}
{ } {
}
{ } {
}
}
{}{}{}
untuk semua $i,j,k$.
}{Berlaku

\mavergleichskettedisp
{\vergleichskette
{ \nabla_{\partial_i }  \partial_j
}
{ =} { 0
}
{ } {
}
{ } {
}
{ } {
}
}
{}{}{}
untuk
\definitionsverweis {medan-medan vektor standar}{}{}
$\partial_i$.
}{Berlaku

\mavergleichskettedisp
{\vergleichskette
{ \nabla_{\partial_i }  \left(  \sum_{j = 1}^n f_j \partial_j  \right)
}
{ =} { \sum_{j = 1}^n \partial_i(f_j) \partial_j
}
{ } {
}
{ } {
}
{ } {
}
}
{}{}{}
untuk fungsi-fungsi diferensiabel $f_j$.
}{Untuk suatu medan vektor $V$ dan suatu medan vektor diferensiabel $W$ pada $U$ berlaku

\mavergleichskettedisp
{\vergleichskette
{ ( \nabla_V W)(P)
}
{ =} { \left(  D W   \right)_{ P } \left(   V(P)  \right)
}
{ } {
}
{ } {
}
{ } {
}
}
{}{}{.}
}

}
{} {}

\inputaufgabegibtloesung
{}
{

Pada
\definitionsverweis {bundel vektor trivial}{}{}

\mathl{\R \times \R^2}{} di atas $\R$, kita tinjau
\definitionsverweis {koneksi trivial}{}{}
$\nabla$.
\aufzaehlungdrei{Tunjukkan bahwa

\mavergleichskette
{\vergleichskette
{ f_1(t)
}
{ = }{ \left(  1+ t^2  , \,  t^3                   \right)
}
{  }{
}
{    }{
}
{   }{
}
}
{}{}{}
dan

\mavergleichskette
{\vergleichskette
{ f_2(t)
}
{ = }{ \left(  t  , \,   1+t^2                   \right)
}
{  }{
}
{    }{
}
{   }{
}
}
{}{}{}
merupakan
\definitionsverweis {penampang-penampang basis}{}{}
dari bundel vektor tersebut.
}{Nyatakan penampang basis standar $e_1,e_2$ sebagai kombinasi linear dari penampang basis $f_1,f_2$.
}{Tentukan
\definitionsverweis {simbol-simbol Christoffel}{}{}
dari $\nabla$ terhadap penampang basis tersebut.
}

}
{} {}

\inputaufgabe
{}
{

Pada
\definitionsverweis {bundel vektor trivial}{}{}

\mathl{\R^2 \times \R^2}{} di atas $\R^2$, kita tinjau
\definitionsverweis {koneksi trivial}{}{}
$\nabla$.
\aufzaehlungzwei {Tunjukkan bahwa

\mavergleichskette
{\vergleichskette
{ f_1(x,y)
}
{ = }{ \left(  x   , \,   y                   \right)
}
{  }{
}
{    }{
}
{   }{
}
}
{}{}{}
dan

\mavergleichskette
{\vergleichskette
{ f_2(x,y)
}
{ = }{ \left(  - y  , \,   x                   \right)
}
{  }{
}
{    }{
}
{   }{
}
}
{}{}{}
merupakan
\definitionsverweis {penampang-penampang basis}{}{}
dari bundel vektor pada

\mavergleichskette
{\vergleichskette
{ U
}
{ = }{ \R^2 \setminus \{ (0,0)\}
}
{  }{
}
{    }{
}
{   }{
}
}
{}{}{.}
} {Tentukan
\definitionsverweis {simbol-simbol Christoffel}{}{}
dari $\nabla$ terhadap penampang basis tersebut pada $U$.
}

}
{} {}

\inputaufgabegibtloesung
{}
{

Misalkan diberikan suatu
\definitionsverweis {sistem persamaan diferensial linear homogen}{}{}

\mavergleichskettedisp
{\vergleichskette
{ u'
}
{ =} { Mu
}
{ } {
}
{ } {
}
{ } {
}
}
{}{}{}
pada suatu
\definitionsverweis {interval terbuka}{}{}

\mavergleichskette
{\vergleichskette
{ I
}
{ \subseteq }{ \R
}
{  }{
}
{    }{
}
{   }{
}
}
{}{}{}
dengan

\mavergleichskettedisp
{\vergleichskette
{ M(t)
}
{ =} { \left(  a_{kj} (t)  \right)_{kj}
}
{ } {
}
{ } {
}
{ } {
}
}
{}{}{}
\zusatzklammer {di sini $k$ adalah indeks baris dan $j$ adalah indeks kolom} {} {}
serta fungsi-fungsi kontinu
\maabbdisp {a_{kj}} { I } {\R
} {}
diberikan.
Melalui persamaan

\mavergleichskettedisp
{\vergleichskette
{ \Gamma_{j}^k (t)
}
{ \defeq} { -a_{kj}
}
{ } {
}
{ } {
}
{ } {
}
}
{}{}{}
kita memandang fungsi-fungsi ini sebagai
\definitionsverweis {simbol-simbol Christoffel}{}{}
\zusatzklammer {dengan

\mavergleichskettek
{\vergleichskettek
{ i
}
{ = }{ 1
}
{  }{
}
{    }{
}
{   }{
}
}
{}{}{}} {} {}
dari
\definitionsverweis {koneksi linear}{}{}
$\nabla$ pada bundel vektor trivial
\maabb {p} {I \times \R^n } { I
} {}
dalam pengertian
Catatan 25.8.
Tunjukkan bahwa suatu
\definitionsverweis {kurva diferensiabel}{}{}
\maabbdisp {u} { I } { \R^n
} {}
merupakan solusi sistem persamaan diferensial jika dan hanya jika $u$ merupakan
\definitionsverweis {penampang horizontal}{}{}
terhadap koneksi tersebut.

}
{} {}

\inputaufgabe
{}
{

Misalkan

\mavergleichskette
{\vergleichskette
{ Y
}
{ \subseteq }{ W
}
{ \subseteq }{ \R^n
}
{    }{
}
{   }{
}
}
{}{}{,}
$W$
\definitionsverweis {terbuka}{}{,}
suatu
\definitionsverweis {hipermuka terdiferensialkan}{}{,}
dan
\maabb {\gamma} { I } { Y
} {}
suatu
\definitionsverweis {kurva diferensiabel}{}{.}
Misalkan
\maabbdisp {F} { I } { TY
} {}
suatu medan vektor diferensiabel sepanjang $\gamma$; dengan kata lain, terdapat diagram komutatif

\mathdisp {\begin{matrix}  &   &   & TY  \\ & & F \nearrow & \downarrow  \\  &  I  & \stackrel{ \gamma }{\longrightarrow}  &  Y \! \\ \end{matrix}} {  }

Tunjukkan bahwa $F$
\definitionsverweis {paralel sepanjang}{}{}
$\gamma$ jika dan hanya jika $F$ merupakan
\definitionsverweis {penampang horizontal}{}{}
terhadap
\zusatzklammer {yang didefinisikan menggunakan proyeksi ortogonal} {} {}
\definitionsverweis {koneksi}{}{}
pada $TY$.

}
{} {}

\inputaufgabe
{}
{

Misalkan

\mavergleichskette
{\vergleichskette
{ Y
}
{ \subseteq }{ W
}
{ \subseteq }{ \R^n
}
{    }{
}
{   }{
}
}
{}{}{,}
$W$
\definitionsverweis {terbuka}{}{,}
suatu
\definitionsverweis {hipermuka terdiferensialkan}{}{,}
dan misalkan $\nabla$
\zusatzklammer {yang didefinisikan menggunakan proyeksi ortogonal} {} {}
\definitionsverweis {koneksi}{}{}
pada $TY$. Tunjukkan bahwa $\nabla$
\definitionsverweis {linear}{}{.}

}
{} {}

\inputaufgabe
{}
{

Misalkan
\mathkor {} {\nabla_1} {dan} {\nabla_2} {}
\definitionsverweis {koneksi-koneksi linear}{}{}
pada suatu
\definitionsverweis {bundel vektor diferensiabel}{}{}
$E$ di atas suatu
\definitionsverweis {manifold diferensiabel}{}{}
$X$. Misalkan
\maabbdisp {g_1,g_2} { X } { \R
} {}
fungsi-fungsi yang
\definitionsverweis {terdiferensialkan secara kontinu}{}{}
dengan

\mavergleichskette
{\vergleichskette
{ g_1+g_2
}
{ = }{ 1
}
{  }{
}
{    }{
}
{   }{
}
}
{}{}{.}
Tunjukkan bahwa
\mathl{g_1 \nabla_1+ g_2 \nabla_2}{} juga merupakan suatu koneksi linear
\zusatzklammer {dalam pengertian yang telah didefinisikan} {?} {}
pada $E$.

}
{} {}

\inputaufgabe
{}
{

Misalkan

\mavergleichskette
{\vergleichskette
{ X
}
{ = }{ \bigcup_{i \in I} U_i
}
{  }{
}
{    }{
}
{   }{
}
}
{}{}{}
suatu
\definitionsverweis {tutupan terbuka}{}{}
dari suatu
\definitionsverweis {manifold diferensiabel}{}{}
$X$ yang mentrivialkan
\definitionsverweis {bundel vektor diferensiabel}{}{}
$E$. Misalkan $\nabla_i$
\definitionsverweis {koneksi trivial}{}{}
pada $E\vert_{U_i }$ dan misalkan
\mathkor {} {h_j} {} {j \in J} {,}
suatu
\definitionsverweis {partisi satuan}{}{}
yang subordinat terhadap tutupan tersebut,

\mathbed {h_j} {}
 {j \in J} {}
 {} {} {} {.}
Tunjukkan bahwa $\sum_{j \in J} h_j \nabla_{i(j)}$ merupakan suatu koneksi linear yang terdefinisi dengan baik pada $E$.

}
{} {}

\inputaufgabe
{}
{

Misalkan $X$ suatu
\definitionsverweis {manifold diferensiabel}{}{}
dan $\omega$ suatu
$1$-\definitionsverweis {bentuk diferensial}{}{}
kontinu pada $X$. Tunjukkan bahwa pemetaan
\maabbeledisp {} {C^1(X,\R)} { C^0(X,T^*X)
} { f } { \nabla_\omega (f) \defeq f \omega + df
} {,}
memenuhi aturan Leibniz dari
Teorema 25.5.

}
{} {}

\inputaufgabe
{}
{

Misalkan
\maabb {p} { E } { X
} {}
suatu
\definitionsverweis {bundel vektor}{}{}
diferensiabel di atas suatu
\definitionsverweis {manifold diferensiabel}{}{}
$X$. Misalkan $E$ sendiri merupakan suatu
\definitionsverweis {manifold Riemann}{}{.}
Tunjukkan bahwa suatu
\definitionsverweis {koneksi linear}{}{}
pada $E$ diperoleh dengan mengambil, untuk
\definitionsverweis {bundel vertikal}{}{}

\mavergleichskette
{\vergleichskette
{ V
}
{ \subseteq }{ T E
}
{  }{
}
{    }{
}
{   }{
}
}
{}{}{}
\definitionsverweis {komplemen ortogonal}{}{}
sebagai suatu
\definitionsverweis {bundel horizontal}{}{.}

}
{} {}

\inputaufgabe
{}
{

Misalkan
\maabb {p} { E } { X
} {}
suatu
\definitionsverweis {bundel vektor}{}{}
di atas
\definitionsverweis {manifold diferensiabel}{}{}
$X$ dengan
\definitionsverweis {basis topologi terhitung}{}{.}
Gunakan
Soal 22.15
dan
Soal 25.18
untuk menunjukkan bahwa terdapat suatu
\definitionsverweis {koneksi linear}{}{}
pada $E$.

}
{} {}

  \zwischenueberschrift{Soal yang dikumpulkan}

\inputaufgabe
{5 (3+2)}
{

Misalkan

\mavergleichskette
{\vergleichskette
{ X
}
{ = }{ S^1
}
{  }{
}
{    }{
}
{   }{
}
}
{}{}{}
\definitionsverweis {lingkaran satuan}{}{}
dengan
\definitionsverweis {bundel tangen}{}{}
\maabb {p} {TS^1 = S^1 \times \R} { S^1
} {.}
\aufzaehlungzwei {Deskripsikan suatu
\definitionsverweis {koneksi}{}{}
pada $TS^1$ sedemikian sehingga tidak terdapat
\definitionsverweis {penampang horizontal}{}{}
pada seluruh $S^1$.
} {Tentukan suatu penampang horizontal dari koneksi pada (1) sepanjang pemetaan
\maabbeledisp {} {\R} { S^1
} { t } { \left(  \cos  t   , \,   \sin  t                   \right)
} {.}
}

}
{} {}

\inputaufgabe
{4}
{

Misalkan $X$ suatu
\definitionsverweis {manifold diferensiabel}{}{}
dan $E$ suatu
\definitionsverweis {bundel vektor diferensiabel}{}{}
di atas $X$ yang dilengkapi dengan suatu
\definitionsverweis {koneksi}{}{.}
Misalkan $Z$ suatu manifold diferensiabel lain dengan pemetaan diferensiabel
\maabb {\psi} { Z } { X
} {.}
Misalkan $Z$
\definitionsverweis {terhubung}{}{}
dan misalkan
\maabb {s_1,s_2} { Z } { E
} {}
\definitionsverweis {penampang-penampang horizontal}{}{}
di atas $\psi$ dengan

\mavergleichskette
{\vergleichskette
{ s_1(z)
}
{ = }{ s_2(z)
}
{  }{
}
{    }{
}
{   }{
}
}
{}{}{}
untuk suatu titik

\mavergleichskette
{\vergleichskette
{ z
}
{ \in }{ Z
}
{  }{
}
{    }{
}
{   }{
}
}
{}{}{.}
Tunjukkan bahwa

\mavergleichskette
{\vergleichskette
{ s_1
}
{ = }{ s_2
}
{  }{
}
{    }{
}
{   }{
}
}
{}{}{.}

}
{} {}

\inputaufgabe
{3}
{

Misalkan $X$ suatu
\definitionsverweis {manifold diferensiabel}{}{}
dan $E$ suatu
\definitionsverweis {bundel vektor diferensiabel}{}{}
di atas $X$ yang dilengkapi dengan suatu
\definitionsverweis {koneksi linear}{}{,}
dan misalkan $G$ suatu
\definitionsverweis {medan vektor}{}{}
pada $X$. Tunjukkan bahwa untuk setiap
\definitionsverweis {penampang diferensiabel}{}{}
$s$ di $E$ dan setiap
\definitionsverweis {fungsi diferensiabel}{}{}
$f$
\zusatzklammer {yang keduanya terdefinisi pada suatu himpunan terbuka

\mavergleichskette
{\vergleichskette
{ U
}
{ \subseteq }{ X
}
{  }{
}
{    }{
}
{   }{
}
}
{}{}{}} {} {}
berlaku hubungan

\mavergleichskettedisp
{\vergleichskette
{ \nabla_G (fs)
}
{ =} { s \otimes D_G (f) + f \nabla_G (s)
}
{ } {
}
{ } {
}
{ } {
}
}
{}{}{.}

}
{} {}

\inputaufgabe
{3}
{

Berikan suatu contoh
\definitionsverweis {koneksi linear}{}{}
pada suatu
\definitionsverweis {bundel vektor}{}{}
\maabb {p} { E } { X
} {}
di atas suatu
\definitionsverweis {manifold diferensiabel}{}{}
$X$ sedemikian sehingga ruang vektor
\mathl{H(\nabla,X)}{} dari
\definitionsverweis {penampang-penampang horizontal}{}{}
memiliki
\definitionsverweis {dimensi ruang vektor}{}{}
yang lebih besar daripada
\definitionsverweis {rank}{}{}
$E$.

}
{} {}


\section*{Solusi yang disediakan oleh sumber}
\addcontentsline{toc}{section}{Solusi yang disediakan oleh sumber}
\subsection*{Solusi Soal 25.1}
\aufzaehlungfuenf{Pemetaan
\maabbeledisp {\pi_{\text{vert} }} {TE= TX \times \R \times \R } { p^*E  = X \times \R \times \R
} {(P,u,v,w)} { (P,u,  \omega(P,v) + w)
} {}
bersifat linear dalam
\mathkor {} {v} {dan} {w} {,}
untuk $(P,u)$ yang tetap, sehingga merupakan suatu homomorfisme bundel vektor. Ketergantungannya pada basis $E$ terdiferensialkan secara kontinu karena $\omega$ terdiferensialkan secara kontinu. Pemetaan ini merupakan suatu pemisahan bagi penyertaan alami
\maabb {} {p^*E} { TE
} {}
karena
\mathl{(P,u,w)}{} dipetakan ke
\mathl{(P,u,0,w)}{} dan kemudian ke

\mavergleichskette
{\vergleichskette
{ (P,u, \omega(P,0) + w)
}
{ = }{ (P,u,  w)
}
{  }{
}
{    }{
}
{   }{
}
}
{}{}{.}
}{Untuk
\maabb {f} { U } {\R
} {}
kita perlu meninjau komposisi

\mathdisp {TX \stackrel{T(f)}{  \longrightarrow } T E = TX \times \R \times \R \stackrel{ \pi_{\text{vert} } }{ \longrightarrow } X \times \R \longrightarrow \R} {  }

yang menghasilkan

\mathdisp {(P,v) \longmapsto  (P,f(P), v, (df)_P(v) ) \longmapsto (P, f(P), \omega(P,v) + (df)_P(v) ) \longmapsto  \omega(P,v) + (df)_P(v)} { . }

}{Suatu penampang horizontal ada jika dan hanya jika turunan vertikalnya lenyap; dalam hal ini, syarat tersebut berarti

\mavergleichskettedisp
{\vergleichskette
{ \omega(P,v) + (df)_P(v)
}
{ =} { 0
}
{ } {
}
{ } {
}
{ } {
}
}
{}{}{}
untuk semua

\mavergleichskette
{\vergleichskette
{ (P,v)
}
{ \in }{ TX
}
{  }{
}
{    }{
}
{   }{
}
}
{}{}{.}
}{Hal ini mengikuti dari kenyataan bahwa suatu bentuk eksak jika dan hanya jika secara lokal bentuk tersebut mempunyai suatu bentuk primitif.
}{Hal ini mengikuti dari hasil-hasil yang telah dibuktikan.
}

\subsection*{Solusi Soal 25.7}
Menurut
Soal,
terdapat suatu isomorfisme bundel vektor. Menurut
Lema,
untuk suatu penampang horizontal $s$ berlaku

\mavergleichskette
{\vergleichskette
{ \nabla s
}
{ = }{ 0
}
{  }{
}
{    }{
}
{   }{
}
}
{}{}{.}
Oleh karena itu, pada suatu himpunan terbuka

\mavergleichskette
{\vergleichskette
{ U
}
{ \cong }{ V
}
{ \subseteq }{  \R^n
}
{    }{
}
{   }{
}
}
{}{}{}
secara khusus berlaku

\mavergleichskettedisp
{\vergleichskette
{ \nabla_{\partial_i} (s)
}
{ =} { 0
}
{ } {
}
{ } {
}
{ } {
}
}
{}{}{}
untuk semua $i$. Ini berarti bahwa semua simbol Christoffel dari
\definitionsverweis {penampang basis}{}{}
$s_j$ terhadap medan vektor standar $\partial_i$ bernilai nol.

\subsection*{Solusi Soal 25.8}
\input{generated/worksheet25_exercise08_solution.id.build.tex}
\subsection*{Solusi Soal 25.11}
\aufzaehlungvier{Hal ini mengikuti dari deskripsi eksplisit turunan vertikal suatu koneksi linear; lihat
Catatan.
}{Jelas dari (1).
}{Hal ini mengikuti dari (2) dan
Lema.
}{Kita tuliskan

\mavergleichskette
{\vergleichskette
{ V
}
{ = }{ \sum_{i = 1}^n {h_i } \partial_i
}
{  }{
}
{    }{
}
{   }{
}
}
{}{}{}
dan

\mavergleichskette
{\vergleichskette
{ W
}
{ = }{ \sum_{j = 1}^n {f_j } \partial_j
}
{  }{
}
{    }{
}
{   }{
}
}
{}{}{.}
Menurut
Lema
dan bagian-bagian yang telah dibuktikan,

\mavergleichskettealign
{\vergleichskettealign
{ \nabla_V W
}
{ =} { \nabla_{\sum_{i = 1}^n {h_i } \partial_i }   \left(  \sum_{j = 1}^n {f_j } \partial_j   \right)
}
{ =} { \sum_{j = 1}^n \left(  \sum_{i = 1}^n h_i  \partial_i f_j   \right) \partial_j
}
{ } {
}
{ } {
}
}
{}
{}{.}
Pada suatu titik

\mavergleichskette
{\vergleichskette
{ P
}
{ \in }{ U
}
{  }{
}
{    }{
}
{   }{
}
}
{}{}{}
nilai ini adalah

\mathdisp {\sum_{j  = 1}^n \left(  \sum_{i = 1}^n h_i (P) ( \partial_i f_j)(P)   \right) \partial_j} { . }

Di sisi lain,

\mavergleichskettealign
{\vergleichskettealign
{ \left(  D W   \right)_{ P } \left(   V(P)  \right)
}
{ =} { \left(  (\partial_j f_i)(P)  \right)_{ji} \begin{pmatrix} h_1(P) \\ \vdots\\  h_n(P)   \end{pmatrix}
}
{ =} { \begin{pmatrix}  \sum_{i = 1}^n h_i (P) \left(  \partial_i f_1  \right)(P)  \\ \vdots\\  \sum_{i = 1}^n h_i (P)  \left(  \partial_i f_n  \right) (P)   \end{pmatrix}
}
{ } {
}
{ } {
}
}
{}
{}{.}
}

\subsection*{Solusi Soal 25.12}
\input{generated/worksheet25_exercise12_solution.id.build.tex}
\subsection*{Solusi Soal 25.14}
Pemetaan
\maabbdisp {u} { I } {\R^n
} {}
dinyatakan terhadap penampang basis
\mathl{e_1 , \ldots , e_n}{}
oleh
\mavergleichskettedisp
{\vergleichskette
{ u
}
{ =} { u_1e_1 + \cdots + u_ne_n
}
{ } {
}
{ } {
}
{ } {
}
}
{}{}{}
.
Menurut
\definitionsverweis {turunan vertikal}{}{}
dari koneksi tersebut, turunan penampang ini dalam satu-satunya arah turunan $\partial$, berdasarkan
Lema,
adalah

\mavergleichskettedisp
{\vergleichskette
{  \sum_{k = 1}^n \left(    \partial u_k  +  \sum_{j = 1}^n u_j \Gamma_{j}^k  \right) e_k
}
{ =} { \sum_{k = 1}^n \left(   u_k'  +  \sum_{j = 1}^n u_j \Gamma_{j}^k  \right) e_k
}
{ } {
}
{ } {
}
{ } {
}
}
{}{}{.}
Menurut
Lema,
$u$ merupakan penampang horizontal jika dan hanya jika turunan vertikalnya sama dengan $0$. Untuk setiap komponen, ini berarti

\mavergleichskettedisp
{\vergleichskette
{ u_k' + \sum_{j = 1}^n \Gamma_j^k u_j
}
{ =} { 0
}
{ } {
}
{ } {
}
{ } {
}
}
{}{}{}
untuk semua $k$. Kondisi ini ekuivalen dengan

\mavergleichskettedisp
{\vergleichskette
{ u_k'
}
{ =} { \sum_{j = 1}^n a_{kj} u_j
}
{ } {
}
{ } {
}
{ } {
}
}
{}{}{}
untuk semua $k$, yang tepat berarti bahwa $u$ memenuhi sistem matriks

\mavergleichskettedisp
{\vergleichskette
{ u'
}
{ =} { \left(  a_{kj}    \right) u
}
{ } {
}
{ } {
}
{ } {
}
}
{}{}{}
;
dengan demikian, pemetaan itu merupakan suatu solusi sistem persamaan diferensial linear tersebut.


\part{Unit 26}
\chapter[Kuliah 26]{Kuliah 26: Koneksi Levi--Civita}
\setcounter{section}{26}

  \zwischenueberschrift{Koneksi Levi--Civita pada himpunan terbuka}

Misalkan

\mavergleichskette
{\vergleichskette
{ U
}
{ \subseteq }{ \R^n
}
{  }{
}
{    }{
}
{   }{
}
}
{}{}{}
terbuka dan padanya diberikan suatu struktur Riemann melalui fungsi-fungsi $g_{ij}$. Kita mempunyai fungsi koordinat $x_i$ dan medan vektor konstan terkait $\partial_i$. Suatu koneksi linear pada bundel tangen
\mathl{U \times \R^n}{} dijelaskan oleh turunan

\mavergleichskettedisp
{\vergleichskette
{ \nabla_{\partial_i } \partial_j
}
{ =} { \sum_{k = 1}^n\Gamma_{ij}^k\,\partial_k
}
{ } {
}
{ } {
}
{ } {
}
}
{}{}{}
Berikut ini akan terlihat bahwa pilihan berikut menghasilkan suatu koneksi yang berkaitan erat dengan struktur Riemann.

\inputdefinition
{}
{

Misalkan

\mavergleichskette
{\vergleichskette
{ U
}
{ \subseteq }{ \R^n
}
{  }{
}
{    }{
}
{   }{
}
}
{}{}{}
terbuka dan padanya diberikan suatu struktur Riemann melalui fungsi-fungsi $\left(  g_{ij}  \right)$ dengan matriks invers
\mathl{\left(  g^{k l}  \right)}{.} Fungsi-fungsi bernilai riil yang didefinisikan pada $U$,

\mavergleichskettedisp
{\vergleichskette
{ \Gamma_{ij}^k
}
{ =} { { \frac{   1  }{  2 } } \sum_{l = 1}^n g^{kl}(\partial_ig_{jl}+\partial_jg_{il}-\partial_lg_{ij})
}
{ } {
}
{ } {
}
{ } {
}
}
{}{}{}
disebut
\definitionswort {simbol Christoffel}{}
untuk koneksi Levi--Civita.

}

\inputdefinition
{}
{

Misalkan

\mavergleichskette
{\vergleichskette
{ U
}
{ \subseteq }{ \R^n
}
{  }{
}
{    }{
}
{   }{
}
}
{}{}{}
terbuka dan padanya diberikan suatu struktur Riemann melalui fungsi-fungsi $\left(  g_{ij}  \right)$ dengan matriks invers
\mathl{\left(  g^{k l}  \right)}{.} Pada $U$,
\definitionsverweis {simbol Christoffel}{}{}

\mavergleichskettedisp
{\vergleichskette
{ \Gamma_{ij}^k
}
{ =} { { \frac{   1  }{  2 } } \sum_{l = 1}^n g^{kl}(\partial_ig_{jl}+\partial_jg_{il}-\partial_lg_{ij})
}
{ } {
}
{ } {
}
{ } {
}
}
{}{}{}
menentukan suatu
\definitionsverweis {koneksi linear}{}{}
pada

\mavergleichskette
{\vergleichskette
{ TU
}
{ = }{ U \times \R^n
}
{  }{
}
{    }{
}
{   }{
}
}
{}{}{}
yang disebut
\definitionswort {koneksi Levi--Civita}{}
pada $U$.

}



\inputfaktbeweis
{R^n/Offene Menge/Levi-Civita-Zusammenhang/Trivial/Fakt}
{Lemma}
{}
{

\faktsituation {Misalkan

\mavergleichskette
{\vergleichskette
{ U
}
{ \subseteq }{ \R^n
}
{  }{
}
{    }{
}
{   }{
}
}
{}{}{}
suatu
\definitionsverweis {himpunan bagian terbuka}{}{,}
yang dilengkapi dengan
\definitionsverweis {struktur Riemann}{}{}
terinduksi dari $\R^n$.}
\faktfolgerung {Maka
\definitionsverweis {koneksi Levi--Civita}{}{}
pada $U$
\definitionsverweis {trivial}{}{.}
Secara khusus, berlaku pernyataan-pernyataan berikut.
\aufzaehlungvier{Semua
\definitionsverweis {simbol Christoffel}{}{}
memenuhi

\mavergleichskettedisp
{\vergleichskette
{ \Gamma_{ij}^k
}
{ =} { 0
}
{ } {
}
{ } {
}
{ } {
}
}
{}{}{}
untuk setiap $i,j,k$.
}{Berlaku

\mavergleichskettedisp
{\vergleichskette
{ \nabla_{\partial_i }  \partial_j
}
{ =} { 0
}
{ } {
}
{ } {
}
{ } {
}
}
{}{}{}
untuk
\definitionsverweis {medan vektor standar}{}{}
$\partial_i$.
}{Untuk fungsi-fungsi terdiferensialkan $f_j$ berlaku

\mavergleichskettedisp
{\vergleichskette
{ \nabla_{\partial_i }  \left(  \sum_{j = 1}^n f_j \partial_j  \right)
}
{ =} { \sum_{j = 1}^n \partial_i(f_j) \partial_j
}
{ } {
}
{ } {
}
{ } {
}
}
{}{}{.}
}{Untuk suatu medan vektor $V$ dan suatu medan vektor terdiferensialkan $W$ pada $U$, berlaku

\mavergleichskettedisp
{\vergleichskette
{ ( \nabla_V W)(P)
}
{ =} { \left(  D W   \right)_{ P } \left(   V(P)  \right)
}
{ } {
}
{ } {
}
{ } {
}
}
{}{}{.}
}}
\faktzusatz {}
\faktzusatz {}

}
{

Fungsi-fungsi fundamental Riemann $g_{ij}$ semuanya konstan, sehingga turunan parsialnya yang muncul dalam definisi simbol Christoffel lenyap. Karena itu, menurut
Catatan 25.8,
koneksi tersebut adalah koneksi trivial. Sifat-sifat lainnya mengikuti dari
Latihan 25.11.

}

\inputbeispiel{}
{

Misalkan

\mavergleichskette
{\vergleichskette
{ I
}
{ \subseteq }{ \R
}
{  }{
}
{    }{
}
{   }{
}
}
{}{}{}
suatu
\definitionsverweis {interval riil}{}{}
dan misalkan
\maabbdisp {g} { I } {\R_+
} {}
suatu fungsi positif terdiferensialkan, yang kita tafsirkan sebagai
\definitionsverweis {metrik Riemann}{}{}
pada $I$; lihat
Contoh 16.10.
Satu-satunya
\definitionsverweis {simbol Christoffel}{}{}
untuk koneksi Levi--Civita ialah

\mavergleichskettedisp
{\vergleichskette
{ \Gamma
}
{ =} { { \frac{   1  }{  2 } }  g^{-1} g'
}
{ } {
}
{ } {
}
{ } {
}
}
{}{}{,}
yakni, selain faktor di depannya, merupakan
\definitionsverweis {turunan logaritmik}{}{}
dari $g$. Turunan vertikalnya adalah

\mavergleichskettedisp
{\vergleichskette
{ \nabla_\partial \partial
}
{ =} { { \frac{  g' }{  2g } } \partial
}
{ } {
}
{ } {
}
{ } {
}
}
{}{}{,}
dengan $\partial$ menyatakan turunan pada $I$. Jika diterapkan pada suatu fungsi yang terdiferensialkan secara kontinu
\maabb {f} { I } {\R
} {}
\zusatzklammer {atau, setara dengan itu, pada medan arah $f \partial$} {} {,}
maka menurut
Teorema 25.5
berlaku

\mavergleichskettedisp
{\vergleichskette
{ \nabla_\partial (f)
}
{ =} { \partial f + f  { \frac{  g' }{  2g } }
}
{ =} { f' + f { \frac{  g' }{  2g } }
}
{ } {
}
{ } {
}
}
{}{}{.}

}



\inputfaktbeweis
{Riemannsche Mannigfaltigkeit/Zweidimensional/Levi-Civita-Zusammenhang/Christoffelsymbole/Beziehung/Fakt}
{Lemma}
{}
{

\faktsituation {Misalkan

\mavergleichskette
{\vergleichskette
{ U
}
{ \subseteq }{ \R^2
}
{  }{
}
{    }{
}
{   }{
}
}
{}{}{}
\definitionsverweis {terbuka}{}{}
dan dilengkapi dengan suatu struktur Riemann yang diberikan oleh
\definitionsverweis {fungsi-fungsi terdiferensialkan}{}{}
\maabbdisp {g_{11},g_{12},g_{22}} { U } {\R
} {}
serta
\definitionsverweis {struktur Riemann}{}{}
tersebut mempunyai matriks invers $\left(  g^{kl}  \right)$.}
\faktfolgerung {Maka
\definitionsverweis {simbol Christoffel}{}{}
untuk koneksi Levi--Civita pada $U$ adalah

\mavergleichskettedisphandlinks
{\vergleichskettedisphandlinks
{ \Gamma_{ij}^k
}
{ =} { { \frac{   1  }{  2 } } \left(  g^{k1}(\partial_ig_{j1}+\partial_jg_{i1}-\partial_1 g_{ij})   +  g^{k2}(\partial_ig_{j2}+\partial_jg_{i2}-\partial_2 g_{ij})  \right)
}
{ } {
}
{ } {
}
{ } {
}
}
{}{}{.}
Jika

\mavergleichskette
{\vergleichskette
{ Y
}
{ \subseteq }{ \R^3
}
{  }{
}
{    }{
}
{   }{
}
}
{}{}{}
suatu permukaan berorientasi yang dua kali terdiferensialkan secara kontinu dan dilengkapi dengan struktur yang diinduksi oleh ruang sekitar $\R^3$ sebagai
\definitionsverweis {struktur Riemann}{}{,}
serta
\maabbdisp {\varphi} { V } { Y
} {,}

\mavergleichskette
{\vergleichskette
{ V
}
{ \subseteq }{ \R^2
}
{  }{
}
{    }{
}
{   }{
}
}
{}{}{}
terbuka, merupakan suatu parametrisasi lokal $Y$ yang dua kali terdiferensialkan, maka simbol-simbol Christoffel ini berimpit dengan
\definitionsverweis {simbol Christoffel}{}{}
yang didefinisikan dengan bantuan ruang sekitar.}
\faktzusatz {}
\faktzusatz {}

}
{

Pernyataan pertama langsung mengikuti dari
Definisi 26.1.
Pernyataan kedua mengikuti dari pernyataan pertama dan
Lemma 19.5.

}



\inputfaktbeweis
{Offene Menge/Riemannsche Struktur/Christoffelsymbole/Eigenschaften/Fakt}
{Lemma}
{}
{

\faktsituation {Misalkan

\mavergleichskette
{\vergleichskette
{ U
}
{ \subseteq }{ \R^n
}
{  }{
}
{    }{
}
{   }{
}
}
{}{}{}
terbuka dan padanya diberikan suatu
\definitionsverweis {struktur Riemann}{}{}
melalui fungsi-fungsi $\left(  g_{ij}  \right)$. Maka
\definitionsverweis {simbol Christoffel}{}{}
untuk koneksi Levi--Civita pada $U$ memenuhi

\mavergleichskettedisp
{\vergleichskette
{ \Gamma^k_{ij}
}
{ =} { \Gamma^k_{ji}
}
{ } {
}
{ } {
}
{ } {
}
}
{}{}{.}}
\faktfolgerung {}
\faktzusatz {}
\faktzusatz {}

}
{

Hal ini langsung mengikuti dari definisi dan kesimetrian
\mathl{\left(  g_{ij}  \right)}{.}

}

\inputbeispiel{}
{

Kita melanjutkan
Contoh 18.8,
yakni kita meninjau setengah bidang Poincaré $\mathbb H$ dengan fungsi-fungsi fundamental Riemann

\mavergleichskette
{\vergleichskette
{ g_{11}
}
{ = }{ g_{22}
}
{ = }{ { \frac{   1  }{  y^2 } }
}
{    }{
}
{   }{
}
}
{}{}{}
dan

\mavergleichskette
{\vergleichskette
{ g_{12}
}
{ = }{ 0
}
{  }{
}
{    }{
}
{   }{
}
}
{}{}{.}
Menurut
Lemma 26.5,
dari

\mathdisp {{ \frac{   1  }{  2 } } \left(  g^{k1}(\partial_ig_{j1}+\partial_jg_{i1}-\partial_1 g_{ij})   +  g^{k2}(\partial_ig_{j2}+\partial_jg_{i2}-\partial_2 g_{ij})  \right)} {  }

diperoleh

\mavergleichskettedisp
{\vergleichskette
{ \Gamma^1_{11}
}
{ =} { { \frac{   1  }{  2 } }  \left(  y^2  \partial_1  \left( \frac{   1  }{  y^2 } \right)     \right)
}
{ =} { 0
}
{ } {
}
{ } {}
}
{}{}{,}

\mavergleichskettedisp
{\vergleichskette
{ \Gamma^2_{11}
}
{ =} { { \frac{   1  }{  2 } }  \left(  y^2 \left(  - \partial_2  \left( \frac{   1  }{  y^2 } \right)  \right)      \right)
}
{ =} { { \frac{   1  }{  2 } }  \left(  y^2 \left(  -  \left( \frac{   -2 }{  y^3 } \right)  \right)      \right)
}
{ =} { { \frac{   1  }{  y } }
}
{ } {
}
}
{}{}{}
dan dengan cara serupa

\mavergleichskette
{\vergleichskette
{ \Gamma^1_{12}
}
{ = }{ - { \frac{   1  }{  y } }
}
{  }{
}
{    }{
}
{   }{
}
}
{}{}{,}

\mavergleichskette
{\vergleichskette
{ \Gamma^2_{12}
}
{ = }{ 0
}
{  }{
}
{    }{
}
{   }{
}
}
{}{}{,}

\mavergleichskette
{\vergleichskette
{ \Gamma^1_{22}
}
{ = }{ 0
}
{  }{
}
{    }{
}
{   }{
}
}
{}{}{,}

\mavergleichskette
{\vergleichskette
{ \Gamma^2_{22}
}
{ = }{ - { \frac{   1  }{  y } }
}
{  }{
}
{    }{
}
{   }{
}
}
{}{}{.}

}

  \zwischenueberschrift{Koneksi Levi--Civita}

Kita hendak memperluas konsep koneksi Levi--Civita dari suatu himpunan terbuka

\mavergleichskette
{\vergleichskette
{ U
}
{ \subseteq }{ \R^n
}
{  }{
}
{    }{
}
{   }{
}
}
{}{}{}
dengan struktur Riemann ke suatu manifold Riemann sembarang. Kesulitan utamanya adalah menunjukkan bahwa konstruksi langsung pada daerah-daerah peta terbuka saling konsisten pada irisan yang tumpang tindih. Untuk itu, kita memperkenalkan konsep-konsep berikut.

\inputdefinition
{}
{

Misalkan $M$ suatu
\definitionsverweis {manifold terdiferensialkan}{}{}
dan diberikan suatu
\definitionsverweis {koneksi linear}{}{}
pada
\definitionsverweis {bundel tangen}{}{}
$TM$. Koneksi tersebut disebut
\definitionswort {bebas torsi}{,}
jika

\mavergleichskettedisp
{\vergleichskette
{ \nabla_V W- \nabla_W V
}
{ =} { [V,W]
}
{ } {
}
{ } {
}
{ } {
}
}
{}{}{}
untuk setiap
\definitionsverweis {medan vektor}{}{}
$V,W$ yang
\definitionsverweis {terdiferensialkan secara kontinu}{}{}
pada setiap himpunan terbuka.

}

\inputdefinition
{}
{

Misalkan $M$ suatu
\definitionsverweis {manifold Riemann}{}{}
dan diberikan suatu
\definitionsverweis {koneksi linear}{}{}
pada
\definitionsverweis {bundel tangen}{}{}
$TM$. Koneksi tersebut disebut
\definitionswort {metrik}{,}
jika

\mavergleichskettedisp
{\vergleichskette
{ D_V \left\langle  W  ,  Z  \right\rangle
}
{ =} { \left\langle \nabla_V W , Z  \right\rangle  + \left\langle  W  ,  \nabla_VZ  \right\rangle
}
{ } {
}
{ } {
}
{ } {
}
}
{}{}{.}
untuk setiap
\definitionsverweis {medan vektor}{}{}
$V,W,Z$ yang
\definitionsverweis {terdiferensialkan secara kontinu}{}{}
pada setiap himpunan terbuka.

}



\inputfaktbeweis
{Riemannsche Mannigfaltigkeit/Tangentialbündel/Linearer Zusammenhang/Metrisch und torsionsfrei/Koszul-Eigenschaft/Eindeutigkeit/Fakt}
{Lemma}
{}
{

\faktsituation {Misalkan $M$ suatu
\definitionsverweis {manifold Riemann}{}{}
dan diberikan suatu
\definitionsverweis {koneksi linear}{}{}
$\nabla$ pada
\definitionsverweis {bundel tangen}{}{}
$TM$ yang
\definitionsverweis {metrik}{}{}
dan
\definitionsverweis {bebas torsi}{}{.}}
\faktfolgerung {Maka untuk setiap
\definitionsverweis {medan vektor}{}{}
$V,W,Z$ yang
\definitionsverweis {terdiferensialkan secara kontinu}{}{}
pada setiap himpunan terbuka berlaku apa yang disebut \stichwort {rumus Koszul} {}

\mavergleichskettedisphandlinks
{\vergleichskettedisphandlinks
{ \left\langle \nabla_V W , Z  \right\rangle
}
{ =} { { \frac{   1  }{  2 } }  \left(  D_V \left\langle  W  ,  Z  \right\rangle +  D_W \left\langle  Z  ,  V  \right\rangle -  D_Z \left\langle  V  ,  W  \right\rangle  - \left\langle  W  , [V,Z]   \right\rangle - \left\langle  Z  , [W,V]   \right\rangle + \left\langle  V  , [Z,W]   \right\rangle   \right)
}
{ } {
}
{ } {
}
{ } {
}
}
{}{}{.}}
\faktzusatz {Secara khusus, paling banyak ada satu koneksi dengan sifat-sifat ini.}
\faktzusatz {}

}
{

Karena koneksi bebas torsi, berlaku

\mavergleichskettedisp
{\vergleichskette
{ \nabla_VW- \nabla_WV
}
{ =} { [V,W]
}
{ } {
}
{ } {
}
{ } {
}
}
{}{}{,}

\mavergleichskettedisp
{\vergleichskette
{ \nabla_VZ- \nabla_ZV
}
{ =} { [V,Z]
}
{ } {
}
{ } {
}
{ } {
}
}
{}{}{}
dan

\mavergleichskettedisp
{\vergleichskette
{ \nabla_WZ- \nabla_ZW
}
{ =} { [W,Z]
}
{ } {
}
{ } {
}
{ } {
}
}
{}{}{,}
dan identitas pertama juga dapat ditulis sebagai

\mavergleichskettedisp
{\vergleichskette
{ \nabla_VW + \nabla_WV
}
{ =} { 2 \nabla_VW -[V,W]
}
{ } {
}
{ } {
}
{ } {
}
}
{}{}{.}
Karena koneksi metrik, berlaku identitas-identitas

\mavergleichskettedisp
{\vergleichskette
{ D_V \left\langle  W  ,  Z  \right\rangle
}
{ =} { \left\langle \nabla_V W , Z  \right\rangle  + \left\langle  W  ,  \nabla_V Z  \right\rangle
}
{ } {
}
{ } {
}
{ } {
}
}
{}{}{,}

\mavergleichskettedisp
{\vergleichskette
{ D_W \left\langle  Z  ,  V  \right\rangle
}
{ =} { \left\langle \nabla_WZ , V  \right\rangle  + \left\langle  Z  ,  \nabla_WV  \right\rangle
}
{ } {
}
{ } {
}
{ } {
}
}
{}{}{}
dan

\mavergleichskettedisp
{\vergleichskette
{ D_Z \left\langle  V  ,  W  \right\rangle
}
{ =} { \left\langle \nabla_Z V , W  \right\rangle  + \left\langle  V  ,  \nabla_Z W   \right\rangle
}
{ } {
}
{ } {
}
{ } {
}
}
{}{}{.}
Dengan menjumlahkan persamaan-persamaan ini, kita memperoleh

\mavergleichskettedisphandlinks
{\vergleichskettedisphandlinks
{ D_V \left\langle  W  ,  Z  \right\rangle+ D_W \left\langle  Z  ,  V  \right\rangle  - D_Z \left\langle  V  ,  W  \right\rangle
}
{ =} { \left\langle \nabla_V W+\nabla_W V  ,  Z  \right\rangle + \left\langle \nabla_W Z-\nabla_Z W   ,  V  \right\rangle+\left\langle  \nabla_V Z-\nabla_Z V   ,  W  \right\rangle
}
{ } {
}
{ } {
}
{ } {
}
}
{}{}{.}
Kita substitusikan ungkapan-ungkapan di atas ke ruas kanan dan memperoleh

\mavergleichskettealignhandlinks
{\vergleichskettealignhandlinks
{ D_V \left\langle  W  ,  Z  \right\rangle+ D_W \left\langle  Z  ,  V  \right\rangle  - D_Z \left\langle  V  ,  W  \right\rangle
}
{ =} { \left\langle  2 \nabla_VW -[V,W]  ,  Z  \right\rangle + \left\langle  [W,Z]  ,  V  \right\rangle + \left\langle  [V,Z]   ,  W  \right\rangle
}
{ =} { 2\left\langle  \nabla_VW   ,  Z  \right\rangle - \left\langle  [V,W]  ,  Z  \right\rangle + \left\langle  [W,Z]  ,  V  \right\rangle + \left\langle  [V,Z]   ,  W  \right\rangle
}
{ } {
}
{ } {
}
}
{}
{}{.}
Dengan menata ulang persamaan tersebut, diperoleh rumus yang dinyatakan.

Berdasarkan persamaan itu,
\mathl{\left\langle  \nabla_V W , Z  \right\rangle}{} untuk medan-medan vektor sembarang telah ditentukan oleh ruas kanan, yang sama sekali tidak memuat koneksi. Karena hal ini berlaku untuk setiap medan vektor $Z$, turunan vertikal $\nabla_V W$, dan dengan demikian koneksi linearnya, juga ditentukan secara unik.

}

\inputfaktbeweis
{R^n/Offene Menge/Riemannsche Struktur/Zusammenhang aus Christoffel-Symbole/Eigenschaften/Fakt}
{Lemma}
{}
{

\faktsituation {Misalkan

\mavergleichskette
{\vergleichskette
{ U
}
{ \subseteq }{ \R^n
}
{  }{
}
{    }{
}
{   }{
}
}
{}{}{}
\definitionsverweis {terbuka}{}{}
dan padanya diberikan suatu
\definitionsverweis {struktur Riemann}{}{}
melalui fungsi-fungsi $\left(  g_{ij}  \right)$ dengan
\definitionsverweis {matriks invers}{}{}

\mathl{\left(  g^{k l}  \right)}{.} Misalkan $\nabla$ adalah koneksi yang diberikan oleh
\definitionsverweis {simbol-simbol Christoffel}{}{}

\mavergleichskettedisp
{\vergleichskette
{ \Gamma_{ij}^k
}
{ =} { { \frac{   1  }{  2 } } \sum_{l = 1}^n g^{kl}(\partial_ig_{jl}+\partial_jg_{il}-\partial_lg_{ij})
}
{ } {
}
{ } {
}
{ } {
}
}
{}{}{}
sebagai
\definitionsverweis {koneksi Levi--Civita}{}{,}
yang dicirikan oleh

\mavergleichskettedisp
{\vergleichskette
{ \nabla_{\partial_i } \partial_j
}
{ =} { \sum_{k = 1}^n\Gamma_{ij}^k\,\partial_k
}
{ } {
}
{ } {
}
{ } {
}
}
{}{}{.}}
\faktuebergang {Maka $\nabla$ memenuhi sifat-sifat berikut.}
\faktfolgerung {\aufzaehlungvier{Berlaku

\mavergleichskettedisphandlinks
{\vergleichskettedisphandlinks
{ \left\langle  \nabla_{\partial_i }  \partial_j  ,  \partial_k   \right\rangle
}
{ =} { { \frac{   1  }{  2 } }  \left(  D_{\partial_i }  \left\langle  \partial_j  ,  \partial_k   \right\rangle   +D_{\partial_j }  \left\langle  \partial_i  ,  \partial_k   \right\rangle -D_{\partial_k }  \left\langle  \partial_i  ,  \partial_j   \right\rangle   \right)
}
{ } {}
{ } {}
{ } {}
}
{}{}{.}
}{Koneksi ini memenuhi rumus Koszul dari
Lemma 26.10.
}{Koneksi ini
\definitionsverweis {bebas torsi}{}{.}
}{Koneksi ini
\definitionsverweis {metrik}{}{.}
}}
\faktzusatz {}
\faktzusatz {}

}
{

\aufzaehlungvier{Berlaku

\mavergleichskettedisp
{\vergleichskette
{ D_{\partial_k }  \left\langle  \partial_i  ,  \partial_j   \right\rangle
}
{ =} { \partial_k  \left(  g_{ij}  \right)
}
{ } {
}
{ } {
}
{ } {
}
}
{}{}{.}
Selanjutnya,

\mavergleichskettealignhandlinks
{\vergleichskettealignhandlinks
{ \left\langle  \nabla_{\partial_i }  \partial_j  ,  \partial_k   \right\rangle
}
{ =} { \left\langle  \sum_{\ell = 1}^n\Gamma_{ij}^\ell \,\partial_\ell  ,  \partial_k   \right\rangle
}
{ =} { \sum_{\ell = 1}^n\Gamma_{ij}^\ell \left\langle  \,\partial_\ell  ,  \partial_k   \right\rangle
}
{ =} { \sum_{\ell = 1}^n\Gamma_{ij}^\ell g_{\ell k}
}
{ =} { { \frac{   1  }{  2 } } \sum_{\ell = 1}^n  \left(  \sum_{r = 1}^n g^{\ell r}  \left(  \partial_ig_{j r}+\partial_jg_{ir}-\partial_r g_{ij}  \right)   \right)  g_{\ell k}
}
}
{
\vergleichskettefortsetzungalign
{ =} { { \frac{   1  }{  2 } } \sum_{r = 1}^n  \left(  \partial_ig_{j r}+\partial_jg_{ir}-\partial_r g_{ij}  \right)  \left(  \sum_{\ell = 1}^n  g^{\ell r} g_{\ell k}   \right)
}
{ =} { { \frac{   1  }{  2 } }  \left(  \partial_ig_{j k}+\partial_jg_{i k}-\partial_k g_{ij}  \right)
}
{ =} { { \frac{   1  }{  2 } }  \left(  D_{\partial_i }  \left\langle  \partial_j  ,  \partial_k   \right\rangle   +D_{\partial_j }  \left\langle  \partial_i  ,  \partial_k   \right\rangle   -D_{\partial_k }  \left\langle  \partial_i  ,  \partial_j   \right\rangle   \right)
}
{ } {}
}
{}{.}
}{Karena kurung Lie
\mathl{[ \partial_i, \partial_j]}{} trivial pada $U$, menurut bagian (1) rumus Koszul berlaku untuk medan-medan basis $\partial_i$. Untuk kasus umum, lihat
Latihan 26.7.
}{Menurut
Lemma 26.6
dan karena

\mavergleichskette
{\vergleichskette
{ [\partial_i, \partial_j]
}
{ = }{ 0
}
{  }{
}
{    }{
}
{   }{
}
}
{}{}{}
berlaku

\mavergleichskettedisp
{\vergleichskette
{ \nabla_{\partial_i }  \partial_j
}
{ =} { \nabla_{\partial_j }  \partial_i
}
{ } {
}
{ } {
}
{ } {
}
}
{}{}{.}
Pernyataan itu sekarang mengikuti dari
Latihan 26.5.
}{Menurut bagian (1),

\mavergleichskettealignhandlinks
{\vergleichskettealignhandlinks
{ \left\langle  \nabla_{\partial_i }  \partial_j  ,  \partial_k   \right\rangle + \left\langle  \nabla_{\partial_i }  \partial_k  ,  \partial_j   \right\rangle
}
{ =} { { \frac{   1  }{  2 } }  \left(  D_{\partial_i }  \left\langle  \partial_j  ,  \partial_k   \right\rangle   +D_{\partial_j }  \left\langle  \partial_i  ,  \partial_k   \right\rangle   -D_{\partial_k }  \left\langle  \partial_i  ,  \partial_j   \right\rangle    +D_{\partial_i }  \left\langle  \partial_k  ,  \partial_j   \right\rangle +D_{\partial_k }  \left\langle  \partial_i  ,  \partial_j   \right\rangle  - D_{\partial_j }  \left\langle  \partial_i  ,  \partial_k   \right\rangle   \right)
}
{ =} { D_{\partial_i }  \left\langle  \partial_j  ,  \partial_k   \right\rangle
}
{ } {}
{ } {}
}
{}
{}{.}
Pernyataan itu sekarang mengikuti dari
Latihan 26.8.
}

}



\inputfaktbeweis
{Riemannsche Mannigfaltigkeit/Levi-Civita-Zusammenhang/Eindeutige Existenz/Fakt}
{Teorema}
{}
{

\faktsituation {Misalkan $X$ suatu
\definitionsverweis {manifold Riemann}{}{.}
Maka pada
\definitionsverweis {bundel tangen}{}{}
$TX$ terdapat suatu
\definitionsverweis {koneksi linear}{}{}
yang
\definitionsverweis {bebas torsi}{}{,}
\definitionsverweis {metrik}{}{,}
dan ditentukan secara unik.}
\faktfolgerung {}
\faktzusatz {}
\faktzusatz {}

}
{

Menurut
Lemma 26.10,
paling banyak ada satu koneksi semacam itu. Menurut
Lemma 26.11,
pada setiap daerah peta kita dapat memberikan secara eksplisit suatu koneksi yang mempunyai sifat-sifat yang disyaratkan. Karena keunikan yang baru saja disebutkan, koneksi-koneksi yang dikonstruksi dengan cara ini berimpit pada irisan daerah-daerah peta.

}

\inputdefinition
{}
{

Misalkan $X$ suatu
\definitionsverweis {manifold Riemann}{}{.}
\definitionsverweis {Koneksi}{}{}
pada
\definitionsverweis {bundel tangen}{}{}
$TX$ yang didefinisikan secara lokal oleh
\definitionsverweis {simbol-simbol Christoffel}{}{}
disebut
\definitionswort {koneksi Levi--Civita}{.}

}



\inputfaktbeweis
{Riemannsche Mannigfaltigkeit/Levi-Civita-Zusammenhang/Eigenschaften/Fakt}
{Teorema}
{}
{

\faktsituation {Misalkan $X$ suatu
\definitionsverweis {manifold Riemann}{}{.}
Maka
\definitionsverweis {koneksi Levi--Civita}{}{}
$\nabla$ memenuhi sifat-sifat berikut.}
\faktfolgerung {\aufzaehlungdrei{ $\nabla$
\definitionsverweis {linear}{}{.}
}{ $\nabla$ metrik, yakni

\mavergleichskettedisp
{\vergleichskette
{ D_V \left\langle  W  ,  Z  \right\rangle
}
{ =} { \left\langle \nabla_V W , Z  \right\rangle  + \left\langle  W  ,  \nabla_VZ  \right\rangle
}
{ } {
}
{ } {
}
{ } {
}
}
{}{}{.}
}{ $\nabla$ bebas torsi, yakni

\mavergleichskettedisp
{\vergleichskette
{ \nabla_V W- \nabla_W V
}
{ =} { [V,W]
}
{ } {
}
{ } {
}
{ } {
}
}
{}{}{.}
}}
\faktzusatz {}
\faktzusatz {}

}
{

Hal ini mengikuti dari
Lemma 26.11.

}



\inputfaktbeweis
{R^3/Riemannsche orientierte Fläche/Levi-Civita-Zusammenhang und induzierter Zusammenhang/Fakt}
{Lemma}
{}
{

\faktsituation {Misalkan

\mavergleichskette
{\vergleichskette
{ Y
}
{ \subseteq }{ \R^3
}
{  }{
}
{    }{
}
{   }{
}
}
{}{}{}
suatu permukaan yang dua kali terdiferensialkan dan dilengkapi oleh ruang sekitar $\R^3$ dengan
\definitionsverweis {struktur Riemann}{}{}
terinduksi.}
\faktfolgerung {Maka
\definitionsverweis {koneksi Levi--Civita}{}{}
pada $Y$ berimpit dengan apa yang
\zusatzklammer {melalui
\definitionsverweis {proyeksi ortogonal}{}{}
diberikan pada ruang-ruang tangen} {} {}
didefinisikan sebagai
\definitionsverweis {koneksi}{}{}
dari
Catatan 24.8.}
\faktzusatz {}
\faktzusatz {}

}
{

Kedua koneksi itu linear. Karena itu, cukup ditunjukkan secara lokal bahwa
\definitionsverweis {turunan vertikal}{}{}
terkait sama pada medan-medan basis. Misalkan

\mavergleichskette
{\vergleichskette
{ V
}
{ \subseteq }{ \R^2
}
{  }{
}
{    }{
}
{   }{
}
}
{}{}{}
\definitionsverweis {terbuka}{}{}
dan
\maabbdisp {\varphi} { V } { U
} {}
suatu parametrisasi lokal yang dua kali terdiferensialkan secara kontinu dari sebuah himpunan bagian terbuka

\mavergleichskette
{\vergleichskette
{ U
}
{ \subseteq }{ Y
}
{  }{
}
{    }{
}
{   }{
}
}
{}{}{.}
Misalkan $\partial_1,\partial_2$ adalah kedua medan arah konstan pada $V$. Untuk koneksi Levi--Civita berlaku

\mavergleichskettedisp
{\vergleichskette
{ \nabla_{\partial_i } \partial_j
}
{ =} { \sum_{k = 1}^2  \Gamma^k_{ij} \partial_k
}
{ } {
}
{ } {
}
{ } {
}
}
{}{}{,}
dengan simbol-simbol Christoffel yang didefinisikan berdasarkan

\mavergleichskettedisp
{\vergleichskette
{ g_{ij}
}
{ =} { \left\langle  \partial_i \varphi ,  \partial_j \varphi  \right\rangle
}
{ } {
}
{ } {
}
{ } {
}
}
{}{}{.}
Menurut
Lemma 26.5,
simbol-simbol Christoffel ini juga memenuhi

\mavergleichskettedisp
{\vergleichskette
{ \partial_i \partial_j \varphi
}
{ =} { \Gamma^1_{ij} \partial_1  \varphi +  \Gamma^2_{ij}  \partial_2  \varphi+   h_{ij} N
}
{ } {
}
{ } {
}
{ } {
}
}
{}{}{,}
dengan $N$ menyatakan
\definitionsverweis {medan normal satuan}{}{}
dan $h_{ij}$ merupakan entri-entri
\definitionsverweis {matriks fundamental kedua}{}{}
\zusatzklammer {semuanya dipandang pada $V$} {} {.}

Medan-medan vektor $\partial_i$ pada $V$ bersesuaian, pada

\mavergleichskette
{\vergleichskette
{ U
}
{ \subseteq }{ Y
}
{ \subseteq }{ \R^3
}
{    }{
}
{   }{
}
}
{}{}{}
dengan medan-medan vektor $\left(  \partial_i \varphi  \right)  \circ \varphi^{-1}$. Menurut definisi
\definitionsverweis {turunan vertikal}{}{}
untuk koneksi terbatas tersebut, kita harus meninjau

\mathdisp {U \stackrel{ \left(  \partial_i\varphi  \right) \circ \varphi^{-1} }{\longrightarrow} TU \stackrel{ T \left(  \left(  \partial_j \varphi  \right) \circ \varphi^{-1}    \right) }{\longrightarrow} TTU \stackrel{ \pi_{\text{vert} } } \longrightarrow TU} {  }

Ini menghasilkan proyeksi ortogonal
\mathl{\left(  \partial_i \partial_j \varphi  \right) \circ \varphi^{-1}}{}
ke bundel tangen pada $Y$, yang berimpit dengan
\mathl{\Gamma^1_{ij} \partial_1  \varphi +  \Gamma^2_{ij}  \partial_2  \varphi}{,}
dipandang di atas $U$
\zusatzklammer {lihat juga
Latihan 26.11} {} {.}

}

Pernyataan ini juga berlaku dengan cara yang sama untuk setiap submanifold tertutup

\mavergleichskette
{\vergleichskette
{ Y
}
{ \subseteq }{ \R^n
}
{  }{
}
{    }{
}
{   }{
}
}
{}{}{}
yang dilengkapi dengan struktur Riemann terinduksi.


\chapter{Lembar Kerja 26}
\input{generated/worksheet26.id.build.tex}

\section*{Solusi yang disediakan oleh sumber}
\addcontentsline{toc}{section}{Solusi yang disediakan oleh sumber}
\subsection*{Solusi Soal 26.3}
Penampang horizontal dicirikan oleh lenyapnya pemetaan vertikalnya. Dalam situasi ini, menurut
contoh tersebut,
turunan vertikal sama dengan

\mavergleichskettedisp
{\vergleichskette
{ \nabla_\partial (f)
}
{ =} { f' + f { \frac{  g' }{  2g } }
}
{ } {
}
{ } {
}
{ } {}
}
{}{}{.}
Pertama-tama,

\mavergleichskette
{\vergleichskette
{ f
}
{ = }{ 0
}
{  }{
}
{    }{
}
{   }{
}
}
{}{}{}
merupakan suatu solusi. Sekarang misalkan

\mavergleichskette
{\vergleichskette
{ f(t_0)
}
{ \neq }{ 0
}
{  }{
}
{    }{
}
{   }{
}
}
{}{}{}
dan karena itu sifat ini juga berlaku pada suatu lingkungan terbuka dari $t_0$. Pada lingkungan tersebut,

\mavergleichskettedisp
{\vergleichskette
{  { \frac{  f' }{ f } }
}
{ =} { { \frac{  g' }{  2g  } }
}
{ } {
}
{ } {
}
{ } {
}
}
{}{}{.}
Metode penyelesaian persamaan diferensial linear menghasilkan

\mavergleichskettedisp
{\vergleichskette
{   \ln  \betrag {  f  }
}
{ =} { c+ { \frac{   1  }{  2 } }  \ln   g
}
{ } {
}
{ } {
}
{ } {
}
}
{}{}{}
dan dengan demikian

\mavergleichskettedisp
{\vergleichskette
{ \betrag {  f  }
}
{ =} { e^c e^{ { \frac{   1  }{  2 } }  \ln  g }
}
{ =} { e^c  \sqrt{g}
}
{ } {
}
{ } {
}
}
{}{}{}
atau

\mavergleichskettedisp
{\vergleichskette
{ f
}
{ =} { d \sqrt{g}
}
{ } {
}
{ } {
}
{ } {
}
}
{}{}{}
\zusatzklammer {dengan

\mavergleichskette
{\vergleichskette
{ d
}
{ \in }{  \R
}
{  }{
}
{    }{
}
{   }{
}
}
{}{}{}} {} {,}
yang juga dapat diperiksa secara langsung.

\subsection*{Solusi Soal 26.6}
\input{generated/worksheet26_exercise06_solution.id.build.tex}
\subsection*{Solusi Soal 26.9}
\input{generated/worksheet26_exercise09_solution.id.build.tex}

\part{Unit 27}
\chapter[Kuliah 27]{Kuliah 27: Turunan Vertikal, Percepatan Tangensial, dan Geodesik}
\setcounter{section}{27}

  \zwischenueberschrift{Turunan vertikal sepanjang suatu pemetaan}

Misalkan
\maabb {p} { E } { M
} {}
suatu
\definitionsverweis {bundel vektor diferensiabel}{}{}
di atas suatu
\definitionsverweis {manifold diferensiabel}{}{}
$M$ yang dilengkapi dengan suatu
\definitionsverweis {koneksi linear}{}{.}
Koneksi ini menghasilkan suatu
\definitionsverweis {turunan vertikal}{}{.}
Untuk setiap medan vektor
\maabb {V} { M } { TM
} {}
dan setiap penampang yang terdiferensialkan secara kontinu
\maabbdisp {s} { M } { E
} {}
\mathl{\nabla_V s}{} adalah penampang kontinu di $E$ yang diberikan oleh

\mathdisp {M \stackrel{V}{ \longrightarrow  } TM  \stackrel{T(s) }{ \longrightarrow } TE \stackrel{  \pi_{\text{vert} }  }{ \longrightarrow} E} {  }

Pada suatu himpunan terbuka

\mavergleichskette
{\vergleichskette
{ U
}
{ \subseteq }{ M
}
{  }{
}
{    }{
}
{   }{
}
}
{}{}{}
dengan

\mavergleichskette
{\vergleichskette
{ U
}
{ \cong }{ V
}
{ \subseteq }{ \R^n
}
{    }{
}
{   }{
}
}
{}{}{}
serta turunan-turunan parsial
\mathl{\partial_1  , \ldots , \partial_n}{} dan suatu trivialisasi

\mavergleichskettedisp
{\vergleichskette
{ E \vert_U
}
{ \cong} { U \times \R^r
}
{ } {
}
{ } {
}
{ } {
}
}
{}{}{}
dengan penampang basis
\mathl{s_1 , \ldots , s_r}{,} turunan vertikal semacam itu dideskripsikan sepenuhnya oleh simbol-simbol Christoffel $\Gamma_{ij}^k$ melalui

\mavergleichskettedisp
{\vergleichskette
{ \nabla_{\partial_i } (s_j)
}
{ =} { \sum_{k= 1}^r  \Gamma_{ij}^k s_k
}
{ } {
}
{ } {
}
{ } {
}
}
{}{}{}
berdasarkan
Lema 25.10.
Kita ingin memiliki konsep ini bukan hanya bagi penampang di atas $M$, melainkan juga bagi penampang
\maabbdisp {s} { L } { E
} {}
sepanjang suatu
\definitionsverweis {pemetaan diferensiabel}{}{}
tetap
\maabbdisp {\psi} { L } { M
} {}
yaitu ketika terdapat diagram komutatif

\mathdisp {\begin{matrix}  &   &   &  E   \\ & & s \nearrow & \downarrow  \\  &  L  & \stackrel{ \psi }{\longrightarrow}  &  M \! \\ \end{matrix}} {  }

Situasi ini telah kita jumpai pada penampang horizontal. Perhatikan bahwa suatu penampang
\maabbdisp {s} { L } { E
} {}
sama dengan suatu penampang
\maabbdisp {\tilde{s}} { L  } { \psi^*(E) = E \times_XY
} {}
di dalam
\definitionsverweis {bundel vektor tarik balik}{}{}
$\psi^*(E)$. Untuk suatu medan vektor $F$ pada $L$ dan suatu penampang semacam itu, kita dapat meninjau komposisi pemetaan

\mathdisp {L \stackrel{F}{ \longrightarrow  } TL  \stackrel{T(s) }{ \longrightarrow } TE \stackrel{  \pi_{\text{vert} }  }{ \longrightarrow} E} {  }

yang kembali kita nyatakan dengan $\nabla_F s$. Koneksi yang ditentukan oleh turunan vertikal ini kita sebut
\stichwort {koneksi tarik balik} {} $\psi^*\nabla$. Koneksi tersebut mempunyai sifat-sifat berikut.



\inputfaktbeweis
{Linearer Zusammenhang/Vertikale Ableitung/Längs Abbildung/Eigenschaften/Fakt}
{Lema}
{}
{

\faktsituation {Misalkan
\maabb {p} { E } { M
} {}
suatu
\definitionsverweis {bundel vektor diferensiabel}{}{}
di atas suatu
\definitionsverweis {manifold diferensiabel}{}{}
$M$ yang dilengkapi dengan suatu
\definitionsverweis {koneksi linear}{}{}
$\nabla$. Misalkan
\maabb {\psi} { L } { M
} {}
suatu
\definitionsverweis {pemetaan diferensiabel}{}{}
dengan
\definitionsverweis {koneksi tarik balik}{}{}
$\psi^*\nabla$ pada
\definitionsverweis {bundel vektor tarik balik}{}{}
$\psi^*E$ di atas $L$. Maka berlaku sifat-sifat berikut.}
\faktfolgerung {\aufzaehlungzwei {Pemetaan
\maabbeledisp {\psi^*\nabla} {C^1(\psi^* E)} { C^0(\psi^* E \otimes T^* L )
} { s } { \left(  \psi^*\nabla  \right) (s)
} {,}
bersifat
$\R$-\definitionsverweis {linear}{}{.}
Secara khusus, $\psi^*\nabla$ merupakan suatu koneksi linear.
} {Untuk suatu pembatasan
\maabbdisp {\psi} {U'} {U
} {}
pada domain-domain peta terbuka dengan koordinat
\mathl{z_1 , \ldots , z_m}{} di $U'$ dan koordinat
\mathl{x_1 , \ldots , x_n}{}
di $U$, bagi
\definitionsverweis {simbol-simbol Christoffel}{}{}
$\tilde{\Gamma}_{\ell j}^k$ dari $\psi^*\nabla$ dan
\definitionsverweis {penampang basis}{}{}

\mathl{s_1  , \ldots , s_r}{} berlaku hubungan

\mavergleichskettedisp
{\vergleichskette
{ \tilde{\Gamma}_{j \ell }^k
}
{ =} { \sum_{i = 1}^n  \left(  \partial_j \psi_i  \right)   \Gamma_{i\ell }^k
}
{ } {
}
{ } {
}
{ } {
}
}
{}{}{.}
}}
\faktzusatz {}
\faktzusatz {}

}
{

\aufzaehlungzwei {Jelas.
} {Kita tinjau situasi lokal secara langsung. Suatu
\definitionsverweis {penampang basis}{}{}
$s_\ell$ sepanjang $L$ dapat dipandang langsung sebagai penampang basis di atas $M$. Pemetaan-pemetaan yang relevan adalah

\mathdisp {U' \stackrel{\partial_j }{\longrightarrow} TU'  \stackrel{T (\psi) }{\longrightarrow} T(U) \stackrel{ T(s_\ell ) }{\longrightarrow} T( U \times \R^r) \stackrel{  \pi_{\text{vert} } }{\longrightarrow} \R^r} { . }

Di sini,

\mavergleichskettedisp
{\vergleichskette
{ T(\psi) \circ \partial_j
}
{ =} { \partial_j \psi
}
{ =} { \sum_{i = 1}^n \left(  \partial_j \psi_i  \right)  \partial_i
}
{ } {
}
{ } {
}
}
{}{}{.}
Dengan demikian, untuk

\mavergleichskette
{\vergleichskette
{ Q
}
{ \in }{ U'
}
{  }{
}
{    }{
}
{   }{
}
}
{}{}{}
dan dengan menggunakan
Lema 25.10,

\mavergleichskettealignhandlinks
{\vergleichskettealignhandlinks
{ \left(  \left(  \psi^*\nabla  \right) _{\partial_j } (s_\ell)  \right) (Q)
}
{ =} { \left(  \psi^*\nabla  \right) _{\partial_j (Q) } (s_\ell) (Q)
}
{ =} { \nabla_{  \left(  T(\psi) \circ \partial_j  \right)  (Q)  }  (s_\ell) ( \psi(Q) )
}
{ =} { \nabla_{  \sum_{i = 1}^n  \left(  \partial_j \psi_i  \right)  (Q) \partial_i  }  (s_\ell) ( \psi(Q) )
}
{ =} { \sum_{i = 1}^n  \left(  \partial_j \psi_i  \right)  (Q)\nabla_{  \partial_i }  (s_\ell) ( \psi(Q) )
}
}
{
\vergleichskettefortsetzungalign
{ =} { \sum_{i = 1}^n  \left(  \partial_j \psi_i  \right)  (Q)  \left(  \sum_{k = 1}^r \Gamma_{i \ell }^k(\psi(Q)) s_k (\psi(Q))  \right)
}
{ =} { \sum_{k = 1}^r  \left(  \sum_{i = 1}^n  \left(  \partial_j \psi_i  \right)  (Q)  \Gamma_{i \ell }^k(\psi(Q))   \right) s_k (\psi(Q))
}
{ } {}
{ } {}
}
{}{.}
Dengan meninjau komponen ke-$k$, kita memperoleh

\mavergleichskettedisp
{\vergleichskette
{ \tilde{\Gamma}_{j \ell }^k (Q)
}
{ =} { \sum_{i = 1}^n \Gamma_{i\ell }^k (\psi(Q))  \left(  \partial_j \psi_i  \right)  (Q)
}
{ } {
}
{ } {
}
{ } {
}
}
{}{}{.}
}

}



\inputfaktbeweis
{Linearer Zusammenhang/Vertikale Ableitung/Längs Abbildung/Metrisch/Fakt}
{Lema}
{}
{

\faktsituation {Misalkan
\maabb {p} { E } { M
} {}
suatu
\definitionsverweis {bundel vektor diferensiabel}{}{}
di atas suatu
\definitionsverweis {manifold diferensiabel}{}{}
$M$ yang dilengkapi dengan suatu
\definitionsverweis {koneksi linear}{}{}
\definitionsverweis {metrik}{}{}
$\nabla$. Misalkan
\maabb {\psi} { L } { M
} {}
suatu
\definitionsverweis {pemetaan diferensiabel}{}{}
dengan
\definitionsverweis {koneksi tarik balik}{}{}
$\psi^*\nabla$ pada
\definitionsverweis {bundel vektor tarik balik}{}{}
$\psi^*E$ di atas $L$.}
\faktfolgerung {Maka
\mathl{\psi^*\nabla}{} juga metrik.}
\faktzusatz {}
\faktzusatz {}

}
{

Pertama, menurut
Soal 16.13,
$\psi^*E$ kembali merupakan suatu bundel Riemann. Kita tinjau situasi lokal
\maabbdisp {\psi} {U'} {U
} {}
dan

\mavergleichskette
{\vergleichskette
{ E
}
{ = }{ U \times \R^r
}
{  }{
}
{    }{
}
{   }{
}
}
{}{}{}
serta, bersesuaian dengan itu,

\mavergleichskette
{\vergleichskette
{ \psi^*E
}
{ = }{ U' \times \R^r
}
{  }{
}
{    }{
}
{   }{
}
}
{}{}{.}
Fungsi-fungsi yang mendeskripsikan struktur Riemann pada $E$,

\mavergleichskettedisp
{\vergleichskette
{ h_{ \ell m}
}
{ =} { \left\langle s_\ell , s_m   \right\rangle
}
{ } {
}
{ } {
}
{ } {
}
}
{}{}{,}
pada $U$ dan $U'$ berhubungan langsung melalui $\psi$. Misalkan $\partial_j$ suatu
\definitionsverweis {medan vektor standar}{}{}
pada $U'$ dan
\mathl{s_\ell, s_m}{} penampang basis di $E$. Maka

\mavergleichskettealign
{\vergleichskettealign
{ \partial_j  \left\langle s_\ell , s_m   \right\rangle
}
{ =} { \partial_j  \left(  h_{\ell , m} \circ\psi   \right)
}
{ =} { \sum_{i = 1}^n  \left(  \partial_j \psi_i  \right) \left(  \partial_i  h_{\ell , m}  \right)
}
{ =} { \sum_{i = 1}^n  \left(  \partial_j \psi_i  \right) \partial_i  \left\langle s_\ell , s_m   \right\rangle
}
{ } {}
}
{}
{}{.}
Menurut
Lema 27.1,

\mavergleichskettealignhandlinks
{\vergleichskettealignhandlinks
{ \left\langle  \left(  \psi^*\nabla  \right)_{\partial_j } s_\ell ,  s_m   \right\rangle
}
{ =} { \left\langle  \sum_{k = 1}^r \tilde{\Gamma}_{j \ell }^k s_ k ,  s_m   \right\rangle
}
{ =} { \sum_{k = 1}^r \left\langle  \tilde{\Gamma}_{j \ell }^k s_ k ,  s_m   \right\rangle
}
{ =} { \sum_{k = 1}^r \left\langle  \sum_{i = 1}^n \Gamma_{i\ell }^k  \left(  \partial_j \psi_i  \right)  s_k  ,  s_m   \right\rangle
}
{ =} { \sum_{i = 1}^n \left(  \partial_j \psi_i  \right) \left\langle  \sum_{k = 1}^r \Gamma_{i\ell }^k s_k  ,  s_m   \right\rangle
}
}
{
\vergleichskettefortsetzungalign
{ =} { \sum_{i = 1}^n \left(  \partial_j \psi_i  \right) \left\langle  \nabla_{\partial_i } s_\ell  ,  s_m   \right\rangle
}
{ } {}
{ } {}
{ } {}
}
{}{}
dan, secara bersesuaian,

\mavergleichskettedisp
{\vergleichskette
{ \left\langle  s_\ell ,  \left(  \psi^*\nabla  \right)_{\partial_j } s_m   \right\rangle
}
{ =} { \sum_{i = 1}^n  \left(  \partial_j \psi_i  \right)   \left\langle  s_\ell  , \nabla_{\partial_i } s_m   \right\rangle
}
{ } {
}
{ } {
}
{ } {
}
}
{}{}{.}
Karena $\nabla$ metrik, diperoleh

\mavergleichskettedisp
{\vergleichskette
{ \partial_j  \left\langle s_\ell , s_m   \right\rangle
}
{ =} { \left\langle  \left(  \psi^*\nabla  \right)_{\partial_j } s_\ell ,  s_m   \right\rangle    + \left\langle  s_\ell ,  \left(  \psi^*\nabla  \right)_{\partial_j } s_m   \right\rangle
}
{ } {
}
{ } {
}
{ } {
}
}
{}{}{}
dan dari sini, bersama
Soal 26.10,
diperoleh bahwa $\psi^*\nabla$ metrik.

}

  \zwischenueberschrift{Percepatan tangensial}

Misalkan $M$ suatu
\definitionsverweis {manifold}{}{.}
Untuk suatu
\definitionsverweis {kurva diferensiabel}{}{}
\maabbdisp {\gamma} { I } { M
} {}
berlaku bahwa
\maabbdisp {\gamma'} { I} {TM
} {}
merupakan kurva di dalam bundel tangen yang menjadi suatu pengangkatan dari $\gamma$. Jika kurva ini pada gilirannya terdiferensialkan, diperoleh suatu kurva
\maabbdisp {\gamma^{\prime \prime}} { I} {TTM
} {,}
namun kurva ini tidak dapat langsung dihubungkan dengan $\gamma'$ karena nilainya berada di ruang lain. Sebaliknya, apabila pada $TM$ diberikan suatu
\definitionsverweis {koneksi}{}{},
maka melalui turunan vertikal
\zusatzklammer {tarik balik} {} {}
turunan kedua dapat didefinisikan sebagai
\mathl{\nabla_{ \partial } \gamma'}{}.

\inputdefinition
{}
{

Misalkan $M$ suatu
\definitionsverweis {manifold Riemann}{}{}
yang dilengkapi dengan
\definitionsverweis {koneksi Levi--Civita}{}{.}
Untuk suatu
\definitionsverweis {kurva diferensiabel}{}{}
\maabbdisp {\gamma} { I } { M
} {}
yang dua kali terdiferensialkan secara kontinu, pemetaan
\mathl{\pi_{\text{vert} }  \circ TT(\gamma)}{} disebut
\definitionswort {percepatan tangensial}{}
dari $\gamma$.

}

Besaran ini juga ditulis sebagai
\mathl{\pi_{\text{vert} } \circ \gamma^{\prime \prime}}{} atau
\zusatzklammer {sebagai turunan vertikal sepanjang $\gamma$} {} {}

\mathl{\nabla_\partial \gamma'}{,} dengan
\maabbdisp {\gamma'} { I } { TM
} {}
dipandang sebagai penampang dalam bundel tangen sepanjang lintasan $\gamma$, dan dengan meninjau rantai pemetaan

\mathdisp {I \stackrel{\partial}{\longrightarrow} I \times \R = TI \stackrel{T(\gamma')}{\longrightarrow} TTM \stackrel{ \pi_{\text{vert} } }{\longrightarrow}  TM} {  }

.



\inputfaktbeweis
{Hyperfläche/Zweite Ableitung einer Kurve/Direkt und über Zusammenhang/Fakt}
{Lema}
{}
{

\faktsituation {Misalkan

\mavergleichskette
{\vergleichskette
{ W
}
{ \subseteq }{ \R^n
}
{  }{
}
{    }{
}
{   }{
}
}
{}{}{}
\definitionsverweis {terbuka}{}{,}
\maabb {h} { W } { \R
} {}
suatu
\definitionsverweis {fungsi yang terdiferensialkan secara kontinu}{}{}
dan

\mavergleichskette
{\vergleichskette
{ Y
}
{ = }{ h^{-1}(c)
}
{  }{
}
{    }{
}
{   }{
}
}
{}{}{}
\definitionsverweis {serat}{}{}
di atas

\mavergleichskette
{\vergleichskette
{ c
}
{ \in }{ \R
}
{  }{
}
{    }{
}
{   }{
}
}
{}{}{,}
dengan $h$ bersifat
\definitionsverweis {reguler}{}{}
di setiap titik $Y$. Misalkan
\maabbdisp {\gamma} { I } { Y
} {}
suatu kurva yang dua kali terdiferensialkan secara kontinu.}
\faktfolgerung {Maka percepatan tangensial $\gamma$ dalam pengertian
Definisi 1.7
sama dengan percepatan tangensial dalam pengertian
Definisi 27.3.}
\faktzusatz {}
\faktzusatz {}

}
{

Menurut
Lema 26.15,
\zusatzklammer {yang juga berlaku dalam dimensi lebih tinggi} {} {},
kedua koneksi tersebut sama.

}

  \zwischenueberschrift{Kurva geodesik}

\inputdefinition
{}
{

Misalkan $M$ suatu
\definitionsverweis {manifold Riemann}{}{}
yang dilengkapi dengan
\definitionsverweis {koneksi Levi--Civita}{}{.}
Suatu
\definitionsverweis {kurva diferensiabel}{}{}
\maabbdisp {\gamma} { I } { M
} {}
yang dua kali terdiferensialkan disebut
\definitionswort {kurva geodesik}{}
\zusatzklammer {atau
\definitionswort {geodesik}{}} {} {,}
apabila

\mavergleichskettedisp
{\vergleichskette
{ \nabla_{d/dt}  \gamma'
}
{ =} { 0
}
{ } {
}
{ } {
}
{ } {
}
}
{}{}{}
berlaku pada $I$.

}

Kadang-kadang digunakan pula istilah \stichwort {geodesik} {.} Perlu ditekankan bahwa kurva geodesik adalah sebuah kurva; dengan demikian, yang dimaksud bukan hanya
\definitionsverweis {citra}{}{}
kurva tersebut, melainkan juga proses geraknya.



\inputfaktbeweis
{Hyperfläche/Geodätische/Direkt und über Zusammenhang/Fakt}
{Lema}
{}
{

\faktsituation {Misalkan

\mavergleichskette
{\vergleichskette
{ W
}
{ \subseteq }{ \R^n
}
{  }{
}
{    }{
}
{   }{
}
}
{}{}{}
\definitionsverweis {terbuka}{}{,}
\maabb {h} { W } { \R
} {}
suatu
\definitionsverweis {fungsi yang terdiferensialkan secara kontinu}{}{}
dan

\mavergleichskette
{\vergleichskette
{ Y
}
{ = }{ h^{-1}(c)
}
{  }{
}
{    }{
}
{   }{
}
}
{}{}{}
\definitionsverweis {serat}{}{}
di atas

\mavergleichskette
{\vergleichskette
{ c
}
{ \in }{ \R
}
{  }{
}
{    }{
}
{   }{
}
}
{}{}{,}
dengan $h$ bersifat
\definitionsverweis {reguler}{}{}
di setiap titik $Y$.}
\faktfolgerung {Maka geodesik-geodesik di $Y$ dalam pengertian
Definisi 1.8
sama dengan geodesik-geodesik dalam pengertian
Definisi 27.5.}
\faktzusatz {}
\faktzusatz {}

}
{

Hal ini mengikuti dari
Lema 27.4.

}



\inputfaktbeweis
{Riemannsche Mannigfaltigkeit/Levi-Civita-Zusammenhang/Geodätische/Geschwindigkeitskonstanz/Fakt}
{Lema}
{}
{

\faktsituation {Misalkan
\maabb {\gamma} { I } { M
} {}
suatu
\definitionsverweis {kurva geodesik}{}{}
pada suatu
\definitionsverweis {manifold Riemann}{}{}
$M$ yang dilengkapi dengan
\definitionsverweis {koneksi Levi--Civita}{}{.}}
\faktfolgerung {Maka
\mathl{\Vert { \gamma'(t) } \Vert}{} konstan pada $I$.}
\faktzusatz {}
\faktzusatz {}

}
{

Pada suatu domain peta terbuka, dengan menggunakan
Teorema 26.14  (2)
dan
Lema 27.2,
diperoleh

\mavergleichskettealign
{\vergleichskettealign
{ { \frac{   \partial }{  \partial t } }  \left\langle  \gamma'(t) ,  \gamma'(t)  \right\rangle
}
{ =} { \left\langle \nabla_{ { \frac{   \partial }{  \partial t } }  }\gamma'(t) ,  \gamma'(t)  \right\rangle + \left\langle  \gamma'(t) , \nabla_{  { \frac{   \partial }{  \partial t } }  }\gamma'(t)  \right\rangle
}
{ =} { 2 \left\langle \nabla_{ { \frac{   \partial }{  \partial t } }  }\gamma'(t) ,  \gamma'(t)  \right\rangle
}
{ =} { 2 \left\langle  0  ,  \gamma'(t)  \right\rangle
}
{ =} { 0
}
}
{}
{}{,}
sehingga
\mathl{\left\langle  \gamma'(t) ,  \gamma'(t)  \right\rangle}{} konstan.

}



\inputfaktbeweis
{Riemannsche Mannigfaltigkeit/Levi-Civita-Zusammenhang/Geodätische/Differentialgleichung/Fakt}
{Teorema}
{}
{

\faktsituation {Suatu
\definitionsverweis {kurva diferensiabel}{}{}
\maabb {\gamma} { I } { M
} {}
yang dua kali terdiferensialkan pada suatu
\definitionsverweis {manifold Riemann}{}{}
$M$ adalah}
\faktfolgerung {suatu
\definitionsverweis {kurva geodesik}{}{}
\zusatzklammer {terhadap
\definitionsverweis {koneksi Levi--Civita}{}{}} {} {,}
jika dan hanya jika pada setiap peta kurva tersebut memenuhi
\definitionsverweis {sistem persamaan diferensial biasa orde kedua}{}{}

\mavergleichskettedisp
{\vergleichskette
{ \gamma_k^{\prime \prime} + \sum_{ij} \Gamma^k_{ij} \gamma_i' \gamma_j'
}
{ =} { 0
}
{ } {
}
{ } {
}
{ } {
}
}
{}{}{}
\zusatzklammer {untuk semua $k$} {} {}.}
\faktzusatz {}
\faktzusatz {}

}
{

Misalkan $n$ dimensi $M$. Kita meninjau situasinya secara langsung pada suatu citra peta terbuka

\mavergleichskette
{\vergleichskette
{ U
}
{ \subseteq }{ \R^n
}
{  }{
}
{    }{
}
{   }{
}
}
{}{}{.}
Menurut
Catatan 25.8,
turunan vertikal diberikan oleh
\maabbeledisp {} {T(U \times \R^n) = U \times \R^n \times \R^n \times \R^n } { U \times \R^n \times  \R^n
} {(P,e_j, \partial_i, w)} {(P,e_j,\sum_{k = 1}^n \Gamma^k_{ij} (P) e_k +w) \cong V(U \times \R^n)
} {},
sedangkan pemetaan ke $TU$ diperoleh dengan menghilangkan komponen tengah. Turunan kedua kurva mula-mula adalah pemetaan tangen kedua, yaitu
\zusatzklammer {dengan mengabaikan perkalian dengan arah berdimensi satu dari ruang tangen kurva} {} {}
\maabbeledisp {T(\gamma)} { I } { TU
} { t } { (\gamma(t), \gamma'(t))
} {,}
dan, sejalan dengan itu,
\maabbeledisp {T(T(\gamma))} { I } { T(TU)
} { t } { ((\gamma(t), \gamma'(t)), (\gamma'(t), \gamma^{\prime \prime} (t)))
} {}
\zusatzklammer {kedua komponen pemetaan tangen diturunkan} {} {.}
Dengan

\mavergleichskette
{\vergleichskette
{ \gamma'(t)
}
{ = }{ \sum_{j = 1}^n \gamma_j'(t) \partial_j
}
{  }{
}
{    }{
}
{   }{
}
}
{}{}{}
hal ini dipetakan oleh proyeksi vertikal ke

\mathdisp {\sum_{k = 1}^n  \left(  \sum_{i,j} \gamma_i'(t) \gamma_j'(t) \Gamma^k_{ij}(\gamma(t))  \right)  \partial_k  +\sum_{k = 1}^n \gamma_k^{\prime \prime}(t) \partial_k} {  }

Syarat bagi suatu geodesik bahwa turunan keduanya
\zusatzklammer {di dalam $TTM$} {} {}
selalu horizontal ekuivalen dengan syarat bahwa proyeksi vertikal yang dihitung di atas sama dengan $0$. Ini berarti setiap komponennya sama dengan $0$, yakni

\mavergleichskettedisp
{\vergleichskette
{ \sum_{i,j} \gamma_i'(t) \gamma_j'(t) \Gamma^k_{ij}(\gamma(t))  + \gamma_k^{\prime \prime}(t)
}
{ =} { 0
}
{ } {
}
{ } {
}
{ } {
}
}
{}{}{}
untuk

\mavergleichskette
{\vergleichskette
{ k
}
{ = }{ 1  , \ldots , n
}
{  }{
}
{    }{
}
{   }{
}
}
{}{}{.}

}

\inputbemerkung
{}
{

Dalam situasi
Teorema 27.8,
sistem persamaan diferensial

\mavergleichskettedisp
{\vergleichskette
{ \gamma_k^{\prime \prime}(t) + \sum_{ij} \Gamma^k_{ij}(\gamma(t)) \gamma_i'(t) \gamma_j'(t)
}
{ =} { 0
}
{ } {
}
{ } {
}
{ } {
}
}
{}{}{}
dengan

\mavergleichskette
{\vergleichskette
{ k
}
{ = }{ 1  , \ldots , n
}
{  }{
}
{    }{
}
{   }{
}
}
{}{}{,}
yang secara lokal mencirikan geodesik, disebut \stichwort {persamaan diferensial geodesik} {.} Ini adalah suatu
\definitionsverweis {sistem persamaan diferensial biasa}{}{}
orde kedua dengan $n$ persamaan. Sistem ini tidak linear dan secara umum sulit diselesaikan.

}

\inputbeispiel{}
{

\bild{ \begin{center}
\includegraphics[width=5.5cm]{ \bildeinlesungsvg {Parallel lines in Poincare's model of hyperbolic geometry} {svg} }
\end{center}
\bildtext {} }

\bildlizenz { Parallel lines in Poincare's model of hyperbolic geometry.svg } {} {Januszkaja} {Commons} {CC-by-sa 3.0} {}

Kita melanjutkan
Contoh 26.7.
Menurut
Teorema 27.8,
suatu geodesik harus memenuhi kedua syarat

\mavergleichskettedisp
{\vergleichskette
{ \gamma_1^{\prime \prime} (t)
}
{ =} { { \frac{   2  }{  \gamma_2(t) } } \gamma_1'(t) \gamma_2'(t)
}
{ } {
}
{ } {
}
{ } {
}
}
{}{}{}
dan

\mavergleichskettedisp
{\vergleichskette
{ \gamma_2^{\prime \prime} (t)
}
{ =} { -  { \frac{   1  }{  \gamma_2(t) } }  \gamma_1'(t) \gamma_1'(t) + { \frac{   1  }{  \gamma_2(t) } }    \gamma_2'(t) \gamma_2'(t)
}
{ } {
}
{ } {
}
{ } {
}
}
{}{}{}.
Untuk

\mavergleichskette
{\vergleichskette
{ \gamma_1
}
{ = }{ x_0
}
{  }{
}
{    }{
}
{   }{
}
}
{}{}{}
yang konstan, sistem ini menjadi

\mavergleichskettedisp
{\vergleichskette
{ \gamma_2^{\prime \prime} (t)
}
{ =} { { \frac{   1  }{  \gamma_2(t) } }  \gamma_2'(t) \gamma_2'(t)
}
{ } {
}
{ } {
}
{ } {
}
}
{}{}{}
dengan solusi komponen

\mavergleichskettedisp
{\vergleichskette
{ y_2(t)
}
{ =} { e^t
}
{ } {
}
{ } {
}
{ } {
}
}
{}{}{}
dan solusi keseluruhan

\mavergleichskettedisp
{\vergleichskette
{ y(t)
}
{ =} { \left(  x_0   , \,  e^t                   \right)
}
{ } {
}
{ } {
}
{ } {
}
}
{}{}{,}
yang bergerak pada suatu garis vertikal. Selain itu, menurut
Soal 27.13,

\mavergleichskettedisp
{\vergleichskette
{ \gamma(t)
}
{ =} { \left(  s   , \,   0                   \right) + r \left(  \tanh  t   , \,   { \frac{   1  }{  \cosh  t  } }                   \right)
}
{ } {
}
{ } {
}
{ } {
}
}
{}{}{}
untuk

\mavergleichskette
{\vergleichskette
{ s
}
{ \in }{ \R
}
{  }{
}
{    }{
}
{   }{
}
}
{}{}{,}

\mavergleichskette
{\vergleichskette
{ r
}
{ \in }{ \R_+
}
{  }{
}
{    }{
}
{   }{
}
}
{}{}{}
merupakan solusi yang, menurut
Soal 20.50 (Analysis (Osnabrück 2021-2023))  (3),
bergerak pada setengah lingkaran dengan pusat
\mathl{\left(  s   , \,   0                   \right)}{} dan jari-jari $r$.

}

\inputbemerkung
{}
{

Apa yang disebut \stichwort {aksioma paralel} {}
\zusatzklammer {atau postulat paralel} {} {}
dalam geometri Euklides bidang menyatakan bahwa untuk setiap garis dan setiap titik yang tidak terletak pada garis tersebut, ada tepat satu garis paralel yang melalui titik itu. Sebuah garis paralel ditentukan oleh sifat bahwa garis tersebut tidak memotong garis yang diberikan. Salah satu pertanyaan penting dalam sejarah adalah apakah postulat paralel dapat diturunkan dari aksioma dan postulat Euklides lainnya, sehingga postulat ini sebenarnya tidak diperlukan. Pada paruh pertama abad ke-19, pertanyaan ini dijawab secara negatif dengan memberikan model-model geometri
\zusatzklammer {yang khususnya memuat konsep titik dan garis} {} {}
yang memenuhi aksioma-aksioma geometri Euklides lainnya, tetapi tidak memenuhi aksioma paralel. Dalam kerangka geometri Riemann, khususnya dengan konsep
\definitionsverweis {geodesik}{}{},
model-model semacam itu muncul secara alami apabila geodesik-geodesik
\zusatzklammer {di sini yang dimaksud adalah citra suatu kurva geodesik} {} {}
dipandang sebagai garis. Perhatikan bahwa dalam pengembangan aksiomatik suatu geometri, garis tidak ditentukan oleh gambaran intuitif, melainkan oleh sifat-sifat yang dimilikinya dalam hubungannya dengan objek-objek lain. Sebagai contoh,
Contoh 27.10
bersama geodesik-geodesik yang dideskripsikan di sana memberikan suatu model geometri non-Euklides: bagi suatu geodesik yang diberikan dan suatu titik lain, selalu terdapat tak terhingga banyak geodesik melalui titik tersebut yang tidak berpotongan dengan geodesik yang diberikan.

}

\inputdefinition
{}
{

Misalkan $M$ suatu
\definitionsverweis {manifold Riemann}{}{}
yang
\definitionsverweis {terhubung}{}{.}
Untuk titik-titik

\mavergleichskette
{\vergleichskette
{ P,Q
}
{ \in }{ M
}
{  }{
}
{    }{
}
{   }{
}
}
{}{}{}
didefinisikan

\mavergleichskettedisp
{\vergleichskette
{ d(P,Q)
}
{ \defeq} { \operatorname{inf} ( L(\gamma), \gamma \text{ adalah lintasan penghubung kontinu yang terdiferensialkan secara kontinu sepotong-sepotong dari } P \text{ ke } Q  )
}
{ } {
}
{ } {
}
{ } {
}
}
{}{}{}
dan besaran ini disebut
\definitionswort {metrik Riemann}{}
pada $M$.

}

Metrik ini benar-benar menjadikan $M$ suatu ruang metrik yang topologinya sama dengan topologi yang diberikan. Selain itu, dapat dibuktikan bahwa realisasi jarak antara dua titik yang diparameterkan oleh panjang busur merupakan suatu geodesik.


\chapter{Lembar Kerja 27}
\input{generated/worksheet27.id.build.tex}

\section*{Solusi yang disediakan oleh sumber}
\addcontentsline{toc}{section}{Solusi yang disediakan oleh sumber}
\subsection*{Solusi Soal 27.4}
\input{generated/worksheet27_exercise04_solution.id.build.tex}
\subsection*{Solusi Soal 27.5}
Satu-satunya turunan standar pada $\R$ adalah $\partial$, sehingga kita tidak menuliskan indeks bawah pada simbol-simbol Christoffel. Menurut
Fakta (2),
untuk

\mavergleichskette
{\vergleichskette
{ \ell,k
}
{ = }{ 1,2
}
{  }{
}
{    }{
}
{   }{
}
}
{}{}{}

berlaku

\mavergleichskettedisp
{\vergleichskette
{ \tilde{\Gamma}_\ell^k (t)
}
{ =} { \psi_1'(t) \Gamma_{1 \ell}^k (\psi(t)) + \psi'_2(t) \Gamma_{2 \ell}^k (\psi(t))
}
{ =} { \Gamma_{1 \ell}^k (t,e^t)   + e^t \Gamma_{2 \ell}^k (t,e^t)
}
{ } {}
{ } {}
}
{}{}{.}
Dari
contoh
kita memperoleh

\mavergleichskettedisp
{\vergleichskette
{ \tilde{\Gamma}_1^1 (t)
}
{ =} { \Gamma_{1 1}^1 (t,e^t) + e^t \Gamma_{2 1}^1 (t,e^t)
}
{ =} {  -e^t e^{-t}
}
{ =} { -1
}
{ } {
}
}
{}{}{,}

\mavergleichskettedisp
{\vergleichskette
{ \tilde{\Gamma}_1^2 (t)
}
{ =} { \Gamma_{1 1}^2 (t,e^t) + e^t \Gamma_{2 1}^2 (t,e^t)
}
{ =} { e^{-t}
}
{ } {
}
{ } {
}
}
{}{}{,}

\mavergleichskettedisp
{\vergleichskette
{ \tilde{\Gamma}_2^1 (t)
}
{ =} { \Gamma_{1 2}^1 (t,e^t) + e^t \Gamma_{2 2}^1 (t,e^t)
}
{ =} {  -e^{-t}
}
{ } {
}
{ } {
}
}
{}{}{}
dan

\mavergleichskettedisp
{\vergleichskette
{ \tilde{\Gamma}_2^2 (t)
}
{ =} { \Gamma_{1 2}^2 (t,e^t) + e^t \Gamma_{2 2}^2 (t,e^t)
}
{ =} { - e^t e^{-t}
}
{ =} { -1
}
{ } {
}
}
{}{}{.}

\subsection*{Solusi Soal 27.9}
\input{generated/worksheet27_exercise09_solution.id.build.tex}
\subsection*{Solusi Soal 27.13}
\input{generated/worksheet27_exercise13_solution.id.build.tex}

\part{Unit 28}
\chapter[Kuliah 28]{Kuliah 28: Kelengkungan dan Teorema Frobenius}
\setcounter{section}{28}

  \zwischenueberschrift{Kelengkungan}

Misalkan
\mathkor {} {V} {dan} {W} {}
adalah medan vektor yang dua kali terdiferensialkan secara kontinu pada suatu manifold diferensiabel $M$, dan misalkan $\nabla$ adalah
\definitionsverweis {koneksi linear}{}{}
pada
\definitionsverweis {bundel vektor}{}{}
\maabb {p} { E } { M
} {}
di atas $M$. Pemetaan
\maabbeledisp {\nabla_W} {C^2(E) } { C^1 (E)
} { s } { \nabla_W (s)
} {,}
adalah
\definitionsverweis {turunan vertikal}{}{.}
Pada hasil $\nabla_W (s)$ kita dapat sekali lagi menerapkan turunan vertikal dalam arah $V$, sehingga diperoleh penampang kontinu
\mathl{\nabla_V  \left(  \nabla_W (s)  \right)}{.} Jika semua objek berkelas $C^\infty$, derajat keterdiferensialannya tidak lagi perlu diperhatikan. Di sini kita meninjau ekspresi

\mavergleichskettedisp
{\vergleichskette
{ R(V,W)
}
{ =} { \nabla_V \nabla_W - \nabla_W \nabla_V- \nabla_{[V,W]}
}
{ } {
}
{ } {
}
{ } {
}
}
{}{}{,}
yang mengirim penampang
\zusatzklammer {yang cukup terdiferensialkan} {} {}
dari $E$ kembali ke penampang dari $E$. Di sini
\mathl{[V,W]}{} adalah
\definitionsverweis {braket Lie}{}{}
dari kedua medan vektor tersebut.

\inputdefinition
{}
{

Misalkan $M$ adalah
\definitionsverweis {manifold diferensiabel}{}{}
yang dua kali terdiferensialkan secara kontinu, dan
\maabb {} { E } { M
} {}
adalah
\definitionsverweis {bundel vektor diferensiabel}{}{}
yang dua kali terdiferensialkan secara kontinu di atas $M$, dilengkapi dengan
\definitionsverweis {koneksi linear}{}{}
$\nabla$. Untuk
\definitionsverweis {medan vektor}{}{}
diferensiabel
\mathkor {} {V} {dan} {W} {}
pada $M$, pemetaan
\maabbeledisp {} {C^2(E)} {C^0(E)
} { s } { \nabla_V \nabla_W (s) - \nabla_W \nabla_V (s) - \nabla_{[V,W]}(s)
} {,}
disebut
\definitionswort {operator kelengkungan}{}
untuk $V$ dan $W$. Operator ini dinotasikan dengan $R(V,W)$.

}



\inputfaktbeweis
{Reelles Vektorbündel/Trivial/Linearer Zusammenhang/Krümmungsoperator/Beschreibung/Fakt}
{Lema}
{}
{

\faktsituation {Misalkan

\mavergleichskette
{\vergleichskette
{ U
}
{ \subseteq }{ \R^d
}
{  }{
}
{    }{
}
{   }{
}
}
{}{}{}
adalah himpunan
\definitionsverweis {terbuka}{}{}
dan $\nabla$ adalah
\definitionsverweis {koneksi linear}{}{}
pada
\definitionsverweis {bundel vektor}{}{}
trivial

\mavergleichskette
{\vergleichskette
{ E
}
{ = }{ U \times \R^r
}
{  }{
}
{    }{
}
{   }{
}
}
{}{}{.}
Misalkan

\mathbed {\partial_i} {}
 {1 \leq i \leq d} {}
 {} {} {} {}
adalah
\definitionsverweis {medan vektor standar}{}{}
dan

\mathbed {s_j} {}
 {1 \leq j \leq r} {}
 {} {} {} {}
adalah penampang standar di $E$.}
\faktfolgerung {Untuk
\definitionsverweis {operator kelengkungan}{}{}
berlaku

\mavergleichskettedisp
{\vergleichskette
{ R( \partial_i , \partial_j)(s_k)
}
{ =} { \sum_{\ell = 1}^r  R^\ell_{ijk} s_\ell
}
{ } {
}
{ } {
}
{ } {
}
}
{}{}{}
dengan

\mavergleichskettedisp
{\vergleichskette
{ R^\ell_{ijk}
}
{ =} { {\partial_i } \Gamma^\ell_{jk} - {\partial_j } \Gamma^\ell_{ik}  +  \sum_{ a  = 1}^r  \Gamma^a_{jk} \Gamma_{i a}^\ell -\sum_{ a  = 1}^r   \Gamma^a_{ik}  \Gamma_{j a}^\ell
}
{ } {
}
{ } {
}
{ } {
}
}
{}{}{,}
di mana $\Gamma^\ell_{ik}$ menyatakan
\definitionsverweis {simbol Christoffel}{}{}
lokal dari koneksi tersebut.}
\faktzusatz {}
\faktzusatz {}

}
{

Kita mempunyai

\mavergleichskettealignhandlinks
{\vergleichskettealignhandlinks
{ R( \partial_i ,\partial_j)
}
{ =} { \nabla_{\partial_i } \circ \nabla_{\partial_j } -  \nabla_{\partial_j } \circ \nabla_{\partial_i } - \nabla_{[\partial_i, \partial_j]}
}
{ =} { \nabla_{\partial_i } \circ \nabla_{\partial_j } -  \nabla_{\partial_j } \circ \nabla_{\partial_i }
}
{ } {
}
{ } {
}
}
{}
{}{,}
sebab turunan parsial saling komutatif dan karena itu braket Lie-nya sama dengan $0$. Maka, untuk penampang standar $s_k$, dengan menggunakan
Lema 25.10,

\mavergleichskettealignhandlinks
{\vergleichskettealignhandlinks
{ R( \partial_i ,\partial_j) (s_k)
}
{ =} { \nabla_{\partial_i } ( \nabla_{\partial_j } (s_k) ) -  \nabla_{\partial_j } ( \nabla_{\partial_i } (s_k))
}
{ =} { \nabla_{\partial_i }  \left(  \sum_{\ell = 1}^r \Gamma^\ell_{jk} s_\ell  \right)  -  \nabla_{\partial_j }  \left(  \sum_{\ell = 1}^r \Gamma^\ell_{ik} s_\ell  \right)
}
{ =} { \sum_{\ell = 1}^r  \left(  \nabla_{\partial_i }  \left(  \Gamma^\ell_{jk} s_\ell  \right) - \nabla_{\partial_j }  \left(  \Gamma^\ell_{ik} s_\ell  \right)   \right)
}
{ =} { \sum_{\ell = 1}^r  \left(  {\partial_i } \Gamma^\ell_{jk} s_\ell +  \Gamma^\ell_{jk} \sum_{ \rho  = 1}^r \Gamma_{i \ell}^\rho  s_\rho  - {\partial_j } \Gamma^\ell_{ik} s_\ell - \Gamma^\ell_{ik} \sum_{ \rho  = 1}^r \Gamma_{j \ell}^\rho  s_\rho   \right)
}
}
{
\vergleichskettefortsetzungalign
{ =} { \sum_{\ell = 1}^r  \left(  {\partial_i } \Gamma^\ell_{jk} - {\partial_j } \Gamma^\ell_{ik}  \right) s_\ell  + \sum_{\ell = 1}^r \sum_{ \rho  = 1}^r  \left(  \Gamma^\ell_{jk} \Gamma_{i \ell}^\rho  s_\rho  - \Gamma^\ell_{ik}  \Gamma_{j \ell}^\rho  s_\rho  \right)
}
{ =} { \sum_{\ell = 1}^r  \left(  {\partial_i } \Gamma^\ell_{jk} - {\partial_j } \Gamma^\ell_{ik}  \right) s_\ell  + \sum_{\ell = 1}^r  \left(  \sum_{ a  = 1}^r  \Gamma^a_{jk} \Gamma_{i a}^\ell -\sum_{ a  = 1}^r   \Gamma^a_{ik}  \Gamma_{j a}^\ell  \right)  s_\ell
}
{ } {
}
{ } {}
}
{}{.}

}



\inputfaktbeweis
{Reelles Vektorbündel/Rang 1/Offene Menge/Krümmungsoperator/Fakt}
{Akibat}
{}
{

\faktsituation {Misalkan

\mavergleichskette
{\vergleichskette
{ U
}
{ \subseteq }{ \R^n
}
{  }{
}
{    }{
}
{   }{
}
}
{}{}{}
adalah
\definitionsverweis {himpunan terbuka}{}{}
dan $\nabla$ adalah
\definitionsverweis {koneksi linear}{}{}
pada bundel vektor trivial
\mathl{U \times \R}{} yang memiliki
\definitionsverweis {rank}{}{}
$1$, dengan
\definitionsverweis {simbol Christoffel}{}{}

\mathbed {\Gamma_i} {}
 {i = 1 , \ldots , n} {}
 {} {} {} {.}}
\faktuebergang {Pernyataan berikut berlaku bagi
\definitionsverweis {operator kelengkungan}{}{}
terhadap
\definitionsverweis {medan vektor standar}{}{}
$\partial_i$.}
\faktfolgerung {\aufzaehlungzwei {Berlaku

\mavergleichskettedisp
{\vergleichskette
{ R(\partial_i, \partial_j)
}
{ =} { \partial_i \Gamma_j - \partial_j \Gamma_i
}
{ } {
}
{ } {
}
{ } {
}
}
{}{}{.}
} {Persamaan

\mavergleichskette
{\vergleichskette
{ R(\partial_i, \partial_j)
}
{ = }{ 0
}
{  }{
}
{    }{
}
{   }{
}
}
{}{}{}
berlaku untuk semua $i,j$ jika dan hanya jika
$1$-\definitionsverweis {bentuk diferensial}{}{}
$\sum_{i = 1}^n \Gamma_i dx_i$
\definitionsverweis {tertutup}{}{.}
}}
\faktzusatz {}
\faktzusatz {}

}
{

\aufzaehlungzwei {Karena rank bundel sama dengan $1$, kita tuliskan

\mavergleichskette
{\vergleichskette
{ \Gamma_i
}
{ = }{ \Gamma_{i1}^1
}
{  }{
}
{    }{
}
{   }{
}
}
{}{}{}
dan, secara serupa,

\mavergleichskette
{\vergleichskette
{ R_{ij}
}
{ = }{ R_{ij1}^1
}
{  }{
}
{    }{
}
{   }{
}
}
{}{}{,}
sehingga pernyataan tersebut harus dipahami sebagai

\mavergleichskettedisp
{\vergleichskette
{ R(\partial_i, \partial_j) (e)
}
{ =} { \left(  \partial_i \Gamma_j - \partial_j \Gamma_i  \right) e
}
{ } {
}
{ } {
}
{ } {
}
}
{}{}{}
dengan $e$ penampang satuan konstan. Berdasarkan deskripsi eksplisit dalam
Lema 28.2,
kita memperoleh

\mavergleichskettedisp
{\vergleichskette
{ R( \partial_i , \partial_j)
}
{ =} { {\partial_i } \Gamma_{j} - {\partial_j } \Gamma_{i}  + \Gamma_{j} \Gamma_{i } - \Gamma_{i}  \Gamma_{j }
}
{ =} { {\partial_i } \Gamma_{j} - {\partial_j } \Gamma_{i}
}
{ } {
}
{ } {
}
}
{}{}{.}
} {Pernyataan ini mengikuti dari (1) dan
Soal 20.8.
}

}



\inputfaktbeweis
{Reelles Vektorbündel/Linearer Zusammenhang/Krümmungsoperator/Eigenschaften/Fakt}
{Lema}
{}
{

\faktsituation {Misalkan $M$ adalah
\definitionsverweis {manifold diferensiabel}{}{}
yang dua kali terdiferensialkan secara kontinu, dan
\maabb {} { E } { M
} {}
adalah
\definitionsverweis {bundel vektor diferensiabel}{}{}
yang dua kali terdiferensialkan secara kontinu di atas $M$, dilengkapi dengan
\definitionsverweis {koneksi linear}{}{}
$\nabla$.}
\faktuebergang {Maka
\definitionsverweis {operator kelengkungan}{}{}
memiliki sifat-sifat berikut.}
\faktfolgerung {\aufzaehlungdrei{Berlaku

\mavergleichskettedisp
{\vergleichskette
{ R(V,W)
}
{ =} { -R(W,V)
}
{ } {
}
{ } {
}
{ } {
}
}
{}{}{.}
}{ $R$ aditif dalam kedua komponennya.
}{Untuk setiap fungsi diferensiabel
\maabb {f} { M } { \R
} {}
berlaku

\mavergleichskettedisp
{\vergleichskette
{ R(fV,W)
}
{ =} { f R(V,W)
}
{ } {
}
{ } {
}
{ } {
}
}
{}{}{.}
}}
\faktzusatz {}
\faktzusatz {}

}
{

Sifat-sifat tersebut lokal pada manifold. Karena itu, cukup membuktikannya pada suatu tutupan terbuka yang setiap anggotanya difeomorfik dengan himpunan terbuka di $\R^d$ dan di atasnya bundel vektor tersebut trivial. Dengan demikian, (1) mengikuti dari
Lema 28.2.
Pernyataan (2) mengikuti dari
Lema 24.10
dan
Lema 11.8.
Untuk (3), tinjau situasi lokal dengan medan vektor $f \partial_i$ dan $g \partial_j$. Kita mempunyai

\mavergleichskettedisp
{\vergleichskette
{ R(f \partial_i, g \partial_j)
}
{ =} { \nabla_{f \partial_i } \circ \nabla_{g \partial_j } - \nabla_{g\partial_j } \circ \nabla_{f \partial_i } - \nabla_{[ f \partial_i, g \partial_j]}
}
{ } {
}
{ } {
}
{ } {
}
}
{}{}{.}
Menurut
Lema 11.8,

\mavergleichskettedisp
{\vergleichskette
{ [f \partial_i, g\partial_j]
}
{ =} { f [\partial_i, g\partial_j] - g \partial_j(f) \partial_i
}
{ } {
}
{ } {
}
{ } {
}
}
{}{}{.}
Oleh karena itu, untuk suatu penampang $s$, dengan menggunakan
Teorema 25.5,

\mavergleichskettealignhandlinks
{\vergleichskettealignhandlinks
{ R(f \partial_i, g\partial_j) (s)
}
{ =} { \nabla_{f \partial_i } \circ \nabla_{g \partial_j } (s)- \nabla_{g \partial_j } \circ \nabla_{f \partial_i } (s)- \nabla_{[ f \partial_i, g \partial_j]}(s)
}
{ =} { f  \nabla_{ \partial_i } ( \nabla_{g\partial_j } (s) ) - \nabla_{g\partial_j } ( \nabla_{f \partial_i } (s))  - \nabla_{ f [\partial_i, g\partial_j] - g \partial_j(f) \partial_i  } (s)
}
{ =} { f  \nabla_{ \partial_i } ( \nabla_{ g \partial_j } (s) )-\nabla_{ g \partial_j } (f \nabla_{ \partial_i } (s)) - \nabla_{f [\partial_i, g\partial_j] } (s)  + \nabla_{ g \partial_j (f) \partial_i }  (s)
}
{ =} { f  \nabla_{ \partial_i } ( \nabla_{ g \partial_j } (s) )- f \nabla_{ g \partial_j } ( \nabla_{ \partial_i } (s))- g \partial_j(f) \nabla_{\partial_i } (s)  - \nabla_{f [\partial_i, g\partial_j] } (s)          + g \partial_j( f)  \nabla_{\partial_i } (s)
}
}
{
\vergleichskettefortsetzungalign
{ =} { f \nabla_{ \partial_i } ( \nabla_{g \partial_j } (s) )- f \nabla_{ g \partial_j } ( \nabla_{ \partial_i } (s)) - f \nabla_{[\partial_i, g \partial_j]} (s)
}
{ =} { f R( \partial_i, \partial_j) (s)
}
{ } {
}
{ } {
}
}
{}{.}
Bersama (1) dan (2), perhitungan ini membuktikan pernyataan tersebut.

}



\inputfaktbeweis
{Eindimensionale Mannigfaltigkeit/Reelles Vektorbündel/Linearer Zusammenhang/Krümmungsoperator trivial/Fakt}
{Akibat}
{}
{

\faktsituation {Misalkan $M$ adalah
\definitionsverweis {manifold diferensiabel}{}{}
berdimensi satu dan berkelas $C^2$, serta
\maabb {} { E } { M
} {}
adalah
\definitionsverweis {bundel vektor diferensiabel}{}{}
yang dua kali terdiferensialkan secara kontinu di atas $M$, dilengkapi dengan
\definitionsverweis {koneksi linear}{}{}
$\nabla$.}
\faktfolgerung {Maka
\definitionsverweis {operator kelengkungan}{}{}
bersifat trivial untuk sebarang
\definitionsverweis {medan vektor}{}{}
$V,W$.}
\faktzusatz {}
\faktzusatz {}

}
{

Karena pernyataan bersifat lokal, kita dapat langsung mengambil

\mavergleichskette
{\vergleichskette
{ M
}
{ = }{ I
}
{  }{
}
{    }{
}
{   }{
}
}
{}{}{}
dengan suatu interval terbuka, serta

\mavergleichskette
{\vergleichskette
{ V
}
{ = }{ f \partial
}
{  }{
}
{    }{
}
{   }{
}
}
{}{}{,}

\mavergleichskette
{\vergleichskette
{ W
}
{ = }{ g \partial
}
{  }{
}
{    }{
}
{   }{
}
}
{}{}{}
dengan
\definitionsverweis {medan vektor standar}{}{}
$\partial$. Berdasarkan
Lema 28.4  (3),

\mavergleichskettedisp
{\vergleichskette
{ R(f \partial, g \partial)
}
{ =} { fg R( \partial, \partial)
}
{ =} { 0
}
{ } {
}
{ } {
}
}
{}{}{}
karena

\mavergleichskette
{\vergleichskette
{ [\partial, \partial ]
}
{ = }{ 0
}
{  }{
}
{    }{
}
{   }{
}
}
{}{}{.}

}



\inputfaktbeweis
{Reelles Vektorbündel/Trivialer Zusammenhang/Trivialer Krümmungsoperator/Fakt}
{Akibat}
{}
{

\faktsituation {Misalkan $M$ adalah
$C^2$-\definitionsverweis {manifold diferensiabel}{}{}
dan

\mavergleichskette
{\vergleichskette
{ E
}
{ = }{ M \times \R^r
}
{  }{
}
{    }{
}
{   }{
}
}
{}{}{}
adalah
\definitionsverweis {bundel vektor trivial}{}{}
dengan
\definitionsverweis {koneksi trivial}{}{.}}
\faktfolgerung {Maka, untuk sebarang
\definitionsverweis {medan vektor}{}{}
$V,W$,
\definitionsverweis {operator kelengkungan}{}{}
$R(V,W)$ adalah trivial.}
\faktzusatz {}
\faktzusatz {}

}
{

Pernyataan ini bersifat lokal, sehingga kita dapat mengambil

\mavergleichskette
{\vergleichskette
{ M
}
{ = }{ U
}
{ \subseteq }{ \R^n
}
{    }{
}
{   }{
}
}
{}{}{}
dan himpunan tersebut terbuka. Menurut
Lema 28.4  (3),
kita dapat mereduksi lebih lanjut ke kasus ketika medan-medannya adalah
\definitionsverweis {medan vektor standar}{}{}
$\partial_i$. Menurut
Soal 25.11,
semua simbol Christoffel koneksi trivial sama dengan nol. Karena itu, pernyataan mengikuti dari
Lema 28.2.

}

  \zwischenueberschrift{Teorema Frobenius}

Kita tidak dapat membuktikan teorema berikut secara lengkap. Perhatikan bahwa dari simbol Christoffel terhadap penampang basis yang diberikan
\zusatzklammer {secara lokal} {} {}
saja, tidak terlihat apakah simbol-simbol itu akan lenyap terhadap pilihan penampang basis lain. Keadaan ini tepat dicirikan oleh lenyapnya kelengkungan.


\inputfaktbeweis
{Differenzierbare Mannigfaltigkeit/Vektorbündel/Linearer Zusammenhang/Lokal integrabel/Krümmung/Fakt}
{Teorema}
{}
{

\faktsituation {Misalkan $M$ adalah
\definitionsverweis {manifold diferensiabel}{}{}
dan
\maabbdisp {p} { E } { M
} {}
adalah
\definitionsverweis {bundel vektor diferensiabel}{}{,}
yang dilengkapi dengan
\definitionsverweis {koneksi linear}{}{}
$\nabla$.}
\faktuebergang {Pernyataan-pernyataan berikut ekuivalen.}
\faktfolgerung {\aufzaehlungvier{Koneksi tersebut
\definitionsverweis {terintegralkan secara lokal}{}{.}
}{Secara lokal, terhadap
\definitionsverweis {penampang basis}{}{}
yang sesuai, koneksi tersebut isomorfik dengan
\definitionsverweis {koneksi trivial}{}{.}
}{Untuk setiap titik

\mavergleichskette
{\vergleichskette
{ P
}
{ \in }{ M
}
{  }{
}
{    }{
}
{   }{
}
}
{}{}{}
terdapat lingkungan peta terbuka

\mavergleichskette
{\vergleichskette
{ P
}
{ \in }{ U
}
{  }{
}
{    }{
}
{   }{
}
}
{}{}{}
sedemikian sehingga $E \vert_U$ trivial dan
\definitionsverweis {simbol Christoffel}{}{}
yang bersesuaian semuanya merupakan fungsi nol.
}{Untuk sebarang
\definitionsverweis {medan vektor}{}{},
\definitionsverweis {operator kelengkungan}{}{}
adalah trivial.
}}
\faktzusatz {}
\faktzusatz {}

}
{

Semua pernyataan bersifat lokal. Karena itu, kita dapat langsung mulai dengan himpunan terbuka

\mavergleichskette
{\vergleichskette
{ U
}
{ \subseteq }{ \R^n
}
{  }{
}
{    }{
}
{   }{
}
}
{}{}{}
dan bundel vektor trivial $U \times \R^r$. Andaikan (1) berlaku. Kita dapat pula mengandaikan bahwa pada $U$ diberikan keluarga yang terdiri dari $r$
\definitionsverweis {penampang horizontal}{}{}

\mathl{s_1 , \ldots , s_r}{} yang bebas linear. Menurut
Soal 25.6
dan
Soal 25.7,
koneksi itu kemudian dapat ditrivialkan secara lokal. Ini membuktikan (2). Ekuivalensi (2) dan (3) mengikuti dari
Soal 25.11
serta
Catatan 25.8.
Jika simbol Christoffel sama dengan nol, maka penampang standar bersifat horizontal dan secara lokal membentuk penampang basis horizontal. Jadi, ketiga rumusan pertama ekuivalen. Jika (2) berlaku, kita dapat menerapkan
Akibat 28.6
secara lokal dan memperoleh bahwa operator kelengkungan trivial.

Arah dari (4) ke (1) jauh lebih sulit.

}

Untuk kasus-kasus khusus dari arah (4) ke (1), lihat
Soal 28.4
dan
Soal 28.5.


\chapter{Lembar Kerja 28}
\input{generated/worksheet28.id.build.tex}

\section*{Solusi yang disediakan oleh sumber}
\addcontentsline{toc}{section}{Solusi yang disediakan oleh sumber}
\subsection*{Solusi Soal 28.2}
\input{generated/worksheet28_exercise02_solution.id.build.tex}
\subsection*{Solusi Soal 28.5}
\input{generated/worksheet28_exercise05_solution.id.build.tex}

\part{Unit 29}
\chapter[Kuliah 29]{Kuliah 29: Kelengkungan Riemann dan Kelengkungan Seksional}
\setcounter{section}{29}

  \zwischenueberschrift{Kelengkungan pada manifold Riemann}

Misalkan $M$ suatu
\definitionsverweis {manifold Riemann}{}{}
dan
\definitionsverweis {bundel tangen}{}{}
$TM$ dilengkapi dengan
\definitionsverweis {koneksi Levi--Civita}{}{.}
Kita hendak memahami
\definitionsverweis {operator kelengkungan}{}{}
$R$ dalam situasi ini secara lebih mendalam. Untuk medan-medan vektor $V,W$ yang diberikan, operator ini merupakan suatu pemetaan
\maabbdisp {R(V,W)} { C^2(M,TM) } { C^0(M,TM)
} {,}
yakni, ketika diterapkan pada suatu medan vektor $Z$
\zusatzklammer {yang dua kali terdiferensialkan} {} {,}
operator ini menghasilkan medan vektor

\mavergleichskettedisp
{\vergleichskette
{ R(V,W)(Z)
}
{ =} { \nabla_V \nabla_W Z
}
{ } {
}
{ } {
}
{ } {
}
}
{}{}{.}
Dengan sudut pandang ini, pada suatu manifold yang terdiferensialkan tak hingga, keseluruhan konstruksi merupakan pemetaan
\maabbdisp {} {C^\infty(TM) \times C^\infty(TM) \times C^\infty(TM)  } { C^\infty(TM)
} {,}
tetapi ketiga argumennya tidak mempunyai kedudukan yang setara.



\inputfaktbeweis
{R^n/Offene Menge/Riemannsche Metrik/Levi-Civita-Zusammenhang/Krümmungsoperator/Berechnung/Fakt}
{Lemma}
{}
{

\faktsituation {Misalkan

\mavergleichskette
{\vergleichskette
{ U
}
{ \subseteq }{ \R^n
}
{  }{
}
{    }{
}
{   }{
}
}
{}{}{}
\definitionsverweis {terbuka}{}{}
dan dilengkapi dengan suatu
\definitionsverweis {metrik Riemann}{}{}
$g_{ij}$. Misalkan $\nabla$ adalah
\definitionsverweis {koneksi Levi--Civita}{}{}
pada
\definitionsverweis {bundel tangen}{}{}

\mavergleichskette
{\vergleichskette
{ E
}
{ = }{ U \times \R^n
}
{  }{
}
{    }{
}
{   }{
}
}
{}{}{.}}
\faktfolgerung {Maka untuk
\definitionsverweis {operator kelengkungan}{}{}
berlaku

\mavergleichskettedisp
{\vergleichskette
{ R( \partial_i , \partial_j)( \partial_k)
}
{ =} { \sum_{\ell = 1}^n  R^\ell_{ijk} \partial_\ell
}
{ } {
}
{ } {
}
{ } {
}
}
{}{}{}
dengan

\mavergleichskettedisp
{\vergleichskette
{ R^\ell_{ijk}
}
{ =} { {\partial_i } \Gamma^\ell_{jk} - {\partial_j } \Gamma^\ell_{ik}  +  \sum_{ a  = 1}^n  \Gamma^a_{jk} \Gamma_{i a}^\ell -\sum_{ a  = 1}^n   \Gamma^a_{ik}  \Gamma_{j a}^\ell
}
{ } {
}
{ } {
}
{ } {
}
}
{}{}{,}
di mana

\mavergleichskettedisp
{\vergleichskette
{ \Gamma_{ij}^\ell
}
{ =} { { \frac{   1  }{  2 } } \sum_{\mu = 1}^n g^{\ell  \mu }(\partial_ig_{j\mu }+\partial_jg_{i\mu }-\partial_\mu g_{ij})
}
{ } {
}
{ } {
}
{ } {
}
}
{}{}{}
menyatakan
\definitionsverweis {simbol-simbol Christoffel}{}{}
dari koneksi tersebut.}
\faktzusatz {}
\faktzusatz {}

}
{

Ini merupakan kasus khusus dari
Lemma 28.2.

}



\inputfaktbeweis
{Riemannsche Mannigfaltigkeit/Levi-Civita-Zusammenhang/Krümmungsoperator/Eigenschaften/Fakt}
{Lemma}
{}
{

\faktsituation {Misalkan $M$ suatu
\definitionsverweis {manifold Riemann}{}{}
dan
\definitionsverweis {bundel tangen}{}{}
dilengkapi dengan
\definitionsverweis {koneksi Levi--Civita}{}{.}}
\faktuebergang {Maka
\definitionsverweis {operator kelengkungan}{}{}
mempunyai sifat-sifat berikut.}
\faktfolgerung {\aufzaehlungsechs{Ekspresi
\mathl{R(V,W)Z}{} linear dalam ketiga komponennya.
}{Untuk

\mavergleichskette
{\vergleichskette
{ P
}
{ \in }{ M
}
{  }{
}
{    }{
}
{   }{
}
}
{}{}{}
nilai
\mathl{(R(V,W)Z) (P)}{} hanya bergantung pada
\mathl{V(P),W(P),Z(P)}{}.
}{Berlaku

\mavergleichskettedisp
{\vergleichskette
{ R(V,W)
}
{ =} { -R(W,V)
}
{ } {
}
{ } {
}
{ } {
}
}
{}{}{}
}{Berlaku

\mavergleichskettedisp
{\vergleichskette
{ R(V,W)Z+ R(W,Z)V+R(Z,V)W
}
{ =} { 0
}
{ } {
}
{ } {
}
{ } {
}
}
{}{}{.}
}{Berlaku

\mavergleichskettedisp
{\vergleichskette
{ \left\langle R(V,W) T  ,  Z   \right\rangle
}
{ =} { - \left\langle R(V,W) Z  ,  T   \right\rangle
}
{ } {
}
{ } {
}
{ } {
}
}
{}{}{.}
}{Berlaku

\mavergleichskettedisp
{\vergleichskette
{ \left\langle R(V,W) T  ,  Z   \right\rangle
}
{ =} { \left\langle R(T,Z) V  ,  W   \right\rangle
}
{ } {
}
{ } {
}
{ } {
}
}
{}{}{.}
}}
\faktzusatz {}
\faktzusatz {}

}
{

\aufzaehlungsechs{Linearitas dalam $Z$ berasal dari linearitas koneksi. Linearitas dalam $V$ dan $W$ mengikuti dari
Lemma 24.10
dan
Lemma 11.8.
}{Untuk ketergantungan terhadap $V$ dan $W$, pernyataan tersebut mengikuti dari
Lemma 28.4.
Untuk menunjukkan bahwa ketergantungan terhadap $Z$ juga hanya bergantung pada $Z(P)$, kita dapat bekerja dalam situasi lokal pada

\mavergleichskette
{\vergleichskette
{ U
}
{ \subseteq }{ \R^n
}
{  }{
}
{    }{
}
{   }{
}
}
{}{}{}
dan mengambil

\mavergleichskette
{\vergleichskette
{ V
}
{ = }{ \partial_i
}
{  }{
}
{    }{
}
{   }{
}
}
{}{}{,}

\mavergleichskette
{\vergleichskette
{ W
}
{ = }{ \partial_j
}
{  }{
}
{    }{
}
{   }{
}
}
{}{}{}
dan

\mavergleichskette
{\vergleichskette
{ Z
}
{ = }{ h \partial_k
}
{  }{
}
{    }{
}
{   }{
}
}
{}{}{}
dengan suatu fungsi
\maabb {h} { U } {\R
} {}
yang dua kali terdiferensialkan secara kontinu. Menurut
Lemma 25.10,
Teorema 25.5,
dan
Teorema Schwarz,
kita memperoleh

\mavergleichskettealigndrucklinks
{\vergleichskettealigndrucklinks
{ R \left(  \partial_i, \partial_j  \right) \left(  h \partial_k  \right)
}
{ =} { \nabla_{\partial_i }  \left(  \nabla_{\partial_j }  \left(  h \partial_k  \right)  \right) - \nabla_{\partial_j } \left(  \nabla_{\partial_i }  \left(  h \partial_k  \right)  \right)
}
{ =} { \nabla_{\partial_i }  \left(  h \nabla_{\partial_j } \partial_k+ \left(  \partial_j h  \right)  \partial_k  \right) - \nabla_{\partial_j } \left( h \nabla_{\partial_i }  \partial_k  + \left(  \partial_i h  \right) \partial_k  \right)
}
{ =} { h \nabla_{\partial_i } \left(  \nabla_{\partial_j }  \partial_k   \right) +  \left(  \partial_i h  \right) \nabla_{\partial_j } \partial_k + \left(  \partial_j h  \right) \nabla_{\partial_i } \partial_k  +  \left(  \partial_i \partial_j h  \right) \partial_k
-h \nabla_{\partial_j } \left(  \nabla_{\partial_i } \partial_k   \right) -  \left(  \partial_j h  \right) \nabla_{\partial_i } \partial_k - \left(  \partial_i h  \right) \nabla_{\partial_j } \partial_k  -\left(  \partial_j  \partial_i h  \right) \partial_k
}
{ =} { h \nabla_{\partial_i } \left(  \nabla_{\partial_j }  \partial_k   \right)-h \nabla_{\partial_j } \left(  \nabla_{\partial_i } \partial_k   \right)
}
}
{}{}{,}
yang menunjukkan bahwa ekspresi ini hanya bergantung pada $h(P)$.
}{Hal ini langsung terlihat dari definisi dan
Lemma 11.8.
}{Karena
\definitionsverweis {koneksi bebas torsi}{}{}
\zusatzklammer {lihat
Teorema 26.14} {} {}
dan
identitas Jacobi,
berlaku

\mavergleichskettealigndrucklinks
{\vergleichskettealigndrucklinks
{ R(V,W)Z+ R(W,Z)V+R(Z,V)W
}
{ =} { \nabla_V  \left(  \nabla_W Z  \right) -\nabla_W  \left(  \nabla_V Z  \right) - \nabla_{[V,W]} Z +\nabla_W  \left(  \nabla_Z V  \right) -\nabla_Z  \left(  \nabla_W V  \right) - \nabla_{[W,Z]} V+\nabla_Z  \left(  \nabla_V W  \right) -\nabla_V  \left(  \nabla_Z W  \right) - \nabla_{[Z,V]} W
}
{ =} { \nabla_V  \left(  [W,Z ]  \right) + \nabla_W  \left(  [Z, V]  \right) + \nabla_Z  \left(  [V, W]  \right) - \nabla_{[V,W]} Z - \nabla_{[W,Z]} V  - \nabla_{[Z,V]} W
}
{ =} {\nabla_V  \left(  [W,Z ]  \right)  - \nabla_{[W,Z]} V + \nabla_W  \left(  [Z, V]  \right) - \nabla_{[Z,V]} W + \nabla_Z  \left(  [V, W]  \right) - \nabla_{[V,W]} Z
}
{ =} { [V, [W,Z]] + [W, [Z,V]] + [Z, [V,W]]
}
}
{
\vergleichskettefortsetzungalign
{ =} { 0
}
{ } {}
{ } {}
{ } {}
}{}{.}
}{Kita dapat membatasi diri pada medan-medan vektor $V,W$ dengan

\mavergleichskette
{\vergleichskette
{ [V,W]
}
{ = }{ 0
}
{  }{
}
{    }{
}
{   }{
}
}
{}{}{.}
Menurut
Teorema 26.14 (2)
dan
Teorema Schwarz,
kita kemudian memperoleh

\mavergleichskettealigndrucklinks
{\vergleichskettealigndrucklinks
{ \left\langle R(V,W)T , Z  \right\rangle
}
{ =} { \left\langle  \nabla_V \nabla_WT -\nabla_W \nabla_VT   ,  Z  \right\rangle
}
{ =} { \left\langle  \nabla_V \nabla_WT , Z  \right\rangle - \left\langle  \nabla_W \nabla_VT  ,  Z  \right\rangle
}
{ =} { - \left\langle  \nabla_WT , \nabla_V Z  \right\rangle + D_V\left\langle  \nabla_W T  ,  Z  \right\rangle + \left\langle  \nabla_VT  , \nabla_W  Z  \right\rangle -  D_W \left\langle  \nabla_V T  ,  Z  \right\rangle
}
{ =} { \left\langle  T ,  \nabla_W  \nabla_V Z  \right\rangle - D_W \left\langle  T ,  \nabla_VZ   \right\rangle  + D_V \left(  - \left\langle  T  ,  \nabla_WZ  \right\rangle +  D_W \left\langle T  ,  Z   \right\rangle  \right)  - \left\langle  T  ,  \nabla_V \nabla_W  Z  \right\rangle +  D_V \left\langle  T  ,  \nabla_W Z  \right\rangle - D_W  \left(  - \left\langle T  ,  \nabla_V Z   \right\rangle + D_V \left\langle T  ,  Z   \right\rangle  \right)
}
}
{
\vergleichskettefortsetzungalign
{ =} { \left\langle  T ,  \nabla_W  \nabla_V  Z  \right\rangle - \left\langle  T  ,  \nabla_V \nabla_W  Z  \right\rangle
}
{ =} { - \left\langle R(V,W)Z , T  \right\rangle
}
{ } {}
{ } {}
}{}{.}
}{Dengan menggunakan (3), (4), dan (5), diperoleh

\mavergleichskettealigndrucklinks
{\vergleichskettealigndrucklinks
{ \left\langle R(V,W)T , Z  \right\rangle
}
{ =} { - \left\langle R(T,V)W , Z  \right\rangle - \left\langle R(W,T)V , Z  \right\rangle
}
{ =} { \left\langle R(T,V) Z , W  \right\rangle + \left\langle R(W,T) Z ,  V  \right\rangle
}
{ =} { - \left\langle R(V,Z) T , W  \right\rangle - \left\langle R(Z,T) V  ,  W  \right\rangle - \left\langle R(T,Z) W ,  V  \right\rangle - \left\langle R(Z,W) T ,  V  \right\rangle
}
{ =} { 2 \left\langle R(T,Z) V  ,  W  \right\rangle - \left\langle R(V,Z) T , W  \right\rangle - \left\langle R(Z,W) T ,  V  \right\rangle
}
}
{
\vergleichskettefortsetzungalign
{ =} { 2 \left\langle R(T,Z) V  ,  W  \right\rangle + \left\langle R(V,Z) W , T  \right\rangle + \left\langle R(Z,W) V  ,  T  \right\rangle
}
{ =} { 2 \left\langle R(T,Z) V  ,  W  \right\rangle - \left\langle R(W,V) Z , T  \right\rangle
}
{ =} { 2 \left\langle R(T,Z) V  ,  W  \right\rangle - \left\langle R(V,W) T , Z  \right\rangle
}
{ } {}
}{}{.}
Ini membuktikan pernyataan tersebut.
}

}

Berdasarkan
Lemma 29.2 (2),
untuk vektor-vektor

\mavergleichskette
{\vergleichskette
{ v,w,z
}
{ \in }{ T_PM
}
{  }{
}
{    }{
}
{   }{
}
}
{}{}{}
di suatu titik

\mavergleichskette
{\vergleichskette
{ P
}
{ \in }{ M
}
{  }{
}
{    }{
}
{   }{
}
}
{}{}{}
terdefinisi dengan baik suatu vektor

\mavergleichskette
{\vergleichskette
{ R(v,w)(z)
}
{ \in }{ T_PM
}
{  }{
}
{    }{
}
{   }{
}
}
{}{}{,}
sebab vektor-vektor tersebut dapat direalisasikan oleh medan-medan vektor terdiferensialkan
\mathl{V,W,Z}{} dengan

\mavergleichskette
{\vergleichskette
{ V(P)
}
{ = }{ v
}
{  }{
}
{    }{
}
{   }{
}
}
{}{}{}
dan seterusnya
\zusatzklammer {secara lokal} {} {}
dan kemudian $R(V,W)(Z)$ dapat dievaluasi di titik tersebut.

  \zwischenueberschrift{Kelengkungan seksional}

\inputdefinition
{}
{

Misalkan $M$ suatu
\definitionsverweis {manifold Riemann}{}{,}
yang dilengkapi dengan
\definitionsverweis {koneksi Levi--Civita}{}{}
pada
\definitionsverweis {bundel tangen}{}{}
$TM$ dan
\definitionsverweis {operator kelengkungan}{}{}
terkait $R$. Untuk vektor-vektor
\definitionsverweis {bebas linear}{}{}

\mavergleichskette
{\vergleichskette
{ v,w
}
{ \in }{ T_PM
}
{  }{
}
{    }{
}
{   }{
}
}
{}{}{}
di atas suatu titik

\mavergleichskette
{\vergleichskette
{ P
}
{ \in }{ M
}
{  }{
}
{    }{
}
{   }{
}
}
{}{}{}
bilangan

\mavergleichskettedisp
{\vergleichskette
{ K(v,w)
}
{ =} { { \frac{   \left\langle R(v,w)w , v  \right\rangle  }{  \left\langle  v  , v  \right\rangle  \left\langle  w  , w  \right\rangle  -\left\langle  v  , w  \right\rangle^2   } }
}
{ } {
}
{ } {
}
{ } {
}
}
{}{}{}
disebut
\definitionswort {kelengkungan seksional}{}
untuk $v,w$ di $P$.

}

Kelengkungan seksional ini terdefinisi dengan baik karena, dalam kasus kebebasan linear, penyebutnya tidak sama dengan $0$.



\inputfaktbeweis
{Riemannsche Mannigfaltigkeit/Levi-Civita-Zusammenhang/Schnittkrümmung/Tangentiale Ebene/Fakt}
{Lemma}
{}
{

\faktsituation {Misalkan $M$ suatu
\definitionsverweis {manifold Riemann}{}{,}
yang dilengkapi dengan
\definitionsverweis {koneksi Levi--Civita}{}{}
pada
\definitionsverweis {bundel tangen}{}{}
$TM$ dan
\definitionsverweis {operator kelengkungan}{}{}
terkait $R$.}
\faktfolgerung {Maka
\definitionsverweis {kelengkungan seksional}{}{}
untuk dua vektor bebas linear

\mavergleichskette
{\vergleichskette
{ v,w
}
{ \in }{ T_PM
}
{  }{
}
{    }{
}
{   }{
}
}
{}{}{}
hanya bergantung pada bidang yang direntang oleh vektor-vektor tersebut,

\mavergleichskettedisp
{\vergleichskette
{ \R v + \R w
}
{ \subseteq} { T_PM
}
{ } {
}
{ } {
}
{ } {
}
}
{}{}{.}}
\faktzusatz {}
\faktzusatz {}

}
{

Kita meninjau ekspresi pendefinisi

\mavergleichskettedisp
{\vergleichskette
{ K(v,w)
}
{ =} { { \frac{   \left\langle R(v,w)w , v  \right\rangle  }{  \left\langle  v  , v  \right\rangle \left\langle  w  , w  \right\rangle  - \left\langle  v  , w  \right\rangle^2   } }
}
{ } {
}
{ } {
}
{ } {
}
}
{}{}{}
untuk kelengkungan seksional. Menurut
Lemma 29.2,
ekspresi
\mathl{R(v,w)z}{} linear dalam semua komponennya. Jika $v$ atau $w$ dikalikan dengan suatu skalar

\mavergleichskette
{\vergleichskette
{ a
}
{ \neq }{ 0
}
{  }{
}
{    }{
}
{   }{
}
}
{}{}{,}
faktor $a^2$ dapat dikeluarkan baik dari pembilang maupun penyebut. Selanjutnya,

\mavergleichskettealignhandlinks
{\vergleichskettealignhandlinks
{ \left\langle  R(v+w,w) w  , v +w   \right\rangle
}
{ =} { \left\langle  R(v,w)w +R(w,w)w  , v+w   \right\rangle
}
{ =} { \left\langle  R(v,w)w   , v   \right\rangle  +  \left\langle  R(w,w)w  ,  v   \right\rangle + \left\langle  R(v,w)w  ,  w   \right\rangle  +  \left\langle  R(w,w)w   , w   \right\rangle
}
{ } {
}
{ } {
}
}
{}
{}{.}
Menurut
Lemma 29.2 (4),
\mathkor {} {\left\langle  R(w,w)w  ,  v   \right\rangle} {dan} {\left\langle  R(w,w)w   , w   \right\rangle} {}
sama dengan $0$, dan menurut
Lemma 29.2 (5,6),
juga berlaku

\mavergleichskette
{\vergleichskette
{ \left\langle  R(v,w)w  ,  w   \right\rangle
}
{ = }{ 0
}
{  }{
}
{    }{
}
{   }{
}
}
{}{}{.}
Jadi pembilang kelengkungan seksional tidak berubah ketika $v$ diganti dengan $v+w$. Karena

\mavergleichskettealigndrucklinks
{\vergleichskettealigndrucklinks
{ \left\langle v+w , v+w  \right\rangle \left\langle  w  , w  \right\rangle -\left\langle v+w , w  \right\rangle^2
}
{ =} { \left\langle  v  , v  \right\rangle  \left\langle  w  , w  \right\rangle+\left\langle  w  , w  \right\rangle \left\langle  w  , w  \right\rangle +2\left\langle  v  , w  \right\rangle \left\langle  w  , w  \right\rangle -\left\langle v  , w  \right\rangle^2-\left\langle w  , w  \right\rangle^2 -2 \left\langle v  , w  \right\rangle \left\langle w  , w  \right\rangle
}
{ =} { \left\langle  v  , v  \right\rangle \left\langle  w  , w  \right\rangle -\left\langle v  , w  \right\rangle^2
}
{ } {
}
{ } {
}
}
{}{}{}
penyebutnya juga tidak berubah.

}

\inputbemerkung
{}
{

Pada suatu manifold Riemann berdimensi dua, satu-satunya subruang vektor berdimensi dua dalam setiap ruang tangen adalah ruang tangen itu sendiri. Ini berarti bahwa, dalam kasus berdimensi dua, menurut
Lemma 29.4,
argumen-argumen dalam
\definitionsverweis {kelengkungan seksional}{}{}
tidak diperlukan. Karena itu, dalam situasi ini kita memandang kelengkungan seksional sebagai suatu fungsi bernilai riil pada manifold tersebut.

}



\inputfaktbeweis
{Riemannsche Mannigfaltigkeit/Zweidimensional/Levi-Civita-Zusammenhang/Schnittkrümmung/Berechnung/Fakt}
{Lemma}
{}
{

\faktsituation {Untuk suatu
\definitionsverweis {manifold Riemann}{}{}
berdimensi dua}
\faktfolgerung {
\definitionsverweis {kelengkungan seksional}{}{}
pada setiap peta sama dengan
\mathl{-  { \frac{   R^2_{121} }{ g_{11}  } }}{.}}
\faktzusatz {}
\faktzusatz {}

}
{

Kita dapat langsung meninjau situasi pada suatu peta. Menurut
Lemma 29.4,
dalam kasus berdimensi dua kelengkungan seksional dapat dihitung menggunakan basis sembarang; kita memilih
\mathkor {} {\partial_1} {dan} {\partial_2} {.}
Dengan menggunakan
Lemma 29.2 (5),
untuk kelengkungan seksional diperoleh

\mavergleichskettealign
{\vergleichskettealign
{ K( \partial_1, \partial_2)
}
{ =} { { \frac{   \left\langle R(\partial_1 ,\partial_2)\partial_2  ,  \partial_1   \right\rangle  }{  \left\langle  \partial_1  ,  \partial_1   \right\rangle  \left\langle  \partial_2  ,  \partial_2   \right\rangle  -\left\langle  \partial_1  ,  \partial_2   \right\rangle^2   } }
}
{ =} { { \frac{   \left\langle R(\partial_1 ,\partial_2)\partial_2  ,  \partial_1   \right\rangle  }{  g_{11} g_{22}  -g_{12}^2   } }
}
{ =} { -  { \frac{   \left\langle R(\partial_1 ,\partial_2)\partial_1  ,  \partial_2   \right\rangle  }{  g_{11} g_{22}  -g_{12}^2   } }
}
{ } {
}
}
{}
{}{}
Menurut
Lemma 29.1,
berlaku

\mavergleichskettedisp
{\vergleichskette
{ R( \partial_1 , \partial_2)( \partial_1)
}
{ =} { R^1_{121} \partial_1+ R^2_{121} \partial_2
}
{ } {
}
{ } {
}
{ } {
}
}
{}{}{.}
Hal ini memberikan

\mavergleichskettealign
{\vergleichskettealign
{ K( \partial_1, \partial_2)
}
{ =} { - { \frac{   \left\langle R(\partial_1 ,\partial_2)\partial_1  ,  \partial_2   \right\rangle  }{  g_{11} g_{22}  -g_{12}^2   } }
}
{ =} { - { \frac{   R^1_{121} g_{12}+ R^2_{121} g_{22}  }{  g_{11} g_{22}  -g_{12}^2   } }
}
{ } {
}
{ } {
}
}
{}
{}{.}
Menurut
Lemma 29.2 (4),

\mavergleichskettedisp
{\vergleichskette
{ \left\langle  R(\partial_1, \partial_1) \partial_1  ,  \partial_2   \right\rangle
}
{ =} { R^1_{121}  g_{11}  + R_{121}^2 g_{12}
}
{ =} { 0
}
{ } {
}
{ } {
}
}
{}{}{.}
Kita mengalikan pembilang kelengkungan seksional dengan $g_{11}$ dan memperoleh

\mavergleichskettealignhandlinks
{\vergleichskettealignhandlinks
{ g_{11}  \left(  R^1_{121} g_{12}+ R^2_{121} g_{22}  \right)
}
{ =} { g_{11}  \left(  R^1_{121} g_{12}+ R^2_{121} g_{22}  \right) - g_{12}  \left(  R^1_{121}  g_{11}  + R_{121}^2 g_{12}   \right)
}
{ =} { g_{11} g_{22} R^2_{121} - g_{12}^2 R^2_{121}
}
{ =} { \left(  g_{11} g_{22}  - g_{12}^2  \right) R^2_{121}
}
{ } {
}
}
{}
{}{.}
Jadi,

\mavergleichskettedisp
{\vergleichskette
{ K( \partial_1, \partial_2)
}
{ =} { -  { \frac{   R^2_{121} }{ g_{11}  } }
}
{ } {
}
{ } {
}
{ } {
}
}
{}{}{.}

}



\inputfaktbeweis
{Fläche im Raum/Schnittkrümmung/Gaußkrümmung/Fakt}
{Lemma}
{}
{

\faktsituation {Misalkan

\mavergleichskette
{\vergleichskette
{ W
}
{ \subseteq }{ \R^3
}
{  }{
}
{    }{
}
{   }{
}
}
{}{}{}
\definitionsverweis {terbuka}{}{}
dan misalkan
\maabb {h} { W } { \R
} {}
suatu
\definitionsverweis {fungsi yang dua kali terdiferensialkan secara kontinu}{}{}
serta

\mavergleichskette
{\vergleichskette
{ Y
}
{ = }{ h^{-1}(c)
}
{  }{
}
{    }{
}
{   }{
}
}
{}{}{}
serat di atas

\mavergleichskette
{\vergleichskette
{ c
}
{ \in }{ \R
}
{  }{
}
{    }{
}
{   }{
}
}
{}{}{,}
dan andaikan $h$
\definitionsverweis {reguler}{}{}
di setiap titik $Y$.}
\faktfolgerung {Maka
\definitionsverweis {kelengkungan Gauss}{}{}
dari $Y$ sama dengan
\definitionsverweis {kelengkungan seksional}{}{}
dari $Y$.}
\faktzusatz {}
\faktzusatz {}

}
{

Menurut
Lemma 29.6,
kelengkungan seksional pada suatu peta sama dengan
\mathl{- { \frac{  R^2_{121}  }{ g_{11}  } }}{.} Karena

\mavergleichskettedisphandlinks
{\vergleichskettedisphandlinks
{ R^2_{121}
}
{ =} { {\partial_1 } \Gamma^2_{12} - {\partial_2 } \Gamma^2_{11} +  \Gamma^1_{12} \Gamma_{1 1}^2 + \Gamma^2_{12} \Gamma_{1 2}^2 - \Gamma^1_{11}  \Gamma_{12 }^2  - \Gamma^2_{11} \Gamma_{2 2}^2
}
{ =} { -g_{11} K
}
{ } {
}
{ } {
}
}
{}{}{.}
maka

\mavergleichskettedisphandlinks
{\vergleichskettedisphandlinks
{ K
}
{ =} { { \frac{   1  }{ g_{11}  } }  \left(  \partial_2 \Gamma_{11}^2 - \partial_1 \Gamma_{12}^2 + \Gamma^1_{11} \Gamma^2_{12}  + \Gamma_{11}^2 \Gamma_{22}^2 - \Gamma_{12}^1 \Gamma_{11}^2-  \left(  \Gamma_{12}^2  \right)^2  \right)
}
{ } {
}
{ } {
}
{ } {
}
}
{}{}{}
yang berimpit dengan kelengkungan Gauss menurut
Teorema 19.7.

}

\inputbeispiel{}
{

Untuk
\definitionsverweis {bidang Euklides}{}{}
$\R^2$ dengan
\definitionsverweis {hasil kali skalar standar}{}{,}
\definitionsverweis {kelengkungan seksional}{}{}
konstan sama dengan $0$.

}

\inputbeispiel{}
{

Untuk
\definitionsverweis {bola satuan}{}{}

\mavergleichskette
{\vergleichskette
{ S^2
}
{ \subseteq }{ \R^3
}
{  }{
}
{    }{
}
{   }{
}
}
{}{}{}
dengan
\definitionsverweis {struktur Riemann}{}{}
terinduksi,
\definitionsverweis {kelengkungan seksional}{}{}
konstan sama dengan $1$. Hal ini mengikuti dari
Contoh 5.6
bersama dengan
Lemma 29.7.

}

\inputbeispiel{}
{

\bild{ \begin{center}
\includegraphics[width=5.5cm]{ \bildeinlesungsvg {Parallel lines in Poincare's model of hyperbolic geometry} {svg} }
\end{center}
\bildtext {} }

\bildlizenz { Parallel lines in Poincare's model of hyperbolic geometry.svg } {} {Januszkaja} {Commons} {CC-by-sa 3.0} {}

Kita melanjutkan
Contoh 18.8
dan
Contoh 26.7.
Menurut
Lemma 29.1,

\mavergleichskettedisphandlinks
{\vergleichskettedisphandlinks
{ R^2_{121}
}
{ =} { {\partial_1 } \Gamma^2_{12} - {\partial_2 } \Gamma^2_{11} + \Gamma^1_{12} \Gamma_{1 1}^2 + \Gamma^2_{12} \Gamma_{1 2}^2 - \Gamma^1_{11} \Gamma_{12 }^2 - \Gamma^2_{11} \Gamma_{2 2}^2
}
{ =} { { \frac{   1  }{  y^2 } } - { \frac{   1  }{  y^2 } }+  { \frac{   1  }{  y^2 } }
}
{ =} { { \frac{   1  }{  y^2 } }
}
{ } {
}
}
{}{}{.}
Menurut
Lemma 29.6,
untuk
\definitionsverweis {kelengkungan seksional}{}{}
berlaku

\mavergleichskettedisp
{\vergleichskette
{ K
}
{ =} { - { \frac{   R^2_{121} }{ g_{11 } }}
}
{ =} { - { \frac{   { \frac{   1  }{  y^2 } }  }{  { \frac{   1  }{  y^2 } }  } }
}
{ =} { -1
}
{ } {
}
}
{}{}{,}
jadi kita memperoleh kelengkungan seksional konstan $-1$.

}


\chapter{Lembar Kerja 29}
\setcounter{section}{29}

  \zwischenueberschrift{Soal latihan}

\inputaufgabe
{}
{

Misalkan $M$ suatu
\definitionsverweis {manifold diferensiabel}{}{}
sedemikian sehingga
\definitionsverweis {bundel tangen}{}{}
$TM$ bersifat trivial. Melalui suatu trivialisasi diferensiabel

\mavergleichskette
{\vergleichskette
{ TM
}
{ \cong }{ M \times \R^n
}
{  }{
}
{    }{
}
{   }{
}
}
{}{}{}
lengkapi $M$
\zusatzklammer {dengan menggunakan
\definitionsverweis {hasil kali dalam standar}{}{}
pada $\R^n$} {} {}
dengan suatu
\definitionsverweis {struktur Riemann}{}{}
\zusatzklammer {lihat
Soal 16.2} {} {.}
Tunjukkan bahwa
\definitionsverweis {operator kelengkungan}{}{}
pada $M$ bersifat trivial.

}
{} {}

\inputaufgabegibtloesung
{}
{

Dalam situasi dua dimensi pada
Lemma 29.1,
hitung simbol-simbol Christoffel yang relevan untuk
\definitionsverweis {operator kelengkungan}{}{.}

}
{} {}

\inputaufgabe
{}
{

Kita tinjau
\definitionsverweis {torus}{}{}
$S^1 \times S^1$. Pertama, lengkapi torus tersebut dengan struktur yang diberikan oleh struktur produk
\zusatzklammer {bersama penyematan

\mavergleichskette
{\vergleichskette
{ S^1
}
{ \subseteq }{ \R^2
}
{  }{
}
{    }{
}
{   }{
}
}
{}{}{,}
lihat
Soal 16.3
dan
Soal 16.4} {} {}
dan
\definitionsverweis {struktur Riemann}{}{,}
kemudian bandingkan dengan realisasi tertanam berikut.

\mavergleichskettedisp
{\vergleichskette
{ T
}
{ =} { \left\{ (x,y,z) \in \R^3  \mid   \left(  \sqrt{x^2+y^2} -R  \right)^2 +z^2  = r^2         \right\}
}
{ \subseteq} { \R^3
}
{ } {
}
{ } {
}
}
{}{}{}
dengan struktur Riemann terinduksi. Bandingkan
\definitionsverweis {operator-operator kelengkungan}{}{}
untuk kedua struktur tersebut.

}
{} {}

  \zwischenueberschrift{Soal untuk dikumpulkan}


\section*{Solusi yang disediakan oleh sumber}
\addcontentsline{toc}{section}{Solusi yang disediakan oleh sumber}
\subsection*{Solusi Soal 29.2}
Berdasarkan berbagai identitas pada
Fakta,
hanya simbol-simbol Christoffel berikut yang mungkin relevan:

\mathdisp {R^1_{112}, R^2_{112},R^1_{122},R^2_{122},} {  }

Menurut
Fakta,

\mavergleichskettedisp
{\vergleichskette
{ R^1_{112}
}
{ =} { {\partial_1} \Gamma^1_{12} - {\partial_1} \Gamma^1_{12}  +  \Gamma^1_{12} \Gamma_{11}^1  +  \Gamma^2_{12} \Gamma_{12}^1        -     \Gamma^1_{12}  \Gamma_{11}^1        -     \Gamma^2_{12}  \Gamma_{12}^1
}
{ =} {  0
}
{ } {
}
{ } {
}
}
{}{}{,}

\mavergleichskettedisp
{\vergleichskette
{ R^2_{112}
}
{ =} { {\partial_1} \Gamma^2_{12} - {\partial_1} \Gamma^2_{12} +  \Gamma^1_{12} \Gamma_{11}^2 +\Gamma^2_{12} \Gamma_{12}^2 - \Gamma^1_{12}  \Gamma_{11}^2 - \Gamma^2_{12}  \Gamma_{12}^2
}
{ =} {  0
}
{ } {
}
{ } {
}
}
{}{}{,}

\mavergleichskettedisp
{\vergleichskette
{ R^1_{121}
}
{ =} { - {\partial_1} \Gamma^1_{12} - {\partial_2} \Gamma^1_{11} -  \Gamma^1_{12} \Gamma_{1 1}^1 - \Gamma^2_{12} \Gamma_{1 2}^1 +  \Gamma^1_{11}  \Gamma_{12 }^1  - \Gamma^2_{11} \Gamma_{2 2}^1
}
{ } {
}
{ } {
}
{ } {
}
}
{}{}{,}

\mavergleichskettedisp
{\vergleichskette
{ R^2_{121}
}
{ =} {  {\partial_1} \Gamma^2_{12} - {\partial_2} \Gamma^2_{11} +  \Gamma^1_{12} \Gamma_{1 1}^2 + \Gamma^2_{12} \Gamma_{1 2}^2 - \Gamma^1_{11}  \Gamma_{12 }^2  - \Gamma^2_{11} \Gamma_{2 2}^2
}
{ } {
}
{ } {
}
{ } {
}
}
{}{}{.}

\mavergleichskettedisp
{\vergleichskette
{ R^1_{122}
}
{ =} { {\partial_1} \Gamma^1_{22} - {\partial_2} \Gamma^1_{22}  +   \Gamma^1_{22} \Gamma_{1 1}^1   +   \Gamma^2_{22} \Gamma_{1 2}^1 -  \Gamma^1_{12}  \Gamma_{12 }^1  - \Gamma^2_{12} \Gamma_{2 2}^1
}
{ } {
}
{ } {
}
{ } {
}
}
{}{}{,}

\mavergleichskettedisp
{\vergleichskette
{ R^2_{122}
}
{ =} { {\partial_1} \Gamma^2_{22} - {\partial_2} \Gamma^2_{22}  +   \Gamma^1_{22} \Gamma_{1 1}^2  -  \Gamma^1_{12}  \Gamma_{12 }^2
}
{ } {
}
{ } {
}
{ } {
}
}
{}{}{.}


\part{Jembatan Asli}
\chapter{Jembatan Grup Lie dan Aljabar Lie}
\noindent\textit{Materi asli edisi ini; bukan solusi atau teks yang disediakan oleh sumber Brenner/Wikiversity.}
% O011-BL — modul asli, bukan terjemahan atau adaptasi dari teks pembanding.
% Judul: Jembatan Grup Lie dan Aljabar Lie
% Lisensi modul: CC BY-SA 4.0.
% Provenans: disusun oleh OpenAI Codex gpt-5.6-sol, Ultra, atas arahan pengguna.
% Prasyarat internal: Unit 7, 9, 10, 11, dan 28 edisi ini.

\setcounter{section}{30}
\section{Jembatan Grup Lie dan Aljabar Lie}
\label{o011-bridge-lie}

\subsection{Tujuan dan hubungan dengan materi sebelumnya}

Modul ini menghubungkan bahasa manifold dengan simetri kontinu. Unit 7
menyediakan manifold mulus, Unit 9 ruang singgung, Unit 10 pemetaan mulus dan
medan vektor, Unit 11 produk manifold, sedangkan Unit 28 menyediakan braket
Lie medan vektor. Semua hasil di bawah dibangun dari perangkat itu. Tujuan
utamanya ialah memahami bagaimana operasi grup yang mulus menghasilkan suatu
aljabar linear pada ruang singgung di unsur identitas, bagaimana eksponensial
menghubungkan kedua struktur tersebut, dan bagaimana aksi grup menghasilkan
orbit serta stabilisator.

Sepanjang modul, semua manifold dianggap berdimensi hingga, Hausdorff, dan
memiliki basis terhitung. Kata ``mulus'' berarti terdiferensialkan tak hingga.
Jika $f\colon M\to N$ mulus dan $p\in M$, diferensialnya ditulis
$T_pf\colon T_pM\to T_{f(p)}N$.

\subsection{Grup Lie dan contoh matriks}

\begin{definition}\label{o011-bl-def-lie-group}
Suatu \emph{grup Lie} ialah sebuah manifold mulus $G$ yang sekaligus merupakan
grup, sedemikian sehingga pemetaan perkalian dan invers
\[
  m\colon G\times G\longrightarrow G,\qquad m(g,h)=gh,
  \qquad
  \iota\colon G\longrightarrow G,\qquad \iota(g)=g^{-1},
\]
bersifat mulus. Unsur identitasnya ditulis $e$. Homomorfisme grup Lie ialah
homomorfisme grup yang juga merupakan pemetaan mulus.
\end{definition}

Syarat tersebut bukan sekadar hiasan topologis. Ia memungkinkan kita
mendiferensialkan identitas grup. Sebagai contoh, dari $gg^{-1}=e$ kita akan
memperoleh rumus diferensial untuk invers; dari keasosiatifan kita akan
memperoleh medan vektor invarian dan braket pada $T_eG$.

\begin{beispiel}\label{o011-bl-ex-vector-groups}
Ruang vektor nyata $\mathbb R^n$ dengan penjumlahan merupakan grup Lie. Unsur
identitasnya $0$, inversnya $x\mapsto -x$, dan kedua operasi itu polinomial.
Torus
\[
  \mathbb T^n=(S^1)^n
\]
dengan perkalian komponen demi komponen juga merupakan grup Lie. Pemetaan
\[
  \mathbb R^n\longrightarrow\mathbb T^n,
  \qquad (t_1,\ldots,t_n)\longmapsto
  (e^{it_1},\ldots,e^{it_n})
\]
ialah homomorfisme grup Lie surjektif dengan kernel $(2\pi\mathbb Z)^n$.
\end{beispiel}

\begin{beispiel}\label{o011-bl-ex-general-linear}
Di dalam ruang semua matriks nyata $M_n(\mathbb R)\cong\mathbb R^{n^2}$,
himpunan
\[
  \mathrm{GL}_n(\mathbb R)=\{A\in M_n(\mathbb R):\det A\ne0\}
\]
terbuka karena determinan kontinu. Jadi ia merupakan manifold. Perkalian
matriks bersifat polinomial, sedangkan
\[
  A^{-1}=\frac{1}{\det A}\operatorname{adj}(A)
\]
mulus pada himpunan tempat $\det A\ne0$. Dengan demikian
$\mathrm{GL}_n(\mathbb R)$ ialah grup Lie.
\end{beispiel}

\begin{beispiel}\label{o011-bl-ex-matrix-subgroups}
Contoh-contoh penting berikut merupakan subgrup tertutup
$\mathrm{GL}_n(\mathbb R)$ dan, dengan struktur manifold terinduksi yang
sesuai, merupakan grup Lie:
\[
\begin{aligned}
 \mathrm{SL}_n(\mathbb R)&=\{A:\det A=1\},\\
 \mathrm O(n)&=\{A:A^{\mathsf T}A=I\},\\
 \mathrm{SO}(n)&=\{A:A^{\mathsf T}A=I,\ \det A=1\}.
\end{aligned}
\]
Untuk $\mathrm{SL}_n$, nilai $1$ merupakan nilai regular determinan karena
$D(\det)_A(H)=\det(A)\operatorname{tr}(A^{-1}H)$ tidak nol sebagai fungsional.
Untuk $\mathrm O(n)$, persamaan $A^{\mathsf T}A=I$ dapat diperlakukan dengan
teorema nilai regular ke ruang matriks simetris. Kita tidak memerlukan teorema
umum bahwa setiap subgrup tertutup dari grup Lie adalah submanifold, meskipun
teorema itu memang benar.
\end{beispiel}

\begin{bemerkung}\label{o011-bl-rem-components}
$\mathrm{GL}_n(\mathbb R)$ mempunyai dua komponen terhubung, dibedakan oleh
tanda determinan. $\mathrm O(n)$ juga mempunyai dua komponen, sedangkan
$\mathrm{SO}(n)$ adalah komponen identitasnya. Aljabar Lie hanya merekam
geometri infinitesimal di sekitar $e$; karena itu ia tidak dengan sendirinya
merekam semua komponen terhubung suatu grup.
\end{bemerkung}

\subsection{Translasi kiri dan medan vektor invarian}

Untuk $g\in G$, definisikan translasi kiri dan kanan
\[
  L_g(h)=gh,
  \qquad
  R_g(h)=hg.
\]
Keduanya difeomorfisme, dengan invers masing-masing $L_{g^{-1}}$ dan
$R_{g^{-1}}$.

\begin{definition}\label{o011-bl-def-left-invariant}
Suatu medan vektor mulus $X$ pada $G$ disebut \emph{invarian-kiri} apabila
\[
  T_hL_g(X_h)=X_{gh}
\]
untuk semua $g,h\in G$.
\end{definition}

\begin{Proposition}\label{o011-bl-prop-left-invariant-correspondence}
Untuk setiap $v\in T_eG$ terdapat tepat satu medan vektor invarian-kiri
$X^v$, yaitu
\[
  X^v_g=T_eL_g(v).
\]
Pemetaan $v\mapsto X^v$ merupakan isomorfisme linear dari $T_eG$ ke ruang
medan vektor invarian-kiri.
\end{Proposition}

\begin{proof}
Keasosiatifan memberi $L_g\circ L_h=L_{gh}$. Maka
\[
 T_hL_g(X^v_h)
 =T_hL_g\bigl(T_eL_h(v)\bigr)
 =T_e(L_g\circ L_h)(v)
 =T_eL_{gh}(v)=X^v_{gh}.
\]
Sebaliknya, jika $X$ invarian-kiri, ambil $h=e$ untuk memperoleh
$X_g=T_eL_g(X_e)$. Jadi $X$ ditentukan secara unik oleh nilainya di $e$.
Linearitas jelas dari linearitas diferensial.
\end{proof}

\begin{Proposition}\label{o011-bl-prop-invariant-bracket}
Jika $X$ dan $Y$ invarian-kiri, maka braket medan vektor $[X,Y]$ juga
invarian-kiri.
\end{Proposition}

\begin{proof}
Untuk setiap difeomorfisme $F$, sifat natural braket menyatakan
$F_*[X,Y]=[F_*X,F_*Y]$. Terapkan ini pada $F=L_g$. Karena
$(L_g)_*X=X$ dan $(L_g)_*Y=Y$, diperoleh $(L_g)_*[X,Y]=[X,Y]$.
\end{proof}

\subsection{Aljabar Lie di unsur identitas}

\begin{definition}\label{o011-bl-def-lie-algebra}
\emph{Aljabar Lie} suatu grup Lie $G$ ialah ruang vektor
\[
  \mathfrak g=T_eG
\]
dengan operasi bilinear
\[
  [v,w]_{\mathfrak g}=[X^v,X^w]_e.
\]
Jika konteksnya jelas, subskrip $\mathfrak g$ dihilangkan.
\end{definition}

Proposisi sebelumnya menjamin bahwa $[X^v,X^w]$ kembali invarian-kiri,
sehingga nilainya di $e$ menentukan seluruh medan. Sifat braket medan vektor
langsung memberi
\[
 [v,w]=-[w,v],
 \qquad
 [u,[v,w]]+[v,[w,u]]+[w,[u,v]]=0.
\]
Identitas kedua disebut identitas Jacobi.

\begin{Proposition}\label{o011-bl-prop-homomorphism-differential}
Jika $F\colon G\to H$ homomorfisme grup Lie, maka
\[
  dF_e=T_eF\colon\mathfrak g\longrightarrow\mathfrak h
\]
mempertahankan braket Lie.
\end{Proposition}

\begin{proof}
Identitas $F\circ L_g=L_{F(g)}\circ F$ menunjukkan bahwa medan
invarian-kiri $X^v$ dan $X^{dF_e(v)}$ saling terkait oleh $F$. Naturalitas
braket untuk medan yang saling terkait menghasilkan
\[
  dF_e([v,w])=[dF_e(v),dF_e(w)].
\]
\end{proof}

\begin{beispiel}\label{o011-bl-ex-matrix-algebras}
Untuk grup matriks, kita mengidentifikasi ruang singgung di $I$ dengan
subruang matriks. Hasilnya ialah
\[
\begin{aligned}
 \mathfrak{gl}_n(\mathbb R)&=M_n(\mathbb R),\\
 \mathfrak{sl}_n(\mathbb R)&=\{X:\operatorname{tr}X=0\},\\
 \mathfrak o(n)=\mathfrak{so}(n)&=\{X:X^{\mathsf T}+X=0\}.
\end{aligned}
\]
Braketnya adalah komutator
\[
  [X,Y]=XY-YX.
\]
Sebagai contoh, diferensiasikan $A(t)^{\mathsf T}A(t)=I$ pada $t=0$ dengan
$A(0)=I$ untuk memperoleh $A'(0)^{\mathsf T}+A'(0)=0$.
\end{beispiel}

\begin{bemerkung}\label{o011-bl-rem-sign}
Konvensi braket di atas memakai medan invarian-kiri. Jika seseorang
mengidentifikasi $T_eG$ dengan medan invarian-kanan tanpa mengubah definisi,
tanda braket akan berbalik. Menyatakan konvensi ini mencegah kesalahan tanda
pada rumus adjoint.
\end{bemerkung}

\subsection{Subgrup satu-parameter dan pemetaan eksponensial}

\begin{definition}\label{o011-bl-def-one-parameter}
Subgrup satu-parameter dalam $G$ ialah homomorfisme grup Lie
\[
  \gamma\colon(\mathbb R,+)\longrightarrow G.
\]
\end{definition}

Jika $\gamma$ subgrup satu-parameter, maka
\[
  \gamma(t+s)=\gamma(t)\gamma(s)
\]
dan, setelah mendiferensialkan terhadap $s$ pada $0$,
\[
  \gamma'(t)=T_eL_{\gamma(t)}\bigl(\gamma'(0)\bigr).
\]
Jadi $\gamma$ adalah kurva integral medan invarian-kiri yang ditentukan oleh
$v=\gamma'(0)$.

\begin{Theorem}\label{o011-bl-thm-one-parameter-existence}
Untuk setiap $v\in\mathfrak g$ terdapat tepat satu subgrup satu-parameter
$\gamma_v$ dengan $\gamma_v'(0)=v$.
\end{Theorem}

\begin{proof}
Ambil kurva integral maksimal $\gamma$ dari medan invarian-kiri $X^v$ dengan
$\gamma(0)=e$. Untuk $s$ tetap, kurva $t\mapsto\gamma(s)\gamma(t)$ dan
$t\mapsto\gamma(s+t)$ memenuhi persamaan diferensial yang sama serta memiliki
nilai awal yang sama pada $t=0$. Ketunggalan solusi memberi
$\gamma(s+t)=\gamma(s)\gamma(t)$ selama kedua ruas terdefinisi.

Identitas ini memperpanjang kurva integral lokal langkah demi langkah ke
seluruh $\mathbb R$: pilih $\varepsilon>0$ dengan solusi pada
$(-\varepsilon,\varepsilon)$, lalu definisikan nilai pada selang yang lebih
panjang dengan hasil kali potongan-potongan kecil. Identitas lokal dan
ketunggalan menjamin bahwa definisi itu konsisten. Kurva global yang diperoleh
adalah homomorfisme. Ketunggalan subgrup satu-parameter juga mengikuti dari
ketunggalan kurva integral.
\end{proof}

\begin{definition}\label{o011-bl-def-exponential}
Pemetaan eksponensial grup Lie adalah
\[
  \exp_G\colon\mathfrak g\longrightarrow G,
  \qquad \exp_G(v)=\gamma_v(1).
\]
\end{definition}

Dari definisi dan penskalaan parameter diperoleh
\[
  \gamma_v(t)=\exp_G(tv),
  \qquad
  \exp_G((s+t)v)=\exp_G(sv)\exp_G(tv).
\]

\begin{Proposition}\label{o011-bl-prop-exp-local}
Pemetaan $\exp_G$ mulus, $\exp_G(0)=e$, dan diferensialnya di $0$ adalah
identitas pada $\mathfrak g$. Karena itu $\exp_G$ merupakan difeomorfisme dari
suatu lingkungan $0$ ke suatu lingkungan $e$.
\end{Proposition}

\begin{proof}
Kebergantungan mulus solusi persamaan diferensial pada nilai awal dan parameter
memberi kemulusan $\exp_G$. Kurva $t\mapsto\exp_G(tv)$ mempunyai kecepatan $v$
pada $0$, sehingga
\[
  T_0\exp_G(v)=\left.\frac{d}{dt}\right|_{0}\exp_G(tv)=v.
\]
Teorema fungsi invers menyelesaikan pernyataan terakhir.
\end{proof}

\begin{beispiel}\label{o011-bl-ex-matrix-exp}
Untuk $G=\mathrm{GL}_n(\mathbb R)$,
\[
  \exp_G(X)=e^X=\sum_{k=0}^{\infty}\frac{X^k}{k!}.
\]
Deret ini konvergen mutlak dalam setiap norma matriks. Turunan
$t\mapsto e^{tX}$ ialah $Xe^{tX}=e^{tX}X$, sehingga kurva tersebut merupakan
subgrup satu-parameter dengan kecepatan awal $X$. Jika $X$ dan $Y$ komutatif,
$e^{X+Y}=e^Xe^Y$; tanpa hipotesis komutatif, identitas ini umumnya salah.
\end{beispiel}

\begin{Proposition}\label{o011-bl-prop-exp-natural}
Untuk homomorfisme grup Lie $F\colon G\to H$ berlaku
\[
  F(\exp_G v)=\exp_H(dF_e v).
\]
\end{Proposition}

\begin{proof}
Kedua ruas, setelah $v$ diganti $tv$, merupakan subgrup satu-parameter di $H$
dengan kecepatan awal $dF_e v$. Gunakan ketunggalan pada Teorema
\ref{o011-bl-thm-one-parameter-existence} dan ambil $t=1$.
\end{proof}

\subsection{Konjugasi dan representasi adjoint}

Untuk $g\in G$, konjugasi oleh $g$ adalah difeomorfisme
\[
  C_g\colon G\longrightarrow G,
  \qquad C_g(h)=ghg^{-1}.
\]
Ia mempertahankan $e$.

\begin{definition}\label{o011-bl-def-adjoint}
Pemetaan adjoint pada tingkat grup adalah
\[
  \operatorname{Ad}\colon G\longrightarrow\mathrm{GL}(\mathfrak g),
  \qquad \operatorname{Ad}_g=T_eC_g.
\]
Diferensialnya di identitas adalah
\[
  \operatorname{ad}=T_e\operatorname{Ad}\colon
  \mathfrak g\longrightarrow\mathfrak{gl}(\mathfrak g).
\]
\end{definition}

Karena $C_{gh}=C_g\circ C_h$, pemetaan $\operatorname{Ad}$ merupakan
homomorfisme grup Lie. Untuk grup matriks,
\[
  \operatorname{Ad}_g(X)=gXg^{-1}.
\]

\begin{Theorem}\label{o011-bl-thm-ad-bracket}
Untuk $v,w\in\mathfrak g$ berlaku
\[
  \operatorname{ad}_v(w)=[v,w].
\]
Selain itu,
\[
  \operatorname{Ad}_g[v,w]
  =[\operatorname{Ad}_gv,\operatorname{Ad}_gw]
\]
dan
\[
  \operatorname{Ad}_{\exp(tv)}
  =\exp\bigl(t\operatorname{ad}_v\bigr).
\]
\end{Theorem}

\begin{proof}
Konjugasi membawa medan invarian-kiri $X^w$ ke medan invarian-kiri yang
bernilai $\operatorname{Ad}_g w$ di identitas. Naturalitas braket memberi
identitas kedua. Untuk identitas pertama, diferensiasikan keluarga
$C_{\exp(tv)}$ pada $t=0$; turunan aksi keluarga difeomorfisme ini pada medan
invarian-kiri adalah braket dengan $X^v$. Evaluasi di $e$ memberi
$\operatorname{ad}_v(w)=[v,w]$.

Terakhir, $t\mapsto\operatorname{Ad}_{\exp(tv)}$ ialah subgrup satu-parameter
di $\mathrm{GL}(\mathfrak g)$ dengan kecepatan awal
$\operatorname{ad}_v$. Contoh \ref{o011-bl-ex-matrix-exp} pada ruang
$\mathfrak g$ memberi identitas ketiga.
\end{proof}

\subsection{Aksi mulus, orbit, dan stabilisator}

\begin{definition}\label{o011-bl-def-action}
Aksi kiri mulus grup Lie $G$ pada manifold $M$ adalah pemetaan mulus
\[
  \Phi\colon G\times M\longrightarrow M,
  \qquad (g,p)\longmapsto g\cdot p,
\]
dengan $e\cdot p=p$ dan $(gh)\cdot p=g\cdot(h\cdot p)$.
Orbit dan stabilisator titik $p$ adalah
\[
  G\cdot p=\{g\cdot p:g\in G\},
  \qquad
  G_p=\{g\in G:g\cdot p=p\}.
\]
\end{definition}

Stabilisator $G_p$ adalah subgrup tertutup, sebab ia adalah prapeta
$\{p\}$ di bawah pemetaan kontinu $g\mapsto g\cdot p$. Pemetaan orbit
\[
  \Phi_p\colon G\longrightarrow M,
  \qquad \Phi_p(g)=g\cdot p
\]
memenuhi $\Phi_p(gh)=g\cdot\Phi_p(h)$.

\begin{definition}\label{o011-bl-def-fundamental-field}
Untuk $v\in\mathfrak g$, medan vektor fundamental yang dihasilkan oleh $v$
adalah
\[
  v_M(p)=\left.\frac{d}{dt}\right|_{0}\exp(tv)\cdot p
  =T_e\Phi_p(v).
\]
\end{definition}

Dengan konvensi aksi kiri dan medan invarian-kiri yang dipakai di sini,
pemetaan $v\mapsto v_M$ merupakan antihomomorfisme aljabar Lie:
\[
  [v_M,w_M]=-[v,w]_M.
\]
Tanda minus berasal dari fakta bahwa $p\mapsto\exp(tv)\cdot p$ membentuk
aksi kiri, sedangkan medan fundamentalnya terkait secara alami dengan medan
invarian-kanan pada $G$. Mengubah salah satu konvensi akan mengubah tanda.

\begin{Proposition}\label{o011-bl-prop-orbit-tangent}
Untuk setiap $p\in M$,
\[
  \operatorname{im}(T_e\Phi_p)
  =\{v_M(p):v\in\mathfrak g\}
\]
adalah ruang arah tangensial orbit di $p$. Kernel
\[
  \mathfrak g_p=\ker(T_e\Phi_p)
\]
adalah aljabar Lie stabilisator. Jika orbit mempunyai struktur manifold
homogen yang lazim, maka
\[
  T_p(G\cdot p)\cong\mathfrak g/\mathfrak g_p,
  \qquad
  \dim(G\cdot p)=\dim G-\dim G_p.
\]
\end{Proposition}

\begin{proof}
Identitas pertama adalah definisi diferensial pemetaan orbit. Kurva
$\exp(tv)$ terletak di $G_p$ tepat secara infinitesimal ketika
$v_M(p)=0$, sehingga kernelnya adalah ruang singgung stabilisator di $e$.
Teorema isomorfisme linear memberi
$\mathfrak g/\ker(T_e\Phi_p)\cong\operatorname{im}(T_e\Phi_p)$.
Apabila orbit direalisasikan sebagai manifold homogen $G/G_p$, diferensial
pemetaan hasil bagi mengidentifikasi ruang terakhir dengan $T_p(G\cdot p)$.
Rumus dimensi mengikuti.
\end{proof}

\begin{beispiel}\label{o011-bl-ex-sphere-action}
$\mathrm{SO}(n)$ bertindak pada $S^{n-1}$ melalui perkalian matriks. Aksi ini
transitif. Stabilisator $e_n$ terdiri atas matriks berbentuk
\[
  \begin{pmatrix} A&0\\0&1\end{pmatrix},
  \qquad A\in\mathrm{SO}(n-1),
\]
sehingga $S^{n-1}\cong\mathrm{SO}(n)/\mathrm{SO}(n-1)$. Rumus dimensi memberi
\[
 \frac{n(n-1)}2-\frac{(n-1)(n-2)}2=n-1,
\]
sesuai dengan dimensi bola satuan.
\end{beispiel}

\subsection{Ringkasan struktural}

Data global dan infinitesimal yang dibangun di atas tersusun sebagai berikut:
\[
\begin{array}{ccl}
G&\longmapsto&\mathfrak g=T_eG,\\
F\colon G\to H&\longmapsto&dF_e\colon\mathfrak g\to\mathfrak h,\\
v\in\mathfrak g&\longmapsto&X^v\text{ dan }\exp(tv),\\
g\in G&\longmapsto&\operatorname{Ad}_g,\\
v\in\mathfrak g&\longmapsto&\operatorname{ad}_v=[v,\cdot],\\
G\curvearrowright M&\longmapsto&v_M(p)=T_e\Phi_p(v).
\end{array}
\]
Eksponensial memberi koordinat lokal di sekitar identitas, tetapi tidak harus
surjektif secara global dan tidak mengubah persoalan global menjadi persoalan
linear. Orbit mengubah simetri menjadi geometri submanifold, sedangkan
stabilisator mengukur simetri yang tidak menggerakkan suatu titik.

% O011-BL-A — latihan, petunjuk, solusi, dan uji penguasaan asli.
% Lisensi: CC BY-SA 4.0.
% Provenans: disusun oleh OpenAI Codex gpt-5.6-sol, Ultra, atas arahan pengguna.

\subsection{Latihan, petunjuk bertahap, dan solusi}

Setiap latihan di bagian ini mempunyai identitas stabil yang tidak bergantung
pada bahasa. Petunjuk disusun bertahap: baca petunjuk kedua hanya jika petunjuk
pertama belum cukup. Solusi yang diberikan merupakan bagian asli modul ini,
bukan solusi yang dinisbahkan kepada sumber Brenner.

\subsubsection*{Latihan L1 — Grup perkalian positif (O011-BL-E01)}
\label{o011-bl-e01}

Tunjukkan bahwa $\mathbb R_{>0}$ dengan perkalian merupakan grup Lie
satu dimensi. Tentukan aljabar Lie dan pemetaan eksponensialnya. Tunjukkan pula
bahwa logaritma memberikan isomorfisme grup Lie
$\mathbb R_{>0}\cong(\mathbb R,+)$.

\paragraph{Petunjuk 1.}
Gunakan koordinat biasa pada himpunan terbuka $\mathbb R_{>0}\subset\mathbb R$.

\paragraph{Petunjuk 2.}
Subgrup satu-parameter berkecepatan awal $v$ harus memenuhi
$\gamma'(t)=v\gamma(t)$ dan $\gamma(0)=1$.

\paragraph{Solusi.}
Perkalian $(x,y)\mapsto xy$ dan invers $x\mapsto x^{-1}$ mulus pada
$\mathbb R_{>0}$. Ruang singgung di identitas $1$ dapat diidentifikasi dengan
$\mathbb R$. Karena grupnya abelian, braket Lie identik nol. Persamaan untuk
medan invarian-kiri yang bernilai $v$ di $1$ ialah
$\gamma'(t)=v\gamma(t)$, sehingga $\gamma_v(t)=e^{tv}$ dan
$\exp(v)=e^v$. Akhirnya $\log(xy)=\log x+\log y$; logaritma dan eksponensial
saling invers serta mulus.

\subsubsection*{Latihan L2 — Ruang singgung grup matriks (O011-BL-E02)}
\label{o011-bl-e02}

Buktikan langsung bahwa
\[
 T_I\mathrm{SL}_n(\mathbb R)=\{X:\operatorname{tr}X=0\}
 \quad\text{dan}\quad
 T_I\mathrm O(n)=\{X:X^{\mathsf T}+X=0\}.
\]
Hitung dimensinya.

\paragraph{Petunjuk 1.}
Diferensiasikan persamaan $\det A(t)=1$ dan
$A(t)^{\mathsf T}A(t)=I$ pada $t=0$.

\paragraph{Petunjuk 2.}
Untuk arah sebaliknya, gunakan $e^{tX}$ dan identitas
$\det(e^{tX})=e^{t\operatorname{tr}X}$; untuk matriks antisimetri gunakan
$(e^{tX})^{\mathsf T}e^{tX}=I$.

\paragraph{Solusi.}
Jika $A(0)=I$ dan $A'(0)=X$, rumus turunan determinan memberi
\[
 0=\left.\frac d{dt}\right|_0\det A(t)=\operatorname{tr}X.
\]
Sebaliknya, jika $\operatorname{tr}X=0$, kurva $e^{tX}$ berada di
$\mathrm{SL}_n$ dan berkecepatan $X$ pada nol. Jadi ruang singgung pertama
tepat seperti dinyatakan dan berdimensi $n^2-1$.

Untuk $A(t)\in\mathrm O(n)$,
\[
 0=\left.\frac d{dt}\right|_0 A(t)^{\mathsf T}A(t)=X^{\mathsf T}+X.
\]
Jika $X^{\mathsf T}=-X$, maka
$(e^{tX})^{\mathsf T}=e^{-tX}$, sehingga $e^{tX}\in\mathrm O(n)$. Matriks
antisimetri ditentukan oleh entri di atas diagonal; dimensinya
$n(n-1)/2$.

\subsubsection*{Latihan L3 — Medan invarian pada grup umum linear (O011-BL-E03)}
\label{o011-bl-e03}

Untuk $A\in M_n(\mathbb R)$, tentukan medan invarian-kiri $X^A$ dan medan
invarian-kanan $Y^A$ pada $\mathrm{GL}_n(\mathbb R)$. Hitung alirannya.

\paragraph{Petunjuk 1.}
Diferensial perkalian kiri $L_g(h)=gh$ pada $I$ adalah $A\mapsto gA$.

\paragraph{Petunjuk 2.}
Periksa kurva $g e^{tA}$ dan $e^{tA}g$.

\paragraph{Solusi.}
Kita memperoleh
\[
  X^A_g=gA,
  \qquad Y^A_g=Ag.
\]
Kurva integral $X^A$ yang berawal di $g$ ialah $t\mapsto ge^{tA}$, sebab
turunannya $ge^{tA}A$. Kurva integral $Y^A$ ialah
$t\mapsto e^{tA}g$. Keduanya lengkap untuk semua $t\in\mathbb R$.

\subsubsection*{Latihan L4 — Braket pada grup afin garis (O011-BL-E04)}
\label{o011-bl-e04}

Pertimbangkan
\[
 G=\left\{\begin{pmatrix}a&b\\0&1\end{pmatrix}:a>0,
 b\in\mathbb R\right\}.
\]
Tentukan aljabar Lie dan braketnya. Apakah aljabar Lie ini abelian?

\paragraph{Petunjuk 1.}
Tuliskan arah singgung di identitas sebagai
$X(x,y)=\left(\begin{smallmatrix}x&y\\0&0\end{smallmatrix}\right)$.

\paragraph{Petunjuk 2.}
Hitung $X(x,y)X(u,v)-X(u,v)X(x,y)$.

\paragraph{Solusi.}
Aljabar Lie terdiri atas semua
\[
 X(x,y)=\begin{pmatrix}x&y\\0&0\end{pmatrix}.
\]
Perkalian langsung memberi
\[
 [X(x,y),X(u,v)]
 =\begin{pmatrix}0&xv-uy\\0&0\end{pmatrix}
 =X(0,xv-uy).
\]
Aljabar ini tidak abelian; misalnya
$[X(1,0),X(0,1)]=X(0,1)$.

\subsubsection*{Latihan L5 — Eksponensial rotasi dan geseran (O011-BL-E05)}
\label{o011-bl-e05}

Hitung eksponensial dari
\[
 J=\begin{pmatrix}0&-1\\1&0\end{pmatrix},
 \qquad
 N=\begin{pmatrix}0&1\\0&0\end{pmatrix}.
\]
Jelaskan subgrup satu-parameter yang dihasilkan.

\paragraph{Petunjuk 1.}
Gunakan $J^2=-I$ dan $N^2=0$.

\paragraph{Petunjuk 2.}
Pisahkan deret pangkat $e^{tJ}$ menjadi suku genap dan ganjil.

\paragraph{Solusi.}
Karena $J^{2k}=(-1)^kI$ dan $J^{2k+1}=(-1)^kJ$,
\[
 e^{tJ}=\cos(t)I+\sin(t)J
 =\begin{pmatrix}\cos t&-\sin t\\\sin t&\cos t\end{pmatrix}.
\]
Ini adalah subgrup rotasi di $\mathrm{SO}(2)$. Karena $N^2=0$,
\[
 e^{tN}=I+tN=\begin{pmatrix}1&t\\0&1\end{pmatrix},
\]
yakni subgrup geseran unipoten.

\subsubsection*{Latihan L6 — Determinan eksponensial (O011-BL-E06)}
\label{o011-bl-e06}

Buktikan bahwa untuk setiap matriks nyata $X$,
\[
  \det(e^X)=e^{\operatorname{tr}X}.
\]
Gunakan hasil ini untuk menunjukkan bahwa eksponensial
$\mathfrak{sl}_n\to\mathrm{SL}_n$ terdefinisi dengan baik.

\paragraph{Petunjuk 1.}
Tinjau $f(t)=\det(e^{tX})$ dan gunakan rumus Jacobi untuk turunan determinan.

\paragraph{Petunjuk 2.}
Tunjukkan $f'(t)=\operatorname{tr}(X)f(t)$ dengan $f(0)=1$.

\paragraph{Solusi.}
Karena $(e^{tX})'=Xe^{tX}$ dan $e^{tX}$ invertibel,
\[
 f'(t)=f(t)\operatorname{tr}\bigl(e^{-tX}Xe^{tX}\bigr)
 =f(t)\operatorname{tr}X.
\]
Solusi persamaan skalar ini dengan $f(0)=1$ ialah
$f(t)=e^{t\operatorname{tr}X}$. Ambil $t=1$. Jika
$\operatorname{tr}X=0$, determinan $e^X$ sama dengan $1$.

\subsubsection*{Latihan L7 — Adjoint grup matriks (O011-BL-E07)}
\label{o011-bl-e07}

Untuk subgrup matriks $G\subseteq\mathrm{GL}_n(\mathbb R)$, buktikan
\[
 \operatorname{Ad}_gX=gXg^{-1},
 \qquad
 \operatorname{ad}_XY=XY-YX.
\]

\paragraph{Petunjuk 1.}
Diferensiasikan $g(I+tX+o(t))g^{-1}$.

\paragraph{Petunjuk 2.}
Kemudian diferensiasikan $e^{tX}Ye^{-tX}$ pada $t=0$.

\paragraph{Solusi.}
Kurva $A(t)$ melalui $I$ dengan $A'(0)=X$ memberi
\[
 \left.\frac d{dt}\right|_0gA(t)g^{-1}=gXg^{-1},
\]
yang membuktikan rumus $\operatorname{Ad}$. Selanjutnya
\[
 \left.\frac d{dt}\right|_0 e^{tX}Ye^{-tX}
 =XY-YX,
\]
sehingga $\operatorname{ad}_X(Y)=[X,Y]$.

\subsubsection*{Latihan L8 — Komutator infinitesimal (O011-BL-E08)}
\label{o011-bl-e08}

Untuk $X,Y\in M_n(\mathbb R)$, tetapkan
\[
 c(t)=e^{tX}e^{tY}e^{-tX}e^{-tY}.
\]
Tunjukkan bahwa
\[
 c(t)=I+t^2[X,Y]+O(t^3).
\]

\paragraph{Petunjuk 1.}
Gunakan $e^{tX}=I+tX+\tfrac12t^2X^2+O(t^3)$ untuk keempat faktor.

\paragraph{Petunjuk 2.}
Semua suku orde satu harus saling menghapus; kelompokkan suku orde dua.

\paragraph{Solusi.}
Kalikan dua faktor pertama dan dua faktor terakhir hingga orde dua:
\[
 e^{tX}e^{tY}
 =I+t(X+Y)+t^2\left(\tfrac12X^2+XY+\tfrac12Y^2\right)+O(t^3),
\]
\[
 e^{-tX}e^{-tY}
 =I-t(X+Y)+t^2\left(\tfrac12X^2+XY+\tfrac12Y^2\right)+O(t^3).
\]
Ketika kedua ekspansi dikalikan, koefisien orde satu hilang dan koefisien
orde dua menyederhana menjadi $XY-YX$. Jadi rumus yang diminta berlaku.
Rumus ini memperlihatkan bahwa braket mengukur kegagalan gerak kecil dalam
dua arah untuk saling komutatif.

\subsubsection*{Latihan L9 — Aksi lingkaran (O011-BL-E09)}
\label{o011-bl-e09}

Biarkan $\mathbb R$ bertindak pada $S^1\subset\mathbb C$ melalui
$t\cdot z=e^{it}z$. Tentukan orbit, stabilisator, medan fundamental untuk
$v\in\mathbb R$, dan kernel diferensial pemetaan orbit di $0$.

\paragraph{Petunjuk 1.}
Untuk $z$ tetap, selesaikan $e^{it}z=z$.

\paragraph{Petunjuk 2.}
Diferensiasikan $e^{itv}z$ pada $t=0$.

\paragraph{Solusi.}
Orbit setiap $z$ adalah seluruh $S^1$. Stabilisatornya
$2\pi\mathbb Z$, suatu subgrup diskret. Medan fundamental ialah
\[
 v_{S^1}(z)=ivz.
\]
Pemetaan $v\mapsto ivz$ dari $\mathbb R$ ke $T_zS^1$ adalah isomorfisme,
sehingga kernelnya nol. Ini sesuai dengan aljabar Lie stabilisator diskret,
yang juga nol, meskipun stabilisator grupnya sendiri tidak trivial.

\subsubsection*{Latihan L10 — Orbit rotasi di ruang Euklides (O011-BL-E10)}
\label{o011-bl-e10}

Tentukan orbit dan stabilisator aksi standar $\mathrm{SO}(3)$ pada
$\mathbb R^3$. Untuk $p\ne0$, buktikan
\[
  T_p(\mathrm{SO}(3)\cdot p)=p^\perp.
\]

\paragraph{Petunjuk 1.}
Rotasi mempertahankan norma dan bertindak transitif pada setiap bola berjari-jari
positif.

\paragraph{Petunjuk 2.}
Medan fundamental yang ditentukan $X\in\mathfrak{so}(3)$ bernilai $Xp$.
Gunakan antisimetri untuk menunjukkan $Xp\perp p$ dan hitung dimensi.

\paragraph{Solusi.}
Orbit $0$ adalah $\{0\}$ dan stabilisatornya seluruh $\mathrm{SO}(3)$. Jika
$p\ne0$, orbitnya bola $S^2_{\|p\|}$ dan stabilisatornya isomorfik dengan
$\mathrm{SO}(2)$, yaitu rotasi di bidang $p^\perp$. Karena
$\langle Xp,p\rangle=-\langle p,Xp\rangle$, setiap $Xp$ tegak lurus pada
$p$. Pemetaan $X\mapsto Xp$ berperingkat dua: kernelnya berdimensi satu,
yakni aljabar stabilisator. Jadi citranya adalah seluruh ruang dua dimensi
$p^\perp$.

\subsubsection*{Latihan L11 — Orbit konjugasi (O011-BL-E11)}
\label{o011-bl-e11}

$\mathrm{GL}_n(\mathbb R)$ bertindak pada $M_n(\mathbb R)$ melalui
$g\cdot A=gAg^{-1}$. Tentukan stabilisator $A$ dan ruang tangensial orbitnya
di $A$.

\paragraph{Petunjuk 1.}
Stabilisator adalah semua matriks invertibel yang komutatif dengan $A$.

\paragraph{Petunjuk 2.}
Diferensiasikan $e^{tX}Ae^{-tX}$.

\paragraph{Solusi.}
Stabilisatornya adalah sentralisator invertibel
\[
 (\mathrm{GL}_n)_A=\{g\in\mathrm{GL}_n:gA=Ag\}.
\]
Medan fundamental arah $X$ bernilai
\[
 \left.\frac d{dt}\right|_0e^{tX}Ae^{-tX}=XA-AX=[X,A].
\]
Oleh karena itu
\[
 T_A(\mathrm{GL}_n\cdot A)=\{[X,A]:X\in M_n(\mathbb R)\}.
\]
Kernelnya adalah aljabar semua matriks yang komutatif dengan $A$.

\subsubsection*{Latihan L12 — Kernel homomorfisme (O011-BL-E12)}
\label{o011-bl-e12}

Misalkan $F\colon G\to H$ homomorfisme grup Lie. Buktikan bahwa
$\ker(dF_e)$ adalah ideal dalam $\mathfrak g$ dan
$\operatorname{im}(dF_e)$ adalah subaljabar Lie dalam $\mathfrak h$.

\paragraph{Petunjuk 1.}
Gunakan $dF_e[v,w]=[dF_ev,dF_ew]$.

\paragraph{Petunjuk 2.}
Untuk sifat ideal, ambil $v$ di kernel dan $w$ sebarang.

\paragraph{Solusi.}
Jika $v\in\ker(dF_e)$ dan $w\in\mathfrak g$, maka
\[
 dF_e[v,w]=[dF_ev,dF_ew]=[0,dF_ew]=0.
\]
Jadi $[v,w]$ kembali berada dalam kernel; kernel adalah ideal. Jika
$a=dF_ev$ dan $b=dF_ew$ berada dalam citra, maka
\[
 [a,b]=[dF_ev,dF_ew]=dF_e[v,w]
\]
juga berada dalam citra. Jadi citra adalah subaljabar Lie.

\subsection{Uji penguasaan kumulatif}

Keempat soal berikut dimaksudkan sebagai penilaian setelah seluruh modul.
Setiap soal mempunyai rubrik, solusi lengkap, dan parameter alternatif untuk
menghasilkan bentuk ekuivalen tanpa mengubah kompetensi yang diuji.

\subsubsection*{Penguasaan M1 — Grup Heisenberg (O011-BL-M01)}
\label{o011-bl-m01}

Untuk
\[
 H=\left\{h(x,y,z)=
 \begin{pmatrix}1&x&z\\0&1&y\\0&0&1\end{pmatrix}:x,y,z\in\mathbb R\right\},
\]
(a) buktikan bahwa $H$ adalah grup Lie; (b) tentukan hukum grup dalam
koordinat $(x,y,z)$; (c) tentukan aljabar Lie dan semua braket basis
$X=E_{12}$, $Y=E_{23}$, $Z=E_{13}$; (d) hitung
$\exp(aX+bY+cZ)$.

\paragraph{Petunjuk.}
Semua matriks $N=aX+bY+cZ$ memenuhi $N^3=0$. Perhatikan $XY=Z$ tetapi
$YX=0$.

\paragraph{Solusi.}
Himpunan ini adalah submanifold afin tiga dimensi dari $M_3(\mathbb R)$.
Perkalian dan invers matriks membatas menjadi pemetaan mulus, dan
\[
 h(x,y,z)h(x',y',z')
 =h(x+x',y+y',z+z'+xy').
\]
Inversnya $h(x,y,z)^{-1}=h(-x,-y,-z+xy)$. Aljabar Lie ialah
$\operatorname{span}\{X,Y,Z\}$ dengan
\[
 [X,Y]=Z,
 \qquad [X,Z]=[Y,Z]=0.
\]
Untuk $N=aX+bY+cZ$, berlaku $N^2=abZ$ dan $N^3=0$, sehingga
\[
 \exp N=I+N+\tfrac12N^2
 =h\left(a,b,c+\tfrac12ab\right).
\]

\paragraph{Rubrik (12 poin).}
Dua poin untuk manifold dan kemulusan operasi; tiga poin untuk hukum grup dan
invers; tiga poin untuk ruang aljabar serta semua braket; empat poin untuk
perhitungan eksponensial beserta alasan deret berhenti.

\paragraph{Parameter alternatif.}
Ganti koordinat entri kanan atas dengan $\lambda z$ untuk suatu
$\lambda\ne0$. Peserta harus menyesuaikan hukum grup, basis pusat, dan
koefisien eksponensial; kompetensi yang diuji tetap sama.

\subsubsection*{Penguasaan M2 — Grup afin dan eksponensial (O011-BL-M02)}
\label{o011-bl-m02}

Untuk grup $G$ pada Latihan L4, (a) hitung
$\exp X(x,y)$; (b) tentukan $\operatorname{Ad}_{g(a,b)}X(x,y)$; (c) verifikasi
langsung $\operatorname{ad}_{X(x,y)}X(u,v)=[X(x,y),X(u,v)]$.

\paragraph{Petunjuk.}
Jika $x\ne0$, jumlahkan deret entri kanan atas; perlakukan $x=0$ sebagai
limit. Gunakan
$g(a,b)^{-1}=\left(\begin{smallmatrix}a^{-1}&-a^{-1}b\\0&1\end{smallmatrix}\right)$.

\paragraph{Solusi.}
Pangkat-pangkat $X=X(x,y)$ memenuhi, untuk $k\ge1$,
\[
 X^k=\begin{pmatrix}x^k&x^{k-1}y\\0&0\end{pmatrix}.
\]
Maka
\[
 \exp X(x,y)=
 \begin{cases}
 \left(\begin{smallmatrix}e^x&y(e^x-1)/x\\0&1\end{smallmatrix}\right),&x\ne0,\\[2mm]
 \left(\begin{smallmatrix}1&y\\0&1\end{smallmatrix}\right),&x=0.
 \end{cases}
\]
Perkalian matriks memberi
\[
 \operatorname{Ad}_{g(a,b)}X(x,y)=X(x,ay-bx).
\]
Diferensiasikan rumus ini pada $g=\exp(tX(x,y))$ dan $t=0$, atau hitung
komutator langsung, untuk memperoleh
\[
 \operatorname{ad}_{X(x,y)}X(u,v)=X(0,xv-uy).
\]

\paragraph{Rubrik (12 poin).}
Lima poin untuk eksponensial termasuk kasus $x=0$; empat poin untuk adjoint;
tiga poin untuk diferensiasi atau komutator yang memverifikasi rumus
$\operatorname{ad}$.

\paragraph{Parameter alternatif.}
Gunakan representasi afin dengan matriks
$\left(\begin{smallmatrix}a&0&b\\0&1&0\\0&0&1\end{smallmatrix}\right)$.
Peserta harus mengenali bahwa perhitungan esensial tetap dua dimensi.

\subsubsection*{Penguasaan M3 — Bola sebagai ruang homogen (O011-BL-M03)}
\label{o011-bl-m03}

Untuk aksi $\mathrm{SO}(n)$ pada $S^{n-1}$ dengan $n\ge3$, (a) buktikan
transitivitas; (b) tentukan stabilisator $e_n$; (c) hitung diferensial pemetaan
orbit di identitas; (d) deduksi
$S^{n-1}\cong\mathrm{SO}(n)/\mathrm{SO}(n-1)$ dan periksa dimensinya.

\paragraph{Petunjuk.}
Lengkapi setiap vektor satuan menjadi basis ortonormal berorientasi. Untuk
$X\in\mathfrak{so}(n)$, diferensial pemetaan orbit adalah $X\mapsto Xe_n$.

\paragraph{Solusi.}
Jika $p,q\in S^{n-1}$, lengkapi masing-masing menjadi basis ortonormal
berorientasi. Pemetaan yang membawa basis pertama ke basis kedua berada dalam
$\mathrm{SO}(n)$ dan membawa $p$ ke $q$, sehingga aksi transitif. Matriks yang
menetapkan $e_n$ mempunyai bentuk blok
$\operatorname{diag}(A,1)$ dengan $A\in\mathrm{SO}(n-1)$.

Pemetaan orbit $\Phi_{e_n}(g)=ge_n$ mempunyai diferensial
$T_I\Phi_{e_n}(X)=Xe_n$. Citranya adalah $e_n^\perp$, sedangkan kernelnya
terdiri atas matriks antisimetri dengan baris dan kolom terakhir nol, yaitu
$\mathfrak{so}(n-1)$. Karena itu ruang homogen adalah $S^{n-1}$ dan
\[
 \dim\mathrm{SO}(n)-\dim\mathrm{SO}(n-1)
 =\frac{n(n-1)-(n-1)(n-2)}2=n-1.
\]

\paragraph{Rubrik (12 poin).}
Tiga poin untuk transitivitas; dua poin untuk stabilisator; empat poin untuk
diferensial, kernel, dan citra; tiga poin untuk identifikasi hasil bagi dan
perhitungan dimensi.

\paragraph{Parameter alternatif.}
Ganti bola satuan dengan bola berjari-jari $r>0$ dan titik dasar $re_1$.
Stabilisator dan hasil bagi sama hingga pilihan titik dasar; diferensial
pemetaan orbit memperoleh faktor $r$.

\subsubsection*{Penguasaan M4 — Naturalisasi eksponensial dan kegagalan komutatif (O011-BL-M04)}
\label{o011-bl-m04}

Misalkan $F\colon G\to H$ homomorfisme grup Lie. (a) buktikan
$F(\exp_Gv)=\exp_H(dF_ev)$; (b) deduksi
$F(\exp_G(tv))=\exp_H(t,dF_ev)$ untuk semua $t$; (c) berikan matriks
$X,Y\in\mathfrak{gl}_2$ yang menunjukkan bahwa umumnya
$e^{X+Y}\ne e^Xe^Y$; (d) nyatakan hipotesis sederhana yang menjamin kesamaan.

\paragraph{Petunjuk.}
Untuk (a), bandingkan dua subgrup satu-parameter dengan kecepatan awal yang
sama. Untuk (c), gunakan kelipatan kecil matriks $E_{12}$ dan $E_{21}$ atau
bandingkan suku orde dua.

\paragraph{Solusi.}
Kurva $t\mapsto F(\exp_G(tv))$ adalah homomorfisme
$\mathbb R\to H$ dan turunannya pada nol ialah $dF_ev$. Berdasarkan ketunggalan
subgrup satu-parameter, kurva itu sama dengan
$t\mapsto\exp_H(t,dF_ev)$. Pernyataan (a) adalah kasus $t=1$, sedangkan (b)
adalah identitas kurva penuh.

Ambil $X=sE_{12}$ dan $Y=sE_{21}$ dengan $s\ne0$ cukup kecil. Ekspansi hingga
orde dua memberi
\[
 e^Xe^Y=I+X+Y+XY+O(s^3),
\]
sedangkan
\[
 e^{X+Y}=I+X+Y+\tfrac12(XY+YX)+O(s^3).
\]
Karena $XY\ne YX$, kedua matriks berbeda untuk $s$ cukup kecil. Jika
$[X,Y]=0$, deret pangkat dapat dikelompokkan seperti bilangan komutatif dan
$e^{X+Y}=e^Xe^Y$.

\paragraph{Rubrik (12 poin).}
Empat poin untuk argumen subgrup satu-parameter; dua poin untuk identitas
berparameter; empat poin untuk contoh tandingan yang benar beserta verifikasi;
dua poin untuk hipotesis komutatif dan alasannya.

\paragraph{Parameter alternatif.}
Ganti pasangan $E_{12},E_{21}$ dengan
$X=\left(\begin{smallmatrix}0&s\\0&0\end{smallmatrix}\right)$ dan
$Y=\left(\begin{smallmatrix}r&0\\0&0\end{smallmatrix}\right)$ untuk $rs\ne0$.
Peserta harus mendeteksi komutator tak nol dan membandingkan eksponensial.


\chapter{Jembatan Kohomologi de Rham dan Topologi Diferensial}
\noindent\textit{Materi asli edisi ini; bukan solusi atau teks yang disediakan oleh sumber Brenner/Wikiversity.}
% O011-BR — modul asli, bukan terjemahan atau adaptasi dari teks pembanding.
% Judul: Jembatan Kohomologi de Rham dan Topologi Diferensial
% Lisensi modul: CC BY-SA 4.0.
% Provenans: disusun oleh OpenAI Codex gpt-5.6-sol, Ultra, atas arahan pengguna.
% Prasyarat internal: Unit 14, 20, 22, dan 23 edisi ini.

\setcounter{section}{31}
\section{Jembatan Kohomologi de Rham dan Topologi Diferensial}
\label{o011-bridge-de-rham}

\subsection{Tujuan dan prasyarat}

Unit 14 memperkenalkan bentuk diferensial, Unit 20 turunan eksterior, Unit 22
partisi kesatuan, dan Unit 23 teorema Stokes. Modul ini menyatukan keempatnya
menjadi sebuah invariant topologis. Gagasan pokoknya sederhana: bentuk tertutup
yang bukan turunan eksterior bentuk lain mendeteksi ``lubang'' pada manifold.
Kita membuktikan hasil lokal, membangun invariansi homotopi, memakai barisan
Mayer--Vietoris untuk perhitungan global, lalu menandai jalan menuju derajat,
transversalitas, dan teori Morse tanpa berpura-pura membahas teori-teori itu
secara lengkap.

Semua manifold di bagian ini dianggap mulus, berdimensi hingga, Hausdorff, dan
memiliki basis terhitung. Ruang bentuk-$k$ mulus pada $M$ ditulis
$\Omega^k(M)$; untuk $k<0$ atau $k>\dim M$, ruang itu adalah nol.

\subsection{Kompleks de Rham}

\begin{definition}\label{o011-br-def-closed-exact}
Suatu bentuk $\omega\in\Omega^k(M)$ disebut \emph{tertutup} jika
$d\omega=0$, dan disebut \emph{eksak} jika terdapat
$\eta\in\Omega^{k-1}(M)$ dengan $\omega=d\eta$.
\end{definition}

Karena $d\circ d=0$, setiap bentuk eksak tertutup. Jadi turunan eksterior
membentuk kompleks rantai-kobalik
\[
  0\longrightarrow\Omega^0(M)\xrightarrow{d}\Omega^1(M)
  \xrightarrow{d}\cdots\xrightarrow{d}\Omega^n(M)
  \longrightarrow0.
\]

\begin{definition}\label{o011-br-def-cohomology}
Kohomologi de Rham derajat $k$ adalah ruang vektor
\[
  H^k_{\mathrm{dR}}(M)
  =\frac{\ker(d\colon\Omega^k(M)\to\Omega^{k+1}(M))}
  {\operatorname{im}(d\colon\Omega^{k-1}(M)\to\Omega^k(M))}.
\]
Kelas bentuk tertutup $\omega$ ditulis $[\omega]$.
\end{definition}

\begin{Proposition}\label{o011-br-prop-wedge-ring}
Produk wedge menurunkan produk bergradasi
\[
 H^p_{\mathrm{dR}}(M)\times H^q_{\mathrm{dR}}(M)
 \longrightarrow H^{p+q}_{\mathrm{dR}}(M),
 \qquad ([\alpha],[\beta])\longmapsto[\alpha\wedge\beta].
\]
Produk ini memenuhi
$[\alpha]\wedge[\beta]=(-1)^{pq}[\beta]\wedge[\alpha]$.
\end{Proposition}

\begin{proof}
Aturan Leibniz bergradasi memberi
\[
 d(\alpha\wedge\beta)=d\alpha\wedge\beta
 +(-1)^p\alpha\wedge d\beta.
\]
Jika $\alpha$ dan $\beta$ tertutup, wedge-nya tertutup. Jika
$\alpha$ diganti $\alpha+d\mu$, maka
\[
 (\alpha+d\mu)\wedge\beta-\alpha\wedge\beta
 =d(\mu\wedge\beta)
\]
karena $d\beta=0$. Argumen yang sama berlaku pada faktor kedua. Jadi produk
tidak bergantung pada wakil. Komutativitas bergradasi diwarisi dari bentuk.
\end{proof}

\begin{Proposition}\label{o011-br-prop-h0-components}
$H^0_{\mathrm{dR}}(M)$ adalah ruang fungsi yang konstan pada setiap komponen
terhubung. Khususnya, jika $M$ terhubung, $H^0_{\mathrm{dR}}(M)\cong\mathbb R$.
\end{Proposition}

\begin{proof}
Untuk fungsi mulus $f$, persamaan $df=0$ berarti semua turunan arah lokalnya
nol. Maka $f$ konstan pada setiap komponen terhubung. Tidak ada bentuk
berderajat $-1$, sehingga pada derajat nol tidak ada hasil bagi tambahan.
\end{proof}

\subsection{Tarikan balik dan funktorialitas}

Jika $f\colon M\to N$ mulus, tarikan balik memenuhi
\[
 f^*(\alpha\wedge\beta)=f^*\alpha\wedge f^*\beta,
 \qquad d(f^*\alpha)=f^*(d\alpha).
\]

\begin{Proposition}\label{o011-br-prop-pullback-cohomology}
Pemetaan $f$ menginduksi homomorfisme aljabar bergradasi
\[
 f^*\colon H^*_{\mathrm{dR}}(N)\longrightarrow H^*_{\mathrm{dR}}(M).
\]
Jika $g\colon N\to P$ mulus, maka $(g\circ f)^*=f^*\circ g^*$, dan
$\operatorname{id}_M^*$ adalah identitas.
\end{Proposition}

\begin{proof}
Komutasi dengan $d$ menunjukkan bahwa tarikan balik membawa bentuk tertutup
ke bentuk tertutup dan bentuk eksak ke bentuk eksak. Karena itu ia terdefinisi
pada kelas kohomologi. Semua identitas lain sudah berlaku pada tingkat bentuk.
\end{proof}

Perhatikan pembalikan arah: pemetaan $M\to N$ menghasilkan pemetaan
kohomologi dari $N$ ke $M$. Kohomologi de Rham adalah funktor kontravarian.

\subsection{Lema Poincar\'e pada himpunan berbentuk bintang}

Suatu himpunan terbuka $U\subseteq\mathbb R^n$ disebut berbentuk bintang
terhadap $0$ jika $tx\in U$ untuk semua $x\in U$ dan $t\in[0,1]$. Misalkan
$E=\sum_i x_i\partial_{x_i}$ adalah medan Euler dan $\iota_E$ kontraksi oleh
$E$.

Untuk $\omega\in\Omega^k(U)$ dengan $k\ge1$, definisikan operator homotopi oleh
\[
 (K\omega)_x(v_1,\ldots,v_{k-1})
 =\int_0^1 t^{k-1}\omega_{tx}(x,v_1,\ldots,v_{k-1})\,dt.
\]
Dengan medan Euler $E_y=y$, integran yang sama dapat ditulis
$t^{k-2}(\iota_E\omega)_{tx}$ untuk $t>0$; rumus pertama tetap teratur di
$t=0$ dan karena itu kita pakai sebagai definisi. Rumus yang lebih intrinsik
akan muncul pada homotopi umum di bawah.

\begin{Theorem}[Lema Poincar\'e]\label{o011-br-thm-poincare}
Jika $U\subseteq\mathbb R^n$ terbuka dan berbentuk bintang, setiap bentuk
tertutup $\omega\in\Omega^k(U)$ dengan $k\ge1$ adalah eksak. Dengan kata lain,
$H^k_{\mathrm{dR}}(U)=0$ untuk $k\ge1$.
\end{Theorem}

\begin{proof}
Ambil homotopi $F\colon U\times[0,1]\to U$, $F(x,t)=tx$. Perhitungan di
koordinat, atau rumus homotopi yang dibuktikan pada Teorema
\ref{o011-br-thm-chain-homotopy}, memberi
\[
 dK\omega+K(d\omega)=F_1^*\omega-F_0^*\omega.
\]
Untuk $k\ge1$, tarikan balik bentuk-$k$ oleh pemetaan konstan $F_0$ adalah nol,
sedangkan $F_1$ adalah identitas. Jika $d\omega=0$, maka
$\omega=d(K\omega)$.
\end{proof}

\begin{bemerkung}\label{o011-br-rem-local-global}
Lema Poincar\'e bersifat lokal: setiap titik manifold mempunyai lingkungan
koordinat yang dapat diperkecil menjadi bola, sehingga setiap bentuk tertutup
berderajat positif lokalnya eksak. Hambatan untuk memperoleh primitif global
adalah tepat informasi yang diukur kohomologi de Rham.
\end{bemerkung}

\subsection{Homotopi rantai dan invariansi homotopi}

Misalkan $F\colon M\times[0,1]\to N$ suatu homotopi mulus dan
$F_t(p)=F(p,t)$. Untuk $\omega\in\Omega^k(N)$, definisikan
\[
 K_F\omega
 =\int_0^1 \iota_t^*\bigl(\iota_{\partial_t}F^*\omega\bigr)\,dt
 \in\Omega^{k-1}(M),
\]
di mana $\iota_t(p)=(p,t)$.

\begin{Theorem}[Rumus homotopi rantai]\label{o011-br-thm-chain-homotopy}
Operator $K_F$ memenuhi
\[
 dK_F+K_Fd=F_1^*-F_0^*.
\]
\end{Theorem}

\begin{proof}
Pada $M\times[0,1]$, turunan terhadap $t$ dari tarikan balik sepanjang
$\iota_t$ memenuhi rumus Cartan
\[
 \frac d{dt}\iota_t^*(F^*\omega)
 =\iota_t^*\mathcal L_{\partial_t}(F^*\omega)
 =\iota_t^*\bigl(d\iota_{\partial_t}F^*\omega
 +\iota_{\partial_t}dF^*\omega\bigr).
\]
Integrasikan dari $0$ sampai $1$, gunakan bahwa tarikan balik komutatif dengan
$d$, dan kenali kedua suku sebagai $dK_F\omega$ serta $K_Fd\omega$.
\end{proof}

\begin{Korollar}\label{o011-br-cor-homotopy-invariance}
Pemetaan mulus yang homotopik menginduksi pemetaan kohomologi yang sama. Jika
$M$ dan $N$ ekuivalen-homotopi, maka
$H^*_{\mathrm{dR}}(M)\cong H^*_{\mathrm{dR}}(N)$ sebagai aljabar bergradasi.
\end{Korollar}

\begin{proof}
Jika $d\omega=0$, rumus homotopi memberi
$F_1^*\omega-F_0^*\omega=d(K_F\omega)$, sehingga kedua tarikan balik menentukan
kelas yang sama. Untuk ekuivalensi homotopi, terapkan funktorialitas pada
komposisi yang homotopik dengan identitas.
\end{proof}

\begin{Korollar}\label{o011-br-cor-contractible}
Jika $M$ kontraktibel, maka $H^0_{\mathrm{dR}}(M)\cong\mathbb R$ dan
$H^k_{\mathrm{dR}}(M)=0$ untuk $k>0$.
\end{Korollar}

\subsection{Barisan Mayer--Vietoris}

Misalkan $M=U\cup V$ dengan $U,V$ terbuka. Untuk setiap $k$ terdapat barisan
pendek kompleks
\[
0\longrightarrow\Omega^k(M)
\xrightarrow{r}(\Omega^k(U)\oplus\Omega^k(V))
\xrightarrow{\delta}\Omega^k(U\cap V)
\longrightarrow0,
\]
dengan
\[
 r(\omega)=(\omega|_U,\omega|_V),
 \qquad \delta(\alpha,\beta)=\alpha|_{U\cap V}-\beta|_{U\cap V}.
\]
Injektivitas dan ketepatan di tengah mengikuti sifat pengeleman bentuk. Untuk
surjektivitas $\delta$, pilih partisi kesatuan $\rho_U+\rho_V=1$ yang
tersubordinasi pada $U,V$; suatu bentuk $\gamma$ di $U\cap V$ diangkat oleh
perpanjangan nol yang dibangun dari $\rho_V\gamma$ pada $U$ dan
$-\rho_U\gamma$ pada $V$.

\begin{Theorem}[Mayer--Vietoris]\label{o011-br-thm-mv}
Penutup $M=U\cup V$ menghasilkan barisan eksak panjang
\[
\cdots\longrightarrow H^{k-1}_{\mathrm{dR}}(U\cap V)
\xrightarrow{\partial}H^k_{\mathrm{dR}}(M)
\xrightarrow{r^*}H^k_{\mathrm{dR}}(U)\oplus H^k_{\mathrm{dR}}(V)
\xrightarrow{\delta^*}H^k_{\mathrm{dR}}(U\cap V)
\xrightarrow{\partial}H^{k+1}_{\mathrm{dR}}(M)
\longrightarrow\cdots.
\]
\end{Theorem}

\begin{proof}
Barisan pendek kompleks di atas menghasilkan barisan eksak panjang kohomologi
melalui konstruksi standar pemetaan penghubung. Secara konkret, jika
$\gamma$ tertutup pada $U\cap V$, pilih $(\alpha,\beta)$ dengan
$\delta(\alpha,\beta)=\gamma$. Karena $d\gamma=0$, pasangan
$(d\alpha,d\beta)$ saling cocok pada irisan, sehingga menempel menjadi bentuk
tertutup $\omega$ pada $M$. Tetapkan $\partial[\gamma]=[\omega]$. Perubahan
pilihan hanya mengubah $\omega$ dengan bentuk eksak, dan pemeriksaan langsung
memberi ketepatan pada setiap suku.
\end{proof}

\subsection{Perhitungan dasar}

\begin{Proposition}\label{o011-br-prop-circle}
Kohomologi de Rham lingkaran adalah
\[
 H^0_{\mathrm{dR}}(S^1)\cong\mathbb R,
 \qquad H^1_{\mathrm{dR}}(S^1)\cong\mathbb R,
 \qquad H^k_{\mathrm{dR}}(S^1)=0\quad(k\ge2).
\]
Isomorfisme pada derajat satu dapat diberikan oleh periode
\[
 [\omega]\longmapsto\int_{S^1}\omega.
\]
\end{Proposition}

\begin{proof}
Derajat nol mengikuti keterhubungan, dan derajat di atas satu nol karena
dimensi. Tutupi $S^1$ oleh dua busur terbuka kontraktibel $U,V$ yang irisannya
mempunyai dua komponen. Bagian relevan Mayer--Vietoris adalah
\[
 0\to H^0(S^1)\to\mathbb R\oplus\mathbb R
 \to\mathbb R^2\to H^1(S^1)\to0.
\]
Pemetaan tengah mengirim $(a,b)$ ke $(a-b,a-b)$, sehingga cokernelnya satu
dimensi. Integral lenyap pada bentuk eksak dan tidak nol pada bentuk sudut
ternormalisasi; karena itu periode memberi isomorfisme yang dinyatakan.
\end{proof}

Pada $\mathbb R^2\setminus\{0\}$, bentuk
\[
 \eta=\frac{1}{2\pi}\frac{x\,dy-y\,dx}{x^2+y^2}
\]
tertutup dan integralnya pada lingkaran satuan adalah $1$. Jadi ia tidak eksak
dan mewakili generator kelas derajat satu.

\begin{Theorem}\label{o011-br-thm-sphere}
Untuk $n\ge1$,
\[
 H^k_{\mathrm{dR}}(S^n)\cong
 \begin{cases}
 \mathbb R,&k=0\text{ atau }k=n,\\
 0,&\text{selain itu}.
 \end{cases}
\]
\end{Theorem}

\begin{proof}
Kasus $n=1$ sudah dibuktikan. Untuk $n\ge2$, tutupi $S^n$ oleh lingkungan
terbuka kutub utara dan selatan yang masing-masing kontraktibel dan yang
irisannya ekuivalen-homotopi dengan $S^{n-1}$. Mayer--Vietoris, di luar
derajat nol, memberi isomorfisme penghubung
\[
 H^{k-1}_{\mathrm{dR}}(S^{n-1})\cong H^k_{\mathrm{dR}}(S^n)
 \qquad(1<k\le n).
\]
Pada derajat satu, keterhubungan penutup dan irisan memberi nol. Induksi pada
$n$ menghasilkan satu kelas di derajat teratas dan nol pada derajat antara.
Derajat nol kembali mengikuti keterhubungan.
\end{proof}

Jika $S^n$ diberi orientasi, kelas teratas dapat dinormalisasi oleh
$\int_{S^n}\omega=1$. Teorema Stokes memastikan bahwa integral hanya
bergantung pada kelas kohomologi.

\subsection{Teorema de Rham dan gerbang topologi diferensial}

\begin{Theorem}[Teorema de Rham — pernyataan]\label{o011-br-thm-de-rham}
Integrasi bentuk pada rantai singular mulus menginduksi isomorfisme alami
\[
 H^k_{\mathrm{dR}}(M)\xrightarrow{\ \cong\ }
 H^k_{\mathrm{sing}}(M;\mathbb R)
\]
untuk setiap manifold mulus $M$ dan setiap $k$.
\end{Theorem}

Teorema ini dinyatakan, bukan dibuktikan, karena pembuktian lengkap memerlukan
kompleks singular, subdivisi, dan argumen berkas atau resolusi halus yang berada
di luar jembatan ini. Isi konseptualnya ialah bahwa perhitungan dengan bentuk
diferensial menghasilkan invariant topologis yang sama dengan kohomologi
singular berkoefisien nyata.

\paragraph{Derajat.}
Jika $f\colon M\to N$ mulus antara manifold tertutup, terhubung, berorientasi,
dan berdimensi sama $n$, tindakan $f^*$ pada kohomologi teratas adalah
perkalian oleh suatu bilangan bulat $\deg f$. Secara ekuivalen,
\[
 \int_M f^*\omega=(\deg f)\int_N\omega
\]
untuk setiap bentuk volume tertutup $\omega$. Dengan memilih nilai regular,
derajat dapat dihitung sebagai jumlah bertanda titik-titik prapeta.

\paragraph{Transversalitas.}
Jika $f\colon M\to N$ transversal terhadap submanifold $Z\subseteq N$, maka
$f^{-1}(Z)$ adalah submanifold berdimensi
$\dim M-\operatorname{codim}Z$. Transversalitas memungkinkan kelas kohomologi,
bilangan perpotongan, dan derajat direpresentasikan oleh geometri prapeta yang
stabil di bawah perturbasi kecil. Modul ini hanya memakai pernyataan gerbang
tersebut; teorema aproksimasi transversal memerlukan pembahasan tersendiri.

\paragraph{Teori Morse.}
Fungsi Morse mempunyai titik kritis nondegenerat. Ketika nilai melewati titik
kritis berindeks $\lambda$, sublevel berubah melalui pemasangan pegangan
$\lambda$. Dari sini lahir kompleks Morse, ketaksamaan Morse, dan hubungan
antara jumlah titik kritis dengan kohomologi. Pembuktian memerlukan lema Morse,
aliran gradien, dan analisis kekompakan; tidak satu pun disamarkan sebagai
telah dibuktikan di sini.

\begin{bemerkung}\label{o011-br-rem-roadmap}
Jembatan ini menutup rantai logis yang diperlukan untuk kurikulum dasar:
bentuk dan $d$ $\to$ kohomologi $\to$ invariansi homotopi $\to$
Mayer--Vietoris dan contoh $\to$ teorema de Rham sebagai identifikasi global
$\to$ gerbang derajat, transversalitas, serta Morse. Teori lanjutan dapat
memulai tepat dari tiga gerbang terakhir tanpa mengulang bagian sebelumnya.
\end{bemerkung}

\input{generated/bridges/bridge-de-rham-assessment.build.tex}

\part{Formulir Ujian}
\chapter*{Provenans formulir ujian}
\addcontentsline{toc}{chapter}{Provenans formulir ujian}
Sepuluh formulir peserta berikut mempertahankan teks peserta resmi. Sepuluh formulir solusi resmi sesudahnya hanya memuat solusi yang tersedia pada sumber resmi yang dibekukan. Kekosongan sumber tetap kosong dan tidak diisi pada permukaan resmi ini.
% Lokalisasi khusus permukaan pembaca untuk formulir ujian dalam edisi lengkap.
% Nama perintah sumber dipertahankan; hanya teks yang dicetak yang dilokalkan.

\renewcommand{\fachbereich}{Bidang studi}
\renewcommand{\dozent}{Pengajar}
\renewcommand{\klausurdatum}{Tanggal}
\renewcommand{\klausurgebiet}{Matematika}
\renewcommand{\klausurtyp}{Ujian}

\renewcommand{\klausurmodus}{%
Durasi: sepuluh menit untuk membaca dan dua jam penuh untuk mengerjakan.
Alat bantu tidak diperbolehkan. Semua jawaban harus diberi alasan.

Jumlah poin seluruhnya $64$. Berlaku aturan ambang dasar, yakni penilaian
untuk setiap soal atau bagian soal pada umumnya dimulai pada separuh jumlah
poin. Untuk lulus diperlukan $16$ poin; mulai $32$ poin diperoleh nilai satu.

Tuliskan nama dan nomor mahasiswa Anda dengan jelas pada lembar sampul dan
setiap lembar berikutnya. Semoga berhasil!}

\renewcommand{\testklausurmodus}{\klausurmodus}

\renewcommand{\klausurname}{%
\renewcommand{\arraystretch}{1.5}%
\begin{tabular}{@{}p{3.2cm}p{10cm}}
Nama lengkap: &
..................................................................................\\
Nomor mahasiswa: &
..................................................................................\\
\end{tabular}}

\renewcommand{\klausurvorspann}[5]{%
\begin{tabular}{@{}p{3.2cm}p{10cm}}
Bidang studi: & #1\\
Tanggal: & #2\\
Pengajar: & #3
\end{tabular}
\vspace{3ex}

\begin{center}
{\bfseries\large Ujian\\#4\ #5}
\end{center}

\vspace{2ex}
\klausurmodus
\vspace{2ex}
\klausurname\par}

\renewcommand{\testklausurvorspann}[5]{\klausurvorspann{#1}{#2}{#3}{#4}{#5}}
\renewcommand{\klausurmitloesungenvorspann}[5]{%
\klausurvorspann{#1}{#2}{#3}{#4}{#5}
\begin{center}\bfseries Formulir dengan solusi\end{center}}

\renewcommand{\klausurnote}{%
\begin{tabular}{@{}p{3.2cm}p{10cm}}Nilai: & \end{tabular}}

\renewcommand{\klausurtaeuschungwarnung}{%
\vspace{2ex}%
Saya memahami bahwa pelanggaran berat terhadap aturan ujian dapat membatalkan
seluruh komponen penilaian dalam mata kuliah ini. Penggunaan literatur atau
alat bantu elektronik termasuk pelanggaran berat.
\vspace{1ex}}

\renewcommand{\klausureinverstaendnis}{%
\vspace{2ex}
Dengan tanda tangan saya, saya menyetujui bahwa hasil ujian dapat diumumkan
bersama nomor mahasiswa saya.
\vspace{1ex}

\begin{tabular}{@{}p{3.2cm}p{10cm}}
Tanda tangan: &
.....................................................................................
\end{tabular}}

\renewcommand{\inputaufgabeklausurloesung}[3]{%
\par\newpage
\begin{aufgabe}
{\ifthenelse{\equal{#1}{}}{}{\ifthenelse{\equal{#1}{1}}{(1 poin)}{(#1 poin)}}}%
\par\bigskip #2
\end{aufgabe}
\underline{Solusi:}\par\bigskip #3}

\renewcommand{\inputaufgabeloesung}[2]{%
\begin{aufgabe}#1\end{aufgabe}\aufgabenachskip
\bigskip
Solusi
\bigskip #2}

\renewcommand{\inputaufgabeloesungvar}[2]{%
\aufgabevorskip Soal: #1\aufgabenachskip Solusi: #2}

\renewcommand{\inputaufgabepunkteloesung}[3]{%
\aufgabevorskip
\begin{aufgabe}
{\ifthenelse{\equal{#1}{}}{}{\ifthenelse{\equal{#1}{1}}{(1 poin)}{(#1 poin)}}}%
\aufgabepunktskip #2
\end{aufgabe}
\aufgabenachskip\loesung{#3}}

\renewcommand{\inputaufgabepunktenummervonhand}[3]{%
\aufgabevorskip{\bfseries Soal #1} (#2 poin)\aufgabepunktskip #3\aufgabenachskip}

\renewcommand{\loesung}[1]{\underline{Solusi:} #1}

% Delapan formulir sumber memakai huruf penampung K sebagai nilai penghitung
% bagian. Pertahankan byte formulir dan tafsirkan penampung itu sebagai awal
% penomoran hanya pada lapisan kompilasi pembaca.
\let\OExamOriginalSetCounter\setcounter
\renewcommand{\setcounter}[2]{%
\ifthenelse{\equal{#1}{bagian}}%
  {\OExamOriginalSetCounter{section}{#2}}%
  {\ifthenelse{\equal{#1}{section}}%
  {\ifthenelse{\equal{#2}{K}}{\OExamOriginalSetCounter{section}{0}}{\OExamOriginalSetCounter{#1}{#2}}}%
  {\OExamOriginalSetCounter{#1}{#2}}}}

\newcommand{\OExamPointTable}[4]{%
\vspace{3ex}
\begin{center}
\begingroup
\scriptsize
\renewcommand{\arraystretch}{1.45}%
\setlength{\tabcolsep}{0.35ex}%
\begin{tabular}{@{}|l|*{#1}{c|}c|}
\hline
Soal & #2 & Jumlah\\ \hline
Maks. & #3 & #4\\ \hline
Diperoleh & \multicolumn{#1}{c|}{} & \\ \hline
\end{tabular}
\endgroup
\end{center}
\vspace{2ex}}

\renewcommand{\punktetabelleelf}{%
\OExamPointTable{11}%
{1&2&3&4&5&6&7&8&9&10&11}%
{\aeins&\azwei&\adrei&\avier&\afuenf&\asechs&\asieben&\aacht&\aneun&\azehn&\aelf}%
{\azwoelf}}

\renewcommand{\punktetabelledreizehn}{%
\OExamPointTable{13}%
{1&2&3&4&5&6&7&8&9&10&11&12&13}%
{\aeins&\azwei&\adrei&\avier&\afuenf&\asechs&\asieben&\aacht&\aneun&\azehn&\aelf&\azwoelf&\adreizehn}%
{\avierzehn}}

\renewcommand{\punktetabellevierzehn}{%
\OExamPointTable{14}%
{1&2&3&4&5&6&7&8&9&10&11&12&13&14}%
{\aeins&\azwei&\adrei&\avier&\afuenf&\asechs&\asieben&\aacht&\aneun&\azehn&\aelf&\azwoelf&\adreizehn&\avierzehn}%
{\afuenfzehn}}

\renewcommand{\punktetabellesechzehn}{%
\OExamPointTable{16}%
{1&2&3&4&5&6&7&8&9&10&11&12&13&14&15&16}%
{\aeins&\azwei&\adrei&\avier&\afuenf&\asechs&\asieben&\aacht&\aneun&\azehn&\aelf&\azwoelf&\adreizehn&\avierzehn&\afuenfzehn&\asechzehn}%
{\asiebzehn}}

\renewcommand{\punktetabellesiebzehn}{%
\OExamPointTable{17}%
{1&2&3&4&5&6&7&8&9&10&11&12&13&14&15&16&17}%
{\aeins&\azwei&\adrei&\avier&\afuenf&\asechs&\asieben&\aacht&\aneun&\azehn&\aelf&\azwoelf&\adreizehn&\avierzehn&\afuenfzehn&\asechzehn&\asiebzehn}%
{\aachtzehn}}

\providecommand{\theHfakt}{}

\chapter{Formulir Peserta Ujian 1}
\noindent\textit{Formulir peserta dari sumber resmi; teks soal dipertahankan sesuai terjemahan yang telah diverifikasi.}
\begingroup
\renewcommand{\theHfakt}{o011.exam.learner.01.\arabic{fakt}}
%Data institusi

%\input{Dozentdaten}

%\renewcommand{\fachbereich}{Bidang studi}

%\renewcommand{\dozent}{Prof. Dr. . }

%Data ujian

\renewcommand{\klausurgebiet}{ }

\renewcommand{\klausurtyp}{ }

\renewcommand{\klausurdatum}{. 20}

\klausurvorspann {\fachbereich} {\klausurdatum} {\dozent} {\klausurgebiet} {\klausurtyp}

%Data untuk tabel poin berikut

\renewcommand{\aeins}{  3  }

\renewcommand{\azwei}{  3   }

\renewcommand{\adrei}{  6  }

\renewcommand{\avier}{  5  }

\renewcommand{\afuenf}{  5  }

\renewcommand{\asechs}{  2  }

\renewcommand{\asieben}{  5  }

\renewcommand{\aacht}{  5  }

\renewcommand{\aneun}{  7  }

\renewcommand{\azehn}{  3  }

\renewcommand{\aelf}{  2  }

\renewcommand{\azwoelf}{  5  }

\renewcommand{\adreizehn}{  9   }

\renewcommand{\avierzehn}{  5  }

\renewcommand{\afuenfzehn}{  65  }

\renewcommand{\asechzehn}{    }

\renewcommand{\asiebzehn}{    }

\renewcommand{\aachtzehn}{    }

\renewcommand{\aneunzehn}{    }

\renewcommand{\azwanzig}{    }

\renewcommand{\aeinundzwanzig}{    }

\renewcommand{\azweiundzwanzig}{    }

\renewcommand{\adreiundzwanzig}{    }

\renewcommand{\avierundzwanzig}{    }

\renewcommand{\afuenfundzwanzig}{    }

\renewcommand{\asechsundzwanzig}{    }

\punktetabellevierzehn

\klausurnote

\newpage

\setcounter{section}{0}

\inputaufgabegibtloesung
{3}
{

Definisikan istilah-istilah berikut
\zusatzklammer {yang dicetak miring} {} {}.
\aufzaehlungsechs{\stichwort {Lingkaran kelengkungan} {}
dari suatu
\definitionsverweis {kurva}{}{}
yang dua kali
\definitionsverweis {terdiferensialkan secara kontinu}{}{}
dan
\definitionsverweis {berparametrisasi panjang busur}{}{}.
Misalkan
\maabbdisp {\gamma} {I} {\R^2
} {}
dan ambil titik

\mavergleichskette
{\vergleichskette
{ t_0
}
{ \in }{ I
}
{  }{
}
{    }{
}
{   }{
}
}
{}{}{.}

}{Suatu \stichwort {submanifold tertutup} {}
\mathl{M \subseteq N}{} dari manifold diferensiabel $N$.

}{\stichwort {Ruang singgung} {} pada titik
\mathl{P\in M}{} dari manifold diferensiabel $M$.

}{Istilah
\stichwort {berorientasi} {}
untuk suatu
\definitionsverweis {manifold diferensiabel}{}{}
$M$.

}{\stichwort {Batas} {}
suatu
\definitionsverweis {manifold berbatas}{}{.}

}{\stichwort {Simbol Christoffel} {}

\mathl{\Gamma_{ij}^k}{} bagi koneksi Levi-Civita dari struktur Riemann $\left(  g_{ij}  \right)$ pada himpunan terbuka

\mavergleichskette
{\vergleichskette
{ U
}
{ \subseteq }{ \R^n
}
{  }{
}
{    }{
}
{   }{
}
}
{}{}{.}
}

}
{} {}

\inputaufgabegibtloesung
{3}
{

Nyatakan teorema-teorema berikut.
\aufzaehlungdrei{Teorema tentang citra
\definitionsverweis {pemetaan Gauss}{}{}
dari suatu hipermuka.}{Rumus luas permukaan dari permukaan putar yang dibentuk oleh
\definitionsverweis {kurva diferensiabel}{}{}
\maabbeledisp {\gamma} {]a,b[} {\R^2
} { t } {(x(t),y(t))
} {,}
dengan

\mavergleichskette
{\vergleichskette
{  y(t)
}
{ > }{  0
}
{  }{
}
{    }{
}
{   }{
}
}
{}{}{.}}{\stichwort {Aturan hasil kali} {}
untuk turunan eksterior.}

}
{} {}

\inputaufgabegibtloesung
{6 (1+2+3)}
{

Misalkan $Y$ adalah hipermuka yang diberikan oleh persamaan

\mavergleichskettedisp
{\vergleichskette
{ 2x^2 +3y^2+5z^2
}
{ =} { 10
}
{ } {
}
{ } {
}
{ } {
}
}
{}{}{}
di $\R^3$. Misalkan

\mavergleichskette
{\vergleichskette
{ P
}
{ = }{ \begin{pmatrix}  -1 \\ -1\\  1   \end{pmatrix}
}
{ \in }{ Y
}
{    }{
}
{   }{
}
}
{}{}{}
dan

\mavergleichskette
{\vergleichskette
{ w
}
{ = }{ \begin{pmatrix}  -3 \\ 2 \\  4   \end{pmatrix}
}
{ \in }{ \R^3
}
{    }{
}
{   }{
}
}
{}{}{.}
\aufzaehlungdrei{Tentukan
\definitionsverweis {ruang singgung}{}{}

\mathl{T_PY}{.}
}{Tentukan
\definitionsverweis {dekomposisi ortogonal}{}{}
$w$ menjadi komponen tangensial dan komponen ortogonal pada titik $P$.
}{Realisasikan komponen tangensial $w$ di $T_PY$ melalui suatu kurva diferensiabel yang seluruhnya terletak pada $Y$.
}

}
{} {}

\inputaufgabegibtloesung
{5}
{

Misalkan
\maabbeledisp {\gamma} {\R} { \R^2
} { t } { \begin{pmatrix}  \cos \left(  \omega t \right)  \\ \sin  \left(  \omega t \right)    \end{pmatrix}
} {,}
dengan

\mavergleichskette
{\vergleichskette
{ \omega
}
{ > }{ 0
}
{  }{
}
{    }{
}
{   }{
}
}
{}{}{.}
Tunjukkan bahwa tidak ada gerak melingkar lain dengan norma kecepatan konstan yang memiliki nilai dan turunan hingga orde kedua yang sama dengan $\gamma$ di titik parameter

\mavergleichskette
{\vergleichskette
{ t
}
{ = }{ 0
}
{  }{
}
{    }{
}
{   }{
}
}
{}{}{}
tersebut.

}
{} {}

\inputaufgabegibtloesung
{5}
{

Buktikan teorema tentang kelengkungan suatu kurva bidang yang diberikan secara implisit.

}
{} {}

\inputaufgabegibtloesung
{2}
{

Berikan contoh manifold diferensiabel
\mathkor {} {M} {dan} {N} {}
dengan

\mavergleichskette
{\vergleichskette
{ \operatorname{ dim}_{  } ^{  }  \left(   M   \right), \, \operatorname{ dim}_{  } ^{  }  \left(   N   \right)
}
{ \geq }{ 1
}
{  }{
}
{    }{
}
{   }{
}
}
{}{}{}
serta pemetaan diferensiabel
\maabb {\varphi} { M } { N
} {}
dan
\maabb {\psi} { N } { M
} {}
sedemikian sehingga

\mavergleichskette
{\vergleichskette
{ \psi \circ \varphi
}
{ = }{
\operatorname{Id}_{ M }
}
{  }{
}
{    }{
}
{   }{
}
}
{}{}{}
dan

\mavergleichskette
{\vergleichskette
{\varphi \circ  \psi
}
{ \neq }{
\operatorname{Id}_{ N }
}
{  }{
}
{    }{
}
{   }{
}
}
{}{}{}
berlaku.

}
{} {}

\inputaufgabegibtloesung
{5}
{

Buktikan bahwa ekuivalensi tangensial lintasan-lintasan pada suatu manifold di titik $P$ dapat diperiksa dengan menggunakan sebarang peta koordinat.

}
{} {}

\inputaufgabegibtloesung
{5}
{

Tunjukkan bahwa pemetaan singgung $T(\varphi)$ dari
\maabbeledisp {\varphi} {\R^1} {S^1
} { t } {( \cos  t, \sin  t  )
} {,}
bersifat surjektif.

}
{} {}

\inputaufgabegibtloesung
{7 (3+4)}
{

Tinjau pemetaan diferensiabel
\maabbeledisp {\gamma} { [1,c] } {\R^2
} { t } {(t,t^{3})
} {,}
\zusatzklammer {dengan \mathlk{c \geq 1}{}} {} {}
dan
\maabbeledisp {\varphi} {\R_+ \times \R_+ } {\R^3
} {(u,v)} {(u^3,u^2+v^2,u^{-1}v^{-1} )
} {,}
serta bentuk diferensial

\mavergleichskettedisp
{\vergleichskette
{ \omega
}
{ =} { (x-y)dx-z^2dy+dz
}
{ } {
}
{ } {
}
{ } {
}
}
{}{}{}
pada $\R^3$.

a) Hitung bentuk diferensial tarik balik
\mathl{\varphi^* \omega}{} pada $\R_+ \times \R_+$.

b) Hitung integral lintasan dari bentuk diferensial
\mathl{\varphi^* \omega}{} sepanjang lintasan $\gamma$ sebagai fungsi dari $c$.

}
{} {}

\inputaufgabegibtloesung
{3}
{

Misalkan
\maabbdisp {\psi} { L } { M
} {}
adalah
\definitionsverweis {pemetaan diferensiabel}{}{}
antara
\definitionsverweis {manifold diferensiabel}{}{.}
Selanjutnya, misalkan
\maabbdisp {f} { M } {\R
} {}
adalah fungsi pada $M$, dan $\omega$ adalah
$k$-\definitionsverweis {bentuk}{}{}
pada $M$. Tunjukkan bahwa

\mavergleichskettedisp
{\vergleichskette
{ \psi^*(f \omega)
}
{ =} { \psi^*(f)  \psi^* \omega
}
{ } {
}
{ } {
}
{ } {
}
}
{}{}{.}

}
{} {}

\inputaufgabe
{2}
{

Deskripsikan sebanyak mungkin
\definitionsverweis {manifold berbatas}{}{}
berdimensi satu yang
\definitionsverweis {terhubung}{}{.}

}
{} {}

\inputaufgabegibtloesung
{5}
{

Misalkan
\maabbdisp {f} { [a,b] } { \R_{+}
} {}
adalah
\definitionsverweis {fungsi yang terdiferensialkan secara kontinu}{}{.}
Kita pandang
\definitionsverweis {subgraf}{}{}
sebagai
\definitionsverweis {manifold berbatas}{}{}.
Batasnya terdiri dari grafik, interval dasar, dan kedua sisi tegak, tetapi tidak mencakup keempat titik sudut. Buktikan secara langsung berlakunya teorema Stokes untuk
\definitionsverweis {bentuk diferensial}{}{}

\mavergleichskette
{\vergleichskette
{ \omega
}
{ = }{ xdy
}
{  }{
}
{    }{
}
{   }{
}
}
{}{}{.}

}
{} {}

\inputaufgabegibtloesung
{9}
{

Buktikan teorema titik tetap Brouwer.

}
{} {}

\inputaufgabegibtloesung
{5 (3+1+1)}
{

Tinjau silinder

\mavergleichskette
{\vergleichskette
{ Z
}
{ = }{ S^1 \times \R
}
{ \subseteq }{ \R^2 \times \R
}
{    }{
}
{   }{
}
}
{}{}{}
dengan
\definitionsverweis {koneksi trivial}{}{}
pada
\definitionsverweis {bundel tangen}{}{.}
\aufzaehlungdrei{Tentukan, dalam suatu peta isometrik yang sesuai,
\definitionsverweis {persamaan diferensial geodesik}{}{}
untuk kurva geodesik $\psi$ pada $M$ yang memenuhi

\mavergleichskettedisp
{\vergleichskette
{ \psi(0)
}
{ =} { (1,0,3)
}
{ } {
}
{ } {
}
{ } {
}
}
{}{}{}
dan

\mavergleichskettedisp
{\vergleichskette
{ \psi'(0)
}
{ =} { (0,1,-4)
}
{ } {
}
{ } {
}
{ } {
}
}
{}{}{}
.
}{Selesaikan persamaan diferensial dari (1) dalam peta tersebut.
}{Carilah suatu kurva geodesik di $Z$ yang memenuhi syarat-syarat pada (1).
}

}
{} {}

\endgroup

\chapter{Formulir Peserta Ujian 2}
\noindent\textit{Formulir peserta dari sumber resmi; teks soal dipertahankan sesuai terjemahan yang telah diverifikasi.}
\begingroup
\renewcommand{\theHfakt}{o011.exam.learner.02.\arabic{fakt}}
%Data institusi

%\input{Dozentdaten}

%\renewcommand{\fachbereich}{Fachbereich}

%\renewcommand{\dozent}{Prof. Dr. . }

%Data ujian

\renewcommand{\klausurgebiet}{ }

\renewcommand{\klausurtyp}{ }

\renewcommand{\klausurdatum}{ . 20}

\klausurvorspann {\fachbereich} {\klausurdatum} {\dozent} {\klausurgebiet} {\klausurtyp}

%Data untuk tabel nilai berikut

\renewcommand{\aeins}{  3  }

\renewcommand{\azwei}{  3   }

\renewcommand{\adrei}{  8  }

\renewcommand{\avier}{  7  }

\renewcommand{\afuenf}{  8  }

\renewcommand{\asechs}{  3  }

\renewcommand{\asieben}{  5  }

\renewcommand{\aacht}{  5  }

\renewcommand{\aneun}{  6  }

\renewcommand{\azehn}{  4  }

\renewcommand{\aelf}{  12  }

\renewcommand{\azwoelf}{  64  }

\renewcommand{\adreizehn}{     }

\renewcommand{\avierzehn}{    }

\renewcommand{\afuenfzehn}{    }

\renewcommand{\asechzehn}{    }

\renewcommand{\asiebzehn}{    }

\renewcommand{\aachtzehn}{    }

\renewcommand{\aneunzehn}{    }

\renewcommand{\azwanzig}{    }

\renewcommand{\aeinundzwanzig}{    }

\renewcommand{\azweiundzwanzig}{    }

\renewcommand{\adreiundzwanzig}{    }

\renewcommand{\avierundzwanzig}{    }

\renewcommand{\afuenfundzwanzig}{    }

\renewcommand{\asechsundzwanzig}{    }

\punktetabelleelf

\klausurnote

\newpage

\setcounter{section}{0}

\inputaufgabegibtloesung
{3}
{

Definisikan istilah-istilah berikut
\zusatzklammer {yang dicetak miring} {} {.}
\aufzaehlungsechs{Suatu
\stichwort {medan normal} {}
$N$ pada hiperpermukaan diferensiabel

\mavergleichskette
{\vergleichskette
{ Y
}
{ \subseteq }{ \R^n
}
{  }{
}
{    }{
}
{   }{
}
}
{}{}{.}

}{Suatu ruang topologis yang
\stichwort {parakompak}{}
.

}{
\stichwort {Ruang kotangen} {}
di suatu titik

\mavergleichskette
{\vergleichskette
{  P
}
{ \in }{ M
}
{  }{}
{    }{}
{   }{}
}
{}{}{}
dari
\definitionsverweis {manifold diferensiabel}{}{}
$M$.

}{Suatu
\stichwort {titik reguler} {}

\mathl{Q \in L}{} bagi
\definitionsverweis {pemetaan diferensiabel}{}{}
\maabbdisp {\varphi} { L } { M
} {}
antara
\definitionsverweis {manifold-manifold diferensiabel}{}{}
\mathkor {} {L} {dan} {M} {.}

}{
\stichwort {Produk} {}
dari manifold-manifold diferensiabel
\mathkor {} {M} {dan} {N} {.}

}{Suatu
\stichwort {koneksi} {}
pada bundel vektor diferensiabel
\maabb {} {E} {M
} {}
di atas manifold diferensiabel $M$.
}

}
{} {}

\inputaufgabegibtloesung
{3}
{

Nyatakan teorema-teorema berikut.
\aufzaehlungdrei{
\stichwort {Teorema tentang hubungan antara kelengkungan dan vektor normal satuan suatu kurva bidang berparameter panjang busur} {.}}{
\stichwort {Theorema egregium} {.}}{
\stichwort {Teorema Green} {}
untuk luas.}

}
{} {}

\inputaufgabegibtloesung
{8}
{

Misalkan

\mavergleichskettedisp
{\vergleichskette
{ f(x,y)
}
{ =} { ax^2+by^2+cxy
}
{ } {
}
{ } {
}
{ } {
}
}
{}{}{.}
Tentukan
\definitionsverweis {kelengkungan-kelengkungan utama}{}{}
dan
\definitionsverweis {arah-arah kelengkungan utama}{}{}
dari
\definitionsverweis {grafik}{}{}
$f$ di titik asal.

}
{} {}

\inputaufgabegibtloesung
{7}
{

Buktikan teorema eksistensi untuk transpor paralel pada suatu hiperpermukaan

\mavergleichskette
{\vergleichskette
{ Y
}
{ \subseteq }{ \R^n
}
{  }{
}
{    }{
}
{   }{
}
}
{}{}{.}

}
{} {}

\inputaufgabegibtloesung
{8}
{

Tunjukkan bahwa suatu
\definitionsverweis {manifold diferensiabel}{}{}
$M$ berdimensi

\mavergleichskette
{\vergleichskette
{ n
}
{ \geq }{ 1
}
{  }{
}
{    }{
}
{   }{
}
}
{}{}{}
mempunyai tak hingga banyak
\definitionsverweis {difeomorfisme}{}{}
\maabbdisp {\varphi} { M } { M
} {}
.

}
{} {}

\inputaufgabegibtloesung
{3}
{

Misalkan $M$ suatu
\definitionsverweis {manifold diferensiabel}{}{}
dan
\maabbdisp {f} { M } {\R
} {}
suatu
\definitionsverweis {fungsi terdiferensialkan secara kontinu}{}{.}
Misalkan di titik

\mavergleichskette
{\vergleichskette
{ P
}
{ \in }{ M
}
{  }{
}
{    }{
}
{   }{
}
}
{}{}{}
fungsi tersebut mempunyai
\definitionsverweis {ekstremum lokal}{}{.}
Tunjukkan bahwa

\mavergleichskettedisp
{\vergleichskette
{ (df)_P
}
{ =} { 0
}
{ } {
}
{ } {
}
{ } {
}
}
{}{}{.}

}
{} {}

\inputaufgabegibtloesung
{5}
{

Kita tinjau
$1$-\definitionsverweis {bentuk diferensial}{}{}

\mavergleichskettedisp
{\vergleichskette
{ \omega
}
{ =} { -vdu+udv
}
{ } {
}
{ } {
}
{ } {
}
}
{}{}{}
pada
$1$-\definitionsverweis {bola}{}{}

\mavergleichskette
{\vergleichskette
{ S^1
}
{ \subset }{ \R^2
}
{  }{
}
{    }{
}
{   }{
}
}
{}{}{,}
dengan koordinat ruang sekitar $\R^2$ dinotasikan oleh
\mathkor {} {u} {dan} {v} {.}
Tentukan
\mathl{\pi^* \omega}{} di bawah
\definitionsverweis {pemetaan terdiferensialkan secara kontinu}{}{}
\maabbeledisp {\pi} {\R^2 \setminus \{0\}} { S^{1}
} {(x , y) } { { \frac{   1  }{  \sqrt{x^2 + y^2}  } } (x , y)
} {.}

}
{} {}

\inputaufgabegibtloesung
{5}
{

Buktikan teorema tentang perhitungan volume kanonik pada manifold Riemann $M$.

}
{} {}

\inputaufgabegibtloesung
{6 (1+1+3+1)}
{

Jelaskan mengapa $\R^2$ tidak membawa struktur
\definitionsverweis {manifold berbatas}{}{}
sedemikian sehingga subhimpunan $T$ yang diberikan merupakan
\definitionsverweis {batas}{}{.}
\aufzaehlungvierabc{
\mavergleichskette
{\vergleichskette
{ T
}
{ = }{ {]0,1[}
}
{  }{
}
{    }{
}
{   }{
}
}
{}{}{,}
dengan interval tersebut terletak pada sumbu $x$.
}{
\mavergleichskette
{\vergleichskette
{ T
}
{ = }{ [0,1]
}
{  }{
}
{    }{
}
{   }{
}
}
{}{}{,}
dengan interval tersebut terletak pada sumbu $x$.
}{
\mavergleichskette
{\vergleichskette
{ T
}
{ = }{ \R
}
{  }{
}
{    }{
}
{   }{
}
}
{}{}{,}
dengan $\R$ adalah sumbu $x$.
}{
\mavergleichskette
{\vergleichskette
{ T
}
{ = }{ S^1
}
{  }{
}
{    }{
}
{   }{
}
}
{}{}{.}
}

}
{} {}

\inputaufgabegibtloesung
{4}
{

Misalkan $\omega$ suatu
$n-1$-\definitionsverweis {bentuk diferensial}{}{}
yang terdiferensialkan secara kontinu pada
\definitionsverweis {himpunan terbuka}{}{}

\mavergleichskette
{\vergleichskette
{ U
}
{ \subseteq }{ \R^n
}
{  }{
}
{    }{
}
{   }{
}
}
{}{}{}
dan misalkan
\definitionsverweis {turunan eksterior}{}{}
sama dengan

\mavergleichskettedisp
{\vergleichskette
{d \omega
}
{ =} { fdx_1 \wedge \ldots \wedge dx_n
}
{ } {
}
{ } {
}
{ } {
}
}
{}{}{}
dengan suatu fungsi
\maabb {f} { U } {\R
} {.}
Tunjukkan untuk

\mavergleichskette
{\vergleichskette
{ P
}
{ \in }{ U
}
{  }{
}
{    }{
}
{   }{
}
}
{}{}{}
kesamaan

\mavergleichskettedisp
{\vergleichskette
{ f(P)
}
{ =} { { \frac{   1  }{  \beta_n  } } \cdot \operatorname{lim}_{   \epsilon \rightarrow  0  } \, { \frac{   1  }{  \epsilon^n } } \int_{S^{n-1}(P, \epsilon)} \omega
}
{ } {
}
{ } {
}
{ } {
}
}
{}{}{,}
dengan $\beta_n$ menyatakan
volume
bola satuan berdimensi $n$.

}
{} {}

\inputaufgabegibtloesung
{12}
{

Buktikan teorema Stokes untuk manifold berbatas.

}
{} {}

\endgroup

\chapter{Formulir Peserta Ujian 3}
\noindent\textit{Formulir peserta dari sumber resmi; teks soal dipertahankan sesuai terjemahan yang telah diverifikasi.}
\begingroup
\renewcommand{\theHfakt}{o011.exam.learner.03.\arabic{fakt}}
%Data institusi

%\input{Dozentdaten}

%\renewcommand{\fachbereich}{Fachbereich}

%\renewcommand{\dozent}{Prof. Dr. . }

%Data ujian

\renewcommand{\klausurgebiet}{ }

\renewcommand{\klausurtyp}{ }

\renewcommand{\klausurdatum}{ . 20}

\klausurvorspann {\fachbereich} {\klausurdatum} {\dozent} {\klausurgebiet} {\klausurtyp}

%Data untuk tabel nilai berikut

\renewcommand{\aeins}{  3  }

\renewcommand{\azwei}{  3   }

\renewcommand{\adrei}{  0  }

\renewcommand{\avier}{  0  }

\renewcommand{\afuenf}{  6  }

\renewcommand{\asechs}{  6  }

\renewcommand{\asieben}{  0  }

\renewcommand{\aacht}{  6  }

\renewcommand{\aneun}{  11  }

\renewcommand{\azehn}{  5  }

\renewcommand{\aelf}{  5  }

\renewcommand{\azwoelf}{  3  }

\renewcommand{\adreizehn}{  7   }

\renewcommand{\avierzehn}{  2  }

\renewcommand{\afuenfzehn}{  0  }

\renewcommand{\asechzehn}{  6  }

\renewcommand{\asiebzehn}{  63  }

\renewcommand{\aachtzehn}{    }

\renewcommand{\aneunzehn}{    }

\renewcommand{\azwanzig}{    }

\renewcommand{\aeinundzwanzig}{    }

\renewcommand{\azweiundzwanzig}{    }

\renewcommand{\adreiundzwanzig}{    }

\renewcommand{\avierundzwanzig}{    }

\renewcommand{\afuenfundzwanzig}{    }

\renewcommand{\asechsundzwanzig}{    }

\punktetabellesechzehn

\klausurnote

\newpage

\setcounter{section}{K}

\inputaufgabegibtloesung
{3}
{

Definisikan istilah-istilah berikut
\zusatzklammer {yang dicetak miring} {} {.}
\aufzaehlungsechs{Suatu
\stichwort {kurva geodesik} {}
pada hiperpermukaan diferensiabel

\mavergleichskette
{\vergleichskette
{ Y
}
{ \subseteq }{ \R^n
}
{  }{
}
{    }{
}
{   }{
}
}
{}{}{.}

}{
\stichwort {Himpunan rotasi} {}
\zusatzklammer {mengelilingi sumbu $x$} {} {}
dari suatu subhimpunan

\mavergleichskette
{\vergleichskette
{ T
}
{ \subseteq }{  \R \times \R_{\geq 0}
}
{  }{
}
{    }{
}
{   }{
}
}
{}{}{.}

}{
\stichwort {Pemetaan tangen} {}
\maabbdisp {T (\varphi)} {TM} {TN
} {}
dari
\definitionsverweis {pemetaan diferensiabel}{}{}
\maabbdisp {\varphi} { M } { N
} {}
antara
\definitionsverweis {manifold-manifold diferensiabel}{}{}
\mathkor {} {M} {dan} {N} {.}

}{Suatu perubahan peta yang
\stichwort {mempertahankan orientasi} {}
pada
\definitionsverweis {manifold diferensiabel}{}{}
$M$.

}{
\stichwort {Bentuk diferensial tarikan balik} {} $\varphi^* \omega$ dari suatu bentuk diferensial
\mathl{\omega \in
{ \mathcal E }^{ k } (  M   )}{} terhadap suatu pemetaan terdiferensialkan secara kontinu
\maabb {\varphi} { L } { M
} {}
antara dua manifold diferensiabel
\mathkor {} {L} {dan} {M} {.}

}{Suatu koneksi linear yang
\stichwort {metrik} {}
pada bundel tangen suatu manifold Riemann $M$.
}

}
{} {}

\inputaufgabegibtloesung
{3}
{

Nyatakan teorema-teorema berikut.
\aufzaehlungdrei{
\stichwort {Teorema tentang solusi persamaan diferensial biasa pada hiperpermukaan diferensiabel} {}

\mavergleichskette
{\vergleichskette
{  Y
}
{ \subseteq }{  \R^n
}
{  }{
}
{    }{
}
{   }{
}
}
{}{}{.}}{Rumus untuk menghitung volume kanonik suatu manifold Riemann dalam sebuah peta.}{Teorema tentang partisi kesatuan.}

}
{} {}

\inputaufgabe
{0}
{

}
{} {}

\inputaufgabe
{0}
{

}
{} {}

\inputaufgabegibtloesung
{6}
{

Untuk permukaan $Y$ yang diberikan oleh

\mavergleichskettedisp
{\vergleichskette
{ x^2+y^2+2z^2
}
{ =} { 4
}
{ } {
}
{ } {
}
{ } {
}
}
{}{}{}
dan titik

\mavergleichskettedisp
{\vergleichskette
{ P
}
{ =} { \left(  1  , \,   1  , \,   1                 \right)
}
{ \in} { Y
}
{ } {
}
{ } {
}
}
{}{}{,}
tentukan suatu matriks diagonal bagi
\definitionsverweis {pemetaan Weingarten}{}{}
$L_P$.

}
{} {}

\inputaufgabegibtloesung
{6}
{

Buktikan teorema isometri untuk transpor paralel pada hiperpermukaan

\mavergleichskette
{\vergleichskette
{ Y
}
{ \subseteq }{ \R^n
}
{  }{
}
{    }{
}
{   }{
}
}
{}{}{.}

}
{} {}

\inputaufgabe
{0}
{

}
{} {}

\inputaufgabegibtloesung
{6}
{

Misalkan $M$ suatu manifold topologis kompak berdimensi $d$,

\mavergleichskette
{\vergleichskette
{ d
}
{ \geq }{ 1
}
{  }{
}
{    }{
}
{   }{
}
}
{}{}{.}
Tunjukkan bahwa terdapat suatu subhimpunan terbuka terbatas

\mavergleichskette
{\vergleichskette
{ U
}
{ \subseteq }{ \R^d
}
{  }{
}
{    }{
}
{   }{
}
}
{}{}{}
dan suatu pemetaan kontinu surjektif
\maabbdisp {\varphi} { U } { M
} {}
.

}
{} {}

\inputaufgabegibtloesung
{11 (1+3+1+2+4)}
{

Kita tinjau himpunan

\mavergleichskettedisp
{\vergleichskette
{N
}
{ =} { \left\{ A= \begin{pmatrix}  a   &   b   \\  c  & d \end{pmatrix}   \mid   A \text{ nilpoten}         \right\}
}
{ \subseteq} {\R^4
}
{ } {
}
{ } {
}
}
{}{}{}
dari
$2\times 2$-\definitionsverweis {matriks}{}{}
real nilpoten, serta himpunan

\mavergleichskettedisp
{\vergleichskette
{M
}
{ =} { N \setminus \{ \begin{pmatrix}  0   &   0  \\  0  &  0 \end{pmatrix}  \}
}
{ } {
}
{ } {
}
{ } {
}
}
{}{}{.}

a) Apakah $N$ terhubung?

b) Tunjukkan bahwa $M$ suatu
\definitionsverweis {submanifold tertutup}{}{}
dari suatu subhimpunan terbuka

\mavergleichskette
{\vergleichskette
{ G
}
{ \subseteq }{ \R^4
}
{  }{
}
{    }{
}
{   }{
}
}
{}{}{.}

c) Tentukan
\definitionsverweis {dimensi}{}{}
$M$.

d) Apakah $M$ terhubung?

e) Tutupi
\mathl{M}{} dengan peta-peta topologis eksplisit.

}
{} {}

\inputaufgabe
{5}
{

Katakan sesuatu yang cerdas tentang pita Möbius.

}
{} {}

\inputaufgabegibtloesung
{5}
{

Kita tinjau pemetaan
\maabbeledisp {\varphi} { \R^3 } { \R^2
} { (x,y,z) } { \left(  xyz^2,xy-z^3  \right)
} {.}
Hitung matriks pemetaan
\maabbdisp {\bigwedge^2 T_P (\varphi)} { \bigwedge^2 T_P \R^3 } { \bigwedge^2 T_{\varphi(P)}  \R^2
} {}
di titik

\mavergleichskette
{\vergleichskette
{ P
}
{ = }{ (1,3,5)
}
{  }{
}
{    }{
}
{   }{
}
}
{}{}{}
terhadap suatu basis yang sesuai.

}
{} {}

\inputaufgabegibtloesung
{3}
{

Misalkan

\mavergleichskette
{\vergleichskette
{ W
}
{ \subseteq }{ \R^m
}
{  }{
}
{    }{
}
{   }{
}
}
{}{}{}
dan

\mavergleichskette
{\vergleichskette
{ U
}
{ \subseteq }{ \R^n
}
{  }{
}
{    }{
}
{   }{
}
}
{}{}{}
merupakan
\definitionsverweis {subhimpunan-subhimpunan terbuka}{}{}
dan misalkan
\maabbdisp {\psi} { W } { U
} {}
suatu
\definitionsverweis {pemetaan}{}{}
yang
\definitionsverweis {terdiferensialkan secara kontinu}{}{.}
Misalkan
\maabbdisp {f} { U } {\R
} {}
suatu fungsi terdiferensialkan secara kontinu. Simpulkan dari
aturan rantai
bahwa

\mavergleichskettedisp
{\vergleichskette
{ d (\psi^*f)
}
{ =} { \psi^*(df)
}
{ } {
}
{ } {
}
{ } {
}
}
{}{}{,}
dengan $\psi^*$ menyatakan
\definitionsverweis {penarikan balik bentuk diferensial}{}{.}

}
{} {}

\inputaufgabegibtloesung
{7}
{

Misalkan $M$ suatu
\definitionsverweis {manifold diferensiabel}{}{}
dengan
\definitionsverweis {basis terhitung bagi topologinya}{}{}
dan misalkan $\omega$ suatu
\definitionsverweis {bentuk volume positif}{}{}
pada $M$. Misalkan
\maabbdisp {\varphi} { L } { M
} {}
suatu
\definitionsverweis {difeomorfisme}{}{}
dengan manifold $L$ dan

\mavergleichskette
{\vergleichskette
{ S
}
{ \subseteq }{ L
}
{  }{
}
{    }{
}
{   }{
}
}
{}{}{}
suatu
\definitionsverweis {subhimpunan terukur}{}{.}
Tunjukkan bahwa

\mavergleichskettedisp
{\vergleichskette
{ \int_S \varphi^* \omega
}
{ =} { \int_{ \varphi(S)} \omega
}
{ } {
}
{ } {
}
{ } {
}
}
{}{}{.}

}
{} {}

\inputaufgabegibtloesung
{2}
{

Pada
\definitionsverweis {bundel vektor trivial}{}{}
\maabbdisp {p} {\R^2 \times \R} { \R^2
} {}
ber-
\definitionsverweis {rank}{}{}
$1$ di atas $\R^2$, kita tinjau
\definitionsverweis {koneksi linear}{}{,}
yang diberikan oleh
\definitionsverweis {simbol-simbol Christoffel}{}{}

\mavergleichskette
{\vergleichskette
{ \Gamma_1 (x,y)
}
{ = }{ x^5-x^2y^2+3y^4
}
{  }{
}
{    }{
}
{   }{
}
}
{}{}{}
dan

\mavergleichskette
{\vergleichskette
{ \Gamma_2 (x,y)
}
{ = }{ 7x^3 +xy^2-4y^5
}
{  }{
}
{    }{
}
{   }{
}
}
{}{}{.}
Hitung
\definitionsverweis {operator kelengkungan}{}{}
$R \left(  \partial_1, \partial_2  \right)$.

}
{} {}

\inputaufgabe
{0}
{

}
{} {}

\inputaufgabegibtloesung
{6}
{

Buktikan teorema tentang retraksi ke batas pada manifold.

}
{} {}

\endgroup

\chapter{Formulir Peserta Ujian 4}
\noindent\textit{Formulir peserta dari sumber resmi; teks soal dipertahankan sesuai terjemahan yang telah diverifikasi.}
\begingroup
\renewcommand{\theHfakt}{o011.exam.learner.04.\arabic{fakt}}
%Data institusi

%\input{Dozentdaten}

%\renewcommand{\fachbereich}{Fachbereich}

%\renewcommand{\dozent}{Prof. Dr. . }

%Data ujian

\renewcommand{\klausurgebiet}{ }

\renewcommand{\klausurtyp}{ }

\renewcommand{\klausurdatum}{ . 20}

\klausurvorspann {\fachbereich} {\klausurdatum} {\dozent} {\klausurgebiet} {\klausurtyp}

%Data untuk tabel nilai berikut

\renewcommand{\aeins}{  3  }

\renewcommand{\azwei}{  3   }

\renewcommand{\adrei}{  0  }

\renewcommand{\avier}{  0  }

\renewcommand{\afuenf}{  4  }

\renewcommand{\asechs}{  5  }

\renewcommand{\asieben}{  4  }

\renewcommand{\aacht}{  0  }

\renewcommand{\aneun}{  4  }

\renewcommand{\azehn}{  6  }

\renewcommand{\aelf}{  2  }

\renewcommand{\azwoelf}{  6  }

\renewcommand{\adreizehn}{  4   }

\renewcommand{\avierzehn}{  10  }

\renewcommand{\afuenfzehn}{  10  }

\renewcommand{\asechzehn}{  0  }

\renewcommand{\asiebzehn}{  61  }

\renewcommand{\aachtzehn}{    }

\renewcommand{\aneunzehn}{    }

\renewcommand{\azwanzig}{    }

\renewcommand{\aeinundzwanzig}{    }

\renewcommand{\azweiundzwanzig}{    }

\renewcommand{\adreiundzwanzig}{    }

\renewcommand{\avierundzwanzig}{    }

\renewcommand{\afuenfundzwanzig}{    }

\renewcommand{\asechsundzwanzig}{    }

\punktetabellesechzehn

\klausurnote

\newpage

\setcounter{section}{K}

\inputaufgabegibtloesung
{3}
{

Definisikan istilah-istilah berikut
\zusatzklammer {yang dicetak miring} {} {}.
\aufzaehlungsechs{Pemetaan
\stichwort {Weingarten} {}
pada $T_PY$ di suatu titik

\mavergleichskette
{\vergleichskette
{ P
}
{ \in }{ Y
}
{  }{
}
{    }{
}
{   }{
}
}
{}{}{}
dari suatu hiperpermukaan diferensiabel berorientasi

\mavergleichskette
{\vergleichskette
{ Y
}
{ \subseteq }{ \R^n
}
{  }{
}
{    }{
}
{   }{
}
}
{}{}{.}

}{Suatu
\stichwort {bundel vektor riil} {}
$V$ ber-rank $r$ di atas suatu
\definitionsverweis {ruang topologis}{}{}
$X$.

}{\stichwort {Integral lintasan} {} dari suatu bentuk diferensial-$1$
\mathl{\omega \in
{ \mathcal E }^{  1 } (  M   )}{} pada suatu manifold diferensiabel $M$ terhadap suatu kurva yang terdiferensialkan secara kontinu
\maabb {\gamma} {[a,b]} {M
} {.}

}{Suatu
\stichwort {bentuk volume positif} {}
$\omega$ pada suatu manifold diferensiabel $M$.

}{Suatu \stichwort {bentuk diferensial tertutup} {} $\omega$ pada suatu manifold diferensiabel $M$.

}{\stichwort {Percepatan tangensial} {}
dari suatu kurva yang dua kali terdiferensialkan secara kontinu
\maabb {\gamma} {I} {M
} {}
pada suatu
\definitionsverweis {manifold Riemann}{}{}
$M$.
}

}
{} {}

\inputaufgabegibtloesung
{3}
{

Rumuskan teorema-teorema berikut.
\aufzaehlungdrei{\stichwort {Teorema tentang kelengkungan normal} {.}}{Rumus untuk menghitung volume kanonik pada suatu manifold Riemann di dalam suatu peta.}{\stichwort {Teorema Stokes} {} untuk manifold dengan batas.}

}
{} {}

\inputaufgabe
{0}
{

}
{} {}

\inputaufgabe
{0}
{

}
{} {}

\inputaufgabegibtloesung
{4}
{

Buktikan teorema tentang hubungan antara kelengkungan dan vektor normal satuan dari suatu kurva bidang yang diparameterkan oleh panjang busur.

}
{} {}

\inputaufgabegibtloesung
{5 (1+2+2)}
{

Pada
\definitionsverweis {permukaan bola satuan}{}{}
$Y$ berdimensi dua, kita meninjau
\definitionsverweis {medan vektor tangensial}{}{}

\mavergleichskettedisp
{\vergleichskette
{ F \begin{pmatrix}  x  \\ y \\ z   \end{pmatrix}
}
{ =} { \begin{pmatrix}  y  \\ -x\\  0   \end{pmatrix}
}
{ } {
}
{ } {
}
{ } {
}
}
{}{}{.}
\aufzaehlungdrei{Tentukan
\definitionsverweis {matriks Jacobi}{}{}
dari $F$ pada $\R^3$.
}{Untuk titik

\mavergleichskette
{\vergleichskette
{ P
}
{ = }{ \begin{pmatrix}  0  \\ 1 \\  0   \end{pmatrix}
}
{ \in }{ Y
}
{    }{
}
{   }{
}
}
{}{}{}
tentukan suatu matriks untuk pemetaan
\maabbeledisp {} {T_PY} { T_PY
} { v } { \nabla_v F
} {,}
terhadap suatu
\definitionsverweis {basis}{}{}
yang sesuai.
}{Untuk titik

\mavergleichskette
{\vergleichskette
{ Q
}
{ = }{ \begin{pmatrix}  0  \\ 0\\  1   \end{pmatrix}
}
{ \in }{ Y
}
{    }{
}
{   }{
}
}
{}{}{}
tentukan suatu matriks untuk pemetaan
\maabbeledisp {} {T_QY} { T_QY
} { v } { \nabla_v F
} {,}
terhadap suatu
\definitionsverweis {basis}{}{}
yang sesuai.
}

}
{} {}

\inputaufgabegibtloesung
{4 (2+2)}
{

Kita meninjau elips

\mavergleichskettedisp
{\vergleichskette
{ E
}
{ =} { \left\{ (u,v)  \mid   2u^2 +5v^2 = 4         \right\}
}
{ } {
}
{ } {
}
{ } {
}
}
{}{}{.}
\aufzaehlungdrei{Definisikan suatu pemetaan surjektif yang terdiferensialkan secara kontinu
\maabb {} {\R} {E
} {.}
}{Deskripsikan suatu
\definitionsverweis {difeomorfisme}{}{}
antara lingkaran $S^1$ dan $E$.
}{
}

}
{} {}

\inputaufgabe
{0}
{

}
{} {}

\inputaufgabegibtloesung
{4}
{

Berikan suatu
\definitionsverweis {medan vektor}{}{} yang
\definitionsverweis {kontinu}{}{}
pada $S^2$ dan hanya mempunyai satu titik nol.

}
{} {}

\inputaufgabegibtloesung
{6}
{

Misalkan diberikan suatu
\definitionsverweis {data perekatan}{}{}

\mathbed {U_i} {}
 {i \in I} {}
 {} {} {} {,}
untuk
\definitionsverweis {ruang-ruang topologis}{}{.}
Buktikan bahwa terdapat suatu ruang topologis $X$ yang ditentukan secara unik, suatu penutup terbuka

\mavergleichskette
{\vergleichskette
{ X
}
{ = }{ \bigcup_{i \in I} V_i
}
{  }{
}
{    }{
}
{   }{
}
}
{}{}{}
dan
\definitionsverweis {homeomorfisme}{}{}
\maabb {\psi_i} {U_i } { V_i
} {}
sedemikian sehingga

\mavergleichskettedisp
{\vergleichskette
{ \psi_i  \left(  U_{ij}  \right)
}
{ =} { V_i \cap V_j
}
{ } {
}
{ } {
}
{ } {
}
}
{}{}{}
dan

\mavergleichskettedisp
{\vergleichskette
{ \psi_i \vert_{U_{ij} }
}
{ =} { \psi_j \vert_{U_{ji} }  \circ   \varphi_{ji}
}
{ } {
}
{ } {
}
{ } {
}
}
{}{}{}.

}
{} {}

\inputaufgabegibtloesung
{2}
{

Misalkan $M$ suatu
\definitionsverweis {manifold diferensiabel}{}{,}
\maabb {\gamma} { I } { M
} {}
suatu
\definitionsverweis {kurva diferensiabel}{}{}
di $M$, dan $\omega$ suatu
$1$-\definitionsverweis {bentuk}{}{}
terdiferensialkan pada $M$. Buktikan bahwa

\mavergleichskettedisp
{\vergleichskette
{ \gamma^* (d \omega)
}
{ =} { 0
}
{ } {
}
{ } {
}
{ } {
}
}
{}{}{}.

}
{} {}

\inputaufgabegibtloesung
{6 (3+3)}
{

Misalkan $M$ suatu
\definitionsverweis {manifold diferensiabel}{}{}
yang
\definitionsverweis {kompak}{}{}
dan mempunyai sedikitnya dua titik.
\aufzaehlungzwei {Buktikan bahwa suatu
$1$-\definitionsverweis {bentuk diferensial}{}{}
kontinu yang
\definitionsverweis {eksak}{}{}
pada $M$ mempunyai sedikitnya dua titik nol.
} {Buktikan bahwa hal ini tidak harus berlaku bagi suatu bentuk diferensial yang
\definitionsverweis {tertutup}{}{.}
}

}
{} {}

\inputaufgabegibtloesung
{4}
{

Berikan suatu contoh
\definitionsverweis {manifold dengan batas}{}{,}
yang
\definitionsverweis {terhubung}{}{}
dan batasnya terdiri atas
\zusatzklammer {gabungan saling lepas dari} {} {}
tak terhingga banyak salinan $\R^1$.

}
{} {}

\inputaufgabegibtloesung
{10}
{

Buktikan \stichwort {teorema tentang partisi kesatuan} {.}

}
{} {}

\inputaufgabegibtloesung
{10 (2+6+2)}
{

Kita meninjau
$2$-\definitionsverweis {bentuk diferensial}{}{}

\mavergleichskettedisp
{\vergleichskette
{ \omega
}
{ =} { x dy \wedge dz -y dx \wedge dz + z dx \wedge dy
}
{ } {
}
{ } {
}
{ } {
}
}
{}{}{}
pada bola satuan
\mathl{B  \left(  0 , 1  \right)}{.}
\aufzaehlungdreiabc{Buktikan bahwa $d \omega$ adalah tiga kali bentuk volume standar pada bola tersebut.
}{Buktikan bahwa
\mathl{\omega\vert_{S^2}}{} merupakan bentuk luas permukaan standar pada permukaan bola satuan.
}{Hitung luas permukaan bola dari
volume bola
dengan menggunakan
Teorema Stokes.
}

}
{} {}

\inputaufgabe
{0}
{

}
{} {}

\endgroup

\chapter{Formulir Peserta Ujian 5}
\noindent\textit{Formulir peserta dari sumber resmi; teks soal dipertahankan sesuai terjemahan yang telah diverifikasi.}
\begingroup
\renewcommand{\theHfakt}{o011.exam.learner.05.\arabic{fakt}}
%Data institusi

%\input{Dozentdaten}

%\renewcommand{\fachbereich}{Bidang studi}

%\renewcommand{\dozent}{Prof. Dr. . }

%Data ujian

\renewcommand{\klausurgebiet}{ }

\renewcommand{\klausurtyp}{ }

\renewcommand{\klausurdatum}{. 20}

\klausurvorspann {\fachbereich} {\klausurdatum} {\dozent} {\klausurgebiet} {\klausurtyp}

%Data untuk tabel poin berikut

\renewcommand{\aeins}{  3  }

\renewcommand{\azwei}{  3   }

\renewcommand{\adrei}{  5  }

\renewcommand{\avier}{  2  }

\renewcommand{\afuenf}{  4  }

\renewcommand{\asechs}{  8  }

\renewcommand{\asieben}{  4  }

\renewcommand{\aacht}{  5  }

\renewcommand{\aneun}{  3  }

\renewcommand{\azehn}{  2  }

\renewcommand{\aelf}{  6  }

\renewcommand{\azwoelf}{  7  }

\renewcommand{\adreizehn}{  4   }

\renewcommand{\avierzehn}{  8  }

\renewcommand{\afuenfzehn}{  64  }

\renewcommand{\asechzehn}{    }

\renewcommand{\asiebzehn}{    }

\renewcommand{\aachtzehn}{    }

\renewcommand{\aneunzehn}{    }

\renewcommand{\azwanzig}{    }

\renewcommand{\aeinundzwanzig}{    }

\renewcommand{\azweiundzwanzig}{    }

\renewcommand{\adreiundzwanzig}{    }

\renewcommand{\avierundzwanzig}{    }

\renewcommand{\afuenfundzwanzig}{    }

\renewcommand{\asechsundzwanzig}{    }

\punktetabellevierzehn

\klausurnote

\newpage

\setcounter{section}{0}

\inputaufgabegibtloesung
{3}
{

Definisikan istilah-istilah berikut
\zusatzklammer {yang dicetak miring} {} {}.
\aufzaehlungsechs{\stichwort {turunan tangensial kedua} {}
\zusatzklammer {atau
\stichwort {percepatan tangensial} {}} {} {}
\zusatzklammer {sebagai pemetaan ke $\R^n$} {} {}
dari suatu
\definitionsverweis {kurva yang dua kali terdiferensialkan secara kontinu}{}{}
\maabbdisp {\gamma} {I} {Y
} {}
yang bernilai dalam suatu
\definitionsverweis {hipermuka terdiferensialkan}{}{}

\mavergleichskette
{\vergleichskette
{ Y
}
{ \subseteq }{ \R^n
}
{  }{
}
{    }{
}
{   }{
}
}
{}{}{.}

}{Suatu ruang topologis $X$ yang
\stichwort {terhubung} {}.

}{\stichwort {Pemetaan singgung di titik} {}

\mavergleichskette
{\vergleichskette
{  P
}
{ \in }{ M
}
{  }{}
{    }{}
{   }{}
}
{}{}{}
untuk suatu pemetaan diferensiabel
\maabbdisp {\varphi} { M } { N
} {,}
dengan
\mathkor {} {M} {dan} {N} {}
manifold diferensiabel.

}{\stichwort {Produk} {}
dari manifold diferensiabel
\mathkor {} {M} {dan} {N} {.}

}{\stichwort {Ukuran volume} {}
yang terkait dengan bentuk volume positif $\omega$ pada suatu manifold diferensiabel berdimensi $n$.

}{Suatu
\stichwort {kurva geodesik} {}
pada manifold Riemann $M$.
}

}
{} {}

\inputaufgabegibtloesung
{3}
{

Nyatakan teorema-teorema berikut.
\aufzaehlungdrei{\stichwort {Teorema tentang kelengkungan kurva bidang yang diberikan secara implisit} {}

\mavergleichskette
{\vergleichskette
{  Y
}
{ \subseteq }{  \R^2
}
{  }{
}
{    }{
}
{   }{
}
}
{}{}{.}}{\stichwort {Rumus transformasi untuk bentuk volume positif pada manifold} {.}}{Teorema tentang kelengkungan Gauss dan kelengkungan seksional pada suatu permukaan diferensiabel

\mavergleichskette
{\vergleichskette
{  Y
}
{ \subseteq }{  \R^3
}
{  }{
}
{    }{
}
{   }{
}
}
{}{}{}}

}
{} {}

\inputaufgabegibtloesung
{5}
{

Misalkan

\mavergleichskette
{\vergleichskette
{ G
}
{ \subseteq }{ \R^n
}
{  }{
}
{    }{
}
{   }{
}
}
{}{}{}
\definitionsverweis {terbuka}{}{}
dan
\maabbdisp {\varphi} { G } { \R^m
} {}
adalah
\definitionsverweis {pemetaan yang terdiferensialkan secara kontinu}{}{,}
yang pada titik

\mavergleichskette
{\vergleichskette
{ P
}
{ \in }{ G
}
{  }{
}
{    }{
}
{   }{
}
}
{}{}{}
memiliki
\definitionsverweis {diferensial total}{}{}
yang surjektif. Misalkan

\mavergleichskette
{\vergleichskette
{ v
}
{ \in }{ T_PZ
}
{  }{
}
{    }{
}
{   }{
}
}
{}{}{}
adalah suatu vektor dalam
\definitionsverweis {ruang singgung serat}{}{}
$Z$ dari $\varphi$ yang melalui $P$. Tunjukkan bahwa terdapat
\definitionsverweis {kurva yang terdiferensialkan secara kontinu}{}{}
\maabbdisp {\gamma} { ]- \epsilon, \epsilon[ } { Z
} {}
\zusatzklammer {untuk suatu

\mavergleichskettek
{\vergleichskettek
{ \epsilon
}
{ > }{ 0
}
{  }{
}
{    }{
}
{   }{
}
}
{}{}{}
yang sesuai} {} {}
dengan

\mavergleichskette
{\vergleichskette
{ \gamma(0)
}
{ = }{ P
}
{  }{
}
{    }{
}
{   }{
}
}
{}{}{}
dan

\mavergleichskettedisp
{\vergleichskette
{ \gamma'(0)
}
{ =} { v
}
{ } {
}
{ } {
}
{ } {
}
}
{}{}{}.

}
{} {}

\inputaufgabegibtloesung
{2}
{

Misalkan
\mathkor {} {r} {dan} {R} {}
adalah bilangan real positif dengan

\mavergleichskette
{\vergleichskette
{ 0
}
{ < }{ r
}
{ < }{ R
}
{    }{
}
{   }{
}
}
{}{}{.}
Tentukan suatu
\definitionsverweis {medan normal satuan}{}{}
untuk torus
\zusatzklammer {tertanam} {} {}

\mavergleichskettedisp
{\vergleichskette
{ T
}
{ =} { \left\{ (x,y,z) \in \R^3  \mid   \left(  \sqrt{x^2+y^2} -R  \right)^2 +z^2  = r^2         \right\}
}
{ } {
}
{ } {
}
{ } {
}
}
{}{}{.}

}
{} {}

\inputaufgabe
{4}
{

Jelaskan peran teorema tentang pemetaan implisit dalam pengembangan teori manifold.

}
{} {}

\inputaufgabegibtloesung
{8}
{

Buktikan teorema tentang sifat adjoin-diri (self-adjoint) dari pemetaan Weingarten.

}
{} {}

\inputaufgabegibtloesung
{4}
{

Misalkan $X$ adalah suatu
\definitionsverweis {ruang kompak}{}{}
dan misalkan

\mavergleichskette
{\vergleichskette
{Y
}
{ \subseteq }{ X
}
{  }{
}
{    }{
}
{   }{
}
}
{}{}{}
adalah
\definitionsverweis {himpunan bagian tertutup}{}{,}
yang dilengkapi dengan
\definitionsverweis {topologi terinduksi}{}{.}
Tunjukkan bahwa $Y$ juga kompak.

}
{} {}

\inputaufgabegibtloesung
{5}
{

Berikan contoh suatu
\definitionsverweis {manifold diferensiabel}{}{}
$M$ yang
\definitionsverweis {berdimensi dua}{}{}
dan
\definitionsverweis {terhubung}{}{,}
beserta suatu titik

\mavergleichskette
{\vergleichskette
{ P
}
{ \in }{ M
}
{  }{
}
{    }{
}
{   }{
}
}
{}{}{}
sedemikian sehingga
\mathkor {} {M} {dan} {M \setminus \{P\}} {}
saling
\definitionsverweis {difeomorfik}{}{.}

}
{} {}

\inputaufgabegibtloesung
{3}
{

Untuk setiap
\definitionsverweis {vektor singgung}{}{}

\mavergleichskette
{\vergleichskette
{ v
}
{ \in }{ T_PS^2
}
{ \subseteq }{ \R^3
}
{    }{
}
{   }{
}
}
{}{}{}
dengan

\mavergleichskette
{\vergleichskette
{ \Vert {v} \Vert
}
{ = }{ 1
}
{  }{
}
{    }{
}
{   }{
}
}
{}{}{}
di suatu titik

\mavergleichskette
{\vergleichskette
{ P
}
{ \in }{ S^2
}
{ \subset }{ \R^3
}
{    }{
}
{   }{
}
}
{}{}{}
pada
\definitionsverweis {permukaan bola satuan}{}{,}
berikan suatu wakil
\definitionsverweis {diferensiabel}{}{}
\maabbdisp {\gamma} { I } { S^2
} {}
dengan

\mavergleichskette
{\vergleichskette
{ \gamma(0)
}
{ = }{ P
}
{  }{
}
{    }{
}
{   }{
}
}
{}{}{.}

}
{} {}

\inputaufgabegibtloesung
{2}
{

Tinjau parabola standar

\mavergleichskette
{\vergleichskette
{ Y
}
{ \subseteq }{ \R^2
}
{  }{
}
{    }{
}
{   }{
}
}
{}{}{}
dengan struktur Riemann terinduksi dan parametrisasi
\maabbeledisp {\varphi} {\R} { Y
} { t } { \left(  t   , \,  t^2                   \right)
} {,}
yang kita pandang sebagai suatu peta
\zusatzklammer {invers} {} {}.
Tentukan
\definitionsverweis {fungsi fundamental}{}{}
Riemann

\mavergleichskette
{\vergleichskette
{ g(t)
}
{ = }{ g_{11} (t)
}
{  }{
}
{    }{
}
{   }{
}
}
{}{}{.}

}
{} {}

\inputaufgabegibtloesung
{6}
{

Hitung bentuk diferensial tarik balik
\mathl{\varphi^* \tau}{} dari

\mavergleichskettedisp
{\vergleichskette
{ \tau
}
{ =} { dx \wedge dy \wedge dz - w dx \wedge dy \wedge dw + \cos (xy)  dx \wedge dz \wedge dw-yw dy \wedge dz \wedge dw
}
{ } {
}
{ } {
}
{ } {
}
}
{}{}{}
oleh pemetaan
\maabbeledisp {\varphi} { \R^3 } { \R^4
} { (r,s,t) } {  \left(  r^2s,t, \sin  r ,e^{st}  \right)  = (x,y,z,w)
} {.}

}
{} {}

\inputaufgabegibtloesung
{7 (2+3+2)}
{

Tinjau pemetaan-pemetaan diferensiabel
\maabbeledisp {\gamma} { [1,2] } {\R^2
} { t } {(t,t^{-1})
} {,}
dan
\maabbeledisp {\varphi} {\R^2} {\R^3
} {(u,v)} {(u^2,uv,-u+v^2)
} {,}
serta bentuk diferensial

\mavergleichskettedisp
{\vergleichskette
{ \omega
}
{ =} { xdx-zdy+dz
}
{ } {
}
{ } {
}
{ } {
}
}
{}{}{}
pada $\R^3$.

a) Hitung bentuk diferensial tarik balik
\mathl{\varphi^* \omega}{} pada $\R^2$.

b) Hitung integral lintasan bentuk diferensial
\mathl{\varphi^* \omega}{} sepanjang lintasan $\gamma$.

c) Hitung
\zusatzklammer {tanpa menggunakan hasil b)} {} {}
integral lintasan bentuk diferensial
\mathl{\omega}{} sepanjang lintasan $\varphi \circ \gamma$.

}
{} {}

\inputaufgabegibtloesung
{4}
{

Misalkan $M$ adalah suatu
\definitionsverweis {manifold berbatas}{}{.}
Tunjukkan bahwa setiap himpunan bagian terbuka

\mavergleichskette
{\vergleichskette
{ N
}
{ \subseteq }{ M
}
{  }{
}
{    }{
}
{   }{
}
}
{}{}{}
juga merupakan manifold dengan batas
\zusatzklammer {yang mungkin kosong} {} {},
dan bahwa

\mavergleichskettedisp
{\vergleichskette
{ \partial N
}
{ =} { \partial M \cap N
}
{ } {
}
{ } {
}
{ } {
}
}
{}{}{}
berlaku.

}
{} {}

\inputaufgabegibtloesung
{8}
{

Buktikan teorema tentang karakterisasi kurva geodesik terhadap koneksi Levi-Civita.

}
{} {}

\endgroup

\chapter{Formulir Peserta Ujian 6}
\noindent\textit{Formulir peserta dari sumber resmi; teks soal dipertahankan sesuai terjemahan yang telah diverifikasi.}
\begingroup
\renewcommand{\theHfakt}{o011.exam.learner.06.\arabic{fakt}}
\input{generated/exams/exam06_learner.id.build.tex}
\endgroup

\chapter{Formulir Peserta Ujian 7}
\noindent\textit{Formulir peserta dari sumber resmi; teks soal dipertahankan sesuai terjemahan yang telah diverifikasi.}
\begingroup
\renewcommand{\theHfakt}{o011.exam.learner.07.\arabic{fakt}}
%Data institusi

%\input{Dozentdaten}

%\renewcommand{\fachbereich}{Fachbereich}

%\renewcommand{\dozent}{Prof. Dr. . }

%Data ujian

\renewcommand{\klausurgebiet}{ }

\renewcommand{\klausurtyp}{ }

\renewcommand{\klausurdatum}{ . 20}

\klausurvorspann {\fachbereich} {\klausurdatum} {\dozent} {\klausurgebiet} {\klausurtyp}

%Data untuk tabel nilai berikut

\renewcommand{\aeins}{  3  }

\renewcommand{\azwei}{  3   }

\renewcommand{\adrei}{  0  }

\renewcommand{\avier}{  4  }

\renewcommand{\afuenf}{  5  }

\renewcommand{\asechs}{  4  }

\renewcommand{\asieben}{  9  }

\renewcommand{\aacht}{  0  }

\renewcommand{\aneun}{  8  }

\renewcommand{\azehn}{  6  }

\renewcommand{\aelf}{  0  }

\renewcommand{\azwoelf}{  3  }

\renewcommand{\adreizehn}{  3   }

\renewcommand{\avierzehn}{  0  }

\renewcommand{\afuenfzehn}{  0  }

\renewcommand{\asechzehn}{  7  }

\renewcommand{\asiebzehn}{  55  }

\renewcommand{\aachtzehn}{    }

\renewcommand{\aneunzehn}{    }

\renewcommand{\azwanzig}{    }

\renewcommand{\aeinundzwanzig}{    }

\renewcommand{\azweiundzwanzig}{    }

\renewcommand{\adreiundzwanzig}{    }

\renewcommand{\avierundzwanzig}{    }

\renewcommand{\afuenfundzwanzig}{    }

\renewcommand{\asechsundzwanzig}{    }

\punktetabellesechzehn

\klausurnote

\newpage

\setcounter{section}{K}

\inputaufgabegibtloesung
{3}
{

Definisikan istilah-istilah berikut
\zusatzklammer {yang dicetak miring} {} {.}
\aufzaehlungsechs{
\stichwort {Kelengkungan Gauss} {}
dari suatu permukaan diferensiabel

\mavergleichskette
{\vergleichskette
{ Y
}
{ \subseteq }{ \R^3
}
{  }{
}
{    }{
}
{   }{
}
}
{}{}{}
di suatu titik

\mavergleichskette
{\vergleichskette
{ P
}
{ \in }{ Y
}
{  }{
}
{    }{
}
{   }{
}
}
{}{}{.}

}{Suatu
\stichwort {manifold topologis} {}
berdimensi $n$.

}{Suatu
\stichwort {medan vektor tak bergantung waktu} {}
pada manifold diferensiabel $M$.

}{Suatu
\stichwort {bentuk diferensial berderajat $k$} {}
pada manifold diferensiabel $M$.

}{Suatu \stichwort {manifold Riemann} {.}

}{Suatu koneksi linear yang
\stichwort {bebas torsi} {}
pada bundel tangen manifold Riemann $M$.
}

}
{} {}

\inputaufgabegibtloesung
{3}
{

Nyatakan teorema-teorema berikut.
\aufzaehlungdrei{
\stichwort {Teorema eksistensi transpor paralel pada hiperpermukaan} {.}}{Teorema tentang deskripsi lokal bentuk diferensial.}{
\stichwort {Versi balok} {}
dari teorema Stokes.}

}
{} {}

\inputaufgabe
{0}
{

}
{} {}

\inputaufgabegibtloesung
{4 (1+3)}
{

Misalkan $M$ elips yang diberikan oleh

\mavergleichskettedisp
{\vergleichskette
{ x^2+3y^2
}
{ =} { 2
}
{ } {
}
{ } {
}
{ } {
}
}
{}{}{}
dalam $\R^2$,

\mavergleichskette
{\vergleichskette
{ P
}
{ = }{ \begin{pmatrix}  1  \\ { \frac{   1  }{  \sqrt{3}  } }    \end{pmatrix}
}
{  }{
}
{    }{
}
{   }{
}
}
{}{}{}
dan

\mavergleichskette
{\vergleichskette
{ Q
}
{ = }{ \begin{pmatrix}  -1 \\ { \frac{   1  }{  \sqrt{3}  } }    \end{pmatrix}
}
{  }{
}
{    }{
}
{   }{
}
}
{}{}{}
adalah titik-titik pada $M$.
\aufzaehlungdrei{Tunjukkan bahwa

\mavergleichskettedisp
{\vergleichskette
{ v
}
{ =} { \begin{pmatrix}  -  \sqrt{3}  \\ 1    \end{pmatrix}
}
{ } {
}
{ } {
}
{ } {
}
}
{}{}{}
merupakan
\definitionsverweis {vektor tangen}{}{}
di $P$ pada $M$.
}{Tentukan citra $v$ di bawah suatu transpor paralel pada $M$ dari $P$ ke $Q$.
}{
}

}
{} {}

\inputaufgabegibtloesung
{5}
{

Buktikan teorema tentang orientasi-orientasi pada hiperpermukaan diferensiabel.

}
{} {}

\inputaufgabegibtloesung
{4}
{

Misalkan

\mavergleichskette
{\vergleichskette
{ Y
}
{ \subseteq }{ W
}
{ \subseteq }{ \R^3
}
{    }{
}
{   }{
}
}
{}{}{,}
$W$
\definitionsverweis {terbuka}{}{,}
dan suatu
\definitionsverweis {permukaan diferensiabel}{}{}
yang dilengkapi dengan
\definitionsverweis {medan normal satuan}{}{}
$N$. Misalkan
\maabb {\gamma} { I } { Y
} {}
suatu kurva terdiferensialkan secara kontinu dan
\maabbdisp {F} { I } { \R^3
} {}
suatu
\definitionsverweis {medan vektor tangensial}{}{}
diferensiabel sepanjang $\gamma$ yang
\definitionsverweis {paralel}{}{.}
Tunjukkan bahwa medan vektor
\maabbeledisp {} { I } { \R^3
} { t } { N(\gamma(t)) \times F(t)
} {,}
juga paralel, dengan $\times$ menyatakan
\definitionsverweis {hasil kali silang}{}{}
dalam $\R^3$.

}
{} {}

\inputaufgabegibtloesung
{9}
{

Misalkan

\mavergleichskette
{\vergleichskette
{ M
}
{ \neq }{ \emptyset
}
{  }{
}
{    }{
}
{   }{
}
}
{}{}{}
suatu
\definitionsverweis {manifold diferensiabel}{}{}
berdimensi $n$. Tunjukkan bahwa terdapat suatu rantai
\definitionsverweis {submanifold-submanifold tertutup}{}{}

\mavergleichskettedisp
{\vergleichskette
{ M_0
}
{ \subseteq} { M_1
}
{ \subseteq} { M_2
}
{ \subseteq  \ldots \subseteq} { M_{n-1}
}
{ \subseteq} { M_n
}
}
{
\vergleichskettefortsetzung
{ =} { M
}
{ } {}
{ } {}
{ } {}
}{}{}
sedemikian sehingga submanifold tertutup $M_i$ berdimensi $i$.

}
{} {}

\inputaufgabe
{0}
{

}
{} {}

\inputaufgabegibtloesung
{8}
{

Misalkan
\maabbdisp {\psi} {\R^n } {\R
} {}
suatu
\definitionsverweis {fungsi diferensiabel}{}{}
dan

\mavergleichskettedisphandlinks
{\vergleichskettedisphandlinks
{M
}
{ =} { \left\{  \left(  x_1   , \,   , \ldots ,  , \,   x_n , \,   \psi(x_1 , \ldots , x_n)                \right)   \mid  (x_1 , \ldots , x_n) \in \R^n         \right\}
}
{ \subseteq} { \R^n \times \R
}
{ } {
}
{ } {
}
}
{}{}{}
merupakan
\definitionsverweis {grafik}{}{}
dari $\psi$. Tunjukkan bahwa $M$ suatu
\definitionsverweis {manifold Riemann}{}{}
yang
\definitionsverweis {berorientasi}{}{}
dan bahwa untuk
\definitionsverweis {bentuk volume kanonik}{}{}
$\omega$ pada $M$ berlaku rumus

\mavergleichskettedisp
{\vergleichskette
{ (\operatorname{Id} \times \psi)^*(\omega)
}
{ =} { \left(  1+ \sum_{i = 1}^n \left(  { \frac{   \partial   \psi  }{  \partial   x_i   } }  \right)^2  \right)^{1/2} dx_1 \wedge \ldots \wedge dx_n
}
{ } {
}
{ } {
}
{ } {
}
}
{}{}{.}

}
{} {}

\inputaufgabe
{6}
{

Hitung luas permukaan bola satuan.

}
{} {}

\inputaufgabe
{0}
{

}
{} {}

\inputaufgabegibtloesung
{3}
{

Hitung
\definitionsverweis {turunan eksterior}{}{}
$d \omega$ dari
$1$-\definitionsverweis {bentuk diferensial}{}{}

\mavergleichskettedisp
{\vergleichskette
{ \omega
}
{ =} { { \frac{   x^2 }{  y } } dx- { \frac{   x  }{  y^2 } } dy
}
{ } {
}
{ } {
}
{ } {
}
}
{}{}{}
pada

\mavergleichskette
{\vergleichskette
{ U
}
{ = }{ \left\{ (x,y) \in \R^2  \mid   y \neq 0         \right\}
}
{  }{
}
{    }{
}
{   }{
}
}
{}{}{.}

}
{} {}

\inputaufgabegibtloesung
{3}
{

Berikan contoh suatu
\definitionsverweis {manifold berbatas}{}{}
yang
\definitionsverweis {terhubung}{}{,}
dan batasnya terdiri atas
\zusatzklammer {gabungan saling lepas dari} {} {}
tak hingga banyak lingkaran $S^1$.

}
{} {}

\inputaufgabe
{0}
{

}
{} {}

\inputaufgabe
{0}
{

}
{} {}

\inputaufgabegibtloesung
{7}
{

Misalkan $M$ suatu
\definitionsverweis {manifold Riemann}{}{}
dan diberikan suatu
\definitionsverweis {koneksi linear}{}{}
pada
\definitionsverweis {bundel tangen}{}{}
$TM$ yang
\definitionsverweis {metrik}{}{}
dan
\definitionsverweis {bebas torsi}{}{.}
Tunjukkan bahwa untuk medan-medan vektor $X,Y,Z$ berlaku rumus

\mavergleichskettedisp
{\vergleichskette
{ \left\langle \nabla_X Y , Z  \right\rangle
}
{ =} { { \frac{   1  }{  2 } }  \left(  D_X \left\langle  Y  ,  Z  \right\rangle +  D_Y \left\langle  Z  ,  X  \right\rangle -  D_Z \left\langle  X  ,  Y  \right\rangle  - \left\langle  Y  , [X,Z]   \right\rangle - \left\langle  Z  , [Y,X]   \right\rangle + \left\langle  X  , [Z,Y]   \right\rangle   \right)
}
{ } {
}
{ } {
}
{ } {
}
}
{}{}{.}

}
{} {}

\endgroup

\chapter{Formulir Peserta Ujian 8}
\noindent\textit{Formulir peserta dari sumber resmi; teks soal dipertahankan sesuai terjemahan yang telah diverifikasi.}
\begingroup
\renewcommand{\theHfakt}{o011.exam.learner.08.\arabic{fakt}}
\input{generated/exams/exam08_learner.id.build.tex}
\endgroup

\chapter{Formulir Peserta Ujian 9}
\noindent\textit{Formulir peserta dari sumber resmi; teks soal dipertahankan sesuai terjemahan yang telah diverifikasi.}
\begingroup
\renewcommand{\theHfakt}{o011.exam.learner.09.\arabic{fakt}}
\input{generated/exams/exam09_learner.id.build.tex}
\endgroup

\chapter{Formulir Peserta Ujian 10}
\noindent\textit{Formulir peserta dari sumber resmi; teks soal dipertahankan sesuai terjemahan yang telah diverifikasi.}
\begingroup
\renewcommand{\theHfakt}{o011.exam.learner.10.\arabic{fakt}}
\input{generated/exams/exam10_learner.id.build.tex}
\endgroup

\part{Formulir Solusi Resmi}
\chapter*{Batas solusi resmi}
\addcontentsline{toc}{chapter}{Batas solusi resmi}
Setiap formulir pada bagian ini berasal dari permukaan solusi resmi yang dibekukan. Solusi yang tidak tersedia pada sumber tetap tidak tersedia di sini; tidak ada materi asli yang dimasukkan ke dalam formulir resmi.
\chapter{Formulir Solusi Resmi Ujian 1}
\noindent\textit{Solusi yang disediakan oleh sumber resmi; kekosongan sumber dipertahankan.}
\begingroup
\renewcommand{\theHfakt}{o011.exam.official.01.\arabic{fakt}}
%Data tentang institusi

%\input{Data guru}

%\renewcommand{\fachbereich}{Departemen}

%\renewcommand{\dozent}{Prof. Dr. }

%Tanggal ujian

\renewcommand{\klausurgebiet}{ }

\renewcommand{\klausurtyp}{ }

\renewcommand{\klausurdatum}{. 20}

\klausurvorspann {\fachbereich} {\klausurdatum} {\dozent} {\klausurgebiet} {\klausurtyp}

%Data untuk tabel poin berikut

\renewcommand{\aeins}{  3  }

\renewcommand{\azwei}{  3   }

\renewcommand{\adrei}{  6  }

\renewcommand{\avier}{  5  }

\renewcommand{\afuenf}{  5  }

\renewcommand{\asechs}{  2  }

\renewcommand{\asieben}{  5  }

\renewcommand{\aacht}{  5  }

\renewcommand{\aneun}{  7  }

\renewcommand{\azehn}{  3  }

\renewcommand{\aelf}{  2  }

\renewcommand{\azwoelf}{  5  }

\renewcommand{\adreizehn}{  9   }

\renewcommand{\avierzehn}{  5  }

\renewcommand{\afuenfzehn}{  65  }

\renewcommand{\asechzehn}{    }

\renewcommand{\asiebzehn}{    }

\renewcommand{\aachtzehn}{    }

\renewcommand{\aneunzehn}{    }

\renewcommand{\azwanzig}{    }

\renewcommand{\aeinundzwanzig}{    }

\renewcommand{\azweiundzwanzig}{    }

\renewcommand{\adreiundzwanzig}{    }

\renewcommand{\avierundzwanzig}{    }

\renewcommand{\afuenfundzwanzig}{    }

\renewcommand{\asechsundzwanzig}{    }

\punktetabellevierzehn

\klausurnote

\newpage

\setcounter{section}{0}



\inputaufgabeklausurloesung
{3}

{

Definisikan istilah-istilah berikut
\zusatzklammer {yang dicetak miring} {} {}.
\aufzaehlungsechs{\stichwort {Lingkaran kelengkungan} {}
dari suatu
\definitionsverweis {kurva}{}{}
yang dua kali
\definitionsverweis {terdiferensialkan secara kontinu}{}{}
dan
\definitionsverweis {berparametrisasi panjang busur}{}{}.
Misalkan
\maabbdisp {\gamma} {I} {\R^2
} {}
dan ambil titik

\mavergleichskette
{\vergleichskette
{ t_0
}
{ \in }{ I
}
{  }{
}
{    }{
}
{   }{
}
}
{}{}{.}

}{Suatu \stichwort {submanifold tertutup} {}
\mathl{M \subseteq N}{} dari manifold diferensiabel $N$.

}{\stichwort {Ruang singgung} {} pada titik
\mathl{P\in M}{} dari manifold diferensiabel $M$.

}{Istilah
\stichwort {berorientasi} {}
untuk suatu
\definitionsverweis {manifold diferensiabel}{}{}
$M$.

}{\stichwort {Batas} {}
suatu
\definitionsverweis {manifold berbatas}{}{.}

}{\stichwort {Simbol Christoffel} {}

\mathl{\Gamma_{ij}^k}{} bagi koneksi Levi-Civita dari struktur Riemann $\left(  g_{ij}  \right)$ pada himpunan terbuka

\mavergleichskette
{\vergleichskette
{ U
}
{ \subseteq }{ \R^n
}
{  }{
}
{    }{
}
{   }{
}
}
{}{}{.}
}

}
{

\aufzaehlungsechs{Lingkaran dengan jari-jari

\mavergleichskettedisp
{\vergleichskette
{ r
}
{ =} {  { \frac{  1 }{ \betrag {  \gamma_1'(t_0) \gamma_2^{\prime \prime} (t_0) - \gamma_2'(t_0) \gamma_1^{\prime \prime} (t_0)  }  } }
}
{ } {
}
{ } {
}
{ } {
}
}
{}{}{}
dan pusat

\mavergleichskettedisp
{\vergleichskette
{ M
}
{ =} { \gamma(t_0) +  { \frac{  1 }{  \gamma_1'(t_0) \gamma_2^{\prime \prime} (t_0) - \gamma_2'(t_0) \gamma_1^{\prime \prime} (t_0)  } }  \begin{pmatrix} - \gamma_2'(t_0) \\ \gamma_1'(t_0)   \end{pmatrix}
}
{ } {
}
{ } {
}
{ } {
}
}
{}{}{}
dinamakan lingkaran kelengkungan dari $\gamma$ di $t_0$.
}{Suatu himpunan bagian tertutup
\mathl{M \subseteq N}{} dinamakan submanifold tertutup jika untuk setiap titik
\mathl{P \in M}{} terdapat suatu
\definitionsverweis {peta}{}{}
sedemikian sehingga
\mathl{P \in W \subseteq N}{},
\definitionsverweis {terbuka}{}{,}
\maabb {\theta} { W } { W'
} {,}

\mathl{W' \subseteq \R^n}{} terbuka, dan

\mavergleichskettedisp
{\vergleichskette
{ M \cap W
}
{ =} {\theta^{-1} (( \R^m \times\{0\}) \cap W' )
}
{ } {
}
{ } {
}
{ } {
}
}
{}{}{.}
}{Ruang singgung
\mathl{T_PM}{} terdiri atas semua kelas ekuivalensi lintasan terdiferensialkan yang melalui titik tersebut dan ekuivalen secara tangensial.
}{Suatu
\definitionsverweis {manifold diferensiabel}{}{}
$M$ dengan sebuah
\definitionsverweis {atlas}{}{}

\mathl{(U_i, V_i, \alpha_i)}{} disebut berorientasi jika setiap petanya
\definitionsverweis {berorientasi}{}{}
dan semua pergantian petanya
\definitionsverweis {mempertahankan orientasi}{}{}
.
}{Batas $M$ didefinisikan oleh

\mavergleichskettedisp
{\vergleichskette
{ \partial M
}
{ =} {  \left\{ x \in M  \mid   \alpha_i(x) \in  \partial H \cap V_i  \text{ untuk } \text{suatu } i \in I         \right\}
}
{ } {
}
{ } {
}
{ } {
}
}
{}{}{}
di mana $\alpha_i$ adalah peta-peta koordinat.
}{Misalkan
\mathl{\left(  g^{k l}  \right)}{} adalah matriks invers dari
\mathl{\left(  g_{ij}  \right)}{.}
\zusatzklammer {Fungsi-fungsi bernilai riil pada $U$} {} {}

\mavergleichskettedisp
{\vergleichskette
{\Gamma_{ij}^k
}
{ =} {  { \frac{  1 }{ 2 } } \sum_{l = 1}^n g^{kl}(\partial_ig_{jl}+\partial_jg_{il}-\partial_lg_{ij})
}
{ } {
}
{ } {
}
{ } {
}
}
{}{}{}
dinamakan simbol-simbol Christoffel untuk koneksi Levi-Civita.
}

 }



\inputaufgabeklausurloesung
{3}

{

Nyatakan teorema-teorema berikut.
\aufzaehlungdrei{Teorema tentang citra
\definitionsverweis {pemetaan Gauss}{}{}
dari suatu hipermuka.}{Rumus luas permukaan dari permukaan putar yang dibentuk oleh
\definitionsverweis {kurva diferensiabel}{}{}
\maabbeledisp {\gamma} {]a,b[} {\R^2
} { t } {(x(t),y(t))
} {,}
dengan

\mavergleichskette
{\vergleichskette
{  y(t)
}
{ > }{  0
}
{  }{
}
{    }{
}
{   }{
}
}
{}{}{.}}{\stichwort {Aturan hasil kali} {}
untuk turunan eksterior.}

}
{

\aufzaehlungdrei{\faktsituation {Misalkan

\mavergleichskette
{\vergleichskette
{ W
}
{ \subseteq }{\R^n
}
{ }{
}
{ }{
}
{ }{
}
}
{}{}{}
adalah suatu
\definitionsverweis {himpunan bagian terbuka}{}{}
dan
\maabbdisp {h} { W } {\R
} {}
adalah suatu
\definitionsverweis {fungsi terdiferensialkan secara kontinu}{}{.}}
\faktvoraussetzung {
\definitionsverweis {serat}{}{}

\mavergleichskette
{\vergleichskette
{ Y
}
{ = }{ h^{-1}(c)
}
{ }{
}
{ }{
}
{ }{
}
}
{}{}{}
dari $h$ di

\mavergleichskette
{\vergleichskette
{ c
}
{ \in }{\R
}
{  }{
}
{    }{
}
{   }{
}
}
{}{}{}
bersifat
\definitionsverweis {kompakt}{}{}
dan
\definitionsverweis {reguler}{}{.}}
\faktfolgerung {Maka
\definitionsverweis {pemetaan Gauss}{}{}
\maabbdisp {} { Y } { S^{n-1}
} {}
bersifat surjektif.}
\faktzusatz {}
\faktzusatz {}}{\faktsituation {Misalkan
\maabbeledisp {\gamma} {]a,b[} {\R^2
} { t } {(x(t),y(t))
} {,}
adalah suatu
\definitionsverweis {kurva terdiferensialkan}{}{}
dengan

\mavergleichskette
{\vergleichskette
{ y(t)
}
{ > }{ 0
}
{ }{
}
{ }{
}
{ }{
}
}
{}{}{,}}
\faktvoraussetzung {yang menginduksi suatu
\definitionsverweis {difeomorfisme}{}{}
ke

\mavergleichskette
{\vergleichskette
{M
}
{ = }{ \gamma(I)
}
{ }{
}
{ }{
}
{ }{
}
}
{}{}{}
, dengan

\mavergleichskette
{\vergleichskette
{M
}
{ \subseteq }{ G
}
{ }{
}
{ }{
}
{ }{
}
}
{}{}{}
sebagai
\definitionsverweis {submanifold tertutup}{}{}
satu dimensi dalam suatu himpunan terbuka

\mavergleichskette
{\vergleichskette
{G
}
{ \subseteq }{\R \times \R_+
}
{ }{
}
{ }{
}
{ }{
}
}
{}{}{}
.}
\faktfolgerung {Maka
\definitionsverweis {permukaan putar}{}{}
yang bersesuaian merupakan submanifold tertutup dari
\mathl{\R^3}{} tanpa sumbu $x$, dan luas permukaannya adalah

\mathdisp {2 \pi \int_{  a  }^{  b  } \sqrt{ (x'(t))^2+ (y'(t))^2 } \cdot y(t) \, d t} { . }
}
\faktzusatz {}
\faktzusatz {}}{\faktsituation {Misalkan

\mavergleichskette
{\vergleichskette
{ U
}
{ \subseteq }{ \R^n
}
{ }{
}
{ }{
}
{ }{
}
}
{}{}{}
\definitionsverweis {terbuka}{}{}
dan misalkan

\mavergleichskette
{\vergleichskette
{ \omega
}
{ \in }{
{ \mathcal E }^{  k  }_{ 1 } (  U   )
}
{  }{
}
{    }{
}
{   }{
}
}
{}{}{}
serta

\mavergleichskette
{\vergleichskette
{ \tau
}
{ \in }{
{ \mathcal E }^{  \ell }_{ 1 } (  U   )
}
{  }{
}
{    }{
}
{   }{
}
}
{}{}{}
adalah bentuk-bentuk diferensial terdiferensialkan pada $U$.}
\faktfolgerung {Maka berlaku aturan hasil kali

\mavergleichskettedisp
{\vergleichskette
{ d( \omega \wedge \tau)
}
{ =} { ( d \omega) \wedge \tau  + (-1)^k \omega  \wedge  (d \tau)
}
{ } {
}
{ } {
}
{ } {
}
}
{}{}{.}}
\faktzusatz {}
\faktzusatz {}}

 }



\inputaufgabeklausurloesung
{6 (1+2+3)}

{

Misalkan $Y$ adalah hipermuka yang diberikan oleh persamaan

\mavergleichskettedisp
{\vergleichskette
{ 2x^2 +3y^2+5z^2
}
{ =} { 10
}
{ } {
}
{ } {
}
{ } {
}
}
{}{}{}
di $\R^3$. Misalkan

\mavergleichskette
{\vergleichskette
{ P
}
{ = }{ \begin{pmatrix}  -1 \\ -1\\  1   \end{pmatrix}
}
{ \in }{ Y
}
{    }{
}
{   }{
}
}
{}{}{}
dan

\mavergleichskette
{\vergleichskette
{ w
}
{ = }{ \begin{pmatrix}  -3 \\ 2 \\  4   \end{pmatrix}
}
{ \in }{ \R^3
}
{    }{
}
{   }{
}
}
{}{}{.}
\aufzaehlungdrei{Tentukan
\definitionsverweis {ruang singgung}{}{}

\mathl{T_PY}{.}
}{Tentukan
\definitionsverweis {dekomposisi ortogonal}{}{}
$w$ menjadi komponen tangensial dan komponen ortogonal pada titik $P$.
}{Realisasikan komponen tangensial $w$ di $T_PY$ melalui suatu kurva diferensiabel yang seluruhnya terletak pada $Y$.
}

}
{

\aufzaehlungdrei{Total diferensialnya adalah

\mathdisp {\left(  4x  , \,   6y  , \,   10z                 \right)} { , }

sehingga pada $P$ diperoleh

\mavergleichskettedisp
{\vergleichskette
{ T_PY
}
{ =} { \left\{  \begin{pmatrix}  a  \\ b \\ c   \end{pmatrix}   \mid   -4a-6b+10c = 0         \right\}
}
{ =} { \left\{  \begin{pmatrix}  a  \\ b \\ c   \end{pmatrix}   \mid      \left\langle  \begin{pmatrix}  a  \\ b \\ c   \end{pmatrix}   , \begin{pmatrix}  -4 \\ -6\\  10   \end{pmatrix}     \right\rangle = 0         \right\}
}
{ } {
}
{ } {
}
}
{}{}{.}
}{Kita gunakan ansatz

\mavergleichskettedisp
{\vergleichskette
{  \begin{pmatrix}  -3 \\ 2\\  4   \end{pmatrix}
}
{ =} { s  \begin{pmatrix}  -4 \\ -6\\  10   \end{pmatrix} + v
}
{ } {
}
{ } {
}
{ } {
}
}
{}{}{}
dengan skalar yang dicari

\mavergleichskette
{\vergleichskette
{ s
}
{ \in }{ \R
}
{ }{
}
{ }{
}
{ }{
}
}
{}{}{}
dan vektor tangensial yang dicari $v$. Kita peroleh

\mavergleichskettedisp
{\vergleichskette
{ \left\langle \begin{pmatrix} -3 \\ 2\\ 4 \end{pmatrix} , \begin{pmatrix} -4 \\ -6\\ 10 \end{pmatrix} \right\rangle
}
{ =} { 12-12+40
}
{ =} { 40
}
{ } {
}
{ } {
}
}
{}{}{.}
Karena

\mavergleichskettedisp
{\vergleichskette
{ \left\langle  \begin{pmatrix}  -4 \\ -6\\  10   \end{pmatrix}  ,  \begin{pmatrix}  -4 \\ -6\\  10   \end{pmatrix}   \right\rangle
}
{ =} { 16 +36 + 100
}
{ =} { 152
}
{ } {
}
{ } {
}
}
{}{}{}
maka

\mavergleichskettedisp
{\vergleichskette
{ s
}
{ =} { { \frac{   40 }{  152 } }
}
{ =} { { \frac{   5 }{  19 } }
}
{ } {
}
{ } {
}
}
{}{}{}
dan dengan demikian

\mavergleichskettedisp
{\vergleichskette
{ v
}
{ =} { \begin{pmatrix}  -3 \\ 2\\  4   \end{pmatrix}   - { \frac{   5 }{  19 } }  \begin{pmatrix}  -4 \\ -6\\  10   \end{pmatrix}
}
{ =} { { \frac{   1  }{  19 } }  \begin{pmatrix}  -57+20 \\ 38 +30\\  76 -50   \end{pmatrix}
}
{ =} { { \frac{   1  }{  19 } }  \begin{pmatrix}  -37 \\ 68\\  26   \end{pmatrix}
}
{ } {
}
}
{}{}{.}
Jadi, dekomposisi ortogonalnya adalah

\mavergleichskettedisp
{\vergleichskette
{  w
}
{ =} { \begin{pmatrix}  -3 \\ 2\\  4   \end{pmatrix}
}
{ =} { { \frac{   5 }{  19 } }  \begin{pmatrix}  -4 \\ -6\\  10   \end{pmatrix} + { \frac{   1  }{  19 } }  \begin{pmatrix}  -37 \\ 68\\  26   \end{pmatrix}
}
{ } {
}
{ } {
}
}
{}{}{.}
}{Kita tulis persamaan permukaan

\mavergleichskette
{\vergleichskette
{ 2x^2 +3y^2+5z^2
}
{ = }{ 10
}
{ }{
}
{ }{
}
{ }{
}
}
{}{}{}
sebagai

\mavergleichskettedisp
{\vergleichskette
{ z
}
{ =} { \pm \sqrt{ 10- 2x^2 - 3y^2 }
}
{ } {
}
{ } {
}
{ } {
}
}
{}{}{,}
di mana akar positif harus dipilih pada titik $P$. Ini merupakan parameterisasi permukaan pada suatu lingkungan terbuka dari $P$; ruang singgung di
\mathl{\begin{pmatrix} -1 \\ -1 \end{pmatrix}}{} dipetakan ke ruang singgung
\mathl{T_PY}{}.
}

 }



\inputaufgabeklausurloesung
{5}

{

Misalkan
\maabbeledisp {\gamma} {\R} { \R^2
} { t } { \begin{pmatrix}  \cos \left(  \omega t \right)  \\ \sin  \left(  \omega t \right)    \end{pmatrix}
} {,}
dengan

\mavergleichskette
{\vergleichskette
{ \omega
}
{ > }{ 0
}
{  }{
}
{    }{
}
{   }{
}
}
{}{}{.}
Tunjukkan bahwa tidak ada gerak melingkar lain dengan norma kecepatan konstan yang memiliki nilai dan turunan hingga orde kedua yang sama dengan $\gamma$ di titik parameter

\mavergleichskette
{\vergleichskette
{ t
}
{ = }{ 0
}
{  }{
}
{    }{
}
{   }{
}
}
{}{}{}
tersebut.

}
{

Kita peroleh

\mavergleichskettedisp
{\vergleichskette
{ \gamma(0)
}
{ =} { \begin{pmatrix} 1 \\ 0 \end{pmatrix}
}
{ } {
}
{ } {
}
{ } {
}
}
{}{}{,}

\mavergleichskettedisp
{\vergleichskette
{ \gamma'(0)
}
{ =} { \begin{pmatrix} 0 \\ \omega \end{pmatrix}
}
{ } {
}
{ } {
}
{ } {
}
}
{}{}{}
dan

\mavergleichskettedisp
{\vergleichskette
{ \gamma^{\prime \prime} (0)
}
{ =} { \begin{pmatrix} - \omega^2 \\ 0 \end{pmatrix}
}
{ } {
}
{ } {
}
{ } {
}
}
{}{}{.}
Suatu gerak melingkar yang sama dengan gerak yang diberikan pada

\mavergleichskette
{\vergleichskette
{ t
}
{ = }{ 0
}
{ }{
}
{ }{
}
{ }{
}
}
{}{}{}
hingga turunan kedua khususnya harus memiliki garis singgung yang sama di titik tersebut. Oleh karena itu, pusat lingkarannya harus terletak pada sumbu $x$. Selain itu, pusat itu harus terletak di sebelah kiri
\mathl{\begin{pmatrix} 1 \\ 0 \end{pmatrix}}{} karena percepatannya mengarah ke kiri. Karena itu, gerak melingkar beraturan yang dicari dengan pusat
\mathl{\begin{pmatrix} x \\ 0 \end{pmatrix}}{}
\zusatzklammer {
\mavergleichskettek
{\vergleichskettek
{ x
}
{ < }{ 1
}
{ }{
}
{ }{
}
{ }{
}
}
{}{}{}} {} {}
dan jari-jari
\mathl{1-x}{} dapat dituliskan sebagai

\mavergleichskettedisp
{\vergleichskette
{  \delta (t)
}
{ =} {  \begin{pmatrix}  x  \\ 0   \end{pmatrix} + (1-x) \begin{pmatrix}  \cos \left(  u t \right)  \\ \sin  \left(  u t \right)     \end{pmatrix}
}
{ } {
}
{ } {
}
{ } {
}
}
{}{}{}
. Kita mempunyai

\mavergleichskettedisp
{\vergleichskette
{ \delta(0)
}
{ =} { \gamma(0)
}
{ } {
}
{ } {
}
{ } {
}
}
{}{}{.}
Selanjutnya,

\mavergleichskettedisp
{\vergleichskette
{ \delta'( 0)
}
{ =} { (1-x) \begin{pmatrix} 0 \\ u \end{pmatrix}
}
{ } {
}
{ } {
}
{ } {
}
}
{}{}{}
dan

\mavergleichskettedisp
{\vergleichskette
{ \delta^{\prime \prime}( 0)
}
{ =} { (1-x) \begin{pmatrix} -u^2 \\ 0 \end{pmatrix}
}
{ } {
}
{ } {
}
{ } {
}
}
{}{}{.}
Agar syarat-syarat yang diminta terpenuhi, harus berlaku

\mavergleichskettedisp
{\vergleichskette
{ \omega
}
{ =} { (1-x) u
}
{ } {
}
{ } {
}
{ } {
}
}
{}{}{}
dan

\mavergleichskettedisp
{\vergleichskette
{ \omega^2
}
{ =} { (1-x) u^2
}
{ } {
}
{ } {
}
{ } {
}
}
{}{}{}
. Dari sini diperoleh

\mavergleichskettedisp
{\vergleichskette
{ \omega^2
}
{ =} { \omega u
}
{ } {
}
{ } {
}
{ } {
}
}
{}{}{}
dan dengan demikian

\mavergleichskettedisp
{\vergleichskette
{ u
}
{ =} { \omega
}
{ } {
}
{ } {
}
{ } {
}
}
{}{}{}
serta

\mavergleichskettedisp
{\vergleichskette
{ x
}
{ =} { 0
}
{ } {
}
{ } {
}
{ } {
}
}
{}{}{.}

 }



\inputaufgabeklausurloesung
{5}

{

Buktikan teorema tentang kelengkungan suatu kurva bidang yang diberikan secara implisit.

}
{

Misalkan
\maabb {\gamma} { I } { \R^2
} {}
adalah parameterisasi panjang busur kurva $Y$ pada suatu lingkungan dari $P$

\mavergleichskette
{\vergleichskette
{ \gamma (0)
}
{ = }{ P
}
{ }{
}
{ }{
}
{ }{
}
}
{}{}{,}
yang sesuai dengan orientasi yang diberikan. Diferensial total
\maabbdisp {\left(D N \right)_{ P }} { \R^2} {\R^2
} {}
bersifat linier, sehingga cukup membuktikan pernyataan bagi vektor $\gamma'(0)$. Berdasarkan
aturan rantai,

\mavergleichskettedisp
{\vergleichskette
{ \left(  D N   \right)_{ P } \left(   \gamma'(0)  \right)
}
{ =} { (N \circ \gamma )'(0)
}
{ } {
}
{ } {
}
{ } {
}
}
{}{}{.}
Vektor ini merupakan kelipatan dari $\gamma'(0)$, sehingga sama dengan

\mathdisp {\left\langle  (N \circ  \gamma)'(0) ,  \gamma' (0)   \right\rangle  \gamma'(0)} { . }

Karena syarat ortogonalitas memberikan

\mavergleichskettedisp
{\vergleichskette
{ \left\langle (N \circ \gamma)(t) , \gamma' (t) \right\rangle
}
{ =} { 0
}
{ } {
}
{ } {
}
{ } {
}
}
{}{}{}
untuk semua

\mavergleichskette
{\vergleichskette
{ t
}
{ \in }{ I
}
{ }{
}
{ }{
}
{ }{
}
}
{}{}{}
dan karenanya

\mavergleichskettedisp
{\vergleichskette
{ \left\langle (N \circ \gamma)'(t) , \gamma' (t) \right\rangle + \left\langle (N \circ \gamma)(t) , \gamma^{\prime \prime} (t) \right\rangle
}
{ =} { 0
}
{ } {
}
{ } {
}
{ } {
}
}
{}{}{.}
Dengan demikian,

\mavergleichskettedisp
{\vergleichskette
{ \left(  D N   \right)_{ P } \left(   \gamma'(0)  \right)
}
{ =} { - \left\langle  (N \circ  \gamma)(0) ,  \gamma^{\prime \prime} (0)   \right\rangle  \gamma'(0)
}
{ =} { - \kappa( \gamma(0)) \gamma'(0)
}
{ } {
}
{ } {
}
}
{}{}{}
berdasarkan
Lemma 3.8 (Differentialgeometrie (Osnabrück 2023)).

 }



\inputaufgabeklausurloesung
{2}

{

Berikan contoh manifold diferensiabel
\mathkor {} {M} {dan} {N} {}
dengan

\mavergleichskette
{\vergleichskette
{ \operatorname{ dim}_{  } ^{  }  \left(   M   \right), \, \operatorname{ dim}_{  } ^{  }  \left(   N   \right)
}
{ \geq }{ 1
}
{  }{
}
{    }{
}
{   }{
}
}
{}{}{}
serta pemetaan diferensiabel
\maabb {\varphi} { M } { N
} {}
dan
\maabb {\psi} { N } { M
} {}
sedemikian sehingga

\mavergleichskette
{\vergleichskette
{ \psi \circ \varphi
}
{ = }{
\operatorname{Id}_{ M }
}
{  }{
}
{    }{
}
{   }{
}
}
{}{}{}
dan

\mavergleichskette
{\vergleichskette
{\varphi \circ  \psi
}
{ \neq }{
\operatorname{Id}_{ N }
}
{  }{
}
{    }{
}
{   }{
}
}
{}{}{}
berlaku.

}
{

Kita pilih

\mavergleichskette
{\vergleichskette
{M
}
{ = }{\R
}
{ }{
}
{ }{
}
{ }{
}
}
{}{}{}
dan

\mavergleichskette
{\vergleichskette
{N
}
{ = }{\R^2
}
{ }{
}
{ }{
}
{ }{
}
}
{}{}{}
dan
\maabbeledisp {\varphi} {\R} {\R^2
} { x } {(x,0)
} {,}
dan
\maabbeledisp {\psi} {\R^2} {\R
} {(x,y)} { x
} {.}
Maka

\mavergleichskettedisp
{\vergleichskette
{ \left( \psi \circ \varphi \right) (x)
}
{ =} { \psi((x,0))
}
{ =} { x
}
{ } {
}
{ } {
}
}
{}{}{}
dan

\mavergleichskettedisp
{\vergleichskette
{ \left(  \varphi  \circ \psi  \right) ((x,y))
}
{ =} { \varphi (x)
}
{ =} {(x,0)
}
{ } {
}
{ } {
}
}
{}{}{.}

 }



\inputaufgabeklausurloesung
{5}

{

Buktikan bahwa ekuivalensi tangensial lintasan-lintasan pada suatu manifold di titik $P$ dapat diperiksa dengan menggunakan sebarang peta koordinat.

}
{

Untuk suatu
\definitionsverweis {kurva terdiferensialkan}{}{}
\maabbdisp {\gamma} { I } { M
} {}
dengan

\mavergleichskette
{\vergleichskette
{ 0
}
{ \in }{ I
}
{ }{
}
{ }{
}
{ }{
}
}
{}{}{}
dan

\mavergleichskette
{\vergleichskette
{ \gamma(0)
}
{ = }{ P
}
{  }{
}
{    }{
}
{   }{
}
}
{}{}{}
dan suatu
\definitionsverweis {peta}{}{}
\maabbdisp {\alpha} { U } { V
} {}
\zusatzklammer {dengan

\mavergleichskettek
{\vergleichskettek
{P
}
{ \in }{ U
}
{ }{
}
{ }{
}
{ }{
}
}
{}{}{}
dan

\mavergleichskette
{\vergleichskette
{V
}
{ \subseteq }{\R^n
}
{ }{
}
{ }{
}
{ }{
}
}
{}{}{}} {} {}
ekspresi

\mathdisp {\left(  \alpha \circ \left(  \gamma \vert_{\gamma^{-1} (U)}  \right)  \right)'(0)} {  }

tidak berubah apabila kita beralih ke interval terbuka yang lebih kecil

\mavergleichskette
{\vergleichskette
{ 0
}
{ \in }{ I'
}
{ \subseteq }{ I
}
{ }{
}
{ }{
}
}
{}{}{}
dan himpunan terbuka yang lebih kecil

\mavergleichskette
{\vergleichskette
{P
}
{ \in }{ U'
}
{ \subseteq }{ U
}
{ }{
}
{ }{
}
}
{}{}{}
\zusatzklammer {dengan peta yang diinduksi} {} {}
. Karena itu, kita dapat mengasumsikan bahwa
\mathkor {} {\gamma_1} {dan} {\gamma_2} {}
didefinisikan pada interval yang sama, citranya terletak di $U$, dan pada $U$ terdapat dua peta
\maabbdisp {\alpha_1} { U } { V_1
} {}
dan
\maabbdisp {\alpha_2} { U } { V_2
} {}
. Selanjutnya, dari

\mavergleichskettedisp
{\vergleichskette
{ ( \alpha_1 \circ \gamma_1 )'(0)
}
{ =} { ( \alpha_1 \circ \gamma_2 )'(0)
}
{ } {
}
{ } {
}
{ } {
}
}
{}{}{}
berdasarkan
aturan rantai
dan keterdiferensialan
\definitionsverweis {pemetaan transisi}{}{}

\mathl{\alpha_2 \circ \alpha_1^{-1}}{}, langsung diperoleh

\mavergleichskettealign
{\vergleichskettealign
{ \left(  \alpha_2 \circ \gamma_1  \right)'(0)
}
{ =} { \left(  \left(  \alpha_2 \circ \alpha_1^{-1}  \right) \circ \left(  \alpha_1 \circ \gamma_1  \right)  \right)'(0)
}
{ =} { \left(  D   \left(   \alpha_2 \circ \alpha_1^{-1}   \right)      \right)_{\alpha_1(P)} \left(   \left(  \alpha_1 \circ \gamma_1  \right)'(0)   \right)
}
{ =} { \left(  D   \left(   \alpha_2 \circ \alpha_1^{-1}   \right)      \right)_{\alpha_1(P)} \left(   \left(  \alpha_1 \circ \gamma_2  \right)'(0)   \right)
}
{ =} { \left(  \alpha_2 \circ \gamma_2  \right)'(0)
}
}
{}
{}{.}

 }



\inputaufgabeklausurloesung
{5}

{

Tunjukkan bahwa pemetaan singgung $T(\varphi)$ dari
\maabbeledisp {\varphi} {\R^1} {S^1
} { t } {( \cos  t, \sin  t  )
} {,}
bersifat surjektif.

}
{

Pemetaan $\varphi$ bersifat surjektif, sehingga cukup ditunjukkan bahwa untuk setiap
\mathl{t \in \R}{} pemetaan singgung linier
\maabbdisp {} {T_t \R \cong \R} { T_{\varphi(t)} S^1
} {}
bersifat surjektif. Karena kedua ruang itu berdimensi satu, cukup ditunjukkan bahwa suatu vektor yang berbeda dari $0$ tidak dipetakan ke $0$. Vektor singgung di $t$ direalisasikan oleh lintasan terdiferensialkan
\maabbeledisp {\gamma} {\R} {\R
} { s } {t+s
} {.}
Lintasan komposisi
\maabbeledisp {\varphi \circ \gamma} {\R} {S^1
} { s } { ( \cos (t+s) , \sin (t+s) )
} {,}
mewujudkan vektor singgung citra; pada bidang ambien,
\zusatzklammer {
\mathl{T_{\varphi(t)}S^1 \subseteq \R^2}{}} {} {}

\mavergleichskettedisp
{\vergleichskette
{( \varphi \circ \gamma)' (0)
}
{ =} {(- \sin t , \cos t )
}
{ } {
}
{ } {
}
{ } {
}
}
{}{}{,}
dan vektor ini tidak nol.

 }



\inputaufgabeklausurloesung
{7 (3+4)}

{

Tinjau pemetaan diferensiabel
\maabbeledisp {\gamma} { [1,c] } {\R^2
} { t } {(t,t^{3})
} {,}
\zusatzklammer {dengan \mathlk{c \geq 1}{}} {} {}
dan
\maabbeledisp {\varphi} {\R_+ \times \R_+ } {\R^3
} {(u,v)} {(u^3,u^2+v^2,u^{-1}v^{-1} )
} {,}
serta bentuk diferensial

\mavergleichskettedisp
{\vergleichskette
{ \omega
}
{ =} { (x-y)dx-z^2dy+dz
}
{ } {
}
{ } {
}
{ } {
}
}
{}{}{}
pada $\R^3$.

a) Hitung bentuk diferensial tarik balik
\mathl{\varphi^* \omega}{} pada $\R_+ \times \R_+$.

b) Hitung integral lintasan dari bentuk diferensial
\mathl{\varphi^* \omega}{} sepanjang lintasan $\gamma$ sebagai fungsi dari $c$.

}
{

a) Bentuk diferensial tarik baliknya adalah

\mavergleichskettealigndrucklinks
{\vergleichskettealigndrucklinks
{\varphi^* (  (x-y)dx-z^2dy+dz   )
}
{ =} { (u^3-u^2-v^2 ) d(u^3)  -u^{-2}v^{-2} d(u^2+v^2) +d (u^{-1}v^{-1})
}
{ =} { 3 (u^5-u^4-u^2v^2 ) du  -2u^{-1}v^{-2} du -2u^{-2}v^{-1} dv - u^{-2} v^{-1}du - u^{-1} v^{-2}dv
}
{ =} { (3 u^5-3u^4-3u^2v^2 -2u^{-1}v^{-2}- u^{-2} v^{-1} )  du + (  -2u^{-2}v^{-1} - u^{-1} v^{-2})dv
}
{ } {
}
}
{}{}{.}

b) Integral lintasannya adalah

\mavergleichskettealigndrucklinks
{\vergleichskettealigndrucklinks
{ \int_1^c (3 t^5-3t^4-3t^2t^6 -2t^{-1}t^{-6}- t^{-2} t^{-3} )dt  +(  -2t^{-2}t^{-3} - t^{-1} t^{-6}) d t^3
}
{ =} { \int_1^c   (3 t^5-3t^4-3t^8-2t^{-7} - t^{-5})dt + ( -2t^{-5} - t^{-7})3t^2 dt
}
{ =} { \int_1^c (3 t^5-3t^4-3t^8-2t^{-7} - t^{-5}   -6t^{-3} - 3 t^{-5}      ) dt
}
{ =} { \int_1^c (-3t^8 + 3 t^5-3t^4 -6t^{-3} - 4 t^{-5}  -2t^{-7}       )  dt
}
{ =} { (-{ \frac{   1 }{  3 } } t^9 + { \frac{   1 }{  2 } } t^{6}  - { \frac{   3 }{  5 } } t^5  +  3 t^{-2} +  t^{-4} + { \frac{   1 }{  3 } }  t^{-6}     ) \vert_{  1 }^{  c  }
}
}
{
\vergleichskettefortsetzungalign
{ =} { -{ \frac{   1 }{  3 } } c^9 + { \frac{   1 }{  2 } } c^{6}  - { \frac{   3 }{  5 } } c^5  +  3c^{-2} +   c^{-4} + { \frac{   1 }{  3 } }  c^{-6}
-(-{ \frac{   1 }{  3 } } + { \frac{   1 }{  2 } } - { \frac{   3 }{  5 } } +3 + 1 + { \frac{   1 }{  3 } }  )
}
{ =} { -{ \frac{   1 }{  3 } } c^9 + { \frac{   1 }{  2 } } c^{6}  - { \frac{   3 }{  5 } } c^5 + 3 c^{-2} +  c^{-4} + { \frac{   1 }{  3 } } c^{-6} - { \frac{   39 }{  10 } }  }
{ } {}
{ } {}
}{}{.}

 }



\inputaufgabeklausurloesung
{3}

{

Misalkan
\maabbdisp {\psi} { L } { M
} {}
adalah
\definitionsverweis {pemetaan diferensiabel}{}{}
antara
\definitionsverweis {manifold diferensiabel}{}{.}
Selanjutnya, misalkan
\maabbdisp {f} { M } {\R
} {}
adalah fungsi pada $M$, dan $\omega$ adalah
$k$-\definitionsverweis {bentuk}{}{}
pada $M$. Tunjukkan bahwa

\mavergleichskettedisp
{\vergleichskette
{ \psi^*(f \omega)
}
{ =} { \psi^*(f)  \psi^* \omega
}
{ } {
}
{ } {
}
{ } {
}
}
{}{}{.}

}
{

Misalkan
\mathl{P \in L}{} dan
\mathl{v_1 , \ldots , v_k \in T_PL}{.} Maka

\mavergleichskettealign
{\vergleichskettealign
{ \psi^* \left(  f \omega) (P,v_1 , \ldots , v_k  \right)
}
{ =} { (f \omega) \left(  \psi(P), T_P\psi (v_1) , \ldots , T_P\psi (v_k )       \right)
}
{ =} { f(\psi(P))   \cdot \omega \left(  \psi(P), T_P\psi (v_1) , \ldots , T_P\psi (v_k )       \right)
}
{ =} { \left(  \psi^* (f)  \right) (P)  \cdot \omega \left(  \psi(P), T_P\psi (v_1) , \ldots , T_P\psi (v_k )       \right)
}
{ =} { \left(  \psi^* (f)  \right) (P)   \cdot \left(  \psi^*( \omega)  \right)  \left(  P,v_1 , \ldots , v_k  \right)
}
}
{
\vergleichskettefortsetzungalign
{ =} { \left(  \psi^* (f)  \psi^*( \omega)  \right) \left(  P,v_1 , \ldots , v_k  \right)
}
{ } {}
{ } {}
{ } {}
}
{}{.}

 }



\inputaufgabeklausurloesung
{2}

{

Deskripsikan sebanyak mungkin
\definitionsverweis {manifold berbatas}{}{}
berdimensi satu yang
\definitionsverweis {terhubung}{}{.}

}
{  }



\inputaufgabeklausurloesung
{5}

{

Misalkan
\maabbdisp {f} { [a,b] } { \R_{+}
} {}
adalah
\definitionsverweis {fungsi yang terdiferensialkan secara kontinu}{}{.}
Kita pandang
\definitionsverweis {subgraf}{}{}
sebagai
\definitionsverweis {manifold berbatas}{}{}.
Batasnya terdiri dari grafik, interval dasar, dan kedua sisi tegak, tetapi tidak mencakup keempat titik sudut. Buktikan secara langsung berlakunya teorema Stokes untuk
\definitionsverweis {bentuk diferensial}{}{}

\mavergleichskette
{\vergleichskette
{ \omega
}
{ = }{ xdy
}
{  }{
}
{    }{
}
{   }{
}
}
{}{}{.}

}
{

Karena

\mavergleichskettedisp
{\vergleichskette
{d(xdy)
}
{ =} { dx \wedge dy
}
{ } {
}
{ } {
}
{ } {
}
}
{}{}{}
maka, di satu pihak,

\mavergleichskettedisp
{\vergleichskette
{ \int_M dx \wedge dy
}
{ =} { \int_a^b f(t) dt
}
{ =} { F(b)-F(a)
}
{ } {
}
{ } {
}
}
{}{}{}
adalah luas subgraf, dengan $F$ suatu antiturunan dari $f$. Untuk integral lintasan, batas $M$ harus diparameterkan berlawanan arah jarum jam. Interval dasarnya diparameterkan oleh
\mathl{i: t \mapsto \begin{pmatrix} t \\ 0 \end{pmatrix}}{}, dan

\mavergleichskettedisp
{\vergleichskette
{ i^*\omega
}
{ =} { i^*xdy
}
{ =} { x d0
}
{ =} { 0
}
{ } {
}
}
{}{}{,}
sehingga bagian integral lintasan ini bernilai $0$. Sisi kanan diparameterkan oleh
$i: t \mapsto \begin{pmatrix}  b  \\t   \end{pmatrix}$,
; di sini $t$ bergerak antara
\mathkor {} {0} {dan} {f(b)} {}
. Maka

\mavergleichskettedisp
{\vergleichskette
{ i^*xdy
}
{ =} { b dt
}
{ } {
}
{ } {
}
{ } {
}
}
{}{}{}
dan integral lintasan pada sisi ini adalah
\mathl{b f(b)}{.} Kontribusi sisi kiri adalah
\mathl{-af(a)}{.} Grafik diparameterkan oleh
\maabbeledisp {i} {[a,b]} {
} { t } {(t, f(t))
} {,}
. Maka

\mavergleichskettedisp
{\vergleichskette
{ i^*xdy
}
{ =} { t df
}
{ =} { tf'(t) dt
}
{ } {
}
{ } {
}
}
{}{}{.}
Dengan orientasi yang benar, integralnya adalah

\mavergleichskettealign
{\vergleichskettealign
{ - \int_a^b tf'(t)dt
}
{ =} { - ( tf(t) - F (t) ) \mid _a^b
}
{ =} { - ( bf(b)-F(b) -af(a) +F(a) )
}
{ =} { F(b)-F(a) +af(a) - bf(b)
}
{ } {
}
}
{}
{}{.}
Jadi, seluruh integral lintasan juga bernilai

\mavergleichskettealign
{\vergleichskettealign
{bf(b) - af(a)+ F(b)-F(a) +af(a) - bf(b)
}
{ =} { F(b)-F(a)
}
{ } {
}
{ } {
}
{ } {
}
}
{}
{}{.}

 }



\inputaufgabeklausurloesung
{9}

{

Buktikan teorema titik tetap Brouwer.

}
{

Untuk menyederhanakan notasi, misalkan

\mavergleichskette
{\vergleichskette
{ r
}
{ = }{ 1
}
{ }{
}
{ }{
}
{ }{
}
}
{}{}{.}  Andaikan terdapat pemetaan terdiferensialkan secara kontinu $\psi$ tanpa titik tetap. Maka selalu berlaku

\mavergleichskettedisp
{\vergleichskette
{ x
}
{ \neq} { \psi(x)
}
{ } {
}
{ } {
}
{ } {
}
}
{}{}{,}
sehingga kedua titik itu menentukan sebuah garis. Gagasannya adalah menggunakan garis ini untuk memilih
\zusatzklammer {dari keduanya} {} {}
salah satu dari kedua titik potongnya dengan sfera sebagai titik citra suatu retraksi ke batas. Dengan fungsi bantu

\mavergleichskettedisp
{\vergleichskette
{ h(x)
}
{ =} { { \frac{ \psi(x)-x }{ \Vert { \psi (x)-x} \Vert } }
}
{ } {
}
{ } {
}
{ } {
}
}
{}{}{}
kita definisikan pemetaan
\maabbdisp {\varphi} { B \left( 0 , 1 \right) } { \R^n
} {}
melalui

\mavergleichskettedisphandlinks
{\vergleichskettedisphandlinks
{ \varphi(x)
}
{ =} { x + \left(  - \left\langle  x  , h(x)  \right\rangle + \sqrt{1 + \left\langle  x  , h(x)  \right\rangle^2 - \Vert { x} \Vert^2  }  \right) \cdot h(x)
}
{ } {
}
{ } {
}
{ } {
}
}
{}{}{.}
Ekspresi di bawah tanda akar bernilai positif. Untuk

\mavergleichskette
{\vergleichskette
{ \Vert { x } \Vert
}
{ < }{ 1
}
{ }{
}
{ }{
}
{ }{
}
}
{}{}{}
hal ini jelas; sedangkan untuk

\mavergleichskette
{\vergleichskette
{ \Vert { x } \Vert
}
{ = }{ 1
}
{ }{
}
{ }{
}
{ }{
}
}
{}{}{}
terdapat titik pada sfera yang garis penghubungnya dengan titik dalam bola
\mathl{\psi(x)}{} tidak tegak lurus terhadap $x$
\zusatzklammer {ruang singgung afin pada suatu titik sfera memotong bola hanya di satu titik} {} {,}
sehingga

\mavergleichskettedisp
{\vergleichskette
{ \left\langle x , h(x) \right\rangle
}
{ \neq} { 0
}
{ } {
}
{ } {
}
{ } {
}
}
{}{}{}
berlaku. Karena fungsi akar kuadrat dan norma di luar titik nol
\definitionsverweis {terdiferensialkan secara kontinu}{}{}
bersifat demikian, maka
\mathkor {} {h(x)} {dan} {\varphi(x)} {}
merupakan pemetaan-pemetaan terdiferensialkan secara kontinu. Menurut
Soal 23.23 (Differentialgeometrie (Osnabrück 2023)),
pemetaan $\varphi$ memetakan bola ke sfera dan pembatasannya pada sfera adalah identitas. Jadi terdapat suatu
\definitionsverweis {Retraksi}{}{}
bola pejal tertutup ke batasnya, yang menurut
Teorema 23.9 (Differentialgeometrie (Osnabrück 2023))
mustahil ada.

 }



\inputaufgabeklausurloesung
{5 (3+1+1)}

{

Tinjau silinder

\mavergleichskette
{\vergleichskette
{ Z
}
{ = }{ S^1 \times \R
}
{ \subseteq }{ \R^2 \times \R
}
{    }{
}
{   }{
}
}
{}{}{}
dengan
\definitionsverweis {koneksi trivial}{}{}
pada
\definitionsverweis {bundel tangen}{}{.}
\aufzaehlungdrei{Tentukan, dalam suatu peta isometrik yang sesuai,
\definitionsverweis {persamaan diferensial geodesik}{}{}
untuk kurva geodesik $\psi$ pada $M$ yang memenuhi

\mavergleichskettedisp
{\vergleichskette
{ \psi(0)
}
{ =} { (1,0,3)
}
{ } {
}
{ } {
}
{ } {
}
}
{}{}{}
dan

\mavergleichskettedisp
{\vergleichskette
{ \psi'(0)
}
{ =} { (0,1,-4)
}
{ } {
}
{ } {
}
{ } {
}
}
{}{}{}
.
}{Selesaikan persamaan diferensial dari (1) dalam peta tersebut.
}{Carilah suatu kurva geodesik di $Z$ yang memenuhi syarat-syarat pada (1).
}

}
{

\aufzaehlungdrei{Kita tinjau peta isometrik
\zusatzklammer {inverse} {} {}
peta
\maabbeledisp {\varphi} { ]- \pi, \pi[  \times \R } { S^1 \setminus \{  (-1,0) \} \times \R
} {(t,v)} { \left(  \cos  t   , \,   \sin  t  , \,  v                 \right)
} {.}
Karena

\mavergleichskette
{\vergleichskette
{ \varphi(0,3)
}
{ = }{ (1,0,3)
}
{ }{
}
{ }{
}
{ }{
}
}
{}{}{.}
simbol-simbol Christoffelnya nol, persamaan diferensial geodesiknya adalah

\mavergleichskettedisp
{\vergleichskette
{ \gamma^{\prime \prime}
}
{ =} { 0
}
{ } {
}
{ } {
}
{ } {
}
}
{}{}{.}
Syarat-syarat yang diberikan dalam peta ini menjadi

\mavergleichskettedisp
{\vergleichskette
{ \gamma(0)
}
{ =} { (0,3)
}
{ } {
}
{ } {
}
{ } {
}
}
{}{}{}
dan

\mavergleichskettedisp
{\vergleichskette
{ \gamma'(0)
}
{ =} { (1,-4)
}
{ } {
}
{ } {
}
{ } {
}
}
{}{}{.}
}{Suatu solusi dalam peta tersebut adalah

\mavergleichskettedisp
{\vergleichskette
{ \gamma(t)
}
{ =} { (0,3) + t (1,-4)
}
{ =} { (t,3-4t)
}
{ } {
}
{ } {
}
}
{}{}{.}
}{Karena itu, kurva geodesik dengan syarat awal tersebut adalah

\mavergleichskettedisp
{\vergleichskette
{  \psi(t)
}
{ =} { \varphi( \gamma(t))
}
{ =} { \left(  \cos  t   , \,   \sin  t  , \,   3-4t                 \right)
}
{ } {
}
{ } {
}
}
{}{}{.}
}

 }

\endgroup

\chapter{Formulir Solusi Resmi Ujian 2}
\noindent\textit{Solusi yang disediakan oleh sumber resmi; kekosongan sumber dipertahankan.}
\begingroup
\renewcommand{\theHfakt}{o011.exam.official.02.\arabic{fakt}}
%Data institusi

%\input{Dozentdaten}

%\renewcommand{\fachbereich}{Fachbereich}

%\renewcommand{\dozent}{Prof. Dr. . }

%Data ujian

\renewcommand{\klausurgebiet}{ }

\renewcommand{\klausurtyp}{ }

\renewcommand{\klausurdatum}{ . 20}

\klausurvorspann {\fachbereich} {\klausurdatum} {\dozent} {\klausurgebiet} {\klausurtyp}

%Data untuk tabel nilai berikut

\renewcommand{\aeins}{  3  }

\renewcommand{\azwei}{  3   }

\renewcommand{\adrei}{  8  }

\renewcommand{\avier}{  7  }

\renewcommand{\afuenf}{  8  }

\renewcommand{\asechs}{  3  }

\renewcommand{\asieben}{  5  }

\renewcommand{\aacht}{  5  }

\renewcommand{\aneun}{  6  }

\renewcommand{\azehn}{  4  }

\renewcommand{\aelf}{  12  }

\renewcommand{\azwoelf}{  64  }

\renewcommand{\adreizehn}{     }

\renewcommand{\avierzehn}{    }

\renewcommand{\afuenfzehn}{    }

\renewcommand{\asechzehn}{    }

\renewcommand{\asiebzehn}{    }

\renewcommand{\aachtzehn}{    }

\renewcommand{\aneunzehn}{    }

\renewcommand{\azwanzig}{    }

\renewcommand{\aeinundzwanzig}{    }

\renewcommand{\azweiundzwanzig}{    }

\renewcommand{\adreiundzwanzig}{    }

\renewcommand{\avierundzwanzig}{    }

\renewcommand{\afuenfundzwanzig}{    }

\renewcommand{\asechsundzwanzig}{    }

\punktetabelleelf

\klausurnote

\newpage

\setcounter{bagian}{0}



\inputaufgabeklausurloesung
{3}

{

Definisikan istilah-istilah berikut
\zusatzklammer {yang dicetak miring} {} {.}
\aufzaehlungsechs{Suatu
\stichwort {medan normal} {}
$N$ pada hiperpermukaan diferensiabel

\mavergleichskette
{\vergleichskette
{ Y
}
{ \subseteq }{ \R^n
}
{  }{
}
{    }{
}
{   }{
}
}
{}{}{.}

}{Suatu ruang topologis yang
\stichwort {parakompak}{}
.

}{
\stichwort {Ruang kotangen} {}
di suatu titik

\mavergleichskette
{\vergleichskette
{  P
}
{ \in }{ M
}
{  }{}
{    }{}
{   }{}
}
{}{}{}
dari
\definitionsverweis {manifold diferensiabel}{}{}
$M$.

}{Suatu
\stichwort {titik reguler} {}

\mathl{Q \in L}{} bagi
\definitionsverweis {pemetaan diferensiabel}{}{}
\maabbdisp {\varphi} { L } { M
} {}
antara
\definitionsverweis {manifold-manifold diferensiabel}{}{}
\mathkor {} {L} {dan} {M} {.}

}{
\stichwort {Produk} {}
dari manifold-manifold diferensiabel
\mathkor {} {M} {dan} {N} {.}

}{Suatu
\stichwort {koneksi} {}
pada bundel vektor diferensiabel
\maabb {} {E} {M
} {}
di atas manifold diferensiabel $M$.
}

}
{

\aufzaehlungsechs{Suatu
medan normal
pada $Y$ adalah objek yang didefinisikan pada suatu lingkungan terbuka

\mavergleichskette
{\vergleichskette
{ Y
}
{ \subseteq }{ U
}
{ \subseteq }{ \R^n
}
{ }{
}
{ }{
}
}
{}{}{}
dan bersifat
\definitionsverweis {kontinu}{}{}
sebagai suatu
\definitionsverweis {medan vektor}{}{}
\maabb {F} {U} {\R^n
} {}
yang memenuhi

\mavergleichskettedisp
{\vergleichskette
{ F(P)
}
{ \in} { N_PY
}
{ } {
}
{ } {
}
{ } {
}
}
{}{}{}
untuk setiap

\mavergleichskette
{\vergleichskette
{ P
}
{ \in }{ Y
}
{  }{
}
{    }{
}
{   }{
}
}
{}{}{.}
}{Suatu
\definitionsverweis {ruang topologis}{}{}
$X$ disebut kompak terhadap penutup jika untuk setiap penutup terbuka

\mathdisp {X= \bigcup_{i \in I} U_i \, \, \, \text{ dengan } U_i \text{ terbuka dan himpunan indeks sembarang }} {  }

terdapat suatu himpunan bagian berhingga
\mathl{J \subseteq I}{} sedemikian sehingga

\mavergleichskettedisp
{\vergleichskette
{ X
}
{ =} {\bigcup_{i \in J} U_i
}
{ } {
}
{ } {
}
{ } {
}
}
{}{}{}
berlaku.
}{Ruang kotangen di $P$ adalah
\definitionsverweis {ruang dual}{}{}
dari
\definitionsverweis {ruang singgung}{}{}

\mathl{T_PM}{} di $P$.
}{Titik

\mavergleichskette
{\vergleichskette
{ Q
}
{ \in }{ L
}
{ }{
}
{ }{
}
{ }{
}
}
{}{}{}
disebut reguler bagi $\varphi$ jika
\definitionsverweis {pemetaan tangen}{}{}
\maabbdisp {T_Q(\varphi)} {T_QL} {T_{\varphi(Q)}M
} {}
pada titik $Q$ mempunyai
\definitionsverweis {rank maksimal}{}{.}
}{Misalkan
\mathkor {} {(U_i,U_i',\alpha_i, i \in I)} {dan} {(V_j,V_j',\beta_j, j \in J)} {}
merupakan
\definitionsverweis {atlas-atlas}{}{}
dari
\mathkor {} {M} {dan} {N} {.}
Maka
\definitionsverweis {ruang produk}{}{}

\mathl{M \times N}{} yang dilengkapi dengan
\definitionsverweis {peta-peta}{}{}
\maabbdisp {\alpha_i \times \beta_j} {U_i \times V_j } { U_i' \times V_j'
} {}
\zusatzklammer {dengan \mathlk{(i,j) \in I \times J}{} dan \mathlk{U_i' \times V_j' \subseteq \R^m \times \R^n}{}} {} {}
disebut produk manifold
\mathkor {} {M} {dan} {N} {.}
}{Suatu
koneksi
pada $E$ adalah dekomposisi jumlah langsung dari
\definitionsverweis {bundel tangen}{}{}

\mavergleichskette
{\vergleichskette
{ TE
}
{ = }{ V \oplus H
}
{ }{
}
{ }{
}
{ }{
}
}
{}{}{}
menjadi dua subbundel vektor
\mathkor {} {V} {dan} {H} {,}
dengan

\mavergleichskette
{\vergleichskette
{ V
}
{ \subseteq }{ TE
}
{ }{
}
{ }{
}
{ }{
}
}
{}{}{}
merupakan
\definitionsverweis {bundel vertikal}{}{.}
}

 }



\inputaufgabeklausurloesung
{3}

{

Nyatakan teorema-teorema berikut.
\aufzaehlungdrei{
\stichwort {Teorema tentang hubungan antara kelengkungan dan vektor normal satuan suatu kurva bidang berparameter panjang busur} {.}}{
\stichwort {Theorema egregium} {.}}{
\stichwort {Teorema Green} {}
untuk luas.}

}
{

\aufzaehlungdrei{\faktsituation {Misalkan
\maabb {\gamma} { I } {\R^2
} {}
suatu
\definitionsverweis {kurva}{}{}
\definitionsverweis {berparameter panjang busur}{}{}
yang dua kali
\definitionsverweis {terdiferensialkan secara kontinu}{}{.}}
\faktfolgerung {Maka
\definitionsverweis {kelengkungan}{}{}
dari $\gamma$ di $t$ adalah

\mavergleichskettedisp
{\vergleichskette
{ \kappa(t)
}
{ =} { \left\langle \gamma^{\prime \prime}(t) , N(t) \right\rangle
}
{ = } { \pm \Vert { \gamma^{\prime \prime}(t) } \Vert
}
{ } {
}
{ } {
}
}
{}{}{,}
dengan

\mavergleichskettedisp
{\vergleichskette
{ N(t)
}
{ =} { \begin{pmatrix} -\gamma_2'(t) \\ \gamma_1'(t) \end{pmatrix}
}
{ } {
}
{ } {
}
{ } {
}
}
{}{}{}
suatu vektor normal satuan di $\gamma(t)$.}
\faktzusatz {}
\faktzusatz {}}{\faktsituation {Misalkan

\mavergleichskette
{\vergleichskette
{ Y
}
{ \subseteq }{ \R^3
}
{  }{
}
{    }{
}
{   }{
}
}
{}{}{}
suatu
\definitionsverweis {permukaan berorientasi}{}{}
dan misalkan
\maabbdisp {\varphi} { V } { Y
} {,}

\mavergleichskette
{\vergleichskette
{ V
}
{ \subseteq }{ \R^2
}
{  }{
}
{    }{
}
{   }{
}
}
{}{}{,}
suatu parametrisasi lokal $Y$ yang dua kali terdiferensialkan, dengan parameter $u,v$. Misalkan
\mathl{\left(  g_{ij}  \right)}{}
\definitionsverweis {matriks fundamental pertama}{}{}
pada $V$.}
\faktfolgerung {Maka, dengan menggunakan
\definitionsverweis {simbol-simbol Christoffel}{}{,}
\definitionsverweis {kelengkungan Gauss}{}{}
memenuhi hubungan

\mavergleichskettedisphandlinks
{\vergleichskettedisphandlinks
{ K
}
{ =} { { \frac{ 1 }{ g_{11} } } \left( \partial_2 \Gamma_{11}^2 - \partial_1 \Gamma_{12}^2 + \Gamma^1_{11} \Gamma^2_{12} + \Gamma_{11}^2 \Gamma_{22}^2 - \Gamma_{12}^1 \Gamma_{11}^2- \left( \Gamma_{12}^2 \right)^2 \right)
}
{ } {
}
{ } {
}
{ } {
}
}
{}{}{.}}
\faktzusatz {}
\faktzusatz {}}{Misalkan
\mathl{M \subseteq \R^2}{} suatu
\definitionsverweis {manifold dengan batas}{}{}
yang
\definitionsverweis {kompak}{}{}

\mathl{\delta M}{.} Maka luas $M$ adalah

\mavergleichskettedisp
{\vergleichskette
{  \lambda^2(M)
}
{ =} {
\int_{ \partial M  }  x dy
}
{ =} { -
\int_{ \partial M  }  y dx
}
{ } {
}
{ } {
}
}
{}{}{.}}

 }



\inputaufgabeklausurloesung
{8}

{

Misalkan

\mavergleichskettedisp
{\vergleichskette
{ f(x,y)
}
{ =} { ax^2+by^2+cxy
}
{ } {
}
{ } {
}
{ } {
}
}
{}{}{.}
Tentukan
\definitionsverweis {kelengkungan-kelengkungan utama}{}{}
dan
\definitionsverweis {arah-arah kelengkungan utama}{}{}
dari
\definitionsverweis {grafik}{}{}
$f$ di titik asal.

}
{

Kita menggunakan
Lemma 4.9 (Differentialgeometrie (Osnabrück 2023)),
dengan faktor pendahulu di titik nol sama dengan $1$, sehingga
\definitionsverweis {pemetaan Weingarten}{}{}
dideskripsikan langsung oleh
\definitionsverweis {matriks Hesse}{}{}
. Matriks Hesse fungsi tersebut adalah

\mathdisp {\begin{pmatrix} 2a & c \\ c & 2b \end{pmatrix}} { , }

dan
\definitionsverweis {polinomial karakteristik}{}{}
matriks tersebut adalah

\mavergleichskettealign
{\vergleichskettealign
{ (t-2a)(t-2b)-c^2
}
{ =} { t^2 -2(a+b)t +4ab -c^2
}
{ =} { \left(  t-(a+b)  \right)^2 - (a+b)^2 +4ab -c^2
}
{ =} { \left(  t-(a+b)  \right)^2  - a^2 -b^2 +2ab -c^2
}
{ =} { \left(  t-(a+b)  \right)^2  - (a-b)^2 -c^2
}
}
{}
{}{.}
Dengan demikian, nilai-nilai eigen
\zusatzklammer {yang merupakan kelengkungan-kelengkungan utama} {} {}
adalah

\mavergleichskettedisp
{\vergleichskette
{ \lambda_{1,2}
}
{ =} { \pm \sqrt{ (a-b)^2+c^2 } +a+b
}
{ } {
}
{ } {
}
{ } {
}
}
{}{}{.}
Dalam kasus

\mavergleichskette
{\vergleichskette
{ a
}
{ = }{ b
}
{  }{
}
{    }{
}
{   }{
}
}
{}{}{}
dan

\mavergleichskette
{\vergleichskette
{c
}
{ = }{ 0
}
{ }{
}
{ }{
}
{ }{
}
}
{}{}{}
pemetaan Weingarten adalah perkalian dengan $2a$ dan setiap vektor merupakan vektor eigen; mulai sekarang kasus ini dikecualikan. Untuk

\mavergleichskettedisp
{\vergleichskette
{ \lambda_{1}
}
{ =} { \sqrt{ (a-b)^2+c^2 } +a+b
}
{ } {
}
{ } {
}
{ } {
}
}
{}{}{}
perlu ditentukan kernel matriks

\mathdisp {\begin{pmatrix}  \sqrt{(a-b)^2 +c^2} + b-a  &   -c  \\  -c &  \sqrt{(a-b)^2 +c^2} + a-b \end{pmatrix}} {  }

dan kernel tersebut adalah

\mathdisp {\R \begin{pmatrix}  \sqrt{(a-b)^2 +c^2} + a-b  \\c   \end{pmatrix}} { . }

Jadi, salah satu arah kelengkungan utama adalah
\mathl{\begin{pmatrix}  \sqrt{(a-b)^2 +c^2} + a-b  \\c   \end{pmatrix}}{} dengan kelengkungan utama
\mathl{\sqrt{ (a-b)^2+c^2 } +a+b}{.}

Dengan cara yang sama diperoleh arah kelengkungan utama
\mathl{\begin{pmatrix} -c \\ \sqrt{(a-b)^2 +c^2} + a-b \end{pmatrix}}{} dengan kelengkungan utama
\mathl{-\sqrt{ (a-b)^2+c^2 } +a+b}{.}

 }



\inputaufgabeklausurloesung
{7}

{

Buktikan teorema eksistensi untuk transpor paralel pada suatu hiperpermukaan

\mavergleichskette
{\vergleichskette
{ Y
}
{ \subseteq }{ \R^n
}
{  }{
}
{    }{
}
{   }{
}
}
{}{}{.}

}
{

Kita mencari suatu pemetaan diferensiabel
\maabbdisp {F} { I } { \R^n
} {,}
yang memenuhi syarat-syarat
\aufzaehlungdrei{
\mavergleichskettedisp
{\vergleichskette
{ F(t)
}
{ \in} { T_{\gamma(t)} Y
}
{ } {
}
{ } {
}
{ } {
}
}
{}{}{}
untuk setiap

\mavergleichskette
{\vergleichskette
{ t
}
{ \in }{ I
}
{  }{
}
{    }{
}
{   }{
}
}
{}{}{,}
}{
\mavergleichskettedisp
{\vergleichskette
{ F'(t)
}
{ \in} { N_{\gamma(t)} Y
}
{ } {
}
{ } {
}
{ } {
}
}
{}{}{}
untuk setiap

\mavergleichskette
{\vergleichskette
{ t
}
{ \in }{ I
}
{ }{
}
{ }{
}
{ }{
}
}
{}{}{,}
}{
\mavergleichskettedisp
{\vergleichskette
{ F (t_0)
}
{ =} { v
}
{ } {
}
{ } {
}
{ } {
}
}
{}{}{}
}
Syarat pertama berarti bahwa $F(t)$ tegak lurus terhadap vektor normal satuan $N (\gamma(t))$, sehingga

\mavergleichskettedisp
{\vergleichskette
{ \left\langle F( t) , N( \gamma(t)) \right\rangle
}
{ =} { 0
}
{ } {
}
{ } {
}
{ } {
}
}
{}{}{.}
Oleh karena itu, juga berlaku

\mavergleichskettedisp
{\vergleichskette
{ 0
}
{ =} { \left\langle F(t) , N( \gamma(t)) \right\rangle'
}
{ =} { \left\langle F'(t) , N( \gamma(t)) \right\rangle + \left\langle F(t) , (N( \gamma(t)))' \right\rangle
}
{ } {
}
{ } {
}
}
{}{}{.}
Syarat kedua, yakni bahwa $F'(t)$ terletak di ruang normal, berarti

\mavergleichskettedisp
{\vergleichskette
{ F'(t) - \left\langle F'( t) , N( \gamma(t)) \right\rangle N(\gamma(t) )
}
{ =} { 0
}
{ } {
}
{ } {
}
{ } {
}
}
{}{}{}
untuk setiap

\mavergleichskette
{\vergleichskette
{ t
}
{ \in }{ I
}
{  }{
}
{    }{
}
{   }{
}
}
{}{}{,}
Dengan menggunakan persamaan sebelumnya, syarat ini dapat diubah menjadi

\mavergleichskettedisp
{\vergleichskette
{ F'(t) + \left\langle F( t) , (N( \gamma(t)))' \right\rangle N(\gamma(t) )
}
{ =} { 0
}
{ } {
}
{ } {
}
{ } {
}
}
{}{}{}
untuk setiap

\mavergleichskette
{\vergleichskette
{ t
}
{ \in }{ I
}
{  }{
}
{    }{
}
{   }{
}
}
{}{}{}
atau menjadi

\mavergleichskettedisp
{\vergleichskette
{ F'(t)
}
{ =} { - \left\langle F( t) ,  (N( \gamma(t)))'   \right\rangle N(\gamma(t) )
}
{ } {
}
{ } {
}
{ } {
}
}
{}{}{}
untuk setiap

\mavergleichskette
{\vergleichskette
{ t
}
{ \in }{ I
}
{  }{
}
{    }{
}
{   }{
}
}
{}{}{}
. Ini adalah persamaan diferensial linear biasa orde pertama untuk $F$. Bersama syarat awal

\mavergleichskette
{\vergleichskette
{ F(t_0)
}
{ = }{ v
}
{ }{
}
{ }{
}
{ }{
}
}
{}{}{}
menurut
Satz 56.2 (Analysis (Osnabrück 2021-2023))
persamaan tersebut mempunyai solusi yang tunggal. Jadi, sistem syarat semula mempunyai paling banyak satu solusi. Jika $F$ adalah solusi tunggal persamaan diferensial itu, maka pemetaan tersebut memang merupakan solusi sistem syarat semula.

 }



\inputaufgabeklausurloesung
{8}

{

Tunjukkan bahwa suatu
\definitionsverweis {manifold diferensiabel}{}{}
$M$ berdimensi

\mavergleichskette
{\vergleichskette
{ n
}
{ \geq }{ 1
}
{  }{
}
{    }{
}
{   }{
}
}
{}{}{}
mempunyai tak hingga banyak
\definitionsverweis {difeomorfisme}{}{}
\maabbdisp {\varphi} { M } { M
} {}
.

}
{

Kita meninjau fungsi-fungsi berbentuk
\maabbeledisp {\psi_c} {[0,1]} {[0,1]
} { x } { x + \varphi_c (x)
} {,}
dengan

\mavergleichskettedisp
{\vergleichskette
{  \varphi_c(x)
}
{ =} { \begin{cases}  0 \text{ jika } x \leq { \frac{   1  }{  3 } } \, , \\  c \left(  x- { \frac{   1  }{  3 } }   \right)^2 \left(  x- { \frac{   2 }{  3 } }   \right)^2 \text{ jika } { \frac{   1  }{  3 } } < x < { \frac{   2 }{  3 } }  \, ,  \\  0 \text{ jika } x \geq { \frac{   2 }{  3 } }  \, .  \end{cases}
}
{ } {
}
{ } {
}
{ } {
}
}
{}{}{}
dengan parameter

\mavergleichskette
{\vergleichskette
{ c
}
{ \in }{  \R_{\geq 0}
}
{  }{
}
{    }{
}
{   }{
}
}
{}{}{.}
Fungsi $\varphi_c$ kontinu dan juga terdiferensialkan secara kontinu karena akar-akar polinomial tersebut merupakan akar ganda. Dengan demikian, $\psi_c$ juga terdiferensialkan secara kontinu. Untuk $c$ yang cukup kecil, $\psi_c$ meningkat secara ketat dan mendefinisikan suatu pemetaan bijektif yang terdiferensialkan secara kontinu dari interval satuan ke dirinya sendiri, serta merupakan identitas pada sepertiga bawah dan sepertiga atas.

Dengan $\psi_c$, kita mendefinisikan difeomorfisme bola satuan terbuka ke dirinya sendiri melalui
\maabbeledisp {} { U \left(   0 , 1  \right) } { U \left(   0 , 1  \right)
} { \begin{pmatrix}  x_1  \\\vdots \\  x_n   \end{pmatrix} } { \psi_c \left(  \sqrt{  x_1^2 + \cdots + x_n^2 }  \right) \begin{pmatrix}  x_1  \\\vdots \\  x_n   \end{pmatrix}
} {.}
Jadi, titik-titik diregangkan bergantung pada jari-jarinya, sedangkan pemetaan tersebut merupakan identitas pada
\mathl{U \left(   0 , { \frac{   1  }{  3 } }   \right)}{} dan pada
\mathl{U \left(   0 , 1  \right) \setminus  U \left(   0 ,   { \frac{   2 }{  3 } }   \right)}{}.

Jika sekarang $M$ suatu manifold berdimensi $\geq 1$, kita memilih suatu peta
\maabbdisp {\alpha} { U } {V
} {}
dengan

\mavergleichskette
{\vergleichskette
{ V
}
{ \subseteq }{ \R^n
}
{ }{
}
{ }{
}
{ }{
}
}
{}{}{,}
dan, dengan
\zusatzklammer {memperkecil daerahnya} {} {,}
kita dapat menganggap bahwa $V$
\zusatzklammer {merupakan suatu bola terbuka dan kemudian} {} {}
adalah bola satuan terbuka. Difeomorfisme yang dikonstruksi pada $V$ menghasilkan difeomorfisme pada $U$. Karena difeomorfisme ini merupakan identitas pada bagian luar bola
\zusatzklammer {mulai dari jari-jari ${ \frac{   2 }{  3 } }$} {} {,}
kita dapat memperluasnya secara diferensiabel ke $M$ dengan menggunakan identitas pada $M \setminus \alpha^{-1} \left(  U \left(   0 , { \frac{   2 }{  3 } }   \right)  \right)$.

 }



\inputaufgabeklausurloesung
{3}

{

Misalkan $M$ suatu
\definitionsverweis {manifold diferensiabel}{}{}
dan
\maabbdisp {f} { M } {\R
} {}
suatu
\definitionsverweis {fungsi terdiferensialkan secara kontinu}{}{.}
Misalkan di titik

\mavergleichskette
{\vergleichskette
{ P
}
{ \in }{ M
}
{  }{
}
{    }{
}
{   }{
}
}
{}{}{}
fungsi tersebut mempunyai
\definitionsverweis {ekstremum lokal}{}{.}
Tunjukkan bahwa

\mavergleichskettedisp
{\vergleichskette
{ (df)_P
}
{ =} { 0
}
{ } {
}
{ } {
}
{ } {
}
}
{}{}{.}

}
{

Misalkan

\mavergleichskette
{\vergleichskette
{ P
}
{ \in }{ U
}
{ \subseteq }{ M
}
{ }{
}
{ }{
}
}
{}{}{}
suatu
\definitionsverweis {daerah peta}{}{}
dengan peta
\maabb {\alpha} { U } {V \subseteq \R^n
} {.}
Maka
\mathl{f \circ \alpha^{-1}}{} juga mempunyai ekstremum lokal di $\alpha(P)$.
Oleh karena itu, menurut
Fakt ***** (2)
\definitionsverweis {diferensial total}{}{}
\maabbdisp {\left(D f \circ \alpha^{-1} \right)_{\alpha(P) }} {\R^n } {\R
} {}
adalah pemetaan nol. Karena itu,

\mavergleichskette
{\vergleichskette
{ (df)_P
}
{ = }{ 0
}
{ }{
}
{ }{
}
{ }{
}
}
{}{}{.}

 }



\inputaufgabeklausurloesung
{5}

{

Kita tinjau
$1$-\definitionsverweis {bentuk diferensial}{}{}

\mavergleichskettedisp
{\vergleichskette
{ \omega
}
{ =} { -vdu+udv
}
{ } {
}
{ } {
}
{ } {
}
}
{}{}{}
pada
$1$-\definitionsverweis {bola}{}{}

\mavergleichskette
{\vergleichskette
{ S^1
}
{ \subset }{ \R^2
}
{  }{
}
{    }{
}
{   }{
}
}
{}{}{,}
dengan koordinat ruang sekitar $\R^2$ dinotasikan oleh
\mathkor {} {u} {dan} {v} {.}
Tentukan
\mathl{\pi^* \omega}{} di bawah
\definitionsverweis {pemetaan terdiferensialkan secara kontinu}{}{}
\maabbeledisp {\pi} {\R^2 \setminus \{0\}} { S^{1}
} {(x , y) } { { \frac{   1  }{  \sqrt{x^2 + y^2}  } } (x , y)
} {.}

}
{

Berlaku

\mavergleichskettealign
{\vergleichskettealign
{ \pi^*du
}
{ =} { d (u \circ \pi)
}
{ =} { d   \left(  { \frac{   x  }{  \sqrt{x^2 + y^2}  } }  \right)
}
{ =} { { \frac{   \partial   \left(  { \frac{   x  }{  \sqrt{x^2 + y^2}  } }  \right)    }{  \partial   x  } }   dx + { \frac{   \partial   \left(  { \frac{   x  }{  \sqrt{x^2 + y^2}  } }  \right)    }{  \partial   y  } } dy
}
{ =} { { \frac{   \sqrt{x^2+y^2}     -x^2 (x^2 + y^2)^{-1/2}  }{  x^2+y^2  } }  dx   - { \frac{   xy }{  \sqrt{x^2+y^2}^3  } }  dy
}
}
{
\vergleichskettefortsetzungalign
{ =} { { \frac{   x^2+y^2 -x^2    }{  \sqrt{ x^2+y^2}^3  } }  dx  - { \frac{   xy }{  \sqrt{x^2+y^2}^3  } }  dy
}
{ =} { { \frac{   y^2  }{  \sqrt{ x^2+y^2}^3  } }  dx  - { \frac{   xy }{  \sqrt{x^2+y^2}^3  } }  dy
}
{ } {}
{ } {}
}
{}{}
dan

\mavergleichskettealign
{\vergleichskettealign
{ \pi^*dv
}
{ =} { d (v \circ \pi)
}
{ =} { d \left(  { \frac{   y  }{  \sqrt{x^2 + y^2}  } }  \right)
}
{ =} { { \frac{   \partial   \left(  { \frac{   y  }{  \sqrt{x^2 + y^2}  } }  \right)   }{  \partial   x  } } dx + { \frac{   \partial   \left(  { \frac{   y  }{  \sqrt{x^2 + y^2}  } }  \right)    }{  \partial   y  } } dy
}
{ =} { { \frac{   -xy }{  \sqrt{x^2+y^2}^3  } } dx + { \frac{   \sqrt{x^2+y^2} -y^2 (x^2 + y^2)^{-1/2}  }{  x^2+y^2  } } dy
}
}
{
\vergleichskettefortsetzungalign
{ =} { { \frac{   -xy }{  \sqrt{x^2+y^2}^3  } } dx + { \frac{   x^2  }{  \sqrt{x^2+y^2}^3  } }  dy
}
{ } {
}
{ } {}
{ } {}
}
{}{.}
Dengan demikian,

\mavergleichskettealigndrucklinks
{\vergleichskettealigndrucklinks
{ \pi^*(-vdu+udv)
}
{ =} { - (v \circ \pi)  \pi^*du + (u \circ \pi) \pi^*dv
}
{ =} { - { \frac{   y  }{  \sqrt{x^2+y^2}  } } \left(  { \frac{   y^2  }{  \sqrt{ x^2+y^2}^3  } }  dx  - { \frac{   xy }{  \sqrt{x^2+y^2}^3  } }  dy   \right) + { \frac{   x  }{  \sqrt{x^2+y^2}  } }  \left(  { \frac{   -xy }{  \sqrt{x^2+y^2}^3  } } dx + { \frac{   x^2  }{  \sqrt{x^2+y^2}^3  } } dy  \right)
}
{ =} { { \frac{   1 }{  \left(  x^2+y^2  \right)^2 } }  \left(  \left(  -y^3-x^2y  \right) dx + \left(  xy^2+x^3  \right) dy \right)
}
{ =} { { \frac{   1 }{  x^2+y^2  } }  \left(  -  y dx + x dy \right)
}
}
{}{}{.}

 }



\inputaufgabeklausurloesung
{5}

{

Buktikan teorema tentang perhitungan volume kanonik pada manifold Riemann $M$.

}
{

Menurut
Definition  .
kita harus menghitung bentuk diferensial
\mathl{\left(  \alpha^{-1}  \right)^*\omega}{} untuk setiap titik

\mavergleichskette
{\vergleichskette
{ Q
}
{ \in }{ V
}
{  }{
}
{    }{
}
{   }{
}
}
{}{}{}
. Bentuk ini mempunyai bentuk
\mathl{c_Q dx_1 \wedge \ldots \wedge dx_n}{} dan ditentukan oleh nilainya pada
\mathl{e_1 \wedge \ldots \wedge e_n}{}. Berlaku

\mavergleichskettedisphandlinks
{\vergleichskettedisphandlinks
{ \left(  \alpha^{-1}  \right)^*\omega \left(  Q,  e_1 \wedge \ldots \wedge e_n  \right)
}
{ =} { \omega \left(  \alpha^{-1} (Q) ,T_Q \left(  \alpha^{-1}  \right)(e_1) \wedge \ldots \wedge T_Q \left(  \alpha^{-1}  \right)( e_n)  \right)
}
{ } {
}
{ } {
}
{ } {
}
}
{}{}{.}
Menurut definisi matriks fundamental metrik,

\mavergleichskettedisp
{\vergleichskette
{ g_{ij}(Q)
}
{ =} { \left\langle T_Q \left( \alpha^{-1} \right) (e_i) , T_Q \left( \alpha^{-1} \right) (e_j) \right\rangle_{\alpha^{-1}(Q)}
}
{ } {
}
{ } {
}
{ } {
}
}
{}{}{.}
Menurut
Satz 7.8 (Maß- und Integrationstheorie (Osnabrück 2022-2023)),
berlaku

\mavergleichskettealigndrucklinks
{\vergleichskettealigndrucklinks
{ \omega \left(  \alpha^{-1}(Q), T_Q \left(  \alpha^{-1}  \right) (e_1) \wedge \ldots \wedge T_Q \left(  \alpha^{-1}  \right) ( e_n)  \right)
}
{ =} { \left(  \det  \left(  \left\langle T_Q \left(  \alpha^{-1}  \right)  (e_i)  ,  T_Q \left(  \alpha^{-1}  \right) (e_j)   \right\rangle  \right)_{1 \leq i,j \leq n}  \right)^{1/2}
}
{ =} { \left(  \det  \left(  g_{ij}(Q)   \right)_{1 \leq i,j \leq n}  \right)^{1/2}
}
{ =} { \sqrt{g(Q)}
}
{ } {
}
}
{}{}{.}

 }



\inputaufgabeklausurloesung
{6 (1+1+3+1)}

{

Jelaskan mengapa $\R^2$ tidak membawa struktur
\definitionsverweis {manifold berbatas}{}{}
sedemikian sehingga subhimpunan $T$ yang diberikan merupakan
\definitionsverweis {batas}{}{.}
\aufzaehlungvierabc{
\mavergleichskette
{\vergleichskette
{ T
}
{ = }{ {]0,1[}
}
{  }{
}
{    }{
}
{   }{
}
}
{}{}{,}
dengan interval tersebut terletak pada sumbu $x$.
}{
\mavergleichskette
{\vergleichskette
{ T
}
{ = }{ [0,1]
}
{  }{
}
{    }{
}
{   }{
}
}
{}{}{,}
dengan interval tersebut terletak pada sumbu $x$.
}{
\mavergleichskette
{\vergleichskette
{ T
}
{ = }{ \R
}
{  }{
}
{    }{
}
{   }{
}
}
{}{}{,}
dengan $\R$ adalah sumbu $x$.
}{
\mavergleichskette
{\vergleichskette
{ T
}
{ = }{ S^1
}
{  }{
}
{    }{
}
{   }{
}
}
{}{}{.}
}

}
{

a) Batas
\mathl{\partial M}{} dari suatu manifold dengan batas, menurut
Lemma 21.3 (Differentialgeometrie (Osnabrück 2023))  (2)
\definitionsverweis {tertutup}{}{,}
sedangkan interval terbuka sebagai himpunan bagian dari $\R^2$ tidak tertutup.

b) Batas
\mathl{\partial M}{} dari suatu manifold dengan batas, menurut
Satz 21.4 (Differentialgeometrie (Osnabrück 2023)),
merupakan manifold tanpa batas, sedangkan interval tertutup mempunyai titik-titik batas.

c) Andaikan $\R^2$ merupakan suatu manifold dengan batas
\mathl{\R = \R \times \{0\}}{}. Untuk
\mathl{P \in \R}{}, terdapat suatu lingkungan terbuka
\mathl{P\in U}{} dan suatu homeomorfisme
\maabbdisp {\alpha} { U } {V
} {}
dengan suatu himpunan terbuka $V \subseteq H$, dengan $H$ menyatakan suatu setengah bidang. Kita dapat mengganti $V$ dengan lingkungan setengah bola terbuka yang lebih kecil di sekitar $P$. Menurut
Lemma 21.3 (Differentialgeometrie (Osnabrück 2023))  (2)
peta seperti ini memetakan titik batas ke titik batas; artinya,

\mavergleichskettedisp
{\vergleichskette
{\alpha ( U \cap \R )
}
{ = } { V \cap \delta H
}
{ } {
}
{ } {
}
{ } {
}
}
{}{}{.}
Dengan demikian, diperoleh pula suatu homeomorfisme di antara komplemennya, yaitu antara
\mathkor {} {U \setminus U \cap \R} {dan} {V \setminus V \cap \delta H} {.}
Lingkungan setengah bola di sebelah kanan
\definitionsverweis {terhubung lintasan}{}{,}
sedangkan himpunan di sebelah kiri merupakan gabungan saling lepas dari irisannya dengan setengah bagian positif dan negatif. Kedua irisan tersebut terbuka dan tidak kosong karena $U$ memuat suatu lingkungan bola dari $P$. Jadi, himpunan di sebelah kiri tidak terhubung dan homeomorfisme semacam itu tidak mungkin ada.

d) Sama seperti c).

 }



\inputaufgabeklausurloesung
{4}

{

Misalkan $\omega$ suatu
$n-1$-\definitionsverweis {bentuk diferensial}{}{}
yang terdiferensialkan secara kontinu pada
\definitionsverweis {himpunan terbuka}{}{}

\mavergleichskette
{\vergleichskette
{ U
}
{ \subseteq }{ \R^n
}
{  }{
}
{    }{
}
{   }{
}
}
{}{}{}
dan misalkan
\definitionsverweis {turunan eksterior}{}{}
sama dengan

\mavergleichskettedisp
{\vergleichskette
{d \omega
}
{ =} { fdx_1 \wedge \ldots \wedge dx_n
}
{ } {
}
{ } {
}
{ } {
}
}
{}{}{}
dengan suatu fungsi
\maabb {f} { U } {\R
} {.}
Tunjukkan untuk

\mavergleichskette
{\vergleichskette
{ P
}
{ \in }{ U
}
{  }{
}
{    }{
}
{   }{
}
}
{}{}{}
kesamaan

\mavergleichskettedisp
{\vergleichskette
{ f(P)
}
{ =} { { \frac{   1  }{  \beta_n  } } \cdot \operatorname{lim}_{   \epsilon \rightarrow  0  } \, { \frac{   1  }{  \epsilon^n } } \int_{S^{n-1}(P, \epsilon)} \omega
}
{ } {
}
{ } {
}
{ } {
}
}
{}{}{,}
dengan $\beta_n$ menyatakan
volume
bola satuan berdimensi $n$.

}
{

Kita menerapkan
Satz von Stokes
pada bola-bola
\mathl{B^n(P, \epsilon)}{} dengan batas
\mathl{S^{n-1}(P, \epsilon)}{}
\zusatzklammer {dengan pusat $P$ dan jari-jari $\epsilon$} {} {}
dan memperoleh

\mavergleichskettealign
{\vergleichskettealign
{ \operatorname{lim}_{   \epsilon \rightarrow  0  } \, { \frac{   1 }{  \epsilon^n } }     \int_{S^{n-1}(P, \epsilon)} \omega
}
{ =} { \operatorname{lim}_{   \epsilon \rightarrow  0  } \, { \frac{   1 }{  \epsilon^n } }     \int_{ \partial B^{n}(P, \epsilon)}   \omega
}
{ =} { \operatorname{lim}_{   \epsilon \rightarrow  0  } \, { \frac{   1 }{  \epsilon^n } }     \int_{B^{n}(P, \epsilon)} d \omega
}
{ =} { \operatorname{lim}_{   \epsilon \rightarrow  0  } \, { \frac{   1 }{  \epsilon^n } }     \int_{B^{n}(P, \epsilon)} f d \lambda^n
}
{ =} {\beta_n \cdot f(P)
}
}
{}
{}{,}
dengan kesamaan terakhir menggunakan kekontinuan $f$.

 }



\inputaufgabeklausurloesung
{12}

{

Buktikan teorema Stokes untuk manifold berbatas.

}
{

\teilbeweis {}{}{}
{Misalkan

\mathbed {U_i} {}
 {i \in I} {}
 {} {} {} {,}
suatu
\definitionsverweis {penutup terbuka}{}{}
dari $M$ dengan
\definitionsverweis {peta-peta berorientasi}{}{,}
dan misalkan

\mathbed {h_j} {}
 {j \in J} {}
 {} {} {} {,}
suatu
\definitionsverweis {partisi kesatuan yang terdiferensialkan secara kontinu dan tersubordinasi pada penutup tersebut}{}{,}
yang ada menurut
Satz 22.10 (Differentialgeometrie (Osnabrück 2023)). Untuk setiap

\mavergleichskette
{\vergleichskette
{ P
}
{ \in }{ M
}
{  }{
}
{    }{
}
{   }{
}
}
{}{}{}
terdapat suatu lingkungan terbuka

\mavergleichskette
{\vergleichskette
{ P
}
{ \in }{ V_P
}
{ \subseteq }{ M
}
{    }{
}
{   }{
}
}
{}{}{}
sedemikian sehingga

\mavergleichskette
{\vergleichskette
{ h_j \vert_{V_P }
}
{ = }{ 0
}
{ }{
}
{ }{
}
{ }{
}
}
{}{}{}
berlaku kecuali untuk berhingga banyak indeks. Misalkan $Y$ dukungan dari $\omega$. Penutup

\mavergleichskette
{\vergleichskette
{ Y
}
{ \subseteq }{ \bigcup_{P \in M} V_P
}
{  }{
}
{    }{
}
{   }{
}
}
{}{}{}
karena
\definitionsverweis {kekompakan}{}{}
yang diasumsikan, mempunyai suatu subpenutup berhingga, katakanlah

\mavergleichskettedisp
{\vergleichskette
{ Y
}
{ \subseteq} { V_{1} \cup \ldots \cup V_{r}
}
{ =} { V
}
{ } {
}
{ } {
}
}
{}{}{.}
Oleh karena itu, hanya berhingga banyak $h_j$ pada $V$ yang berbeda dari $0$. Kita menetapkan

\mavergleichskette
{\vergleichskette
{ \omega_j
}
{ = }{ h_j \omega
}
{  }{
}
{    }{
}
{   }{
}
}
{}{}{;}
bentuk-bentuk diferensial ini juga terdiferensialkan secara kontinu. Dukungan $h_j$ merupakan
\definitionsverweis {himpunan bagian tertutup}{}{}
dari $M$ yang terletak di dalam
\mathl{U_{i(j)}}{}. Karena itu, dukungan $\omega_j$ terletak di dalam
\mathl{Y \cap U_{i(j)}}{} dan kompak menurut
Aufgabe 13.10 (Differentialgeometrie (Osnabrück 2023)).
Berlaku

\mavergleichskettedisp
{\vergleichskette
{ \omega
}
{ =} { \sum_{j \in J} \omega_j
}
{ } {
}
{ } {
}
{ } {
}
}
{}{}{,}
dan hanya berhingga banyak bentuk diferensial
\mathl{\omega_j}{} yang berbeda dari $0$ karena

\mavergleichskette
{\vergleichskette
{ \omega_j \vert_{M \setminus Y}
}
{ = }{ 0
}
{ }{
}
{ }{
}
{ }{
}
}
{}{}{}
untuk setiap $j$, serta

\mavergleichskettedisp
{\vergleichskette
{ \omega_j \vert_V
}
{ =} { 0
}
{ } {
}
{ } {
}
{ } {
}
}
{}{}{}
untuk setiap $j$ kecuali berhingga banyak indeks. Berdasarkan
Additivität des Integrals von Differentialformen
dan
Additivität der äußeren Ableitung
pernyataan dapat dibuktikan secara terpisah untuk setiap $\omega_j$. Jadi, kita dapat menganggap bahwa diberikan suatu bentuk diferensial-$(n-1)$ yang terdiferensialkan secara kontinu dan mempunyai dukungan kompak yang seluruhnya terletak dalam suatu lingkungan peta $U$.}
{}
\teilbeweis {}{}{}
{Terdapat suatu
\definitionsverweis {difeomorfisme}{}{}
\maabb {\alpha} { U } { V
} {}
dengan

\mavergleichskette
{\vergleichskette
{ V
}
{ \subseteq }{ H
}
{ }{
}
{ }{
}
{ }{
}
}
{}{}{}
terbuka, yang sekaligus menginduksi suatu difeomorfisme di antara batas-batas

\mavergleichskette
{\vergleichskette
{ \partial U
}
{ = }{ \partial M \cap U
}
{  }{
}
{    }{
}
{   }{
}
}
{}{}{}
dan

\mavergleichskette
{\vergleichskette
{ \partial V
}
{ = }{ \partial H \cap V
}
{  }{
}
{    }{
}
{   }{
}
}
{}{}{}
. Berlaku

\mavergleichskettedisp
{\vergleichskette
{
\int_{ \partial U } \omega
}
{ = } {
\int_{ \partial V } \alpha^{-1 *}\omega
}
{ } {
}
{ } {
}
{ } {
}
}
{}{}{}
dan

\mavergleichskettedisp
{\vergleichskette
{
\int_{ U  }  d \omega
}
{ =} {
\int_{ V  }  \alpha^{-1 *} (  d \omega)
}
{ =} {
\int_{  V  }  d \left(  \alpha^{-1 *}\omega  \right)
}
{ } {
}
{ } {
}
}
{}{}{}
menurut
Lemma 20.2 (Differentialgeometrie (Osnabrück 2023))  (5).}
{\leerzeichen{}Jadi, kita dapat mulai dengan suatu bentuk diferensial yang terdiferensialkan secara kontinu, mempunyai dukungan kompak, dan didefinisikan pada

\mavergleichskette
{\vergleichskette
{ V
}
{ \subseteq }{ H
}
{ }{
}
{ }{
}
{ }{
}
}
{}{}{}
serta memperluasnya menjadi bentuk nol pada seluruh $H$ di luar dukungannya.}
\teilbeweis {}{}{}
{Karena kekompakan dukungan, terdapat suatu balok yang cukup besar

\mavergleichskette
{\vergleichskette
{ Q
}
{ \subseteq }{ H
}
{ }{
}
{ }{
}
{ }{
}
}
{}{}{,}
yang salah satu sisinya $S$ terletak pada $\partial H$ dan yang berpotongan dengan dukungan $\omega$ hanya di $S$. Pada semua sisi lain dari $Q$, bentuk $\omega$
\zusatzklammer {dan dengan demikian juga $d\omega$} {} {}
merupakan bentuk nol. Oleh karena itu, di satu pihak

\mavergleichskettedisp
{\vergleichskette
{
\int_{ H } d\omega
}
{ =} {
\int_{ Q } d\omega
}
{ } {
}
{ } {
}
{ } {
}
}
{}{}{}
dan di pihak lain

\mavergleichskettedisp
{\vergleichskette
{
\int_{ \partial H } \omega
}
{ = } {
\int_{ S } \omega
}
{ = } {
\int_{ \partial Q } \omega
}
{ } {
}
{ } {
}
}
{}{}{.}
Dengan demikian, pernyataan tersebut mengikuti dari Quaderversion des Satzes von Stokes.}
{}

 }

\endgroup

\chapter{Formulir Solusi Resmi Ujian 3}
\noindent\textit{Solusi yang disediakan oleh sumber resmi; kekosongan sumber dipertahankan.}
\begingroup
\renewcommand{\theHfakt}{o011.exam.official.03.\arabic{fakt}}
%Data institusi

%\input{Dozentdaten}

%\renewcommand{\fachbereich}{Fachbereich}

%\renewcommand{\dozent}{Prof. Dr. . }

%Data ujian

\renewcommand{\klausurgebiet}{ }

\renewcommand{\klausurtyp}{ }

\renewcommand{\klausurdatum}{ . 20}

\klausurvorspann {\fachbereich} {\klausurdatum} {\dozent} {\klausurgebiet} {\klausurtyp}

%Data untuk tabel nilai berikut

\renewcommand{\aeins}{  3  }

\renewcommand{\azwei}{  3   }

\renewcommand{\adrei}{  0  }

\renewcommand{\avier}{  0  }

\renewcommand{\afuenf}{  6  }

\renewcommand{\asechs}{  6  }

\renewcommand{\asieben}{  0  }

\renewcommand{\aacht}{  6  }

\renewcommand{\aneun}{  11  }

\renewcommand{\azehn}{  5  }

\renewcommand{\aelf}{  5  }

\renewcommand{\azwoelf}{  3  }

\renewcommand{\adreizehn}{  7   }

\renewcommand{\avierzehn}{  2  }

\renewcommand{\afuenfzehn}{  0  }

\renewcommand{\asechzehn}{  6  }

\renewcommand{\asiebzehn}{  63  }

\renewcommand{\aachtzehn}{    }

\renewcommand{\aneunzehn}{    }

\renewcommand{\azwanzig}{    }

\renewcommand{\aeinundzwanzig}{    }

\renewcommand{\azweiundzwanzig}{    }

\renewcommand{\adreiundzwanzig}{    }

\renewcommand{\avierundzwanzig}{    }

\renewcommand{\afuenfundzwanzig}{    }

\renewcommand{\asechsundzwanzig}{    }

\punktetabellesechzehn

\klausurnote

\newpage

\setcounter{section}{K}



\inputaufgabeklausurloesung
{3}

{

Definisikan istilah-istilah berikut
\zusatzklammer {yang dicetak miring} {} {.}
\aufzaehlungsechs{Suatu
\stichwort {kurva geodesik} {}
pada hiperpermukaan diferensiabel

\mavergleichskette
{\vergleichskette
{ Y
}
{ \subseteq }{ \R^n
}
{  }{
}
{    }{
}
{   }{
}
}
{}{}{.}

}{
\stichwort {Himpunan rotasi} {}
\zusatzklammer {mengelilingi sumbu $x$} {} {}
dari suatu subhimpunan

\mavergleichskette
{\vergleichskette
{ T
}
{ \subseteq }{  \R \times \R_{\geq 0}
}
{  }{
}
{    }{
}
{   }{
}
}
{}{}{.}

}{
\stichwort {Pemetaan tangen} {}
\maabbdisp {T (\varphi)} {TM} {TN
} {}
dari
\definitionsverweis {pemetaan diferensiabel}{}{}
\maabbdisp {\varphi} { M } { N
} {}
antara
\definitionsverweis {manifold-manifold diferensiabel}{}{}
\mathkor {} {M} {dan} {N} {.}

}{Suatu perubahan peta yang
\stichwort {mempertahankan orientasi} {}
pada
\definitionsverweis {manifold diferensiabel}{}{}
$M$.

}{
\stichwort {Bentuk diferensial tarikan balik} {} $\varphi^* \omega$ dari suatu bentuk diferensial
\mathl{\omega \in
{ \mathcal E }^{ k } (  M   )}{} terhadap suatu pemetaan terdiferensialkan secara kontinu
\maabb {\varphi} { L } { M
} {}
antara dua manifold diferensiabel
\mathkor {} {L} {dan} {M} {.}

}{Suatu koneksi linear yang
\stichwort {metrik} {}
pada bundel tangen suatu manifold Riemann $M$.
}

}
{

\aufzaehlungsechs{Suatu
\zusatzklammer {sebagai pemetaan ke $\R^n$} {} {}
\definitionsverweis {kurva yang terdiferensialkan secara kontinu dua kali}{}{}
\maabbdisp {\gamma} {I} {Y
} {}
disebut
kurva geodesik
pada $Y$ jika
\definitionsverweis {percepatan tangensial}{}{}
kurva tersebut sama dengan $0$ di setiap titik.
}{Himpunan rotasi dari $T$ adalah

\mathdisp {\left\{ (x, y \cos \alpha, y \sin \alpha  ) \in \R^3  \mid   (x,y) \in T , \, \alpha \in [0, 2 \pi]       \right\}} { . }

}{Yang dimaksud dengan pemetaan tangen
\maabbdisp {T (\varphi)} {TM} {TN
} {}
adalah gabungan saling lepas dari
\definitionsverweis {pemetaan-pemetaan tangen}{}{}
di setiap titik, yakni

\mavergleichskettedisp
{\vergleichskette
{ T (\varphi)
}
{ =} { \biguplus_{P \in M} T_P(\varphi)
}
{ } {
}
{ } {
}
{ } {
}
}
{}{}{.}
}{Misalkan
\maabbdisp {\alpha_1} {U_1 } { V_1
} {}
dan
\maabbdisp {\alpha_2} {U_2 } { V_2
} {}
adalah peta-peta berorientasi dari $M$. Perubahan peta terkait
\definitionsverweis {perubahan peta}{}{}
\maabbdisp {\psi=\alpha_2 \circ \alpha_1^{-1}} {V_1 \cap \alpha_1 (U_1 \cap U_2)} {V_2 \cap \alpha_2 (U_1 \cap U_2 )
} {}
disebut mempertahankan orientasi jika untuk setiap titik
\mathl{Q \in V_1 \cap \alpha_1 (U_1 \cap U_2)}{},
\definitionsverweis {diferensial total}{}{}
\maabbdisp {\left(D\psi\right)_{Q}} {\R^n } {\R^n
} {}
\definitionsverweis {mempertahankan orientasi}{}{}.
}{Bentuk diferensial tarikan balik
\mathl{\varphi^* \omega}{} didefinisikan, untuk
\mathl{P\in L}{} dan
\mathl{v_1 , \ldots , v_k \in T_PL}{}, oleh

\mavergleichskettedisp
{\vergleichskette
{  ( \varphi^* \omega )(P,v_1 , \ldots , v_k )
}
{ \defeq} { \omega(\varphi(P), T_P\varphi(v_1) , \ldots ,  T_P\varphi(v_k) )
}
{ } {
}
{ } {
}
{ } {
}
}
{}{}{}
.
}{Koneksi tersebut disebut
metrik
jika

\mavergleichskettedisp
{\vergleichskette
{ D_V \left\langle W , Z  \right\rangle
}
{ =} { \left\langle \nabla_V W , Z  \right\rangle  + \left\langle W ,  \nabla_VZ  \right\rangle
}
{ } {
}
{ } {
}
{ } {
}
}
{}{}{.}
berlaku untuk sebarang medan vektor $V,W ,Z$.
}

 }



\inputaufgabeklausurloesung
{3}

{

Nyatakan teorema-teorema berikut.
\aufzaehlungdrei{
\stichwort {Teorema tentang solusi persamaan diferensial biasa pada hiperpermukaan diferensiabel} {}

\mavergleichskette
{\vergleichskette
{  Y
}
{ \subseteq }{  \R^n
}
{  }{
}
{    }{
}
{   }{
}
}
{}{}{.}}{Rumus untuk menghitung volume kanonik suatu manifold Riemann dalam sebuah peta.}{Teorema tentang partisi kesatuan.}

}
{

\aufzaehlungdrei{\faktsituation {Misalkan

\mavergleichskette
{\vergleichskette
{ W
}
{ \subseteq }{ \R^n
}
{  }{
}
{    }{
}
{   }{
}
}
{}{}{}
\definitionsverweis {terbuka}{}{,}
\maabb {h} { W } { \R
} {}
suatu
\definitionsverweis {fungsi yang terdiferensialkan secara kontinu}{}{}
dan

\mavergleichskette
{\vergleichskette
{ Y
}
{ = }{ h^{-1}(c)
}
{  }{
}
{    }{
}
{   }{
}
}
{}{}{}
\definitionsverweis {serat}{}{}
atas

\mavergleichskette
{\vergleichskette
{ c
}
{ \in }{ \R
}
{  }{
}
{    }{
}
{   }{
}
}
{}{}{,}
dengan $h$ bersifat
\definitionsverweis {reguler}{}{}
di setiap titik $Y$.
Misalkan

\mavergleichskette
{\vergleichskette
{ Y
}
{ \subseteq }{ U
}
{ \subseteq }{ \R^n
}
{    }{
}
{   }{
}
}
{}{}{}
\definitionsverweis {terbuka}{}{}
dan misalkan $F$ suatu
\definitionsverweis {medan vektor}{}{}
terdiferensialkan pada $U$}
\faktvoraussetzung {dengan sifat bahwa untuk setiap titik

\mavergleichskette
{\vergleichskette
{ P
}
{ \in }{ Y
}
{  }{
}
{    }{
}
{   }{
}
}
{}{}{}
vektor $F(P)$ termasuk dalam
\definitionsverweis {ruang tangen serat}{}{}
di $P$ pada $Y$.}
\faktfolgerung {Maka solusi $v(t)$ dari
\definitionsverweis {masalah nilai awal}{}{}
untuk $F$ dengan syarat awal

\mavergleichskette
{\vergleichskette
{ v(t_0)
}
{ \in }{ Y
}
{  }{
}
{    }{
}
{   }{
}
}
{}{}{}
seluruhnya terletak di $Y$.}
\faktzusatz {}
\faktzusatz {}}{Misalkan $M$ suatu
\definitionsverweis {manifold Riemann}{}{}
\definitionsverweis {berorientasi}{}{}
dan $\omega$ adalah
\definitionsverweis {bentuk volume kanonik}{}{.}
Misalkan
\maabbdisp {\alpha} { U } { V
} {}
suatu
\definitionsverweis {peta berorientasi}{}{}
dengan

\mavergleichskettedisp
{\vergleichskette
{ V
}
{ \subseteq} { \R^n
}
{ } {
}
{ } {
}
{ } {
}
}
{}{}{}
terbuka, berkoordinat
\mathl{x_1 , \ldots , x_n}{}, dengan
\definitionsverweis {matriks fundamental metrik}{}{}

\mathl{G=(g_{ij})_{1 \leq i,j \leq n}}{} dan
\mathl{g= \det  G}{.} Maka

\mavergleichskettedisp
{\vergleichskette
{  \alpha_* \omega
}
{ =} { \sqrt{ g} dx_1 \wedge \ldots \wedge dx_n
}
{ } {
}
{ } {
}
{ } {
}
}
{}{}{.}
Untuk suatu
\definitionsverweis {subhimpunan terukur}{}{}

\mavergleichskette
{\vergleichskette
{ T
}
{ \subseteq }{ U
}
{  }{
}
{    }{
}
{   }{
}
}
{}{}{}
berlaku

\mavergleichskettedisp
{\vergleichskette
{
\int_{  T   }  \omega
}
{ =} {
\int_{ \alpha(T)  }   \sqrt{g} dx_1 \wedge \ldots \wedge dx_n
}
{ =} { \int_{ \alpha(T)  }   \sqrt{g}    \,  d \lambda^n
}
{ } {
}
{ } {
}
}
{}{}{.}}{Misalkan $M$ suatu
\definitionsverweis {manifold diferensiabel}{}{}
dengan
\definitionsverweis {basis terhitung}{}{}
bagi topologinya. Maka untuk setiap penutup terbuka terdapat
\definitionsverweis {partisi kesatuan}{}{}
yang
\definitionsverweis {terdiferensialkan secara kontinu}{}{}
dan subordinat terhadap penutup tersebut.}

 }



\inputaufgabeklausurloesung
{0}

{

}
{  }



\inputaufgabeklausurloesung
{0}

{

}
{  }



\inputaufgabeklausurloesung
{6}

{

Untuk permukaan $Y$ yang diberikan oleh

\mavergleichskettedisp
{\vergleichskette
{ x^2+y^2+2z^2
}
{ =} { 4
}
{ } {
}
{ } {
}
{ } {
}
}
{}{}{}
dan titik

\mavergleichskettedisp
{\vergleichskette
{ P
}
{ =} { \left(  1  , \,   1  , \,   1                 \right)
}
{ \in} { Y
}
{ } {
}
{ } {
}
}
{}{}{,}
tentukan suatu matriks diagonal bagi
\definitionsverweis {pemetaan Weingarten}{}{}
$L_P$.

}
{

Medan gradien ternormalisasi adalah

\mavergleichskettealign
{\vergleichskettealign
{ N(x,y,z)
}
{ =} { { \frac{   1  }{  \sqrt{ 4x^2+4y^2+16 z^2}  } } \begin{pmatrix}  2x \\ 2y\\  4z   \end{pmatrix}
}
{ =} { { \frac{   1  }{  \sqrt{ x^2+y^2+4 z^2}  } } \begin{pmatrix}  x  \\ y \\  2z   \end{pmatrix}
}
{ } {
}
{ } {
}
}
{}
{}{.}
Kita bekerja dengan medan normal satuan

\mavergleichskettedisp
{\vergleichskette
{ M(x,y,z)
}
{ =} { { \frac{   1  }{  \sqrt{ 4+2 z^2}  } } \begin{pmatrix}  x  \\ y \\  2z   \end{pmatrix}
}
{ } {
}
{ } {
}
{ } {
}
}
{}{}{,}
yang pada $Y$ berimpit dengan $N$. Diferensial total dari $M$ adalah

\mathdisp {\begin{pmatrix}  { \frac{   1  }{  \sqrt{4+2z^2}  } }   &   0 &  { \frac{   -2 xz }{  \left(  4+2z^2  \right)^{3/2}  } }  \\  0  &  { \frac{   1  }{  \sqrt{4+2z^2}  } }   &  { \frac{   -2 yz }{  \left(  4+2z^2  \right)^{3/2}  } }   \\ 0   &  0  &  { \frac{   8 }{  \left(  4+2z^2  \right)^{3/2}  } }  \end{pmatrix}} { . }

Di titik $P$ yang diberikan, gradiennya adalah $\begin{pmatrix}  2  \\ 2\\  4   \end{pmatrix}$ dan kedua vektor
\mathkor {} {\begin{pmatrix}  1  \\ -1\\  0   \end{pmatrix}} {dan} {\begin{pmatrix}  2  \\ 0 \\  -1   \end{pmatrix}} {}
membentuk suatu basis ruang tangen $T_PY$. Diferensial total dari $M$ di titik ini sama dengan

\mathdisp {\begin{pmatrix}  { \frac{   1  }{  \sqrt{6}  } }   &   0 &  { \frac{   -1 }{  3 \sqrt{6}  } }  \\  0  &  { \frac{   1  }{  \sqrt{6}  } }   &  { \frac{   -1  }{  3 \sqrt{6}  } }   \\ 0   &  0  &  { \frac{   4  }{  3 \sqrt{6}  } }   \end{pmatrix}} { . }

Jika diterapkan pada vektor basis pertama $\begin{pmatrix}  1  \\ -1\\  0   \end{pmatrix}$, hasilnya adalah $\begin{pmatrix}  { \frac{   1  }{  \sqrt{6}  } }  \\ - { \frac{   1  }{  \sqrt{6}  } } \\  0   \end{pmatrix}$, jadi ini merupakan vektor eigen dengan nilai eigen ${ \frac{   1  }{  \sqrt{6}  } }$. Jika diterapkan pada vektor basis kedua $\begin{pmatrix}  2  \\ 0 \\  -1   \end{pmatrix}$, hasilnya adalah

\mavergleichskettedisp
{\vergleichskette
{ \begin{pmatrix}  { \frac{   7 }{  3 \sqrt{6}  } }  \\ { \frac{   1 }{  3 \sqrt{6 } }}  \\  - { \frac{   4  }{  3 \sqrt{6}  } }   \end{pmatrix}
}
{ =} { - { \frac{   1 }{  3 \sqrt{6 } }} \begin{pmatrix}  1  \\ -1\\  0   \end{pmatrix} +   { \frac{   4  }{  3 \sqrt{6}  } } \begin{pmatrix}  2  \\ 0 \\  -1   \end{pmatrix}
}
{ } {
}
{ } {
}
{ } {
}
}
{}{}{.}
Karena itu, nilai eigen yang lain sama dengan ${ \frac{   4  }{  3 \sqrt{6}  } }$, dan matriks diagonal yang merepresentasikannya adalah

\mathdisp {\begin{pmatrix}  { \frac{   1  }{  \sqrt{6}  } }   &   0  \\  0  &  { \frac{   4  }{  3 \sqrt{6}  } }  \end{pmatrix}} { . }

 }



\inputaufgabeklausurloesung
{6}

{

Buktikan teorema isometri untuk transpor paralel pada hiperpermukaan

\mavergleichskette
{\vergleichskette
{ Y
}
{ \subseteq }{ \R^n
}
{  }{
}
{    }{
}
{   }{
}
}
{}{}{.}

}
{

Untuk membuktikan linearitas, misalkan diberikan

\mavergleichskette
{\vergleichskette
{ v,w
}
{ \in }{ T_PY
}
{  }{
}
{    }{
}
{   }{
}
}
{}{}{}
dan

\mavergleichskette
{\vergleichskette
{ r,s
}
{ \in }{ \R
}
{  }{
}
{    }{
}
{   }{
}
}
{}{}{}.
Misalkan
\mathkor {} {F} {dan} {G} {}
adalah medan vektor paralel sepanjang $\gamma$ yang ditentukan secara unik menurut
Teorema 6.9 (Geometri diferensial (Osnabrück 2023))
dengan

\mavergleichskette
{\vergleichskette
{ F(a)
}
{ = }{ v
}
{  }{
}
{    }{
}
{   }{
}
}
{}{}{}
dan

\mavergleichskette
{\vergleichskette
{ G(a)
}
{ = }{ w
}
{  }{
}
{    }{
}
{   }{
}
}
{}{}{.}
Menurut
Lema 6.7 (Geometri diferensial (Osnabrück 2023)),
\mathl{rF+sG}{} adalah medan vektor paralel dengan

\mavergleichskette
{\vergleichskette
{ (rF+sG)(a)
}
{ = }{ v+ w
}
{  }{
}
{    }{
}
{   }{
}
}
{}{}{.}
Karena keunikan dalam
Teorema 6.9 (Geometri diferensial (Osnabrück 2023)),
\mathl{rF+sG}{} dengan demikian adalah medan vektor paralel untuk vektor tangen
\mathl{rv+sw}{.} Oleh karena itu,

\mavergleichskettedisp
{\vergleichskette
{ \Psi_{\gamma} (rv+sw)
}
{ =} { (rF+sG)(b)
}
{ =} { r F(b) +s G(b)
}
{ =} { r \Psi_{\gamma} (v) + s \Psi_{\gamma} (w)
}
{ } {
}
}
{}{}{.}

Untuk membuktikan kesesuaian dengan hasil kali skalar, misalkan kembali diberikan

\mavergleichskette
{\vergleichskette
{ v,w
}
{ \in }{ T_PY
}
{  }{
}
{    }{
}
{   }{
}
}
{}{}{}
dan misalkan $F,G$ adalah medan-medan vektor paralel yang terkait. Kita peroleh

\mavergleichskettedisp
{\vergleichskette
{ \left\langle  F(t) , G(t)  \right\rangle'
}
{ =} { \left\langle  F'(t) , G(t)  \right\rangle  +  \left\langle  F(t) , G'(t)  \right\rangle
}
{ =} { 0
}
{ } {
}
{ } {
}
}
{}{}{,}
karena $F,G$ tangensial dan $F',G'$ ortogonal terhadap ruang tangen. Jadi
\mathl{\left\langle  F(t) , G(t)  \right\rangle}{} konstan sepanjang lintasan. Oleh karena itu,

\mavergleichskettedisp
{\vergleichskette
{ \left\langle \Psi_\gamma(v)  , \Psi_\gamma(w)  \right\rangle
}
{ =} { \left\langle  F(b)  ,  G(b)   \right\rangle
}
{ =} { \left\langle  F(a)  ,  F(a)   \right\rangle
}
{ =} { \left\langle  v  , w  \right\rangle
}
{ } {
}
}
{}{}{.}
Bijektivitasnya juga jelas.

 }



\inputaufgabeklausurloesung
{0}

{

}
{  }



\inputaufgabeklausurloesung
{6}

{

Misalkan $M$ suatu manifold topologis kompak berdimensi $d$,

\mavergleichskette
{\vergleichskette
{ d
}
{ \geq }{ 1
}
{  }{
}
{    }{
}
{   }{
}
}
{}{}{.}
Tunjukkan bahwa terdapat suatu subhimpunan terbuka terbatas

\mavergleichskette
{\vergleichskette
{ U
}
{ \subseteq }{ \R^d
}
{  }{
}
{    }{
}
{   }{
}
}
{}{}{}
dan suatu pemetaan kontinu surjektif
\maabbdisp {\varphi} { U } { M
} {}
.

}
{

Untuk setiap titik
\mathl{P \in M}{}, pilih suatu lingkungan peta terbuka
\mathl{P \in U_P}{} dengan peta
\maabbdisp {\alpha_P} {U_P } {V_P
} {}
dengan
\mathl{V_P \subseteq \R^d}{.} Dengan beralih ke suatu lingkungan bola dari
\mathl{\alpha_P(P) \in V_P}{} beserta prapetanya, kita dapat mengasumsikan bahwa $V_p$ adalah bola-bola terbuka yang jari-jarinya paling besar
\mathl{{ \frac{   1 }{  3 } }}{}. Himpunan-himpunan

\mathbed {U_P} {}
 {P\in M} {}
 {} {} {} {,}
menutupi manifold tersebut. Karena kekompakan $M$, terdapat berhingga banyak titik
\mathl{P_1 , \ldots , P_n}{} sedemikian sehingga

\mathbed {U_i=U_{P_i}} {}
 {i=1 , \ldots , n} {}
 {} {} {} {,}
juga menutupi manifold tersebut. Kita tempatkan bola-bola terbuka
\mathl{V_i=V_{P_i}}{} dalam
\zusatzklammer {\anfuehrung{$\R^d$ yang baru}{}} {} {},
dengan titik-titik pusat

\mathdisp {(1,0 , \ldots , 0), (2,0 , \ldots , 0) , \ldots , (n,0 , \ldots , 0)} { . }

Karena syarat jari-jari, bola-bola ini saling lepas. Kita tinjau himpunan terbuka terbatas
\mathl{U= \bigcup_{i=1}^n V_i}{.} Pemetaan-pemetaan peta $\alpha_i$ menghasilkan pemetaan kontinu

\mathdisp {\varphi_i :V_i \stackrel{ \left(  \alpha_i  \right)^{-1} }{\longrightarrow} U_i \longrightarrow M} { . }

Karena saling lepas, pemetaan-pemetaan ini mendefinisikan, melalui
\mathl{\varphi(z)= \varphi_i(z)}{,} apabila
\mathl{z \in V_i}{}, suatu pemetaan kontinu
\maabbdisp {\varphi} { U } {M
} {.}
Karena sifat penutupnya, pemetaan ini surjektif.

 }



\inputaufgabeklausurloesung
{11 (1+3+1+2+4)}

{

Kita tinjau himpunan

\mavergleichskettedisp
{\vergleichskette
{N
}
{ =} { \left\{ A= \begin{pmatrix}  a   &   b   \\  c  & d \end{pmatrix}   \mid   A \text{ nilpoten}         \right\}
}
{ \subseteq} {\R^4
}
{ } {
}
{ } {
}
}
{}{}{}
dari
$2\times 2$-\definitionsverweis {matriks}{}{}
real nilpoten, serta himpunan

\mavergleichskettedisp
{\vergleichskette
{M
}
{ =} { N \setminus \{ \begin{pmatrix}  0   &   0  \\  0  &  0 \end{pmatrix}  \}
}
{ } {
}
{ } {
}
{ } {
}
}
{}{}{.}

a) Apakah $N$ terhubung?

b) Tunjukkan bahwa $M$ suatu
\definitionsverweis {submanifold tertutup}{}{}
dari suatu subhimpunan terbuka

\mavergleichskette
{\vergleichskette
{ G
}
{ \subseteq }{ \R^4
}
{  }{
}
{    }{
}
{   }{
}
}
{}{}{.}

c) Tentukan
\definitionsverweis {dimensi}{}{}
$M$.

d) Apakah $M$ terhubung?

e) Tutupi
\mathl{M}{} dengan peta-peta topologis eksplisit.

}
{

a) Setiap matriks nilpoten
\mathl{\begin{pmatrix}  a   &   b   \\  c  & d \end{pmatrix}}{} dapat dihubungkan dengan matriks nol di dalam himpunan matriks nilpoten melalui lintasan linear
\maabbeledisp {} {[0,1]} {N
} { t } {t \begin{pmatrix}  a   &   b   \\  c  & d \end{pmatrix}
} {,}.
Karena itu, $N$ terhubung-lintasan dan dengan demikian juga terhubung.

b) Suatu
$2 \times 2$-\definitionsverweis {matriks}{}{}
bersifat nilpoten tepat ketika jejak dan determinannya sama dengan $0$. Jadi, himpunan $N$ dari matriks-matriks nilpoten dapat dipandang sebagai

\mathdisp {\left\{  (a,b,c,d)  \mid   a+d = 0 \text{ dan } ad-bc = 0         \right\}} {  }

Kita tinjau pemetaan
\maabbeledisp {} {\R^4} {\R^2
} {(a,b,c,d)} { (a+d, ad-bc)
} {.}
Matriks Jacobi pemetaan ini adalah

\mathdisp {\begin{pmatrix}  1 &  0 &  0 &  1 \\  d  &  -c &  -b &  a  \end{pmatrix}} { . }

Pemetaan ini tidak
\definitionsverweis {reguler}{}{,}
di titik nol, tetapi reguler di setiap titik lain pada serat $N$. Memang, jika
\mathl{P \in M}{}, maka dari

\mavergleichskettedisp
{\vergleichskette
{a^2
}
{ =} {bc
}
{ } {
}
{ } {
}
{ } {
}
}
{}{}{}
dan

\mavergleichskettedisp
{\vergleichskette
{b
}
{ =} { c
}
{ =} { 0
}
{ } {
}
{ } {
}
}
{}{}{}
langsung diperoleh

\mavergleichskettedisp
{\vergleichskette
{a
}
{ =} {b
}
{ =} { c
}
{ =} { d
}
{ =} { 0
}
}
{}{}{.}
Jadi, matriks Jacobi memiliki rank maksimal di titik-titik dalam $M$. Dengan

\mavergleichskettedisp
{\vergleichskette
{G
}
{ =} {\R^4 \setminus \{0\}
}
{ } {
}
{ } {
}
{ } {
}
}
{}{}{}
kita dapat memandang $M$ sebagai serat pemetaan terbatas yang reguler di seluruh serat tersebut. Oleh karena itu, menurut
Teorema 8.2 (Geometri diferensial (Osnabrück 2023)),
$M$ adalah submanifold tertutup dari $G$.

c) Menurut
Teorema 8.2 (Geometri diferensial (Osnabrück 2023)),
dimensi $M$ sama dengan

\mavergleichskette
{\vergleichskette
{ 4-2
}
{ = }{ 2
}
{  }{
}
{    }{
}
{   }{
}
}
{}{}{.}

d) Kita tulis

\mavergleichskettedisp
{\vergleichskette
{M_+
}
{ =} { \left\{ (a,b,c,d) \in M  \mid   b>c        \right\}
}
{ } {
}
{ } {
}
{ } {
}
}
{}{}{}
dan

\mavergleichskettedisp
{\vergleichskette
{M_-
}
{ =} { \left\{ (a,b,c,d) \in M  \mid   b<c        \right\}
}
{ } {
}
{ } {
}
{ } {
}
}
{}{}{,}
Keduanya,
\zusatzklammer {sebagai irisan $M$ dengan himpunan yang diberikan oleh
\mathl{b>c}{} dan terbuka di $\R^4$} {} {},
merupakan himpunan terbuka dalam $M$. Matriks-matriks
\mathkor {} {\begin{pmatrix}  0  &   1  \\  0 &  0 \end{pmatrix}} {dan} {\begin{pmatrix}  0  &   0  \\  1 &  0 \end{pmatrix}} {}
menunjukkan bahwa keduanya tidak kosong. Selain itu, keduanya menutupi seluruh $M$. Memang, jika

\mavergleichskette
{\vergleichskette
{b
}
{ = }{ c
}
{  }{
}
{    }{
}
{   }{
}
}
{}{}{}
maka dari

\mavergleichskettedisp
{\vergleichskette
{ 0
}
{ =} { ad-bc
}
{ =} { -a^2 -b^2
}
{ } {
}
{ } {
}
}
{}{}{}
langsung diperoleh

\mavergleichskette
{\vergleichskette
{a
}
{ = }{b
}
{ = }{ c
}
{ =   }{ d
}
{ =  }{ 0
}
}
{}{}{,}
dan titik tersebut tidak termasuk dalam $M$. Jadi terdapat penutup oleh dua himpunan terbuka tak kosong yang saling lepas; karena itu $M$ tidak terhubung.

e) Kita bekerja dengan pemetaan
\maabbeledisp {\psi} {\R^2 \setminus \{(0,0)\}} { M_+
} {(a,u)} { (a, u + \sqrt{a^2+u^2}, u - \sqrt{a^2 +u^2},-a)
} {.}
Karena

\mavergleichskettealign
{\vergleichskettealign
{ad-bc
}
{ =} { -a^2 - \left(  u + \sqrt{a^2+u^2}  \right) \left( u - \sqrt{a^2+u^2}  \right)
}
{ =} { -a^2 - \left(  u^2 - \left(  a^2+u^2 \right)   \right)
}
{ =} { 0
}
{ } {
}
}
{}
{}{}
syarat determinan dipenuhi, dan karena

\mavergleichskettedisp
{\vergleichskette
{b-c
}
{ =} { u + \sqrt{a^2+u^2} - \left(  u- \sqrt{a^2 +u^2})  \right)
}
{ =} { 2 \sqrt{a^2+u^2}
}
{ >} { 0
}
{ } {
}
}
{}{}{}
citranya termasuk dalam $M_+$. Pemetaan
\maabbeledisp {\varphi} {M_+} { \R^2 \setminus \{(0,0)\}
} {(a,b,c,d)} { \left(  a  , \,   { \frac{  b+c }{  2 } }                   \right)
} {,}
merupakan invers dari pemetaan yang diberikan. Di sini,

\mavergleichskettedisp
{\vergleichskette
{ \varphi \circ \psi
}
{ =} {
\operatorname{Id}_{ \R^2 \setminus \{(0,0)\}  }
}
{ } {
}
{ } {
}
{ } {
}
}
{}{}{}
jelas. Identitas yang lain diperoleh dari

\mavergleichskettealign
{\vergleichskettealign
{ { \frac{  b+c }{  2 } } + \sqrt{a^2 + \left( { \frac{  b+c }{  2 } }  \right)^2 }
}
{ =} { { \frac{  b+c }{  2 } } + { \frac{   1 }{  2 } }  \sqrt{4 a^2 + b^2 +c^2 +2bc }
}
{ =} { { \frac{  b+c }{  2 } } + { \frac{   1 }{  2 } }  \sqrt{ b^2 +c^2 -2bc }
}
{ =} { { \frac{  b+c }{  2 } } + { \frac{  b-c }{  2 } }
}
{ =} {b
}
}
{}
{}{}
dan

\mavergleichskettedisp
{\vergleichskette
{ { \frac{  b+c }{  2 } } - \sqrt{a^2 + \left(   { \frac{  b+c }{  2 } }   \right)^2 }
}
{ =} { { \frac{  b+c }{  2 } } - { \frac{  b-c }{  2 } }
}
{ =} { c
}
{ } {}
{ } {}
}
{}{}{.}
Kedua pemetaan kontinu, sehingga diperoleh suatu
\definitionsverweis {homeomorfisme}{}{}.
Untuk
\mathl{M_-}{}, tukarkan peran
\mathkor {} {b} {dan} {c} {.}

 }



\inputaufgabeklausurloesung
{5}

{

Katakan sesuatu yang cerdas tentang pita Möbius.

}
{  }



\inputaufgabeklausurloesung
{5}

{

Kita tinjau pemetaan
\maabbeledisp {\varphi} { \R^3 } { \R^2
} { (x,y,z) } { \left(  xyz^2,xy-z^3  \right)
} {.}
Hitung matriks pemetaan
\maabbdisp {\bigwedge^2 T_P (\varphi)} { \bigwedge^2 T_P \R^3 } { \bigwedge^2 T_{\varphi(P)}  \R^2
} {}
di titik

\mavergleichskette
{\vergleichskette
{ P
}
{ = }{ (1,3,5)
}
{  }{
}
{    }{
}
{   }{
}
}
{}{}{}
terhadap suatu basis yang sesuai.

}
{

Secara umum, matriks Jacobi dari $\varphi$ adalah

\mathdisp {\begin{pmatrix}  yz^2 &  xz^2 &  2xyz \\  y  &  x  &  -3z^2 \end{pmatrix}} { . }

Untuk titik

\mavergleichskette
{\vergleichskette
{ P
}
{ = }{ (1,3,5)
}
{  }{
}
{    }{
}
{   }{
}
}
{}{}{}
matriks Jacobinya adalah

\mathdisp {\begin{pmatrix}  75 &  25 &  30 \\  3 &  1 &  -75  \end{pmatrix}} {  }

Matriks ini merepresentasikan pemetaan linear
\maabbdisp {L} {T_P\R^3 \cong \R^3} { T_{\varphi(P)} \R^2 \cong \R^2
} {}
terhadap basis-basis standar. Kita tentukan representasi matriks untuk pemetaan
\maabbdisp {\bigwedge^2 L} {\bigwedge^2 \R^3 } { \bigwedge^2 \R^2
} {}
terhadap basis
$e_1 \wedge e_2,\, e_1 \wedge e_3$ , $e_2 \wedge e_3$
\zusatzklammer {di ruas kiri} {} {}
dan
\mathl{e_1 \wedge e_2}{}
\zusatzklammer {di ruas kanan} {} {.}
Untuk itu, kita harus menghitung citra hasil kali-hasil kali baji tersebut. Kita peroleh

\mavergleichskettealign
{\vergleichskettealign
{ (\bigwedge^2 L ) (e_1 \wedge e_2)
}
{ =} { L(e_1) \wedge L(e_2)
}
{ =} { (75e_1 + 3e_2) \wedge (  25 e_1 + 1 e_2 )
}
{ =} { 75 e_1 \wedge e_2 + 75 e_2 \wedge  e_1
}
{ =} { 75 e_1 \wedge e_2 - 75 e_1 \wedge  e_2
}
}
{
\vergleichskettefortsetzungalign
{ =} { 0
}
{ } {}
{ } {}
{ } {}
}
{}{,}

\mavergleichskettealign
{\vergleichskettealign
{ (\bigwedge^2 L ) (e_1 \wedge e_3)
}
{ =} { L(e_1) \wedge L(e_3)
}
{ =} { (75e_1 + 3e_2) \wedge (  30 e_1 - 75 e_2 )
}
{ =} { - 75^2 e_1 \wedge e_2 + 90 e_2 \wedge  e_1
}
{ =} { (-5625 - 90) e_1 \wedge e_2
}
}
{
\vergleichskettefortsetzungalign
{ =} { -5715 e_1 \wedge e_2
}
{ } {}
{ } {}
{ } {}
}
{}{,}
dan

\mavergleichskettealign
{\vergleichskettealign
{ (\bigwedge^2 L ) (e_2 \wedge e_3)
}
{ =} { L(e_2) \wedge L(e_3)
}
{ =} { (  25 e_1 + 1 e_2 )   \wedge ( 30 e_1 - 75 e_2)
}
{ =} { - 1875 e_1 \wedge e_2 + 30 e_2 \wedge  e_1
}
{ =} { -  1905 e_1 \wedge  e_2
}
}
{}
{}{.}
Jadi, matriks yang merepresentasikannya adalah

\mathdisp {\left(  0,-5715,-1905                     \right)} { . }

 }



\inputaufgabeklausurloesung
{3}

{

Misalkan

\mavergleichskette
{\vergleichskette
{ W
}
{ \subseteq }{ \R^m
}
{  }{
}
{    }{
}
{   }{
}
}
{}{}{}
dan

\mavergleichskette
{\vergleichskette
{ U
}
{ \subseteq }{ \R^n
}
{  }{
}
{    }{
}
{   }{
}
}
{}{}{}
merupakan
\definitionsverweis {subhimpunan-subhimpunan terbuka}{}{}
dan misalkan
\maabbdisp {\psi} { W } { U
} {}
suatu
\definitionsverweis {pemetaan}{}{}
yang
\definitionsverweis {terdiferensialkan secara kontinu}{}{.}
Misalkan
\maabbdisp {f} { U } {\R
} {}
suatu fungsi terdiferensialkan secara kontinu. Simpulkan dari
aturan rantai
bahwa

\mavergleichskettedisp
{\vergleichskette
{ d (\psi^*f)
}
{ =} { \psi^*(df)
}
{ } {
}
{ } {
}
{ } {
}
}
{}{}{,}
dengan $\psi^*$ menyatakan
\definitionsverweis {penarikan balik bentuk diferensial}{}{.}

}
{

Misalkan

\mavergleichskette
{\vergleichskette
{ P
}
{ \in }{  W
}
{  }{
}
{    }{
}
{   }{
}
}
{}{}{}
dan

\mavergleichskette
{\vergleichskette
{ v
}
{ \in }{  \R^m
}
{  }{
}
{    }{
}
{   }{
}
}
{}{}{.}
Di satu pihak,

\mavergleichskettedisp
{\vergleichskette
{ d( \psi^* f) (P,v)
}
{ =} { D_P( f \circ \psi) (v)
}
{ =} { \left(  \left(  D_{\psi(P)} f  \right)  \circ \left(  D_P  \psi  \right)  \right) (v)
}
{ =} { D_{\psi(P)} (f) \left(  D_P \psi(v)  \right)
}
{ } {
}
}
{}{}{.}
Di pihak lain, juga berlaku

\mavergleichskettedisp
{\vergleichskette
{ \left(  \psi^*(df)  \right) (P,v)
}
{ =} { (df) \left(  \psi(P), (D_P \psi )(v)  \right)
}
{ =} { D_{\psi(P)} (f) \left(  D_P \psi (v)  \right)
}
{ } {
}
{ } {
}
}
{}{}{.}

 }



\inputaufgabeklausurloesung
{7}

{

Misalkan $M$ suatu
\definitionsverweis {manifold diferensiabel}{}{}
dengan
\definitionsverweis {basis terhitung bagi topologinya}{}{}
dan misalkan $\omega$ suatu
\definitionsverweis {bentuk volume positif}{}{}
pada $M$. Misalkan
\maabbdisp {\varphi} { L } { M
} {}
suatu
\definitionsverweis {difeomorfisme}{}{}
dengan manifold $L$ dan

\mavergleichskette
{\vergleichskette
{ S
}
{ \subseteq }{ L
}
{  }{
}
{    }{
}
{   }{
}
}
{}{}{}
suatu
\definitionsverweis {subhimpunan terukur}{}{.}
Tunjukkan bahwa

\mavergleichskettedisp
{\vergleichskette
{ \int_S \varphi^* \omega
}
{ =} { \int_{ \varphi(S)} \omega
}
{ } {
}
{ } {
}
{ } {
}
}
{}{}{.}

}
{

Menurut
Lema 15.2 (Geometri diferensial (Osnabrück 2023)),
kita dapat mengasumsikan bahwa
\mathkor {} {S} {dan} {\varphi(S)} {}
masing-masing seluruhnya terletak dalam satu peta dari
\mathkor {} {L} {dan} {M} {},
katakanlah

\mavergleichskette
{\vergleichskette
{ S
}
{ \subseteq }{ U
}
{  }{
}
{    }{
}
{   }{
}
}
{}{}{}
dengan peta
\maabbdisp {\alpha} { U } { U' \subseteq \R^n
} {}
dan

\mavergleichskette
{\vergleichskette
{ \varphi(S)
}
{ \subseteq }{ V
}
{  }{
}
{    }{
}
{   }{
}
}
{}{}{}
dengan peta
\maabbdisp {\beta} { V } { V' \subseteq \R^n
} {.}
Dengan memperkecil peta-peta itu lebih lanjut, kita dapat mengasumsikan bahwa
\mathkor {} {U} {dan} {V} {}
difeomorfik satu sama lain melalui $\varphi$. Dengan demikian terdapat diagram komutatif difeomorfisme

\mathdisp {\begin{matrix}  U  & \stackrel{ \varphi }{\longrightarrow} &  V  &  \\ \!\!\!\!\! \alpha \downarrow & & \downarrow \beta \!\!\!\!\!  & \\  U'  & \stackrel{ \beta \circ \varphi \circ \alpha^{-1} }{\longrightarrow} &  V' & \! \! \! \!\!\!\!\!   \\ \end{matrix}} {  }

yang menginduksi diagram komutatif

\mathdisp {\begin{matrix}  S  & \stackrel{ \varphi }{\longrightarrow} &  \varphi (S)  &  \\ \!\!\!\!\! \alpha \downarrow & & \downarrow \beta \!\!\!\!\!  & \\  \alpha (S) & \stackrel{ \beta \circ \varphi \circ \alpha^{-1} }{\longrightarrow} &  \beta ( \varphi (S) ) & \! \! \! \!\!\!\!\!   \\ \end{matrix}} {  }

atas himpunan-himpunan terukur. Untuk bentuk volume $\omega$ pada $V$ berlaku

\mavergleichskettedisp
{\vergleichskette
{ \alpha^{ {-1}*} ( \varphi^* \omega)
}
{ =} { \left(  \beta \circ \varphi \circ \alpha^{-1}  \right)^*  ( \beta^{ {-1}*} \omega)
}
{ } {
}
{ } {
}
{ } {
}
}
{}{}{.}
menurut
Lema 14.7 (Geometri diferensial (Osnabrück 2023))  (4).
Bentuk volume $\beta^{ {-1}*} \omega$ pada $V'$ memiliki representasi
\mathl{g dy_1  \wedge \ldots \wedge dy_n}{} dengan suatu fungsi terukur
\maabb {g} {V'} {\R
} {},
dan menurut definisi, $\int_{\varphi(S)} \omega$ sama dengan $\int_{\beta( \varphi(S))} g d \lambda^n$. Demikian pula,
\mathl{\alpha^{ {-1}*} ( \varphi^* \omega)}{} pada $U'$ memiliki representasi berbentuk
\mathl{f dx_1 \wedge \ldots \wedge dx_n}{.} Bentuk volume pada $U'$ ini sama dengan tarikan balik dari
\mathl{g dy_1  \wedge \ldots \wedge dy_n}{}, dan menurut
Korolari 14.9 (Geometri diferensial (Osnabrück 2023)),
tarikan balik tersebut sama dengan

\mathdisp {(g \circ \psi) \cdot \det  \left(  \left(  { \frac{   \partial   \psi_i   }{  \partial   x_j   } }  \right)_{1 \leq  i, j \leq n}  \right) dx_1 \wedge \ldots \wedge dx_n} {  }

\zusatzklammer {dengan

\mavergleichskettek
{\vergleichskettek
{ \psi
}
{ = }{ \beta \circ \varphi \circ \alpha^{-1}
}
{  }{
}
{    }{
}
{   }{
}
}
{}{}{}} {} {.}
Dengan demikian, pernyataan tersebut mengikuti dari
Ujian dengan solusi (Teori ukuran dan integrasi (Osnabrück 2022-2023)) (Teori ukuran dan integrasi (Osnabrück 2022-2023))
(Geometri diferensial (Osnabrück 2023)).

 }



\inputaufgabeklausurloesung
{2}

{

Pada
\definitionsverweis {bundel vektor trivial}{}{}
\maabbdisp {p} {\R^2 \times \R} { \R^2
} {}
ber-
\definitionsverweis {rank}{}{}
$1$ di atas $\R^2$, kita tinjau
\definitionsverweis {koneksi linear}{}{,}
yang diberikan oleh
\definitionsverweis {simbol-simbol Christoffel}{}{}

\mavergleichskette
{\vergleichskette
{ \Gamma_1 (x,y)
}
{ = }{ x^5-x^2y^2+3y^4
}
{  }{
}
{    }{
}
{   }{
}
}
{}{}{}
dan

\mavergleichskette
{\vergleichskette
{ \Gamma_2 (x,y)
}
{ = }{ 7x^3 +xy^2-4y^5
}
{  }{
}
{    }{
}
{   }{
}
}
{}{}{.}
Hitung
\definitionsverweis {operator kelengkungan}{}{}
$R \left(  \partial_1, \partial_2  \right)$.

}
{

Menurut
Fakta *****,
berlaku

\mavergleichskettealign
{\vergleichskettealign
{ R \left(  \partial_1, \partial_2  \right)
}
{ =} { \partial_1 \Gamma_2 - \partial_2 \Gamma_1
}
{ =} { \partial_1 \left(  7x^3 +xy^2-4y^5  \right) -  \partial_2 \left(   x^5-x^2y^2+3y^4   \right)
}
{ =} {  21x^2 +y^2 - 5x^4 -2xy^2
}
{ } {
}
}
{}
{}{.}

 }



\inputaufgabeklausurloesung
{0}

{

}
{  }



\inputaufgabeklausurloesung
{6}

{

Buktikan teorema tentang retraksi ke batas pada manifold.

}
{

Menurut
Teorema 21.8 (Geometri diferensial (Osnabrück 2023)),
batas
\mathl{\partial M}{} merupakan
\definitionsverweis {manifold diferensiabel}{}{}
\definitionsverweis {berorientasi}{}{}
\zusatzklammer {tanpa batas} {} {.}
Oleh karena itu, menurut
Teorema 22.11 (Geometri diferensial (Osnabrück 2023)),
terdapat suatu
\definitionsverweis {bentuk volume positif}{}{}
\definitionsverweis {yang terdiferensialkan secara kontinu}{}{}
$\tau$ pada
\mathl{\partial M}{.} Kita mempunyai
\mathl{\int_{  \partial M  }  \tau >0}{.}
\definitionsverweis {Turunan eksterior}{}{}
dari bentuk volume $\tau$ adalah $0$. Andaikan terdapat suatu
\definitionsverweis {pemetaan yang terdiferensialkan secara kontinu}{}{}
\maabbdisp {\varphi} { M } { \partial M
} {}
dengan
\mathl{\varphi \vert_{\partial M} = \operatorname{Id}_{\partial M}}{}. Maka
\definitionsverweis {bentuk tarikan balik}{}{}

\mathl{\varphi^* \tau}{} adalah
$(n-1)$-\definitionsverweis {bentuk diferensial}{}{}
pada $M$ yang pembatasannya pada batas berimpit dengan $\tau$. Dengan menggunakan
Teorema 23.2 (Geometri diferensial (Osnabrück 2023))
dan
Teorema 20.4 (Geometri diferensial (Osnabrück 2023))  (5),
kita memperoleh

\mavergleichskettealign
{\vergleichskettealign
{
\int_{  \partial M  }  \tau
}
{ =} {
\int_{  \partial  M  }  \varphi^* \tau
}
{ =} {
\int_{   M  }  d(\varphi^* \tau)
}
{ =} {
\int_{   M  }  \varphi^* (d\tau)
}
{ =} {
\int_{   M  }  \varphi^* (0)
}
}
{
\vergleichskettefortsetzungalign
{ =} { 0
}
{ } {}
{ } {}
{ } {}
}
{}{.}
Ini merupakan suatu kontradiksi.

 }

\endgroup

\chapter{Formulir Solusi Resmi Ujian 4}
\noindent\textit{Solusi yang disediakan oleh sumber resmi; kekosongan sumber dipertahankan.}
\begingroup
\renewcommand{\theHfakt}{o011.exam.official.04.\arabic{fakt}}
\input{generated/exams/exam04_solutions.id.build.tex}
\endgroup

\chapter{Formulir Solusi Resmi Ujian 5}
\noindent\textit{Solusi yang disediakan oleh sumber resmi; kekosongan sumber dipertahankan.}
\begingroup
\renewcommand{\theHfakt}{o011.exam.official.05.\arabic{fakt}}
\input{generated/exams/exam05_solutions.id.build.tex}
\endgroup

\chapter{Formulir Solusi Resmi Ujian 6}
\noindent\textit{Solusi yang disediakan oleh sumber resmi; kekosongan sumber dipertahankan.}
\begingroup
\renewcommand{\theHfakt}{o011.exam.official.06.\arabic{fakt}}
\input{generated/exams/exam06_solutions.id.build.tex}
\endgroup

\chapter{Formulir Solusi Resmi Ujian 7}
\noindent\textit{Solusi yang disediakan oleh sumber resmi; kekosongan sumber dipertahankan.}
\begingroup
\renewcommand{\theHfakt}{o011.exam.official.07.\arabic{fakt}}
\input{generated/exams/exam07_solutions.id.build.tex}
\endgroup

\chapter{Formulir Solusi Resmi Ujian 8}
\noindent\textit{Solusi yang disediakan oleh sumber resmi; kekosongan sumber dipertahankan.}
\begingroup
\renewcommand{\theHfakt}{o011.exam.official.08.\arabic{fakt}}
%Data institusi

%\input{Dozentdaten}

%\renewcommand{\fachbereich}{Fachbereich}

%\renewcommand{\dozent}{Prof. Dr. . }

%Data ujian

\renewcommand{\klausurgebiet}{ }

\renewcommand{\klausurtyp}{ }

\renewcommand{\klausurdatum}{ . 20}

\klausurvorspann {\fachbereich} {\klausurdatum} {\dozent} {\klausurgebiet} {\klausurtyp}

%Data untuk tabel nilai berikut

\renewcommand{\aeins}{  3  }

\renewcommand{\azwei}{  3   }

\renewcommand{\adrei}{  3  }

\renewcommand{\avier}{  0  }

\renewcommand{\afuenf}{  2  }

\renewcommand{\asechs}{  0  }

\renewcommand{\asieben}{  3  }

\renewcommand{\aacht}{  7  }

\renewcommand{\aneun}{  0  }

\renewcommand{\azehn}{  7  }

\renewcommand{\aelf}{  4  }

\renewcommand{\azwoelf}{  3  }

\renewcommand{\adreizehn}{  3   }

\renewcommand{\avierzehn}{  7  }

\renewcommand{\afuenfzehn}{  2  }

\renewcommand{\asechzehn}{  9  }

\renewcommand{\asiebzehn}{  7  }

\renewcommand{\aachtzehn}{  63  }

\renewcommand{\aneunzehn}{    }

\renewcommand{\azwanzig}{    }

\renewcommand{\aeinundzwanzig}{    }

\renewcommand{\azweiundzwanzig}{    }

\renewcommand{\adreiundzwanzig}{    }

\renewcommand{\avierundzwanzig}{    }

\renewcommand{\afuenfundzwanzig}{    }

\renewcommand{\asechsundzwanzig}{    }

\punktetabellesiebzehn

\klausurnote

\newpage

\setcounter{section}{0}



\inputaufgabeklausurloesung
{3}

{

Definisikan istilah-istilah berikut
\zusatzklammer {yang dicetak miring} {} {}.
\aufzaehlungsechs{\stichwort {Pemetaan Gauss} {}
dari suatu hiperpermukaan diferensiabel

\mavergleichskette
{\vergleichskette
{ Y
}
{ \subseteq }{ \R^n
}
{  }{
}
{    }{
}
{   }{
}
}
{}{}{.}

}{Suatu
\stichwort {peta topologis} {}
pada suatu
\definitionsverweis {manifold topologis}{}{}
$M$.

}{Suatu $C^1$-\stichwort {manifold diferensiabel} {} $M$.

}{\stichwort {Pangkat eksterior} {} ke-$n$ dari suatu
$K$-\definitionsverweis {ruang vektor}{}{}
$V$
\zusatzklammer {cukup tuliskan simbolnya} {} {.}

}{Suatu
\stichwort {isometri lokal} {}
\maabb {\varphi} {L} {M
} {}
antara manifold-manifold Riemann.

}{Suatu koneksi
\stichwort {terintegralkan secara lokal} {}
pada suatu bundel vektor diferensiabel
\maabb {p} {E} {X
} {}
di atas suatu manifold diferensiabel $X$.
}

}
{

\aufzaehlungsechs{Misalkan suatu
\definitionsverweis {orientasi}{}{}
pada $Y$ telah ditetapkan. Maka pemetaan
\maabbdisp {} {Y} { S^{n-1}
} {,}
yang memetakan setiap titik

\mavergleichskette
{\vergleichskette
{ P
}
{ \in }{ Y
}
{  }{
}
{    }{
}
{   }{
}
}
{}{}{}
ke vektor normal satuan yang ditentukan oleh orientasi tersebut disebut
pemetaan Gauss
dari $Y$.
}{Suatu peta topologis adalah setiap
\definitionsverweis {homeomorfisme}{}{}
\maabbdisp {\varphi} {U} {V
} {,}
dengan

\mavergleichskette
{\vergleichskette
{ U
}
{ \subseteq }{ M
}
{  }{
}
{    }{
}
{   }{
}
}
{}{}{}
dan

\mavergleichskette
{\vergleichskette
{ V
}
{ \subseteq }{ \R^n
}
{  }{
}
{    }{
}
{   }{
}
}
{}{}{}
\definitionsverweis {terbuka}{}{.}
}{Suatu
\definitionsverweis {ruang Hausdorff}{}{}
\definitionsverweis {topologis}{}{}
$M$ bersama suatu
\definitionsverweis {penutup terbuka}{}{}

\mathl{M= \bigcup_{i \in I} U_i}{} dan
\definitionsverweis {peta-peta}{}{}
\maabbdisp {\alpha_i} {U_i } { V_i
} {}
dengan
\mathl{V_i \subseteq \R^n}{} terbuka, sedemikian rupa sehingga
\definitionsverweis {pemetaan-pemetaan transisi}{}{}
\maabbdisp {\alpha_j \circ (\alpha_i)^{-1}} {V_i \cap \alpha_i(U_i \cap U_j)   } { V_j \cap \alpha_j(U_i \cap U_j)
} {}
$C^1$-\definitionsverweis {difeomorfisme}{}{}
untuk semua
\mathl{i,j \in  I}{}. Struktur ini disebut manifold diferensiabel.
}{Ruang vektor-$K$
\mathl{\bigwedge^n V}{} disebut pangkat eksterior ke-$n$ dari $V$.
}{Suatu
\definitionsverweis {pemetaan diferensiabel}{}{}
\maabb {\varphi} {L} {M
} {}
disebut
isometri lokal
jika untuk setiap titik

\mavergleichskette
{\vergleichskette
{ P
}
{ \in }{ L
}
{  }{
}
{    }{
}
{   }{
}
}
{}{}{}
pemetaan
\definitionsverweis {tangen}{}{}
\maabbdisp {T_P \varphi} {T_PL} { T_{\varphi(P) } M
} {}
merupakan suatu
\definitionsverweis {isometri}{}{}
terhadap hasil kali skalar yang diberikan.
}{Koneksi tersebut disebut
terintegralkan secara lokal
jika untuk setiap titik

\mavergleichskette
{\vergleichskette
{ e
}
{ \in }{ E
}
{  }{
}
{    }{
}
{   }{
}
}
{}{}{}
terdapat suatu

\mavergleichskette
{\vergleichskette
{ p(e)
}
{ \in }{ U
}
{ \subseteq }{ X
}
{    }{
}
{   }{
}
}
{}{}{}
\definitionsverweis {penampang horizontal}{}{}
\maabbdisp {s} {U} {E \vert_U
} {}
yang didefinisikan pada lingkungan terbuka tersebut dan melalui $e$.
}

 }



\inputaufgabeklausurloesung
{3}

{

Rumuskan teorema-teorema berikut.
\aufzaehlungdrei{\stichwort {Teorema isometri untuk transpor paralel pada suatu hiperpermukaan} {}

\mavergleichskette
{\vergleichskette
{  Y
}
{ \subseteq }{  \R^n
}
{  }{
}
{    }{
}
{   }{
}
}
{}{}{.}}{\stichwort {Theorema egregium} {.}}{Teorema tentang partisi kesatuan.}

}
{

\aufzaehlungdrei{\faktsituation {Misalkan

\mavergleichskette
{\vergleichskette
{ Y
}
{ \subseteq }{ W
}
{ \subseteq }{ \R^n
}
{    }{
}
{   }{
}
}
{}{}{,}
$W$
\definitionsverweis {terbuka}{}{,}
dan merupakan suatu
\definitionsverweis {hiperpermukaan diferensiabel}{}{.}
Misalkan
\maabb {\gamma} {[a,b]} {Y
} {}
suatu
\definitionsverweis {kurva diferensiabel}{}{}
dengan

\mavergleichskette
{\vergleichskette
{ \gamma(a)
}
{ = }{ P
}
{  }{
}
{    }{
}
{   }{
}
}
{}{}{}
dan

\mavergleichskette
{\vergleichskette
{ \gamma(b)
}
{ = }{ Q
}
{  }{
}
{    }{
}
{   }{
}
}
{}{}{.}}
\faktfolgerung {Maka
\definitionsverweis {transpor paralel}{}{}
sepanjang $\gamma$ merupakan suatu
\definitionsverweis {isometri}{}{}
\maabbdisp {} {T_PY} {T_QY
} {.}}
\faktzusatz {}
\faktzusatz {}}{\faktsituation {Misalkan

\mavergleichskette
{\vergleichskette
{ Y
}
{ \subseteq }{ \R^3
}
{  }{
}
{    }{
}
{   }{
}
}
{}{}{}
suatu
\definitionsverweis {permukaan berorientasi}{}{}
dan misalkan
\maabbdisp {\varphi} { V } { Y
} {,}

\mavergleichskette
{\vergleichskette
{ V
}
{ \subseteq }{ \R^2
}
{  }{
}
{    }{
}
{   }{
}
}
{}{}{,}
suatu parametrisasi lokal $Y$ yang dua kali terdiferensialkan, dengan parameter $u,v$. Misalkan
\mathl{\left(  g_{ij}  \right)}{}
\definitionsverweis {matriks fundamental pertama}{}{}
pada $V$.}
\faktfolgerung {Maka, dengan menggunakan
\definitionsverweis {simbol-simbol Christoffel}{}{,}
\definitionsverweis {kelengkungan Gauss}{}{}
memenuhi hubungan

\mavergleichskettedisphandlinks
{\vergleichskettedisphandlinks
{ K
}
{ =} { { \frac{   1  }{ g_{11}  } }  \left(  \partial_2 \Gamma_{11}^2 - \partial_1 \Gamma_{12}^2 + \Gamma^1_{11} \Gamma^2_{12}  + \Gamma_{11}^2 \Gamma_{22}^2 - \Gamma_{12}^1 \Gamma_{11}^2-  \left(  \Gamma_{12}^2  \right)^2  \right)
}
{ } {
}
{ } {
}
{ } {
}
}
{}{}{.}}
\faktzusatz {}
\faktzusatz {}}{Misalkan $M$ suatu
\definitionsverweis {manifold diferensiabel}{}{}
dengan
\definitionsverweis {basis topologi yang terhitung}{}{.}
Maka, untuk setiap penutup terbuka, terdapat suatu
\definitionsverweis {partisi kesatuan}{}{}
\definitionsverweis {yang terdiferensialkan secara kontinu}{}{}
dan tersubordinasi pada penutup tersebut.}

 }



\inputaufgabeklausurloesung
{3}

{

Misalkan
\maabbdisp {f,g} {\R} { \R^n
} {}
\definitionsverweis {kurva-kurva diferensiabel}{}{.}
Hitung
\definitionsverweis {turunan}{}{}
dari fungsi
\maabbeledisp {} {\R} {\R
} { t } { \left\langle f(t) , g(t)  \right\rangle
} {.}
Rumuskan hasilnya dengan menggunakan hasil kali dalam.

}
{

Misalkan $f_i(t)$ dan $g_i(t)$ masing-masing fungsi komponen dari
\mathkor {} {f} {dan} {g} {.}
Fungsi yang dimaksud adalah

\mavergleichskettedisp
{\vergleichskette
{  \left\langle f(t) , g(t)  \right\rangle
}
{ =} { \sum_{i = 1}^n  f_i(t) g_i(t)
}
{ } {
}
{ } {
}
{ } {
}
}
{}{}{}
dengan turunan
\zusatzklammer {aturan hasil kali diterapkan pada setiap suku} {} {}

\mathdisp {\sum_{i = 1}^n  f'_i(t) g_i(t)  + \sum_{i = 1}^n  f_i(t) g'_i(t)} { . }

Dengan hasil kali skalar, bentuk ini dapat ditulis sebagai

\mathdisp {\left\langle f'(t) ,  g(t)  \right\rangle + \left\langle f(t) ,  g'(t)  \right\rangle} {  }

.

 }



\inputaufgabeklausurloesung
{0}

{

}
{  }



\inputaufgabeklausurloesung
{2}

{

Misalkan

\mavergleichskette
{\vergleichskette
{ Y
}
{ \subseteq }{ W
}
{ \subseteq }{ \R^n
}
{    }{
}
{   }{
}
}
{}{}{,}
$W$
\definitionsverweis {terbuka}{}{,}
suatu
\definitionsverweis {hiperpermukaan diferensiabel}{}{},
dan
\maabb {\gamma} {I =[a,b]} {Y
} {}
suatu
\definitionsverweis {kurva geodesik}{}{.}
Buktikan bahwa
\definitionsverweis {transpor paralel}{}{}
sepanjang $\gamma$ memetakan vektor tangensial

\mavergleichskette
{\vergleichskette
{ \gamma'(a)
}
{ \in }{ T_{\gamma(a)}Y
}
{  }{
}
{    }{
}
{   }{
}
}
{}{}{}
ke vektor tangensial

\mavergleichskette
{\vergleichskette
{ \gamma'(b)
}
{ \in }{ T_{\gamma(b)}Y
}
{  }{
}
{    }{
}
{   }{
}
}
{}{}{.}

}
{

Menurut
Lemma 6.8 (Differentialgeometrie (Osnabrück 2023))
medan kecepatan $\gamma'$ merupakan suatu
\definitionsverweis {medan vektor paralel}{}{}
sepanjang $\gamma$. Oleh karena itu,
\definitionsverweis {transpor paralel}{}{}
$\Psi_\gamma$ sepanjang $\gamma$ memetakan vektor

\mavergleichskette
{\vergleichskette
{ \gamma'(a)
}
{ \in }{  T_{\gamma(a)} Y
}
{  }{
}
{    }{
}
{   }{
}
}
{}{}{}
ke vektor

\mavergleichskette
{\vergleichskette
{ \gamma'(b)
}
{ \in }{  T_{\gamma(b)} Y
}
{  }{
}
{    }{
}
{   }{
}
}
{}{}{}
.

 }



\inputaufgabeklausurloesung
{0}

{

}
{  }



\inputaufgabeklausurloesung
{3}

{

Apakah pemetaan
\maabbeledisp {\varphi} {S^1} {S^1
} {(x,y)} {( \betrag {  x  } ,y)
} {,}
\definitionsverweis {terdiferensialkan}{}{?}

}
{

Kita meninjau komposisi pemetaan-pemetaan

\mathdisp {\R \stackrel{ \theta }{\longrightarrow} S^1 \stackrel{\varphi}{\longrightarrow} S^1 \subset \R^2 \stackrel{p_1}{\longrightarrow} \R} { , }

dengan

\mavergleichskettedisp
{\vergleichskette
{ \theta(t)
}
{ =} {(\cos  t , \sin  t )
}
{ } {
}
{ } {
}
{ } {
}
}
{}{}{}
dan $p_1$ merupakan proyeksi pertama. Parametrisasi trigonometrik, inklusi
\zusatzklammer {suatu submanifold tertutup} {} {,}
dan proyeksi tersebut semuanya diferensiabel. Jika $\varphi$ diferensiabel, maka komposisi keseluruhannya juga harus diferensiabel. Akan tetapi, komposisi itu adalah
\maabbeledisp {} {\R} {\R
} { t } {\betrag {  \cos  t  }
} {.}
Pemetaan ini tidak diferensiabel di
\mathl{t= \pi/2}{} karena gradien dari kiri di titik tersebut adalah $-1$, sedangkan gradien dari kanan adalah $1$. Jadi, pemetaan $\varphi$ tidak diferensiabel.

 }



\inputaufgabeklausurloesung
{7}

{

Misalkan $M$ suatu manifold diferensiabel kompak dengan bentuk volume positif kontinu $\omega$. Buktikan bahwa

\mavergleichskettedisp
{\vergleichskette
{ \int_M \omega
}
{ <} { \infty
}
{ } {
}
{ } {
}
{ } {
}
}
{}{}{}.

}
{

Untuk setiap titik

\mavergleichskette
{\vergleichskette
{ P
}
{ \in }{ M
}
{  }{
}
{    }{
}
{   }{
}
}
{}{}{}
terdapat suatu lingkungan peta terbuka

\mavergleichskette
{\vergleichskette
{ P
}
{ \in }{ U
}
{  }{
}
{    }{
}
{   }{
}
}
{}{}{}
dan suatu pemetaan peta
\maabbdisp {\alpha} { U } { V
} {}
dengan

\mavergleichskette
{\vergleichskette
{ V
}
{ \subseteq }{ \R^n
}
{  }{
}
{    }{
}
{   }{
}
}
{}{}{}
terbuka, sedemikian rupa sehingga

\mavergleichskette
{\vergleichskette
{ \alpha^{-1 *} \omega
}
{ = }{ fdx_1 \wedge \ldots \wedge dx_n
}
{  }{
}
{    }{
}
{   }{
}
}
{}{}{}
berlaku dengan
\mathl{f}{} kontinu dan positif. Kita juga memperoleh suatu lingkungan terbuka

\mavergleichskette
{\vergleichskette
{ P
}
{ \in }{ U'
}
{ \subseteq }{ U
}
{    }{
}
{   }{
}
}
{}{}{,}
yang homeomorf dengan suatu bola terbuka

\mavergleichskette
{\vergleichskette
{ U'
}
{ \cong }{ B'
}
{ \subseteq }{ V
}
{    }{
}
{   }{
}
}
{}{}{}
dan kita dapat menganggap bahwa penutupan bola tersebut seluruhnya terletak di $V$. Bola tertutup itu tertutup dan terbatas; karena itu, fungsi kontinu $f$ terbatas pada bola tersebut dan dengan demikian juga pada $U'$. Maka
\mathl{\int_{U'} \omega}{} berhingga, dengan $U'$ suatu lingkungan terbuka dari $P$.

Himpunan-himpunan terbuka

\mavergleichskette
{\vergleichskette
{ U'
}
{ = }{ U'(P)
}
{  }{
}
{    }{
}
{   }{
}
}
{}{}{}
ini menutupi $M$. Karena kekompakan, terdapat suatu penutup berhingga

\mavergleichskettedisp
{\vergleichskette
{ M
}
{ =} { \bigcup_{i = 1} ^n U_i
}
{ } {
}
{ } {
}
{ } {
}
}
{}{}{}
dengan

\mavergleichskettedisp
{\vergleichskette
{ \int_{U_i } \omega
}
{ <} { \infty
}
{ } {
}
{ } {
}
{ } {
}
}
{}{}{.}
Karena kepositifan, diperoleh

\mavergleichskettedisp
{\vergleichskette
{ \int_{M} \omega
}
{ \leq} { \sum_{i = 1}^n  \left(  \int_{U_i } \omega  \right)
}
{ <} { \infty
}
{ } {
}
{ } {
}
}
{}{}{.}

 }



\inputaufgabeklausurloesung
{0}

{

}
{  }



\inputaufgabeklausurloesung
{7 (1+2+4)}

{

\aufzaehlungdreiabc{ Berikan suatu sifat topologis yang menunjukkan bahwa
\definitionsverweis {interval satuan terbuka}{}{}
dan
\definitionsverweis {lingkaran}{}{}
$S^1$ tidak
\definitionsverweis {homeomorfik}{}{.}
}{Buktikan bahwa setiap
$1$-\definitionsverweis {bentuk diferensial}{}{}
kontinu pada
\mathl{]0,1[}{}
bersifat
\definitionsverweis {eksak}{}{.}
}{ Berikan suatu
$1$-bentuk konkret pada $S^1$ yang
\definitionsverweis {tertutup}{}{}
tetapi tidak eksak.
}

}
{

a) Menurut
Satz 13.14 (Differentialgeometrie (Osnabrück 2023))
$S^1$, sebagai himpunan bagian tertutup dan terbatas dari $\R^2$, bersifat
\definitionsverweis {kompak}{}{,}
sedangkan interval terbuka tersebut tidak kompak.

b) Berlaku

\mavergleichskettedisp
{\vergleichskette
{ \omega
}
{ =} { hdt
}
{ } {
}
{ } {
}
{ } {
}
}
{}{}{}
dengan suatu fungsi kontinu
\maabbdisp {h} {]0,1[} {\R
} {.}
Menurut
Hauptsatz der Infinitesimalrechnung,
$h$ mempunyai suatu antiturunan $H$. Dalam bahasa bentuk diferensial, hal ini berarti

\mavergleichskettedisp
{\vergleichskette
{H'
}
{ =} { hdt
}
{ =} { \omega
}
{ } {
}
{ } {
}
}
{}{}{,}
sehingga $\omega$ eksak.

c) Misalkan
\mathl{S^1 \subset \R^2}{} diberikan, dengan koordinat pada $\R^2$ dinotasikan oleh
\mathkor {} {u} {dan} {v} {}
. Kita meninjau bentuk diferensial

\mavergleichskettedisp
{\vergleichskette
{ \omega
}
{ =} { udv-vdu
}
{ } {
}
{ } {
}
{ } {
}
}
{}{}{}
pada $S^1$. Karena $S^1$ berdimensi satu, bentuk ini tertutup. Untuk lintasan
\maabbeledisp {\gamma} {[0, 2 \pi]} {S^1
} { t } { \begin{pmatrix}  \cos  t  \\ \sin  t     \end{pmatrix}
} {,}
berlaku

\mavergleichskettedisp
{\vergleichskette
{ \int_\gamma \omega
}
{ =} { \int_0^{2 \pi} \gamma^* \omega  dt
}
{ =} { \int_0^{2 \pi}  \cos  t  \cdot \cos  t  +  \sin  t   \sin  t   dt
}
{ =} { \int_0^{2 \pi} 1 dt
}
{ =} { 2 \pi
}
}
{}{}{.}
Menurut
Korollar 23.3 (Differentialgeometrie (Osnabrück 2023))
$\omega$ tidak mungkin eksak.

 }



\inputaufgabeklausurloesung
{4}

{

Misalkan

\mavergleichskette
{\vergleichskette
{ U
}
{ \subseteq }{ \R^n
}
{  }{
}
{    }{
}
{   }{
}
}
{}{}{}
\definitionsverweis {terbuka}{}{}
dan
\maabb {f} { U } { \R
} {}
suatu
\definitionsverweis {fungsi yang terdiferensialkan secara kontinu}{}{}
dua kali. Buktikan bahwa

\mavergleichskettedisp
{\vergleichskette
{ d(df)
}
{ =} { 0
}
{ } {
}
{ } {
}
{ } {
}
}
{}{}{,}
dengan $d$ menyatakan
\definitionsverweis {turunan eksterior}{}{.}

}
{

Untuk suatu bentuk-$1$

\mavergleichskette
{\vergleichskette
{ \omega
}
{ = }{ \sum_{j = 1}^n g_jdx_j
}
{  }{
}
{    }{
}
{   }{
}
}
{}{}{}
dengan menggunakan

\mavergleichskette
{\vergleichskette
{ dx_i \wedge dx_j
}
{ = }{ -dx_j \wedge dx_i
}
{  }{
}
{    }{
}
{   }{
}
}
{}{}{}

\mavergleichskettealign
{\vergleichskettealign
{d \omega
}
{ =} { \sum_{j = 1 }^n  d(g_j dx_j)
}
{ =} { \sum_{j = 1 }^n   \left(  \sum_{i = 1}^n { \frac{   \partial  g_j   }{  \partial   x_i   } } dx_i \wedge dx_j  \right)
}
{ =} { \sum _{1 \leq i <j\leq n} \left(  { \frac{   \partial  g_j   }{  \partial   x_i   } } - { \frac{   \partial  g_i   }{  \partial   x_j   } }  \right) dx_i \wedge dx_j
}
{ } {
}
}
{}
{}{.}
Untuk suatu fungsi $f$ yang dua kali terdiferensialkan secara kontinu, berlaku

\mavergleichskette
{\vergleichskette
{ df
}
{ = }{ \sum_{j = 1}^n g_j dx_j
}
{  }{
}
{    }{
}
{   }{
}
}
{}{}{}
dengan
\definitionsverweis {turunan-turunan parsial}{}{}

\mavergleichskette
{\vergleichskette
{ g_j
}
{ = }{ { \frac{   \partial   f   }{  \partial   x_j   } }
}
{  }{
}
{    }{
}
{   }{
}
}
{}{}{,}
dan karena itu

\mavergleichskette
{\vergleichskette
{ d(df)
}
{ = }{ 0
}
{  }{
}
{    }{
}
{   }{
}
}
{}{}{}
menurut
Satz von Schwarz.

 }



\inputaufgabeklausurloesung
{3}

{

Deskripsikan
\definitionsverweis {pemetaan}{}{}
\maabbeledisp {\psi} { \Complex \setminus \{1\} } { \Complex
} {w} { { \frac{   (w+1)  { \mathrm i}  }{ 1-w  } }
} {}
dalam koordinat riil.

}
{

Kita menulis

\mavergleichskette
{\vergleichskette
{ w
}
{ = }{ a+b { \mathrm i}
}
{  }{
}
{    }{
}
{   }{
}
}
{}{}{.}
Maka

\mavergleichskettealign
{\vergleichskettealign
{  \psi ( a+b { \mathrm i}   )
}
{ =} { { \frac{   \left(  a+b { \mathrm i} +1  \right) { \mathrm i}      }{  1-  a-b { \mathrm i}      } }
}
{ =} { { \frac{     a { \mathrm i} +{ \mathrm i}  -b  }{  1-  a-b { \mathrm i}      } }
}
{ =} { { \frac{     \left(  a { \mathrm i} +{ \mathrm i}  -b  \right) \left(  1-  a+b { \mathrm i}  \right)  }{  \left(  1-  a-b { \mathrm i}  \right) \left(  1-  a+b { \mathrm i}  \right)      } }
}
{ =} { { \frac{     \left(  a+1 \right) \left(  1-a  \right) { \mathrm i}  -b^2 { \mathrm i} -b \left(  1-  a \right) -b\left(  a+1  \right)  }{  \left(  1-  a \right)^2 +b^2  } }
}
}
{
\vergleichskettefortsetzungalign
{ =} { { \frac{     \left(  1-a ^2-b^2  \right) { \mathrm i}   -2 b  }{  \left(  1-  a \right)^2 +b^2  } }
}
{ } {
}
{ } {}
{ } {}
}
{}{.}
Dalam koordinat riil, dengan demikian,

\mavergleichskettedisp
{\vergleichskette
{ \psi \begin{pmatrix}  a  \\b   \end{pmatrix}
}
{ =} { { \frac{   1  }{  -2a+1+a^2+b^2 } } \begin{pmatrix}  -2b \\ 1-a^2-b^2   \end{pmatrix}
}
{ } {
}
{ } {
}
{ } {
}
}
{}{}{.}

 }



\inputaufgabeklausurloesung
{3}

{

Misalkan

\mavergleichskette
{\vergleichskette
{ H
}
{ \subseteq }{ \R^n
}
{  }{
}
{    }{
}
{   }{
}
}
{}{}{}
suatu
\definitionsverweis {setengah ruang Euklides}{}{}
dan

\mavergleichskette
{\vergleichskette
{ P
}
{ \in }{ H
}
{  }{
}
{    }{
}
{   }{
}
}
{}{}{.}
Andaikan terdapat suatu himpunan yang terbuka di $\R^n$, yaitu $V$, dengan

\mavergleichskette
{\vergleichskette
{ P
}
{ \in }{ V
}
{ \subseteq }{ H
}
{    }{
}
{   }{
}
}
{}{}{.}
Buktikan bahwa
\mathl{P}{} bukan titik batas dari $H$.

}
{

Misalkan

\mavergleichskettedisp
{\vergleichskette
{ P
}
{ =} { \begin{pmatrix}  x_1  \\ x_2 \\  \vdots\\ x_n   \end{pmatrix}
}
{ } {
}
{ } {
}
{ } {
}
}
{}{}{.}
Kita dapat mengganti, di $\R^n$, lingkungan terbuka $V$ dengan suatu bola terbuka
\zusatzklammer {di $\R^n$} {} {}
\mathl{U \left(   P , \epsilon  \right)}{} dengan

\mavergleichskette
{\vergleichskette
{ P
}
{ \in }{   U \left(   P , \epsilon  \right)
}
{ \subseteq }{  V
}
{ \subseteq   }{ H
}
{   }{
}
}
{}{}{}
. Jika $P$ suatu titik batas, maka

\mavergleichskette
{\vergleichskette
{ x_1
}
{ = }{ 0
}
{  }{
}
{    }{
}
{   }{
}
}
{}{}{.}
Akan tetapi, dalam hal itu

\mavergleichskette
{\vergleichskette
{ \begin{pmatrix}  - { \frac{   \epsilon }{  2 } }  \\ x_2 \\  \vdots\\ x_n   \end{pmatrix}
}
{ \in }{ U \left(   P , \epsilon  \right)
}
{  }{
}
{    }{
}
{   }{
}
}
{}{}{,}
padahal titik ini tidak termasuk dalam $H$.

 }



\inputaufgabeklausurloesung
{7}

{

Misalkan $M$ suatu
\definitionsverweis {manifold diferensiabel}{}{}
berdimensi $n$ dan

\mavergleichskette
{\vergleichskette
{ P
}
{ \in }{ M
}
{  }{
}
{    }{
}
{   }{
}
}
{}{}{}
suatu titik. Buktikan bahwa terdapat suatu
\definitionsverweis {manifold dengan batas}{}{}
$N$ yang batasnya $\partial N$ difeomorfik dengan permukaan bola berdimensi $(n-1)$, yaitu $S^{n-1}$, sedemikian sehingga
\mathkor {} {M \setminus \{ P \}} {dan} {N \setminus \partial N} {}
saling difeomorfik.

}
{

Misalkan

\mavergleichskette
{\vergleichskette
{ P
}
{ \in }{  U
}
{ \subseteq }{  M
}
{    }{
}
{   }{
}
}
{}{}{}
suatu daerah peta terbuka dengan peta
\maabbdisp {\alpha} { U } {V \subseteq \R^n
} {.}
Kita dapat menganggap bahwa $\alpha(P)$ adalah titik nol dan $V$ adalah bola terbuka berjari-jari $3$. Misalkan
\maabbdisp {\varphi} { U \left(   0 , 3  \right) \setminus \{0\}  } { U \left(   0 , 3  \right) \setminus B  \left(  0 , 1 \right)
} {}
suatu difeomorfisme yang merupakan identitas pada
\mathl{U \left(   0 , 3  \right) \setminus U \left(   0 , 2  \right)}{} dan memetakan
\mathl{U \left(   0 , 2  \right) \setminus \{0\}}{} ke
\mathl{U \left(   0 , 2  \right) \setminus B  \left(  0 , 1 \right)}{}. Pemetaan seperti ini diperoleh dengan meregangkan titik-titik menggunakan fungsi
\maabbdisp {h} {[0,3]} { [0,3]
} {}
dengan

\mavergleichskettedisp
{\vergleichskette
{ h(r)
}
{ =} { \begin{cases} 1+ { \frac{   1  }{  4 } } r^2 \text{ jika } 0 \leq r \leq 2 \, , \\ r \text{ jika } r >  2  \, , \end{cases}
}
{ } {
}
{ } {
}
{ } {
}
}
{}{}{}
. Selanjutnya, citra difeomorfisme $\varphi$, yaitu
\mathl{U \left(   0 , 3  \right) \setminus B  \left(  0 , 1 \right)}{,} dapat dipandang sebagai

\mavergleichskettedisp
{\vergleichskette
{  U \left(   0 , 3  \right) \setminus B  \left(  0 , 1 \right)
}
{ \subset} {  U \left(   0 , 3  \right) \setminus U \left(   0 , 1  \right)
}
{ } {
}
{ } {
}
{ } {
}
}
{}{}{}
. Himpunan yang lebih besar merupakan suatu manifold dengan batas $S \left(   0 ,  1  \right)$. Misalkan $N$ manifold yang diperoleh dengan mengganti $U$ oleh
\mathl{U \left(   0 , 3  \right) \setminus U \left(   0 , 1  \right)}{}. Permukaan bola tersebut menjadi batas dari $N$. Difeomorfisme $\varphi$ dapat diperluas menjadi suatu difeomorfisme
\maabbdisp {} {M \setminus \{P\} } { N \setminus \partial N
} {}
karena pada daerah transisi terbuka
\mathl{U \left(   0 , 3  \right) \setminus U \left(   0 , 2  \right)}{} pemetaan tersebut merupakan identitas.

 }



\inputaufgabeklausurloesung
{2}

{

Berikan suatu
\definitionsverweis {penghabisan kompak}{}{}
untuk $\R^d$.

}
{

Kita meninjau
\definitionsverweis {bola-bola tertutup}{}{}

\mathbed {B  \left(  0 , n  \right)} {}
 {n \in \N_+} {}
 {} {} {} {,}
yang
\definitionsverweis {kompak}{}{}
dan menutupi $\R^d$. Interior terbuka dari
\mathl{B  \left(  0 , n+1  \right)}{}
adalah
\definitionsverweis {bola terbuka}{}{}

\mathl{U \left(   0,n+1  \right)}{}, dan karena

\mavergleichskettedisp
{\vergleichskette
{ B  \left(  0,n \right)
}
{ \subseteq} { U \left(   0,n+1  \right)
}
{ } {
}
{ } {
}
{ } {
}
}
{}{}{}
diperoleh suatu
\definitionsverweis {penghabisan kompak}{}{.}

 }



\inputaufgabeklausurloesung
{9}

{

Buktikan teorema titik tetap Brouwer.

}
{

Untuk menyederhanakan notasi, misalkan

\mavergleichskette
{\vergleichskette
{ r
}
{ = }{ 1
}
{  }{
}
{    }{
}
{   }{
}
}
{}{}{.}  Andaikan terdapat suatu pemetaan $\psi$ tanpa titik tetap yang terdiferensialkan secara kontinu. Maka selalu

\mavergleichskettedisp
{\vergleichskette
{ x
}
{ \neq} { \psi(x)
}
{ } {
}
{ } {
}
{ } {
}
}
{}{}{,}
sehingga kedua titik tersebut menentukan suatu garis. Gagasannya adalah menggunakan garis ini untuk mengambil salah satu
\zusatzklammer {dari kedua} {} {}
titik potong dengan permukaan bola sebagai titik citra suatu retraksi ke batas. Dengan fungsi bantu

\mavergleichskettedisp
{\vergleichskette
{ h(x)
}
{ =} { { \frac{   \psi(x)-x }{  \Vert { \psi (x)-x} \Vert  } }
}
{ } {
}
{ } {
}
{ } {
}
}
{}{}{}
kita mendefinisikan suatu pemetaan
\maabbdisp {\varphi} { B  \left(  0 , 1 \right) } { \R^n
} {}
melalui

\mavergleichskettedisphandlinks
{\vergleichskettedisphandlinks
{ \varphi(x)
}
{ =} { x + \left(  - \left\langle  x  , h(x)  \right\rangle + \sqrt{1 + \left\langle  x  , h(x)  \right\rangle^2 - \Vert { x} \Vert^2  }  \right) \cdot h(x)
}
{ } {
}
{ } {
}
{ } {
}
}
{}{}{.}
Ekspresi di bawah akar tersebut positif. Hal ini jelas untuk

\mavergleichskette
{\vergleichskette
{ \Vert { x } \Vert
}
{ < }{ 1
}
{  }{
}
{    }{
}
{   }{
}
}
{}{}{}
dan untuk

\mavergleichskette
{\vergleichskette
{ \Vert { x } \Vert
}
{ = }{ 1
}
{  }{
}
{    }{
}
{   }{
}
}
{}{}{}
diperoleh suatu titik pada permukaan bola yang garis penghubungnya dengan titik bola
\mathl{\psi(x)}{} tidak tegak lurus terhadap $x$
\zusatzklammer {ruang singgung afin pada suatu titik permukaan bola hanya berpotongan dengan bola di satu titik} {} {,}
sehingga

\mavergleichskettedisp
{\vergleichskette
{ \left\langle  x  , h(x)  \right\rangle
}
{ \neq} { 0
}
{ } {
}
{ } {
}
{ } {
}
}
{}{}{}
berlaku. Karena akar kuadrat dan nilai mutlak di luar titik nol
\definitionsverweis {terdiferensialkan secara kontinu}{}{,}
baik
\mathkor {} {h(x)} {maupun} {\varphi(x)} {}
merupakan pemetaan yang terdiferensialkan secara kontinu. Menurut
Aufgabe 23.23 (Differentialgeometrie (Osnabrück 2023))
pemetaan $\varphi$ memetakan bola ke permukaan bola dan pembatasannya pada permukaan bola adalah identitas. Dengan demikian, diperoleh suatu
\definitionsverweis {retraksi}{}{}
yang terdiferensialkan secara kontinu dari bola penuh tertutup ke batasnya, yang menurut
Satz 23.9 (Differentialgeometrie (Osnabrück 2023))
tidak mungkin ada.

 }



\inputaufgabeklausurloesung
{7 (1+2+4)}

{

Pada
\definitionsverweis {bundel vektor trivial}{}{}

\mathl{\R \times \R^2}{} di atas $\R$, kita meninjau
\definitionsverweis {koneksi trivial}{}{}
$\nabla$.
\aufzaehlungdrei{Buktikan bahwa

\mavergleichskette
{\vergleichskette
{ f_1(t)
}
{ = }{ \left(  1+ t^2  , \,  t^3                   \right)
}
{  }{
}
{    }{
}
{   }{
}
}
{}{}{}
dan

\mavergleichskette
{\vergleichskette
{ f_2(t)
}
{ = }{ \left(  t  , \,   1+t^2                   \right)
}
{  }{
}
{    }{
}
{   }{
}
}
{}{}{}
merupakan
\definitionsverweis {penampang basis}{}{}
dari bundel vektor tersebut.
}{Nyatakan penampang basis standar $e_1,e_2$ sebagai kombinasi linear dari penampang basis $f_1,f_2$.
}{Tentukan
\definitionsverweis {simbol-simbol Christoffel}{}{}
dari $\nabla$ terhadap penampang-penampang basis tersebut.
}

}
{

\aufzaehlungdrei{Berlaku

\mavergleichskettedisp
{\vergleichskette
{ \det  \begin{pmatrix}  1+t^2   &  t^3  \\  t  &  1+t^2  \end{pmatrix}
}
{ =} { 1+2t^2+t^4 -t^4
}
{ =} { 1 +2t^2
}
{ >} {  0
}
{ } {
}
}
{}{}{}
selalu positif. Oleh karena itu, kedua vektor
\mathkor {} {f_1(t)} {dan} {f_2(t)} {}
untuk setiap

\mavergleichskette
{\vergleichskette
{ t
}
{ \in }{  \R
}
{  }{
}
{    }{
}
{   }{
}
}
{}{}{}
membentuk suatu basis.
}{Dari hubungan

\mavergleichskettedisp
{\vergleichskette
{ \begin{pmatrix} f_1  \\f_2   \end{pmatrix}
}
{ =} { \begin{pmatrix}  1+t^2  &   t^3  \\  t  &  1+t^2 \end{pmatrix}  \begin{pmatrix} e_1  \\e_2   \end{pmatrix}
}
{ } {
}
{ } {
}
{ } {
}
}
{}{}{}
diperoleh

\mavergleichskettedisp
{\vergleichskette
{ \begin{pmatrix} e_1  \\e_2   \end{pmatrix}
}
{ =} { { \frac{   1  }{  1+2t^2 } } \begin{pmatrix}  1+t^2  &   -t^3  \\  -t &  1+t^2 \end{pmatrix}  \begin{pmatrix} f_1  \\f_2   \end{pmatrix}
}
{ } {
}
{ } {
}
{ } {
}
}
{}{}{,}
jadi

\mavergleichskettedisp
{\vergleichskette
{ e_1
}
{ =} { { \frac{   1+t^2 }{  1+2t^2 } }   f_1 - { \frac{  t^3 }{  1+2t^2 } } f_2
}
{ } {
}
{ } {
}
{ } {
}
}
{}{}{}
dan

\mavergleichskettedisp
{\vergleichskette
{ e_2
}
{ =} { -{ \frac{   t  }{  1+2t^2 } }   f_1 + { \frac{   1+t^2 }{  1+2t^2 } } f_2
}
{ } {
}
{ } {
}
{ } {
}
}
{}{}{.}
}{Untuk koneksi trivial, berlaku

\mavergleichskettedisp
{\vergleichskette
{ \nabla_\partial (f)
}
{ =} { \partial f
}
{ } {
}
{ } {
}
{ } {
}
}
{}{}{.}
Dengan demikian,

\mavergleichskettealign
{\vergleichskettealign
{ \nabla_\partial (f_1)
}
{ =} { \nabla_\partial \left(  \left(  1+t^2  \right) e_1 + t^3e_2  \right)
}
{ =} { 2t e_1 + 3t^2 e_2
}
{ =} { 2t  \left(  { \frac{   1+t^2 }{  1+2t^2 } }   f_1 - { \frac{  t^3 }{  1+2t^2 } } f_2   \right)  + 3t^2  \left(  -{ \frac{   t  }{  1+2t^2 } }   f_1 + { \frac{   1+t^2 }{  1+2t^2 } } f_2   \right)
}
{ =} {  \left(   { \frac{   2t+2t^3 }{  1+2t^2 } } - { \frac{   3t^3 }{  1+2t^2 } }  \right) f_1  +  \left(  - { \frac{  t^3 }{  1+2t^2 } }   + { \frac{   3t^2+3t^4 }{  1+2t^2 } }     \right)   f_2
}
}
{
\vergleichskettefortsetzungalign
{ =} { { \frac{   2t - t^3 }{  1+2t^2 } }  f_1  +  { \frac{   3t^2-t^3+3t^4 }{  1+2t^2 } }      f_2
}
{ } {}
{ } {}
{ } {}
}
{}{}
dan

\mavergleichskettealign
{\vergleichskettealign
{ \nabla_\partial (f_2)
}
{ =} { \nabla_\partial \left(  t e_1 + \left(  1+t^2  \right) e_2  \right)
}
{ =} { e_1 + 2t  e_2
}
{ =} { { \frac{   1+t^2 }{  1+2t^2 } } f_1 - { \frac{  t^3 }{  1+2t^2 } } f_2 + 2t \left(  -{ \frac{   t  }{  1+2t^2 } } f_1 + { \frac{   1+t^2 }{  1+2t^2 } } f_2   \right)
}
{ =} { \left( { \frac{   1+t^2 }{  1+2t^2 } }  - { \frac{   2t^2 }{  1+2t^2 } }  \right) f_1  + \left(  - { \frac{   t^3  }{  1+2t^2  } } +  { \frac{   2t+2t^3  }{  1+2t^2 } }  \right) f_2
}
}
{
\vergleichskettefortsetzungalign
{ =} { { \frac{   1 - t^2  }{  1+2t^2 } } f_1 + { \frac{   2t   +t^3 }{  1+2t^2 } }      f_2
}
{ } {}
{ } {}
{ } {}
}
{}{.}
Fungsi-fungsi koefisien tersebut adalah simbol-simbol Christoffel.
}

 }

\endgroup

\chapter{Formulir Solusi Resmi Ujian 9}
\noindent\textit{Solusi yang disediakan oleh sumber resmi; kekosongan sumber dipertahankan.}
\begingroup
\renewcommand{\theHfakt}{o011.exam.official.09.\arabic{fakt}}
%Data institusi

%\input{Dozentdaten}

%\renewcommand{\fachbereich}{Bidang studi}

%\renewcommand{\dozent}{Prof. Dr. . }

%Data ujian

\renewcommand{\klausurgebiet}{ }

\renewcommand{\klausurtyp}{ }

\renewcommand{\klausurdatum}{ . 20}

\klausurvorspann {\fachbereich} {\klausurdatum} {\dozent} {\klausurgebiet} {\klausurtyp}

%Data untuk tabel poin berikut

\renewcommand{\aeins}{  3  }

\renewcommand{\azwei}{  3   }

\renewcommand{\adrei}{  6  }

\renewcommand{\avier}{  3  }

\renewcommand{\afuenf}{  3  }

\renewcommand{\asechs}{  5  }

\renewcommand{\asieben}{  13  }

\renewcommand{\aacht}{  4  }

\renewcommand{\aneun}{  3  }

\renewcommand{\azehn}{  7  }

\renewcommand{\aelf}{  6  }

\renewcommand{\azwoelf}{  3  }

\renewcommand{\adreizehn}{  6   }

\renewcommand{\avierzehn}{  65  }

\renewcommand{\afuenfzehn}{    }

\renewcommand{\asechzehn}{    }

\renewcommand{\asiebzehn}{    }

\renewcommand{\aachtzehn}{    }

\renewcommand{\aneunzehn}{    }

\renewcommand{\azwanzig}{    }

\renewcommand{\aeinundzwanzig}{    }

\renewcommand{\azweiundzwanzig}{    }

\renewcommand{\adreiundzwanzig}{    }

\renewcommand{\avierundzwanzig}{    }

\renewcommand{\afuenfundzwanzig}{    }

\renewcommand{\asechsundzwanzig}{    }

\punktetabelledreizehn

\klausurnote

\newpage

\setcounter{section}{0}



\inputaufgabeklausurloesung
{3}

{

Definisikan istilah-istilah berikut
\zusatzklammer {yang dicetak miring} {} {.}
\aufzaehlungsechs{Suatu
\stichwort {ruang Hausdorff} {}
$X$.

}{Suatu
\stichwort {pemetaan diferensiabel} {}
\maabbdisp {\varphi} { L } { M
} {}
antara dua manifold diferensiabel
\mathkor {} {L} {dan} {M} {.}

}{
\stichwort {Bundel tangen} {}
dari suatu manifold diferensiabel $M$.

}{
\stichwort {Matriks fundamental kedua} {}
dari parametrisasi lokal
\maabbdisp {} {V} {U \subseteq Y
} {}
\zusatzklammer {dengan

\mavergleichskette
{\vergleichskette
{ V
}
{ \subseteq }{ \R^{n-1}
}
{  }{
}
{    }{
}
{   }{
}
}
{}{}{}
terbuka} {} {}
suatu hiperpermukaan diferensiabel

\mavergleichskette
{\vergleichskette
{ Y
}
{ \subseteq }{ \R^n
}
{  }{
}
{    }{
}
{   }{
}
}
{}{}{.}

}{
\stichwort {Turunan eksterior} {} dari suatu bentuk diferensial terdiferensialkan secara kontinu
\mathl{\omega \in
{ \mathcal E }^{ k } (  M   )}{} pada manifold diferensiabel $M$.

}{Suatu \stichwort {partisi kesatuan yang tersubordinasi pada penutup} {,} dengan

\mavergleichskette
{\vergleichskette
{  X
}
{ = }{ \bigcup_{i \in I} W_i
}
{  }{
}
{    }{
}
{   }{
}
}
{}{}{}
suatu
\definitionsverweis {penutup terbuka}{}{}
dari
\definitionsverweis {ruang topologis}{}{}
$X$.
}

}
{

\aufzaehlungsechs{Suatu
\definitionsverweis {ruang topologis}{}{}
$X$ disebut Hausdorff jika untuk setiap dua titik berbeda
\mathl{x, y \in X}{} terdapat dua
\definitionsverweis {himpunan terbuka}{}{}
\mathkor {} {U} {dan} {V} {}
yang memenuhi
\mathl{x \in U, \, y \in V}{} dan
\mathl{U \cap V = \emptyset}{.}
}{Misalkan
\mathkor {} {(U_i,U_i', \alpha_i, i \in I)} {dan} {(V_j,V_j', \beta_j, j \in J)} {}
adalah
\definitionsverweis {atlas}{}{}
dari
\mathkor {} {M} {dan} {N} {.}
Pemetaan
\maabbdisp {\varphi} { L } { M
} {}
disebut diferensiabel jika kontinu dan jika untuk setiap
\mathkor {} {i \in I} {dan setiap} {j \in J} {}
pemetaan-pemetaan
\maabbdisp {\beta_j \circ  \varphi \circ (\alpha_i)^{-1}} {\alpha_i( \varphi^{-1}(V_j) \cap U_i) } { V_j'
} {}
\definitionsverweis {terdiferensialkan secara kontinu}{}{.}
}{Himpunan

\mavergleichskettedisp
{\vergleichskette
{TM
}
{ =} {\biguplus_{P \in M} T_PM
}
{ } {
}
{ } {
}
{ } {
}
}
{}{}{,}
yang dilengkapi dengan pemetaan proyeksi
\maabbeledisp {\pi} {TM} {M
} {(P,v)} {P
} {,}
disebut bundel tangen dari $M$.
}{Misalkan $Y$
\definitionsverweis {diorientasikan}{}{}
oleh suatu
\definitionsverweis {medan normal satuan}{}{}
$N$, dengan $N$ dipandang sebagai medan pada $V$. Definisikan

\mavergleichskettedisp
{\vergleichskette
{ h_{i j}
}
{ =} {  \left\langle  \partial_i \partial_j \varphi , N  \right\rangle
}
{ } {
}
{ } {
}
{ } {
}
}
{}{}{.}
Matriks fundamental kedua
dari $\varphi$ adalah matriks
\zusatzklammer {yang bergantung pada

\mavergleichskettek
{\vergleichskettek
{ Q
}
{ \in }{ V
}
{  }{
}
{    }{
}
{   }{
}
}
{}{}{}} {} {}

\mavergleichskettedisp
{\vergleichskette
{ H
}
{ =} {  \left(  h_{ij}   \right)_{1 \leq i,j \leq n-1}
}
{ } {
}
{ } {
}
{ } {
}
}
{}{}{.}
}{Turunan eksterior dari $\omega$, secara lokal pada suatu peta tempat $\omega$ berbentuk

\mavergleichskettedisp
{\vergleichskette
{ \omega
}
{ =} {\sum_{ { \# \left(   I  \right) } = k , \, I \subseteq \{1 , \ldots ,  \operatorname{dim}\, M \}  }f_I dx_I
}
{ } {
}
{ } {
}
{ } {
}
}
{}{}{}
didefinisikan oleh

\mavergleichskettedisp
{\vergleichskette
{  d \omega
}
{ =} {\sum_{ { \# \left(   I  \right) } = k   } d(f_I ) \wedge dx_I
}
{ } {
}
{ } {
}
{ } {
}
}
{}{}{}
rumus di atas.
}{Suatu keluarga fungsi
\maabbdisp {h_j} {X} {\R
} {}
dengan

\mavergleichskette
{\vergleichskette
{ j
}
{ \in }{ J
}
{  }{
}
{    }{
}
{   }{
}
}
{}{}{}
disebut partisi kesatuan yang tersubordinasi pada penutup jika sifat-sifat berikut berlaku.
\aufzaehlungvier{Berlaku

\mavergleichskette
{\vergleichskette
{ h_j(X)
}
{ \subseteq }{ [0,1]
}
{  }{
}
{    }{
}
{   }{
}
}
{}{}{}
untuk setiap

\mavergleichskette
{\vergleichskette
{ j
}
{ \in }{ J
}
{  }{
}
{    }{
}
{   }{
}
}
{}{}{.}
}{Setiap titik

\mavergleichskette
{\vergleichskette
{ P
}
{ \in }{ X
}
{  }{
}
{    }{
}
{   }{
}
}
{}{}{}
memiliki suatu
\definitionsverweis {lingkungan terbuka}{}{}

\mavergleichskette
{\vergleichskette
{ P
}
{ \in }{  U
}
{  }{
}
{    }{
}
{   }{
}
}
{}{}{}
sedemikian sehingga
\definitionsverweis {fungsi-fungsi yang dibatasi}{}{}

\mathl{h_j  \vert_{ U }}{} merupakan fungsi nol, kecuali untuk sejumlah berhingga indeks.
}{Berlaku

\mavergleichskette
{\vergleichskette
{  \sum_{j \in J} h_j
}
{ = }{ 1
}
{  }{
}
{    }{
}
{   }{
}
}
{}{}{.}
}{Untuk setiap

\mavergleichskette
{\vergleichskette
{ j
}
{ \in }{ J
}
{  }{
}
{    }{
}
{   }{
}
}
{}{}{}
terdapat suatu himpunan terbuka
\mathl{W_{i(j)}}{} dari penutup sedemikian sehingga
\definitionsverweis {tumpuan}{}{}
dari $h_j$ termuat dalam
\mathl{W_{i(j)}}{.}
}
}

 }



\inputaufgabeklausurloesung
{3}

{

Nyatakan teorema-teorema berikut.
\aufzaehlungdrei{
\stichwort {Teorema tentang sifat aljabar pemetaan Weingarten} {.}}{Teorema tentang keterorientasian dan bentuk volume positif.}{
\stichwort {Teorema titik tetap Brouwer} {.}}

}
{

\aufzaehlungdrei{\faktsituation {Misalkan

\mavergleichskette
{\vergleichskette
{ W
}
{ \subseteq }{ \R^n
}
{  }{
}
{    }{
}
{   }{
}
}
{}{}{}
\definitionsverweis {terbuka}{}{,}
\maabb {h} { W } { \R
} {}
adalah suatu
\definitionsverweis {fungsi yang dua kali terdiferensialkan secara kontinu}{}{}
dan

\mavergleichskette
{\vergleichskette
{ Y
}
{ = }{ h^{-1}(c)
}
{  }{
}
{    }{
}
{   }{
}
}
{}{}{}
adalah serat di atas

\mavergleichskette
{\vergleichskette
{ c
}
{ \in }{ \R
}
{  }{
}
{    }{
}
{   }{
}
}
{}{}{,}
dengan $h$
\definitionsverweis {reguler}{}{}
di setiap titik $Y$.}
\faktfolgerung {Maka
\definitionsverweis {pemetaan Weingarten}{}{}
$L_P$ di setiap titik

\mavergleichskette
{\vergleichskette
{ P
}
{ \in }{ Y
}
{  }{
}
{    }{
}
{   }{
}
}
{}{}{}
bersifat
\definitionsverweis {adjoin-diri}{}{.}}
\faktzusatz {}
\faktzusatz {}}{\faktsituation {Misalkan $M$ adalah suatu
\definitionsverweis {manifold diferensiabel
}{}{}
dengan suatu
\definitionsverweis {basis topologi terhitung}{}{.}}
\faktfolgerung {Terdapat suatu
\definitionsverweis {bentuk volume}{}{}
\definitionsverweis {kontinu}{}{}
tanpa titik nol pada $M$ jika dan hanya jika $M$
\definitionsverweis {dapat diorientasikan}{}{.}}
\faktzusatz {Bentuk volume ini juga dapat dipilih terdiferensialkan secara kontinu dan positif.}
\faktzusatz {}}{Misalkan
\maabbdisp {\psi} { B  \left( 0,r \right) } { B  \left( 0,r \right)
} {}
adalah suatu pemetaan yang terdiferensialkan secara kontinu dari bola tertutup di $\R^n$ ke dirinya sendiri. Maka $\psi$ memiliki suatu titik tetap.}

 }



\inputaufgabeklausurloesung
{6}

{

Tentukan
\definitionsverweis {kelengkungan}{}{}
di setiap titik lingkaran satuan yang diberikan secara implisit oleh

\mavergleichskettedisp
{\vergleichskette
{ h(x,y)
}
{ =} { x^2+y^2
}
{ =} { 1
}
{ } {
}
{ } {
}
}
{}{}{}
dengan
Lemma 3.11 (Geometri Diferensial (Osnabrück 2023))
dan medan gradien $h$.

}
{

\definitionsverweis {Medan gradien}{}{}
dari $h$ adalah $\begin{pmatrix}  2x \\ 2y   \end{pmatrix}$, sehingga suatu medan normal satuan adalah

\mavergleichskettedisp
{\vergleichskette
{ N \begin{pmatrix}  x  \\ y   \end{pmatrix}
}
{ =} { { \frac{   1 }{  \sqrt{x^2+y^2}  } } \begin{pmatrix}  x  \\ y   \end{pmatrix}
}
{ } {
}
{ } {
}
{ } {
}
}
{}{}{.}
\definitionsverweis {Diferensial total}{}{}
dari $N$ di suatu titik

\mavergleichskettedisp
{\vergleichskette
{ P
}
{ =} { \begin{pmatrix}  x  \\ y   \end{pmatrix}
}
{ } {
}
{ } {
}
{ } {
}
}
{}{}{}
adalah

\mavergleichskettedisp
{\vergleichskette
{  \left(D N \right)_{P}
}
{ =} { \begin{pmatrix}  { \frac{   y^2 }{   \left(  x^2+y^2  \right)^{3/2}  } }   &   - { \frac{   xy }{   \left(  x^2+y^2  \right)^{3/2}  } }   \\    - { \frac{   xy }{   \left(  x^2+y^2  \right)^{3/2}  } }    &  { \frac{   x^2 }{   \left(  x^2+y^2  \right)^{3/2}  } }  \end{pmatrix}
}
{ } {
}
{ } {
}
{ } {
}
}
{}{}{.}
Jika diterapkan pada vektor singgung
\mathl{\begin{pmatrix}  y  \\ -x   \end{pmatrix}}{}, diperoleh

\mavergleichskettealign
{\vergleichskettealign
{ \left(D N \right)_{P} \begin{pmatrix}  y  \\ -x   \end{pmatrix}
}
{ =} { \begin{pmatrix}  y^2    &   - xy    \\  - xy   &  x^2  \end{pmatrix} \begin{pmatrix}  y  \\ -x   \end{pmatrix}
}
{ =} { \begin{pmatrix}  y \left(  y^2 +x^2  \right)  \\ -x \left(  y^2 +x^2  \right)    \end{pmatrix}
}
{ =} { \begin{pmatrix}  y   \\ -x    \end{pmatrix}
}
{ } {
}
}
{}
{}{.}
Oleh karena itu, menurut
Lemma 3.11 (Differentialgeometrie (Osnabrück 2023))
kelengkungan kurva di $P$ sama dengan $-1$.
\zusatzklammer {Hasil ini sesuai karena
\mathl{\begin{pmatrix}  y  \\ -x   \end{pmatrix}}{} dan medan gradien merepresentasikan orientasi standar, sedangkan
\mathl{\begin{pmatrix}  y  \\ -x   \end{pmatrix}}{} adalah arah singgung untuk gerak searah jarum jam.} {} {}

 }



\inputaufgabeklausurloesung
{3}

{

Misalkan $X$ suatu
\definitionsverweis {ruang Hausdorff}{}{.}
Tunjukkan bahwa
\definitionsverweis {diagonal}{}{}

\mavergleichskettedisp
{\vergleichskette
{ \triangle
}
{ =} { \left\{ (x,y) \in X \times X  \mid   x = y        \right\}
}
{ } {
}
{ } {
}
{ } {
}
}
{}{}{}
merupakan
\definitionsverweis {subhimpunan tertutup}{}{}
dalam
\definitionsverweis {ruang produk}{}{}

\mathl{X \times X}{}.

}
{

Kita harus menunjukkan bahwa komplemen

\mavergleichskettedisp
{\vergleichskette
{W
}
{ =} { X \times X \setminus \triangle
}
{ =} { \left\{  (x,y)   \mid   x,y \in X , \,  x \neq y       \right\}
}
{ } {
}
{ } {
}
}
{}{}{}
terbuka. Jadi, misalkan
\mathl{(x,y)}{} adalah suatu pasangan dengan

\mavergleichskette
{\vergleichskette
{ x
}
{ \neq }{ y
}
{  }{
}
{    }{
}
{   }{
}
}
{}{}{.}
Berdasarkan sifat Hausdorff yang diasumsikan, terdapat himpunan-himpunan terbuka saling lepas
\mathl{U,V \subseteq X}{} dengan
\mathl{x \in U}{} dan
\mathl{y \in V}{.} Berlaku
\mathl{(x,y) \in U \times V}{,} dan menurut definisi topologi produk,
\mathl{U \times V}{} adalah suatu himpunan terbuka dalam
\mathl{X \times X}{.} Karena keduanya saling lepas, dari
\mathl{(z,w) \in U\times V}{} langsung diperoleh

\mavergleichskette
{\vergleichskette
{z
}
{ \neq }{w
}
{  }{
}
{    }{
}
{   }{
}
}
{}{}{.}
Jadi,

\mavergleichskettedisp
{\vergleichskette
{(x,y)
}
{ \in} { U \times V
}
{ \subseteq} {  W
}
{ } {
}
{ } {
}
}
{}{}{}
dan $W$ merupakan gabungan himpunan-himpunan produk terbuka seperti itu, sehingga himpunan tersebut sendiri terbuka.

 }



\inputaufgabeklausurloesung
{3}

{

Tunjukkan bahwa ketiga manifold satu dimensi
\mathlistdisp {S^1 = \left\{ (x,y) \in \R^2  \mid   \betrag { (x,y) } = 1         \right\}} {} {\R} {dan} {\R \setminus \{0\}} {}
tidak ada dua di antaranya yang homeomorfik.

}
{

Sebagai himpunan bagian tertutup dan terbatas dari $\R^2$, lingkaran satuan bersifat kompak
\zusatzklammer {Satz von Heine-Borel} {} {.} Sebaliknya, garis real $\R$ dan $\R \setminus \{0\}$ tidak kompak karena keduanya merupakan himpunan bagian tak terbatas dari $\R$. Jadi,
\mathl{S^1 \not \cong \R,\, \R \setminus \{0\}}{.}

Garis real terhubung, sebagaimana mengikuti dari
Zwischenwertsatz
. Sebaliknya, $\R \setminus \{0\}$ tidak terhubung karena dapat ditulis

\mavergleichskettedisp
{\vergleichskette
{\R_+
}
{ =} {(\R \setminus \{0\}) \cap [0, \infty]
}
{ =} {(\R \setminus \{0\}) \cap \, ]0, \infty]
}
{ } {
}
{ } {
}
}
{}{}{}
sehingga diperoleh suatu himpunan bagian tak kosong yang sekaligus terbuka dan tertutup. Jadi,
\mathl{\R \,\not \cong \R \setminus \{0\}}{.}

 }



\inputaufgabeklausurloesung
{5}

{

Misalkan

\mavergleichskette
{\vergleichskette
{ G
}
{ \subseteq }{ \R^m
}
{  }{
}
{    }{
}
{   }{
}
}
{}{}{}
\definitionsverweis {terbuka}{}{}
dan misalkan

\mavergleichskette
{\vergleichskette
{ M
}
{ \subseteq }{ G
}
{  }{
}
{    }{
}
{   }{
}
}
{}{}{}
suatu
$n$-\definitionsverweis {dimensional}{}{}
\definitionsverweis {submanifold tertutup}{}{,}
yang
\definitionsverweis {berorientasi}{}{}
dan dilengkapi dengan struktur Riemann terinduksi serta
\definitionsverweis {bentuk volume kanonik}{}{}
$\omega$. Misalkan

\mavergleichskette
{\vergleichskette
{ W
}
{ \subseteq }{ \R^n
}
{  }{
}
{    }{
}
{   }{
}
}
{}{}{}
terbuka dan misalkan
\maabbdisp {\varphi} { W } { U
} {}
suatu
\definitionsverweis {difeomorfisme}{}{}
dengan himpunan terbuka

\mavergleichskette
{\vergleichskette
{ U
}
{ \subseteq }{ M
}
{  }{
}
{    }{
}
{   }{
}
}
{}{}.
Tunjukkan bahwa pada $W$ berlaku rumus

\mavergleichskettealignhandlinks
{\vergleichskettealignhandlinks
{\varphi^* \left(  \omega \vert_U  \right)
}
{ =} { \left( \det  \left( \left\langle  \begin{pmatrix}  { \frac{   \partial  \varphi_1   }{  \partial   x_i   } }  \\ \vdots \\  { \frac{   \partial  \varphi_m   }{  \partial   x_i   } }   \end{pmatrix}  ,  \begin{pmatrix}  { \frac{   \partial  \varphi_1   }{  \partial   x_j   } }  \\ \vdots \\  { \frac{   \partial  \varphi_m   }{  \partial   x_j   } }   \end{pmatrix}   \right\rangle \right)_{1 \leq i,j \leq n} \right)^{1/2}  dx_1 \wedge \ldots \wedge dx_n
}
{ =} { \left( \det  \left(  { \frac{   \partial  \varphi_1   }{  \partial   x_i   } } { \frac{   \partial  \varphi_1   }{  \partial   x_j   } }  + \cdots +  { \frac{   \partial  \varphi_m   }{  \partial   x_i   } } { \frac{   \partial  \varphi_m   }{  \partial   x_j   } }  \right)_{1 \leq i,j \leq n} \right)^{1/2}  dx_1 \wedge \ldots \wedge dx_n
}
{ } {
}
{ } {
}
}
{}
{}{}{.}

}
{

Persamaan kedua langsung diperoleh dengan mengevaluasi hasil kali dalam standar pada $\R^m$. Menurut definisi
\definitionsverweis {matriks fundamental metrik}{}{,}
untuk

\mavergleichskette
{\vergleichskette
{ Q
}
{ \in }{ W
}
{  }{
}
{    }{
}
{   }{
}
}
{}{}{}

\mavergleichskettealign
{\vergleichskettealign
{ g_{ij} (Q)
}
{ =} { \left\langle  T_Q(\varphi)(e_i)  ,  T_Q(\varphi)(e_j)   \right\rangle_{\varphi(Q)}
}
{ =} { \left\langle  T_Q(\varphi)(e_i)  ,  T_Q(\varphi)(e_j)   \right\rangle
}
{ =} { \left\langle  \left(  D \varphi  \right)_{ Q } \left(   e_i   \right)  ,  \left(  D\varphi  \right)_{ Q } \left(   e_j   \right)   \right\rangle
}
{ =} { \left\langle  \begin{pmatrix}  { \frac{   \partial  \varphi_1   }{  \partial   x_i   } } ( Q)  \\ \vdots \\  { \frac{   \partial  \varphi_m   }{  \partial   x_i   } } ( Q)   \end{pmatrix}  ,  \begin{pmatrix}  { \frac{   \partial   \varphi_1   }{  \partial   x_j   } } ( Q )  \\ \vdots \\  { \frac{   \partial  \varphi_m   }{  \partial   x_j   } } ( Q)   \end{pmatrix}   \right\rangle
}
}
{}
{}{,}
karena ruang singgung
\mathl{T_{\varphi(Q)}M}{} dilengkapi dengan hasil kali dalam yang diinduksi dari $\R^m$, pemetaan singgung dalam kasus lokal adalah diferensial total, dan entri-entrinya dapat dinyatakan dengan turunan parsial. Oleh karena itu, pernyataan tersebut mengikuti dari
Lemma 16.7 (Differentialgeometrie (Osnabrück 2023)).

 }



\inputaufgabeklausurloesung
{weiter}

{

Kita tinjau

\mavergleichskette
{\vergleichskette
{ \R^6
}
{ = }{ \left(  \R^2  \right)^3
}
{  }{
}
{    }{
}
{   }{
}
}
{}{}{}
sebagai himpunan semua segitiga
\zusatzklammer {termasuk yang degenerat} {} {}
dengan mengidentifikasi segitiga yang mempunyai titik-titik sudut
\zusatzklammer {terurut} {} {}

\mavergleichskette
{\vergleichskette
{ P_1
}
{ = }{ \begin{pmatrix}  x_1  \\ y_1    \end{pmatrix}
}
{  }{
}
{    }{
}
{   }{
}
}
{}{}{,}

\mavergleichskette
{\vergleichskette
{ P_2
}
{ = }{ \begin{pmatrix}  x_2  \\ y_2    \end{pmatrix}
}
{  }{
}
{    }{
}
{   }{
}
}
{}{}{}
dan

\mavergleichskette
{\vergleichskette
{ P_3
}
{ = }{ \begin{pmatrix}  x_3  \\ y_3    \end{pmatrix}
}
{  }{
}
{    }{
}
{   }{
}
}
{}{}{,}
dengan tupel koordinat

\mathdisp {(x_1,y_1,x_2,y_2,x_3,y_3)} {  }

.
\aufzaehlungsechs{Tunjukkan bahwa himpunan segitiga dengan dua titik sudut berimpit merupakan subhimpunan tertutup $\R^6$
\zusatzklammer {jadi komplemennya merupakan himpunan terbuka dalam $\R^6$ yang kita namai $U$} {} {.}
}{Tunjukkan bahwa himpunan segitiga dengan ketiga titik sudut segaris merupakan subhimpunan tertutup $\R^6$
\zusatzklammer {jadi komplemennya, yang terdiri atas semua segitiga tak degenerat, merupakan himpunan terbuka dalam $\R^6$ yang kita namai $V$} {} {.}
}{Buat aturan fungsi yang mendeskripsikan pemetaan
\maabbdisp {F} {\R^6} {\R
} {}
yang memetakan suatu segitiga ke kelilingnya.
}{Tunjukkan bahwa fungsi $F$ dari bagian (3) terdiferensialkan secara kontinu pada himpunan $U$.
}{Hitung turunan parsial $F$ terhadap $x_1$ pada $U$.
}{Tunjukkan bahwa himpunan segitiga tak degenerat berkeliling $1$ membentuk
\definitionsverweis {submanifold tertutup}{}{}
dari $V$. Berapakah dimensinya?
}

}
{

\aufzaehlungsechs{Titik-titik sudut
\mathkor {} {P_1} {dan} {P_2} {}
berimpit jika dan hanya jika kedua koordinatnya sama, yakni jika

\mavergleichskette
{\vergleichskette
{ x_1
}
{ = }{ x_2
}
{  }{
}
{    }{
}
{   }{
}
}
{}{}{}
dan

\mavergleichskette
{\vergleichskette
{ y_1
}
{ = }{ y_2
}
{  }{
}
{    }{
}
{   }{
}
}
{}{}{}
berlaku. Syarat seperti ini mendefinisikan suatu himpunan bagian tertutup karena merupakan serat dari suatu pemetaan linear
\zusatzklammer {kontinu} {} {.}
Karena gabungan berhingga dari himpunan-himpunan bagian tertutup kembali tertutup, himpunan yang dimaksud bersifat tertutup.
}{Tinjau vektor-vektor penghubung

\mathdisp {P_2-P_1 = \begin{pmatrix}  x_2-x_1  \\ y_2-y_1   \end{pmatrix} \text{   und } P_3-P_1 = \begin{pmatrix}  x_3-x_1  \\ y_3-y_1   \end{pmatrix}} {  }

Ketiga titik kolinear jika dan hanya jika kedua vektor tersebut bergantung linear, dan hal ini berlaku jika dan hanya jika determinannya sama dengan $0$. Dengan demikian, syarat kolinearitasnya adalah

\mavergleichskettedisp
{\vergleichskette
{ H(x_1,x_2,x_3,y_1,y_2,y_3)
}
{ =} { (x_2-x_1) (y_3-y_1) - (x_3-x_1) (y_2-y_1)
}
{ =} { 0
}
{ } {
}
{ } {
}
}
{}{}{.}
Karena $H$ adalah pemetaan polinomial dan karenanya kontinu, serat di atas $0$ merupakan himpunan bagian tertutup.
}{Karena keliling merupakan jumlah dari ketiga sisi segitiga, berlaku

\mavergleichskettealigndrucklinks
{\vergleichskettealigndrucklinks
{ F(x_1,x_2,x_3,y_1,y_2,y_3)
}
{ =} { d(P_1,P_2) + d(P_1,P_3) +d(P_2,P_3)
}
{ =} { \sqrt{(x_2-x_1)^2 + (y_2-y_1)^2} +  \sqrt{(x_3-x_1)^2 + (y_3-y_1)^2} +  \sqrt{(x_3-x_2)^2 + (y_3-y_2)^2}
}
{ } {
}
{ } {
}
}
{}{}{.}
}{Untuk suatu titik di
\mathl{U}{,} setiap radikan, sebagai fungsi polinomial, terdiferensialkan secara kontinu dan bernilai positif ketat. Oleh karena itu, akar kuadratnya juga terdiferensialkan secara kontinu.
}{Turunan parsialnya adalah

\mavergleichskettealigndrucklinks
{\vergleichskettealigndrucklinks
{ { \frac{   \partial   F   }{  \partial   x_1  } }
}
{ =} { { \frac{   1 }{  2 } }   \left(  (  x_2-x_1  )^2 + (y_2-y_1)^2  \right)^{-1/2} (2x_1-2x_2)    + { \frac{   1 }{  2 } }   \left(  (  x_3-x_1  )^2 + (y_3-y_1)^2  \right)^{-1/2} (2x_1-2x_3)
}
{ =} {  \left(  (  x_2-x_1  )^2 + (y_2-y_1)^2  \right)^{-1/2} (x_1-x_2)    +   \left(  (  x_3-x_1  )^2 + (y_3-y_1)^2  \right)^{-1/2} (x_1-x_3)
}
{ } {
}
{ } {
}
}
{}{}{.}
}{Pada $U$, semua radikan bernilai positif, sehingga turunan parsial terhadap $x_1$ hanya dapat bernilai nol jika
\mathl{x_1=x_2=x_3}{.} Namun, segitiga yang memenuhi syarat ini bersifat kolinear. Ini berarti bahwa pada $V$ turunan parsial terhadap $x_1$ tidak pernah bernilai $0$; khususnya, gradien pada $V$ tidak pernah lenyap, sehingga bersifat reguler. Menurut
Satz 8.2 (Differentialgeometrie (Osnabrück 2023))
serat tersebut merupakan submanifold tertutup, dengan dimensi

\mavergleichskettedisp
{\vergleichskette
{ 6-1
}
{ =} { 5
}
{ } {
}
{ } {
}
{ } {
}
}
{}{}{.}
}

 }



\inputaufgabeklausurloesung
{4}

{

Misalkan
\mathkor {} {r} {dan} {R} {}
bilangan-bilangan real positif dengan

\mavergleichskette
{\vergleichskette
{ 0
}
{ < }{ r
}
{ < }{ R
}
{    }{
}
{   }{
}
}
{}{}{,}
dan tinjau torus
\zusatzklammer {tertanam} {} {}

\mavergleichskettedisp
{\vergleichskette
{ T
}
{ =} { \left\{ (x,y,z) \in \R^3  \mid   \left(  \sqrt{x^2+y^2} -R  \right)^2 +z^2  = r^2         \right\}
}
{ \subseteq} { \R^3
}
{ } {
}
{ } {
}
}
{}{}{}
dengan
\definitionsverweis {struktur Riemann terinduksi}{}{.}
Tunjukkan bahwa untuk setiap

\mavergleichskette
{\vergleichskette
{ \alpha
}
{ \in }{ [0,2 \pi [
}
{  }{
}
{    }{
}
{   }{
}
}
{}{}{}
pemetaan
\maabbeledisp {\Psi_\alpha} {\R^3} {\R^3
} { \left(  x   , \,   y  , \,  z                 \right) } { \left(  x \cos \alpha - y \sin  \alpha   , \,   x \sin \alpha + y \cos  \alpha  , \,  z                 \right)
} {,}
menginduksi suatu
\definitionsverweis {isometri}{}{}
dari $T$ ke dirinya sendiri.

}
{

Pada

\mavergleichskettedisp
{\vergleichskette
{ \R^3
}
{ =} { \R^2 \times \R
}
{ } {
}
{ } {
}
{ } {
}
}
{}{}{}
pemetaan tersebut merupakan produk suatu
\definitionsverweis {rotasi bidang}{}{}
dengan identitas. Oleh karena itu, pemetaan tersebut adalah
\definitionsverweis {isometri}{}{}
Euklides pada $\R^3$. Misalkan

\mavergleichskette
{\vergleichskette
{ (x,y,z)
}
{ \in }{  T
}
{  }{
}
{    }{
}
{   }{
}
}
{}{}{}
dan

\mavergleichskettedisp
{\vergleichskette
{ \Psi_\alpha(x,y,z)
}
{ =} { ( u,v,z)
}
{ } {
}
{ } {
}
{ } {
}
}
{}{}{.}
Maka

\mavergleichskettealigndrucklinks
{\vergleichskettealigndrucklinks
{ \left(  \sqrt{u^2+v^2} -R  \right)^2 +z^2
}
{ =} { \left(  \sqrt{ \left(  x \cos \alpha - y \sin  \alpha   \right)^2+ \left(  x \sin \alpha + y \cos  \alpha  \right)^2} -R  \right)^2 +z^2
}
{ =} { \left(  \sqrt{  x^2 \cos^{ 2  } \alpha + y^2 \sin^{ 2  }  \alpha  -2xy \cos \alpha \sin  \alpha    + x^2 \sin^{ 2  } \alpha + y^2 \cos^{ 2  }  \alpha  +2xy \cos \alpha \sin  \alpha     } -R  \right)^2 +z^2
}
{ =} { \left(  \sqrt{  x^2   + y^2 } -R  \right)^2 +z^2
}
{ =} { r^2
}
}
{}{}{,}
sehingga

\mavergleichskette
{\vergleichskette
{ \Psi_\alpha(x,y,z)
}
{ \in }{  T
}
{  }{
}
{    }{
}
{   }{
}
}
{}{}{.}
Jadi,
\maabbdisp {\Psi_\alpha} { T} {T
} {}
adalah suatu difeomorfisme dan isometri karena $T$ dilengkapi dengan struktur Riemann terinduksi.

 }



\inputaufgabeklausurloesung
{3}

{

Deskripsikan suatu model untuk bidang hiperbolik.

}
{  }



\inputaufgabeklausurloesung
{7}

{

Buktikan Theorema egregium.

}
{

Kita mendiferensialkan persamaan penentu pertama untuk simbol-simbol Christoffel, yaitu

\mavergleichskettedisp
{\vergleichskette
{ \partial_1 \partial_1 \varphi
}
{ =} { \Gamma^1_{11}  \partial_1 \varphi  + \Gamma^2_{11}  \partial_2 \varphi +h_{11} N
}
{ } {
}
{ } {
}
{ } {
}
}
{}{}{,}
terhadap variabel kedua, serta persamaan penentu kedua, yaitu

\mavergleichskettedisp
{\vergleichskette
{ \partial_1 \partial_2 \varphi
}
{ =} { \Gamma^1_{12}  \partial_1 \varphi + \Gamma^2_{12}  \partial_2  \varphi + h_{12} N
}
{ } {
}
{ } {
}
{ } {
}
}
{}{}{,}
terhadap variabel pertama. Berdasarkan teorema Schwarz, diperoleh

\mavergleichskettealigndrucklinks
{\vergleichskettealigndrucklinks
{ \left(  \partial_2 \Gamma^1_{11}  \right) \partial_1 \varphi +  \left(  \partial_2 \Gamma^2_{11}  \right)  \partial_2 \varphi+ \left(  \partial_2 h_{11}  \right)  N + \Gamma^1_{11} \partial_2 \partial_1 \varphi + \Gamma^2_{11}  \partial_2 \partial_2 \varphi  +h_{11} \partial_2 N
}
{ =} { \left(  \partial_1 \Gamma^1_{12}  \right)  \partial_1  \varphi+ \left(  \partial_1 \Gamma^2_{12}  \right)\partial_2 \varphi + \left(  \partial_1 h_{12}  \right)  N+ \Gamma^1_{12}  \partial_1 \partial_1 \varphi + \Gamma^2_{12}  \partial_1 \partial_2 \varphi + h_{12} \partial_1 N
}
{ } {
}
{ } {
}
{ } {
}
}
{}{}{.}
Selisih kedua ekspresi ini adalah $0$, dan kita menentukan bagian yang berkaitan dengan vektor basis $\partial_2 \varphi$. Untuk itu, kita menggunakan persamaan-persamaan penentu bagi simbol Christoffel dan
Lemma 19.3 (Differentialgeometrie (Osnabrück 2023))  (3)
, lalu memperoleh

\mavergleichskettedisphandlinks
{\vergleichskettedisphandlinks
{ 0
}
{ =} { \partial_2 \Gamma^2_{11} - \partial_1 \Gamma^2_{12} +\Gamma^1_{11} \Gamma^2_{12} - \Gamma^1_{12} \Gamma^2_{11} + \Gamma_{11}^2 \Gamma_{22}^2- \left(  \Gamma_{12}^2  \right)^2 + h_{11}  { \frac{   g_{12} h_{12} - g_{11} h_{22}   }{  g_{11}g_{22} -g_{12}^2 } } - h_{12}  { \frac{   g_{12} h_{11} - g_{11}h_{12}  }{  g_{11}g_{22} -g_{12}^2 } }
}
{ } {
}
{ } {
}
{ } {
}
}
{}{}{.}
Dengan

\mavergleichskettedisphandlinks
{\vergleichskettedisphandlinks
{ - g_{11} \left(  h_{11}h_{22} -h_{12}^2  \right)
}
{ =} { h_{11} \left(  g_{12} h_{12}- g_{11} h_{22}  \right) - h_{12} \left(  g_{12} h_{11} - g_{11}h_{12}  \right)
}
{ } {
}
{ } {
}
{ } {
}
}
{}{}{}
kita dapat mengganti kedua suku terakhir. Selanjutnya, menurut
Lemma 19.3 (Differentialgeometrie (Osnabrück 2023))  (4)
, diperoleh

\mavergleichskettealignhandlinks
{\vergleichskettealignhandlinks
{ g_{11} K
}
{ =} { g_{11} { \frac{   h_{11}h_{22} -h_{12}^2  }{ g_{11}g_{22} -g_{12}^2 } }
}
{ =} { \left(  \partial_2 \Gamma_{11}^2 - \partial_1 \Gamma_{12}^2 + \Gamma^1_{11} \Gamma^2_{12}  + \Gamma_{11}^2 \Gamma_{22}^2 - \Gamma_{12}^1 \Gamma_{11}^2-  \left(  \Gamma_{12}^2  \right)^2  \right)
}
{ } {
}
{ } {
}
}
{}
{}{.}

 }



\inputaufgabeklausurloesung
{6 (1+2+2+1)}

{

Kita tinjau
\definitionsverweis {bentuk diferensial}{}{}

\mavergleichskettedisp
{\vergleichskette
{ \omega
}
{ =} { ydx+zdy +xdz
}
{ } {
}
{ } {
}
{ } {
}
}
{}{}{}
pada $\R^3$ dan pemetaan
\maabbeledisp {\varphi} {\R^2} {\R^3
} { \left(  u   , \,  v                   \right) } { \left( u^2  , \,  v^2  , \,  uv                 \right)
} {.}
\aufzaehlungvier{Hitung turunan eksterior $\omega$.
}{Hitung tarikan balik $\varphi^* \omega$ dari $\omega$ di bawah $\varphi$.
}{Hitung turunan eksterior $\varphi^* \omega$ pada $\R^2$.
}{Hitung tarikan balik $\varphi^* d \omega$ dari $d\omega$ di bawah $\varphi$ secara independen dari bagian (3).
}

}
{

\aufzaehlungvier{Berlaku

\mavergleichskettedisp
{\vergleichskette
{ d \omega
}
{ =} { d \left(  ydx+zdy +xdz  \right)
}
{ =} { dy \wedge dx +dz \wedge dy +dx \wedge dz
}
{ =} { - dx \wedge dy -dy \wedge dz +dx \wedge dz
}
{ } {
}
}
{}{}{.}
}{Berlaku

\mavergleichskettealign
{\vergleichskettealign
{  \varphi^* \omega
}
{ =} { v^2 d u^2 + uvdv^2+ u^2duv
}
{ =} { 2v^2udu +2uv^2 dv + u^2(udv +vdu)
}
{ =} { \left(  2uv^2 +u^2v  \right) du + \left(  2uv^2 +u^3  \right) dv
}
{ } {
}
}
{}
{}{.}
}{Berlaku

\mavergleichskettealign
{\vergleichskettealign
{  d \left(   \left(  2uv^2 +u^2v  \right) du + \left(  2uv^2 +u^3  \right) dv     \right)
}
{ =} {  d \left(  2uv^2 +u^2v  \right) du  +  d\left(  2uv^2 +u^3  \right) dv
}
{ =} { \left(  4uv  + u^2      \right)  dv \wedge du  +   \left(  2v^2   +3u^2  \right)   du \wedge dv
}
{ =} { \left(  -4uv +  2v^2   +2u^2  \right)   du \wedge dv
}
{ } {
}
}
{}
{}{.}
}{Tarikan balik $\varphi^* d \omega$ adalah

\mavergleichskettealign
{\vergleichskettealign
{ \varphi^* \left(  - dx \wedge dy -dy \wedge dz +dx \wedge dz   \right)
}
{ =} { - d u^2 \wedge dv^2 -d v^2 \wedge duv +d u^2 \wedge duv
}
{ =} { -4uv  du \wedge dv -2v^2 dv \wedge du + 2   u^2 du \wedge dv
}
{ =} { \left(  -4uv+2v^2 + 2 u^2  \right)  du \wedge dv
}
{ } {
}
}
{}
{}{.}
}

 }



\inputaufgabeklausurloesung
{3}

{

Misalkan $M$ suatu
\definitionsverweis {manifold diferensiabel}{}{}
dan $\omega$ suatu
\definitionsverweis {bentuk diferensial}{}{}
yang
\definitionsverweis {eksak}{}{}
pada $M$. Misalkan $S$ suatu
\definitionsverweis {manifold diferensiabel}{}{}
yang
\definitionsverweis {kompak}{}{}
dan
\definitionsverweis {berorientasi}{}{}
\zusatzklammer {tanpa batas} {} {}
dengan
\definitionsverweis {basis terhitung bagi topologinya}{}{}
dan misalkan
\maabbdisp {\varphi} { S } { M
} {}
suatu
\definitionsverweis {pemetaan terdiferensialkan secara kontinu}{}{.}
Tunjukkan bahwa

\mavergleichskettedisp
{\vergleichskette
{
\int_{  S  }   \varphi^*\omega
}
{ =} { 0
}
{ } {
}
{ } {
}
{ } {
}
}
{}{}{.}

}
{

Karena menurut
Satz 20.4 (Differentialgeometrie (Osnabrück 2023))  (5)
operasi tarikan balik bentuk kompatibel dengan turunan eksterior,
\mathl{\varphi^*\omega}{} juga eksak. Misalkan $\tau$ adalah suatu bentuk diferensiabel pada $S$ dengan

\mavergleichskettedisp
{\vergleichskette
{d \tau
}
{ =} { \varphi^* \omega
}
{ } {
}
{ } {
}
{ } {
}
}
{}{}{.}
Menurut
Satz von Stokes,
berlaku

\mavergleichskettedisp
{\vergleichskette
{ \int_S \varphi^* \omega
}
{ =} { \int_S d \tau
}
{ =} { \int_{ \partial S } \tau
}
{ =} { \int_\emptyset \tau
}
{ =} { 0
}
}
{}{}{.}

 }



\inputaufgabeklausurloesung
{6}

{

Buktikan teorema tentang retraksi ke batas pada manifold.

}
{

Batas
\mathl{\partial M}{}, menurut
Satz 21.8 (Differentialgeometrie (Osnabrück 2023))
, merupakan suatu
\definitionsverweis {manifold diferensiabel}{}{}
\definitionsverweis {berorientasi}{}{}
\zusatzklammer {tanpa batas} {} {.}
Oleh karena itu, menurut
Satz 22.11 (Differentialgeometrie (Osnabrück 2023))
, terdapat suatu
\definitionsverweis {bentuk volume positif}{}{}
\definitionsverweis {yang terdiferensialkan secara kontinu}{}{}
$\tau$ pada
\mathl{\partial M}{.} Kita mempunyai
\mathl{\int_{  \partial M  }  \tau >0}{.}
\definitionsverweis {Turunan eksterior}{}{}
dari bentuk volume $\tau$ adalah $0$. Andaikan terdapat suatu
\definitionsverweis {pemetaan yang terdiferensialkan secara kontinu}{}{}
\maabbdisp {\varphi} { M } { \partial M
} {}
dengan
\mathl{\varphi \vert_{\partial M} = \operatorname{Id}_{\partial M}}{.} Maka
\definitionsverweis {bentuk tarikan balik}{}{}

\mathl{\varphi^* \tau}{} adalah
$(n-1)$-\definitionsverweis {bentuk diferensial}{}{}
pada $M$ yang pembatasannya pada batas berimpit dengan $\tau$. Dengan menggunakan
Satz 23.2 (Differentialgeometrie (Osnabrück 2023))
dan
Satz 20.4 (Differentialgeometrie (Osnabrück 2023))  (5)
kita memperoleh

\mavergleichskettealign
{\vergleichskettealign
{
\int_{  \partial M  }  \tau
}
{ =} {
\int_{  \partial  M  }  \varphi^* \tau
}
{ =} {
\int_{   M  }  d(\varphi^* \tau)
}
{ =} {
\int_{   M  }  \varphi^* (d\tau)
}
{ =} {
\int_{   M  }  \varphi^* (0)
}
}
{
\vergleichskettefortsetzungalign
{ =} { 0
}
{ } {}
{ } {}
{ } {}
}
{}{.}
Ini merupakan suatu kontradiksi.

 }

\endgroup

\chapter{Formulir Solusi Resmi Ujian 10}
\noindent\textit{Solusi yang disediakan oleh sumber resmi; kekosongan sumber dipertahankan.}
\begingroup
\renewcommand{\theHfakt}{o011.exam.official.10.\arabic{fakt}}
%Data institusi

%\input{Dozentdaten}

%\renewcommand{\fachbereich}{Bidang studi}

%\renewcommand{\dozent}{Prof. Dr. . }

%Data ujian

\renewcommand{\klausurgebiet}{ }

\renewcommand{\klausurtyp}{ }

\renewcommand{\klausurdatum}{. 20}

\klausurvorspann {\fachbereich} {\klausurdatum} {\dozent} {\klausurgebiet} {\klausurtyp}

%Data untuk tabel poin berikut

\renewcommand{\aeins}{  3  }

\renewcommand{\azwei}{  3   }

\renewcommand{\adrei}{  6  }

\renewcommand{\avier}{  4  }

\renewcommand{\afuenf}{  5  }

\renewcommand{\asechs}{  0  }

\renewcommand{\asieben}{  5  }

\renewcommand{\aacht}{  6  }

\renewcommand{\aneun}{  6  }

\renewcommand{\azehn}{  0  }

\renewcommand{\aelf}{  6  }

\renewcommand{\azwoelf}{  12  }

\renewcommand{\adreizehn}{  0   }

\renewcommand{\avierzehn}{  9  }

\renewcommand{\afuenfzehn}{  65  }

\renewcommand{\asechzehn}{    }

\renewcommand{\asiebzehn}{    }

\renewcommand{\aachtzehn}{    }

\renewcommand{\aneunzehn}{    }

\renewcommand{\azwanzig}{    }

\renewcommand{\aeinundzwanzig}{    }

\renewcommand{\azweiundzwanzig}{    }

\renewcommand{\adreiundzwanzig}{    }

\renewcommand{\avierundzwanzig}{    }

\renewcommand{\afuenfundzwanzig}{    }

\renewcommand{\asechsundzwanzig}{    }

\punktetabellevierzehn

\klausurnote

\newpage

\setcounter{section}{0}



\inputaufgabeklausurloesung
{3}

{

Definisikan istilah-istilah berikut
\zusatzklammer {yang dicetak miring} {} {}.
\aufzaehlungsechs{Suatu
\stichwort {kelengkungan utama} {}
pada suatu permukaan diferensiabel

\mavergleichskette
{\vergleichskette
{ Y
}
{ \subseteq }{ \R^3
}
{  }{
}
{    }{
}
{   }{
}
}
{}{}{}
pada titik

\mavergleichskette
{\vergleichskette
{ P
}
{ \in }{ Y
}
{  }{
}
{    }{
}
{   }{
}
}
{}{}{.}

}{Suatu
\stichwort {difeomorfisme} {}
antara manifold diferensiabel
\mathkor {} {L} {dan} {M} {.}

}{Suatu
\stichwort {vektor singgung} {}
di titik

\mavergleichskette
{\vergleichskette
{  P
}
{ \in }{  M
}
{  }{
}
{    }{
}
{   }{
}
}
{}{}{}
dari manifold diferensiabel $M$.

}{Suatu
\stichwort {operator diferensial orde pertama} {}
pada manifold diferensiabel $M$.

}{Suatu
\stichwort {setengah ruang Euklides} {.}

}{Suatu
\definitionsverweis {penghabisan kompak}{}{}
bagi suatu
\definitionsverweis {ruang topologis}{}{}
$X$.
}

}
{

\aufzaehlungsechs{Setiap
\definitionsverweis {nilai eigen}{}{}
dari
\definitionsverweis {pemetaan Weingarten}{}{}
\maabbeledisp {L_P} {T_P Y} { T_PY
} {v} {- \left(  D_{v} N  \right) \left(  P  \right)
} {,}
disebut
\definitionswort {kelengkungan utama}{}
dari $Y$ di $P$.
}{Suatu
\definitionsverweis {homeomorfisme}{}{}
\maabbdisp {\varphi} {L} {M
} {}
disebut
\stichwortpraemath {C^1} {difeomorfisme}{,}
jika
\mathkor {} {\varphi} {maupun} {\varphi^{-1}} {}
merupakan
$C^1$-\definitionsverweis {pemetaan}{}{}.
}{Yang dimaksud dengan vektor tangen di $P$ adalah suatu
\definitionsverweis {kelas ekuivalensi}{}{}
dari
\definitionsverweis {kurva-kurva diferensiabel}{}{}
yang
\definitionsverweis {ekuivalen secara tangensial}{}{}
melalui $P$.
}{Suatu pemetaan
\maabbdisp {D} {C^1(M,\R) } {C^0(M,\R)
} {}
disebut
operator diferensial orde pertama
jika memenuhi sifat-sifat berikut.
\aufzaehlungzwei {$D$ bersifat
$\R$-\definitionsverweis {linear}{}{.}
} {Berlaku

\mavergleichskette
{\vergleichskette
{ D(fg)
}
{ = }{ fD(g)+gD(f)
}
{  }{
}
{    }{
}
{   }{
}
}
{}{}{.}
}
}{Yang dimaksud dengan setengah ruang Euklides berdimensi $n$ adalah himpunan

\mavergleichskettedisp
{\vergleichskette
{ H
}
{ =} { \left\{ x \in \R^n   \mid  x_1 \geq 0        \right\}
}
{ } {
}
{ } {
}
{ } {
}
}
{}{}{}
dengan
\definitionsverweis {topologi terinduksi}{}{.}
}{Suatu penghabisan kompak

\mathbed {A_n} {}
 {n \in \N} {}
 {} {} {} {,}
dari $X$ adalah suatu
\definitionsverweis {barisan}{}{}
dari
\definitionsverweis {subhimpunan-subhimpunan kompak}{}{}

\mavergleichskette
{\vergleichskette
{ A_n
}
{ \subseteq }{ X
}
{  }{
}
{    }{
}
{   }{
}
}
{}{}{}
dengan

\mathdisp {A_n \subseteq A_{n+1}^{o} \text{ dan } \bigcup_{n=0}^\infty A_n = X} { . }

}

}



\inputaufgabeklausurloesung
{3}

{

Nyatakan teorema-teorema berikut.
\aufzaehlungdrei{Teorema tentang
\definitionsverweis {ruang singgung}{}{}
dari suatu
\definitionsverweis {submanifold tertutup}{}{}

\mavergleichskette
{\vergleichskette
{  M
}
{ \subseteq }{  N
}
{  }{
}
{    }{
}
{   }{
}
}
{}{}{.}}{Teorema tentang bentuk eksplisit dari bentuk diferensial tarik balik $\varphi^* \omega$ oleh pemetaan diferensiabel
\maabbdisp {\varphi} { U } { V
} {}
\zusatzklammer {dengan
\definitionsverweis {himpunan bagian terbuka}{}{}
\mathkor {} {U \subseteq \R^n} {dan} {V \subseteq \R^m} {,}
yang koordinatnya masing-masing dinyatakan dengan
\mathl{x_1  , \ldots , x_n}{} dan
\mathl{y_1  , \ldots , y_m}{}} {} {}
dari suatu
$k$-\definitionsverweis {bentuk diferensial}{}{}
pada $V$ dengan representasi

\mavergleichskettedisp
{\vergleichskette
{  \omega
}
{ =} {\sum_{I,\,  { \# \left(   I  \right) } = k}  f_I  dy_I
}
{ } {
}
{ } {
}
{ } {
}
}
{}{}{,}
dengan
\maabb {f_I} { V } {\R
} {}
adalah
\definitionsverweis {fungsi}{}{.}}{Teorema tentang karakterisasi kurva geodesik terhadap koneksi Levi-Civita pada manifold Riemann $M$.}

}
{

\aufzaehlungdrei{\faktsituation {Misalkan $N$ suatu
\definitionsverweis {manifold diferensiabel}{}{}
ber-
\definitionsverweis {dimensi}{}{}
$n$ dan misalkan

\mavergleichskette
{\vergleichskette
{ M
}
{ \subseteq }{ N
}
{  }{
}
{    }{
}
{   }{
}
}
{}{}{}
suatu
\definitionsverweis {submanifold tertutup}{}{}
berdimensi $m$.}
\faktfolgerung {Maka untuk setiap titik

\mavergleichskette
{\vergleichskette
{ P
}
{ \in }{ M
}
{  }{
}
{    }{
}
{   }{
}
}
{}{}{}
\definitionsverweis {pemetaan tangen}{}{}
\maabbdisp {} { T_PM } { T_PN
} {}
bersifat
\definitionsverweis {injektif}{}{.}}
\faktzusatz {Artinya, ruang tangen
\mathl{T_PM}{} merupakan
\definitionsverweis {subruang vektor}{}{}
berdimensi $m$ dari
\mathl{T_PN}{.}}
\faktzusatz {}}{\faktsituation {Misalkan
\mathkor {} {U \subseteq \R^n} {dan} {V \subseteq \R^m} {}
merupakan
\definitionsverweis {subhimpunan-subhimpunan terbuka}{}{,}
yang koordinatnya masing-masing dinyatakan dengan
\mathl{x_1 , \ldots , x_n}{} dan
\mathl{y_1 , \ldots , y_m}{}. Misalkan
\maabbdisp {\varphi} { U } { V
} {}
suatu
\definitionsverweis {pemetaan diferensiabel}{}{}
dan misalkan $\omega$ suatu
$k$-\definitionsverweis {bentuk diferensial}{}{}
pada $V$ dengan representasi

\mavergleichskettedisp
{\vergleichskette
{ \omega
}
{ =} { \sum_{I,\, { \# \left(   I  \right) } = k}  f_I  dy_I
}
{ } {
}
{ } {
}
{ } {
}
}
{}{}{,}
dengan
\maabb {f_I} { V } {\R
} {}
adalah
\definitionsverweis {
fungsi-fungsi}{}{.}}
\faktfolgerung {Maka
\definitionsverweis {bentuk tarikan balik}{}{}
memiliki representasi

\mavergleichskettealignhandlinks
{\vergleichskettealignhandlinks
{ \varphi^* \omega
}
{ =} { \sum_{I,\, { \# \left(   I  \right) } = k}  \left(  f_I \circ \varphi  \right)  \left( \sum_{J,\, { \# \left(   J  \right) } = k} \det  \left(  \left(  { \frac{   \partial  \varphi_i   }{  \partial   x_j   } }  \right)_{ i \in I, j \in J}  \right)  dx_J \right)
}
{ =} { \sum_{J,\, { \# \left(   J  \right) } = k}   \left( \sum_{I,\, { \# \left(   I  \right) } = k} \left(  f_I \circ \varphi  \right)  \det  \left(  \left(  { \frac{   \partial  \varphi_i   }{  \partial   x_j   } }  \right)_{ i \in I, j \in J}  \right)  \right)  dx_J
}
{ } {
}
{ } {
}
}
{}
{}{.}}
\faktzusatz {}
\faktzusatz {}}{\faktsituation {Suatu kurva yang dua kali
\definitionsverweis {diferensiabel}{}{}
\maabb {\gamma} { I } { M
} {}
pada suatu
\definitionsverweis {manifold Riemann}{}{}
$M$ adalah}
\faktfolgerung {suatu
\definitionsverweis {kurva geodesik}{}{}
\zusatzklammer {terhadap
\definitionsverweis {koneksi Levi-Civita}{}{}} {} {,}
tepat ketika dalam setiap peta kurva tersebut memenuhi
\definitionsverweis {sistem persamaan diferensial biasa orde dua}{}{}

\mavergleichskettedisp
{\vergleichskette
{ \gamma_k^{\prime \prime} + \sum_{ij} \Gamma^k_{ij} \gamma_i' \gamma_j'
}
{ =} { 0
}
{ } {
}
{ } {
}
{ } {
}
}
{}{}{}
\zusatzklammer {untuk semua $k$} {} {}
yang diberikan.}
\faktzusatz {}
\faktzusatz {}}

}



\inputaufgabeklausurloesung
{6 (2+4)}

{

Misalkan

\mavergleichskette
{\vergleichskette
{ \alpha,\beta
}
{ \in }{ \R_+
}
{  }{
}
{    }{
}
{   }{
}
}
{}{}{}
dan misalkan

\mavergleichskettedisp
{\vergleichskette
{ Y
}
{ =} { \left\{ (x,y)\in\R^2  \mid   \alpha x^2+\beta y^2 = 1        \right\}
}
{ } {
}
{ } {
}
{ } {
}
}
{}{}{.}
\aufzaehlungzwei {Tentukan
\definitionsverweis {pemetaan Gauss}{}{}
untuk $Y$.
} { Berikan secara eksplisit suatu
\definitionsverweis {pemetaan invers}{}{}
dari pemetaan Gauss untuk $Y$.
}

}
{

\aufzaehlungzwei {Sebagai
\definitionsverweis {medan normal satuan}{}{},
kita dapat mulai dari gradien fungsi pendefinisi, yaitu

\mavergleichskettedisp
{\vergleichskette
{ G(x,y)
}
{ =} { \begin{pmatrix}  2 \alpha x  \\ 2 \beta y    \end{pmatrix}
}
{ } {
}
{ } {
}
{ } {
}
}
{}{}{}
dan setelah normalisasi diperoleh

\mavergleichskettedisp
{\vergleichskette
{ N(x,y)
}
{ =} {  { \frac{   1 }{  \sqrt{ \alpha^2 x^2 + \beta^2 y^2 }  } } \begin{pmatrix}   \alpha x  \\  \beta y    \end{pmatrix}
}
{ } {
}
{ } {
}
{ } {
}
}
{}{}{.}
Jadi, pemetaan Gauss adalah
\maabbeledisp {\varphi} {Y = \left\{  \begin{pmatrix}  x  \\ y   \end{pmatrix}   \in \R^2  \mid   \alpha x^2+\beta y^2 = 1        \right\}  } { S^1
} { \begin{pmatrix}  x  \\ y   \end{pmatrix} } { { \frac{   1 }{  \sqrt{ \alpha^2 x^2 + \beta^2 y^2 }  } } \begin{pmatrix}   \alpha x  \\  \beta y    \end{pmatrix}
} {.}
} {Kita nyatakan bahwa
\maabbeledisp {\psi} {S^1} {Y
} { \begin{pmatrix}  u  \\v   \end{pmatrix} } { { \frac{   1 }{  \sqrt{ { \frac{  u^2 }{ \alpha } }   + { \frac{  v^2 }{ \beta } } }  } }       \begin{pmatrix}   { \frac{   u  }{ \alpha } }   \\ { \frac{   v  }{ \beta } }    \end{pmatrix}
} {,}
adalah fungsi invers. Karena

\mavergleichskettedisp
{\vergleichskette
{ { \frac{   1 }{   { \frac{  u^2 }{ \alpha } }   + { \frac{  v^2 }{ \beta } }  } } \left(  \alpha { \frac{  u^2 }{ \alpha^2 } } + \beta{ \frac{  v^2 }{ \beta^2 } }  \right)
}
{ =} { 1
}
{ } {
}
{ } {
}
{ } {
}
}
{}{}{}
citra $\psi$ terletak dalam $Y$. Karena

\mavergleichskettealign
{\vergleichskettealign
{  \varphi \left(  \psi \begin{pmatrix}  u  \\v   \end{pmatrix}  \right)
}
{ =} { \varphi \left(  { \frac{   1 }{  \sqrt{ { \frac{  u^2 }{ \alpha } }   + { \frac{  v^2 }{ \beta } } }  } }       \begin{pmatrix}   { \frac{   u  }{ \alpha } }   \\ { \frac{   v  }{ \beta } }    \end{pmatrix}    \right)
}
{ =} { \varphi \begin{pmatrix}   { \frac{   1 }{  \alpha    \sqrt{ { \frac{  u^2 }{ \alpha } }   + { \frac{  v^2 }{ \beta } } }  } }   u    \\ { \frac{   1 }{ \beta  \sqrt{ { \frac{  u^2 }{ \alpha } }  + { \frac{  v^2 }{ \beta } } }  } }  v       \end{pmatrix}
}
{ =} { { \frac{   1 }{  \sqrt{  \alpha^2 { \frac{  u^2 }{  \alpha^2   \left(   { \frac{  u^2 }{ \alpha } }  +  { \frac{  v^2 }{ \beta } }   \right)  } }  + \beta^2 { \frac{  v^2 }{  \beta^2 \left(   { \frac{  u^2 }{ \alpha } }  +  { \frac{  v^2 }{ \beta } }   \right)  } }       }  } }  \begin{pmatrix}  { \frac{   1 }{  \sqrt{ { \frac{  u^2 }{ \alpha } }   + { \frac{  v^2 }{ \beta } } }  } }   u    \\   { \frac{   1 }{   \sqrt{ { \frac{  u^2 }{ \alpha } } + { \frac{  v^2 }{ \beta } } }  } } v    \end{pmatrix}
}
{ =} { { \frac{   1 }{  \sqrt{ { \frac{  u^2 }{  \left(   { \frac{  u^2 }{ \alpha } }  +  { \frac{  v^2 }{ \beta } }   \right)  } }   + { \frac{  v^2 }{  \left(   { \frac{  u^2 }{ \alpha } }  +  { \frac{  v^2 }{ \beta } }   \right)  } }       }  } }  \begin{pmatrix}   { \frac{   1 }{  \sqrt{ { \frac{  u^2 }{ \alpha } }   + { \frac{  v^2 }{ \beta } } }  } }   u    \\   { \frac{   1 }{   \sqrt{ { \frac{  u^2 }{ \alpha } }   + { \frac{  v^2 }{ \beta } } }  } } v    \end{pmatrix}
}
}
{
\vergleichskettefortsetzungalign
{ =} { { \frac{    \sqrt{ { \frac{  u^2 }{ \alpha } }  +  { \frac{  v^2 }{ \beta } } } }{  \sqrt{    u^2+ v^2       }  } }  \begin{pmatrix}   { \frac{   1 }{  \sqrt{ { \frac{  u^2 }{ \alpha } }   + { \frac{  v^2 }{ \beta } } }  } }   u    \\   { \frac{   1 }{   \sqrt{ { \frac{  u^2 }{ \alpha } }   + { \frac{  v^2 }{ \beta } } }  } } v    \end{pmatrix}
}
{ =} { \begin{pmatrix}  u  \\v   \end{pmatrix}
}
{ } {
}
{ } {}
}
{}{}
dan

\mavergleichskettealign
{\vergleichskettealign
{  \psi \left(  \varphi \begin{pmatrix}  x  \\ y   \end{pmatrix}  \right)
}
{ =} { \psi \begin{pmatrix}  { \frac{    \alpha x }{  \sqrt{\alpha^2 x^2+ \beta^2 y^2 }  } }   \\ { \frac{   \beta y }{  \sqrt{\alpha^2 x^2+ \beta^2 y^2  } }}    \end{pmatrix}
}
{ =} { { \frac{   1 }{  \sqrt{ { \frac{   \alpha x^2 }{  \alpha^2x^2 + \beta^2y^2  } }  + { \frac{   \beta y^2 }{  \alpha^2x^2 + \beta^2y^2  } }  }  } }  \begin{pmatrix}  { \frac{    x }{  \sqrt{\alpha^2 x^2+ \beta^2 y^2 }  } }   \\ { \frac{   y }{  \sqrt{\alpha^2 x^2+ \beta^2 y^2  } }}    \end{pmatrix}
}
{ =} { { \frac{   \sqrt{\alpha^2 x^2+ \beta^2 y^2 } }{  \sqrt{    \alpha x^2 +  \beta y^2  }  } }  \begin{pmatrix}  { \frac{    x }{  \sqrt{\alpha^2 x^2+ \beta^2 y^2 }  } }   \\ { \frac{   y }{  \sqrt{\alpha^2 x^2+ \beta^2 y^2  } }}    \end{pmatrix}
}
{ =} { \begin{pmatrix}  x  \\ y   \end{pmatrix}
}
}
{}
{}{}
kedua pemetaan tersebut saling invers.
}

}



\inputaufgabeklausurloesung
{4}

{

Jelaskan pasangan istilah \anfuehrung{ekstrinsik dan intrinsik}{} dalam konteks manifold.

}
{  }



\inputaufgabeklausurloesung
{5}

{

Buktikan rumus luas permukaan suatu benda putar yang dihasilkan oleh
\definitionsverweis {kurva diferensiabel}{}{}
\maabbeledisp {\gamma} {]a,b[} {\R^2
} { t } {(x(t),y(t))
} {,}
dengan

\mavergleichskette
{\vergleichskette
{ y(t)
}
{ > }{ 0
}
{  }{
}
{    }{
}
{   }{
}
}
{}{}{.}

}
{

Misalkan $S$ permukaan putar, yang merupakan submanifold tertutup berdimensi dua dalam suatu himpunan terbuka di $\R^3$. Kita terapkan
Korolar 17.2 (Geometri diferensial (Osnabrück 2023))
pada parametrisasi
\maabbeledisp {} { ]a,b[ \times ]0, 2 \pi[} {S
} { (t, \alpha) } { (x(t),  y(t) \cos \alpha  , y(t) \sin \alpha  )
} {,}.
Turunan-turunan parsialnya adalah

\mathdisp {\begin{pmatrix}  x'(t) \\ y'(t) \cos \alpha \\  y'(t) \sin \alpha   \end{pmatrix} \text{   dan } \begin{pmatrix}  0  \\ - y(t) \sin \alpha \\  y(t) \cos \alpha   \end{pmatrix}} {  }

dan karena itu

\mathdisp {E(t, \alpha) = (x'(t))^2 + (y'(t))^2,\, F(t, \alpha)=0,\, G(t, \alpha) = (y(t))^2} { . }

Jadi, luas permukaannya sama dengan

\mavergleichskettedisphandlinks
{\vergleichskettedisphandlinks
{ \int_{  a  }^{  b  } \int_{  0  }^{  2 \pi } \sqrt{(x'(t))^2+(y'(t))^2} \cdot   y(t) \, d \alpha \, d t
}
{ =} { 2 \pi \int_{  a  }^{  b  } \sqrt{(x'(t))^2+(y'(t))^2} \cdot   y(t) \, d t
}
{ } {
}
{ } {
}
{ } {
}
}
{}{}{.}

}



\inputaufgabeklausurloesung
{0}

{

}
{  }



\inputaufgabeklausurloesung
{5}

{

Misalkan $X$ adalah suatu
\definitionsverweis {torus}{}{.}
Berikan suatu
\definitionsverweis {pemetaan diferensiabel}{}{}
surjektif
\maabbdisp {\varphi} {\R^2} {X
} {}
sedemikian sehingga
\definitionsverweis {pemetaan singgung}{}{}
\maabbdisp {T_P(\varphi)} { T_P\R^2 } { T_{\varphi(P)}X
} {}
juga surjektif di setiap titik

\mavergleichskette
{\vergleichskette
{ P
}
{ \in }{ \R^2
}
{  }{
}
{    }{
}
{   }{
}
}
{}{}{.}

}
{

Kita tulis torus sebagai $X=S^1 \times S^1$ dengan lingkaran satuan
\mathl{S^1= \left\{ (x,y) \in \R^2  \mid   \betrag { (x,y) } = 1        \right\}}{.}
Pemetaan
\maabbeledisp {\varphi} {\R^2} {S^1 \times S^1
} {(u,v)} {( \cos  u, \sin  u , \cos  v , \sin  v )
} {,}
bersifat surjektif karena kedua komponennya merupakan parametrisasi standar lingkaran satuan. Surjektivitas pemetaan tangen juga dapat diperiksa per komponen karena ruang tangen manifold produk adalah hasil kali ruang-ruang tangennya. Turunan dari
\maabbeledisp {} {\R} { \R^2
} { t } { ( \cos  t, \sin  t)
} {,}
adalah
\mathl{(- \sin  t, \cos  t) \neq (0,0)}{,} sehingga vektor citra ini selalu merentang ruang tangen satu-dimensional dari lingkaran satuan di titik citra.

}



\inputaufgabeklausurloesung
{6 (2+2+2)}

{

Tinjau grafik $M$ dari pemetaan
\maabbeledisp {\varphi} {\R^2} {\R
} {(u,v)} { u^2+uv-v^3
} {,}
sebagai submanifold tertutup berdimensi dua di $\R^3$, yaitu

\mavergleichskettedisp
{\vergleichskette
{ M
}
{ =} { \left\{ (u,v,u^2+uv-v^3)  \mid  (u,v) \in \R^2         \right\}
}
{ } {
}
{ } {
}
{ } {
}
}
{}{}{}
dengan metrik Riemann yang diinduksi dari $\R^3$. Misalkan
\maabbeledisp {\psi} {\R^2} { M
} {(u,v)} { (u,v,u^2+uv-v^3)
} {,}
adalah difeomorfisme yang bersesuaian.

a) Tentukan diferensial total dari $\psi$ serta vektor-vektor citra
\mathl{T_P(\psi) (e_1)}{} dan
\mathl{T_P(\psi) (e_2)}{} di
\mathl{T_{\psi(P)}M}{.}

b) Untuk setiap titik berbentuk

\mavergleichskette
{\vergleichskette
{ P
}
{ = }{ (u,0)
}
{  }{
}
{    }{
}
{   }{
}
}
{}{}{}
tentukan luas jajaran genjang yang direntang oleh
\mathl{T_P(\psi) (e_1)}{} dan
\mathl{T_P(\psi) (e_2)}{} di
\mathl{T_{\psi(P)}M}{}.

c) Untuk setiap titik berbentuk

\mavergleichskette
{\vergleichskette
{ P
}
{ = }{ (0,v)
}
{  }{
}
{    }{
}
{   }{
}
}
{}{}{}
tentukan luas jajaran genjang yang direntang oleh
\mathl{T_P(\psi) (e_1)}{} dan
\mathl{T_P(\psi) (e_2)}{} di
\mathl{T_{\psi(P)}M}{}.

}
{

a) Diferensial total dari $\psi$ di titik
\mathl{P=(u,v)}{} adalah

\mathdisp {\left(D \psi\right)_{P} = \begin{pmatrix}  1 &  0 \\  0 &  1 \\   2u+v & u-3v^2 \end{pmatrix}} {  }

dan berlaku

\mathdisp {T_P(\psi) (e_1) = \left(  D \psi  \right)_{ P } \left(  e_1  \right) = \begin{pmatrix}  1 \\ 0\\  2u+v   \end{pmatrix} \text{   dan } T_P(\psi) (e_2)   = \left(  D \psi  \right)_{ P } \left(  e_2  \right) = \begin{pmatrix}  0 \\ 1\\ u-3v^2   \end{pmatrix}} { . }

b) dan c) Untuk menentukan luas, pertama-tama kita hitung hasil kali skalar dari kedua vektor
\mathkor {} {\begin{pmatrix}  1 \\ 0\\  2u+v   \end{pmatrix}} {dan} {\begin{pmatrix}  0 \\ 1\\ u-3v^2   \end{pmatrix}} {.}
Kita peroleh

\mathdisp {\left\langle \begin{pmatrix}  1 \\ 0\\  2u+v   \end{pmatrix}  , \begin{pmatrix}  1 \\ 0\\  2u+v   \end{pmatrix}   \right\rangle = 1+(2u+v)^2 = 1 +4u^2+4uv +v^2} { , }

\mathdisp {\left\langle \begin{pmatrix}  1 \\ 0\\  2u+v   \end{pmatrix}  , \begin{pmatrix}  0 \\ 1\\ u-3v^2   \end{pmatrix}    \right\rangle = (2u+v)(u-3v^2) = 2u^2 -6uv^2 +uv -3v^3} {  }

dan

\mathdisp {\left\langle \begin{pmatrix}  0 \\ 1\\ u-3v^2   \end{pmatrix}   , \begin{pmatrix}  0 \\ 1\\ u-3v^2   \end{pmatrix}    \right\rangle = 1+(u-3v^2)^2 = 1 +u^2 -6uv^2  + 9v^4} { . }

b) Untuk
\mathl{P=(u,0)}{}, luas jajaran genjang yang direntang oleh
\mathl{T_P(\psi) (e_1)}{} dan
\mathl{T_P(\psi) (e_2)}{} di
\mathl{T_{\psi(P)}M}{} adalah

\mavergleichskettedisp
{\vergleichskette
{ \sqrt { \det  \begin{pmatrix}  1+4u^2  &   2u^2   \\  2 u^2 &  1+ u^2  \end{pmatrix} }
}
{ =} { \sqrt{  (1+4u^2)(1+u^2) - 4u^4 }
}
{ =} { \sqrt{  1+5u^2 }
}
{ } {
}
{ } {
}
}
{}{}{.}

c) Untuk
\mathl{P=(0,v)}{}, luas jajaran genjang yang direntang oleh
\mathl{T_P(\psi) (e_1)}{} dan
\mathl{T_P(\psi) (e_2)}{} di
\mathl{T_{\psi(P)}M}{} adalah

\mavergleichskettedisp
{\vergleichskette
{ \sqrt { \det  \begin{pmatrix}  1+v^2  &   -3v^3   \\  -3  v^3 &  1+ 9v^4  \end{pmatrix} }
}
{ =} { \sqrt{  (1+v^2)(1+9v^4) - 9v^6 }
}
{ =} { \sqrt{  1+v^2 +9v^4  }
}
{ } {
}
{ } {
}
}
{}{}{.}

}



\inputaufgabeklausurloesung
{6}

{

Buktikan aturan hasil kali untuk bentuk diferensial yang diferensiabel pada suatu himpunan terbuka di $\R^n$.

}
{

Misalkan
\mathl{x_1 , \ldots , x_n}{} adalah koordinat pada $\R^n$. Karena linearitas $d$ dan
multilinearitas hasil kali baji,
kita dapat menulis kedua bentuk diferensial sebagai
\mathkor {} {\omega=f dx_I} {dan} {\tau=g dx_J} {}
dengan himpunan indeks
\mathkor {} {I=\{i_1 , \ldots , i_k\}} {dan} {J=\{ j_1 , \ldots , j_\ell \}} {}.
Maka

\mavergleichskettealignhandlinks
{\vergleichskettealignhandlinks
{ d( \omega \wedge \tau)
}
{ =} { d \left(  fdx_I \wedge gdx_J  \right)
}
{ =} { d \left(  (f g) dx_I \wedge dx_J  \right)
}
{ =} { \sum_{s = 1}^n { \frac{   \partial  (fg)  }{  \partial   x_s  } }  dx_s  \wedge  dx_I  \wedge dx_J
}
{ =} { \sum_{s = 1}^n \left(  g { \frac{   \partial   f   }{  \partial   x_s  } } + f { \frac{   \partial   g   }{  \partial   x_s  } }  \right)  dx_s \wedge  dx_I  \wedge dx_J
}
}
{
\vergleichskettefortsetzungalign
{ =} { \sum_{s = 1}^n  g { \frac{   \partial   f   }{  \partial   x_s  } } dx_s \wedge  dx_I \wedge dx_J + \sum_{s = 1}^n  f { \frac{   \partial   g   }{  \partial   x_s  } }  dx_s \wedge  dx_I  \wedge dx_J
}
{ =} { \sum_{s = 1}^n  { \frac{   \partial   f   }{  \partial   x_s  } } dx_s \wedge  dx_I \wedge gdx_J + \sum_{s = 1}^n  { \frac{   \partial   g   }{  \partial   x_s  } }  dx_s \wedge f dx_I \wedge dx_J
}
{ =} { d(fdx_I ) \wedge g dx_J  + \sum_{s = 1}^n (-1)^k  f dx_I \wedge { \frac{   \partial   g   }{  \partial   x_s  } }  dx_s \wedge dx_J
}
{ =} { d(fdx_I ) \wedge g dx_J + (-1)^k  f dx_I \wedge  d( g dx_J)
}
}
{}{.}

}



\inputaufgabeklausurloesung
{0}

{

}
{  }



\inputaufgabeklausurloesung
{6}

{

Misalkan $M$ adalah suatu
\definitionsverweis {manifold berbatas}{}{}
dan misalkan

\mavergleichskette
{\vergleichskette
{ P
}
{ \in }{ \partial M
}
{  }{
}
{    }{
}
{   }{
}
}
{}{}{}
adalah titik batas. Tunjukkan bahwa sifat-sifat berikut ekuivalen bagi suatu vektor singgung

\mavergleichskette
{\vergleichskette
{ v
}
{ \in }{ T_PM
}
{  }{
}
{    }{
}
{   }{
}
}
{}{}{.}
\aufzaehlungzwei {Terdapat suatu lintasan yang terdiferensialkan secara kontinu
\maabb {\gamma} { [- \epsilon, 0]} { M
} {}
dengan

\mavergleichskette
{\vergleichskette
{ \gamma(0)
}
{ = }{ P
}
{  }{
}
{    }{
}
{   }{
}
}
{}{}{}
dan

\mavergleichskette
{\vergleichskette
{ \gamma'(0)
}
{ = }{ v
}
{  }{
}
{    }{
}
{   }{
}
}
{}{}{.}
} { Pada setiap peta bermodel setengah ruang negatif, $v$ dipetakan oleh pemetaan singgung ke setengah ruang kanan.
}

}
{

Misalkan $\gamma$ suatu setengah lintasan di $M$ seperti yang diberikan, dan misalkan

\mavergleichskette
{\vergleichskette
{ P
}
{ \in }{  U
}
{ \subseteq }{ M
}
{    }{
}
{   }{
}
}
{}{}{}
suatu daerah peta dengan peta
\maabbdisp {\alpha} { U } {V \subseteq  H_{\leq 0}
} {}
dengan

\mavergleichskettedisp
{\vergleichskette
{  H_{\leq 0}
}
{ =} { \left\{ (x_1,x_2 , \ldots, x_n)  \mid   x_1 \leq 0        \right\}
}
{ } {
}
{ } {
}
{ } {
}
}
{}{}{}
dan

\mavergleichskettedisp
{\vergleichskette
{ Q
}
{ =} { \alpha(P)
}
{ =} { \left(  0   , \,  a_2  , \,   \ldots , \,  a_n                \right)
}
{ } {
}
{ } {
}
}
{}{}{.}
Di bawah pemetaan tangen,

\mavergleichskette
{\vergleichskette
{ v
}
{ = }{ \gamma'(0)
}
{  }{
}
{    }{
}
{   }{
}
}
{}{}{}
dipetakan ke vektor turunan dari

\mavergleichskette
{\vergleichskette
{ \delta
}
{ = }{ \alpha \circ \gamma
}
{  }{
}
{    }{
}
{   }{
}
}
{}{}{}
di titik nol. Karena $\gamma$ berjalan dalam $M$, seluruh $\delta$ berjalan dalam $H_{\leq 0}$. Oleh karena itu, dalam hasil bagi selisih

\mathdisp {{ \frac{   \left(  \delta_1(t)   , \,   \delta_2(t)  , \,   \ldots , \,   \delta_n(t)                \right)  }{ t } }} {  }

baik $t$ maupun $\delta_1(t)$ bernilai negatif, sehingga

\mavergleichskettedisp
{\vergleichskette
{  \delta_1'(0)
}
{ \geq} { 0
}
{ } {
}
{ } {
}
{ } {
}
}
{}{}{,}
dan dengan demikian termasuk dalam setengah ruang positif.

Sebaliknya, misalkan
\mathl{T_P(\alpha)(v)}{} termasuk dalam setengah ruang positif. Maka garis
\maabbeledisp {\delta} {\R} {\R^n
} { t } { Q+tv
} {}
untuk

\mavergleichskette
{\vergleichskette
{ t
}
{ \leq }{  0
}
{  }{
}
{    }{
}
{   }{
}
}
{}{}{}
seluruhnya berjalan dalam setengah ruang negatif. Lintasan
\mathl{\alpha^{-1} \circ \delta}{} didefinisikan pada suatu interval negatif, bernilai dalam $M$, dan merupakan realisasi setengah lintasan dari $v$.

}



\inputaufgabeklausurloesung
{12}

{

Buktikan teorema Stokes untuk manifold berbatas.

}
{

\teilbeweis {}{}{}
{Misalkan

\mathbed {U_i} {}
 {i \in I} {}
 {} {} {} {,}
suatu
\definitionsverweis {penutup terbuka}{}{}
dari $M$ dengan
\definitionsverweis {peta-peta berorientasi}{}{},
dan misalkan

\mathbed {h_j} {}
 {j \in J} {}
 {} {} {} {,}
suatu
\definitionsverweis {partisi kesatuan yang terdiferensialkan secara kontinu dan subordinat terhadap penutup tersebut}{}{,}
yang ada menurut
Teorema 22.10 (Geometri diferensial (Osnabrück 2023)).
Untuk setiap

\mavergleichskette
{\vergleichskette
{ P
}
{ \in }{ M
}
{  }{
}
{    }{
}
{   }{
}
}
{}{}{}
terdapat lingkungan terbuka

\mavergleichskette
{\vergleichskette
{ P
}
{ \in }{ V_P
}
{ \subseteq }{ M
}
{    }{
}
{   }{
}
}
{}{}{}
sedemikian sehingga

\mavergleichskette
{\vergleichskette
{ h_j \vert_{V_P }
}
{ = }{ 0
}
{  }{
}
{    }{
}
{   }{
}
}
{}{}{}
kecuali untuk berhingga banyak indeks. Misalkan $Y$ dukungan dari $\omega$. Penutup

\mavergleichskette
{\vergleichskette
{ Y
}
{ \subseteq }{ \bigcup_{P \in M} V_P
}
{  }{
}
{    }{
}
{   }{
}
}
{}{}{}
memiliki subpenutup berhingga karena
\definitionsverweis {kekompakan}{}{}
yang diasumsikan, katakanlah

\mavergleichskettedisp
{\vergleichskette
{ Y
}
{ \subseteq} { V_{1} \cup \ldots \cup V_{r}
}
{ =} { V
}
{ } {
}
{ } {
}
}
{}{}{.}
Jadi, hanya berhingga banyak $h_j$ pada $V$ yang tidak sama dengan $0$. Kita tetapkan

\mavergleichskette
{\vergleichskette
{ \omega_j
}
{ = }{ h_j \omega
}
{  }{
}
{    }{
}
{   }{
}
}
{}{}{;}
bentuk-bentuk diferensial ini juga terdiferensialkan secara kontinu. Dukungan $h_j$ adalah suatu subhimpunan dalam $M$ yang
\definitionsverweis {tertutup}{}{,}
dan terletak dalam
\mathl{U_{i(j)}}{}; karena itu, dukungan $\omega_j$ terletak dalam
\mathl{Y \cap U_{i(j)}}{} dan bersifat kompak menurut
Soal 13.10 (Geometri diferensial (Osnabrück 2023)).
Berlaku

\mavergleichskettedisp
{\vergleichskette
{ \omega
}
{ =} { \sum_{j \in J} \omega_j
}
{ } {
}
{ } {
}
{ } {
}
}
{}{}{,}
dan hanya berhingga banyak bentuk diferensial
\mathl{\omega_j}{} ini yang tidak sama dengan $0$ karena

\mavergleichskette
{\vergleichskette
{ \omega_j \vert_{M \setminus Y}
}
{ = }{ 0
}
{  }{
}
{    }{
}
{   }{
}
}
{}{}{}
untuk semua $j$, sedangkan

\mavergleichskettedisp
{\vergleichskette
{ \omega_j \vert_V
}
{ =} { 0
}
{ } {
}
{ } {
}
{ } {
}
}
{}{}{}
untuk semua $j$ kecuali berhingga banyak. Berdasarkan
aditivitas integral bentuk diferensial
dan
aditivitas turunan eksterior,
pernyataan dapat dibuktikan secara terpisah untuk setiap $\omega_j$. Jadi, kita dapat mengasumsikan bahwa diberikan suatu bentuk diferensial berderajat $(n-1)$ yang terdiferensialkan secara kontinu, memiliki dukungan kompak, dan seluruh dukungannya terletak dalam suatu lingkungan peta $U$.}
{}
\teilbeweis {}{}{}
{Terdapat suatu
\definitionsverweis {difeomorfisme}{}{}
\maabb {\alpha} { U } { V
} {}
dengan

\mavergleichskette
{\vergleichskette
{ V
}
{ \subseteq }{ H
}
{  }{
}
{    }{
}
{   }{
}
}
{}{}{}
terbuka, yang sekaligus menginduksi difeomorfisme antara batas-batas

\mavergleichskette
{\vergleichskette
{ \partial U
}
{ = }{ \partial M \cap U
}
{  }{
}
{    }{
}
{   }{
}
}
{}{}{}
dan

\mavergleichskette
{\vergleichskette
{ \partial V
}
{ = }{ \partial H \cap V
}
{  }{
}
{    }{
}
{   }{
}
}
{}{}{}.
Di sini berlaku

\mavergleichskettedisp
{\vergleichskette
{
\int_{  \partial U  }   \omega
}
{ =} {
\int_{  \partial V  }  \alpha^{-1 *}\omega
}
{ } {
}
{ } {
}
{ } {
}
}
{}{}{}
dan

\mavergleichskettedisp
{\vergleichskette
{
\int_{ U  }  d \omega
}
{ =} {
\int_{ V  }  \alpha^{-1 *} (  d \omega)
}
{ =} {
\int_{  V  }  d \left(  \alpha^{-1 *}\omega  \right)
}
{ } {
}
{ } {
}
}
{}{}{}
menurut
Lema 20.2 (Geometri diferensial (Osnabrück 2023))  (5).}
{\leerzeichen{}Dengan demikian, kita dapat mulai dari suatu bentuk diferensial yang terdiferensialkan secara kontinu dan memiliki dukungan kompak, didefinisikan pada

\mavergleichskette
{\vergleichskette
{ V
}
{ \subseteq }{ H
}
{  }{
}
{    }{
}
{   }{
}
}
{}{}{},
yang dapat kita perluas menjadi bentuk nol pada seluruh $H$ di luar dukungannya.}
\teilbeweis {}{}{}
{Karena kekompakan dukungannya, terdapat suatu balok yang cukup besar

\mavergleichskette
{\vergleichskette
{ Q
}
{ \subseteq }{ H
}
{  }{
}
{    }{
}
{   }{
}
}
{}{}{,}
yang salah satu sisinya $S$ terletak pada $\partial H$ dan yang bertemu dukungan $\omega$ hanya di $S$. Pada semua sisi lain dari $Q$, $\omega$
\zusatzklammer {dan dengan demikian juga $d\omega$} {} {}
adalah bentuk nol. Oleh karena itu, di satu pihak

\mavergleichskettedisp
{\vergleichskette
{
\int_{ H  }  d\omega
}
{ =} {
\int_{ Q  }  d\omega
}
{ } {
}
{ } {
}
{ } {
}
}
{}{}{}
dan di pihak lain

\mavergleichskettedisp
{\vergleichskette
{
\int_{  \partial H  }   \omega
}
{ =} {
\int_{  S  }   \omega
}
{ =} {
\int_{  \partial Q  }   \omega
}
{ } {
}
{ } {
}
}
{}{}{.}
Jadi, pernyataan tersebut mengikuti dari versi balok teorema Stokes.}
{}

}



\inputaufgabeklausurloesung
{0}

{

}
{  }



\inputaufgabeklausurloesung
{9 (1+1+3+2+2)}

{

Tinjau parametrisasi trigonometrik
\maabbeledisp {\psi} {\R } { S^1
} { t } { \left(  \cos  t   , \,   \sin  t                   \right)
} {}
dari lingkaran satuan. Misalkan

\mavergleichskette
{\vergleichskette
{ c
}
{ \in }{ [0, 2 \pi[
}
{  }{
}
{    }{
}
{   }{
}
}
{}{}{}
dipilih tetap. Pada bundel vektor trivial
\maabbdisp {} {S^1 \times \R^2} { S^1
} {}
diberikan suatu
\definitionsverweis {koneksi linear}{}{}
$\nabla$ sedemikian sehingga, sepanjang $\psi$,
\definitionsverweis {koneksi tarik balik}{}{}
$\psi^* \nabla$ diberikan oleh
\definitionsverweis {simbol Christoffel}{}{}
\zusatzklammer {indeks bagi satu-satunya arah diferensiasi dihilangkan} {} {}

\mavergleichskettedisp
{\vergleichskette
{ \begin{pmatrix} \tilde{\Gamma}_1^1 (t)   &  \tilde{\Gamma}_1^2 (t)   \\ \tilde{\Gamma}_2^1(t)  &  \tilde{\Gamma}_2^2(t)  \end{pmatrix}
}
{ =} { \begin{pmatrix}  0  &   - { \frac{  c  }{  2 \pi   } }   \\  { \frac{  c  }{  2 \pi  } }  &  0  \end{pmatrix}
}
{ } {
}
{ } {
}
{ } {
}
}
{}{}{.}
Untuk suatu titik

\mavergleichskette
{\vergleichskette
{ Q
}
{ = }{ \begin{pmatrix}  a  \\b   \end{pmatrix}
}
{ \in }{ \R^2
}
{    }{
}
{   }{
}
}
{}{}{}
tinjau pemetaan
\maabbeledisp {\Psi_{Q}} {\R} { S^1 \times \R^2
} { t } { \left(  \cos  t   , \,   \sin  t  , \,   M(t)  \begin{pmatrix}  a  \\b   \end{pmatrix}                 \right)
} {}
dengan

\mavergleichskettedisp
{\vergleichskette
{ M(t)
}
{ =} { \begin{pmatrix}  \cos  { \frac{  ct  }{  2 \pi   } }   &   -  \sin  { \frac{  ct  }{  2 \pi   } }    \\  \sin  { \frac{  ct  }{  2 \pi  } }   &  \cos  { \frac{  ct  }{  2 \pi   } }  \end{pmatrix}
}
{ } {
}
{ } {
}
{ } {
}
}
{}{}{.}
\aufzaehlungfuenf{Tentukan $\Psi_Q(0)$.
}{Tentukan $\Psi_Q(2 \pi)$.
}{Tunjukkan bahwa $\Psi_Q$ adalah suatu
\definitionsverweis {pengangkatan horizontal}{}{}
sepanjang $\psi$.
}{Tunjukkan bahwa
\maabbeledisp {} {\R^2} { \R^2
} { Q } { \Psi_Q(2 \pi)
} {,}
merupakan
\definitionsverweis {isometri linear}{}{}
\zusatzklammer {titik pangkal $(1,0)$ tidak dicantumkan di sini} {} {.}
}{Tunjukkan bahwa
\maabbeledisp {} {\Z} { \operatorname{SL}_{  2  } \! \left(  \R  \right)
} { k } { \left(  Q \mapsto \Psi_Q(2k \pi )  \right)
} {,}
merupakan
\definitionsverweis {homomorfisme grup}{}{.}
}

}
{

\aufzaehlungfuenf{Kita peroleh

\mavergleichskettealign
{\vergleichskettealign
{ \Psi_Q(0)
}
{ =} { \left(  \cos  0   , \,   \sin  0  , \,   M(0) \begin{pmatrix}  a  \\b   \end{pmatrix}                 \right)
}
{ =} { \left(  1   , \,   0  , \,   \begin{pmatrix}   1   &   0   \\  0   &  1  \end{pmatrix}    \begin{pmatrix}  a  \\b   \end{pmatrix}                 \right)
}
{ =} { \left(  1   , \,   0  , \,     \begin{pmatrix}  a  \\b   \end{pmatrix}                 \right)
}
{ } {
}
}
{}
{}{.}
}{Kita peroleh

\mavergleichskettealign
{\vergleichskettealign
{ \Psi_Q(2 \pi )
}
{ =} { \left(  \cos  2 \pi   , \,   \sin  2 \pi  , \,   M(2 \pi) \begin{pmatrix}  a  \\b   \end{pmatrix}                 \right)
}
{ =} { \left(  1   , \,   0  , \,   \begin{pmatrix}  \cos  c   &   - \sin c    \\  \sin c   &   \cos  c     \end{pmatrix} \begin{pmatrix}  a  \\b   \end{pmatrix}                 \right)
}
{ =} { \left(  1   , \,   0  , \,     \begin{pmatrix} a \cos  c   - b \sin c   \\ a \sin  c   - b \cos c    \end{pmatrix}                 \right)
}
{ } {
}
}
{}
{}{.}
}{Kita peroleh

\mavergleichskettealign
{\vergleichskettealign
{  \Psi_Q(t)
}
{ =} { \left(  \cos  t    , \,   \sin  t  , \,   \begin{pmatrix}   \cos  { \frac{  ct  }{  2 \pi   } }   &   - \sin  { \frac{  ct  }{  2 \pi   } }    \\  \sin  { \frac{  ct  }{  2 \pi  } }   &   \cos  { \frac{  ct  }{  2 \pi   } }  \end{pmatrix}  \begin{pmatrix}  a  \\b   \end{pmatrix}                 \right)
}
{ =} { \left(  \cos  t   , \,   \sin  t  , \,   \begin{pmatrix}  a \cos  { \frac{  ct  }{  2 \pi  } }  - b \sin  { \frac{  ct  }{  2 \pi  } }   \\ a \sin  { \frac{  ct  }{  2 \pi  } }  +   b  \cos  { \frac{  ct  }{  2 \pi } }    \end{pmatrix}                 \right)
}
{ =} { \left(  \cos  t   , \,   \sin  t  , \,     \left(  a \cos  { \frac{  ct  }{  2 \pi  } }  - b \sin  { \frac{  ct  }{  2 \pi  } }  \right)   e_1 +  \left(  a \sin  { \frac{  ct  }{  2 \pi  } }  +   b  \cos  { \frac{  ct  }{  2 \pi } }  \right)   e_2                 \right)
}
{ } {
}
}
{}
{}{,}
khususnya, ini merupakan suatu pengangkatan. Turunan vertikal dari penampang ini dalam $\R^2$ di atas $\R$ dapat ditentukan menggunakan simbol-simbol Christoffel yang diberikan; kita peroleh

\mavergleichskettealigndrucklinks
{\vergleichskettealigndrucklinks
{ \left(  \psi^* \nabla  \right)_\partial  \left(    \left(  a \cos  { \frac{  ct  }{  2 \pi  } }  - b \sin  { \frac{  ct  }{  2 \pi  } }  \right)   e_1 +  \left(  a \sin  { \frac{  ct  }{  2 \pi  } }  +   b  \cos  { \frac{  ct  }{  2 \pi } }  \right)   e_2  \right)
}
{ =} {   \left(  a \cos  { \frac{  ct  }{  2 \pi  } }  - b \sin  { \frac{  ct  }{  2 \pi  } }  \right)   \left(  \psi^* \nabla  \right)_\partial \left(  e_1  \right)  + \partial    \left(  a \cos  { \frac{  ct  }{  2 \pi  } }  - b \sin  { \frac{  ct  }{  2 \pi  } }  \right)   e_1    +   \left(  a \sin  { \frac{  ct  }{  2 \pi  } }  +   b  \cos  { \frac{  ct  }{  2 \pi } }  \right)      \left(  \psi^* \nabla  \right)_\partial \left(  e_2  \right) + \partial  \left(  a \sin  { \frac{  ct  }{  2 \pi  } }  +   b  \cos  { \frac{  ct  }{  2 \pi } }  \right)   e_2
}
{ =} { -  \left(  a \cos  { \frac{  ct  }{  2 \pi  } }  - b \sin  { \frac{  ct  }{  2 \pi  } }  \right)  { \frac{   c  }{  2 \pi } }  e_2  +   { \frac{   c  }{  2 \pi } } \left(  -a \sin  { \frac{  ct  }{  2 \pi  } }  - b \cos  { \frac{  ct  }{  2 \pi  } }  \right)   e_1    +   \left(  a \sin  { \frac{  ct  }{  2 \pi  } }  +   b  \cos  { \frac{  ct  }{  2 \pi } }  \right) { \frac{   c  }{  2 \pi } }   e_1    +  { \frac{   c  }{  2 \pi } }   \left(  a \cos  { \frac{  ct  }{  2 \pi  } } -  b  \sin  { \frac{  ct  }{  2 \pi } }  \right)   e_2
}
{ =} { 0
}
{ } {
}
}
{}{}{,}
sehingga terdapat suatu
\definitionsverweis {pengangkatan horizontal}{}{}.
}{Pemetaan yang dimaksud adalah
\maabbeledisp {} {\R^2} { \R^2
} { \begin{pmatrix}  a  \\b   \end{pmatrix} } {    \begin{pmatrix}  \cos  c   &   - \sin  c   \\  \sin  c  &  \cos  c  \end{pmatrix} \begin{pmatrix}  a  \\b   \end{pmatrix}
} {,}
yang linear. Matriks yang merepresentasikannya ortogonal
\zusatzklammer {suatu matriks rotasi} {} {,}
sehingga diperoleh suatu isometri.
}{Pemetaan
\mathl{\Psi (2k \pi)}{} adalah pemetaan linear
\maabbeledisp {} {\R^2} { \R^2
} {\begin{pmatrix}  a  \\b   \end{pmatrix} } { \begin{pmatrix}  \cos kc   &   - \sin kc   \\  \sin kc  &  \cos kc  \end{pmatrix} \begin{pmatrix}  a  \\b   \end{pmatrix}
} {,}
yaitu rotasi sebesar sudut $kc$. Matriks rotasi untuk jumlah dua sudut adalah hasil kali matriks dari kedua rotasi tersebut
\zusatzklammer {lihat
Fakta *****} {} {,}
sehingga diperoleh suatu
\definitionsverweis {homomorfisme grup}{}{}.
}

}

\endgroup

\part{Perbaikan Asli yang Terpisah}
\chapter{Enam solusi asli untuk kekosongan sumber}
\noindent\textit{Bagian ini berisi tepat enam solusi asli edisi, berlabel terpisah, dan bukan bagian dari sumber Brenner/Wikiversity maupun formulir solusi resmi.}
% Enam solusi asli untuk kejadian soal yang tidak mempunyai tautan solusi sumber.
% Bukan bagian dari sumber Brenner dan tidak boleh dinisbahkan kepada Brenner/Wikiversity.
% Lisensi: CC BY-SA 4.0.
% Provenans: disusun oleh OpenAI Codex gpt-5.6-sol, Ultra, atas arahan pengguna.

\section{Solusi asli untuk enam kekosongan pada bank ujian}
\label{o011-exam-original-repairs}

Bagian ini memuat tepat enam solusi yang tidak tersedia pada permukaan sumber
yang dibekukan. Setiap solusi diberi ID terpisah dan secara eksplisit ditandai
sebagai materi asli. Teks soal tetap berasal dari formulir peserta resmi;
solusinya tidak dinisbahkan kepada pengarang atau Wikiversity.

\subsection*{Ujian 1, soal tanpa solusi sumber (O011-EXAM01-ORIG-SOL-01)}
\label{o011-exam01-orig-sol-01}

\paragraph{Soal.}
Deskripsikan sebanyak mungkin manifold berbatas berdimensi satu yang
terhubung.

\paragraph{Solusi asli.}
Dengan konvensi manifold Hausdorff dan berbasiskan terhitung yang dipakai
dalam edisi ini, klasifikasi hingga difeomorfisme adalah
\[
 S^1,\qquad \mathbb R,\qquad [0,\infty),\qquad [0,1].
\]
Alasannya dapat dibaca dari banyaknya titik batas dan kekompakan. Jika
batasnya kosong, manifold satu dimensi terhubung adalah lingkaran apabila
kompak dan garis nyata apabila tidak kompak. Jika batasnya mempunyai dua titik,
manifoldnya difeomorfik dengan interval kompak $[0,1]$. Jika batasnya mempunyai
satu titik, manifoldnya difeomorfik dengan setengah garis $[0,\infty)$.
Interval $[0,1)$ tidak memberikan tipe baru karena
$x\mapsto x/(1-x)$ adalah difeomorfisme ke $[0,\infty)$; semua interval terbuka
tak kosong difeomorfik dengan $\mathbb R$. Jadi daftar tersebut juga menjelaskan
mengapa tidak ada tipe tambahan di bawah hipotesis standar.

\subsection*{Ujian 3, soal tanpa solusi sumber (O011-EXAM03-ORIG-SOL-01)}
\label{o011-exam03-orig-sol-01}

\paragraph{Soal.}
Katakan sesuatu yang cerdas tentang pita M\"obius.

\paragraph{Solusi asli.}
Pita M\"obius dapat dibangun sebagai hasil bagi
\[
 \bigl([0,1]\times[-1,1]\bigr)/\bigl((0,t)\sim(1,-t)\bigr).
\]
Ia merupakan manifold mulus kompak berdimensi dua dengan batas; kedua sisi
vertikal direkatkan dengan pembalikan arah. Dua ruas horizontal bergabung
menjadi satu komponen batas. Garis tengah $[0,1]\times\{0\}$ menjadi sebuah
lingkaran dan merupakan retrak deformasi, sehingga pita M\"obius mempunyai
tipe homotopi $S^1$.

Pita ini tidak dapat diorientasikan. Jika suatu basis berorientasi diangkut
sekali mengelilingi garis tengah, identifikasi $(0,t)\sim(1,-t)$ membalik satu
arah basis dan mengembalikannya dengan orientasi berlawanan. Penutup ganda
berorientasinya adalah silinder. Secara setara, pita M\"obius adalah ruang
total bundel garis nyata tak trivial di atas $S^1$; fakta bahwa ia tertanam
sebagai permukaan di $\mathbb R^3$ tidak membuatnya berorientasi.

\subsection*{Ujian 5, soal tanpa solusi sumber (O011-EXAM05-ORIG-SOL-01)}
\label{o011-exam05-orig-sol-01}

\paragraph{Soal.}
Jelaskan peran teorema tentang pemetaan implisit dalam pengembangan teori
manifold.

\paragraph{Solusi asli.}
Teorema pemetaan implisit adalah mesin lokal yang mengubah persamaan menjadi
koordinat. Jika $F\colon U\subseteq\mathbb R^n\to\mathbb R^r$ mulus dan
$D F_p$ surjektif, maka di sekitar $p$ terdapat koordinat mulus yang di
dalamnya $F$ hanyalah proyeksi pada $r$ koordinat. Akibatnya, himpunan aras
regular
\[
 F^{-1}(F(p))
\]
secara lokal adalah $\mathbb R^{n-r}\times\{0\}$ dan karena itu merupakan
submanifold berdimensi $n-r$. Ruang singgungnya diberikan secara langsung oleh
\[
 T_pF^{-1}(F(p))=\ker D F_p.
\]

Hasil ini memasok banyak contoh dasar: hipermuka sebagai aras regular satu
fungsi, grup matriks yang ditentukan oleh persamaan regular, dan grafik solusi
lokal persamaan. Ia juga menjamin bahwa definisi abstrak manifold cocok dengan
gambaran sebagai ruang solusi lokal. Bersama teorema fungsi invers, ia
menyediakan koordinat, membuktikan kemulusan perubahan koordinat yang sesuai,
dan mendasari teorema rank konstan serta kriteria submanifold. Dengan demikian,
teorema itu menghubungkan kalkulus multivariabel dengan teori manifold.

\subsection*{Ujian 7, soal tanpa solusi sumber (O011-EXAM07-ORIG-SOL-01)}
\label{o011-exam07-orig-sol-01}

\paragraph{Soal.}
Hitung luas permukaan bola satuan.

\paragraph{Solusi asli.}
Parametrisasi bola satuan $S^2\subseteq\mathbb R^3$, kecuali garis bujur dan
kedua kutub yang tidak memengaruhi integral, dengan
\[
 \Phi(\theta,\varphi)
 =(\sin\varphi\cos\theta,\sin\varphi\sin\theta,\cos\varphi),
 \quad 0<\theta<2\pi,\quad0<\varphi<\pi.
\]
Vektor turunannya memenuhi
\[
 \left\|\partial_\theta\Phi\times\partial_\varphi\Phi\right\|
 =\sin\varphi.
\]
Maka luasnya adalah
\[
 \operatorname{luas}(S^2)
 =\int_0^{2\pi}\int_0^\pi\sin\varphi\,d\varphi\,d\theta
 =2\pi\,[ -\cos\varphi ]_0^\pi
 =4\pi.
\]

\subsection*{Ujian 9, soal tanpa solusi sumber (O011-EXAM09-ORIG-SOL-01)}
\label{o011-exam09-orig-sol-01}

\paragraph{Soal.}
Deskripsikan suatu model untuk bidang hiperbolik.

\paragraph{Solusi asli.}
Model setengah bidang atas adalah manifold
\[
 \mathbb H^2=\{(x,y)\in\mathbb R^2:y>0\}
\]
dengan metrik Riemann
\[
 g=\frac{dx^2+dy^2}{y^2}.
\]
Ia lengkap, terhubung, dan mempunyai kelengkungan Gaussian konstan $-1$.
Untuk memeriksa kelengkungan, tulis $g=e^{2u}(dx^2+dy^2)$ dengan
$u=-\log y$. Rumus kelengkungan metrik konformal memberi
\[
 K=-e^{-2u}(u_{xx}+u_{yy})
 =-y^2\left(0+\frac1{y^2}\right)=-1.
\]
Geodesiknya adalah garis-garis vertikal dan busur-busur lingkaran yang
memotong batas ideal $y=0$ secara ortogonal. Transformasi M\"obius nyata
\[
 z\longmapsto\frac{az+b}{cz+d},
 \qquad ad-bc=1,
\]
bertindak sebagai isometri; setelah mengidentifikasi matriks yang berbeda
tanda, grup isometri berorientasi adalah $\mathrm{PSL}_2(\mathbb R)$. Model
cakram Poincar\'e bersifat isometrik dengan model ini melalui transformasi
pecahan linear.

\subsection*{Ujian 10, soal tanpa solusi sumber (O011-EXAM10-ORIG-SOL-01)}
\label{o011-exam10-orig-sol-01}

\paragraph{Soal.}
Jelaskan pasangan istilah ``ekstrinsik dan intrinsik'' dalam konteks manifold.

\paragraph{Solusi asli.}
Suatu konsep bersifat \emph{intrinsik} jika dapat didefinisikan dari manifold
dan struktur yang dibawanya tanpa memilih suatu penanaman ke ruang luar.
Contohnya ialah atlas mulus, ruang singgung, metrik Riemann, koneksi
Levi--Civita, panjang kurva, jarak geodesik, dan tensor kelengkungan Riemann.
Isometri manifold mempertahankan semua data tersebut.

Suatu konsep bersifat \emph{ekstrinsik} jika bergantung pada cara manifold
direalisasikan sebagai submanifold ruang ambien. Untuk
$M\subseteq\mathbb R^N$, medan normal, bentuk fundamental kedua, pemetaan
Weingarten, dan kelengkungan rata-rata memakai ruang ambien dan dapat berubah
ketika penanaman diubah, walaupun metrik intrinsik tetap sama.

Perbedaan ini tidak selalu tampak dari rumus awal. Kelengkungan Gaussian suatu
permukaan pertama-tama dapat dihitung melalui bentuk fundamental kedua, tetapi
Theorema egregium menunjukkan bahwa nilainya sebenarnya ditentukan sepenuhnya
oleh metrik intrinsik. Jadi ``intrinsik'' atau ``ekstrinsik'' merujuk pada
ketergantungan geometris konsep, bukan sekadar pada bentuk rumus yang dipakai
untuk menghitungnya.

% O011-COMPLETE-EXTENSION-END

\backmatter
\listoffigures
\chapter*{Atribusi dan Hak Media}
\addcontentsline{toc}{chapter}{Atribusi dan Hak Media}
\section*{Unit 1}
\addcontentsline{toc}{section}{Unit 1}
Empat berkas berikut digunakan dalam Unit 1. Setiap berkas tetap mengikuti lisensinya sendiri; tautan sumber dan lisensi dapat diklik.
\begin{description}
\item[\texttt{3d-function-6.svg}] MartinThoma (karya asli).
\href{https://commons.wikimedia.org/wiki/File:3d-function-6.svg}{Halaman sumber di Wikimedia Commons}; lisensi \href{https://creativecommons.org/licenses/by/3.0}{CC BY 3.0}.
\item[\texttt{Great circle passing through two points.svg}] HaEr48; karya turunan dari Polar angle to spherical side.svg oleh Episcophagus.
\href{https://commons.wikimedia.org/wiki/File:Great_circle_passing_through_two_points.svg}{Halaman sumber di Wikimedia Commons}; lisensi \href{https://creativecommons.org/licenses/by-sa/4.0}{CC BY-SA 4.0}.
\item[\texttt{2019-07-Helix.jpg}] Ag2gaeh (karya asli).
\href{https://commons.wikimedia.org/wiki/File:2019-07-Helix.jpg}{Halaman sumber di Wikimedia Commons}; lisensi \href{https://creativecommons.org/licenses/by-sa/4.0}{CC BY-SA 4.0}.
\item[\texttt{Planned flight map of the Oiseau Blanc.svg}] Pethrus; karya turunan dari BlankMap-World8.svg oleh AMK1211.
\href{https://commons.wikimedia.org/wiki/File:Planned_flight_map_of_the_Oiseau_Blanc.svg}{Halaman sumber di Wikimedia Commons}; lisensi \href{https://creativecommons.org/licenses/by-sa/3.0}{CC BY-SA 3.0}.
\end{description}
Identitas revisi, ukuran, SHA-1 Wikimedia, SHA-256, URL berkas asli, dan hash turunan cetak tercatat dalam manifest media yang disertakan.

\section*{Unit 2}
\addcontentsline{toc}{section}{Unit 2}
Dua berkas berikut digunakan dalam Unit 2. Setiap berkas tetap mengikuti status hak atau lisensinya sendiri; tautan sumber dapat diklik dan tautan lisensi dicantumkan jika tersedia.
\begin{description}
\item[]\texttt{Integral} \texttt{apl} \texttt{rot} \texttt{obsah1.svg}
Pajs; dialihkan dari Wikipedia bahasa Ceko ke Wikimedia Commons.
\href{https://commons.wikimedia.org/wiki/File:Integral_apl_rot_obsah1.svg}{Halaman sumber di Wikimedia Commons}; status hak: domain publik.
\item[]\texttt{Hyperboloid1.png}
Lars H. Rohwedder (RokerHRO; karya asli).
\href{https://commons.wikimedia.org/wiki/File:Hyperboloid1.png}{Halaman sumber di Wikimedia Commons}; status hak: domain publik.
\end{description}
Identitas revisi, ukuran, SHA-1 Wikimedia, SHA-256, URL berkas asli, dan hash turunan cetak tercatat dalam manifest media yang disertakan.

\section*{Unit 3}
\addcontentsline{toc}{section}{Unit 3}
Tiga berkas berikut digunakan dalam Unit 3. Setiap berkas tetap mengikuti status hak atau lisensinya sendiri; tautan sumber dapat diklik dan tautan lisensi dicantumkan jika tersedia.
\begin{description}
\item[]\texttt{Parabola} \texttt{circle.svg}
IkamusumeFan (karya sendiri).
\href{https://commons.wikimedia.org/wiki/File:Parabola_circle.svg}{Halaman sumber di Wikimedia Commons}; lisensi \href{https://creativecommons.org/licenses/by-sa/4.0}{CC BY-SA 4.0}.
\item[]\texttt{Euler} \texttt{spiral.svg}
AdiJapan (karya sendiri).
\href{https://commons.wikimedia.org/wiki/File:Euler_spiral.svg}{Halaman sumber di Wikimedia Commons}; lisensi \href{https://creativecommons.org/licenses/by-sa/3.0}{CC BY-SA 3.0}.
\item[]\texttt{Evolute-parab.svg}
Ag2gaeh (karya sendiri).
\href{https://commons.wikimedia.org/wiki/File:Evolute-parab.svg}{Halaman sumber di Wikimedia Commons}; lisensi \href{https://creativecommons.org/licenses/by-sa/4.0}{CC BY-SA 4.0}.
\end{description}
Identitas revisi, ukuran, SHA-1 Wikimedia, SHA-256, URL berkas asli, dan hash turunan cetak tercatat dalam manifest media yang disertakan.

\section*{Unit 4}
\addcontentsline{toc}{section}{Unit 4}
Unit 4 tidak menggunakan berkas media. Penutupan nol-aset ini diverifikasi terhadap kuliah, lembar kerja, dan semua solusi yang disediakan oleh sumber.

\section*{Unit 5}
\addcontentsline{toc}{section}{Unit 5}
Satu berkas berikut digunakan dalam Unit 5. Setiap berkas tetap mengikuti status hak atau lisensinya sendiri; tautan sumber dapat diklik dan tautan lisensi dicantumkan jika tersedia.
\begin{description}
\item[]\texttt{Minimal} \texttt{surface} \texttt{curvature} \texttt{planes-de.svg}
Eric Gaba (Sting); berdasarkan sebuah gambar dalam buku.
\href{https://commons.wikimedia.org/wiki/File:Minimal_surface_curvature_planes-de.svg}{Halaman sumber di Wikimedia Commons}; lisensi \href{https://creativecommons.org/licenses/by-sa/3.0/}{CC BY-SA 3.0}.
\end{description}
Identitas revisi, ukuran, SHA-1 Wikimedia, SHA-256, URL berkas asli, dan hash turunan cetak tercatat dalam manifest media yang disertakan.

\section*{Unit 6}
\addcontentsline{toc}{section}{Unit 6}
Satu berkas berikut digunakan dalam Unit 6. Setiap berkas tetap mengikuti status hak atau lisensinya sendiri; tautan sumber dapat diklik dan tautan lisensi dicantumkan jika tersedia.
\begin{description}
\item[]\texttt{Parallel} \texttt{transport} \texttt{sphere2.svg}
Silly rabbit di Wikipedia bahasa Inggris (karya sendiri).
\href{https://commons.wikimedia.org/wiki/File:Parallel_transport_sphere2.svg}{Halaman sumber di Wikimedia Commons}; lisensi \href{https://creativecommons.org/licenses/by-sa/3.0/}{CC BY-SA 3.0}.
\end{description}
Identitas revisi, ukuran, SHA-1 Wikimedia, SHA-256, URL berkas asli, dan hash turunan cetak tercatat dalam manifest media yang disertakan.

\section*{Unit 7}
\addcontentsline{toc}{section}{Unit 7}
Tiga berkas berikut digunakan dalam Unit 7. Setiap berkas tetap mengikuti status hak atau lisensinya sendiri; tautan sumber dapat diklik dan tautan lisensi dicantumkan jika tersedia.
\begin{description}
\item[]\texttt{Stereographic} \texttt{projection} \texttt{in} \texttt{3D.png}
Mark.Howison (karya sendiri; domain publik).
\href{https://commons.wikimedia.org/wiki/File:Stereographic_projection_in_3D.png}{Halaman sumber di Wikimedia Commons}; status hak: Public domain.
\item[]\texttt{Manifold} \texttt{zahyou3.png}
132ninme di Wikipedia bahasa Jepang (karya sendiri; nama Unicode persis dicatat dalam manifes).
\href{https://commons.wikimedia.org/wiki/File:Manifold_zahyou3.png}{Halaman sumber di Wikimedia Commons}; lisensi \href{https://creativecommons.org/licenses/by-sa/3.0/}{CC BY-SA 3.0}.
\item[]\texttt{Circle} \texttt{-} \texttt{black} \texttt{simple.svg}
Dakdada (karya sendiri; domain publik).
\href{https://commons.wikimedia.org/wiki/File:Circle_-_black_simple.svg}{Halaman sumber di Wikimedia Commons}; status hak: Public domain.
\end{description}
Identitas revisi, ukuran, SHA-1 Wikimedia, SHA-256, URL berkas asli, dan hash turunan cetak tercatat dalam manifest media yang disertakan.

\section*{Unit 8}
\addcontentsline{toc}{section}{Unit 8}
Unit 8 tidak menggunakan berkas media. Penutupan nol-aset ini diverifikasi terhadap kuliah, lembar kerja, dan semua solusi yang disediakan oleh sumber.

\section*{Unit 9}
\addcontentsline{toc}{section}{Unit 9}
Satu berkas berikut digunakan dalam Unit 9. Setiap berkas tetap mengikuti status hak atau lisensinya sendiri; tautan sumber dapat diklik dan tautan lisensi dicantumkan jika tersedia.
\begin{description}
\item[]\texttt{Tangentialvektor.svg}
McSush (karya turunan dari Tangentialvektor.png oleh TN; domain publik).
\href{https://commons.wikimedia.org/wiki/File:Tangentialvektor.svg}{Halaman sumber di Wikimedia Commons}; status hak: domain publik.
\end{description}
Identitas revisi, ukuran, SHA-1 Wikimedia, SHA-256, URL berkas asli, dan hash turunan cetak tercatat dalam manifest media yang disertakan.

\section*{Unit 10}
\addcontentsline{toc}{section}{Unit 10}
Dua berkas berikut digunakan dalam Unit 10. Setiap berkas tetap mengikuti status hak atau lisensinya sendiri; tautan sumber dapat diklik dan tautan lisensi dicantumkan jika tersedia.
\begin{description}
\item[]\texttt{Tangent} \texttt{bundle.svg}
Oleg Alexandrov (karya sendiri dengan Matlab dan Inkscape; domain publik).
\href{https://commons.wikimedia.org/wiki/File:Tangent_bundle.svg}{Halaman sumber di Wikimedia Commons}; status hak: domain publik.
\item[]\texttt{Torus} \texttt{vectors} \texttt{oblique.jpg}
RokerHRO (karya sendiri).
\href{https://commons.wikimedia.org/wiki/File:Torus_vectors_oblique.jpg}{Halaman sumber di Wikimedia Commons}; lisensi \href{https://creativecommons.org/licenses/by-sa/3.0/}{CC BY-SA 3.0}.
\end{description}
Identitas revisi, ukuran, SHA-1 Wikimedia, SHA-256, URL berkas asli, dan hash turunan cetak tercatat dalam manifest media yang disertakan.

\section*{Unit 11}
\addcontentsline{toc}{section}{Unit 11}
Satu berkas berikut digunakan dalam Unit 11. Setiap berkas tetap mengikuti status hak atau lisensinya sendiri; tautan sumber dapat diklik dan tautan lisensi dicantumkan jika tersedia.
\begin{description}
\item[]\texttt{Toroidal} \texttt{coord.png}
DaveBurke (karya sendiri).
\href{https://commons.wikimedia.org/wiki/File:Toroidal_coord.png}{Halaman sumber di Wikimedia Commons}; lisensi \href{https://creativecommons.org/licenses/by/2.5/}{CC BY 2.5}.
\end{description}
Identitas revisi, ukuran, SHA-1 Wikimedia, SHA-256, URL berkas asli, dan hash turunan cetak tercatat dalam manifest media yang disertakan.

\section*{Unit 12}
\addcontentsline{toc}{section}{Unit 12}
Dua berkas berikut digunakan dalam Unit 12. Setiap berkas tetap mengikuti status hak atau lisensinya sendiri; tautan sumber dapat diklik dan tautan lisensi dicantumkan jika tersedia.
\begin{description}
\item[]\texttt{Fiddler} \texttt{crab} \texttt{mobius} \texttt{strip.gif}
Hamishtodd1 (karya sendiri).
\href{https://commons.wikimedia.org/wiki/File:Fiddler_crab_mobius_strip.gif}{Halaman sumber di Wikimedia Commons}; lisensi \href{https://creativecommons.org/licenses/by-sa/4.0/}{CC BY-SA 4.0}.
\item[]\texttt{Inclusion-exclusion.svg}
Pembuat tidak diketahui; domain publik.
\href{https://commons.wikimedia.org/wiki/File:Inclusion-exclusion.svg}{Halaman sumber di Wikimedia Commons}; status hak: domain publik.
\end{description}
Identitas revisi, ukuran, SHA-1 Wikimedia, SHA-256, URL berkas asli, dan hash turunan cetak tercatat dalam manifest media yang disertakan.

\section*{Unit 13}
\addcontentsline{toc}{section}{Unit 13}
Satu berkas berikut digunakan dalam Unit 13. Setiap berkas tetap mengikuti status hak atau lisensinya sendiri; tautan sumber dapat diklik dan tautan lisensi dicantumkan jika tersedia.
\begin{description}
\item[]\texttt{Möbius} \texttt{strip.jpg}
David Benbennick (karya sendiri).
\href{https://commons.wikimedia.org/wiki/File:M%C3%B6bius_strip.jpg}{Halaman sumber di Wikimedia Commons}; lisensi \href{https://creativecommons.org/licenses/by-sa/3.0/}{CC BY-SA 3.0}.
\end{description}
Identitas revisi, ukuran, SHA-1 Wikimedia, SHA-256, URL berkas asli, dan hash turunan cetak tercatat dalam manifest media yang disertakan.

\section*{Unit 14}
\addcontentsline{toc}{section}{Unit 14}
Unit 14 tidak menggunakan berkas media. Penutupan nol-aset ini diverifikasi terhadap kuliah, lembar kerja, dan semua solusi yang disediakan oleh sumber.

\section*{Unit 15}
\addcontentsline{toc}{section}{Unit 15}
Unit 15 tidak menggunakan berkas media. Penutupan nol-aset ini diverifikasi terhadap kuliah, lembar kerja, dan semua solusi yang disediakan oleh sumber.

\section*{Unit 16}
\addcontentsline{toc}{section}{Unit 16}
Dua berkas berikut digunakan dalam Unit 16. Setiap berkas tetap mengikuti status hak atau lisensinya sendiri; tautan sumber dapat diklik dan tautan lisensi dicantumkan jika tersedia.
\begin{description}
\item[]\texttt{Georg} \texttt{Friedrich} \texttt{Bernhard} \texttt{Riemann.jpeg}
Potret Georg Friedrich Bernhard Riemann; pembuat tidak dicantumkan dalam metadata Commons.
\href{https://commons.wikimedia.org/wiki/File:Georg_Friedrich_Bernhard_Riemann.jpeg}{Halaman sumber di Wikimedia Commons}; status hak: domain publik.
\item[]\texttt{Sphere} \texttt{with} \texttt{three} \texttt{handles.png}
Oleg Alexandrov (karya sendiri; domain publik).
\href{https://commons.wikimedia.org/wiki/File:Sphere_with_three_handles.png}{Halaman sumber di Wikimedia Commons}; status hak: domain publik.
\end{description}
Identitas revisi, ukuran, SHA-1 Wikimedia, SHA-256, URL berkas asli, dan hash turunan cetak tercatat dalam manifest media yang disertakan.

\section*{Unit 17}
\addcontentsline{toc}{section}{Unit 17}
Satu berkas berikut digunakan dalam Unit 17. Setiap berkas tetap mengikuti status hak atau lisensinya sendiri; tautan sumber dapat diklik dan tautan lisensi dicantumkan jika tersedia.
\begin{description}
\item[]\texttt{Cilinderprojectie-constructie.jpg}
KoenB (karya sendiri; domain publik).
\href{https://commons.wikimedia.org/wiki/File:Cilinderprojectie-constructie.jpg}{Halaman sumber di Wikimedia Commons}; status hak: domain publik.
\end{description}
Identitas revisi, ukuran, SHA-1 Wikimedia, SHA-256, URL berkas asli, dan hash turunan cetak tercatat dalam manifest media yang disertakan.

\section*{Unit 18}
\addcontentsline{toc}{section}{Unit 18}
Dua berkas berikut digunakan dalam Unit 18. Setiap berkas tetap mengikuti status hak atau lisensinya sendiri; tautan sumber dapat diklik dan tautan lisensi dicantumkan jika tersedia.
\begin{description}
\item[]\texttt{Poincarehalfplaneconform.gif}
HB (karya sendiri).
\href{https://commons.wikimedia.org/wiki/File:Poincarehalfplaneconform.gif}{Halaman sumber di Wikimedia Commons}; lisensi \href{https://creativecommons.org/licenses/by-sa/3.0/}{CC BY-SA 3.0}.
\item[]\texttt{Hyperboloid2.png}
Lars H. Rohwedder (RokerHRO; karya asli).
\href{https://commons.wikimedia.org/wiki/File:Hyperboloid2.png}{Halaman sumber di Wikimedia Commons}; status hak: domain publik.
\end{description}
Identitas revisi, ukuran, SHA-1 Wikimedia, SHA-256, URL berkas asli, dan hash turunan cetak tercatat dalam manifest media yang disertakan.

\section*{Unit 19}
\addcontentsline{toc}{section}{Unit 19}
Unit 19 tidak menggunakan berkas media. Penutupan nol-aset ini diverifikasi terhadap kuliah, lembar kerja, dan semua solusi yang disediakan oleh sumber.

\section*{Unit 20}
\addcontentsline{toc}{section}{Unit 20}
Unit 20 tidak menggunakan berkas media. Penutupan nol-aset ini diverifikasi terhadap kuliah, lembar kerja, dan semua solusi yang disediakan oleh sumber.

\section*{Unit 21}
\addcontentsline{toc}{section}{Unit 21}
Tiga berkas berikut digunakan dalam Unit 21. Setiap berkas tetap mengikuti status hak atau lisensinya sendiri; tautan sumber dapat diklik dan tautan lisensi dicantumkan jika tersedia.
\begin{description}
\item[]\texttt{Runge} \texttt{theorem.svg}
Oleg Alexandrov (domain publik).
\href{https://commons.wikimedia.org/wiki/File:Runge_theorem.svg}{Halaman sumber di Wikimedia Commons}; status hak: domain publik.
\item[]\texttt{Circle} \texttt{on} \texttt{sphere} \texttt{wireframe} \texttt{10deg} \texttt{6r.svg}
Itai (karya sendiri).
\href{https://commons.wikimedia.org/wiki/File:Circle_on_sphere_wireframe_10deg_6r.svg}{Halaman sumber di Wikimedia Commons}; lisensi \href{https://creativecommons.org/licenses/by/3.0/}{CC BY 3.0}.
\item[]\texttt{Not-star-shaped.svg}
WillT.Net; dialihkan dari Wikipedia bahasa Inggris (domain publik).
\href{https://commons.wikimedia.org/wiki/File:Not-star-shaped.svg}{Halaman sumber di Wikimedia Commons}; status hak: domain publik.
\end{description}
Identitas revisi, ukuran, SHA-1 Wikimedia, SHA-256, URL berkas asli, dan hash turunan cetak tercatat dalam manifest media yang disertakan.

\section*{Unit 22}
\addcontentsline{toc}{section}{Unit 22}
Dua berkas berikut digunakan dalam Unit 22. Setiap berkas tetap mengikuti status hak atau lisensinya sendiri; tautan sumber dapat diklik dan tautan lisensi dicantumkan jika tersedia.
\begin{description}
\item[]\texttt{Inner} \texttt{point.png}
Zasdfgbnm (karya sendiri; domain publik).
\href{https://commons.wikimedia.org/wiki/File:Inner_point.png}{Halaman sumber di Wikimedia Commons}; status hak: domain publik.
\item[]\texttt{Partition} \texttt{of} \texttt{unity} \texttt{illustration.svg}
Oleg Alexandrov (dibuat sendiri dengan MATLAB dan disempurnakan di Inkscape; domain publik).
\href{https://commons.wikimedia.org/wiki/File:Partition_of_unity_illustration.svg}{Halaman sumber di Wikimedia Commons}; status hak: domain publik.
\end{description}
Identitas revisi, ukuran, SHA-1 Wikimedia, SHA-256, URL berkas asli, dan hash turunan cetak tercatat dalam manifest media yang disertakan.

\section*{Unit 23}
\addcontentsline{toc}{section}{Unit 23}
Tiga berkas berikut digunakan dalam Unit 23. Setiap berkas tetap mengikuti status hak atau lisensinya sendiri; tautan sumber dapat diklik dan tautan lisensi dicantumkan jika tersedia.
\begin{description}
\item[]\texttt{SS-stokes.jpg}
Kelson (pengunggah di Wikimedia Commons; domain publik).
\href{https://commons.wikimedia.org/wiki/File:SS-stokes.jpg}{Halaman sumber di Wikimedia Commons}; status hak: domain publik.
\item[]\texttt{Théorème-de-Brouwer-(cond-2).jpg}
Jean-Luc W (karya sendiri).
\href{https://commons.wikimedia.org/wiki/File:Th%C3%A9or%C3%A8me-de-Brouwer-(cond-2).jpg}{Halaman sumber di Wikimedia Commons}; lisensi \href{https://creativecommons.org/licenses/by-sa/3.0/}{CC BY-SA 3.0}.
\item[]\texttt{Théorème-de-Brouwer-(cond-1).jpg}
Jean-Luc W (karya sendiri).
\href{https://commons.wikimedia.org/wiki/File:Th%C3%A9or%C3%A8me-de-Brouwer-(cond-1).jpg}{Halaman sumber di Wikimedia Commons}; lisensi \href{https://creativecommons.org/licenses/by-sa/3.0/}{CC BY-SA 3.0}.
\end{description}
Identitas revisi, ukuran, SHA-1 Wikimedia, SHA-256, URL berkas asli, dan hash turunan cetak tercatat dalam manifest media yang disertakan.

\section*{Unit 24}
\addcontentsline{toc}{section}{Unit 24}
Unit 24 tidak menggunakan berkas media. Penutupan nol-aset ini diverifikasi terhadap kuliah, lembar kerja, dan semua solusi yang disediakan oleh sumber.

\section*{Unit 25}
\addcontentsline{toc}{section}{Unit 25}
Unit 25 tidak menggunakan berkas media. Penutupan nol-aset ini diverifikasi terhadap kuliah, lembar kerja, dan semua solusi yang disediakan oleh sumber.

\section*{Unit 26}
\addcontentsline{toc}{section}{Unit 26}
Unit 26 tidak menggunakan berkas media. Penutupan nol-aset ini diverifikasi terhadap kuliah, lembar kerja, dan semua solusi yang disediakan oleh sumber.

\input{generated/unit27-media-attribution-complete.tex}
\section*{Unit 28}
\addcontentsline{toc}{section}{Unit 28}
Unit 28 tidak menggunakan berkas media. Penutupan nol-aset ini diverifikasi terhadap kuliah, lembar kerja, dan semua solusi yang disediakan oleh sumber.

\section*{Unit 29}
\addcontentsline{toc}{section}{Unit 29}
Satu berkas berikut digunakan dalam Unit 29. Setiap berkas tetap mengikuti status hak atau lisensinya sendiri; tautan sumber dapat diklik dan tautan lisensi dicantumkan jika tersedia.
\begin{description}
\item[]\texttt{Parallel} \texttt{lines} \texttt{in} \texttt{Poincare's} \texttt{model} \texttt{of} \texttt{hyperbolic} \texttt{geometry.svg}
Januszkaja (karya sendiri).
\href{https://commons.wikimedia.org/wiki/File:Parallel_lines_in_Poincare%27s_model_of_hyperbolic_geometry.svg}{Halaman sumber di Wikimedia Commons}; lisensi \href{https://creativecommons.org/licenses/by-sa/3.0/}{CC BY-SA 3.0}.
\end{description}
Identitas revisi, ukuran, SHA-1 Wikimedia, SHA-256, URL berkas asli, dan hash turunan cetak tercatat dalam manifest media yang disertakan.

\chapter*{Lisensi}
\Lizenztext
\end{document}
